
\documentclass[11pt,a4paper,openany]{ctexbook}
\usepackage[a4paper,margin=30mm,headheight=15pt,footskip=12mm]{geometry}
\usepackage{fontspec}
\setmainfont{Times New Roman}
\setsansfont{Noto Sans}
\setCJKmainfont{SimSun}
\setCJKsansfont{Noto Sans SC}
\setCJKmonofont{Noto Sans SC}
\usepackage{setspace,microtype}
\setstretch{1.16}
\setlength{\parindent}{2em}
\setlength{\parskip}{0.18em}
\usepackage{amsmath,amssymb,mathtools,bm,mathrsfs}
\allowdisplaybreaks[2]
\usepackage{enumitem,array,booktabs,xcolor,titlesec}
\usepackage[most]{tcolorbox}
\usepackage{hyperref}
\usepackage{fancyhdr}
\definecolor{InkBlue}{RGB}{33,72,112}
\definecolor{SoftGreen}{RGB}{64,109,84}
\definecolor{SoftRed}{RGB}{143,76,68}
\definecolor{PaperGray}{RGB}{248,248,246}
\definecolor{RuleGray}{RGB}{126,126,126}
\hypersetup{colorlinks=true,linkcolor=InkBlue,urlcolor=SoftRed,citecolor=SoftGreen}
\pagestyle{fancy}
\fancyhf{}
\fancyhead[LE,RO]{\thepage}
\fancyhead[LO]{\small 组合学研讨班讲义}
\fancyhead[RE]{\small 详细解读版}
\renewcommand{\headrulewidth}{0.35pt}
\titleformat{\chapter}[display]{\normalfont\huge\bfseries}{第\thechapter 章}{0.6em}{}
\titlespacing*{\chapter}{0pt}{-0.4em}{1.4em}
\titleformat{\section}{\normalfont\Large\bfseries}{\thesection}{0.8em}{}
\titlespacing*{\section}{0pt}{1.2em}{0.55em}
\tcbset{common/.style={enhanced,breakable,sharp corners,boxrule=0.32pt,left=2.2mm,right=1.8mm,top=1mm,bottom=1mm,before skip=0.65em,after skip=0.75em}}
\newtcolorbox{roadmapblock}{common,colback=white,colframe=InkBlue!55!black,borderline west={1.2pt}{0pt}{InkBlue!70!black},fonttitle=\bfseries,title={阅读地图}}
\newtcolorbox{sourceblock}[1]{common,colback=PaperGray,colframe=RuleGray!55,borderline west={1pt}{0pt}{RuleGray},fonttitle=\bfseries,title={原书数学骨架：#1}}
\newtcolorbox{explainblock}{common,colback=white,colframe=SoftGreen!45!black,borderline west={1.2pt}{0pt}{SoftGreen!70!black},fonttitle=\bfseries,title={给初学者的拆解}}

\title{组合学研讨班讲义\\\large 基于《欧拉与计数组合学》的详细解读版}
\author{面向初学者的逐段讲解}
\date{}
\begin{document}
\maketitle
\frontmatter
\chapter*{怎么读这份讲义}
这份讲义不是把原 PDF 页面贴进来，也不保留封面、目录、插图和原书版面。它保留原书正文里的定义、定理、例题、习题和关键论述作为“数学骨架”，然后在每个骨架后面加一层面向初学者的拆解。

推荐读法是：先读灰色的“原书数学骨架”，知道原书在说什么；再读绿色的“给初学者的拆解”，确认这一段为什么这样数、容易错在哪里。读组合学最怕跳步，这份讲义的目标就是把跳过去的判断补出来。
\tableofcontents
\mainmatter

\chapter{排列与组合}
\begin{roadmapblock}
本章保留原书正文的主要数学骨架，并在每个定义、定理、例题和论述片段后加入解读。
阅读时不要把目标放在“背下答案”上，而是放在“看懂这一步为什么这样数”上。

\smallskip
\textbf{本章最低目标：} 能说出每个公式左边在数什么、右边在数什么，以及为什么两边相等。
\end{roadmapblock}


\begin{sourceblock}{第 1 章正文}
{\small\sloppy
\noindent 排列与组合\par
\noindent 排列与组合是组合学中最基本的两个组合概念. 本章主要介绍了与排列和组合有\par
\noindent 关的几个计数公式, 并利用加法原理和乘法原理及归纳法和构造法给出相关证明.\par
}
\end{sourceblock}


\begin{explainblock}
这一段是原书论述链的一部分。读的时候把它当作桥梁：它通常负责把上一个定义、例子或定理，引到下一个计数方法。

\smallskip
\textbf{初学者检查点：}
\begin{itemize}[leftmargin=2em,itemsep=0.25em]
\item 找出这一段承接的上一个概念。
\item 找出它准备引出的下一个概念。
\item 把中间的关键等式或构造单独抄出来。
\end{itemize}

\smallskip
\textbf{本章读法提醒：} 第一章所有题都先判断“顺序是否重要”。这个判断错了，后面公式一定错。
\end{explainblock}


\begin{sourceblock}{§ 1.1 排列}
{\small\sloppy
\noindent § 1.1. 排列\par
}
\end{sourceblock}


\begin{explainblock}
这一节先不要急着背公式。先问三个问题：对象是什么，哪些限制条件不能丢，最后到底是在数所有对象、还是在数某个统计量的分布。把这三件事说清楚，后面的公式才会有位置。

\smallskip
\textbf{初学者检查点：}
\begin{itemize}[leftmargin=2em,itemsep=0.25em]
\item 把本节出现的新对象写成一句话定义。
\item 把每个公式左边和右边分别翻译成“在数什么”。
\item 找一个最小非平凡例子，手工列出所有对象。
\end{itemize}

\smallskip
\textbf{本章读法提醒：} 第一章所有题都先判断“顺序是否重要”。这个判断错了，后面公式一定错。
\end{explainblock}


\begin{sourceblock}{定义 1.1.1 (排列) n 元集合 A 的 k-排列是从 A 中选出 k 个元素排成一列 (考虑}
{\small\sloppy
\noindent 定义 1.1.1 (排列) n 元集合 A 的 k-排列是从 A 中选出 k 个元素排成一列 (考虑\par
\noindent 元素先后出现次序)，这样排列的总数记为 Pnk . 特别地，当 k = n 时，这个排列被\par
\noindent 称作集合 A 的全排列 (permutation), 该排列总数记为 Pn .\par
}
\end{sourceblock}


\begin{explainblock}
定义段的任务是定边界。这里要特别注意参数范围、是否允许重复、是否考虑顺序、对象有没有标签。后面的每一道例题和证明，本质上都在反复使用这些边界。

\smallskip
\textbf{初学者检查点：}
\begin{itemize}[leftmargin=2em,itemsep=0.25em]
\item 圈出被定义的对象名称。
\item 圈出参数范围，例如 n、k 是否要求正整数。
\item 判断它是有序对象、无序对象、带标签对象还是结构对象。
\end{itemize}

\smallskip
\textbf{本章读法提醒：} 第一章所有题都先判断“顺序是否重要”。这个判断错了，后面公式一定错。
\end{explainblock}


\begin{sourceblock}{例 1.1.1 (1) 集合 \{1, 2, 3\} 的全排列有:}
{\small\sloppy
\noindent 例 1.1.1 (1) 集合 \{1, 2, 3\} 的全排列有:\par
\noindent 123, 132, 213, 231, 312, 321,\par
\noindent 所以 P3 = 6;\par
\noindent (2) 集合 \{1, 2, 3, 4\} 的 2-排列有:\par
\noindent 12, 21, 13, 31, 14, 41, 23, 32, 24, 42, 34, 43,\par
\noindent 故 P42 = 12;\par
\noindent 记 [n] = \{1, 2, . . . , n\}, 那么 [n] 上的一个排列 x1 x2 · · · xn 可以看成是 [n] 到自身\par
\noindent 的一个双射 π, 其中 π(i) = xi .\par
}
\end{sourceblock}


\begin{explainblock}
例题是学习这本书最重要的部分。不要只看答案，要把每个限制条件变成一个计数动作：先安排谁、再插入谁、有没有重复、需不需要除以对称性。

\smallskip
\textbf{初学者检查点：}
\begin{itemize}[leftmargin=2em,itemsep=0.25em]
\item 先写出没有限制时的总数。
\item 逐条加入限制条件，观察总数怎样变化。
\item 最后检查有没有把同一种方案重复算了多次。
\end{itemize}

\smallskip
\textbf{本章读法提醒：} 第一章所有题都先判断“顺序是否重要”。这个判断错了，后面公式一定错。
\end{explainblock}


\begin{sourceblock}{定理 1.1.2 设 1 ≤ k ≤ n, 则}
{\small\sloppy
\noindent 定理 1.1.2 设 1 ≤ k ≤ n, 则\par
\noindent n!\par
\noindent Pnk = n(n − 1)(n − 2) · · · (n − k + 1) = . (1.1)\par
\noindent (n − k)!\par
\noindent 证明: 从 n 元集合中选取 k 个元素, 若不允许重复, 则排列的第 1 项有 n 种取\par
\noindent 法, 第 2 项有 n − 1 种取法. 依此下去, 只要选定前 k − 1 项, 则排列的第 k 项就有\par
\noindent n − k + 1 种取法, 应用乘法原理得\par
\noindent n!\par
\noindent Pnk = n(n − 1) · · · (n − k + 1) = = (n)k .\par
\noindent (n − k)!\par
\noindent 特别地, 当 k = n 时,\par
\noindent Pnn = Pn = n!.\par
}
\end{sourceblock}


\begin{explainblock}
定理段不要只看结论，要看它把哪两个计数方式连在一起。组合学证明最常见的技巧是：左边按一种标准数，右边按另一种标准数，然后说明两边数的是同一批对象。

\smallskip
\textbf{初学者检查点：}
\begin{itemize}[leftmargin=2em,itemsep=0.25em]
\item 先复述定理的假设，不要把适用范围扩大。
\item 找出证明中的基本计数原则：乘法、加法、双射、递推、容斥或生成函数。
\item 检查最后的等式化简是否和前面的对象解释一致。
\end{itemize}

\smallskip
\textbf{本章读法提醒：} 第一章所有题都先判断“顺序是否重要”。这个判断错了，后面公式一定错。
\end{explainblock}


\begin{sourceblock}{注记 1.1.3 这里记}
{\small\sloppy
\noindent 注记 1.1.3 这里记\par
\noindent (n)k = n(n − 1) · · · (n − k + 1),\par
\noindent 读作 “ n 的 k 次下阶乘 (n lower factorial k) .”\par
\noindent n! = n(n − 1) · · · 2 · 1,\par
\noindent 读作 “ n 的阶乘 ( n factorial) ”, 并规定 0! = 1;\par
}
\end{sourceblock}


\begin{explainblock}
注记通常是在补充术语、记号、历史或一个容易忽略的约定。它未必给新公式，但常常决定你后面会不会把符号读错。

\smallskip
\textbf{初学者检查点：}
\begin{itemize}[leftmargin=2em,itemsep=0.25em]
\item 记下新符号的读法。
\item 确认这个约定是否只在本章使用。
\item 把注记里的限制条件补到前面的定义旁边。
\end{itemize}

\smallskip
\textbf{本章读法提醒：} 第一章所有题都先判断“顺序是否重要”。这个判断错了，后面公式一定错。
\end{explainblock}


\begin{sourceblock}{例 1.1.2 12 个人从左至右排成一排, 其中小 A 不能在队首, 也不能排在队尾, 问}
{\small\sloppy
\noindent 例 1.1.2 12 个人从左至右排成一排, 其中小 A 不能在队首, 也不能排在队尾, 问\par
\noindent 有多少种排法?\par
\noindent 解: 12 个人排成一排共有 12! 种排法, 然后减去小 A 排在队首和队尾的排法共\par
\noindent 有 2 × 11!, 故满足条件的排法有:\par
\noindent 12! − 2 × 11! = 10 × 11!.\par
}
\end{sourceblock}


\begin{explainblock}
例题是学习这本书最重要的部分。不要只看答案，要把每个限制条件变成一个计数动作：先安排谁、再插入谁、有没有重复、需不需要除以对称性。

\smallskip
\textbf{初学者检查点：}
\begin{itemize}[leftmargin=2em,itemsep=0.25em]
\item 先写出没有限制时的总数。
\item 逐条加入限制条件，观察总数怎样变化。
\item 最后检查有没有把同一种方案重复算了多次。
\end{itemize}

\smallskip
\textbf{本章读法提醒：} 第一章所有题都先判断“顺序是否重要”。这个判断错了，后面公式一定错。
\end{explainblock}


\begin{sourceblock}{定义 1.1.4 (圆排列) 从 n 元集合中选取 k 个元素按照某种顺序 (顺时针或逆时}
{\small\sloppy
\noindent 定义 1.1.4 (圆排列) 从 n 元集合中选取 k 个元素按照某种顺序 (顺时针或逆时\par
\noindent 针) 排成一个圆圈, 称这样的排列为圆排列.\par
}
\end{sourceblock}


\begin{explainblock}
定义段的任务是定边界。这里要特别注意参数范围、是否允许重复、是否考虑顺序、对象有没有标签。后面的每一道例题和证明，本质上都在反复使用这些边界。

\smallskip
\textbf{初学者检查点：}
\begin{itemize}[leftmargin=2em,itemsep=0.25em]
\item 圈出被定义的对象名称。
\item 圈出参数范围，例如 n、k 是否要求正整数。
\item 判断它是有序对象、无序对象、带标签对象还是结构对象。
\end{itemize}

\smallskip
\textbf{本章读法提醒：} 第一章所有题都先判断“顺序是否重要”。这个判断错了，后面公式一定错。
\end{explainblock}


\begin{sourceblock}{定理 1.1.5 n 元集合 A 中的 k 个元素的圆排列的个数为}
{\small\sloppy
\noindent 定理 1.1.5 n 元集合 A 中的 k 个元素的圆排列的个数为\par
\noindent Pnk (n)k\par
\noindent = . (1.2)\par
\noindent k k\par
\noindent 证明: 不妨设 A = a1 , a2 , . . . , an , 由于把一个圆排列旋转所得到的另一个\par
\noindent 圆排列视为相同的圆排列, 因此这些 k-排列: aj1 aj2 · · · ajk , aj2 aj3 · · · ajk aj1 , . . .,\par
\noindent ajk aj1 · · · ajk−1 在圆排列中是同一个, 即一个圆排列可以产生 k 个 k-排列, 由定理\par
\noindent 1.1.2 知 k-排列个数有 (n)k 个, 故圆排列的个数为 (n)\par
\noindent k .\par
\noindent k\par
}
\end{sourceblock}


\begin{explainblock}
定理段不要只看结论，要看它把哪两个计数方式连在一起。组合学证明最常见的技巧是：左边按一种标准数，右边按另一种标准数，然后说明两边数的是同一批对象。

\smallskip
\textbf{初学者检查点：}
\begin{itemize}[leftmargin=2em,itemsep=0.25em]
\item 先复述定理的假设，不要把适用范围扩大。
\item 找出证明中的基本计数原则：乘法、加法、双射、递推、容斥或生成函数。
\item 检查最后的等式化简是否和前面的对象解释一致。
\end{itemize}

\smallskip
\textbf{本章读法提醒：} 第一章所有题都先判断“顺序是否重要”。这个判断错了，后面公式一定错。
\end{explainblock}


\begin{sourceblock}{例 1.1.3 10 个男生和 5 个女生聚餐, 围坐在圆桌旁, 任意两个女生不相邻的坐法}
{\small\sloppy
\noindent 例 1.1.3 10 个男生和 5 个女生聚餐, 围坐在圆桌旁, 任意两个女生不相邻的坐法\par
\noindent 有多少种?\par
\noindent 解: 先把 10 个男生排成圆形, 有 9! 种方法, 固定一个男生, 把 5 个女生插在 10\par
\noindent 个男生之间, 每两个男生之间, 只有一个女生, 而且 5 个女生之间存在的排序问题,\par
\noindent 5 种排法, 由乘法原理得有 9! × (10) 种坐法.\par
\noindent 故有 P10 5\par
}
\end{sourceblock}


\begin{explainblock}
例题是学习这本书最重要的部分。不要只看答案，要把每个限制条件变成一个计数动作：先安排谁、再插入谁、有没有重复、需不需要除以对称性。

\smallskip
\textbf{初学者检查点：}
\begin{itemize}[leftmargin=2em,itemsep=0.25em]
\item 先写出没有限制时的总数。
\item 逐条加入限制条件，观察总数怎样变化。
\item 最后检查有没有把同一种方案重复算了多次。
\end{itemize}

\smallskip
\textbf{本章读法提醒：} 第一章所有题都先判断“顺序是否重要”。这个判断错了，后面公式一定错。
\end{explainblock}


\begin{sourceblock}{定义 1.1.6 (重排列) 从 n 元集合 A, 可重复地选取 k 个元素按照一定的顺序排}
{\small\sloppy
\noindent 定义 1.1.6 (重排列) 从 n 元集合 A, 可重复地选取 k 个元素按照一定的顺序排\par
\noindent 列, 称这种 k-排列为可重排列.\par
}
\end{sourceblock}


\begin{explainblock}
定义段的任务是定边界。这里要特别注意参数范围、是否允许重复、是否考虑顺序、对象有没有标签。后面的每一道例题和证明，本质上都在反复使用这些边界。

\smallskip
\textbf{初学者检查点：}
\begin{itemize}[leftmargin=2em,itemsep=0.25em]
\item 圈出被定义的对象名称。
\item 圈出参数范围，例如 n、k 是否要求正整数。
\item 判断它是有序对象、无序对象、带标签对象还是结构对象。
\end{itemize}

\smallskip
\textbf{本章读法提醒：} 第一章所有题都先判断“顺序是否重要”。这个判断错了，后面公式一定错。
\end{explainblock}


\begin{sourceblock}{例 1.1.4 集合 \{1, 2, 3\} 中的 2-可重排列有 9 个:}
{\small\sloppy
\noindent 例 1.1.4 集合 \{1, 2, 3\} 中的 2-可重排列有 9 个:\par
\noindent 11, 12, 13, 21, 22, 23, 31, 32, 33.\par
}
\end{sourceblock}


\begin{explainblock}
例题是学习这本书最重要的部分。不要只看答案，要把每个限制条件变成一个计数动作：先安排谁、再插入谁、有没有重复、需不需要除以对称性。

\smallskip
\textbf{初学者检查点：}
\begin{itemize}[leftmargin=2em,itemsep=0.25em]
\item 先写出没有限制时的总数。
\item 逐条加入限制条件，观察总数怎样变化。
\item 最后检查有没有把同一种方案重复算了多次。
\end{itemize}

\smallskip
\textbf{本章读法提醒：} 第一章所有题都先判断“顺序是否重要”。这个判断错了，后面公式一定错。
\end{explainblock}


\begin{sourceblock}{定理 1.1.7 n 元集合的 k-可重排列个数为 nk .}
{\small\sloppy
\noindent 定理 1.1.7 n 元集合的 k-可重排列个数为 nk .\par
\noindent 证明: 因为排列的第 1 个位置有 n 种取法, 由于可重复选取, 所以第 2 个位置\par
\noindent 也有 n 种取法. 同理, 该排列的任意一个位置均有 n 种取法, 由乘法原理可得 n 元\par
\noindent 集合的 k-可重排列的个数有 nk .\par
}
\end{sourceblock}


\begin{explainblock}
定理段不要只看结论，要看它把哪两个计数方式连在一起。组合学证明最常见的技巧是：左边按一种标准数，右边按另一种标准数，然后说明两边数的是同一批对象。

\smallskip
\textbf{初学者检查点：}
\begin{itemize}[leftmargin=2em,itemsep=0.25em]
\item 先复述定理的假设，不要把适用范围扩大。
\item 找出证明中的基本计数原则：乘法、加法、双射、递推、容斥或生成函数。
\item 检查最后的等式化简是否和前面的对象解释一致。
\end{itemize}

\smallskip
\textbf{本章读法提醒：} 第一章所有题都先判断“顺序是否重要”。这个判断错了，后面公式一定错。
\end{explainblock}


\begin{sourceblock}{定理 1.1.8 设集合 A = \{a1 , a2 , . . . , ak \}, 则使得 A 中每一个元素 ai 恰出现 ni 次}
{\small\sloppy
\noindent 定理 1.1.8 设集合 A = \{a1 , a2 , . . . , ak \}, 则使得 A 中每一个元素 ai 恰出现 ni 次\par
\noindent 的排列的个数为\par
\noindent n!\par
\noindent ,\par
\noindent n1 !n2 ! · · · nk !\par
\noindent Pk\par
\noindent 其中 n = i=1 ni .\par
\noindent 证明: 设这样的排列方式有 pn 种. 因为这样的排列实际上是由 ni (1 ≤ i ≤ k)\par
\noindent 个 ai 组成的 n 长排列, 不妨给这 ni 个 ai 分别赋予上标 1, 2, . . . , ni , 这样就变成了\par
\noindent 对 n1 + n2 + · · · + nk = n 个元素的全排列, 共计 n! 个. 另一方面, 因为对 ai 赋于上\par
\noindent 标的方式恰是 ni 个元素全排列, 有 ni ! 种赋于方式, 由乘法原理知\par
\noindent n! = pn · n1 ! · n2 ! · · · nk !,\par
\noindent 所以\par
\noindent n!\par
\noindent pn = .\par
\noindent n1 · n2 ! · · · nk !\par
}
\end{sourceblock}


\begin{explainblock}
定理段不要只看结论，要看它把哪两个计数方式连在一起。组合学证明最常见的技巧是：左边按一种标准数，右边按另一种标准数，然后说明两边数的是同一批对象。

\smallskip
\textbf{初学者检查点：}
\begin{itemize}[leftmargin=2em,itemsep=0.25em]
\item 先复述定理的假设，不要把适用范围扩大。
\item 找出证明中的基本计数原则：乘法、加法、双射、递推、容斥或生成函数。
\item 检查最后的等式化简是否和前面的对象解释一致。
\end{itemize}

\smallskip
\textbf{本章读法提醒：} 第一章所有题都先判断“顺序是否重要”。这个判断错了，后面公式一定错。
\end{explainblock}


\begin{sourceblock}{例 1.1.5 由 2 面红旗, 3 面蓝旗, 4 面黄旗, 5 面绿旗可以组成多少种由 14 面旗}
{\small\sloppy
\noindent 例 1.1.5 由 2 面红旗, 3 面蓝旗, 4 面黄旗, 5 面绿旗可以组成多少种由 14 面旗\par
\noindent 帜排成的一列彩旗?\par
\noindent 解: 由定理 1.1.8 知彩旗列的排列方式有:\par
\noindent (2 + 3 + 4 + 5)!\par
\noindent .\par
\noindent 2! · 3! · 4! · 5!\par
}
\end{sourceblock}


\begin{explainblock}
例题是学习这本书最重要的部分。不要只看答案，要把每个限制条件变成一个计数动作：先安排谁、再插入谁、有没有重复、需不需要除以对称性。

\smallskip
\textbf{初学者检查点：}
\begin{itemize}[leftmargin=2em,itemsep=0.25em]
\item 先写出没有限制时的总数。
\item 逐条加入限制条件，观察总数怎样变化。
\item 最后检查有没有把同一种方案重复算了多次。
\end{itemize}

\smallskip
\textbf{本章读法提醒：} 第一章所有题都先判断“顺序是否重要”。这个判断错了，后面公式一定错。
\end{explainblock}


\begin{sourceblock}{§ 1.2 组合}
{\small\sloppy
\noindent § 1.2. 组合\par
}
\end{sourceblock}


\begin{explainblock}
这一节先不要急着背公式。先问三个问题：对象是什么，哪些限制条件不能丢，最后到底是在数所有对象、还是在数某个统计量的分布。把这三件事说清楚，后面的公式才会有位置。

\smallskip
\textbf{初学者检查点：}
\begin{itemize}[leftmargin=2em,itemsep=0.25em]
\item 把本节出现的新对象写成一句话定义。
\item 把每个公式左边和右边分别翻译成“在数什么”。
\item 找一个最小非平凡例子，手工列出所有对象。
\end{itemize}

\smallskip
\textbf{本章读法提醒：} 第一章所有题都先判断“顺序是否重要”。这个判断错了，后面公式一定错。
\end{explainblock}


\begin{sourceblock}{定义 1.2.1 (组合) n 元集合 A 的一个 k-组合是从 A 中选出 k 个元素为组成的一}
{\small\sloppy
\noindent 定义 1.2.1 (组合) n 元集合 A 的一个 k-组合是从 A 中选出 k 个元素为组成的一\par
\noindent 个集合, 所以也称为 n 元集合的 k 元子集. n 元集合产生的 k 组合的个数或者 n 元\par
\noindent 集合 k 元子集的个数记为 nk .\par
}
\end{sourceblock}


\begin{explainblock}
定义段的任务是定边界。这里要特别注意参数范围、是否允许重复、是否考虑顺序、对象有没有标签。后面的每一道例题和证明，本质上都在反复使用这些边界。

\smallskip
\textbf{初学者检查点：}
\begin{itemize}[leftmargin=2em,itemsep=0.25em]
\item 圈出被定义的对象名称。
\item 圈出参数范围，例如 n、k 是否要求正整数。
\item 判断它是有序对象、无序对象、带标签对象还是结构对象。
\end{itemize}

\smallskip
\textbf{本章读法提醒：} 第一章所有题都先判断“顺序是否重要”。这个判断错了，后面公式一定错。
\end{explainblock}


\begin{sourceblock}{例 1.2.1 集合 \{1, 2, 3, 4\} 的 2-组合有:}
{\small\sloppy
\noindent 例 1.2.1 集合 \{1, 2, 3, 4\} 的 2-组合有:\par
\noindent 12, 13, 14, 23, 24, 34,\par
\noindent 共计 6 个, 即 2 = 6.\par
}
\end{sourceblock}


\begin{explainblock}
例题是学习这本书最重要的部分。不要只看答案，要把每个限制条件变成一个计数动作：先安排谁、再插入谁、有没有重复、需不需要除以对称性。

\smallskip
\textbf{初学者检查点：}
\begin{itemize}[leftmargin=2em,itemsep=0.25em]
\item 先写出没有限制时的总数。
\item 逐条加入限制条件，观察总数怎样变化。
\item 最后检查有没有把同一种方案重复算了多次。
\end{itemize}

\smallskip
\textbf{本章读法提醒：} 第一章所有题都先判断“顺序是否重要”。这个判断错了，后面公式一定错。
\end{explainblock}


\begin{sourceblock}{注记 1.2.2 n 元集合 A 的 k-组合也可看作是集合 A 的 k 元子集. 由该定义, 易}
{\small\sloppy
\noindent 注记 1.2.2 n 元集合 A 的 k-组合也可看作是集合 A 的 k 元子集. 由该定义, 易\par
\noindent 知, 当 n < k 时, nk = 0, 且 n0 = 1, n1 = n.\par
}
\end{sourceblock}


\begin{explainblock}
注记通常是在补充术语、记号、历史或一个容易忽略的约定。它未必给新公式，但常常决定你后面会不会把符号读错。

\smallskip
\textbf{初学者检查点：}
\begin{itemize}[leftmargin=2em,itemsep=0.25em]
\item 记下新符号的读法。
\item 确认这个约定是否只在本章使用。
\item 把注记里的限制条件补到前面的定义旁边。
\end{itemize}

\smallskip
\textbf{本章读法提醒：} 第一章所有题都先判断“顺序是否重要”。这个判断错了，后面公式一定错。
\end{explainblock}


\begin{sourceblock}{定理 1.2.3 设 1 ≤ k ≤ n, 则}
{\small\sloppy
\noindent 定理 1.2.3 设 1 ≤ k ≤ n, 则\par
\noindent n n!\par
\noindent = . (1.3)\par
\noindent k k!(n − k)!\par
\noindent 证明:\par
\noindent 由定理 1.1.2 可知, n 元集合 A 的 k-排列个数为 (n)k . 另一方面, 可按\par
\noindent 照如下方式计数该排列个数, 先从 n 元集合 A 中取出 k 个元素, 有 nk 方法, 再对\par
\noindent 这 k 个元素进行排列, 有 k! 种方式, 由乘法原理可知, 按这种计数方式计数 n 元集\par
\noindent 合 A 的 k-排列共有 k! nk 个数, 故\par
\noindent n\par
\noindent k! = (n)k ,\par
\noindent k\par
\noindent 所以\par
\noindent n (n)k n!\par
\noindent = = .\par
\noindent k k! k!(n − k)!\par
\noindent 由定理 1.2.3 可得下面的推论.\par
}
\end{sourceblock}


\begin{explainblock}
定理段不要只看结论，要看它把哪两个计数方式连在一起。组合学证明最常见的技巧是：左边按一种标准数，右边按另一种标准数，然后说明两边数的是同一批对象。

\smallskip
\textbf{初学者检查点：}
\begin{itemize}[leftmargin=2em,itemsep=0.25em]
\item 先复述定理的假设，不要把适用范围扩大。
\item 找出证明中的基本计数原则：乘法、加法、双射、递推、容斥或生成函数。
\item 检查最后的等式化简是否和前面的对象解释一致。
\end{itemize}

\smallskip
\textbf{本章读法提醒：} 第一章所有题都先判断“顺序是否重要”。这个判断错了，后面公式一定错。
\end{explainblock}


\begin{sourceblock}{推论 1.2.4 设 1 ≤ k ≤ n, 则}
{\small\sloppy
\noindent 推论 1.2.4 设 1 ≤ k ≤ n, 则\par
\noindent n n\par
\noindent = . (1.4)\par
\noindent k n−k\par
\noindent n\par
\noindent 证明: 该结论也可利用组合数的组合定义加以证明. 由定义 1.2.1, 可知 k 计\par
\noindent 数了 n 元集合的 k 元子集的个数, 而每一个 k 元子集对应一个 n − k 元的余集, 所\par
\noindent 以 n 元集合的 k 元子集的个数与 n 元集合的 n − k 元子集的个数相同. 再由定义\par
\noindent n\par
\noindent 1.2.1, 可知 n 元集合的 n − k 元子集的个数为 n−k , 故等式 (1.4) 成立.\par
}
\end{sourceblock}


\begin{explainblock}
定理段不要只看结论，要看它把哪两个计数方式连在一起。组合学证明最常见的技巧是：左边按一种标准数，右边按另一种标准数，然后说明两边数的是同一批对象。

\smallskip
\textbf{初学者检查点：}
\begin{itemize}[leftmargin=2em,itemsep=0.25em]
\item 先复述定理的假设，不要把适用范围扩大。
\item 找出证明中的基本计数原则：乘法、加法、双射、递推、容斥或生成函数。
\item 检查最后的等式化简是否和前面的对象解释一致。
\end{itemize}

\smallskip
\textbf{本章读法提醒：} 第一章所有题都先判断“顺序是否重要”。这个判断错了，后面公式一定错。
\end{explainblock}


\begin{sourceblock}{定理 1.2.5 设 1 ≤ k ≤ n − 1, 则}
{\small\sloppy
\noindent 定理 1.2.5 设 1 ≤ k ≤ n − 1, 则\par
\noindent n n−1 n−1\par
\noindent = + . (1.5)\par
\noindent k k k−1\par
\noindent n\par
\noindent 证 明: 方 法 一 (组 合 证 明): 由 k 的 定 义 知, 它 可 以 表 示 集 合 [n] =\par
\noindent \{1, 2, 3, . . . , n\} 的 k 元子集的个数. 这个 k 元子集可以根据元素 n 是否存在分\par
\noindent 成两类. 若 n 存在, 可以看作是将 n 放入到集合 [n − 1] 的所有 k − 1 元子集得到,\par
\noindent 共计 n−1\par
\noindent k−1 个; 若 n 不存在, 则可以看作是集合 [n − 1] = \{1, 2, 3 . . . , n − 1\} 的所有\par
\noindent k 元子集, 共计 n−1\par
\noindent k . 故等式 (1.5) 成立.\par
\noindent 方法二 (代数证明): 由定理 1.2.3 可知:\par
\noindent n−1 n−1 (n − 1)! (n − 1)!\par
\noindent + = +\par
\noindent k k−1 k!(n − 1 − k)! (k − 1)!(n − k)!\par
\noindent (n − 1)! · (n − k) + (n − 1)! · k\par
\noindent =\par
}
\end{sourceblock}


\begin{explainblock}
定理段不要只看结论，要看它把哪两个计数方式连在一起。组合学证明最常见的技巧是：左边按一种标准数，右边按另一种标准数，然后说明两边数的是同一批对象。

\smallskip
\textbf{初学者检查点：}
\begin{itemize}[leftmargin=2em,itemsep=0.25em]
\item 先复述定理的假设，不要把适用范围扩大。
\item 找出证明中的基本计数原则：乘法、加法、双射、递推、容斥或生成函数。
\item 检查最后的等式化简是否和前面的对象解释一致。
\end{itemize}

\smallskip
\textbf{本章读法提醒：} 第一章所有题都先判断“顺序是否重要”。这个判断错了，后面公式一定错。
\end{explainblock}


\begin{sourceblock}{第 1 章续读}
{\small\sloppy
\noindent k!(n − k)!\par
\noindent n!\par
\noindent =\par
\noindent k!(n − k)!\par
\noindent n\par
\noindent = .\par
\noindent k\par
\noindent 公 式 (1.5) 通 常 称 为 帕 斯 卡∗ 公 式 (Pascal’s formula). 事 实 上, 在 公 元 1303\par
\noindent ∗\par
\noindent 布莱士 · 帕斯卡 ( Blaise Pascal, 1623–1662)，法国数学家.\par
\noindent 年, 朱 世 杰 在 《四 元 玉 鉴》 的 前 页 就 用 三 角 形 数 表 记 载 了 该 关 系, 称 之 为 “古\par
\noindent 法七乘方图”（此处原书引用了一个图表，这里改用文字说明，不复现图表。）, 并在书中说明此表引自约公元 1050 年前贾宪的《释锁\par
\noindent 算术》. 而在公元 1261 年, 杨辉所著的《详解九章算法》一书中也辑录了这个三角\par
\noindent 形数表, 直到公元 1653 年帕斯卡在著作 Traité du triangle arithmétique [14] 中才提\par
\noindent 到这个三角形数表, 比杨辉要晚了近 400 年, 比贾宪晚了近 600 年. 现在称该三角形\par
\noindent 数列表为贾宪三角或杨辉三角或帕斯卡三角 （此处原书引用了一个图表，这里改用文字说明，不复现图表。）.\par
\noindent k\par
\noindent 0 1 2 3 4 5 6 7 8 ···\par
}
\end{sourceblock}


\begin{explainblock}
这一段是原书论述链的一部分。读的时候把它当作桥梁：它通常负责把上一个定义、例子或定理，引到下一个计数方法。

\smallskip
\textbf{初学者检查点：}
\begin{itemize}[leftmargin=2em,itemsep=0.25em]
\item 找出这一段承接的上一个概念。
\item 找出它准备引出的下一个概念。
\item 把中间的关键等式或构造单独抄出来。
\end{itemize}

\smallskip
\textbf{本章读法提醒：} 第一章所有题都先判断“顺序是否重要”。这个判断错了，后面公式一定错。
\end{explainblock}


\begin{sourceblock}{第 1 章续读}
{\small\sloppy
\noindent n\par
\noindent 7 1 7 21 35 35 21 7 1\par
\noindent 8 1 8 28 56 70 56 28 8 1\par
\noindent .. .. .. .. .. .. .. .. .. .. ..\par
\noindent . . . . . . . . . . .\par
}
\end{sourceblock}


\begin{explainblock}
这一段是原书论述链的一部分。读的时候把它当作桥梁：它通常负责把上一个定义、例子或定理，引到下一个计数方法。

\smallskip
\textbf{初学者检查点：}
\begin{itemize}[leftmargin=2em,itemsep=0.25em]
\item 找出这一段承接的上一个概念。
\item 找出它准备引出的下一个概念。
\item 把中间的关键等式或构造单独抄出来。
\end{itemize}

\smallskip
\textbf{本章读法提醒：} 第一章所有题都先判断“顺序是否重要”。这个判断错了，后面公式一定错。
\end{explainblock}


\begin{sourceblock}{定义 1.2.6从 n 元集合中可重复的选取 k 个元素不考虑次序的合成一组称为 n}
{\small\sloppy
\noindent 定义 1.2.6从 n 元集合中可重复的选取 k 个元素不考虑次序的合成一组称为 n\par
\noindent 元集合的 k-可重组合, 其个数记为 nk .\par
}
\end{sourceblock}


\begin{explainblock}
定义段的任务是定边界。这里要特别注意参数范围、是否允许重复、是否考虑顺序、对象有没有标签。后面的每一道例题和证明，本质上都在反复使用这些边界。

\smallskip
\textbf{初学者检查点：}
\begin{itemize}[leftmargin=2em,itemsep=0.25em]
\item 圈出被定义的对象名称。
\item 圈出参数范围，例如 n、k 是否要求正整数。
\item 判断它是有序对象、无序对象、带标签对象还是结构对象。
\end{itemize}

\smallskip
\textbf{本章读法提醒：} 第一章所有题都先判断“顺序是否重要”。这个判断错了，后面公式一定错。
\end{explainblock}


\begin{sourceblock}{例 1.2.2 设 A = \{1, 2, 3\}, 则 A 的 2-可重组合有 6 个:}
{\small\sloppy
\noindent 例 1.2.2 设 A = \{1, 2, 3\}, 则 A 的 2-可重组合有 6 个:\par
\noindent 11, 12, 13, 22, 23, 33.\par
\noindent 故 2 = 6.\par
}
\end{sourceblock}


\begin{explainblock}
例题是学习这本书最重要的部分。不要只看答案，要把每个限制条件变成一个计数动作：先安排谁、再插入谁、有没有重复、需不需要除以对称性。

\smallskip
\textbf{初学者检查点：}
\begin{itemize}[leftmargin=2em,itemsep=0.25em]
\item 先写出没有限制时的总数。
\item 逐条加入限制条件，观察总数怎样变化。
\item 最后检查有没有把同一种方案重复算了多次。
\end{itemize}

\smallskip
\textbf{本章读法提醒：} 第一章所有题都先判断“顺序是否重要”。这个判断错了，后面公式一定错。
\end{explainblock}


\begin{sourceblock}{定理 1.2.7 设 n, k 为正整数, 则}
{\small\sloppy
\noindent 定理 1.2.7 设 n, k 为正整数, 则\par
\noindent n n+k−1\par
\noindent = . (1.6)\par
\noindent k k\par
\noindent 证明: 设 A = \{1, 2, . . . , n\}, B = \{1, 2, . . . , n + k − 1\}. 由定义 1.2.1 和定义\par
\noindent n+k−1\par
\noindent 1.2.6 可知, 计数了集合 B 上的 k-组合的个数, nk 计数了集合 A 上的\par
\noindent k\par
\noindent k-可重组合的个数. 记 P 是集合 A 上所有 k-可重组合构成的集合, Q 是集合 B 上\par
\noindent 的 k-无重组合构成的集合, 则\par
\noindent n n+k−1\par
\noindent |P | = , |Q| = .\par
\noindent k k\par
\noindent 为证明等式 (1.6) 成立, 只需建立集合 P 与集合 Q 之间的双射. 现设 \{a1 , a2 , . . . , ak \}\par
\noindent 是集合 A 的一个 k-可重组合, 即 \{a1 , a2 , . . . , ak \} ∈ P , 不妨设\par
\noindent a1 ≤ a2 ≤ · · · ≤ ak .\par
\noindent 令\par
\noindent bi = ai + i − 1, 1 ≤ i ≤ k,\par
}
\end{sourceblock}


\begin{explainblock}
定理段不要只看结论，要看它把哪两个计数方式连在一起。组合学证明最常见的技巧是：左边按一种标准数，右边按另一种标准数，然后说明两边数的是同一批对象。

\smallskip
\textbf{初学者检查点：}
\begin{itemize}[leftmargin=2em,itemsep=0.25em]
\item 先复述定理的假设，不要把适用范围扩大。
\item 找出证明中的基本计数原则：乘法、加法、双射、递推、容斥或生成函数。
\item 检查最后的等式化简是否和前面的对象解释一致。
\end{itemize}

\smallskip
\textbf{本章读法提醒：} 第一章所有题都先判断“顺序是否重要”。这个判断错了，后面公式一定错。
\end{explainblock}


\begin{sourceblock}{第 1 章续读}
{\small\sloppy
\noindent 则\par
\noindent 1 ≤ b1 < b2 < · · · < bk ≤ n + k − 1,\par
\noindent 所以 \{b1 , b2 , . . . , bk \} 是集合 B 上的 k-组合, 即 \{b1 , b2 , . . . , bk \} ∈ Q; 反之, 若设\par
\noindent \{b1 , b2 , . . . , bk \} ∈ Q, 且满足\par
\noindent b1 < b2 < · · · < bk ,\par
\noindent 令\par
\noindent ai = bi − i + 1, 1 ≤ i ≤ k,\par
\noindent 则\par
\noindent 1 ≤ a1 ≤ a2 ≤ · · · ≤ ak ≤ n,\par
\noindent 可知 \{a1 , a2 , . . . , ak \} 是集合 A 的一个 k-可重组合, 即 \{a1 , a2 , . . . , ak \} ∈ P . 事实上,\par
\noindent 这样就建立了集合 P 与集合 Q 之间的双射, 故等式 (1.6) 成立.\par
}
\end{sourceblock}


\begin{explainblock}
这一段是原书论述链的一部分。读的时候把它当作桥梁：它通常负责把上一个定义、例子或定理，引到下一个计数方法。

\smallskip
\textbf{初学者检查点：}
\begin{itemize}[leftmargin=2em,itemsep=0.25em]
\item 找出这一段承接的上一个概念。
\item 找出它准备引出的下一个概念。
\item 把中间的关键等式或构造单独抄出来。
\end{itemize}

\smallskip
\textbf{本章读法提醒：} 第一章所有题都先判断“顺序是否重要”。这个判断错了，后面公式一定错。
\end{explainblock}


\begin{sourceblock}{例 1.2.3 设 n 与 k 为正整数, 求方程}
{\small\sloppy
\noindent 例 1.2.3 设 n 与 k 为正整数, 求方程\par
\noindent x1 + x2 + · · · + xn = k (1.7)\par
\noindent 的非负整数解的个数.\par
\noindent 解: 设 A = \{a1 , a2 , . . . , an \}, 则集合 A 的任何 k-可重组合都具有 x1 个 a1 , x2 个\par
\noindent a2 , . . ., xn 个 an , 其中, xi (i = 1, 2, . . . , n) 是非负整数且满足方程 (1.7) . 反之, 对于\par
\noindent 满足方程 (1.7) 的非负整数解 x1 , x2 , . . . , xn 都对应集合 A 的一个 k-可重组合. 所以\par
\noindent 方程 (1.7) 的非负解的个数为 n+k−1k .\par
}
\end{sourceblock}


\begin{explainblock}
例题是学习这本书最重要的部分。不要只看答案，要把每个限制条件变成一个计数动作：先安排谁、再插入谁、有没有重复、需不需要除以对称性。

\smallskip
\textbf{初学者检查点：}
\begin{itemize}[leftmargin=2em,itemsep=0.25em]
\item 先写出没有限制时的总数。
\item 逐条加入限制条件，观察总数怎样变化。
\item 最后检查有没有把同一种方案重复算了多次。
\end{itemize}

\smallskip
\textbf{本章读法提醒：} 第一章所有题都先判断“顺序是否重要”。这个判断错了，后面公式一定错。
\end{explainblock}


\begin{sourceblock}{例 1.2.4 设将 n 个元素绕成一个圈, 从中选取 k 个元素, 这 k 个元素恰好互不}
{\small\sloppy
\noindent 例 1.2.4 设将 n 个元素绕成一个圈, 从中选取 k 个元素, 这 k 个元素恰好互不\par
\noindent 相邻的方法数记为 g(n, k), 求 g(n, k).\par
\noindent 解: 设 f (n, k) 表示从排成一行的 n 个元素中选择 k 个互不相邻的元素的方法\par
\noindent 数, 则它相当于先将 n − k 元素排好, 然后将剩余的 k 个元素插入到产生的 n − k + 1\par
\noindent 个空中, 因此有\par
\noindent n−k+1\par
\noindent f (n, k) = . (1.8)\par
\noindent k\par
\noindent 设 g(n, k) 表示从排成一圈的 n 个元素中选择 k 个互不相邻的元素的方法数. 固\par
\noindent 定一个元素, 根据所选的 k 个元素中是否包含这个元素, 可以将选择方法分为两类,\par
\noindent 所以有\par
\noindent n n−k\par
\noindent g(n, k) = f (n − 1, k) + f (n − 3, k − 1) = . (1.9)\par
\noindent n−k k\par
}
\end{sourceblock}


\begin{explainblock}
例题是学习这本书最重要的部分。不要只看答案，要把每个限制条件变成一个计数动作：先安排谁、再插入谁、有没有重复、需不需要除以对称性。

\smallskip
\textbf{初学者检查点：}
\begin{itemize}[leftmargin=2em,itemsep=0.25em]
\item 先写出没有限制时的总数。
\item 逐条加入限制条件，观察总数怎样变化。
\item 最后检查有没有把同一种方案重复算了多次。
\end{itemize}

\smallskip
\textbf{本章读法提醒：} 第一章所有题都先判断“顺序是否重要”。这个判断错了，后面公式一定错。
\end{explainblock}


\begin{sourceblock}{§ 1.3 二项式定理}
{\small\sloppy
\noindent § 1.3. 二项式定理\par
\noindent 设 n 为正整数, k 为非负整数. 上节提到 n 元集合的 k 元子集个数称为组合数,\par
\noindent 记为 nk . 在这节我们会发现 nk 也是二项式 (x + y)n 展开后单项式 xk y n−k 的系\par
\noindent 数, 因此组合数也被称作二项式系数 (Binomial Coefficients).\par
}
\end{sourceblock}


\begin{explainblock}
这一节先不要急着背公式。先问三个问题：对象是什么，哪些限制条件不能丢，最后到底是在数所有对象、还是在数某个统计量的分布。把这三件事说清楚，后面的公式才会有位置。

\smallskip
\textbf{初学者检查点：}
\begin{itemize}[leftmargin=2em,itemsep=0.25em]
\item 把本节出现的新对象写成一句话定义。
\item 把每个公式左边和右边分别翻译成“在数什么”。
\item 找一个最小非平凡例子，手工列出所有对象。
\end{itemize}

\smallskip
\textbf{本章读法提醒：} 第一章所有题都先判断“顺序是否重要”。这个判断错了，后面公式一定错。
\end{explainblock}


\begin{sourceblock}{定理 1.3.1 (二项式定理) 设 n 是正整数, 对所有的 x 和 y ,}
{\small\sloppy
\noindent 定理 1.3.1 (二项式定理) 设 n 是正整数, 对所有的 x 和 y ,\par
\noindent n\par
\noindent X\par
\noindent n n k n−k\par
\noindent (x + y) = x y . (1.10)\par
\noindent k\par
\noindent k=0\par
\noindent 证明: (方法一) 下面我们用归纳法证明. 当 n = 1 时,\par
\noindent x y + x y = y + x,\par
\noindent 所以 n = 1 时, 等式 (1.10) 成立; 设 n = m 时, 等式 (1.10) 成立, 即\par
\noindent m\par
\noindent X\par
\noindent m m k m−k\par
\noindent (x + y) = x y ,\par
\noindent k\par
\noindent k=0\par
\noindent 则当 n = m + 1 时, 考虑\par
\noindent X\par
}
\end{sourceblock}


\begin{explainblock}
定理段不要只看结论，要看它把哪两个计数方式连在一起。组合学证明最常见的技巧是：左边按一种标准数，右边按另一种标准数，然后说明两边数的是同一批对象。

\smallskip
\textbf{初学者检查点：}
\begin{itemize}[leftmargin=2em,itemsep=0.25em]
\item 先复述定理的假设，不要把适用范围扩大。
\item 找出证明中的基本计数原则：乘法、加法、双射、递推、容斥或生成函数。
\item 检查最后的等式化简是否和前面的对象解释一致。
\end{itemize}

\smallskip
\textbf{本章读法提醒：} 第一章所有题都先判断“顺序是否重要”。这个判断错了，后面公式一定错。
\end{explainblock}


\begin{sourceblock}{第 1 章续读}
{\small\sloppy
\noindent m+1\par
\noindent m + 1 k m+1−k\par
\noindent m\par
\noindent X\par
\noindent m+1 k\par
\noindent m+1 m+1\par
\noindent x y =x +y + x + y m+1−k , (1.11)\par
\noindent k k\par
\noindent k=0 k=1\par
\noindent 又由定理 1.2.5 知\par
\noindent m+1 m m\par
\noindent = + ,\par
\noindent k k−1 k\par
\noindent 所以\par
\noindent X\par
\noindent m+1\par
\noindent m + 1 k m+1−k\par
\noindent m\par
}
\end{sourceblock}


\begin{explainblock}
这一段是原书论述链的一部分。读的时候把它当作桥梁：它通常负责把上一个定义、例子或定理，引到下一个计数方法。

\smallskip
\textbf{初学者检查点：}
\begin{itemize}[leftmargin=2em,itemsep=0.25em]
\item 找出这一段承接的上一个概念。
\item 找出它准备引出的下一个概念。
\item 把中间的关键等式或构造单独抄出来。
\end{itemize}

\smallskip
\textbf{本章读法提醒：} 第一章所有题都先判断“顺序是否重要”。这个判断错了，后面公式一定错。
\end{explainblock}


\begin{sourceblock}{第 1 章续读}
{\small\sloppy
\noindent X m\par
\noindent m\par
\noindent m+1 m+1\par
\noindent x y =x +y + + xk y m+1−k\par
\noindent k k−1 k\par
\noindent k=0 k=1\par
\noindent m\par
\noindent X m\par
\noindent X\par
\noindent m+1 m+1 m k m+1−k m\par
\noindent =x +y + x y + xk y m+1−k\par
\noindent k−1 k\par
\noindent k=1 k=1\par
\noindent X\par
\noindent m−1 m\par
\noindent m k+1 m−k X m k m+1−k\par
\noindent m+1 m+1\par
\noindent =x +y + x y + x y\par
}
\end{sourceblock}


\begin{explainblock}
这一段是原书论述链的一部分。读的时候把它当作桥梁：它通常负责把上一个定义、例子或定理，引到下一个计数方法。

\smallskip
\textbf{初学者检查点：}
\begin{itemize}[leftmargin=2em,itemsep=0.25em]
\item 找出这一段承接的上一个概念。
\item 找出它准备引出的下一个概念。
\item 把中间的关键等式或构造单独抄出来。
\end{itemize}

\smallskip
\textbf{本章读法提醒：} 第一章所有题都先判断“顺序是否重要”。这个判断错了，后面公式一定错。
\end{explainblock}


\begin{sourceblock}{第 1 章续读}
{\small\sloppy
\noindent k k\par
\noindent k=0 k=1\par
\noindent X m\par
\noindent m−1 m\par
\noindent X\par
\noindent m m k m−k m m m k m−k\par
\noindent =x x +x x y +y x +y x y\par
\noindent m k 0 k\par
\noindent k=0 k=1\par
\noindent m\par
\noindent X m\par
\noindent X\par
\noindent m k m−k m k m−k\par
\noindent =x x y +y x y ,\par
\noindent k k\par
\noindent k=0 k=0\par
\noindent 由等式 (1.11) 可得:\par
\noindent X\par
}
\end{sourceblock}


\begin{explainblock}
这一段是原书论述链的一部分。读的时候把它当作桥梁：它通常负责把上一个定义、例子或定理，引到下一个计数方法。

\smallskip
\textbf{初学者检查点：}
\begin{itemize}[leftmargin=2em,itemsep=0.25em]
\item 找出这一段承接的上一个概念。
\item 找出它准备引出的下一个概念。
\item 把中间的关键等式或构造单独抄出来。
\end{itemize}

\smallskip
\textbf{本章读法提醒：} 第一章所有题都先判断“顺序是否重要”。这个判断错了，后面公式一定错。
\end{explainblock}


\begin{sourceblock}{第 1 章续读}
{\small\sloppy
\noindent m+1\par
\noindent m + 1 k m+1−k\par
\noindent x y = x(x + y)m + y(x + y)m\par
\noindent k\par
\noindent k=0\par
\noindent = (x + y)m+1 .\par
\noindent 即对 n = m + 1 等式 (1.10) 也成立. 由归纳法知, 对任意正整数 n 等式 (1.10) 成立.\par
\noindent (方法二) 将 (x + y)n 写成 n 个 x + y 因子的乘积形式, 即\par
\noindent (x + y)n = (x + y)(x + y) · · · (x + y).\par
\noindent 将等式右边连乘展开, 显然展开的每一项都可以写成 xk y n−k 的形式, 其中 k =\par
\noindent 0, 1, . . . , n. 而 xn−k y k 可以从任意 k 个因式中选取 k 个 x 并在剩余的 n − k 个因式\par
\noindent 中选择 y 相乘得到, 共计有 nk 种选择. 故等式 (1.10) 成立.\par
\noindent 特别地, 在等式 (1.10) 中令 y = 1 可得\par
\noindent n\par
\noindent X\par
\noindent n n\par
\noindent (1 + x) = xk . (1.12)\par
\noindent k\par
}
\end{sourceblock}


\begin{explainblock}
这一段是原书论述链的一部分。读的时候把它当作桥梁：它通常负责把上一个定义、例子或定理，引到下一个计数方法。

\smallskip
\textbf{初学者检查点：}
\begin{itemize}[leftmargin=2em,itemsep=0.25em]
\item 找出这一段承接的上一个概念。
\item 找出它准备引出的下一个概念。
\item 把中间的关键等式或构造单独抄出来。
\end{itemize}

\smallskip
\textbf{本章读法提醒：} 第一章所有题都先判断“顺序是否重要”。这个判断错了，后面公式一定错。
\end{explainblock}


\begin{sourceblock}{第 1 章续读}
{\small\sloppy
\noindent k=0\par
\noindent 通过对等式 (1.12) 两边的 x 赋值或关于 x 求导后再赋值可得到一些关于二项式系\par
\noindent 数的恒等式. 比如在等式 (1.12) 令 x = 1 可得到下面的恒等式.\par
}
\end{sourceblock}


\begin{explainblock}
这一段是原书论述链的一部分。读的时候把它当作桥梁：它通常负责把上一个定义、例子或定理，引到下一个计数方法。

\smallskip
\textbf{初学者检查点：}
\begin{itemize}[leftmargin=2em,itemsep=0.25em]
\item 找出这一段承接的上一个概念。
\item 找出它准备引出的下一个概念。
\item 把中间的关键等式或构造单独抄出来。
\end{itemize}

\smallskip
\textbf{本章读法提醒：} 第一章所有题都先判断“顺序是否重要”。这个判断错了，后面公式一定错。
\end{explainblock}


\begin{sourceblock}{推论 1.3.2 设 n ≥ 1, 则}
{\small\sloppy
\noindent 推论 1.3.2 设 n ≥ 1, 则\par
\noindent n\par
\noindent X n\par
\noindent = 2n . (1.13)\par
\noindent k\par
\noindent k=0\par
\noindent 这里给出该推论的一种组合证明.\par
\noindent 证明: 因为 nk 计数了 n 元集合的 k 元子集的个数, 所以等式左边的和式是 n\par
\noindent 元集合所有子集的个数; 另一方面, 对于 n 元集合的任意子集, 我们都可以根据集合\par
\noindent 中的元素是否属于这个子集合来确定, 每个元素都有两种选择方式, 共计 2n 种方式,\par
\noindent 每种方式唯一确定一个子集, 所以共有 2n 个子集, 综上分析知等式 (1.13) 成立.\par
\noindent 在等式 (1.12) 两边令 x = −1 可得到下面的恒等式.\par
}
\end{sourceblock}


\begin{explainblock}
定理段不要只看结论，要看它把哪两个计数方式连在一起。组合学证明最常见的技巧是：左边按一种标准数，右边按另一种标准数，然后说明两边数的是同一批对象。

\smallskip
\textbf{初学者检查点：}
\begin{itemize}[leftmargin=2em,itemsep=0.25em]
\item 先复述定理的假设，不要把适用范围扩大。
\item 找出证明中的基本计数原则：乘法、加法、双射、递推、容斥或生成函数。
\item 检查最后的等式化简是否和前面的对象解释一致。
\end{itemize}

\smallskip
\textbf{本章读法提醒：} 第一章所有题都先判断“顺序是否重要”。这个判断错了，后面公式一定错。
\end{explainblock}


\begin{sourceblock}{推论 1.3.3 设 n ≥ 0, 则}
{\small\sloppy
\noindent 推论 1.3.3 设 n ≥ 0, 则\par
\noindent X\par
\noindent n\par
\noindent n\par
\noindent (−1)k = 0, (1.14)\par
\noindent k\par
\noindent k=0\par
\noindent 整理可得:\par
\noindent n n n n\par
\noindent + + ··· = + + · · · = 2n−1 . (1.15)\par
\noindent 这里给出该推论的一种组合证明.\par
\noindent 证明: 设 A = \{1, 2, . . . , n\}, 则 n\par
\noindent k 表示集合 A 的 k 元子集的个数. 不妨记 P\par
\noindent 为集合 A 的所有子集构成的子集族, 则由恒等式 (1.13) 知 |P | = 2n . 现对 P 中元素\par
\noindent 作如下处理: 设 B ∈ P , 若 B 中有奇数个元素, 则将 B 标记为 “−”; 若 B 中有偶数\par
\noindent 个元素, 则将 B 标记为 “+”. 记标记为 “-” 的元素所构成的集合为 P − , 标记为 “+”\par
\noindent 的元素所构成的集合为 P + , 则等式 (1.14) 表示对 P 作上述处理后, P 中符号为 “-”\par
\noindent 的子集的个数与符号为 “+” 的子集个数相同. 因此, 为证明该结论, 只需建立 P − 与\par
}
\end{sourceblock}


\begin{explainblock}
定理段不要只看结论，要看它把哪两个计数方式连在一起。组合学证明最常见的技巧是：左边按一种标准数，右边按另一种标准数，然后说明两边数的是同一批对象。

\smallskip
\textbf{初学者检查点：}
\begin{itemize}[leftmargin=2em,itemsep=0.25em]
\item 先复述定理的假设，不要把适用范围扩大。
\item 找出证明中的基本计数原则：乘法、加法、双射、递推、容斥或生成函数。
\item 检查最后的等式化简是否和前面的对象解释一致。
\end{itemize}

\smallskip
\textbf{本章读法提醒：} 第一章所有题都先判断“顺序是否重要”。这个判断错了，后面公式一定错。
\end{explainblock}


\begin{sourceblock}{第 1 章续读}
{\small\sloppy
\noindent P + 之间的双射.\par
\noindent 设 B 是 P − 中任意一元素, 对 B 作如下变换 f :\par
\noindent (1) 若 1 ∈ B, 则令\par
\noindent f (B) = B\textbackslash{}\{1\},\par
\noindent (2) 若 1 ∈\par
\noindent / B1 , 则令\par
\noindent f (B) = B ∪ \{1\},\par
\noindent 显然 f (B) ∈ P + , 故 f 是从 P − 到 P + 的映射, 由 f 的定义知其显然可逆, 即 f 是\par
\noindent 从 P − 到 P + 的双射, 所以等式 (1.14) 成立. 又 |P | = 2n , 故等式 (1.15) 也成立.\par
}
\end{sourceblock}


\begin{explainblock}
这一段是原书论述链的一部分。读的时候把它当作桥梁：它通常负责把上一个定义、例子或定理，引到下一个计数方法。

\smallskip
\textbf{初学者检查点：}
\begin{itemize}[leftmargin=2em,itemsep=0.25em]
\item 找出这一段承接的上一个概念。
\item 找出它准备引出的下一个概念。
\item 把中间的关键等式或构造单独抄出来。
\end{itemize}

\smallskip
\textbf{本章读法提醒：} 第一章所有题都先判断“顺序是否重要”。这个判断错了，后面公式一定错。
\end{explainblock}


\begin{sourceblock}{注记 1.3.4 由证明过程可知, 等式 (1.14) 与等式 (1.15) 实际上指任意集合的奇子}
{\small\sloppy
\noindent 注记 1.3.4 由证明过程可知, 等式 (1.14) 与等式 (1.15) 实际上指任意集合的奇子\par
\noindent 集与偶子集的个数相同.\par
\noindent 若对 (1.12) 等式两边关于 x 求导可得\par
\noindent X\par
\noindent n\par
\noindent n−1 n k−1\par
\noindent n(1 + x) = k x . (1.16)\par
\noindent k\par
\noindent k=1\par
\noindent 在等式 (1.16) 两边令 x = 1 可得下面的恒等式.\par
}
\end{sourceblock}


\begin{explainblock}
注记通常是在补充术语、记号、历史或一个容易忽略的约定。它未必给新公式，但常常决定你后面会不会把符号读错。

\smallskip
\textbf{初学者检查点：}
\begin{itemize}[leftmargin=2em,itemsep=0.25em]
\item 记下新符号的读法。
\item 确认这个约定是否只在本章使用。
\item 把注记里的限制条件补到前面的定义旁边。
\end{itemize}

\smallskip
\textbf{本章读法提醒：} 第一章所有题都先判断“顺序是否重要”。这个判断错了，后面公式一定错。
\end{explainblock}


\begin{sourceblock}{推论 1.3.5 设 n 为正整数, 则}
{\small\sloppy
\noindent 推论 1.3.5 设 n 为正整数, 则\par
\noindent X\par
\noindent n\par
\noindent n\par
\noindent k = n2n−1 . (1.17)\par
\noindent k\par
\noindent k=1\par
\noindent n\par
\noindent n\par
\noindent 由 k 的组合定义知它对所有非负整数均有定义, 但从 k 的显式表达式 (1.3),\par
\noindent 即\par
\noindent n n(n − 1) · · · (n − k + 1)\par
\noindent = ,\par
\noindent k k(k − 1) · · · 1\par
\noindent n\par
\noindent 可看出 k 对 k 为非负整数及 n 为实数仍有意义, 但此时可能不再具有组合意义.\par
}
\end{sourceblock}


\begin{explainblock}
定理段不要只看结论，要看它把哪两个计数方式连在一起。组合学证明最常见的技巧是：左边按一种标准数，右边按另一种标准数，然后说明两边数的是同一批对象。

\smallskip
\textbf{初学者检查点：}
\begin{itemize}[leftmargin=2em,itemsep=0.25em]
\item 先复述定理的假设，不要把适用范围扩大。
\item 找出证明中的基本计数原则：乘法、加法、双射、递推、容斥或生成函数。
\item 检查最后的等式化简是否和前面的对象解释一致。
\end{itemize}

\smallskip
\textbf{本章读法提醒：} 第一章所有题都先判断“顺序是否重要”。这个判断错了，后面公式一定错。
\end{explainblock}


\begin{sourceblock}{定义 1.3.6 设 α 是实数, k 为任意整数, 定义}
{\small\sloppy
\noindent 定义 1.3.6 设 α 是实数, k 为任意整数, 定义\par
\noindent \par
\noindent  α(α−1)···(α−k+1) ,\par
\noindent \par
\noindent \par
\noindent  k·(k−1)···2·1\par
\noindent k ≥ 1,\par
\noindent α\par
\noindent = 1, k = 0, (1.18)\par
\noindent k \par
\noindent \par
\noindent \par
\noindent  0, k ≤ −1,\par
\noindent α\par
\noindent 并称 k 为广义的二项式系数.\par
\noindent α\par
}
\end{sourceblock}


\begin{explainblock}
定义段的任务是定边界。这里要特别注意参数范围、是否允许重复、是否考虑顺序、对象有没有标签。后面的每一道例题和证明，本质上都在反复使用这些边界。

\smallskip
\textbf{初学者检查点：}
\begin{itemize}[leftmargin=2em,itemsep=0.25em]
\item 圈出被定义的对象名称。
\item 圈出参数范围，例如 n、k 是否要求正整数。
\item 判断它是有序对象、无序对象、带标签对象还是结构对象。
\end{itemize}

\smallskip
\textbf{本章读法提醒：} 第一章所有题都先判断“顺序是否重要”。这个判断错了，后面公式一定错。
\end{explainblock}


\begin{sourceblock}{注记 1.3.7 由定义可知 k 正整数时, k 实事上是 α 的 k 次多项式.}
{\small\sloppy
\noindent 注记 1.3.7 由定义可知 k 正整数时, k 实事上是 α 的 k 次多项式.\par
}
\end{sourceblock}


\begin{explainblock}
注记通常是在补充术语、记号、历史或一个容易忽略的约定。它未必给新公式，但常常决定你后面会不会把符号读错。

\smallskip
\textbf{初学者检查点：}
\begin{itemize}[leftmargin=2em,itemsep=0.25em]
\item 记下新符号的读法。
\item 确认这个约定是否只在本章使用。
\item 把注记里的限制条件补到前面的定义旁边。
\end{itemize}

\smallskip
\textbf{本章读法提醒：} 第一章所有题都先判断“顺序是否重要”。这个判断错了，后面公式一定错。
\end{explainblock}


\begin{sourceblock}{例 1.3.1 (1)}
{\small\sloppy
\noindent 例 1.3.1 (1)\par
\noindent −n −n · (−n − 1) · · · (−n − k + 1)\par
\noindent =\par
\noindent k k(k − 1) · · · 1\par
\noindent (n + k − 1)(n + k − 2) · · · (n − 1)n\par
\noindent = (−1)k\par
\noindent k(k − 1) · · · 1\par
\noindent k n+k−1\par
\noindent = (−1) . (1.19)\par
\noindent k\par
\noindent (2)\par
\noindent 2 −1 · · · 21 − k + 1\par
\noindent =\par
\noindent k k!\par
\noindent 1 1 · 3 · · · (2k − 3)\par
\noindent = (−1)k−1 · ·\par
\noindent 2k k!\par
\noindent = (−1)k−1\par
}
\end{sourceblock}


\begin{explainblock}
例题是学习这本书最重要的部分。不要只看答案，要把每个限制条件变成一个计数动作：先安排谁、再插入谁、有没有重复、需不需要除以对称性。

\smallskip
\textbf{初学者检查点：}
\begin{itemize}[leftmargin=2em,itemsep=0.25em]
\item 先写出没有限制时的总数。
\item 逐条加入限制条件，观察总数怎样变化。
\item 最后检查有没有把同一种方案重复算了多次。
\end{itemize}

\smallskip
\textbf{本章读法提醒：} 第一章所有题都先判断“顺序是否重要”。这个判断错了，后面公式一定错。
\end{explainblock}


\begin{sourceblock}{第 1 章续读}
{\small\sloppy
\noindent · · . (1.20)\par
\noindent k · 22k−1 k−1\par
}
\end{sourceblock}


\begin{explainblock}
这一段是原书论述链的一部分。读的时候把它当作桥梁：它通常负责把上一个定义、例子或定理，引到下一个计数方法。

\smallskip
\textbf{初学者检查点：}
\begin{itemize}[leftmargin=2em,itemsep=0.25em]
\item 找出这一段承接的上一个概念。
\item 找出它准备引出的下一个概念。
\item 把中间的关键等式或构造单独抄出来。
\end{itemize}

\smallskip
\textbf{本章读法提醒：} 第一章所有题都先判断“顺序是否重要”。这个判断错了，后面公式一定错。
\end{explainblock}


\begin{sourceblock}{定理 1.3.8 (牛顿二项式定理) 设 α 和 x 均为实数且 |x| < 1, 则}
{\small\sloppy
\noindent 定理 1.3.8 (牛顿二项式定理) 设 α 和 x 均为实数且 |x| < 1, 则\par
\noindent ∞\par
\noindent X\par
\noindent α α\par
\noindent (1 + x) = xk . (1.21)\par
\noindent k\par
\noindent k=0\par
\noindent 证明: (方法一) 显然 (1 + x)α 任意阶导数存在, 设 k 为任意非负整数, 则\par
\noindent (1 + x)α 的 k 阶导数为\par
\noindent (α)k (1 + x)α−k ,\par
\noindent 所以 (1 + x)α 在 x = 0 处泰勒展开可得\par
\noindent ∞\par
\noindent X ∞\par
\noindent X\par
\noindent α xk α k\par
\noindent (1 + x) = (α)k = x .\par
\noindent k! k\par
\noindent k=0 k=0\par
}
\end{sourceblock}


\begin{explainblock}
定理段不要只看结论，要看它把哪两个计数方式连在一起。组合学证明最常见的技巧是：左边按一种标准数，右边按另一种标准数，然后说明两边数的是同一批对象。

\smallskip
\textbf{初学者检查点：}
\begin{itemize}[leftmargin=2em,itemsep=0.25em]
\item 先复述定理的假设，不要把适用范围扩大。
\item 找出证明中的基本计数原则：乘法、加法、双射、递推、容斥或生成函数。
\item 检查最后的等式化简是否和前面的对象解释一致。
\end{itemize}

\smallskip
\textbf{本章读法提醒：} 第一章所有题都先判断“顺序是否重要”。这个判断错了，后面公式一定错。
\end{explainblock}


\begin{sourceblock}{第 1 章续读}
{\small\sloppy
\noindent (方法二) 令\par
\noindent ∞\par
\noindent X α k\par
\noindent f (x) = x ,\par
\noindent k\par
\noindent k=0\par
\noindent 则 f (x) 的收敛半径为:\par
\noindent α\par
\noindent (α − n)\par
\noindent lim n+1\par
\noindent α\par
\noindent = lim = 1.\par
\noindent n→∞\par
\noindent n\par
\noindent n→∞ n+1\par
\noindent 所以 f (x) 在 |x| < 1 时一致收敛, 可以逐项求导. 对 f (x) 逐项求导:\par
\noindent ∞\par
\noindent X\par
}
\end{sourceblock}


\begin{explainblock}
这一段是原书论述链的一部分。读的时候把它当作桥梁：它通常负责把上一个定义、例子或定理，引到下一个计数方法。

\smallskip
\textbf{初学者检查点：}
\begin{itemize}[leftmargin=2em,itemsep=0.25em]
\item 找出这一段承接的上一个概念。
\item 找出它准备引出的下一个概念。
\item 把中间的关键等式或构造单独抄出来。
\end{itemize}

\smallskip
\textbf{本章读法提醒：} 第一章所有题都先判断“顺序是否重要”。这个判断错了，后面公式一定错。
\end{explainblock}


\begin{sourceblock}{第 1 章续读}
{\small\sloppy
\noindent ′ α\par
\noindent f (x) = kxk−1\par
\noindent k\par
\noindent k=1\par
\noindent X∞\par
\noindent α\par
\noindent = (k + 1)xk\par
\noindent k+1\par
\noindent k=0\par
\noindent X∞\par
\noindent α\par
\noindent = (α − k)xk\par
\noindent k\par
\noindent k=0\par
\noindent X∞ ∞\par
\noindent α k X α\par
\noindent =α x − kxk\par
\noindent k k\par
}
\end{sourceblock}


\begin{explainblock}
这一段是原书论述链的一部分。读的时候把它当作桥梁：它通常负责把上一个定义、例子或定理，引到下一个计数方法。

\smallskip
\textbf{初学者检查点：}
\begin{itemize}[leftmargin=2em,itemsep=0.25em]
\item 找出这一段承接的上一个概念。
\item 找出它准备引出的下一个概念。
\item 把中间的关键等式或构造单独抄出来。
\end{itemize}

\smallskip
\textbf{本章读法提醒：} 第一章所有题都先判断“顺序是否重要”。这个判断错了，后面公式一定错。
\end{explainblock}


\begin{sourceblock}{第 1 章续读}
{\small\sloppy
\noindent k=0 k=1\par
\noindent X α\par
\noindent ∞\par
\noindent = αf (x) − x kxk−1\par
\noindent k\par
\noindent k=1\par
\noindent = αf (x) − xf ′ (x),\par
\noindent 求得微分方程:\par
\noindent (1 + x)f ′ (x) = αf (x),\par
\noindent 即\par
\noindent f ′ (x) α\par
\noindent = ,\par
\noindent f (x) 1+x\par
\noindent 解得:\par
\noindent log(f (x)) = C1 + α log(1 + x)\par
\noindent = C1 + log(1 + x)α ,\par
\noindent 所以\par
\noindent f (x) = C2 · (1 + x)α .\par
}
\end{sourceblock}


\begin{explainblock}
这一段是原书论述链的一部分。读的时候把它当作桥梁：它通常负责把上一个定义、例子或定理，引到下一个计数方法。

\smallskip
\textbf{初学者检查点：}
\begin{itemize}[leftmargin=2em,itemsep=0.25em]
\item 找出这一段承接的上一个概念。
\item 找出它准备引出的下一个概念。
\item 把中间的关键等式或构造单独抄出来。
\end{itemize}

\smallskip
\textbf{本章读法提醒：} 第一章所有题都先判断“顺序是否重要”。这个判断错了，后面公式一定错。
\end{explainblock}


\begin{sourceblock}{第 1 章续读}
{\small\sloppy
\noindent 将 x = 0 带入上式, 由 f (0) = 1 解得 C2 = 1, 所以\par
\noindent f (x) = (1 + x)α .\par
\noindent 也即是:\par
\noindent ∞\par
\noindent X α\par
\noindent (1 + x)α = xk .\par
\noindent k\par
\noindent k=0\par
}
\end{sourceblock}


\begin{explainblock}
这一段是原书论述链的一部分。读的时候把它当作桥梁：它通常负责把上一个定义、例子或定理，引到下一个计数方法。

\smallskip
\textbf{初学者检查点：}
\begin{itemize}[leftmargin=2em,itemsep=0.25em]
\item 找出这一段承接的上一个概念。
\item 找出它准备引出的下一个概念。
\item 把中间的关键等式或构造单独抄出来。
\end{itemize}

\smallskip
\textbf{本章读法提醒：} 第一章所有题都先判断“顺序是否重要”。这个判断错了，后面公式一定错。
\end{explainblock}


\begin{sourceblock}{例 1.3.2 设 |x| < 1, 证明下面两个命题成立:}
{\small\sloppy
\noindent 例 1.3.2 设 |x| < 1, 证明下面两个命题成立:\par
\noindent ∞\par
\noindent X\par
\noindent √ (−1)k−1 2k − 2 k\par
\noindent (1) 1+x=1+ x . (1.22)\par
\noindent k · 22k−1 k − 1\par
\noindent k=1\par
\noindent ∞\par
\noindent X\par
\noindent (2) √ = x . (1.23)\par
\noindent 证明: (1) 由牛顿二项式定理知\par
\noindent ∞ 1\par
\noindent X\par
\noindent √\par
\noindent 1+x= 2 xk ,\par
\noindent k\par
\noindent k=0\par
\noindent 由 (1.18) 知 k = 0 时, 1\par
}
\end{sourceblock}


\begin{explainblock}
例题是学习这本书最重要的部分。不要只看答案，要把每个限制条件变成一个计数动作：先安排谁、再插入谁、有没有重复、需不需要除以对称性。

\smallskip
\textbf{初学者检查点：}
\begin{itemize}[leftmargin=2em,itemsep=0.25em]
\item 先写出没有限制时的总数。
\item 逐条加入限制条件，观察总数怎样变化。
\item 最后检查有没有把同一种方案重复算了多次。
\end{itemize}

\smallskip
\textbf{本章读法提醒：} 第一章所有题都先判断“顺序是否重要”。这个判断错了，后面公式一定错。
\end{explainblock}


\begin{sourceblock}{第 1 章续读}
{\small\sloppy
\noindent k\par
\noindent 当 k 6= 0 时,\par
\noindent 2 · 2 −1 · · · 12 − k + 1\par
\noindent k k!\par
\noindent 1 · (−1) · (−3) · · · (−2k − 3)\par
\noindent =\par
\noindent 2k k!\par
\noindent (2k − 3)!!\par
\noindent = (−1)k−1\par
\noindent 2k k!\par
\noindent (2k − 3)!! · 2k−1 (k − 1)!\par
\noindent = (−1)k−1\par
\noindent 22k−1 · k! · (k − 1)\par
\noindent (−1)k−1 2k − 2\par
\noindent = ,\par
\noindent k · 22k−1 k − 1\par
\noindent 所以 (1.22) 成立.\par
\noindent (2) 同样由牛顿二项式定理可得:\par
}
\end{sourceblock}


\begin{explainblock}
这一段是原书论述链的一部分。读的时候把它当作桥梁：它通常负责把上一个定义、例子或定理，引到下一个计数方法。

\smallskip
\textbf{初学者检查点：}
\begin{itemize}[leftmargin=2em,itemsep=0.25em]
\item 找出这一段承接的上一个概念。
\item 找出它准备引出的下一个概念。
\item 把中间的关键等式或构造单独抄出来。
\end{itemize}

\smallskip
\textbf{本章读法提醒：} 第一章所有题都先判断“顺序是否重要”。这个判断错了，后面公式一定错。
\end{explainblock}


\begin{sourceblock}{第 1 章续读}
{\small\sloppy
\noindent X∞ 1\par
\noindent √ = (1 + x) = ,\par
\noindent 1+x k\par
\noindent k=0\par
\noindent 由 (1.18) 知\par
\noindent −2 (− 21 ) · (− 32 ) · · · (− 2k−1\par
\noindent = x\par
\noindent k k!\par
\noindent 1 · 3 · · · (2k − 1)\par
\noindent = (−1)k\par
\noindent k!\par
\noindent 1 · 3 · · · (2k − 1) · 2k · k!\par
\noindent = (−1)k\par
\noindent 22k · k! · k!\par
\noindent (−1)k 1 · 2 · · · (2k − 1) · 2k\par
\noindent = ·\par
\noindent 22k k! · k!\par
\noindent (−1)k 2k\par
}
\end{sourceblock}


\begin{explainblock}
这一段是原书论述链的一部分。读的时候把它当作桥梁：它通常负责把上一个定义、例子或定理，引到下一个计数方法。

\smallskip
\textbf{初学者检查点：}
\begin{itemize}[leftmargin=2em,itemsep=0.25em]
\item 找出这一段承接的上一个概念。
\item 找出它准备引出的下一个概念。
\item 把中间的关键等式或构造单独抄出来。
\end{itemize}

\smallskip
\textbf{本章读法提醒：} 第一章所有题都先判断“顺序是否重要”。这个判断错了，后面公式一定错。
\end{explainblock}


\begin{sourceblock}{第 1 章续读}
{\small\sloppy
\noindent = 2k ,\par
\noindent 所以 (1.23) 成立.\par
}
\end{sourceblock}


\begin{explainblock}
这一段是原书论述链的一部分。读的时候把它当作桥梁：它通常负责把上一个定义、例子或定理，引到下一个计数方法。

\smallskip
\textbf{初学者检查点：}
\begin{itemize}[leftmargin=2em,itemsep=0.25em]
\item 找出这一段承接的上一个概念。
\item 找出它准备引出的下一个概念。
\item 把中间的关键等式或构造单独抄出来。
\end{itemize}

\smallskip
\textbf{本章读法提醒：} 第一章所有题都先判断“顺序是否重要”。这个判断错了，后面公式一定错。
\end{explainblock}


\begin{sourceblock}{注记 1.3.9 在等式 (1.22) 与 (1.23) 两端对 x 赋值也可得到一些恒等式.}
{\small\sloppy
\noindent 注记 1.3.9 在等式 (1.22) 与 (1.23) 两端对 x 赋值也可得到一些恒等式.\par
}
\end{sourceblock}


\begin{explainblock}
注记通常是在补充术语、记号、历史或一个容易忽略的约定。它未必给新公式，但常常决定你后面会不会把符号读错。

\smallskip
\textbf{初学者检查点：}
\begin{itemize}[leftmargin=2em,itemsep=0.25em]
\item 记下新符号的读法。
\item 确认这个约定是否只在本章使用。
\item 把注记里的限制条件补到前面的定义旁边。
\end{itemize}

\smallskip
\textbf{本章读法提醒：} 第一章所有题都先判断“顺序是否重要”。这个判断错了，后面公式一定错。
\end{explainblock}


\begin{sourceblock}{§ 1.4 错位排列}
{\small\sloppy
\noindent § 1.4. 错位排列\par
\noindent 某人写了 n 封不同的信准备寄到不同的地址, 现将相应的地址写到 n 个信封\par
\noindent 上, 那么他把这 n 封信都装到不是它对应的信封中的方法有多少种? 这个问题\par
\noindent 被欧拉称为 “组合数学中的一个奇妙问题”(拉丁文: a quaestio curiosa ex doctrina\par
\noindent combinationis). 它最早由尼古拉 · 伯努利 † 提出, 最后欧拉和伯努利都独立地给出了\par
\noindent 解答. 这个问题因此也称为伯努利—欧拉错装信封问题. 值得一提的是欧拉在解决\par
\noindent 这个问题的同时也发展了组合数学中一个非常强大的工具—递推方法. 而实事上, 早\par
\noindent 在伯努利之前, 德蒙马特 ‡ 就已经研究过类似的问题—“相遇”(Treize 或 Rencontre)\par
\noindent 的纸牌游戏. 这个游戏规则: 发牌的人洗好 A − K 13 张纸牌, 每次发一张, 发第一张\par
\noindent 的时候同时说 “1”, 发第二张纸牌时说 “2”, 依次类推, 发牌结束后将牌翻开, 若和发\par
\noindent 牌时所说的都一致则发牌者获胜. 德蒙马特因解决了此纸牌游戏发牌者获胜的概率\par
\noindent 而在组合学的历史上得以留名.\par
}
\end{sourceblock}


\begin{explainblock}
这一节先不要急着背公式。先问三个问题：对象是什么，哪些限制条件不能丢，最后到底是在数所有对象、还是在数某个统计量的分布。把这三件事说清楚，后面的公式才会有位置。

\smallskip
\textbf{初学者检查点：}
\begin{itemize}[leftmargin=2em,itemsep=0.25em]
\item 把本节出现的新对象写成一句话定义。
\item 把每个公式左边和右边分别翻译成“在数什么”。
\item 找一个最小非平凡例子，手工列出所有对象。
\end{itemize}

\smallskip
\textbf{本章读法提醒：} 第一章所有题都先判断“顺序是否重要”。这个判断错了，后面公式一定错。
\end{explainblock}


\begin{sourceblock}{定义 1.4.1 (错位排列) 设 π = π1 π2 · · · πn 是集合 [n] 的一个排列, 若任意的}
{\small\sloppy
\noindent 定义 1.4.1 (错位排列) 设 π = π1 π2 · · · πn 是集合 [n] 的一个排列, 若任意的\par
\noindent i(i = 1, 2, . . . , n) 都不在排列的第 i 个位置, 即 π(i) 6= i, 则称 π 为错位排列, 简称错\par
\noindent 排 (derangement). 记 Dn 为 n 元集合错排的个数, 并规定 D0 = 1.\par
}
\end{sourceblock}


\begin{explainblock}
定义段的任务是定边界。这里要特别注意参数范围、是否允许重复、是否考虑顺序、对象有没有标签。后面的每一道例题和证明，本质上都在反复使用这些边界。

\smallskip
\textbf{初学者检查点：}
\begin{itemize}[leftmargin=2em,itemsep=0.25em]
\item 圈出被定义的对象名称。
\item 圈出参数范围，例如 n、k 是否要求正整数。
\item 判断它是有序对象、无序对象、带标签对象还是结构对象。
\end{itemize}

\smallskip
\textbf{本章读法提醒：} 第一章所有题都先判断“顺序是否重要”。这个判断错了，后面公式一定错。
\end{explainblock}


\begin{sourceblock}{例 1.4.1 对于集合 \{1, 2, 3\}, 它全排列有:}
{\small\sloppy
\noindent 例 1.4.1 对于集合 \{1, 2, 3\}, 它全排列有:\par
\noindent 123, 132, 213, 231, 312, 321,\par
\noindent 则全排数为 3! = 6, 其中只有 231 与 312 是错排. 所以 D3 = 2.\par
}
\end{sourceblock}


\begin{explainblock}
例题是学习这本书最重要的部分。不要只看答案，要把每个限制条件变成一个计数动作：先安排谁、再插入谁、有没有重复、需不需要除以对称性。

\smallskip
\textbf{初学者检查点：}
\begin{itemize}[leftmargin=2em,itemsep=0.25em]
\item 先写出没有限制时的总数。
\item 逐条加入限制条件，观察总数怎样变化。
\item 最后检查有没有把同一种方案重复算了多次。
\end{itemize}

\smallskip
\textbf{本章读法提醒：} 第一章所有题都先判断“顺序是否重要”。这个判断错了，后面公式一定错。
\end{explainblock}


\begin{sourceblock}{定理 1.4.2 (递推关系) 设 n ≥ 2, D0 = 1, D1 = 0, 则}
{\small\sloppy
\noindent 定理 1.4.2 (递推关系) 设 n ≥ 2, D0 = 1, D1 = 0, 则\par
\noindent Dn = (n − 1)(Dn−1 + Dn−2 ). (1.24)\par
\noindent 证明: 当 n = 2 时, 有唯一的错排 ‘‘21”, 即 D2 = 1, 满足递推关系 (1.24). 对于\par
\noindent n ≥ 3, 集合 [n] 的任意一个错排 π, 由于 π(1) 有 n−1 种取值. 对于任意一种取值, 不\par
\noindent 妨设为 π(1) = j, 此时若 π(j) = 1, 则显然剩下的就是对 2, 3, . . . j − 1, j + 1, . . . , n 的\par
\noindent †\par
\noindent 尼古拉 · 伯努利 (Nikolaus Bernoulli, 1687–1759), 生于瑞士巴塞尔, 瑞士数学家.\par
\noindent ‡\par
\noindent 皮耶 · 黑蒙 · 德蒙马特 (Pierre Rémond de Montmort, 1678–1719), 法国数学家, 在概率论方面的\par
\noindent 研究颇有盛名.\par
\noindent 错排, 共有 Dn−2 种情形; 若此时 π(j) 6= 1, 因为下面是对 1, 2, · · · , j − 1, j + 1, · · · , n\par
\noindent 的错排, 要求 π(j) 6= 1, 所以该排列相当于对 n − 1 个元素的错排, 故有 Dn−1 种方\par
\noindent 法. 综上, 递推式 (1.24) 得证.\par
}
\end{sourceblock}


\begin{explainblock}
定理段不要只看结论，要看它把哪两个计数方式连在一起。组合学证明最常见的技巧是：左边按一种标准数，右边按另一种标准数，然后说明两边数的是同一批对象。

\smallskip
\textbf{初学者检查点：}
\begin{itemize}[leftmargin=2em,itemsep=0.25em]
\item 先复述定理的假设，不要把适用范围扩大。
\item 找出证明中的基本计数原则：乘法、加法、双射、递推、容斥或生成函数。
\item 检查最后的等式化简是否和前面的对象解释一致。
\end{itemize}

\smallskip
\textbf{本章读法提醒：} 第一章所有题都先判断“顺序是否重要”。这个判断错了，后面公式一定错。
\end{explainblock}


\begin{sourceblock}{定理 1.4.3 设 n 为非负整数, 则}
{\small\sloppy
\noindent 定理 1.4.3 设 n 为非负整数, 则\par
\noindent n\par
\noindent X n\par
\noindent Dn−k = n!. (1.25)\par
\noindent k\par
\noindent k=0\par
\noindent 证明: 设 π = π1 π2 · · · πn 集合 [n] 上的任意排列, 则全排列的个数共计有 n!\par
\noindent 个. 对 1 ≤ i ≤ n, 若 πi = i, 则称 i 是排列 π 的固定点. 下面对集合 [n] 的排列按固\par
\noindent 定点个数的进行分类. 因为以集合 [n] 中特定的 k (k ≤ n) 个元素为固定点的排列的\par
\noindent 个数显然是 Dn−k , 而从集合 [n] 中选定 k 个字母作为固定点的方法数有 nk , 所以\par
\noindent 集合 [n] 的排列中具有 k 个固定点的排列个数共有 nk Dn−k 个, 故 (1.25) 成立.\par
}
\end{sourceblock}


\begin{explainblock}
定理段不要只看结论，要看它把哪两个计数方式连在一起。组合学证明最常见的技巧是：左边按一种标准数，右边按另一种标准数，然后说明两边数的是同一批对象。

\smallskip
\textbf{初学者检查点：}
\begin{itemize}[leftmargin=2em,itemsep=0.25em]
\item 先复述定理的假设，不要把适用范围扩大。
\item 找出证明中的基本计数原则：乘法、加法、双射、递推、容斥或生成函数。
\item 检查最后的等式化简是否和前面的对象解释一致。
\end{itemize}

\smallskip
\textbf{本章读法提醒：} 第一章所有题都先判断“顺序是否重要”。这个判断错了，后面公式一定错。
\end{explainblock}


\begin{sourceblock}{定理 1.4.4 (递推关系§ ) 设 D0 = 1, D1 = 0, 则 n ≥ 2}
{\small\sloppy
\noindent 定理 1.4.4 (递推关系§ ) 设 D0 = 1, D1 = 0, 则 n ≥ 2\par
\noindent Dn = nDn−1 + (−1)n . (1.26)\par
\noindent 证明: 将 (1.24) 两边同时除以 n! 得:\par
\noindent Dn 1 Dn−1 1 Dn−2\par
\noindent = 1− + ,\par
\noindent n! n (n − 1)! n (n − 2)!\par
\noindent 即:\par
\noindent Dn Dn−1 1 Dn−1 Dn−2\par
\noindent − =− − . (1.27)\par
\noindent n! (n − 1)! n (n − 1)! (n − 2)!\par
\noindent 由 (1.27) 式可知:\par
\noindent Dn−1 Dn−2 1 Dn−2 Dn−3\par
\noindent − =− − ,\par
\noindent (n − 1)! (n − 2)! n−1 (n − 2)! (n − 3)!\par
\noindent Dn−2 Dn−3 1 Dn−3 Dn−4\par
\noindent − =− − ,\par
\noindent (n − 2)! (n − 3)! n−2 (n − 3)! (n − 4)!\par
\noindent .........\par
}
\end{sourceblock}


\begin{explainblock}
定理段不要只看结论，要看它把哪两个计数方式连在一起。组合学证明最常见的技巧是：左边按一种标准数，右边按另一种标准数，然后说明两边数的是同一批对象。

\smallskip
\textbf{初学者检查点：}
\begin{itemize}[leftmargin=2em,itemsep=0.25em]
\item 先复述定理的假设，不要把适用范围扩大。
\item 找出证明中的基本计数原则：乘法、加法、双射、递推、容斥或生成函数。
\item 检查最后的等式化简是否和前面的对象解释一致。
\end{itemize}

\smallskip
\textbf{本章读法提醒：} 第一章所有题都先判断“顺序是否重要”。这个判断错了，后面公式一定错。
\end{explainblock}


\begin{sourceblock}{第 1 章续读}
{\small\sloppy
\noindent D3 D2 1 D2 D1 1\par
\noindent − =− − =− .\par
\noindent 3! 2! 3 2! 1! 3!\par
\noindent 通过迭代可得:\par
\noindent Dn Dn−1 (−1)n−2 (−1)n\par
\noindent − = = , n≥2 (1.28)\par
\noindent n! (n − 1)! n! n!\par
\noindent 对上述等式两边同乘以 n! 立得(1.26).\par
\noindent §\par
\noindent 该递推公式由欧拉发现, 被欧拉称为 “ 漂亮的关系”.\par
}
\end{sourceblock}


\begin{explainblock}
这一段是原书论述链的一部分。读的时候把它当作桥梁：它通常负责把上一个定义、例子或定理，引到下一个计数方法。

\smallskip
\textbf{初学者检查点：}
\begin{itemize}[leftmargin=2em,itemsep=0.25em]
\item 找出这一段承接的上一个概念。
\item 找出它准备引出的下一个概念。
\item 把中间的关键等式或构造单独抄出来。
\end{itemize}

\smallskip
\textbf{本章读法提醒：} 第一章所有题都先判断“顺序是否重要”。这个判断错了，后面公式一定错。
\end{explainblock}


\begin{sourceblock}{定理 1.4.5 设 n ≥ 0, 则}
{\small\sloppy
\noindent 定理 1.4.5 设 n ≥ 0, 则\par
\noindent X\par
\noindent n\par
\noindent (−1)k\par
\noindent Dn = n! . (1.29)\par
\noindent k!\par
\noindent k=0\par
\noindent 证明: 方法一: 观察 (1.28) 式可得\par
\noindent Dk Dk−1 (−1)k−2 (−1)k\par
\noindent − = = ,\par
\noindent k! (k − 1)! k! k!\par
\noindent 对上式两边关于 k 求和可得:\par
\noindent X\par
\noindent n n\par
\noindent X\par
\noindent (−1)k Dk Dk−1\par
\noindent = − ,\par
\noindent k! k! (k − 1)!\par
}
\end{sourceblock}


\begin{explainblock}
定理段不要只看结论，要看它把哪两个计数方式连在一起。组合学证明最常见的技巧是：左边按一种标准数，右边按另一种标准数，然后说明两边数的是同一批对象。

\smallskip
\textbf{初学者检查点：}
\begin{itemize}[leftmargin=2em,itemsep=0.25em]
\item 先复述定理的假设，不要把适用范围扩大。
\item 找出证明中的基本计数原则：乘法、加法、双射、递推、容斥或生成函数。
\item 检查最后的等式化简是否和前面的对象解释一致。
\end{itemize}

\smallskip
\textbf{本章读法提醒：} 第一章所有题都先判断“顺序是否重要”。这个判断错了，后面公式一定错。
\end{explainblock}


\begin{sourceblock}{第 1 章续读}
{\small\sloppy
\noindent k=1 k=1\par
\noindent D1 D0 D2 D1 D3 D2\par
\noindent = − + − + −\par
\noindent 1! 0! 2! 1! 3! 2!\par
\noindent Dn−1 Dn−2 Dn Dn−1\par
\noindent + ··· + − + −\par
\noindent (n − 1)! (n − 2)! n! (n − 1)!\par
\noindent Dn\par
\noindent = − 1.\par
\noindent n!\par
\noindent 即\par
\noindent Dn X (−1)k X (−1)k\par
\noindent n n\par
\noindent = = ,\par
\noindent n! k! k!\par
\noindent k=2 k=0\par
\noindent 两边同时乘以 n! 立得 (1.29).\par
}
\end{sourceblock}


\begin{explainblock}
这一段是原书论述链的一部分。读的时候把它当作桥梁：它通常负责把上一个定义、例子或定理，引到下一个计数方法。

\smallskip
\textbf{初学者检查点：}
\begin{itemize}[leftmargin=2em,itemsep=0.25em]
\item 找出这一段承接的上一个概念。
\item 找出它准备引出的下一个概念。
\item 把中间的关键等式或构造单独抄出来。
\end{itemize}

\smallskip
\textbf{本章读法提醒：} 第一章所有题都先判断“顺序是否重要”。这个判断错了，后面公式一定错。
\end{explainblock}


\begin{sourceblock}{定理 1.4.5 也可用对合原理加以证明. 首先阐述对合原理的基本思想 [1, pp.202–}
{\small\sloppy
\noindent 定理 1.4.5 也可用对合原理加以证明. 首先阐述对合原理的基本思想 [1, pp.202–\par
\noindent 203]: 考虑有限集合 S 的 2-部划分: S = S + ∪ S − , 其中 S + 称为 S 的正部子集, S −\par
\noindent 称为 S 的负部子集, S 称为具有符号标记的集合. 设 f 是从 S 到它本身的对合. 若\par
\noindent f (x) = x 则称 x 为不动点. 除不动点外, x 与 f (x) 不同时在 S + 或 S − 中的对合,\par
\noindent 称为变号对合 (sign-reversing involution). 记 S + 的不动点的集合为 Fix+ , S − 的不\par
\noindent 动点的集合为 Fix− , 则显然有\par
\noindent |S + | − |S − | = |Fix+ | − |Fix− |. (1.30)\par
}
\end{sourceblock}


\begin{explainblock}
定理段不要只看结论，要看它把哪两个计数方式连在一起。组合学证明最常见的技巧是：左边按一种标准数，右边按另一种标准数，然后说明两边数的是同一批对象。

\smallskip
\textbf{初学者检查点：}
\begin{itemize}[leftmargin=2em,itemsep=0.25em]
\item 先复述定理的假设，不要把适用范围扩大。
\item 找出证明中的基本计数原则：乘法、加法、双射、递推、容斥或生成函数。
\item 检查最后的等式化简是否和前面的对象解释一致。
\end{itemize}

\smallskip
\textbf{本章读法提醒：} 第一章所有题都先判断“顺序是否重要”。这个判断错了，后面公式一定错。
\end{explainblock}


\begin{sourceblock}{定理 1.4.5 的证明: 对 [n] 中的所有排列进行下面规则的标记: 对任一个排列 π,}
{\small\sloppy
\noindent 定理 1.4.5 的证明: 对 [n] 中的所有排列进行下面规则的标记: 对任一个排列 π,\par
\noindent 若 i 不是固定点, 则不作任何标记; 若 i 是固定点, 则可将 i 标记为 + 或 −, 每个固\par
\noindent 定点只能标记一种符号.\par
\noindent 带标记的排列 π 中被标记为 “+” 的固定点的个数称为 π 的正指标, 记为 π + ; 被\par
\noindent 标记为 “-” 的固定点的个数称为 π 的负指标, 记为 π − .\par
\noindent 将标记后的所有排列的集合记为 P . 定义下面两个集合: Pe 表示有偶数个\par
\noindent 负指标的排列构成的集合, Po 表示有奇数个负指标的排列构成的集合. 考虑有\par
\noindent k(k = 0, 1, 2, . . . , n) 个负指标的排列的个数. 因为有 k 负指标, 所以这个 k 个点是固\par
\noindent 定点, 而从 n 个元素中选取 k 个元素作为固定点的方法数为 nk ; 选定这 k 个固定\par
\noindent 点后, 剩下元素排列的个数为 (n − k)!, 对这些排列, 我们只需对其固定点标记 “+”\par
\noindent 即可. 因此, 有 k 个负指标的排列的个数为:\par
\noindent n n!\par
\noindent (n − k)! = .\par
\noindent k k!\par
\noindent 记集合 S 的元素个数为 \#S, 则我们有\par
\noindent X\par
\noindent [n/2]\par
\noindent n! X\par
}
\end{sourceblock}


\begin{explainblock}
定理段不要只看结论，要看它把哪两个计数方式连在一起。组合学证明最常见的技巧是：左边按一种标准数，右边按另一种标准数，然后说明两边数的是同一批对象。

\smallskip
\textbf{初学者检查点：}
\begin{itemize}[leftmargin=2em,itemsep=0.25em]
\item 先复述定理的假设，不要把适用范围扩大。
\item 找出证明中的基本计数原则：乘法、加法、双射、递推、容斥或生成函数。
\item 检查最后的等式化简是否和前面的对象解释一致。
\end{itemize}

\smallskip
\textbf{本章读法提醒：} 第一章所有题都先判断“顺序是否重要”。这个判断错了，后面公式一定错。
\end{explainblock}


\begin{sourceblock}{第 1 章续读}
{\small\sloppy
\noindent [(n−1)/2]\par
\noindent n!\par
\noindent nX (−1)k\par
\noindent \#Pe − \#Po = − = n! .\par
\noindent (2k)! (2k + 1)! k!\par
\noindent k=0 k=0 k=0\par
\noindent 为了证明 (1.29) 成立, 只需证明\par
\noindent \#Pe − \#Po = Dn . (1.31)\par
\noindent 我们定义映射 Φ : P 7→ P . 它将 P 中没有固定点的排列映到它本身; 通过改变有\par
\noindent 固定点的排列中最小固定点的标记来将有偶 (奇) 数个负指标的排列映到有奇 (偶)\par
\noindent 数个负指标的排列. 具体而言, 设排列 π ∈ P , 若 π 中没有固定点则令 Φ(π) = π; 若\par
\noindent π 中有固定点, 且 i 是 π 中最小固定点. 则 Φ(π) 是改变 i 标记后的排列. 显然若 π\par
\noindent 没有固定点, 则 π 与 Φ(π) 负指标的个数奇偶性不同. 因此, 当 π 没有固定点时, π\par
\noindent 与 Φ(π) 不同时在集合 Pe 或 Po 中. 由 Φ 定义, 显然我们有 Φ(Φ(π)) = π, 即, Φ 是\par
\noindent 一个对合映射. 因此, 由对合原理可知 \#Pe − \#Po 事实上是在 Φ 作用下 Pe 中不动\par
\noindent 点的个数: 即没有固定点的排列的个数, 也即是错位排列的个数. 即 Φ 的所有不动点\par
\noindent 的集合为 Dn . 所以等式(1.31) 成立, 公式 (1.29) 得证.\par
}
\end{sourceblock}


\begin{explainblock}
这一段是原书论述链的一部分。读的时候把它当作桥梁：它通常负责把上一个定义、例子或定理，引到下一个计数方法。

\smallskip
\textbf{初学者检查点：}
\begin{itemize}[leftmargin=2em,itemsep=0.25em]
\item 找出这一段承接的上一个概念。
\item 找出它准备引出的下一个概念。
\item 把中间的关键等式或构造单独抄出来。
\end{itemize}

\smallskip
\textbf{本章读法提醒：} 第一章所有题都先判断“顺序是否重要”。这个判断错了，后面公式一定错。
\end{explainblock}


\begin{sourceblock}{注记 1.4.6 定理 1.4.5 可利用生成函数的方法或容斥原理加以证明, 详见例 2.3.1}
{\small\sloppy
\noindent 注记 1.4.6 定理 1.4.5 可利用生成函数的方法或容斥原理加以证明, 详见例 2.3.1\par
\noindent 与例 3.2.1.\par
}
\end{sourceblock}


\begin{explainblock}
注记通常是在补充术语、记号、历史或一个容易忽略的约定。它未必给新公式，但常常决定你后面会不会把符号读错。

\smallskip
\textbf{初学者检查点：}
\begin{itemize}[leftmargin=2em,itemsep=0.25em]
\item 记下新符号的读法。
\item 确认这个约定是否只在本章使用。
\item 把注记里的限制条件补到前面的定义旁边。
\end{itemize}

\smallskip
\textbf{本章读法提醒：} 第一章所有题都先判断“顺序是否重要”。这个判断错了，后面公式一定错。
\end{explainblock}


\begin{sourceblock}{例 1.4.2 若有 n 位客人来到家里做客, 每个客人都戴了一顶帽子, 来的时候他们}
{\small\sloppy
\noindent 例 1.4.2 若有 n 位客人来到家里做客, 每个客人都戴了一顶帽子, 来的时候他们\par
\noindent 把自己的帽子都挂在衣帽架上. 晚会结束后, 客人们取走帽子回家, 那么所有客人取\par
\noindent 走的不是自己帽子的概率是多少呢? 这个概率与客人的数量有什么关系呢?\par
\noindent 解: 由定理 1.1.2 知 n 个客人在晚会结束后, 取走帽子不同方法的数有 n! 种, 其中所\par
\noindent 有人取走的不是自己帽子的方法数为\par
\noindent X\par
\noindent n\par
\noindent (−1)k\par
\noindent Dn = n! .\par
\noindent k!\par
\noindent k=0\par
\noindent 设所有人取走的不是自己帽子的概率为 Pn , 则\par
\noindent Dn X (−1)k\par
\noindent n\par
\noindent Pn = = . (1.32)\par
\noindent n! k!\par
\noindent k=0\par
\noindent 注意到\par
}
\end{sourceblock}


\begin{explainblock}
例题是学习这本书最重要的部分。不要只看答案，要把每个限制条件变成一个计数动作：先安排谁、再插入谁、有没有重复、需不需要除以对称性。

\smallskip
\textbf{初学者检查点：}
\begin{itemize}[leftmargin=2em,itemsep=0.25em]
\item 先写出没有限制时的总数。
\item 逐条加入限制条件，观察总数怎样变化。
\item 最后检查有没有把同一种方案重复算了多次。
\end{itemize}

\smallskip
\textbf{本章读法提醒：} 第一章所有题都先判断“顺序是否重要”。这个判断错了，后面公式一定错。
\end{explainblock}


\begin{sourceblock}{第 1 章续读}
{\small\sloppy
\noindent ∞\par
\noindent X xk\par
\noindent ex = ,\par
\noindent k!\par
\noindent k=0\par
\noindent 则\par
\noindent ∞\par
\noindent X (−1)k 1\par
\noindent lim Pn = = , (1.33)\par
\noindent n→∞ k! e\par
\noindent k=0\par
\noindent 所以错排的概率随着 n 的增大逐渐趋于稳定, 且极限为 1e .\par
\noindent 自然常数 e 最早由雅各布 · 伯努利 ¶ 于 1683 年引入, 最早用 e 来表示自然底数\par
\noindent 的人是欧拉. 伯努利在 1683 年曾考虑银行利息的问题. 假设银行年利率为 100\%, 且\par
\noindent 可随时存随时取. 现在小 A 在银行存 1 元, 那么一年后银行就要给小 A 2 元; 若小\par
\noindent A 存 1 元, 半年后取出可得 1.5 元, 若将这 1.5 元当时就存入银行, 再过半年取出得\par
\noindent 100\% 100\%\par
\noindent 1∗ 1+ ∗ 1+ = 2.25\par
}
\end{sourceblock}


\begin{explainblock}
这一段是原书论述链的一部分。读的时候把它当作桥梁：它通常负责把上一个定义、例子或定理，引到下一个计数方法。

\smallskip
\textbf{初学者检查点：}
\begin{itemize}[leftmargin=2em,itemsep=0.25em]
\item 找出这一段承接的上一个概念。
\item 找出它准备引出的下一个概念。
\item 把中间的关键等式或构造单独抄出来。
\end{itemize}

\smallskip
\textbf{本章读法提醒：} 第一章所有题都先判断“顺序是否重要”。这个判断错了，后面公式一定错。
\end{explainblock}


\begin{sourceblock}{第 1 章续读}
{\small\sloppy
\noindent 元, 比存一整年多了 0.25 元; 当小 A 每四个月将本金和利息取出再存入, 这样一年\par
\noindent 后得\par
\noindent 100\% 100\% 100\%\par
\noindent 1∗ 1+ 1+ 1+ ≃ 2.37037\par
\noindent 元, 又多出一些; 按这样下去, 如果小 A 在这一年不停的在银行取出本金和利息再存\par
\noindent 进去, 那么年底小 A 将得到\par
\noindent 1+\par
\noindent n\par
\noindent 元, 看上去它会随 n 的增长而一直增长, 那么最终会是无穷大么? 当然不会, 也即是\par
\noindent 说极限\par
\noindent lim 1+\par
\noindent n→∞ n\par
\noindent 是存在的, 是一个常数.\par
\noindent 下面我们不妨再看看欧拉是怎样定义 e 的. 设 a > 1, 考虑指数函数 aω . 当 ω 趋\par
\noindent 于 0 时, 由无穷小理论可知:\par
\noindent aω = 1 + o(ω)\par
\noindent 不妨设 aω = 1 + kω, k 是一个与 a 有关的常数. 令 j = ωx , 则\par
\noindent x ω j kx j j\par
}
\end{sourceblock}


\begin{explainblock}
这一段是原书论述链的一部分。读的时候把它当作桥梁：它通常负责把上一个定义、例子或定理，引到下一个计数方法。

\smallskip
\textbf{初学者检查点：}
\begin{itemize}[leftmargin=2em,itemsep=0.25em]
\item 找出这一段承接的上一个概念。
\item 找出它准备引出的下一个概念。
\item 把中间的关键等式或构造单独抄出来。
\end{itemize}

\smallskip
\textbf{本章读法提醒：} 第一章所有题都先判断“顺序是否重要”。这个判断错了，后面公式一定错。
\end{explainblock}


\begin{sourceblock}{第 1 章续读}
{\small\sloppy
\noindent a = (a ) = (1 + kω) = 1 + . (1.34)\par
\noindent j\par
\noindent kx\par
\noindent 因为当 ω 趋于 0 时, j < 1 , 故等式右边可通过牛顿二项式展开, 得到:\par
\noindent j ∞\par
\noindent X ∞\par
\noindent X\par
\noindent kx (j)n k n xn (j)n k n xn\par
\noindent 1+ = = . (1.35)\par
\noindent j n! jn jn n!\par
\noindent n=0 n=0\par
\noindent 因为 ω 趋于 0 时, j 趋向 ∞, 所以\par
\noindent (j)n\par
\noindent lim = 1. (1.36)\par
\noindent j→∞ j n\par
\noindent ¶\par
\noindent 雅各布 · 伯努利 (Jakob Bernoulli, 1654–1705) 瑞士数学家.\par
\noindent P∞ (j)n kn xn\par
}
\end{sourceblock}


\begin{explainblock}
这一段是原书论述链的一部分。读的时候把它当作桥梁：它通常负责把上一个定义、例子或定理，引到下一个计数方法。

\smallskip
\textbf{初学者检查点：}
\begin{itemize}[leftmargin=2em,itemsep=0.25em]
\item 找出这一段承接的上一个概念。
\item 找出它准备引出的下一个概念。
\item 把中间的关键等式或构造单独抄出来。
\end{itemize}

\smallskip
\textbf{本章读法提醒：} 第一章所有题都先判断“顺序是否重要”。这个判断错了，后面公式一定错。
\end{explainblock}


\begin{sourceblock}{第 1 章续读}
{\small\sloppy
\noindent 又因为级数 n=0 j n n! 一致收敛, 求和和极限可以交换位置, 故而有\par
\noindent ∞\par
\noindent x kx j X k n xn\par
\noindent a = lim 1 + = . (1.37)\par
\noindent j→∞ j n!\par
\noindent n=0\par
\noindent 令 x = 1, 欧拉 [25] 得到\par
\noindent k2 k3\par
\noindent a=1+k+ + + ··· , (1.38)\par
\noindent 2! 3!\par
\noindent 随后欧拉令 k = 1, 并将得到的数定义为 e:\par
\noindent X 1 ∞\par
\noindent e=1+1+ + + ··· = .\par
\noindent 2! 3! n!\par
\noindent n=0\par
\noindent 下面我们将说明\par
\noindent ∞\par
\noindent lim 1 + = , (1.39)\par
}
\end{sourceblock}


\begin{explainblock}
这一段是原书论述链的一部分。读的时候把它当作桥梁：它通常负责把上一个定义、例子或定理，引到下一个计数方法。

\smallskip
\textbf{初学者检查点：}
\begin{itemize}[leftmargin=2em,itemsep=0.25em]
\item 找出这一段承接的上一个概念。
\item 找出它准备引出的下一个概念。
\item 把中间的关键等式或构造单独抄出来。
\end{itemize}

\smallskip
\textbf{本章读法提醒：} 第一章所有题都先判断“顺序是否重要”。这个判断错了，后面公式一定错。
\end{explainblock}


\begin{sourceblock}{第 1 章续读}
{\small\sloppy
\noindent n→∞ n n!\par
\noindent n=0\par
\noindent 也即是说伯努利与欧拉得到的常数是一样的, 都是我们现在称作自然常数的 e. 为\par
\noindent 了说明这个事实, 我们需要引入 Tannery 定理, 它是以 Jules Tannery‖ 的姓命名.\par
\noindent Tannery 定理事实上是 Weierstrass M -判别法的特殊情形.\par
\noindent P n\par
}
\end{sourceblock}


\begin{explainblock}
这一段是原书论述链的一部分。读的时候把它当作桥梁：它通常负责把上一个定义、例子或定理，引到下一个计数方法。

\smallskip
\textbf{初学者检查点：}
\begin{itemize}[leftmargin=2em,itemsep=0.25em]
\item 找出这一段承接的上一个概念。
\item 找出它准备引出的下一个概念。
\item 把中间的关键等式或构造单独抄出来。
\end{itemize}

\smallskip
\textbf{本章读法提醒：} 第一章所有题都先判断“顺序是否重要”。这个判断错了，后面公式一定错。
\end{explainblock}


\begin{sourceblock}{定理 1.4.7 (Tannery) 设 n 为自然数, m k=1 ak (n) 是一个有限和式且当 n → ∞}
{\small\sloppy
\noindent 定理 1.4.7 (Tannery) 设 n 为自然数, m k=1 ak (n) 是一个有限和式且当 n → ∞\par
\noindent P\par
\noindent 时, mn → ∞. 若对任意 k, 极限 limn→∞ ak (n) 都存在, 且存在收敛的级数 ∞ k Mk\par
\noindent 使得\par
\noindent |ak (n)| ≤ Mk\par
\noindent 对所有自然数 n 及 1 ≤ k ≤ mn 都成立, 则\par
\noindent X\par
\noindent mn ∞\par
\noindent X\par
\noindent lim ak (n) = lim ak (n). (1.40)\par
\noindent n→∞ n→∞\par
\noindent k=1 k=1\par
\noindent 证明: 不妨设\par
\noindent lim ak (n) = ak .\par
\noindent n→∞\par
\noindent 则由\par
\noindent |ak (n)| ≤ Mk\par
\noindent 及极限的保号性可得\par
}
\end{sourceblock}


\begin{explainblock}
定理段不要只看结论，要看它把哪两个计数方式连在一起。组合学证明最常见的技巧是：左边按一种标准数，右边按另一种标准数，然后说明两边数的是同一批对象。

\smallskip
\textbf{初学者检查点：}
\begin{itemize}[leftmargin=2em,itemsep=0.25em]
\item 先复述定理的假设，不要把适用范围扩大。
\item 找出证明中的基本计数原则：乘法、加法、双射、递推、容斥或生成函数。
\item 检查最后的等式化简是否和前面的对象解释一致。
\end{itemize}

\smallskip
\textbf{本章读法提醒：} 第一章所有题都先判断“顺序是否重要”。这个判断错了，后面公式一定错。
\end{explainblock}


\begin{sourceblock}{第 1 章续读}
{\small\sloppy
\noindent |ak | ≤ Mk .\par
\noindent 因此, 由比较判别法可知等式右边的级数收敛.\par
\noindent ‖\par
\noindent Jules Tannery(1848–1910), 法国数学家\par
\noindent 由 Cauchy 收敛准则可知, 对任意 ε > 0, 存在 ℓ 使得\par
\noindent ε\par
\noindent Mℓ+1 + Mℓ+2 + · · · < .\par
\noindent 又因为当 n → ∞ 时, mn → ∞, 因此, 存在 N1 , 当 n > N1 时, 我们有\par
\noindent mn > ℓ.\par
\noindent 当 n > N1 , 由\par
\noindent |ak (n)| ≤ Mk , |ak | ≤ Mk ,\par
\noindent 我们可得\par
\noindent X\par
\noindent mn ∞\par
\noindent X X\par
\noindent ℓ X\par
\noindent mn ∞\par
\noindent X\par
}
\end{sourceblock}


\begin{explainblock}
这一段是原书论述链的一部分。读的时候把它当作桥梁：它通常负责把上一个定义、例子或定理，引到下一个计数方法。

\smallskip
\textbf{初学者检查点：}
\begin{itemize}[leftmargin=2em,itemsep=0.25em]
\item 找出这一段承接的上一个概念。
\item 找出它准备引出的下一个概念。
\item 把中间的关键等式或构造单独抄出来。
\end{itemize}

\smallskip
\textbf{本章读法提醒：} 第一章所有题都先判断“顺序是否重要”。这个判断错了，后面公式一定错。
\end{explainblock}


\begin{sourceblock}{第 1 章续读}
{\small\sloppy
\noindent ak (n) − ak = (ak (n) − ak ) + ak (n) −\par
\noindent k=1 k=1 k=1 ℓ+1 k=ℓ+1\par
\noindent X\par
\noindent ℓ X\par
\noindent mn ∞\par
\noindent X\par
\noindent ≤ |ak (n) − ak | + + Mk\par
\noindent k=1 k=ℓ+1 k=ℓ+1\par
\noindent X\par
\noindent ℓ\par
\noindent ε ε\par
\noindent ≤ |ak (n) − ak | + + . (1.41)\par
\noindent k=1\par
\noindent 又对于每一个 k, 我们有\par
\noindent lim ak (n) = ak ,\par
\noindent n→∞\par
\noindent 因此, 存在 N2 , 使得当 n > N2 时, 对于 k = 1, 2, . . . , ℓ 都有\par
\noindent ε\par
}
\end{sourceblock}


\begin{explainblock}
这一段是原书论述链的一部分。读的时候把它当作桥梁：它通常负责把上一个定义、例子或定理，引到下一个计数方法。

\smallskip
\textbf{初学者检查点：}
\begin{itemize}[leftmargin=2em,itemsep=0.25em]
\item 找出这一段承接的上一个概念。
\item 找出它准备引出的下一个概念。
\item 把中间的关键等式或构造单独抄出来。
\end{itemize}

\smallskip
\textbf{本章读法提醒：} 第一章所有题都先判断“顺序是否重要”。这个判断错了，后面公式一定错。
\end{explainblock}


\begin{sourceblock}{第 1 章续读}
{\small\sloppy
\noindent |ak (n) − ak | < .\par
\noindent 3ℓ\par
\noindent 故\par
\noindent X\par
\noindent ℓ\par
\noindent ε\par
\noindent |ak (n) − ak | < .\par
\noindent k=1\par
\noindent 代入 (8.5) 即可.\par
\noindent 特别地, 取 mn = n 时, 我们有\par
\noindent X\par
\noindent n ∞\par
\noindent X\par
\noindent lim ak (n) = lim ak (n). (1.42)\par
\noindent n→∞ n→∞\par
\noindent k=1 k=1\par
\noindent 取 mn 为 +∞ 时, 我们可证明同样有\par
\noindent ∞\par
}
\end{sourceblock}


\begin{explainblock}
这一段是原书论述链的一部分。读的时候把它当作桥梁：它通常负责把上一个定义、例子或定理，引到下一个计数方法。

\smallskip
\textbf{初学者检查点：}
\begin{itemize}[leftmargin=2em,itemsep=0.25em]
\item 找出这一段承接的上一个概念。
\item 找出它准备引出的下一个概念。
\item 把中间的关键等式或构造单独抄出来。
\end{itemize}

\smallskip
\textbf{本章读法提醒：} 第一章所有题都先判断“顺序是否重要”。这个判断错了，后面公式一定错。
\end{explainblock}


\begin{sourceblock}{第 1 章续读}
{\small\sloppy
\noindent X ∞\par
\noindent X\par
\noindent lim ak (n) = lim ak (n). (1.43)\par
\noindent n→∞ n→∞\par
\noindent k=1 k=1\par
\noindent 这实际上表明了极限与无穷求和符号可交换.\par
\noindent 下面我们就证明:\par
\noindent n ∞\par
\noindent X\par
\noindent lim 1+ = .\par
\noindent n→∞ n k!\par
\noindent k=0\par
\noindent 由二项式定理得:\par
\noindent n\par
\noindent n\par
\noindent 1+ = := ak (n),\par
\noindent n k nk\par
\noindent k=0 k=0\par
}
\end{sourceblock}


\begin{explainblock}
这一段是原书论述链的一部分。读的时候把它当作桥梁：它通常负责把上一个定义、例子或定理，引到下一个计数方法。

\smallskip
\textbf{初学者检查点：}
\begin{itemize}[leftmargin=2em,itemsep=0.25em]
\item 找出这一段承接的上一个概念。
\item 找出它准备引出的下一个概念。
\item 把中间的关键等式或构造单独抄出来。
\end{itemize}

\smallskip
\textbf{本章读法提醒：} 第一章所有题都先判断“顺序是否重要”。这个判断错了，后面公式一定错。
\end{explainblock}


\begin{sourceblock}{第 1 章续读}
{\small\sloppy
\noindent 其中\par
\noindent n 1\par
\noindent ak (n) = .\par
\noindent k nk\par
\noindent 下面我们只需证明\par
\noindent X\par
\noindent n ∞\par
\noindent X 1\par
\noindent lim ak (n) = .\par
\noindent n→∞ k!\par
\noindent k=0 k=0\par
\noindent 因为 a0 = a1 , 上式等价于\par
\noindent X\par
\noindent n ∞\par
\noindent X 1\par
\noindent lim ak (n) = .\par
\noindent n→∞ k!\par
\noindent k=2 k=2\par
}
\end{sourceblock}


\begin{explainblock}
这一段是原书论述链的一部分。读的时候把它当作桥梁：它通常负责把上一个定义、例子或定理，引到下一个计数方法。

\smallskip
\textbf{初学者检查点：}
\begin{itemize}[leftmargin=2em,itemsep=0.25em]
\item 找出这一段承接的上一个概念。
\item 找出它准备引出的下一个概念。
\item 把中间的关键等式或构造单独抄出来。
\end{itemize}

\smallskip
\textbf{本章读法提醒：} 第一章所有题都先判断“顺序是否重要”。这个判断错了，后面公式一定错。
\end{explainblock}


\begin{sourceblock}{第 1 章续读}
{\small\sloppy
\noindent 当 2 ≤ k ≤ n 且 k 固定时, 我们有\par
\noindent n 1 n! 1 1 2 k−1\par
\noindent ak (n) = = = 1− 1− · 1− ,\par
\noindent k nk k!(n − k)!nk k! n n n\par
\noindent 所以\par
\noindent lim ak (n) = .\par
\noindent n→∞ k!\par
\noindent 所以对于 2 ≤ k ≤ n,\par
\noindent |ak (n)| = 1− 1− · 1− ≤ := Mk .\par
\noindent k! n n n k!\par
\noindent P∞\par
\noindent 而级数 k=2 Mk 是收敛的, 所以由 Tannery 定理\par
\noindent X\par
\noindent n X\par
\noindent n ∞\par
\noindent X 1\par
\noindent lim ak (n) = lim ak (n) = .\par
\noindent n→∞ n→∞ k!\par
}
\end{sourceblock}


\begin{explainblock}
这一段是原书论述链的一部分。读的时候把它当作桥梁：它通常负责把上一个定义、例子或定理，引到下一个计数方法。

\smallskip
\textbf{初学者检查点：}
\begin{itemize}[leftmargin=2em,itemsep=0.25em]
\item 找出这一段承接的上一个概念。
\item 找出它准备引出的下一个概念。
\item 把中间的关键等式或构造单独抄出来。
\end{itemize}

\smallskip
\textbf{本章读法提醒：} 第一章所有题都先判断“顺序是否重要”。这个判断错了，后面公式一定错。
\end{explainblock}


\begin{sourceblock}{第 1 章续读}
{\small\sloppy
\noindent k=2 k=2 k=2\par
\noindent 因此,\par
\noindent n n\par
\noindent X ∞\par
\noindent X\par
\noindent lim 1+ = lim k\par
\noindent = .\par
\noindent n→∞ n n→∞ k n k!\par
\noindent k=0 k=0\par
\noindent 欧拉在 1737 年通过将 e 表示成连分数的形式证明了自然常数 e 是一个无理数.\par
\noindent 这里给出傅里叶 ∗∗ 在 1790–1800 年间关于 e 的无理性初等的证明. 首先考虑不等式：\par
\noindent 2<e=1+ + + · · · < 1 + 1 + + 2 + · · · = 3,\par
\noindent 1! 2! 2 2\par
\noindent ∗∗\par
\noindent 让 · 巴普蒂斯 · 约瑟夫 · 傅里叶 (Jean Baptiste Joseph Fourier, 1768–1830), 法国数学家.\par
\noindent 这表明 e 不是一个整数. 现在采取反证法. 设 e = pq , 其中 p, q 均为整数. 因为 e 不\par
\noindent 是整数, 所以 q ≥ 2. 由假设知 eq! = p(q − 1)! 应为整数, 又因为\par
\noindent X\par
}
\end{sourceblock}


\begin{explainblock}
这一段是原书论述链的一部分。读的时候把它当作桥梁：它通常负责把上一个定义、例子或定理，引到下一个计数方法。

\smallskip
\textbf{初学者检查点：}
\begin{itemize}[leftmargin=2em,itemsep=0.25em]
\item 找出这一段承接的上一个概念。
\item 找出它准备引出的下一个概念。
\item 把中间的关键等式或构造单独抄出来。
\end{itemize}

\smallskip
\textbf{本章读法提醒：} 第一章所有题都先判断“顺序是否重要”。这个判断错了，后面公式一定错。
\end{explainblock}


\begin{sourceblock}{第 1 章续读}
{\small\sloppy
\noindent q ∞\par
\noindent X\par
\noindent q! q!\par
\noindent eq! = + ,\par
\noindent n! n!\par
\noindent n=0 n=q+1\par
\noindent Pq q!\par
\noindent 其中 n=0 n! 显然是整数. 因为 q ≥ 2, 所以有 q + 1 ≥ 3. 故而有\par
\noindent ∞\par
\noindent X q! 1 1 1 1 1\par
\noindent 0< = + + ··· < + 2 + ··· = ,\par
\noindent n! q + 1 (q + 1)(q + 2) 3 3 2\par
\noindent n=q+1\par
\noindent 不是一个整数, 这就得出了矛盾. 因此 e 是一个无理数.\par
}
\end{sourceblock}


\begin{explainblock}
这一段是原书论述链的一部分。读的时候把它当作桥梁：它通常负责把上一个定义、例子或定理，引到下一个计数方法。

\smallskip
\textbf{初学者检查点：}
\begin{itemize}[leftmargin=2em,itemsep=0.25em]
\item 找出这一段承接的上一个概念。
\item 找出它准备引出的下一个概念。
\item 把中间的关键等式或构造单独抄出来。
\end{itemize}

\smallskip
\textbf{本章读法提醒：} 第一章所有题都先判断“顺序是否重要”。这个判断错了，后面公式一定错。
\end{explainblock}


\begin{sourceblock}{习 题 一}
{\small\sloppy
\noindent 习 题 一\par
}
\end{sourceblock}


\begin{explainblock}
习题段不要先追求技巧。先给题目分类：它是在问排列、组合、递推、生成函数、容斥，还是结构计数。分类正确，工具通常就只剩一两个候选。

\smallskip
\textbf{初学者检查点：}
\begin{itemize}[leftmargin=2em,itemsep=0.25em]
\item 判断题目所属工具。
\item 列出变量和限制条件。
\item 先做小规模 n=2、3、4 的例子，确认直觉。
\end{itemize}

\smallskip
\textbf{本章读法提醒：} 第一章所有题都先判断“顺序是否重要”。这个判断错了，后面公式一定错。
\end{explainblock}


\begin{sourceblock}{习题 1.1 已知一副纸牌共 54 张,}
{\small\sloppy
\noindent 习题 1.1 已知一副纸牌共 54 张,\par
\noindent (1) 洗一副纸牌可以有多少种不同的结果?\par
\noindent (2) 发 5 张纸牌可以有多少种不同的结果?\par
}
\end{sourceblock}


\begin{explainblock}
习题段不要先追求技巧。先给题目分类：它是在问排列、组合、递推、生成函数、容斥，还是结构计数。分类正确，工具通常就只剩一两个候选。

\smallskip
\textbf{初学者检查点：}
\begin{itemize}[leftmargin=2em,itemsep=0.25em]
\item 判断题目所属工具。
\item 列出变量和限制条件。
\item 先做小规模 n=2、3、4 的例子，确认直觉。
\end{itemize}

\smallskip
\textbf{本章读法提醒：} 第一章所有题都先判断“顺序是否重要”。这个判断错了，后面公式一定错。
\end{explainblock}


\begin{sourceblock}{习题 1.2 一个俱乐部有 10 名男成员和 12 名女成员, 从中选出 4 人组成委员}
{\small\sloppy
\noindent 习题 1.2 一个俱乐部有 10 名男成员和 12 名女成员, 从中选出 4 人组成委员\par
\noindent 会, 若要求:\par
\noindent (1) 至少有 2 名男成员;\par
\noindent (2) 除要求满足条件外还要求男成员 A 和女成员 B 不能同时入选,\par
\noindent 分别求出有多少种不同的选法.\par
}
\end{sourceblock}


\begin{explainblock}
习题段不要先追求技巧。先给题目分类：它是在问排列、组合、递推、生成函数、容斥，还是结构计数。分类正确，工具通常就只剩一两个候选。

\smallskip
\textbf{初学者检查点：}
\begin{itemize}[leftmargin=2em,itemsep=0.25em]
\item 判断题目所属工具。
\item 列出变量和限制条件。
\item 先做小规模 n=2、3、4 的例子，确认直觉。
\end{itemize}

\smallskip
\textbf{本章读法提醒：} 第一章所有题都先判断“顺序是否重要”。这个判断错了，后面公式一定错。
\end{explainblock}


\begin{sourceblock}{习题 1.3 有 6 位先生和 6 位女士圆桌就桌, 男女相间, 问有多少种不同的就坐}
{\small\sloppy
\noindent 习题 1.3 有 6 位先生和 6 位女士圆桌就桌, 男女相间, 问有多少种不同的就坐\par
\noindent 方式?\par
}
\end{sourceblock}


\begin{explainblock}
习题段不要先追求技巧。先给题目分类：它是在问排列、组合、递推、生成函数、容斥，还是结构计数。分类正确，工具通常就只剩一两个候选。

\smallskip
\textbf{初学者检查点：}
\begin{itemize}[leftmargin=2em,itemsep=0.25em]
\item 判断题目所属工具。
\item 列出变量和限制条件。
\item 先做小规模 n=2、3、4 的例子，确认直觉。
\end{itemize}

\smallskip
\textbf{本章读法提醒：} 第一章所有题都先判断“顺序是否重要”。这个判断错了，后面公式一定错。
\end{explainblock}


\begin{sourceblock}{习题 1.4 求由 n(n ≥ 2) 个不同的元素 a1 , a2 , . . . , an 组成且 a1 不在第 1 个位}
{\small\sloppy
\noindent 习题 1.4 求由 n(n ≥ 2) 个不同的元素 a1 , a2 , . . . , an 组成且 a1 不在第 1 个位\par
\noindent 置, a2 不在第 2 个位置的全排列的个数.\par
}
\end{sourceblock}


\begin{explainblock}
习题段不要先追求技巧。先给题目分类：它是在问排列、组合、递推、生成函数、容斥，还是结构计数。分类正确，工具通常就只剩一两个候选。

\smallskip
\textbf{初学者检查点：}
\begin{itemize}[leftmargin=2em,itemsep=0.25em]
\item 判断题目所属工具。
\item 列出变量和限制条件。
\item 先做小规模 n=2、3、4 的例子，确认直觉。
\end{itemize}

\smallskip
\textbf{本章读法提醒：} 第一章所有题都先判断“顺序是否重要”。这个判断错了，后面公式一定错。
\end{explainblock}


\begin{sourceblock}{习题 1.5 求重集 S = \{1·a, 4·b, 3·c\} 的 6-排列的个数. (注: 重集 \{k1 ·a, k2 ·b, . . .\}}
{\small\sloppy
\noindent 习题 1.5 求重集 S = \{1·a, 4·b, 3·c\} 的 6-排列的个数. (注: 重集 \{k1 ·a, k2 ·b, . . .\}\par
\noindent 是指有 k1 个 a, k2 个 b 等等).\par
}
\end{sourceblock}


\begin{explainblock}
习题段不要先追求技巧。先给题目分类：它是在问排列、组合、递推、生成函数、容斥，还是结构计数。分类正确，工具通常就只剩一两个候选。

\smallskip
\textbf{初学者检查点：}
\begin{itemize}[leftmargin=2em,itemsep=0.25em]
\item 判断题目所属工具。
\item 列出变量和限制条件。
\item 先做小规模 n=2、3、4 的例子，确认直觉。
\end{itemize}

\smallskip
\textbf{本章读法提醒：} 第一章所有题都先判断“顺序是否重要”。这个判断错了，后面公式一定错。
\end{explainblock}


\begin{sourceblock}{习题 1.6 证明:}
{\small\sloppy
\noindent 习题 1.6 证明:\par
\noindent (1)\par
\noindent n n n−1\par
\noindent = .\par
\noindent k n−k k\par
\noindent (2)\par
\noindent n 2n 2n\par
\noindent = .\par
\noindent n+1 n n−1\par
}
\end{sourceblock}


\begin{explainblock}
习题段不要先追求技巧。先给题目分类：它是在问排列、组合、递推、生成函数、容斥，还是结构计数。分类正确，工具通常就只剩一两个候选。

\smallskip
\textbf{初学者检查点：}
\begin{itemize}[leftmargin=2em,itemsep=0.25em]
\item 判断题目所属工具。
\item 列出变量和限制条件。
\item 先做小规模 n=2、3、4 的例子，确认直觉。
\end{itemize}

\smallskip
\textbf{本章读法提醒：} 第一章所有题都先判断“顺序是否重要”。这个判断错了，后面公式一定错。
\end{explainblock}


\begin{sourceblock}{习题 1.7 求方程 x1 + x2 + · · · + xn = k 的正整数解的个数.}
{\small\sloppy
\noindent 习题 1.7 求方程 x1 + x2 + · · · + xn = k 的正整数解的个数.\par
}
\end{sourceblock}


\begin{explainblock}
习题段不要先追求技巧。先给题目分类：它是在问排列、组合、递推、生成函数、容斥，还是结构计数。分类正确，工具通常就只剩一两个候选。

\smallskip
\textbf{初学者检查点：}
\begin{itemize}[leftmargin=2em,itemsep=0.25em]
\item 判断题目所属工具。
\item 列出变量和限制条件。
\item 先做小规模 n=2、3、4 的例子，确认直觉。
\end{itemize}

\smallskip
\textbf{本章读法提醒：} 第一章所有题都先判断“顺序是否重要”。这个判断错了，后面公式一定错。
\end{explainblock}


\begin{sourceblock}{习题 1.8 设 n ≥ 1, r ≥ 1, 证明}
{\small\sloppy
\noindent 习题 1.8 设 n ≥ 1, r ≥ 1, 证明\par
\noindent r\par
\noindent X\par
\noindent n+k−1 n+r\par
\noindent 1+ = .\par
\noindent k n\par
\noindent k=1\par
}
\end{sourceblock}


\begin{explainblock}
习题段不要先追求技巧。先给题目分类：它是在问排列、组合、递推、生成函数、容斥，还是结构计数。分类正确，工具通常就只剩一两个候选。

\smallskip
\textbf{初学者检查点：}
\begin{itemize}[leftmargin=2em,itemsep=0.25em]
\item 判断题目所属工具。
\item 列出变量和限制条件。
\item 先做小规模 n=2、3、4 的例子，确认直觉。
\end{itemize}

\smallskip
\textbf{本章读法提醒：} 第一章所有题都先判断“顺序是否重要”。这个判断错了，后面公式一定错。
\end{explainblock}


\begin{sourceblock}{习题 1.9 求}
{\small\sloppy
\noindent 习题 1.9 求\par
\noindent X n\par
\noindent .\par
\noindent 3k\par
\noindent k≥0\par
}
\end{sourceblock}


\begin{explainblock}
习题段不要先追求技巧。先给题目分类：它是在问排列、组合、递推、生成函数、容斥，还是结构计数。分类正确，工具通常就只剩一两个候选。

\smallskip
\textbf{初学者检查点：}
\begin{itemize}[leftmargin=2em,itemsep=0.25em]
\item 判断题目所属工具。
\item 列出变量和限制条件。
\item 先做小规模 n=2、3、4 的例子，确认直觉。
\end{itemize}

\smallskip
\textbf{本章读法提醒：} 第一章所有题都先判断“顺序是否重要”。这个判断错了，后面公式一定错。
\end{explainblock}


\begin{sourceblock}{习题 1.10 从 A = \{1, 2, . . . , n\} 中选取 k 个数, 使之没有两个数相邻, 求不同}
{\small\sloppy
\noindent 习题 1.10 从 A = \{1, 2, . . . , n\} 中选取 k 个数, 使之没有两个数相邻, 求不同\par
\noindent 的选法数.\par
}
\end{sourceblock}


\begin{explainblock}
习题段不要先追求技巧。先给题目分类：它是在问排列、组合、递推、生成函数、容斥，还是结构计数。分类正确，工具通常就只剩一两个候选。

\smallskip
\textbf{初学者检查点：}
\begin{itemize}[leftmargin=2em,itemsep=0.25em]
\item 判断题目所属工具。
\item 列出变量和限制条件。
\item 先做小规模 n=2、3、4 的例子，确认直觉。
\end{itemize}

\smallskip
\textbf{本章读法提醒：} 第一章所有题都先判断“顺序是否重要”。这个判断错了，后面公式一定错。
\end{explainblock}


\begin{sourceblock}{习题 1.11 设 r 为实数, k 与 n 为整数, 证明:}
{\small\sloppy
\noindent 习题 1.11 设 r 为实数, k 与 n 为整数, 证明:\par
\noindent r n r r−k\par
\noindent = .\par
\noindent n k k n−k\par
}
\end{sourceblock}


\begin{explainblock}
习题段不要先追求技巧。先给题目分类：它是在问排列、组合、递推、生成函数、容斥，还是结构计数。分类正确，工具通常就只剩一两个候选。

\smallskip
\textbf{初学者检查点：}
\begin{itemize}[leftmargin=2em,itemsep=0.25em]
\item 判断题目所属工具。
\item 列出变量和限制条件。
\item 先做小规模 n=2、3、4 的例子，确认直觉。
\end{itemize}

\smallskip
\textbf{本章读法提醒：} 第一章所有题都先判断“顺序是否重要”。这个判断错了，后面公式一定错。
\end{explainblock}


\begin{sourceblock}{习题 1.12 设 1 ≤ k ≤ n, 用组合方法证明 (1.17), 即:}
{\small\sloppy
\noindent 习题 1.12 设 1 ≤ k ≤ n, 用组合方法证明 (1.17), 即:\par
\noindent X\par
\noindent n\par
\noindent n\par
\noindent k = n2n−1 .\par
\noindent k\par
\noindent k=1\par
}
\end{sourceblock}


\begin{explainblock}
习题段不要先追求技巧。先给题目分类：它是在问排列、组合、递推、生成函数、容斥，还是结构计数。分类正确，工具通常就只剩一两个候选。

\smallskip
\textbf{初学者检查点：}
\begin{itemize}[leftmargin=2em,itemsep=0.25em]
\item 判断题目所属工具。
\item 列出变量和限制条件。
\item 先做小规模 n=2、3、4 的例子，确认直觉。
\end{itemize}

\smallskip
\textbf{本章读法提醒：} 第一章所有题都先判断“顺序是否重要”。这个判断错了，后面公式一定错。
\end{explainblock}


\begin{sourceblock}{习题 1.13 设 n 为大于等于 2 的正整数. 证明:}
{\small\sloppy
\noindent 习题 1.13 设 n 为大于等于 2 的正整数. 证明:\par
\noindent X\par
\noindent n\par
\noindent k = n(n + 1)2n−2 .\par
\noindent k\par
\noindent k=0\par
}
\end{sourceblock}


\begin{explainblock}
习题段不要先追求技巧。先给题目分类：它是在问排列、组合、递推、生成函数、容斥，还是结构计数。分类正确，工具通常就只剩一两个候选。

\smallskip
\textbf{初学者检查点：}
\begin{itemize}[leftmargin=2em,itemsep=0.25em]
\item 判断题目所属工具。
\item 列出变量和限制条件。
\item 先做小规模 n=2、3、4 的例子，确认直觉。
\end{itemize}

\smallskip
\textbf{本章读法提醒：} 第一章所有题都先判断“顺序是否重要”。这个判断错了，后面公式一定错。
\end{explainblock}


\begin{sourceblock}{习题 1.14 设 k, n 为非负整数, 证明}
{\small\sloppy
\noindent 习题 1.14 设 k, n 为非负整数, 证明\par
\noindent n\par
\noindent X\par
\noindent i n+1\par
\noindent = . (1.44)\par
\noindent k k+1\par
\noindent i=0\par
}
\end{sourceblock}


\begin{explainblock}
习题段不要先追求技巧。先给题目分类：它是在问排列、组合、递推、生成函数、容斥，还是结构计数。分类正确，工具通常就只剩一两个候选。

\smallskip
\textbf{初学者检查点：}
\begin{itemize}[leftmargin=2em,itemsep=0.25em]
\item 判断题目所属工具。
\item 列出变量和限制条件。
\item 先做小规模 n=2、3、4 的例子，确认直觉。
\end{itemize}

\smallskip
\textbf{本章读法提醒：} 第一章所有题都先判断“顺序是否重要”。这个判断错了，后面公式一定错。
\end{explainblock}


\begin{sourceblock}{习题 1.15 设 k 是非负整数, n 为正整数. 证明:}
{\small\sloppy
\noindent 习题 1.15 设 k 是非负整数, n 为正整数. 证明:\par
\noindent n 2\par
\noindent X\par
\noindent n 2n\par
\noindent = .\par
\noindent k n\par
\noindent k=0\par
}
\end{sourceblock}


\begin{explainblock}
习题段不要先追求技巧。先给题目分类：它是在问排列、组合、递推、生成函数、容斥，还是结构计数。分类正确，工具通常就只剩一两个候选。

\smallskip
\textbf{初学者检查点：}
\begin{itemize}[leftmargin=2em,itemsep=0.25em]
\item 判断题目所属工具。
\item 列出变量和限制条件。
\item 先做小规模 n=2、3、4 的例子，确认直觉。
\end{itemize}

\smallskip
\textbf{本章读法提醒：} 第一章所有题都先判断“顺序是否重要”。这个判断错了，后面公式一定错。
\end{explainblock}


\begin{sourceblock}{习题 1.16 设 k 是非负整数.}
{\small\sloppy
\noindent 习题 1.16 设 k 是非负整数.\par
\noindent (1) 设 m, n 是非负整数, 证明 Vandermonde 恒等式:\par
\noindent Xk\par
\noindent m+n n m\par
\noindent = .\par
\noindent k i k−i\par
\noindent i=0\par
\noindent (2) 设 s, t 是任意实数, 证明 Chu–Vandermonde 恒等式:\par
\noindent Xk\par
\noindent s+t s t\par
\noindent = .\par
\noindent k i k−i\par
\noindent i=0\par
}
\end{sourceblock}


\begin{explainblock}
习题段不要先追求技巧。先给题目分类：它是在问排列、组合、递推、生成函数、容斥，还是结构计数。分类正确，工具通常就只剩一两个候选。

\smallskip
\textbf{初学者检查点：}
\begin{itemize}[leftmargin=2em,itemsep=0.25em]
\item 判断题目所属工具。
\item 列出变量和限制条件。
\item 先做小规模 n=2、3、4 的例子，确认直觉。
\end{itemize}

\smallskip
\textbf{本章读法提醒：} 第一章所有题都先判断“顺序是否重要”。这个判断错了，后面公式一定错。
\end{explainblock}


\begin{sourceblock}{习题 1.17 设 n 是正整数. 证明}
{\small\sloppy
\noindent 习题 1.17 设 n 是正整数. 证明\par
\noindent X\par
\noindent n\par
\noindent 0, 若 n 是奇数;\par
\noindent k n\par
\noindent (−1) =\par
\noindent k=0\par
\noindent k (−1)m 2m\par
\noindent m , 若 n = 2m.\par
}
\end{sourceblock}


\begin{explainblock}
习题段不要先追求技巧。先给题目分类：它是在问排列、组合、递推、生成函数、容斥，还是结构计数。分类正确，工具通常就只剩一两个候选。

\smallskip
\textbf{初学者检查点：}
\begin{itemize}[leftmargin=2em,itemsep=0.25em]
\item 判断题目所属工具。
\item 列出变量和限制条件。
\item 先做小规模 n=2、3、4 的例子，确认直觉。
\end{itemize}

\smallskip
\textbf{本章读法提醒：} 第一章所有题都先判断“顺序是否重要”。这个判断错了，后面公式一定错。
\end{explainblock}


\begin{sourceblock}{习题 1.18 设 1 ≤ k ≤ n, 用组合方法证明 (1.15), 即:}
{\small\sloppy
\noindent 习题 1.18 设 1 ≤ k ≤ n, 用组合方法证明 (1.15), 即:\par
\noindent n n n n n n\par
\noindent + + + ··· = + + + · · · = 2n−1 .\par
}
\end{sourceblock}


\begin{explainblock}
习题段不要先追求技巧。先给题目分类：它是在问排列、组合、递推、生成函数、容斥，还是结构计数。分类正确，工具通常就只剩一两个候选。

\smallskip
\textbf{初学者检查点：}
\begin{itemize}[leftmargin=2em,itemsep=0.25em]
\item 判断题目所属工具。
\item 列出变量和限制条件。
\item 先做小规模 n=2、3、4 的例子，确认直觉。
\end{itemize}

\smallskip
\textbf{本章读法提醒：} 第一章所有题都先判断“顺序是否重要”。这个判断错了，后面公式一定错。
\end{explainblock}


\begin{sourceblock}{习题 1.19 设 n 为正整数, 用归纳法证明}
{\small\sloppy
\noindent 习题 1.19 设 n 为正整数, 用归纳法证明\par
\noindent ∞\par
\noindent X\par
\noindent = x , |x| < 1.\par
\noindent (1 − x)n k\par
\noindent k=0\par
\noindent (假设\par
\noindent X ∞\par
\noindent = xk , |x| < 1\par
\noindent 1−x\par
\noindent k=0\par
\noindent 成立.)\par
}
\end{sourceblock}


\begin{explainblock}
习题段不要先追求技巧。先给题目分类：它是在问排列、组合、递推、生成函数、容斥，还是结构计数。分类正确，工具通常就只剩一两个候选。

\smallskip
\textbf{初学者检查点：}
\begin{itemize}[leftmargin=2em,itemsep=0.25em]
\item 判断题目所属工具。
\item 列出变量和限制条件。
\item 先做小规模 n=2、3、4 的例子，确认直觉。
\end{itemize}

\smallskip
\textbf{本章读法提醒：} 第一章所有题都先判断“顺序是否重要”。这个判断错了，后面公式一定错。
\end{explainblock}


\begin{sourceblock}{习题 1.20 设 Dn 为 n 元集合的错排个数, 请利用 Dn 的递推关系：}
{\small\sloppy
\noindent 习题 1.20 设 Dn 为 n 元集合的错排个数, 请利用 Dn 的递推关系：\par
\noindent Dn = (n − 1)(Dn−1 + Dn−2 )\par
\noindent 应用数学归纳法证明 Dn 的显示表示式:\par
\noindent X\par
\noindent n\par
\noindent (−1)k\par
\noindent Dn = n! .\par
\noindent k!\par
\noindent k=0\par
}
\end{sourceblock}


\begin{explainblock}
习题段不要先追求技巧。先给题目分类：它是在问排列、组合、递推、生成函数、容斥，还是结构计数。分类正确，工具通常就只剩一两个候选。

\smallskip
\textbf{初学者检查点：}
\begin{itemize}[leftmargin=2em,itemsep=0.25em]
\item 判断题目所属工具。
\item 列出变量和限制条件。
\item 先做小规模 n=2、3、4 的例子，确认直觉。
\end{itemize}

\smallskip
\textbf{本章读法提醒：} 第一章所有题都先判断“顺序是否重要”。这个判断错了，后面公式一定错。
\end{explainblock}


\begin{sourceblock}{习题 1.21 证明错排数 Dn 为偶数的充分必要条件是 n 为奇数.}
{\small\sloppy
\noindent 习题 1.21 证明错排数 Dn 为偶数的充分必要条件是 n 为奇数.\par
}
\end{sourceblock}


\begin{explainblock}
习题段不要先追求技巧。先给题目分类：它是在问排列、组合、递推、生成函数、容斥，还是结构计数。分类正确，工具通常就只剩一两个候选。

\smallskip
\textbf{初学者检查点：}
\begin{itemize}[leftmargin=2em,itemsep=0.25em]
\item 判断题目所属工具。
\item 列出变量和限制条件。
\item 先做小规模 n=2、3、4 的例子，确认直觉。
\end{itemize}

\smallskip
\textbf{本章读法提醒：} 第一章所有题都先判断“顺序是否重要”。这个判断错了，后面公式一定错。
\end{explainblock}


\begin{sourceblock}{习题 1.22 n 个人参加一晚会, 每人寄存一顶帽子, 且各人的帽子互相不同, 会}
{\small\sloppy
\noindent 习题 1.22 n 个人参加一晚会, 每人寄存一顶帽子, 且各人的帽子互相不同, 会\par
\noindent 后每人随便戴其一顶, 求:\par
\noindent (1) 没有任何人戴上自己原来的帽子的概率;\par
\noindent (2) 至少有一个人戴上自己原来的帽子的概率;\par
\noindent (3) 至少有两个人戴上自己原来的帽子的概率.\par
}
\end{sourceblock}


\begin{explainblock}
习题段不要先追求技巧。先给题目分类：它是在问排列、组合、递推、生成函数、容斥，还是结构计数。分类正确，工具通常就只剩一两个候选。

\smallskip
\textbf{初学者检查点：}
\begin{itemize}[leftmargin=2em,itemsep=0.25em]
\item 判断题目所属工具。
\item 列出变量和限制条件。
\item 先做小规模 n=2、3、4 的例子，确认直觉。
\end{itemize}

\smallskip
\textbf{本章读法提醒：} 第一章所有题都先判断“顺序是否重要”。这个判断错了，后面公式一定错。
\end{explainblock}


\begin{sourceblock}{习题 1.23 设 Qn 表示由 1, 2, . . . , n (n ≥ 2) 组成的不含连续数对 (即不含}
{\small\sloppy
\noindent 习题 1.23 设 Qn 表示由 1, 2, . . . , n (n ≥ 2) 组成的不含连续数对 (即不含\par
\noindent i, i + 1) 的全排列的个数, 比如 n = 3 时, 这样的排列有 3 个分别为:\par
\noindent 132, 213, 321.\par
\noindent 事实上, Qn 有下面的计算公式:\par
\noindent X\par
\noindent n−1\par
\noindent n−1\par
\noindent Qn = k\par
\noindent (−1) (n − k)!.\par
\noindent k\par
\noindent k=0\par
\noindent (1) 证明:\par
\noindent Qn = Dn + Dn−1 = (n + 1)Dn−1 + (−1)n .\par
\noindent (2) 证明:\par
\noindent Qn = Dn+1 .\par
\noindent n\par
}
\end{sourceblock}


\begin{explainblock}
习题段不要先追求技巧。先给题目分类：它是在问排列、组合、递推、生成函数、容斥，还是结构计数。分类正确，工具通常就只剩一两个候选。

\smallskip
\textbf{初学者检查点：}
\begin{itemize}[leftmargin=2em,itemsep=0.25em]
\item 判断题目所属工具。
\item 列出变量和限制条件。
\item 先做小规模 n=2、3、4 的例子，确认直觉。
\end{itemize}

\smallskip
\textbf{本章读法提醒：} 第一章所有题都先判断“顺序是否重要”。这个判断错了，后面公式一定错。
\end{explainblock}


\chapter{生成函数及其应用}
\begin{roadmapblock}
本章保留原书正文的主要数学骨架，并在每个定义、定理、例题和论述片段后加入解读。
阅读时不要把目标放在“背下答案”上，而是放在“看懂这一步为什么这样数”上。

\smallskip
\textbf{本章最低目标：} 能说出每个公式左边在数什么、右边在数什么，以及为什么两边相等。
\end{roadmapblock}


\begin{sourceblock}{第 2 章正文}
{\small\sloppy
\noindent 生成函数及其应用\par
\noindent 生成函数 (generating function)(或母函数或发生函数) 方法是解决组合计数问题\par
\noindent 的重要工具之一, 它是组合数学中的基本且重要的方法, 应用很广. 生成函数与分析\par
\noindent 中的幂级数形式上是一致的, 最早将幂级数作为生成函数应用到组合计数上的人是\par
\noindent 欧拉, 他在研究将正整数 n 表示成若干不同的正整数和的方法数 an 时, 天才般地发\par
\noindent 现了 an 与无穷乘积 Q = (1 + x)(1 + x2 )(1 + x3 ) · · · 之间的关系: an 恰是 Q 展开后\par
\noindent xn 对应的系数, 即\par
\noindent Q = a0 + a1 x + a2 x2 + a3 x3 + · · · . (2.1)\par
\noindent 而幂级数 Q 就是我们提及的生成函数, 称为序列 \{an \}n≥ 的生成函数. 比如考虑将\par
\noindent 10 元钱换成 1 元, 2 元和 5 元的方法数对应\par
\noindent (1 + x5 + x10 )(1 + x2 + x4 + x6 + x8 + x10 )(1 + x + x2 + · · · x10 )\par
\noindent 展开后的单项式 x10 的系数为 10.\par
\noindent 对 生 成 函 数 这 套 方 法 作 系 统 的 叙 述 最 早 见 于 拉 普 拉 斯 ∗ 在 1812 年 出 版 的\par
\noindent “Théorie Analytique Des Probabilités”(概率解析理论) 中, 并广泛的使用该方法. 在\par
\noindent 之后, 麦克马洪† 在其两卷组合分析著作 “Combinatory Analysis” (组合分析) 中又补\par
\noindent 充了这种古典的分析技巧. 近代的组合分析家 G.C. Rota 等人也对生成函数的一般\par
\noindent 概念和方法作了进一步的发展. 古典分析学家在使用生成函数方法时, 往往局限于收\par
\noindent 敛的幂级数, 为了克服这个局限性, 只须将这些级数看作代数对象, 采用抽象代数中\par
}
\end{sourceblock}


\begin{explainblock}
这一段是原书论述链的一部分。读的时候把它当作桥梁：它通常负责把上一个定义、例子或定理，引到下一个计数方法。

\smallskip
\textbf{初学者检查点：}
\begin{itemize}[leftmargin=2em,itemsep=0.25em]
\item 找出这一段承接的上一个概念。
\item 找出它准备引出的下一个概念。
\item 把中间的关键等式或构造单独抄出来。
\end{itemize}

\smallskip
\textbf{本章读法提醒：} 第二章看见数列就想“它的生成函数是什么”，看见乘法就想“是否在拼接对象”。
\end{explainblock}


\begin{sourceblock}{第 2 章续读}
{\small\sloppy
\noindent 的环论观点即可.\par
}
\end{sourceblock}


\begin{explainblock}
这一段是原书论述链的一部分。读的时候把它当作桥梁：它通常负责把上一个定义、例子或定理，引到下一个计数方法。

\smallskip
\textbf{初学者检查点：}
\begin{itemize}[leftmargin=2em,itemsep=0.25em]
\item 找出这一段承接的上一个概念。
\item 找出它准备引出的下一个概念。
\item 把中间的关键等式或构造单独抄出来。
\end{itemize}

\smallskip
\textbf{本章读法提醒：} 第二章看见数列就想“它的生成函数是什么”，看见乘法就想“是否在拼接对象”。
\end{explainblock}


\begin{sourceblock}{§ 2.1 生成函数}
{\small\sloppy
\noindent § 2.1. 生成函数\par
\noindent 下面先给出生成函数的精确定义.\par
}
\end{sourceblock}


\begin{explainblock}
这一节先不要急着背公式。先问三个问题：对象是什么，哪些限制条件不能丢，最后到底是在数所有对象、还是在数某个统计量的分布。把这三件事说清楚，后面的公式才会有位置。

\smallskip
\textbf{初学者检查点：}
\begin{itemize}[leftmargin=2em,itemsep=0.25em]
\item 把本节出现的新对象写成一句话定义。
\item 把每个公式左边和右边分别翻译成“在数什么”。
\item 找一个最小非平凡例子，手工列出所有对象。
\end{itemize}

\smallskip
\textbf{本章读法提醒：} 第二章看见数列就想“它的生成函数是什么”，看见乘法就想“是否在拼接对象”。
\end{explainblock}


\begin{sourceblock}{定义 2.1.1 (生成函数) 设 \{ak \}k≥0 是数域 F 上一个序列, 构造如下的函数:}
{\small\sloppy
\noindent 定义 2.1.1 (生成函数) 设 \{ak \}k≥0 是数域 F 上一个序列, 构造如下的函数:\par
\noindent G(x) = a0 + a1 x + a2 x2 + · · · ,\par
\noindent ∗\par
\noindent 皮埃尔-西蒙 · 拉普拉斯 (Pierre-Simon Laplace, 1749–1827), 法国数学家.\par
\noindent †\par
\noindent 珀西 · 亚历山大 · 麦克马洪 (Percy Alexander MacMahon, 1854–1929), 英国数学家.\par
\noindent 与指数型函数:\par
\noindent x x2\par
\noindent F (x) = a0 + a1 + a2 + · · · ,\par
\noindent 1! 2!\par
\noindent 则 G(x) 称为该序列普通型生成函数 (ordinary generating functions), F (x) 称为该\par
\noindent 序列指数型生成函数 (exponential generating functions).\par
}
\end{sourceblock}


\begin{explainblock}
定义段的任务是定边界。这里要特别注意参数范围、是否允许重复、是否考虑顺序、对象有没有标签。后面的每一道例题和证明，本质上都在反复使用这些边界。

\smallskip
\textbf{初学者检查点：}
\begin{itemize}[leftmargin=2em,itemsep=0.25em]
\item 圈出被定义的对象名称。
\item 圈出参数范围，例如 n、k 是否要求正整数。
\item 判断它是有序对象、无序对象、带标签对象还是结构对象。
\end{itemize}

\smallskip
\textbf{本章读法提醒：} 第二章看见数列就想“它的生成函数是什么”，看见乘法就想“是否在拼接对象”。
\end{explainblock}


\begin{sourceblock}{例 2.1.1 (1) 组合数序列 \{ nk \}k≥0 的普通型生成函数为}
{\small\sloppy
\noindent 例 2.1.1 (1) 组合数序列 \{ nk \}k≥0 的普通型生成函数为\par
\noindent n n n 2 n n\par
\noindent + x+ x + ··· + x = (1 + x)n .\par
\noindent (2) k-排列数序列 \{(n)k \}k≥0 的指数型生成函数为\par
\noindent x x2 xn\par
\noindent (n)0 + (n)1 + (n)2 + · · · + (n)n = (1 + x)n .\par
\noindent 1! 2! n!\par
\noindent (3) 全排列序列 \{n!\}n≥0 的指数型生成函数为\par
\noindent ∞\par
\noindent X xn 1\par
\noindent n! = .\par
\noindent n! 1−x\par
\noindent n=0\par
}
\end{sourceblock}


\begin{explainblock}
例题是学习这本书最重要的部分。不要只看答案，要把每个限制条件变成一个计数动作：先安排谁、再插入谁、有没有重复、需不需要除以对称性。

\smallskip
\textbf{初学者检查点：}
\begin{itemize}[leftmargin=2em,itemsep=0.25em]
\item 先写出没有限制时的总数。
\item 逐条加入限制条件，观察总数怎样变化。
\item 最后检查有没有把同一种方案重复算了多次。
\end{itemize}

\smallskip
\textbf{本章读法提醒：} 第二章看见数列就想“它的生成函数是什么”，看见乘法就想“是否在拼接对象”。
\end{explainblock}


\begin{sourceblock}{例 2.1.2 用生成函数方法证明帕斯卡公式:}
{\small\sloppy
\noindent 例 2.1.2 用生成函数方法证明帕斯卡公式:\par
\noindent n n−1 n−1\par
\noindent = + .\par
\noindent k k k−1\par
\noindent 证明: 考虑多项式 (x + 1)n = (x + 1)n−1 (x + 1): 因为\par
\noindent n\par
\noindent X\par
\noindent n n\par
\noindent (x + 1) = xk ,\par
\noindent k\par
\noindent k=0\par
\noindent X\par
\noindent n−1 n−1\par
\noindent n − 1 k X n − 1 k+1\par
\noindent n−1\par
\noindent (x + 1) (1 + x) = x + x ,\par
\noindent k k\par
\noindent k=0 k=0\par
}
\end{sourceblock}


\begin{explainblock}
例题是学习这本书最重要的部分。不要只看答案，要把每个限制条件变成一个计数动作：先安排谁、再插入谁、有没有重复、需不需要除以对称性。

\smallskip
\textbf{初学者检查点：}
\begin{itemize}[leftmargin=2em,itemsep=0.25em]
\item 先写出没有限制时的总数。
\item 逐条加入限制条件，观察总数怎样变化。
\item 最后检查有没有把同一种方案重复算了多次。
\end{itemize}

\smallskip
\textbf{本章读法提醒：} 第二章看见数列就想“它的生成函数是什么”，看见乘法就想“是否在拼接对象”。
\end{explainblock}


\begin{sourceblock}{第 2 章续读}
{\small\sloppy
\noindent 比较 xk 的系数即可.\par
}
\end{sourceblock}


\begin{explainblock}
这一段是原书论述链的一部分。读的时候把它当作桥梁：它通常负责把上一个定义、例子或定理，引到下一个计数方法。

\smallskip
\textbf{初学者检查点：}
\begin{itemize}[leftmargin=2em,itemsep=0.25em]
\item 找出这一段承接的上一个概念。
\item 找出它准备引出的下一个概念。
\item 把中间的关键等式或构造单独抄出来。
\end{itemize}

\smallskip
\textbf{本章读法提醒：} 第二章看见数列就想“它的生成函数是什么”，看见乘法就想“是否在拼接对象”。
\end{explainblock}


\begin{sourceblock}{例 2.1.3 设集合 A = \{1, 2, 3, 4, 5\} 中的 k-可重排列中含有偶数个 1 的排列个数}
{\small\sloppy
\noindent 例 2.1.3 设集合 A = \{1, 2, 3, 4, 5\} 中的 k-可重排列中含有偶数个 1 的排列个数\par
\noindent 为 ak , 求 \{ak \}k≥0 的生成函数和 ak .\par
\noindent 解: 因为只含偶数个 1, 其它的没有限制, 所以 \{ak \}k≥0 的生成函数为:\par
\noindent ∞\par
\noindent X 4\par
\noindent xk x2 x4 x x2\par
\noindent ak = 1+ + + ··· 1+ + + ···\par
\noindent k! 2! 4! 1 2!\par
\noindent k=0\par
\noindent ex + e−x 4x\par
\noindent = e\par
\noindent e5x + e3x\par
\noindent = ,\par
\noindent 展开比较 xk 的系数得:\par
\noindent 5k + 3k\par
\noindent ak = .\par
}
\end{sourceblock}


\begin{explainblock}
例题是学习这本书最重要的部分。不要只看答案，要把每个限制条件变成一个计数动作：先安排谁、再插入谁、有没有重复、需不需要除以对称性。

\smallskip
\textbf{初学者检查点：}
\begin{itemize}[leftmargin=2em,itemsep=0.25em]
\item 先写出没有限制时的总数。
\item 逐条加入限制条件，观察总数怎样变化。
\item 最后检查有没有把同一种方案重复算了多次。
\end{itemize}

\smallskip
\textbf{本章读法提醒：} 第二章看见数列就想“它的生成函数是什么”，看见乘法就想“是否在拼接对象”。
\end{explainblock}


\begin{sourceblock}{例 2.1.4 求连续 10 次投掷一枚骰子出现点数和为 30 的概率.}
{\small\sloppy
\noindent 例 2.1.4 求连续 10 次投掷一枚骰子出现点数和为 30 的概率.\par
\noindent 解: 用多项式\par
\noindent x + x2 + x3 + x4 + x5 + x6\par
\noindent 来表示投掷一次可能出现的点数 1, 2, . . . , 6, 假设第 i(i = 1, 2, . . . , 10) 次投掷出现的\par
\noindent 点数是 mi , 若用 x 个的单项式表示, 则第 i(i = 1, 2, . . . , 10) 次对应单项式为 xmi ,\par
\noindent P\par
\noindent 而这 10 次投掷得到的点数和为 ni=1 mi , 恰对应这 10 个单项式乘积后的 x 的指数,\par
\noindent 故而可知\par
\noindent (x + x2 + x3 + x4 + x5 + x6 )10\par
\noindent 展开得到多项式中 x30 的系数就是这 10 次投掷点数和为 30 所有的情形的个数. 由\par
\noindent 2 3 4 5 6 10 10 1 − x6\par
\noindent (x + x + x + x + x + x ) =x\par
\noindent 1−x\par
\noindent = x10 (1 − x6 )10 (1 − x)−10\par
\noindent X ∞\par
\noindent 10 6i X 10 − 1 + i i\par
\noindent =x (−1) x x,\par
\noindent i i\par
}
\end{sourceblock}


\begin{explainblock}
例题是学习这本书最重要的部分。不要只看答案，要把每个限制条件变成一个计数动作：先安排谁、再插入谁、有没有重复、需不需要除以对称性。

\smallskip
\textbf{初学者检查点：}
\begin{itemize}[leftmargin=2em,itemsep=0.25em]
\item 先写出没有限制时的总数。
\item 逐条加入限制条件，观察总数怎样变化。
\item 最后检查有没有把同一种方案重复算了多次。
\end{itemize}

\smallskip
\textbf{本章读法提醒：} 第二章看见数列就想“它的生成函数是什么”，看见乘法就想“是否在拼接对象”。
\end{explainblock}


\begin{sourceblock}{第 2 章续读}
{\small\sloppy
\noindent i=0 k=0\par
\noindent 所以 x30 的系数为\par
\noindent 10 29 10 23 10 17 10 11\par
\noindent − + − = 2930455.\par
\noindent 故所求概率为\par
\noindent ≈ 0.0485.\par
}
\end{sourceblock}


\begin{explainblock}
这一段是原书论述链的一部分。读的时候把它当作桥梁：它通常负责把上一个定义、例子或定理，引到下一个计数方法。

\smallskip
\textbf{初学者检查点：}
\begin{itemize}[leftmargin=2em,itemsep=0.25em]
\item 找出这一段承接的上一个概念。
\item 找出它准备引出的下一个概念。
\item 把中间的关键等式或构造单独抄出来。
\end{itemize}

\smallskip
\textbf{本章读法提醒：} 第二章看见数列就想“它的生成函数是什么”，看见乘法就想“是否在拼接对象”。
\end{explainblock}


\begin{sourceblock}{§ 2.2 形式幂级数及其性质}
{\small\sloppy
\noindent § 2.2. 形式幂级数及其性质\par
\noindent 从上节叙述知在使用生成函数方法时, 可能会存在收敛性的问题, 比如考虑序列\par
\noindent P\par
\noindent \{n!\}n≥0 的普通型生成函数 n≥ n!xn , 由柯西–阿达马公式知它的收敛半径为 0, 也\par
\noindent 即是说, 若考虑收敛性, 则就连很简单的序列都可能不存在生成函数, 这一问题很大\par
\noindent 地限制了生成函数的应用.\par
\noindent 通过将生成函数看作代数对象, 采用抽象代数中的环论的观点后这种限制就不复\par
\noindent 存在了，这一点在 20 年代 E. T. Bell 的著作中早就给出了详细地叙述. 按照近代的\par
\noindent 观点, 任意序列的生成函数不过是一种形式上的幂级数, 下面我们就从代数的观点引\par
\noindent 入形式幂级数的概念.\par
}
\end{sourceblock}


\begin{explainblock}
这一节先不要急着背公式。先问三个问题：对象是什么，哪些限制条件不能丢，最后到底是在数所有对象、还是在数某个统计量的分布。把这三件事说清楚，后面的公式才会有位置。

\smallskip
\textbf{初学者检查点：}
\begin{itemize}[leftmargin=2em,itemsep=0.25em]
\item 把本节出现的新对象写成一句话定义。
\item 把每个公式左边和右边分别翻译成“在数什么”。
\item 找一个最小非平凡例子，手工列出所有对象。
\end{itemize}

\smallskip
\textbf{本章读法提醒：} 第二章看见数列就想“它的生成函数是什么”，看见乘法就想“是否在拼接对象”。
\end{explainblock}


\begin{sourceblock}{定义 2.2.1 设 \{ak \}k≥0 是数域 F 上一个序列, 设 x 为不定元, 称形如}
{\small\sloppy
\noindent 定义 2.2.1 设 \{ak \}k≥0 是数域 F 上一个序列, 设 x 为不定元, 称形如\par
\noindent ∞\par
\noindent X\par
\noindent an xn = a0 + a1 x + a2 x2 + a3 x3 + · · ·\par
\noindent n=0\par
\noindent 的和式为数域 F 上的形式幂级数.\par
}
\end{sourceblock}


\begin{explainblock}
定义段的任务是定边界。这里要特别注意参数范围、是否允许重复、是否考虑顺序、对象有没有标签。后面的每一道例题和证明，本质上都在反复使用这些边界。

\smallskip
\textbf{初学者检查点：}
\begin{itemize}[leftmargin=2em,itemsep=0.25em]
\item 圈出被定义的对象名称。
\item 圈出参数范围，例如 n、k 是否要求正整数。
\item 判断它是有序对象、无序对象、带标签对象还是结构对象。
\end{itemize}

\smallskip
\textbf{本章读法提醒：} 第二章看见数列就想“它的生成函数是什么”，看见乘法就想“是否在拼接对象”。
\end{explainblock}


\begin{sourceblock}{注记 2.2.2 形式幂级数中的 x 只是一个不定元, 是一个抽象的符号, 并不需要对 x}
{\small\sloppy
\noindent 注记 2.2.2 形式幂级数中的 x 只是一个不定元, 是一个抽象的符号, 并不需要对 x\par
\noindent 进行赋值, 因此也不必考虑它的收敛性. 数域 F 上所有的形式幂级数记为 F[[x]], 在\par
\noindent 集合 F[[x]] 上适当的定义加法和乘法运算可使其构成一个可交换的整环.\par
\noindent 考虑数域 F 上的序列 \{ak \}k≥0 , \{bk \}k≥0 和 \{ck \}k≥0 , 并设这些序列对应的形式幂\par
\noindent 级数分别为 f (x), g(x) 及 h(x), 则 f (x), g(x), h(x) ∈ F[[x]]. 我们称 f (x) 与 g(x) 相\par
\noindent 等当且仅当\par
\noindent ak = bk , k = 0, 1, 2, . . . ,\par
\noindent 并记为 f (x) = g(x). 若 α ∈ F, 称 αf (x) 是 α 与 f (x) 的数乘当且仅当\par
\noindent ∞\par
\noindent X\par
\noindent αf (x) = αak xk .\par
\noindent k=0\par
\noindent 集合 F[[x]] 上的加法与乘法两种运算可按下面的方式定义:\par
\noindent 加法: h(x) = f (x) + g(x),\par
\noindent 当且仅当\par
\noindent ck = ak + bk , k = 0, 1, 2, . . . ;\par
\noindent 乘法: h(x) = f (x)g(x),\par
\noindent 当且仅当\par
}
\end{sourceblock}


\begin{explainblock}
注记通常是在补充术语、记号、历史或一个容易忽略的约定。它未必给新公式，但常常决定你后面会不会把符号读错。

\smallskip
\textbf{初学者检查点：}
\begin{itemize}[leftmargin=2em,itemsep=0.25em]
\item 记下新符号的读法。
\item 确认这个约定是否只在本章使用。
\item 把注记里的限制条件补到前面的定义旁边。
\end{itemize}

\smallskip
\textbf{本章读法提醒：} 第二章看见数列就想“它的生成函数是什么”，看见乘法就想“是否在拼接对象”。
\end{explainblock}


\begin{sourceblock}{第 2 章续读}
{\small\sloppy
\noindent X\par
\noindent k\par
\noindent ck = ai bk−i , k = 0, 1, 2, . . . ,\par
\noindent i=0\par
\noindent 此时称序列 \{ck \}k≥0 是序列 \{ak \}k≥0 与 \{bk \}k≥0 的卷积或柯西乘积.\par
\noindent 可验证集合 F[[x]] 关于上面定义的加法与乘法运算构成一个可交换的整环, 其中\par
\noindent 零元是序列 \{0, 0, 0, . . .\} 的形式幂级数, 记为 0, 单位元是序列 \{1, 0, 0, . . .\} 的形式幂\par
\noindent 级数, 记为 1. 若形式幂级数 f (x) 有可逆元, 则记其逆元为 f −1 (x).\par
}
\end{sourceblock}


\begin{explainblock}
这一段是原书论述链的一部分。读的时候把它当作桥梁：它通常负责把上一个定义、例子或定理，引到下一个计数方法。

\smallskip
\textbf{初学者检查点：}
\begin{itemize}[leftmargin=2em,itemsep=0.25em]
\item 找出这一段承接的上一个概念。
\item 找出它准备引出的下一个概念。
\item 把中间的关键等式或构造单独抄出来。
\end{itemize}

\smallskip
\textbf{本章读法提醒：} 第二章看见数列就想“它的生成函数是什么”，看见乘法就想“是否在拼接对象”。
\end{explainblock}


\begin{sourceblock}{例 2.2.1 由}
{\small\sloppy
\noindent 例 2.2.1 由\par
\noindent (1 − x)(1 + x + x2 + · · · ) = 1\par
\noindent 知形式幂级数 1 − x 与形式幂级数 (1 + x + x2 + · · · ) 在环 F[[x]] 中互为逆元.\par
\noindent 一般地, 下面的定理给出了哪些形式幂级数是具有逆元的.\par
\noindent P\par
}
\end{sourceblock}


\begin{explainblock}
例题是学习这本书最重要的部分。不要只看答案，要把每个限制条件变成一个计数动作：先安排谁、再插入谁、有没有重复、需不需要除以对称性。

\smallskip
\textbf{初学者检查点：}
\begin{itemize}[leftmargin=2em,itemsep=0.25em]
\item 先写出没有限制时的总数。
\item 逐条加入限制条件，观察总数怎样变化。
\item 最后检查有没有把同一种方案重复算了多次。
\end{itemize}

\smallskip
\textbf{本章读法提醒：} 第二章看见数列就想“它的生成函数是什么”，看见乘法就想“是否在拼接对象”。
\end{explainblock}


\begin{sourceblock}{定理 2.2.3 设 f (x) = ∞k=0 ak x ∈ F, 则 f (x) 有乘法逆元当且仅当 a0 6= 0, 即}
{\small\sloppy
\noindent 定理 2.2.3 设 f (x) = ∞k=0 ak x ∈ F, 则 f (x) 有乘法逆元当且仅当 a0 6= 0, 即\par
\noindent k\par
\noindent f (0) 6= 0.\par
\noindent 证明: 必要性显然成立, 现在证明充分性. 即证当 a0 6= 0 时, 存在形式幂级数\par
\noindent P∞\par
\noindent g(x) = k=0 bk xk ∈ F[[x]] 使得\par
\noindent ∞ X\par
\noindent X k\par
\noindent f (x)g(x) = ai bk−i xk = 1,\par
\noindent k=0 i=0\par
\noindent 比较等式两边 xk (k = 0, 1, 2, . . .) 的系数得方程组:\par
\noindent \par
\noindent \par
\noindent a0 b0 = 1,\par
\noindent \par
\noindent \par
\noindent \par
\noindent \par
}
\end{sourceblock}


\begin{explainblock}
定理段不要只看结论，要看它把哪两个计数方式连在一起。组合学证明最常见的技巧是：左边按一种标准数，右边按另一种标准数，然后说明两边数的是同一批对象。

\smallskip
\textbf{初学者检查点：}
\begin{itemize}[leftmargin=2em,itemsep=0.25em]
\item 先复述定理的假设，不要把适用范围扩大。
\item 找出证明中的基本计数原则：乘法、加法、双射、递推、容斥或生成函数。
\item 检查最后的等式化简是否和前面的对象解释一致。
\end{itemize}

\smallskip
\textbf{本章读法提醒：} 第二章看见数列就想“它的生成函数是什么”，看见乘法就想“是否在拼接对象”。
\end{explainblock}


\begin{sourceblock}{第 2 章续读}
{\small\sloppy
\noindent \par
\noindent \par
\noindent \par
\noindent a1 b0 + a0 b1 = 0,\par
\noindent \par
\noindent \par
\noindent \par
\noindent ··· (2.2)\par
\noindent \par
\noindent \par
\noindent \par
\noindent \par
\noindent \par
\noindent \par
\noindent \par
\noindent ak b0 + ak−1 b1 + · · · + a0 bk = 0,\par
\noindent \par
\noindent \par
}
\end{sourceblock}


\begin{explainblock}
这一段是原书论述链的一部分。读的时候把它当作桥梁：它通常负责把上一个定义、例子或定理，引到下一个计数方法。

\smallskip
\textbf{初学者检查点：}
\begin{itemize}[leftmargin=2em,itemsep=0.25em]
\item 找出这一段承接的上一个概念。
\item 找出它准备引出的下一个概念。
\item 把中间的关键等式或构造单独抄出来。
\end{itemize}

\smallskip
\textbf{本章读法提醒：} 第二章看见数列就想“它的生成函数是什么”，看见乘法就想“是否在拼接对象”。
\end{explainblock}


\begin{sourceblock}{第 2 章续读}
{\small\sloppy
\noindent \par
\noindent \par
\noindent  ··· ···\par
\noindent 所以只需证明方程组 (2.2) 有解. 由于 a0 6= 0, b0 , b1 , b2 , . . . 显然可由方程组 (2.2) 依\par
\noindent 次唯一确定, 即 g(x) 被唯一确定, 故 f (x) 可逆且其可逆元由方程组 (2.2) 唯一确\par
\noindent 定.\par
\noindent 由定理 2.2.3 知可交换整环 F[[x]] 上任意一个有非 0 常数项的形式幂级数均可求\par
\noindent 其逆, 因此如收敛级数一样, 在整环 F[[x]] 作四则运算可以畅行无阻. 又因为如上规\par
\noindent 定的运算方式与收敛级数间的运算是一致的, 所以在特别情形下 (即当形式幂级数收\par
\noindent 敛时), 计算的最终结果—即由生成函数最终确定的序列总是一样的. 因此本书中以\par
\noindent 后采用的生成函数均理解为形式幂级数, 而不再考虑其收敛性的问题.\par
\noindent P\par
\noindent 对环 F[[x]] 中形式幂级数还可定义形式导数. 设 f (x) = ∞k=0 an x ∈ F[[x]], 记\par
\noindent k\par
\noindent d\par
\noindent D = dx 表示微分算子, 规定:\par
\noindent ∞\par
\noindent X\par
}
\end{sourceblock}


\begin{explainblock}
这一段是原书论述链的一部分。读的时候把它当作桥梁：它通常负责把上一个定义、例子或定理，引到下一个计数方法。

\smallskip
\textbf{初学者检查点：}
\begin{itemize}[leftmargin=2em,itemsep=0.25em]
\item 找出这一段承接的上一个概念。
\item 找出它准备引出的下一个概念。
\item 把中间的关键等式或构造单独抄出来。
\end{itemize}

\smallskip
\textbf{本章读法提醒：} 第二章看见数列就想“它的生成函数是什么”，看见乘法就想“是否在拼接对象”。
\end{explainblock}


\begin{sourceblock}{第 2 章续读}
{\small\sloppy
\noindent Df (x) := a1 + 2a2 x + 3a3 x + · · · =\par
\noindent kak xk−1 ,\par
\noindent k=1\par
\noindent 则称 Df (x) 为 f (x) 的形式导数; f (x) 的 n 次形式导数可以递归地定义为:\par
\noindent \par
\noindent D0 f (x) := f (x),\par
\noindent (2.3)\par
\noindent  n\par
\noindent D f (x) := D[Dn−1 f (x)].\par
\noindent 为了书写方便, 以后用 f ′ (x), f ′′ (x), . . . 代替 D[f (x)], D2 [f (x)], . . . 来表示 f (x)\par
\noindent 的各阶形式导数.\par
\noindent 形式导数也满足普通导数的性质.\par
}
\end{sourceblock}


\begin{explainblock}
这一段是原书论述链的一部分。读的时候把它当作桥梁：它通常负责把上一个定义、例子或定理，引到下一个计数方法。

\smallskip
\textbf{初学者检查点：}
\begin{itemize}[leftmargin=2em,itemsep=0.25em]
\item 找出这一段承接的上一个概念。
\item 找出它准备引出的下一个概念。
\item 把中间的关键等式或构造单独抄出来。
\end{itemize}

\smallskip
\textbf{本章读法提醒：} 第二章看见数列就想“它的生成函数是什么”，看见乘法就想“是否在拼接对象”。
\end{explainblock}


\begin{sourceblock}{命题 2.2.4 设 f (x) 与 g(x) 为数域 F 上的生成函数. 则}
{\small\sloppy
\noindent 命题 2.2.4 设 f (x) 与 g(x) 为数域 F 上的生成函数. 则\par
\noindent (1)\par
\noindent (f (x) + g(x))′ = f ′ (x) + g ′ (x),\par
\noindent (2)\par
\noindent (f (x)g(x))′ = f ′ (x)g(x) + f (x)g ′ (x).\par
}
\end{sourceblock}


\begin{explainblock}
定理段不要只看结论，要看它把哪两个计数方式连在一起。组合学证明最常见的技巧是：左边按一种标准数，右边按另一种标准数，然后说明两边数的是同一批对象。

\smallskip
\textbf{初学者检查点：}
\begin{itemize}[leftmargin=2em,itemsep=0.25em]
\item 先复述定理的假设，不要把适用范围扩大。
\item 找出证明中的基本计数原则：乘法、加法、双射、递推、容斥或生成函数。
\item 检查最后的等式化简是否和前面的对象解释一致。
\end{itemize}

\smallskip
\textbf{本章读法提醒：} 第二章看见数列就想“它的生成函数是什么”，看见乘法就想“是否在拼接对象”。
\end{explainblock}


\begin{sourceblock}{注记 2.2.5 当形式幂级数在某个范围内收敛时, 形式导数就与微积分中的求导运}
{\small\sloppy
\noindent 注记 2.2.5 当形式幂级数在某个范围内收敛时, 形式导数就与微积分中的求导运\par
\noindent 算一样.\par
\noindent P∞ k ∈ F[[x]], 也可定义积分运算:\par
\noindent 设 f (x) = k=0 an x\par
\noindent Z x X1 ∞\par
\noindent f (t)dt = a0 x + a1 x2 + · · · = ak−1 xk ,\par
\noindent k=1\par
\noindent 而且满足 Z x ′\par
\noindent f (t)dt = f (t).\par
}
\end{sourceblock}


\begin{explainblock}
注记通常是在补充术语、记号、历史或一个容易忽略的约定。它未必给新公式，但常常决定你后面会不会把符号读错。

\smallskip
\textbf{初学者检查点：}
\begin{itemize}[leftmargin=2em,itemsep=0.25em]
\item 记下新符号的读法。
\item 确认这个约定是否只在本章使用。
\item 把注记里的限制条件补到前面的定义旁边。
\end{itemize}

\smallskip
\textbf{本章读法提醒：} 第二章看见数列就想“它的生成函数是什么”，看见乘法就想“是否在拼接对象”。
\end{explainblock}


\begin{sourceblock}{§ 2.3 生成函数的应用}
{\small\sloppy
\noindent § 2.3. 生成函数的应用\par
\noindent 生成函数的应用十分广泛, 有兴趣的读者可以参考 H. Wilf 的专著 [31]. 这节我\par
\noindent 们将举一些相关的例子来说明如何用生成函数的方法解决实际的问题.\par
}
\end{sourceblock}


\begin{explainblock}
这一节先不要急着背公式。先问三个问题：对象是什么，哪些限制条件不能丢，最后到底是在数所有对象、还是在数某个统计量的分布。把这三件事说清楚，后面的公式才会有位置。

\smallskip
\textbf{初学者检查点：}
\begin{itemize}[leftmargin=2em,itemsep=0.25em]
\item 把本节出现的新对象写成一句话定义。
\item 把每个公式左边和右边分别翻译成“在数什么”。
\item 找一个最小非平凡例子，手工列出所有对象。
\end{itemize}

\smallskip
\textbf{本章读法提醒：} 第二章看见数列就想“它的生成函数是什么”，看见乘法就想“是否在拼接对象”。
\end{explainblock}


\begin{sourceblock}{例 2.3.1 设 Dn 为 n 元集合的错排数, 由递推关式}
{\small\sloppy
\noindent 例 2.3.1 设 Dn 为 n 元集合的错排数, 由递推关式\par
\noindent Dn = nDn−1 + (−1)n , D0 = 1, n ≥ 1.\par
\noindent 用生成函数的方法证明\par
\noindent X\par
\noindent n\par
\noindent (−1)k\par
\noindent Dn = n! .\par
\noindent k!\par
\noindent k=0\par
\noindent 证明: 设 \{Dn \}n≥0 的指数型生成函数为 F (x), 即\par
\noindent ∞\par
\noindent X xn\par
\noindent F (x) = Dn . (2.4)\par
\noindent n!\par
\noindent n=0\par
\noindent 则由 Dn 的递推关系:\par
\noindent Dn = nDn−1 + (−1)n , D0 = 1, n ≥ 1,\par
\noindent 得\par
}
\end{sourceblock}


\begin{explainblock}
例题是学习这本书最重要的部分。不要只看答案，要把每个限制条件变成一个计数动作：先安排谁、再插入谁、有没有重复、需不需要除以对称性。

\smallskip
\textbf{初学者检查点：}
\begin{itemize}[leftmargin=2em,itemsep=0.25em]
\item 先写出没有限制时的总数。
\item 逐条加入限制条件，观察总数怎样变化。
\item 最后检查有没有把同一种方案重复算了多次。
\end{itemize}

\smallskip
\textbf{本章读法提醒：} 第二章看见数列就想“它的生成函数是什么”，看见乘法就想“是否在拼接对象”。
\end{explainblock}


\begin{sourceblock}{第 2 章续读}
{\small\sloppy
\noindent ∞\par
\noindent X X ∞\par
\noindent xn xn\par
\noindent F (x) = Dn = D0 + (nDn−1 + (−1)n )\par
\noindent n! n!\par
\noindent n=0 n=1\par
\noindent ∞\par
\noindent X X ∞\par
\noindent xn xn\par
\noindent =1+ Dn−1 + (−1)n\par
\noindent (n − 1)! n!\par
\noindent n=1 n=1\par
\noindent X ∞\par
\noindent xn X (−x)n\par
\noindent =x Dn +\par
\noindent n! n!\par
\noindent n=0 n=0\par
\noindent = xF (x) + e−x ,\par
}
\end{sourceblock}


\begin{explainblock}
这一段是原书论述链的一部分。读的时候把它当作桥梁：它通常负责把上一个定义、例子或定理，引到下一个计数方法。

\smallskip
\textbf{初学者检查点：}
\begin{itemize}[leftmargin=2em,itemsep=0.25em]
\item 找出这一段承接的上一个概念。
\item 找出它准备引出的下一个概念。
\item 把中间的关键等式或构造单独抄出来。
\end{itemize}

\smallskip
\textbf{本章读法提醒：} 第二章看见数列就想“它的生成函数是什么”，看见乘法就想“是否在拼接对象”。
\end{explainblock}


\begin{sourceblock}{第 2 章续读}
{\small\sloppy
\noindent 整理得\par
\noindent ∞\par
\noindent X X (−1)k n\par
\noindent e−x\par
\noindent F (x) = = xn . (2.5)\par
\noindent 1−x k!\par
\noindent n=0 k=0\par
\noindent 比较 (2.4) 和 (2.5), 立得\par
\noindent X\par
\noindent n\par
\noindent (−1)k\par
\noindent Dn = n! .\par
\noindent k!\par
\noindent k=0\par
}
\end{sourceblock}


\begin{explainblock}
这一段是原书论述链的一部分。读的时候把它当作桥梁：它通常负责把上一个定义、例子或定理，引到下一个计数方法。

\smallskip
\textbf{初学者检查点：}
\begin{itemize}[leftmargin=2em,itemsep=0.25em]
\item 找出这一段承接的上一个概念。
\item 找出它准备引出的下一个概念。
\item 把中间的关键等式或构造单独抄出来。
\end{itemize}

\smallskip
\textbf{本章读法提醒：} 第二章看见数列就想“它的生成函数是什么”，看见乘法就想“是否在拼接对象”。
\end{explainblock}


\begin{sourceblock}{例 2.3.2 设 Dn 为 n 元集合的错排数，由下面等式}
{\small\sloppy
\noindent 例 2.3.2 设 Dn 为 n 元集合的错排数，由下面等式\par
\noindent n\par
\noindent X n\par
\noindent Dn−k = n!,\par
\noindent k\par
\noindent k=0\par
\noindent 利用生成函数方法证明\par
\noindent X\par
\noindent n\par
\noindent (−1)k\par
\noindent Dn = n! .\par
\noindent k!\par
\noindent k=0\par
\noindent 证明: 由\par
\noindent ∞\par
\noindent X xn 1\par
\noindent n! = ,\par
\noindent n! 1−x\par
}
\end{sourceblock}


\begin{explainblock}
例题是学习这本书最重要的部分。不要只看答案，要把每个限制条件变成一个计数动作：先安排谁、再插入谁、有没有重复、需不需要除以对称性。

\smallskip
\textbf{初学者检查点：}
\begin{itemize}[leftmargin=2em,itemsep=0.25em]
\item 先写出没有限制时的总数。
\item 逐条加入限制条件，观察总数怎样变化。
\item 最后检查有没有把同一种方案重复算了多次。
\end{itemize}

\smallskip
\textbf{本章读法提醒：} 第二章看见数列就想“它的生成函数是什么”，看见乘法就想“是否在拼接对象”。
\end{explainblock}


\begin{sourceblock}{第 2 章续读}
{\small\sloppy
\noindent n=0\par
\noindent 及\par
\noindent n\par
\noindent X n\par
\noindent Dn−k = n!\par
\noindent k\par
\noindent k=0\par
\noindent 可得\par
\noindent n\par
\noindent ∞ X\par
\noindent X ∞\par
\noindent n xn X xn 1\par
\noindent Dn−k = n! = , (2.6)\par
\noindent k n! n! 1−x\par
\noindent n=0 k=0 n=0\par
\noindent 对等式 (2.6) 左边整理得\par
\noindent n\par
\noindent ∞ X\par
}
\end{sourceblock}


\begin{explainblock}
这一段是原书论述链的一部分。读的时候把它当作桥梁：它通常负责把上一个定义、例子或定理，引到下一个计数方法。

\smallskip
\textbf{初学者检查点：}
\begin{itemize}[leftmargin=2em,itemsep=0.25em]
\item 找出这一段承接的上一个概念。
\item 找出它准备引出的下一个概念。
\item 把中间的关键等式或构造单独抄出来。
\end{itemize}

\smallskip
\textbf{本章读法提醒：} 第二章看见数列就想“它的生成函数是什么”，看见乘法就想“是否在拼接对象”。
\end{explainblock}


\begin{sourceblock}{第 2 章续读}
{\small\sloppy
\noindent X ∞ n\par
\noindent n xn X X Dn−k xn−k xk\par
\noindent Dn−k =\par
\noindent k n! (n − k)! k!\par
\noindent n=0 k=0 n=0 k=0\par
\noindent ∞\par
\noindent X ∞\par
\noindent Dm xm X xk\par
\noindent =\par
\noindent m! k!\par
\noindent m=0 k=0\par
\noindent ∞\par
\noindent X\par
\noindent x Dm xm\par
\noindent =e ,\par
\noindent m!\par
\noindent m=0\par
\noindent 故\par
}
\end{sourceblock}


\begin{explainblock}
这一段是原书论述链的一部分。读的时候把它当作桥梁：它通常负责把上一个定义、例子或定理，引到下一个计数方法。

\smallskip
\textbf{初学者检查点：}
\begin{itemize}[leftmargin=2em,itemsep=0.25em]
\item 找出这一段承接的上一个概念。
\item 找出它准备引出的下一个概念。
\item 把中间的关键等式或构造单独抄出来。
\end{itemize}

\smallskip
\textbf{本章读法提醒：} 第二章看见数列就想“它的生成函数是什么”，看见乘法就想“是否在拼接对象”。
\end{explainblock}


\begin{sourceblock}{第 2 章续读}
{\small\sloppy
\noindent ∞\par
\noindent X ∞\par
\noindent X X (−1)k n\par
\noindent Dn xn e−x\par
\noindent = = xn .\par
\noindent n! 1−x k!\par
\noindent n=0 n=0 k=0\par
\noindent 比较两边 xn 的系数可得\par
\noindent X\par
\noindent n\par
\noindent (−1)k\par
\noindent Dn = n! .\par
\noindent k!\par
\noindent k=0\par
}
\end{sourceblock}


\begin{explainblock}
这一段是原书论述链的一部分。读的时候把它当作桥梁：它通常负责把上一个定义、例子或定理，引到下一个计数方法。

\smallskip
\textbf{初学者检查点：}
\begin{itemize}[leftmargin=2em,itemsep=0.25em]
\item 找出这一段承接的上一个概念。
\item 找出它准备引出的下一个概念。
\item 把中间的关键等式或构造单独抄出来。
\end{itemize}

\smallskip
\textbf{本章读法提醒：} 第二章看见数列就想“它的生成函数是什么”，看见乘法就想“是否在拼接对象”。
\end{explainblock}


\begin{sourceblock}{例 2.3.3 用生成函数方法证明}
{\small\sloppy
\noindent 例 2.3.3 用生成函数方法证明\par
\noindent n n+k−1\par
\noindent = .\par
\noindent k k\par
\noindent 证明: 设 n 元集合 A = \{a1 , a2 , . . . , an \}, 则 n\par
\noindent k 可表示集合 A 的 k 可重组\par
\noindent 合的个数. 可用形式幂级数\par
\noindent 1 + x + x2 + x3 + · · ·\par
\noindent 表示 A 中元素可能在组合中可能出现的次数: 0, 1, 2, . . ., 则序列 \{ n\par
\noindent k \}k≥0 普通型\par
\noindent 生成函数为\par
\noindent Y\par
\noindent n\par
\noindent (1 + x + x2 + x3 + · · · ) = ,\par
\noindent (1 − x)n\par
\noindent i=1\par
\noindent 由牛顿二项式定理得\par
\noindent n\par
}
\end{sourceblock}


\begin{explainblock}
例题是学习这本书最重要的部分。不要只看答案，要把每个限制条件变成一个计数动作：先安排谁、再插入谁、有没有重复、需不需要除以对称性。

\smallskip
\textbf{初学者检查点：}
\begin{itemize}[leftmargin=2em,itemsep=0.25em]
\item 先写出没有限制时的总数。
\item 逐条加入限制条件，观察总数怎样变化。
\item 最后检查有没有把同一种方案重复算了多次。
\end{itemize}

\smallskip
\textbf{本章读法提醒：} 第二章看见数列就想“它的生成函数是什么”，看见乘法就想“是否在拼接对象”。
\end{explainblock}


\begin{sourceblock}{第 2 章续读}
{\small\sloppy
\noindent X\par
\noindent = (−x)k ,\par
\noindent (1 − x)n k\par
\noindent k=0\par
\noindent 所以 A 的 k-可重组合的个数为\par
\noindent −n n+k−1\par
\noindent = ,\par
\noindent k k\par
\noindent 即\par
\noindent n n+k−1\par
\noindent = .\par
\noindent k k\par
}
\end{sourceblock}


\begin{explainblock}
这一段是原书论述链的一部分。读的时候把它当作桥梁：它通常负责把上一个定义、例子或定理，引到下一个计数方法。

\smallskip
\textbf{初学者检查点：}
\begin{itemize}[leftmargin=2em,itemsep=0.25em]
\item 找出这一段承接的上一个概念。
\item 找出它准备引出的下一个概念。
\item 把中间的关键等式或构造单独抄出来。
\end{itemize}

\smallskip
\textbf{本章读法提醒：} 第二章看见数列就想“它的生成函数是什么”，看见乘法就想“是否在拼接对象”。
\end{explainblock}


\begin{sourceblock}{例 2.3.4 试确定不定方程}
{\small\sloppy
\noindent 例 2.3.4 试确定不定方程\par
\noindent x1 + x2 + x3 = 20,\par
\noindent 满足:\par
\noindent 1 ≤ x1 ≤ 10, 0 ≤ x2 ≤ 8, 4 ≤ x3 ≤ 10\par
\noindent 的整数解的个数.\par
\noindent 解: 由题中要求可知 x1 的取值为 1 到 10 的正整数, 不妨用多项式\par
\noindent x + x2 + x3 + · · · x10\par
\noindent 表示 x1 的可能取值, 同理, 用多项式\par
\noindent 1 + x + x2 + · · · + x8 ,\par
\noindent 和多项式\par
\noindent x4 + x5 + · · · + x10\par
\noindent 分别表示 x2 与 x3 的可能取值. 因此不定方程的整数解的个数等于\par
\noindent (x + x2 + x3 + · · · x10 )(1 + x + x2 + · · · + x8 )(x4 + x5 + · · · + x10 )\par
\noindent 展开后对应的 x20 的系数. 又因为\par
\noindent (x + x2 + x3 + · · · x10 )(1 + x + x2 + · · · + x8 )(x4 + x5 + · · · + x10 )\par
\noindent (1 − x10 )(1 − x9 )(1 − x7 )\par
\noindent = x5\par
\noindent (1 − x)3\par
}
\end{sourceblock}


\begin{explainblock}
例题是学习这本书最重要的部分。不要只看答案，要把每个限制条件变成一个计数动作：先安排谁、再插入谁、有没有重复、需不需要除以对称性。

\smallskip
\textbf{初学者检查点：}
\begin{itemize}[leftmargin=2em,itemsep=0.25em]
\item 先写出没有限制时的总数。
\item 逐条加入限制条件，观察总数怎样变化。
\item 最后检查有没有把同一种方案重复算了多次。
\end{itemize}

\smallskip
\textbf{本章读法提醒：} 第二章看见数列就想“它的生成函数是什么”，看见乘法就想“是否在拼接对象”。
\end{explainblock}


\begin{sourceblock}{第 2 章续读}
{\small\sloppy
\noindent ∞\par
\noindent X\par
\noindent k+2 k\par
\noindent = (−x 31\par
\noindent +x 24\par
\noindent +x 22\par
\noindent +x 21\par
\noindent −x 15\par
\noindent −x 14\par
\noindent −x 12 5\par
\noindent +x ) x ,\par
\noindent k\par
\noindent k=0\par
\noindent 所以整数解的个数为\par
\noindent − − − + = 42.\par
}
\end{sourceblock}


\begin{explainblock}
这一段是原书论述链的一部分。读的时候把它当作桥梁：它通常负责把上一个定义、例子或定理，引到下一个计数方法。

\smallskip
\textbf{初学者检查点：}
\begin{itemize}[leftmargin=2em,itemsep=0.25em]
\item 找出这一段承接的上一个概念。
\item 找出它准备引出的下一个概念。
\item 把中间的关键等式或构造单独抄出来。
\end{itemize}

\smallskip
\textbf{本章读法提醒：} 第二章看见数列就想“它的生成函数是什么”，看见乘法就想“是否在拼接对象”。
\end{explainblock}


\begin{sourceblock}{例 2.3.5 设 pk (n) 表示 n 元集合 k-可重排列的个数, 则 pk (n) = nk , 故序列}
{\small\sloppy
\noindent 例 2.3.5 设 pk (n) 表示 n 元集合 k-可重排列的个数, 则 pk (n) = nk , 故序列\par
\noindent \{pk (n)\}k≥0 的指数型生成函数为\par
\noindent X xk X xk\par
\noindent pk (n) = nk = enx .\par
\noindent k! k!\par
\noindent k≥0 k≥0\par
\noindent 从另一角度考虑, 不妨设集合 A = \{a1 , a2 , . . . , an \}. 设 π 是集合 A 的任意一个 k-可\par
\noindent 重排列, 若 π 中 ai (1 ≤ i ≤ n) 出现次数为 ki (ki ∈ N0 ), 则\par
\noindent k1 + k2 + · · · + kn = k, (2.7)\par
\noindent 其排列数为\par
\noindent k!\par
\noindent .\par
\noindent k1 !k2 ! · · · kn !\par
\noindent 因此, 集合 A 的 k-可重排列的个数为\par
\noindent X k!\par
\noindent pk = ,\par
\noindent k +k +···+k =k\par
\noindent k1 !k2 ! · · · kn !\par
}
\end{sourceblock}


\begin{explainblock}
例题是学习这本书最重要的部分。不要只看答案，要把每个限制条件变成一个计数动作：先安排谁、再插入谁、有没有重复、需不需要除以对称性。

\smallskip
\textbf{初学者检查点：}
\begin{itemize}[leftmargin=2em,itemsep=0.25em]
\item 先写出没有限制时的总数。
\item 逐条加入限制条件，观察总数怎样变化。
\item 最后检查有没有把同一种方案重复算了多次。
\end{itemize}

\smallskip
\textbf{本章读法提醒：} 第二章看见数列就想“它的生成函数是什么”，看见乘法就想“是否在拼接对象”。
\end{explainblock}


\begin{sourceblock}{第 2 章续读}
{\small\sloppy
\noindent ki ∈N0\par
\noindent 它恰好是 n\par
\noindent x x2 x3\par
\noindent 1+ + + + ··· = enx\par
\noindent 1! 2! 3!\par
\noindent k\par
\noindent 展开式中 xk! 的系数, 所以\par
\noindent ∞\par
\noindent X\par
\noindent nx xk\par
\noindent e = pk ,\par
\noindent k!\par
\noindent k=0\par
\noindent 这和从排列个数 nk 直接得到指数型生成函数完全一致.\par
}
\end{sourceblock}


\begin{explainblock}
这一段是原书论述链的一部分。读的时候把它当作桥梁：它通常负责把上一个定义、例子或定理，引到下一个计数方法。

\smallskip
\textbf{初学者检查点：}
\begin{itemize}[leftmargin=2em,itemsep=0.25em]
\item 找出这一段承接的上一个概念。
\item 找出它准备引出的下一个概念。
\item 把中间的关键等式或构造单独抄出来。
\end{itemize}

\smallskip
\textbf{本章读法提醒：} 第二章看见数列就想“它的生成函数是什么”，看见乘法就想“是否在拼接对象”。
\end{explainblock}


\begin{sourceblock}{例 2.3.6 用生成函数的方法求集合 [n] 的 k-可重排列中每个元素都至少出现一}
{\small\sloppy
\noindent 例 2.3.6 用生成函数的方法求集合 [n] 的 k-可重排列中每个元素都至少出现一\par
\noindent 次的排列数 qk .\par
\noindent 解: 考虑序列 \{qk \}k≥0 的指数型生成函数为\par
\noindent ∞\par
\noindent X n\par
\noindent xk x x2 x3\par
\noindent qk = + + + ···\par
\noindent k! 1! 2! 3!\par
\noindent k=0\par
\noindent = (ex − 1)n\par
\noindent X\par
\noindent n\par
\noindent i n\par
\noindent = (−1) e(n−i)x\par
\noindent i\par
\noindent i=0\par
\noindent X\par
\noindent n X\par
}
\end{sourceblock}


\begin{explainblock}
例题是学习这本书最重要的部分。不要只看答案，要把每个限制条件变成一个计数动作：先安排谁、再插入谁、有没有重复、需不需要除以对称性。

\smallskip
\textbf{初学者检查点：}
\begin{itemize}[leftmargin=2em,itemsep=0.25em]
\item 先写出没有限制时的总数。
\item 逐条加入限制条件，观察总数怎样变化。
\item 最后检查有没有把同一种方案重复算了多次。
\end{itemize}

\smallskip
\textbf{本章读法提醒：} 第二章看见数列就想“它的生成函数是什么”，看见乘法就想“是否在拼接对象”。
\end{explainblock}


\begin{sourceblock}{第 2 章续读}
{\small\sloppy
\noindent ∞\par
\noindent i n (n − i)k xk\par
\noindent = (−1)\par
\noindent i k!\par
\noindent i=0 k=0\par
\noindent ∞ X\par
\noindent X n\par
\noindent i n xk\par
\noindent = (−1) (n − i)k\par
\noindent i k!\par
\noindent k=0 i=0\par
\noindent 比较两边系数可得\par
\noindent X\par
\noindent n\par
\noindent i n\par
\noindent qk = (−1) (n − i)k .\par
\noindent i\par
\noindent i=0\par
}
\end{sourceblock}


\begin{explainblock}
这一段是原书论述链的一部分。读的时候把它当作桥梁：它通常负责把上一个定义、例子或定理，引到下一个计数方法。

\smallskip
\textbf{初学者检查点：}
\begin{itemize}[leftmargin=2em,itemsep=0.25em]
\item 找出这一段承接的上一个概念。
\item 找出它准备引出的下一个概念。
\item 把中间的关键等式或构造单独抄出来。
\end{itemize}

\smallskip
\textbf{本章读法提醒：} 第二章看见数列就想“它的生成函数是什么”，看见乘法就想“是否在拼接对象”。
\end{explainblock}


\begin{sourceblock}{注记 2.3.1 例 2.3.6 可利用容斥原理加以证明, 详见例 3.2.3.}
{\small\sloppy
\noindent 注记 2.3.1 例 2.3.6 可利用容斥原理加以证明, 详见例 3.2.3.\par
}
\end{sourceblock}


\begin{explainblock}
注记通常是在补充术语、记号、历史或一个容易忽略的约定。它未必给新公式，但常常决定你后面会不会把符号读错。

\smallskip
\textbf{初学者检查点：}
\begin{itemize}[leftmargin=2em,itemsep=0.25em]
\item 记下新符号的读法。
\item 确认这个约定是否只在本章使用。
\item 把注记里的限制条件补到前面的定义旁边。
\end{itemize}

\smallskip
\textbf{本章读法提醒：} 第二章看见数列就想“它的生成函数是什么”，看见乘法就想“是否在拼接对象”。
\end{explainblock}


\begin{sourceblock}{例 2.3.7 利用生成函数证明:}
{\small\sloppy
\noindent 例 2.3.7 利用生成函数证明:\par
\noindent n(n + 1)(2n + 1)\par
\noindent 12 + 2 2 + 3 2 + · · · + n2 = .\par
\noindent 证明: 不妨令\par
\noindent an = 12 + 22 + 32 + · · · + n2 ,\par
\noindent 则序列 \{an \}n≥0 的普通型生成函数为\par
\noindent ∞\par
\noindent X ∞ X\par
\noindent X n\par
\noindent n\par
\noindent an x = j 2 xn\par
\noindent n=0 n=1 j=1\par
\noindent ∞\par
\noindent X ∞\par
\noindent X\par
\noindent = j xn\par
\noindent j=1 n=j\par
\noindent ∞\par
}
\end{sourceblock}


\begin{explainblock}
例题是学习这本书最重要的部分。不要只看答案，要把每个限制条件变成一个计数动作：先安排谁、再插入谁、有没有重复、需不需要除以对称性。

\smallskip
\textbf{初学者检查点：}
\begin{itemize}[leftmargin=2em,itemsep=0.25em]
\item 先写出没有限制时的总数。
\item 逐条加入限制条件，观察总数怎样变化。
\item 最后检查有没有把同一种方案重复算了多次。
\end{itemize}

\smallskip
\textbf{本章读法提醒：} 第二章看见数列就想“它的生成函数是什么”，看见乘法就想“是否在拼接对象”。
\end{explainblock}


\begin{sourceblock}{第 2 章续读}
{\small\sloppy
\noindent X ∞\par
\noindent X\par
\noindent = j x xk\par
\noindent j=1 k=0\par
\noindent ∞\par
\noindent = j x .\par
\noindent 1−x\par
\noindent j=1\par
\noindent P∞ 2 j\par
\noindent 下面求 j=1 j x 的闭形式.\par
\noindent ∞\par
\noindent X ∞\par
\noindent X ∞\par
\noindent X\par
\noindent j 2 xj = (j 2 + j)xj − jxj\par
\noindent j=1 j=1 j=1\par
\noindent ∞\par
\noindent X ∞\par
}
\end{sourceblock}


\begin{explainblock}
这一段是原书论述链的一部分。读的时候把它当作桥梁：它通常负责把上一个定义、例子或定理，引到下一个计数方法。

\smallskip
\textbf{初学者检查点：}
\begin{itemize}[leftmargin=2em,itemsep=0.25em]
\item 找出这一段承接的上一个概念。
\item 找出它准备引出的下一个概念。
\item 把中间的关键等式或构造单独抄出来。
\end{itemize}

\smallskip
\textbf{本章读法提醒：} 第二章看见数列就想“它的生成函数是什么”，看见乘法就想“是否在拼接对象”。
\end{explainblock}


\begin{sourceblock}{第 2 章续读}
{\small\sloppy
\noindent j+1 j X j\par
\noindent =2 x − jx\par
\noindent j=1 j=1\par
\noindent ∞\par
\noindent X ∞\par
\noindent X\par
\noindent −3 −2\par
\noindent =2 (−1) j−1\par
\noindent x −\par
\noindent j\par
\noindent (−1)j−1\par
\noindent xj\par
\noindent j−1 j−1\par
\noindent j=1 j=1\par
\noindent = 2x(1 − x)−3 − x(1 − x)−2\par
\noindent x(1 + x)\par
\noindent = ,\par
\noindent (1 − x)3\par
}
\end{sourceblock}


\begin{explainblock}
这一段是原书论述链的一部分。读的时候把它当作桥梁：它通常负责把上一个定义、例子或定理，引到下一个计数方法。

\smallskip
\textbf{初学者检查点：}
\begin{itemize}[leftmargin=2em,itemsep=0.25em]
\item 找出这一段承接的上一个概念。
\item 找出它准备引出的下一个概念。
\item 把中间的关键等式或构造单独抄出来。
\end{itemize}

\smallskip
\textbf{本章读法提醒：} 第二章看见数列就想“它的生成函数是什么”，看见乘法就想“是否在拼接对象”。
\end{explainblock}


\begin{sourceblock}{第 2 章续读}
{\small\sloppy
\noindent 故\par
\noindent ∞\par
\noindent X 1 x(1 + x)\par
\noindent an xn = ·\par
\noindent n=0\par
\noindent x(1 + x)\par
\noindent =\par
\noindent (1 − x)4\par
\noindent ∞\par
\noindent X\par
\noindent = (x + x) (−1) x\par
\noindent n\par
\noindent n=0\par
\noindent ∞\par
\noindent X\par
\noindent = (x + x) x\par
\noindent n\par
\noindent n=0\par
}
\end{sourceblock}


\begin{explainblock}
这一段是原书论述链的一部分。读的时候把它当作桥梁：它通常负责把上一个定义、例子或定理，引到下一个计数方法。

\smallskip
\textbf{初学者检查点：}
\begin{itemize}[leftmargin=2em,itemsep=0.25em]
\item 找出这一段承接的上一个概念。
\item 找出它准备引出的下一个概念。
\item 把中间的关键等式或构造单独抄出来。
\end{itemize}

\smallskip
\textbf{本章读法提醒：} 第二章看见数列就想“它的生成函数是什么”，看见乘法就想“是否在拼接对象”。
\end{explainblock}


\begin{sourceblock}{第 2 章续读}
{\small\sloppy
\noindent ∞\par
\noindent X ∞\par
\noindent X\par
\noindent n+3 n+2 n + 3 n+1\par
\noindent = x + x\par
\noindent n n\par
\noindent n=0 n=0\par
\noindent ∞\par
\noindent X ∞\par
\noindent n+1 n X n+2 n\par
\noindent = x + x\par
\noindent n=2 n=1\par
\noindent ∞\par
\noindent X\par
\noindent n+1 n+2\par
\noindent = + xn\par
\noindent n=1\par
\noindent ∞\par
}
\end{sourceblock}


\begin{explainblock}
这一段是原书论述链的一部分。读的时候把它当作桥梁：它通常负责把上一个定义、例子或定理，引到下一个计数方法。

\smallskip
\textbf{初学者检查点：}
\begin{itemize}[leftmargin=2em,itemsep=0.25em]
\item 找出这一段承接的上一个概念。
\item 找出它准备引出的下一个概念。
\item 把中间的关键等式或构造单独抄出来。
\end{itemize}

\smallskip
\textbf{本章读法提醒：} 第二章看见数列就想“它的生成函数是什么”，看见乘法就想“是否在拼接对象”。
\end{explainblock}


\begin{sourceblock}{第 2 章续读}
{\small\sloppy
\noindent X n(n + 1)(2n + 1)\par
\noindent = xn ,\par
\noindent n=1\par
\noindent 比较两边 xn 的系数即可.\par
}
\end{sourceblock}


\begin{explainblock}
这一段是原书论述链的一部分。读的时候把它当作桥梁：它通常负责把上一个定义、例子或定理，引到下一个计数方法。

\smallskip
\textbf{初学者检查点：}
\begin{itemize}[leftmargin=2em,itemsep=0.25em]
\item 找出这一段承接的上一个概念。
\item 找出它准备引出的下一个概念。
\item 把中间的关键等式或构造单独抄出来。
\end{itemize}

\smallskip
\textbf{本章读法提醒：} 第二章看见数列就想“它的生成函数是什么”，看见乘法就想“是否在拼接对象”。
\end{explainblock}


\begin{sourceblock}{习 题 二}
{\small\sloppy
\noindent 习 题 二\par
\noindent \}n≥0 的指数型生成函数是 e x−1 ;\par
\noindent x\par
}
\end{sourceblock}


\begin{explainblock}
习题段不要先追求技巧。先给题目分类：它是在问排列、组合、递推、生成函数、容斥，还是结构计数。分类正确，工具通常就只剩一两个候选。

\smallskip
\textbf{初学者检查点：}
\begin{itemize}[leftmargin=2em,itemsep=0.25em]
\item 判断题目所属工具。
\item 列出变量和限制条件。
\item 先做小规模 n=2、3、4 的例子，确认直觉。
\end{itemize}

\smallskip
\textbf{本章读法提醒：} 第二章看见数列就想“它的生成函数是什么”，看见乘法就想“是否在拼接对象”。
\end{explainblock}


\begin{sourceblock}{习题 2.1 (1) 验证序列 \{ n+1}
{\small\sloppy
\noindent 习题 2.1 (1) 验证序列 \{ n+1\par
\noindent P\par
\noindent 的指数型生成函数是 (e x−1)\par
\noindent x 2\par
\noindent (2) 利用 (1) 证明序列 nk=0 (n−k+1)!(k+1)!\par
\noindent n!\par
\noindent Pn n!\par
\noindent (3) 利用 (2) 求出 k=0 (n−k+1)!(k+1)! 闭形式.\par
}
\end{sourceblock}


\begin{explainblock}
习题段不要先追求技巧。先给题目分类：它是在问排列、组合、递推、生成函数、容斥，还是结构计数。分类正确，工具通常就只剩一两个候选。

\smallskip
\textbf{初学者检查点：}
\begin{itemize}[leftmargin=2em,itemsep=0.25em]
\item 判断题目所属工具。
\item 列出变量和限制条件。
\item 先做小规模 n=2、3、4 的例子，确认直觉。
\end{itemize}

\smallskip
\textbf{本章读法提醒：} 第二章看见数列就想“它的生成函数是什么”，看见乘法就想“是否在拼接对象”。
\end{explainblock}


\begin{sourceblock}{习题 2.2 设 f (x) 是 \{an \}∞}
{\small\sloppy
\noindent 习题 2.2 设 f (x) 是 \{an \}∞\par
\noindent n=0 的普通型生成函数, 问: 哪个序列的普通型生成\par
\noindent (x)\par
\noindent 函数是 f1−x ?\par
}
\end{sourceblock}


\begin{explainblock}
习题段不要先追求技巧。先给题目分类：它是在问排列、组合、递推、生成函数、容斥，还是结构计数。分类正确，工具通常就只剩一两个候选。

\smallskip
\textbf{初学者检查点：}
\begin{itemize}[leftmargin=2em,itemsep=0.25em]
\item 判断题目所属工具。
\item 列出变量和限制条件。
\item 先做小规模 n=2、3、4 的例子，确认直觉。
\end{itemize}

\smallskip
\textbf{本章读法提醒：} 第二章看见数列就想“它的生成函数是什么”，看见乘法就想“是否在拼接对象”。
\end{explainblock}


\begin{sourceblock}{习题 2.3 求幂级数 f (x) = (1−x 2 )2 中 x}
{\small\sloppy
\noindent 习题 2.3 求幂级数 f (x) = (1−x 2 )2 中 x\par
\noindent n 的系数.\par
}
\end{sourceblock}


\begin{explainblock}
习题段不要先追求技巧。先给题目分类：它是在问排列、组合、递推、生成函数、容斥，还是结构计数。分类正确，工具通常就只剩一两个候选。

\smallskip
\textbf{初学者检查点：}
\begin{itemize}[leftmargin=2em,itemsep=0.25em]
\item 判断题目所属工具。
\item 列出变量和限制条件。
\item 先做小规模 n=2、3、4 的例子，确认直觉。
\end{itemize}

\smallskip
\textbf{本章读法提醒：} 第二章看见数列就想“它的生成函数是什么”，看见乘法就想“是否在拼接对象”。
\end{explainblock}


\begin{sourceblock}{习题 2.4 证明序列 \{0, 1, 8, 27, . . . , k 3 , . . . , \} 的普通型生成函数是}
{\small\sloppy
\noindent 习题 2.4 证明序列 \{0, 1, 8, 27, . . . , k 3 , . . . , \} 的普通型生成函数是\par
\noindent x(x2 + 4x + 1)\par
\noindent .\par
\noindent (1 − x)4\par
}
\end{sourceblock}


\begin{explainblock}
习题段不要先追求技巧。先给题目分类：它是在问排列、组合、递推、生成函数、容斥，还是结构计数。分类正确，工具通常就只剩一两个候选。

\smallskip
\textbf{初学者检查点：}
\begin{itemize}[leftmargin=2em,itemsep=0.25em]
\item 判断题目所属工具。
\item 列出变量和限制条件。
\item 先做小规模 n=2、3、4 的例子，确认直觉。
\end{itemize}

\smallskip
\textbf{本章读法提醒：} 第二章看见数列就想“它的生成函数是什么”，看见乘法就想“是否在拼接对象”。
\end{explainblock}


\begin{sourceblock}{习题 2.5 用生成函数的方法求不定方程}
{\small\sloppy
\noindent 习题 2.5 用生成函数的方法求不定方程\par
\noindent x1 + x2 + x3 + x4 + x5 + x6 = 30\par
\noindent 满足下列条件的整数解的个数.\par
\noindent (1)\par
\noindent xi ≥ 0, i = 1, 2, 3, 4, 5, 6;\par
\noindent (2)\par
\noindent xi ≥ i, i = 1, 2, 3, 4, 5, 6.\par
}
\end{sourceblock}


\begin{explainblock}
习题段不要先追求技巧。先给题目分类：它是在问排列、组合、递推、生成函数、容斥，还是结构计数。分类正确，工具通常就只剩一两个候选。

\smallskip
\textbf{初学者检查点：}
\begin{itemize}[leftmargin=2em,itemsep=0.25em]
\item 判断题目所属工具。
\item 列出变量和限制条件。
\item 先做小规模 n=2、3、4 的例子，确认直觉。
\end{itemize}

\smallskip
\textbf{本章读法提醒：} 第二章看见数列就想“它的生成函数是什么”，看见乘法就想“是否在拼接对象”。
\end{explainblock}


\begin{sourceblock}{习题 2.6 设 n, k 均为正整数, 用生成函数的方法证明方程}
{\small\sloppy
\noindent 习题 2.6 设 n, k 均为正整数, 用生成函数的方法证明方程\par
\noindent x1 + x2 + · · · + xk + xk+1 = n\par
\noindent 与方程\par
\noindent y1 + y2 + · · · + yn + yn+1 = k\par
\noindent 的非负整数解的个数相同.\par
}
\end{sourceblock}


\begin{explainblock}
习题段不要先追求技巧。先给题目分类：它是在问排列、组合、递推、生成函数、容斥，还是结构计数。分类正确，工具通常就只剩一两个候选。

\smallskip
\textbf{初学者检查点：}
\begin{itemize}[leftmargin=2em,itemsep=0.25em]
\item 判断题目所属工具。
\item 列出变量和限制条件。
\item 先做小规模 n=2、3、4 的例子，确认直觉。
\end{itemize}

\smallskip
\textbf{本章读法提醒：} 第二章看见数列就想“它的生成函数是什么”，看见乘法就想“是否在拼接对象”。
\end{explainblock}


\begin{sourceblock}{习题 2.7 用生成函数的方法证明下列等式:}
{\small\sloppy
\noindent 习题 2.7 用生成函数的方法证明下列等式:\par
\noindent (1)\par
\noindent n+2 n+1 n n\par
\noindent −2 + = ,\par
\noindent k k k k−2\par
\noindent (2)\par
\noindent X\par
\noindent k\par
\noindent j k n+k−j n\par
\noindent (−1) = .\par
\noindent j m m−k\par
\noindent j=0\par
}
\end{sourceblock}


\begin{explainblock}
习题段不要先追求技巧。先给题目分类：它是在问排列、组合、递推、生成函数、容斥，还是结构计数。分类正确，工具通常就只剩一两个候选。

\smallskip
\textbf{初学者检查点：}
\begin{itemize}[leftmargin=2em,itemsep=0.25em]
\item 判断题目所属工具。
\item 列出变量和限制条件。
\item 先做小规模 n=2、3、4 的例子，确认直觉。
\end{itemize}

\smallskip
\textbf{本章读法提醒：} 第二章看见数列就想“它的生成函数是什么”，看见乘法就想“是否在拼接对象”。
\end{explainblock}


\begin{sourceblock}{习题 2.8 设集合 A = \{1, 2, 3, 4, 5\}, 求 A 的 k-可重排列中 1 与 2 出现的次数}
{\small\sloppy
\noindent 习题 2.8 设集合 A = \{1, 2, 3, 4, 5\}, 求 A 的 k-可重排列中 1 与 2 出现的次数\par
\noindent 之和是偶数的排列的个数.\par
}
\end{sourceblock}


\begin{explainblock}
习题段不要先追求技巧。先给题目分类：它是在问排列、组合、递推、生成函数、容斥，还是结构计数。分类正确，工具通常就只剩一两个候选。

\smallskip
\textbf{初学者检查点：}
\begin{itemize}[leftmargin=2em,itemsep=0.25em]
\item 判断题目所属工具。
\item 列出变量和限制条件。
\item 先做小规模 n=2、3、4 的例子，确认直觉。
\end{itemize}

\smallskip
\textbf{本章读法提醒：} 第二章看见数列就想“它的生成函数是什么”，看见乘法就想“是否在拼接对象”。
\end{explainblock}


\begin{sourceblock}{习题 2.9 证明:}
{\small\sloppy
\noindent 习题 2.9 证明:\par
\noindent (1)\par
\noindent 1 n 1 n k1 n n+1 1 n n\par
\noindent − + · · · + (−1) + · · · + (−1) = ,\par
\noindent 2 1 3 2 k k−1 n+1 n n+1\par
\noindent (2)\par
\noindent 1 n 1 n 1 n 1 n 2n+1 − n − 2\par
\noindent + + ··· + + ··· + = .\par
\noindent 2 1 3 2 k k−1 n+1 n n+1\par
}
\end{sourceblock}


\begin{explainblock}
习题段不要先追求技巧。先给题目分类：它是在问排列、组合、递推、生成函数、容斥，还是结构计数。分类正确，工具通常就只剩一两个候选。

\smallskip
\textbf{初学者检查点：}
\begin{itemize}[leftmargin=2em,itemsep=0.25em]
\item 判断题目所属工具。
\item 列出变量和限制条件。
\item 先做小规模 n=2、3、4 的例子，确认直觉。
\end{itemize}

\smallskip
\textbf{本章读法提醒：} 第二章看见数列就想“它的生成函数是什么”，看见乘法就想“是否在拼接对象”。
\end{explainblock}


\begin{sourceblock}{习题 2.10 求下列和式:}
{\small\sloppy
\noindent 习题 2.10 求下列和式:\par
\noindent (1)\par
\noindent 14 + 24 + 34 + · · · + n4 ,\par
\noindent (2)\par
\noindent 1 · 2 + 2 · 3 + · · · + n(n + 1).\par
}
\end{sourceblock}


\begin{explainblock}
习题段不要先追求技巧。先给题目分类：它是在问排列、组合、递推、生成函数、容斥，还是结构计数。分类正确，工具通常就只剩一两个候选。

\smallskip
\textbf{初学者检查点：}
\begin{itemize}[leftmargin=2em,itemsep=0.25em]
\item 判断题目所属工具。
\item 列出变量和限制条件。
\item 先做小规模 n=2、3、4 的例子，确认直觉。
\end{itemize}

\smallskip
\textbf{本章读法提醒：} 第二章看见数列就想“它的生成函数是什么”，看见乘法就想“是否在拼接对象”。
\end{explainblock}


\begin{sourceblock}{习题 2.11 斐波那契 (Leonardo Fibonacci) 在 1202 年出版的 Liber Abaci 里}
{\small\sloppy
\noindent 习题 2.11 斐波那契 (Leonardo Fibonacci) 在 1202 年出版的 Liber Abaci 里\par
\noindent 提到一个有趣的问题: 假定一对刚出生的小兔一个月就能长成大兔, 再过一个月便能\par
\noindent 生下一对小兔, 并且此后每个月都生一对小兔. 若不考虑死亡问题, 设第 n(n ≥ 0) 个\par
\noindent 月时大兔的个数为 fn , 则 f0 = 0, f1 = 1.\par
\noindent (1) 证明 fn 满足下面的递推关系\par
\noindent fn+1 = fn + fn−1 .\par
\noindent (2) 求序列 \{fn \}n≥0 的普通型生成函数.\par
\noindent (3) 利用 (2) 求 fn 的表达式.\par
}
\end{sourceblock}


\begin{explainblock}
习题段不要先追求技巧。先给题目分类：它是在问排列、组合、递推、生成函数、容斥，还是结构计数。分类正确，工具通常就只剩一两个候选。

\smallskip
\textbf{初学者检查点：}
\begin{itemize}[leftmargin=2em,itemsep=0.25em]
\item 判断题目所属工具。
\item 列出变量和限制条件。
\item 先做小规模 n=2、3、4 的例子，确认直觉。
\end{itemize}

\smallskip
\textbf{本章读法提醒：} 第二章看见数列就想“它的生成函数是什么”，看见乘法就想“是否在拼接对象”。
\end{explainblock}


\begin{sourceblock}{习题 2.12 设 b0 = 0,}
{\small\sloppy
\noindent 习题 2.12 设 b0 = 0,\par
\noindent k k+1 k+j 2k\par
\noindent ak = + + ··· + + ··· + , k ≥ 0,\par
\noindent k k+1 k+j 2k − 1\par
\noindent bk = + + ··· + + ··· + , k > 0.\par
\noindent 它们的生成函数分别为\par
\noindent A(t) = a0 + a1 t + a2 t2 + · · · + an tn + · · ·\par
\noindent B(t) = b0 + b1 t + a2 t2 + · · · + bn tn + · · · .\par
\noindent (1) 证明\par
\noindent ak+1 = ak + bk+1 ,\par
\noindent bk+1 = ak + bk ;\par
\noindent (2) 求 A(t) 和 B(t).\par
}
\end{sourceblock}


\begin{explainblock}
习题段不要先追求技巧。先给题目分类：它是在问排列、组合、递推、生成函数、容斥，还是结构计数。分类正确，工具通常就只剩一两个候选。

\smallskip
\textbf{初学者检查点：}
\begin{itemize}[leftmargin=2em,itemsep=0.25em]
\item 判断题目所属工具。
\item 列出变量和限制条件。
\item 先做小规模 n=2、3、4 的例子，确认直觉。
\end{itemize}

\smallskip
\textbf{本章读法提醒：} 第二章看见数列就想“它的生成函数是什么”，看见乘法就想“是否在拼接对象”。
\end{explainblock}


\begin{sourceblock}{习题 2.13 设 α 是三次本原单位根, 因此 α3 = 1. 证明}
{\small\sloppy
\noindent 习题 2.13 设 α 是三次本原单位根, 因此 α3 = 1. 证明\par
\noindent 2 + (1 + α)k α2m + (1 + α2 )k αm\par
\noindent k k k\par
\noindent = + + + ··· , m = 0, 1, 2.\par
\noindent m m+3 m+6\par
\noindent 并验证\par
\noindent k\par
\noindent 2 + 2 cos = + + + ··· ,\par
\noindent k\par
\noindent 2 + 2 cos = + + + ··· ,\par
\noindent k\par
\noindent 2 + 2 cos = + + + ··· .\par
}
\end{sourceblock}


\begin{explainblock}
习题段不要先追求技巧。先给题目分类：它是在问排列、组合、递推、生成函数、容斥，还是结构计数。分类正确，工具通常就只剩一两个候选。

\smallskip
\textbf{初学者检查点：}
\begin{itemize}[leftmargin=2em,itemsep=0.25em]
\item 判断题目所属工具。
\item 列出变量和限制条件。
\item 先做小规模 n=2、3、4 的例子，确认直觉。
\end{itemize}

\smallskip
\textbf{本章读法提醒：} 第二章看见数列就想“它的生成函数是什么”，看见乘法就想“是否在拼接对象”。
\end{explainblock}


\begin{sourceblock}{习题 2.14 设 a0 = a1 = 1, a2 = 2, 且}
{\small\sloppy
\noindent 习题 2.14 设 a0 = a1 = 1, a2 = 2, 且\par
\noindent n\par
\noindent an+1 = (n + 1)an − an−2 , n > 1.\par
\noindent 证明它的指数型生成函数:\par
\noindent A(t) = a0 + a1 t + a2 t2 /2! + · · ·\par
\noindent 满足下面的微分方程:\par
\noindent dA 1 2\par
\noindent (1 − t) = 1 − t A(t),\par
\noindent dt 2\par
\noindent 并求 A(t).\par
}
\end{sourceblock}


\begin{explainblock}
习题段不要先追求技巧。先给题目分类：它是在问排列、组合、递推、生成函数、容斥，还是结构计数。分类正确，工具通常就只剩一两个候选。

\smallskip
\textbf{初学者检查点：}
\begin{itemize}[leftmargin=2em,itemsep=0.25em]
\item 判断题目所属工具。
\item 列出变量和限制条件。
\item 先做小规模 n=2、3、4 的例子，确认直觉。
\end{itemize}

\smallskip
\textbf{本章读法提醒：} 第二章看见数列就想“它的生成函数是什么”，看见乘法就想“是否在拼接对象”。
\end{explainblock}


\begin{sourceblock}{习题 2.15 请用指数型生成函数来证明}
{\small\sloppy
\noindent 习题 2.15 请用指数型生成函数来证明\par
\noindent X n!\par
\noindent (x1 + x2 + x3 )n = xr11 xr22 xr33 .\par
\noindent r +r +r =n\par
\noindent r1 !r2 !r3 !\par
}
\end{sourceblock}


\begin{explainblock}
习题段不要先追求技巧。先给题目分类：它是在问排列、组合、递推、生成函数、容斥，还是结构计数。分类正确，工具通常就只剩一两个候选。

\smallskip
\textbf{初学者检查点：}
\begin{itemize}[leftmargin=2em,itemsep=0.25em]
\item 判断题目所属工具。
\item 列出变量和限制条件。
\item 先做小规模 n=2、3、4 的例子，确认直觉。
\end{itemize}

\smallskip
\textbf{本章读法提醒：} 第二章看见数列就想“它的生成函数是什么”，看见乘法就想“是否在拼接对象”。
\end{explainblock}


\begin{sourceblock}{习题 2.16 求由直线 x + 3y = 12, 直线 x = 0 以及直线 y = 0 所围成的三角形}
{\small\sloppy
\noindent 习题 2.16 求由直线 x + 3y = 12, 直线 x = 0 以及直线 y = 0 所围成的三角形\par
\noindent 内 (包括边界) 的整点个数 (即横坐标与纵坐标均为整数的点).\par
}
\end{sourceblock}


\begin{explainblock}
习题段不要先追求技巧。先给题目分类：它是在问排列、组合、递推、生成函数、容斥，还是结构计数。分类正确，工具通常就只剩一两个候选。

\smallskip
\textbf{初学者检查点：}
\begin{itemize}[leftmargin=2em,itemsep=0.25em]
\item 判断题目所属工具。
\item 列出变量和限制条件。
\item 先做小规模 n=2、3、4 的例子，确认直觉。
\end{itemize}

\smallskip
\textbf{本章读法提醒：} 第二章看见数列就想“它的生成函数是什么”，看见乘法就想“是否在拼接对象”。
\end{explainblock}


\begin{sourceblock}{习题 2.17 从一个 n 元排列中选取 k 个元做组合, 使得组合中的任何两个元 a,}
{\small\sloppy
\noindent 习题 2.17 从一个 n 元排列中选取 k 个元做组合, 使得组合中的任何两个元 a,\par
\noindent b 满足条件: 在原排列中 a 和 b 之间至少有 r(n ≥ (k − 1)r + k) 个元, 以 fr (n, k) 表\par
\noindent 示做成的不同组合的个数. 也即是说考虑 [n] 的任意一个排列 a1 a2 · · · an , 从中选取\par
\noindent k 个元素, 不妨设 \{ai1 , ai2 , . . . , aik \} 其中 1 ≤ i1 < i2 < · · · < ik ≤ n, 则 fr (n, k) 表\par
\noindent 示满足 \par
\noindent \par
\noindent i ≥ 1,\par
\noindent 1\par
\noindent ij+1 − i1 − 1 ≥ r, j = 1, 2, . . . , k − 1,\par
\noindent \par
\noindent \par
\noindent \par
\noindent ik ≤ n,\par
\noindent 的正整数解的个数. 比如 n = 4, k = 2, r = 1 时, 设从 [4] 的排列中选取排列\par
\noindent a1 a2 a3 a4 , 从 a1 , a2 , a3 , a4 取两个元素满足要求的只有 \{a1 , a3 \}, \{a1 , a4 \} 和 \{a2 , a4 \},\par
\noindent 所以 f1 (4, 2) = 3. 请求 fr (n, k).\par
}
\end{sourceblock}


\begin{explainblock}
习题段不要先追求技巧。先给题目分类：它是在问排列、组合、递推、生成函数、容斥，还是结构计数。分类正确，工具通常就只剩一两个候选。

\smallskip
\textbf{初学者检查点：}
\begin{itemize}[leftmargin=2em,itemsep=0.25em]
\item 判断题目所属工具。
\item 列出变量和限制条件。
\item 先做小规模 n=2、3、4 的例子，确认直觉。
\end{itemize}

\smallskip
\textbf{本章读法提醒：} 第二章看见数列就想“它的生成函数是什么”，看见乘法就想“是否在拼接对象”。
\end{explainblock}


\begin{sourceblock}{习题 2.18 设序列 \{an \}n≥0 满足下面的递推关系:}
{\small\sloppy
\noindent 习题 2.18 设序列 \{an \}n≥0 满足下面的递推关系:\par
\noindent an = 5an−1 − 6an−2 + 4n−1 , n ≥ 2,\par
\noindent 其初值为 a0 = 1, a1 = 3. 求 an 的表达式.\par
\noindent 应用公式 (1 − 4t)−1 = [(1 − 4t)− 2 ]2 证明：\par
}
\end{sourceblock}


\begin{explainblock}
习题段不要先追求技巧。先给题目分类：它是在问排列、组合、递推、生成函数、容斥，还是结构计数。分类正确，工具通常就只剩一两个候选。

\smallskip
\textbf{初学者检查点：}
\begin{itemize}[leftmargin=2em,itemsep=0.25em]
\item 判断题目所属工具。
\item 列出变量和限制条件。
\item 先做小规模 n=2、3、4 的例子，确认直觉。
\end{itemize}

\smallskip
\textbf{本章读法提醒：} 第二章看见数列就想“它的生成函数是什么”，看见乘法就想“是否在拼接对象”。
\end{explainblock}


\begin{sourceblock}{习题 2.19}
{\small\sloppy
\noindent 习题 2.19\par
\noindent n\par
\noindent X\par
\noindent n 2k 2n − 2k\par
\noindent k n−k\par
\noindent k=0\par
}
\end{sourceblock}


\begin{explainblock}
习题段不要先追求技巧。先给题目分类：它是在问排列、组合、递推、生成函数、容斥，还是结构计数。分类正确，工具通常就只剩一两个候选。

\smallskip
\textbf{初学者检查点：}
\begin{itemize}[leftmargin=2em,itemsep=0.25em]
\item 判断题目所属工具。
\item 列出变量和限制条件。
\item 先做小规模 n=2、3、4 的例子，确认直觉。
\end{itemize}

\smallskip
\textbf{本章读法提醒：} 第二章看见数列就想“它的生成函数是什么”，看见乘法就想“是否在拼接对象”。
\end{explainblock}


\begin{sourceblock}{习题 2.20 证明:}
{\small\sloppy
\noindent 习题 2.20 证明:\par
\noindent (1)\par
\noindent n 2 n 2 n 2 n\par
\noindent +2 +3 + ··· + n = n(n + 1)2n−2 ,\par
\noindent (2)\par
\noindent n 2 n 2 n n−1 2 n\par
\noindent −2 +3 − · · · + (−1) n = 0, n > 2.\par
}
\end{sourceblock}


\begin{explainblock}
习题段不要先追求技巧。先给题目分类：它是在问排列、组合、递推、生成函数、容斥，还是结构计数。分类正确，工具通常就只剩一两个候选。

\smallskip
\textbf{初学者检查点：}
\begin{itemize}[leftmargin=2em,itemsep=0.25em]
\item 判断题目所属工具。
\item 列出变量和限制条件。
\item 先做小规模 n=2、3、4 的例子，确认直觉。
\end{itemize}

\smallskip
\textbf{本章读法提醒：} 第二章看见数列就想“它的生成函数是什么”，看见乘法就想“是否在拼接对象”。
\end{explainblock}


\chapter{容斥原理及其应用}
\begin{roadmapblock}
本章保留原书正文的主要数学骨架，并在每个定义、定理、例题和论述片段后加入解读。
阅读时不要把目标放在“背下答案”上，而是放在“看懂这一步为什么这样数”上。

\smallskip
\textbf{本章最低目标：} 能说出每个公式左边在数什么、右边在数什么，以及为什么两边相等。
\end{roadmapblock}


\begin{sourceblock}{第 3 章正文}
{\small\sloppy
\noindent 容斥原理及其应用\par
\noindent 前一章介绍了组合对象的计数问题中常用的生成函数的方法及应用, 这一章我们\par
\noindent 介绍另一个非常有用的古典的计数工具—容斥原理. 粗略地讲，容斥原理就是组合\par
\noindent 计数学中的筛法, 它是一种首先从一个较大的集合入手, 以某种方式减掉或删去一些\par
\noindent 不需要的元素, 从而计算出集合基数的方法, 是组合计数学的主要工具之一. 从抽象\par
\noindent 的角度来讲, 容斥原理就是计算某个矩阵的逆, 因此, 它只不过是线性代数中一个很\par
\noindent 平凡的结果, 但这个原理的优美之处远不在结果本身, 而在于它应用的广泛性. 本章\par
\noindent 主要介绍容斥原理及其在组合与数论上的简单应用.\par
}
\end{sourceblock}


\begin{explainblock}
这一段是原书论述链的一部分。读的时候把它当作桥梁：它通常负责把上一个定义、例子或定理，引到下一个计数方法。

\smallskip
\textbf{初学者检查点：}
\begin{itemize}[leftmargin=2em,itemsep=0.25em]
\item 找出这一段承接的上一个概念。
\item 找出它准备引出的下一个概念。
\item 把中间的关键等式或构造单独抄出来。
\end{itemize}

\smallskip
\textbf{本章读法提醒：} 第三章先定义坏事件，再做容斥；坏事件定义不清楚，公式没有意义。
\end{explainblock}


\begin{sourceblock}{§ 3.1 容斥原理}
{\small\sloppy
\noindent § 3.1. 容斥原理\par
\noindent 很多的计数问题都可以抽象成下面的问题: 已知有限的集合 S 及与 S 中元素相\par
\noindent 关的性质集 P = \{P1 , P2 , . . . , Pm \}. 问: S 恰好具有 P 中某 k 个性质的元素的个数.\par
\noindent 特别地, 问: S 不具有 P 中任一性质元素的个数.\par
\noindent 这正是容斥原理所研究的问题. 通常情形下, S 恰好具有 P 中某 k 个性质的元\par
\noindent 素的个数是很难计算的. 相较之下, S 中至少满足 P 中 r 个性质的元素的个数则容\par
\noindent 易直接计算, 而容斥原理的本质就是通过后者元素的个数来计算前者元素的个数.\par
\noindent 在具体介绍容斥原理之前, 我们先从简单的例子谈起.\par
}
\end{sourceblock}


\begin{explainblock}
这一节先不要急着背公式。先问三个问题：对象是什么，哪些限制条件不能丢，最后到底是在数所有对象、还是在数某个统计量的分布。把这三件事说清楚，后面的公式才会有位置。

\smallskip
\textbf{初学者检查点：}
\begin{itemize}[leftmargin=2em,itemsep=0.25em]
\item 把本节出现的新对象写成一句话定义。
\item 把每个公式左边和右边分别翻译成“在数什么”。
\item 找一个最小非平凡例子，手工列出所有对象。
\end{itemize}

\smallskip
\textbf{本章读法提醒：} 第三章先定义坏事件，再做容斥；坏事件定义不清楚，公式没有意义。
\end{explainblock}


\begin{sourceblock}{例 3.1.1 在由 15 人组成的出国访问的棋艺代表团中, 有 9 名围棋手, 7 名象棋}
{\small\sloppy
\noindent 例 3.1.1 在由 15 人组成的出国访问的棋艺代表团中, 有 9 名围棋手, 7 名象棋\par
\noindent 手. 其中有 3 人即是围棋手也是象棋手, 问此代表团中有几名非棋手?\par
\noindent 解: 因为即是围棋手又是象棋手的 3 名棋手在 9 名围棋手和 7 名象棋手中都被\par
\noindent 计算了一次, 所以在 9 + 7 = 16 这个和中, 这 3 名棋手被计算了两次, 也即是说, 从\par
\noindent 16 中减去多计算的 3 名就恰好是代表团中棋手的人数, 即 16 − 3 = 13 人, 所以代表\par
\noindent 团中不是棋手的人数为 15 − 13 = 2 名.\par
}
\end{sourceblock}


\begin{explainblock}
例题是学习这本书最重要的部分。不要只看答案，要把每个限制条件变成一个计数动作：先安排谁、再插入谁、有没有重复、需不需要除以对称性。

\smallskip
\textbf{初学者检查点：}
\begin{itemize}[leftmargin=2em,itemsep=0.25em]
\item 先写出没有限制时的总数。
\item 逐条加入限制条件，观察总数怎样变化。
\item 最后检查有没有把同一种方案重复算了多次。
\end{itemize}

\smallskip
\textbf{本章读法提醒：} 第三章先定义坏事件，再做容斥；坏事件定义不清楚，公式没有意义。
\end{explainblock}


\begin{sourceblock}{例 3.1.2 数学院大一某班共有学生 26 人, 其中参加象棋社的有 12 人, 参加围棋}
{\small\sloppy
\noindent 例 3.1.2 数学院大一某班共有学生 26 人, 其中参加象棋社的有 12 人, 参加围棋\par
\noindent 社及五子棋社的各有 9 人, 同时参加象棋社与围棋社的有 6 人, 同时参加象棋社与\par
\noindent 五子棋社和同时参加了五子棋社与围棋社的各有 5 人, 三个棋社都经参加的有 2 人.\par
\noindent 问该班中没有参加这三个棋社任意一个的人数有几人?\par
\noindent 解: 分别用 A象 , A围 , A五 表示参加象棋社, 围棋社和五子棋社的学生组成的集合,\par
\noindent 则\par
\noindent |A象 | = 12, |A围 | = |A五 | = 9,\par
\noindent |A象 ∩ A围 | = 6, |A象 ∩ A五 | = |A五 ∩ A围 | = 5,\par
\noindent 及\par
\noindent |A象 ∩ A围 ∩ A五 | = 2.\par
\noindent 为求出没参加这三个棋社中任何一个的人数, 只需求出至少参加这三个棋社中任何\par
\noindent 一个的人数, 为此先将参加这三个棋社的人数相加\par
\noindent |A象 | + |A围 | + |A五 | = 12 + 9 + 9 = 30, (3.1)\par
\noindent 这比全班的人数还要多, 但并不奇怪, 因为凡是同时只参加这三个棋社中任何两个棋\par
\noindent 社的人数均都计算了两次, 而三个棋社都参加的人数均都计算了三次, 所以要将多计\par
\noindent 算的给减去. 再考虑同时只参加这三个棋社中任何两个棋社的人数, 也是先将参加这\par
\noindent 三个棋社中两个棋社的人数相加\par
\noindent |A象 ∩ A围 | + |A象 ∩ A五 | + |A五 ∩ A围 | = 6 + 5 + 5 = 16. (3.2)\par
}
\end{sourceblock}


\begin{explainblock}
例题是学习这本书最重要的部分。不要只看答案，要把每个限制条件变成一个计数动作：先安排谁、再插入谁、有没有重复、需不需要除以对称性。

\smallskip
\textbf{初学者检查点：}
\begin{itemize}[leftmargin=2em,itemsep=0.25em]
\item 先写出没有限制时的总数。
\item 逐条加入限制条件，观察总数怎样变化。
\item 最后检查有没有把同一种方案重复算了多次。
\end{itemize}

\smallskip
\textbf{本章读法提醒：} 第三章先定义坏事件，再做容斥；坏事件定义不清楚，公式没有意义。
\end{explainblock}


\begin{sourceblock}{第 3 章续读}
{\small\sloppy
\noindent 上式结果不仅包含了只参加这三个棋社中任何两个棋社的人数, 而且同样将三个棋\par
\noindent 社都参加的人数均计算了三次. 因此 (3.1) 与 (3.2) 的差, 即\par
\noindent (|A象 | + |A围 | + |A五 |) − (|A象 ∩ A围 | + |A象 ∩ A五 | + |A五 ∩ A围 |) = 14\par
\noindent 将在计算参加这三个棋社的人数时多计算的那些只参加这三个棋社中任何两个棋社\par
\noindent 的人数恰好给减去了, 但是同时也减去了三个棋社都参加的人数, 因此需要再给补\par
\noindent 上, 即加上\par
\noindent |A象 ∩ A围 ∩ A五 | = 2.\par
\noindent 这样就得到了至少参加了这三个棋社任何一个的人数:\par
\noindent 所以该班上没参加这三个棋社中任意一个的人数为\par
\noindent 正如上面的两个简单的例子: 通过 “包含” 过多就 “排除”, “排除” 的过多再 “包\par
\noindent 含” 进去的原则, 反复交替地来求得准确的计数结果, 容斥原理就是前人在这样简单\par
\noindent 的原则上得到的一种重要的计数方法.\par
}
\end{sourceblock}


\begin{explainblock}
这一段是原书论述链的一部分。读的时候把它当作桥梁：它通常负责把上一个定义、例子或定理，引到下一个计数方法。

\smallskip
\textbf{初学者检查点：}
\begin{itemize}[leftmargin=2em,itemsep=0.25em]
\item 找出这一段承接的上一个概念。
\item 找出它准备引出的下一个概念。
\item 把中间的关键等式或构造单独抄出来。
\end{itemize}

\smallskip
\textbf{本章读法提醒：} 第三章先定义坏事件，再做容斥；坏事件定义不清楚，公式没有意义。
\end{explainblock}


\begin{sourceblock}{定理 3.1.1 (容斥原理) 设 S 是一个有限集合, P = \{P1 , P2 , . . . , Pm \} 是一组性质}
{\small\sloppy
\noindent 定理 3.1.1 (容斥原理) 设 S 是一个有限集合, P = \{P1 , P2 , . . . , Pm \} 是一组性质\par
\noindent 集, 其中 Pi 是同 S 有关某种性质. 设 Ai 是 S 中具有性质 Pi 的元素构成的集合\par
\noindent (1 ≤ i ≤ m), Ai 表示 S 中不具有性质 Pi 的元素构成的集合, 则 S 中不具有 P 中任\par
\noindent 何性质的元素的个数为:\par
\noindent X\par
\noindent m X X\par
\noindent |A1 ∩ A2 ∩ · · · ∩ Am | = |S| − |Ai | + |Ai ∩ Aj | − |Ai ∩ Aj ∩ Ak |\par
\noindent i=1 1≤i̸=j<m 1≤i<j<k≤m\par
\noindent + · · · + (−1)m |A1 ∩ A2 · · · ∩ Am |\par
\noindent X\par
\noindent m X\par
\noindent = |S| − (−1)k−1 |Ai1 ∩ · · · ∩ Aik |. (3.3)\par
\noindent k=1 1≤i1 <···<ik ≤m\par
\noindent 该定理可通过对性质集合的个数 k 作归纳进行证明, 这里给出另外两种不同的方\par
\noindent 法.\par
\noindent 证明: (方法一) 下面通过其组合意义来给出证明. 等式 (3.3) 左端表示的是 S\par
\noindent 中不具有 P 中任意一个性质的元素的个数, 下面将证明: 对 S 中的每一个元素 x,\par
\noindent 若 x 不具有 P 中任意一个性质则 x 对等式 (3.3) 右端的贡献为 1, 否则其贡献为 0,\par
}
\end{sourceblock}


\begin{explainblock}
定理段不要只看结论，要看它把哪两个计数方式连在一起。组合学证明最常见的技巧是：左边按一种标准数，右边按另一种标准数，然后说明两边数的是同一批对象。

\smallskip
\textbf{初学者检查点：}
\begin{itemize}[leftmargin=2em,itemsep=0.25em]
\item 先复述定理的假设，不要把适用范围扩大。
\item 找出证明中的基本计数原则：乘法、加法、双射、递推、容斥或生成函数。
\item 检查最后的等式化简是否和前面的对象解释一致。
\end{itemize}

\smallskip
\textbf{本章读法提醒：} 第三章先定义坏事件，再做容斥；坏事件定义不清楚，公式没有意义。
\end{explainblock}


\begin{sourceblock}{第 3 章续读}
{\small\sloppy
\noindent 从而证明等式 (3.3).\par
\noindent 设 x ∈ S, 即 x 对 S 的贡献为 1. 下面分两种情形讨论 x 对等式 (3.3) 右端其它\par
\noindent 集合的贡献:\par
\noindent (1) 若 x 不具有 P 中任意一个性质, 则对任意 i(i = 1, 2, . . . , m), x ∈\par
\noindent / Ai . 所以 x\par
\noindent 不在等式 (3.3) 右端除 S 外的任何其它集合中, 所以 x 对等式 (3.3) 的贡献为\par
\noindent 1.\par
\noindent P\par
\noindent (2) 若 x 具有 P 中任意 k(1 ≤ k ≤ m) 个性质, 则 x 对 |Ai | 的贡献为 k1 ; x 对\par
\noindent P P\par
\noindent |Ai ∩ Aj | 的贡献为 k2 . 一般地, x 对 |Ai1 ∩ · · · ∩ Ail | 的贡献实际上就是\par
\noindent k\par
\noindent 从 k 个性质选取 ℓ 个性质的方法数, 即 l . 因此, x 对等式 (3.3) 的贡献为\par
\noindent k k k\par
\noindent 1− + − · · · + (−1)m = 0.\par
\noindent 综上所述, 等式 (3.3) 恰好等于集合 S 中不具有 P 任何一个性质的元素的个数.\par
\noindent (方法二) (特征函数法) 对任意集合 A, 我们定义 χA 为 A 的特征函数 (charac-\par
\noindent teristic function) 或指示函数 (indicator function):\par
}
\end{sourceblock}


\begin{explainblock}
这一段是原书论述链的一部分。读的时候把它当作桥梁：它通常负责把上一个定义、例子或定理，引到下一个计数方法。

\smallskip
\textbf{初学者检查点：}
\begin{itemize}[leftmargin=2em,itemsep=0.25em]
\item 找出这一段承接的上一个概念。
\item 找出它准备引出的下一个概念。
\item 把中间的关键等式或构造单独抄出来。
\end{itemize}

\smallskip
\textbf{本章读法提醒：} 第三章先定义坏事件，再做容斥；坏事件定义不清楚，公式没有意义。
\end{explainblock}


\begin{sourceblock}{第 3 章续读}
{\small\sloppy
\noindent (\par
\noindent 1, 若 a ∈ A,\par
\noindent χA (a) = (3.4)\par
\noindent 0, 若 a ∈\par
\noindent / A.\par
\noindent 设 Ā 是 A 的补集, 则对任意的 a ∈ A,\par
\noindent χĀ (a) = 1 − χA (a).\par
\noindent 下面我们利用特征函数来证明该定理. 由特征函数的定义可知\par
\noindent X\par
\noindent |Ai | = χAi (s), 1 ≤ i ≤ m, (3.5)\par
\noindent s∈S\par
\noindent 一般地, 对任意 1 ≤ i1 < i2 < · · · < ik ≤ m(1 ≤ k ≤ m) 也有:\par
\noindent X Y\par
\noindent |Ai1 ∩ Ai2 ∩ · · · ∩ Aik | = χAiℓ (s). (3.6)\par
\noindent s∈S 1≤ℓ≤k\par
\noindent 因此 S 中不具有 P 任一个性质的元素个数为:\par
\noindent X Y\par
\noindent |Ā1 ∩ Ā2 ∩ · · · ∩ Ām | = χĀi (s)\par
}
\end{sourceblock}


\begin{explainblock}
这一段是原书论述链的一部分。读的时候把它当作桥梁：它通常负责把上一个定义、例子或定理，引到下一个计数方法。

\smallskip
\textbf{初学者检查点：}
\begin{itemize}[leftmargin=2em,itemsep=0.25em]
\item 找出这一段承接的上一个概念。
\item 找出它准备引出的下一个概念。
\item 把中间的关键等式或构造单独抄出来。
\end{itemize}

\smallskip
\textbf{本章读法提醒：} 第三章先定义坏事件，再做容斥；坏事件定义不清楚，公式没有意义。
\end{explainblock}


\begin{sourceblock}{第 3 章续读}
{\small\sloppy
\noindent s∈S 1≤i≤m\par
\noindent X Y\par
\noindent = (1 − χAi (s))\par
\noindent s∈S 1≤i≤m\par
\noindent  \par
\noindent X X\par
\noindent m X Y\par
\noindent = 1 + (−1)k χAil (s)\par
\noindent s∈S k=1 1≤i1 <i2 <···<ik ≤m 1≤l≤k\par
\noindent X\par
\noindent m X X Y\par
\noindent = |S| + (−1)k χAil (s)\par
\noindent k=1 1≤i1 <i2 <···<ik ≤m s∈S 1≤l≤k\par
\noindent X\par
\noindent m X\par
\noindent = |S| + (−1)k |Ai1 ∩ Ai2 ∩ · · · ∩ Aik |.\par
\noindent k=1 1≤i1 <i2 <···<ik ≤m\par
}
\end{sourceblock}


\begin{explainblock}
这一段是原书论述链的一部分。读的时候把它当作桥梁：它通常负责把上一个定义、例子或定理，引到下一个计数方法。

\smallskip
\textbf{初学者检查点：}
\begin{itemize}[leftmargin=2em,itemsep=0.25em]
\item 找出这一段承接的上一个概念。
\item 找出它准备引出的下一个概念。
\item 把中间的关键等式或构造单独抄出来。
\end{itemize}

\smallskip
\textbf{本章读法提醒：} 第三章先定义坏事件，再做容斥；坏事件定义不清楚，公式没有意义。
\end{explainblock}


\begin{sourceblock}{定理得证.}
{\small\sloppy
\noindent 定理得证.\par
\noindent 注意到\par
\noindent |A1 ∪ A2 ∪ · · · ∪ Am | = |S| − |A1 ∩ A2 ∩ · · · ∩ Am |,\par
\noindent 将上式代入定理 3.1.1 中的等式 (3.3), 可得到下面的推论.\par
}
\end{sourceblock}


\begin{explainblock}
定理段不要只看结论，要看它把哪两个计数方式连在一起。组合学证明最常见的技巧是：左边按一种标准数，右边按另一种标准数，然后说明两边数的是同一批对象。

\smallskip
\textbf{初学者检查点：}
\begin{itemize}[leftmargin=2em,itemsep=0.25em]
\item 先复述定理的假设，不要把适用范围扩大。
\item 找出证明中的基本计数原则：乘法、加法、双射、递推、容斥或生成函数。
\item 检查最后的等式化简是否和前面的对象解释一致。
\end{itemize}

\smallskip
\textbf{本章读法提醒：} 第三章先定义坏事件，再做容斥；坏事件定义不清楚，公式没有意义。
\end{explainblock}


\begin{sourceblock}{推论 3.1.2 (容斥原理的对偶形式) 设 S 是一个有限集合, P = \{P1 , P2 , . . . , Pm \}}
{\small\sloppy
\noindent 推论 3.1.2 (容斥原理的对偶形式) 设 S 是一个有限集合, P = \{P1 , P2 , . . . , Pm \}\par
\noindent 是一组性质集, 其中 Pi 是同 S 有关某种性质. 设 Ai 是 S 中具有性质 Pi 的元素构\par
\noindent 成的集合 (1 ≤ i ≤ m), 则 S 中至少具有 P 中一个性质的元素的个数为:\par
\noindent X\par
\noindent m X\par
\noindent |A1 ∪ A2 ∪ · · · ∪ Am | = (−1)k−1 |Ai1 ∩ · · · ∩ Aik |.\par
\noindent k=1 1≤i1 <i2 <···<ik ≤m\par
}
\end{sourceblock}


\begin{explainblock}
定理段不要只看结论，要看它把哪两个计数方式连在一起。组合学证明最常见的技巧是：左边按一种标准数，右边按另一种标准数，然后说明两边数的是同一批对象。

\smallskip
\textbf{初学者检查点：}
\begin{itemize}[leftmargin=2em,itemsep=0.25em]
\item 先复述定理的假设，不要把适用范围扩大。
\item 找出证明中的基本计数原则：乘法、加法、双射、递推、容斥或生成函数。
\item 检查最后的等式化简是否和前面的对象解释一致。
\end{itemize}

\smallskip
\textbf{本章读法提醒：} 第三章先定义坏事件，再做容斥；坏事件定义不清楚，公式没有意义。
\end{explainblock}


\begin{sourceblock}{定理 3.1.3 (一般形式的容斥原理) 给定有限的集合 S 及与 S 中元素相关的性质}
{\small\sloppy
\noindent 定理 3.1.3 (一般形式的容斥原理) 给定有限的集合 S 及与 S 中元素相关的性质\par
\noindent 集 P = \{P1 , P2 , · · · Pn \}. 设 X ⊂ P , 记 N≥ (X) 为 S 中至少具有 X 中所有性质的元\par
\noindent 素个数; N= (X) 为 S 中恰好具有 X 中所有性质的元素个数. 则\par
\noindent X\par
\noindent N= (X) = (−1)|Y |−|X| N≥ (Y ). (3.7)\par
\noindent Y ⊇X\par
\noindent 特别地, 当 X = Φ 时,\par
\noindent X\par
\noindent N= (Φ) = (−1)|Y | N≥ (Y )\par
\noindent Y ⊆P\par
\noindent X\par
\noindent n\par
\noindent = |S| − (−1)|Y | N≥ (Y ). (3.8)\par
\noindent |Y |=1\par
\noindent 证明: 我们只证明 (3.7).\par
\noindent 因为\par
\noindent X\par
\noindent N≥ (Y ) = N= (Z), (3.9)\par
}
\end{sourceblock}


\begin{explainblock}
定理段不要只看结论，要看它把哪两个计数方式连在一起。组合学证明最常见的技巧是：左边按一种标准数，右边按另一种标准数，然后说明两边数的是同一批对象。

\smallskip
\textbf{初学者检查点：}
\begin{itemize}[leftmargin=2em,itemsep=0.25em]
\item 先复述定理的假设，不要把适用范围扩大。
\item 找出证明中的基本计数原则：乘法、加法、双射、递推、容斥或生成函数。
\item 检查最后的等式化简是否和前面的对象解释一致。
\end{itemize}

\smallskip
\textbf{本章读法提醒：} 第三章先定义坏事件，再做容斥；坏事件定义不清楚，公式没有意义。
\end{explainblock}


\begin{sourceblock}{第 3 章续读}
{\small\sloppy
\noindent Z⊇Y\par
\noindent 故\par
\noindent X X X\par
\noindent (−1)|Y |−|X| N≥ (Y ) = (−1)|Y |−|X| N= (Z)\par
\noindent Y ⊇X Y ⊇X Z⊇Y\par
\noindent  \par
\noindent X X\par
\noindent =  (−1)|Y |−|X|  N= (Z)\par
\noindent Z⊇X Z⊇Y ⊇X\par
\noindent X |Z|−|X|\par
\noindent X\par
\noindent |Z| − |X|\par
\noindent = (−1)i N= (Z)\par
\noindent i\par
\noindent Z⊇X i=0\par
\noindent X\par
\noindent = δ0,|Z|−|X| N= (Z)\par
\noindent Z⊇X\par
}
\end{sourceblock}


\begin{explainblock}
这一段是原书论述链的一部分。读的时候把它当作桥梁：它通常负责把上一个定义、例子或定理，引到下一个计数方法。

\smallskip
\textbf{初学者检查点：}
\begin{itemize}[leftmargin=2em,itemsep=0.25em]
\item 找出这一段承接的上一个概念。
\item 找出它准备引出的下一个概念。
\item 把中间的关键等式或构造单独抄出来。
\end{itemize}

\smallskip
\textbf{本章读法提醒：} 第三章先定义坏事件，再做容斥；坏事件定义不清楚，公式没有意义。
\end{explainblock}


\begin{sourceblock}{第 3 章续读}
{\small\sloppy
\noindent = N= (X).\par
\noindent 利用定理 3.1.1 和定理 3.1.3, 我们可以计算集合 S 中不具有性质集 P 上任意一\par
\noindent 个性质的元素个数和 S 中具有性质集 P 中某些性质的元素个数. 下面我们将其推\par
\noindent 广到更一般的情形.\par
}
\end{sourceblock}


\begin{explainblock}
这一段是原书论述链的一部分。读的时候把它当作桥梁：它通常负责把上一个定义、例子或定理，引到下一个计数方法。

\smallskip
\textbf{初学者检查点：}
\begin{itemize}[leftmargin=2em,itemsep=0.25em]
\item 找出这一段承接的上一个概念。
\item 找出它准备引出的下一个概念。
\item 把中间的关键等式或构造单独抄出来。
\end{itemize}

\smallskip
\textbf{本章读法提醒：} 第三章先定义坏事件，再做容斥；坏事件定义不清楚，公式没有意义。
\end{explainblock}


\begin{sourceblock}{定理 3.1.4 (广义的容斥原理) 设 S 是一个有限集合, P = \{P1 , P2 , . . . , Pm \} 是一}
{\small\sloppy
\noindent 定理 3.1.4 (广义的容斥原理) 设 S 是一个有限集合, P = \{P1 , P2 , . . . , Pm \} 是一\par
\noindent 组性质集, 其中 Pi 是同 S 有关某种性质. 设 X ⊆ P , N≥ (X) 为至少具有 X 中所有\par
\noindent 性质的元素的个数, N= (X) 为恰好具有 X 中所有性质的元素的个数. 对任意非负整\par
\noindent 数 k, 设集合 S 中恰好具有 P 中 k 个性质的元素个数为 Ek , 并记\par
\noindent X\par
\noindent Nk = N≥ (X),\par
\noindent |X|=k\par
\noindent X⊆P\par
\noindent 则\par
\noindent |P |\par
\noindent X\par
\noindent n\par
\noindent Ek = (−1)n−k Nn . (3.10)\par
\noindent k\par
\noindent n=k\par
\noindent 证明: (方法一) 显然\par
\noindent X\par
\noindent Ek = N= (X),\par
}
\end{sourceblock}


\begin{explainblock}
定理段不要只看结论，要看它把哪两个计数方式连在一起。组合学证明最常见的技巧是：左边按一种标准数，右边按另一种标准数，然后说明两边数的是同一批对象。

\smallskip
\textbf{初学者检查点：}
\begin{itemize}[leftmargin=2em,itemsep=0.25em]
\item 先复述定理的假设，不要把适用范围扩大。
\item 找出证明中的基本计数原则：乘法、加法、双射、递推、容斥或生成函数。
\item 检查最后的等式化简是否和前面的对象解释一致。
\end{itemize}

\smallskip
\textbf{本章读法提醒：} 第三章先定义坏事件，再做容斥；坏事件定义不清楚，公式没有意义。
\end{explainblock}


\begin{sourceblock}{第 3 章续读}
{\small\sloppy
\noindent |X|=k\par
\noindent X⊆P\par
\noindent 所以由定理 3.1.3 得\par
\noindent X\par
\noindent Ek = N= (X)\par
\noindent |X|=k\par
\noindent X⊆P\par
\noindent X X\par
\noindent = (−1)|Y |−|X| N≥ (Y )\par
\noindent |X|=k Y ⊇X\par
\noindent X⊆P\par
\noindent |P |\par
\noindent X X\par
\noindent n−k n\par
\noindent = (−1) N≥ (Y )\par
\noindent k |Y |=n\par
\noindent n=k\par
\noindent Y ⊂P\par
}
\end{sourceblock}


\begin{explainblock}
这一段是原书论述链的一部分。读的时候把它当作桥梁：它通常负责把上一个定义、例子或定理，引到下一个计数方法。

\smallskip
\textbf{初学者检查点：}
\begin{itemize}[leftmargin=2em,itemsep=0.25em]
\item 找出这一段承接的上一个概念。
\item 找出它准备引出的下一个概念。
\item 把中间的关键等式或构造单独抄出来。
\end{itemize}

\smallskip
\textbf{本章读法提醒：} 第三章先定义坏事件，再做容斥；坏事件定义不清楚，公式没有意义。
\end{explainblock}


\begin{sourceblock}{第 3 章续读}
{\small\sloppy
\noindent |P |\par
\noindent X\par
\noindent n−k n\par
\noindent = (−1) Nn .\par
\noindent k\par
\noindent n=k\par
\noindent (方法二) 设 s ∈ S, 记 P (s) 为 s 所具有的 P 中的性质组成的集合, 则\par
\noindent X\par
\noindent Nk = N≥ (X)\par
\noindent |X|=k\par
\noindent X⊆P\par
\noindent X X\par
\noindent = 1\par
\noindent |X|=k s∈S\par
\noindent X⊆P X⊆P (s)\par
\noindent X X\par
\noindent = 1\par
\noindent s∈S |X|=k\par
}
\end{sourceblock}


\begin{explainblock}
这一段是原书论述链的一部分。读的时候把它当作桥梁：它通常负责把上一个定义、例子或定理，引到下一个计数方法。

\smallskip
\textbf{初学者检查点：}
\begin{itemize}[leftmargin=2em,itemsep=0.25em]
\item 找出这一段承接的上一个概念。
\item 找出它准备引出的下一个概念。
\item 把中间的关键等式或构造单独抄出来。
\end{itemize}

\smallskip
\textbf{本章读法提醒：} 第三章先定义坏事件，再做容斥；坏事件定义不清楚，公式没有意义。
\end{explainblock}


\begin{sourceblock}{第 3 章续读}
{\small\sloppy
\noindent X⊆P (s)\par
\noindent X |P (s)|\par
\noindent = ,\par
\noindent k\par
\noindent s∈S\par
\noindent 因此,\par
\noindent |P |\par
\noindent X n\par
\noindent Nk = En . (3.11)\par
\noindent k\par
\noindent n=k\par
\noindent 设 Nk 的普通生成函数为 N (t), Ek 的生成函数为 E(t), 即:\par
\noindent |P | |P |\par
\noindent X X\par
\noindent k\par
\noindent N (x) = Nk t , E(x) = Ek tk .\par
\noindent k=0 k=0\par
\noindent 由 (3.11) 得\par
}
\end{sourceblock}


\begin{explainblock}
这一段是原书论述链的一部分。读的时候把它当作桥梁：它通常负责把上一个定义、例子或定理，引到下一个计数方法。

\smallskip
\textbf{初学者检查点：}
\begin{itemize}[leftmargin=2em,itemsep=0.25em]
\item 找出这一段承接的上一个概念。
\item 找出它准备引出的下一个概念。
\item 把中间的关键等式或构造单独抄出来。
\end{itemize}

\smallskip
\textbf{本章读法提醒：} 第三章先定义坏事件，再做容斥；坏事件定义不清楚，公式没有意义。
\end{explainblock}


\begin{sourceblock}{第 3 章续读}
{\small\sloppy
\noindent |P |\par
\noindent X\par
\noindent N (t) = Nk tk\par
\noindent k=0\par
\noindent |P | |P |\par
\noindent X X n\par
\noindent = En tk\par
\noindent k\par
\noindent k=0 n=k\par
\noindent |P |\par
\noindent X n\par
\noindent X n k\par
\noindent = En t\par
\noindent k\par
\noindent n=0 k=0\par
\noindent |P |\par
\noindent X\par
\noindent = En (t + 1)n\par
}
\end{sourceblock}


\begin{explainblock}
这一段是原书论述链的一部分。读的时候把它当作桥梁：它通常负责把上一个定义、例子或定理，引到下一个计数方法。

\smallskip
\textbf{初学者检查点：}
\begin{itemize}[leftmargin=2em,itemsep=0.25em]
\item 找出这一段承接的上一个概念。
\item 找出它准备引出的下一个概念。
\item 把中间的关键等式或构造单独抄出来。
\end{itemize}

\smallskip
\textbf{本章读法提醒：} 第三章先定义坏事件，再做容斥；坏事件定义不清楚，公式没有意义。
\end{explainblock}


\begin{sourceblock}{第 3 章续读}
{\small\sloppy
\noindent n=0\par
\noindent = E(t + 1).\par
\noindent 所以\par
\noindent |P |\par
\noindent X\par
\noindent E(t) = N (t − 1) = Nn (t − 1)n\par
\noindent n=0\par
\noindent |P |\par
\noindent X X\par
\noindent n\par
\noindent n−k n k\par
\noindent = Nn (−1) t\par
\noindent k\par
\noindent n=0 k=0\par
\noindent |P | |P |\par
\noindent X X\par
\noindent n−k n\par
\noindent = (−1) Nn tk ,\par
}
\end{sourceblock}


\begin{explainblock}
这一段是原书论述链的一部分。读的时候把它当作桥梁：它通常负责把上一个定义、例子或定理，引到下一个计数方法。

\smallskip
\textbf{初学者检查点：}
\begin{itemize}[leftmargin=2em,itemsep=0.25em]
\item 找出这一段承接的上一个概念。
\item 找出它准备引出的下一个概念。
\item 把中间的关键等式或构造单独抄出来。
\end{itemize}

\smallskip
\textbf{本章读法提醒：} 第三章先定义坏事件，再做容斥；坏事件定义不清楚，公式没有意义。
\end{explainblock}


\begin{sourceblock}{第 3 章续读}
{\small\sloppy
\noindent k\par
\noindent k=0 n=k\par
\noindent 故\par
\noindent |P |\par
\noindent X\par
\noindent n−k n\par
\noindent Ek = (−1) Nn .\par
\noindent k\par
\noindent n=k\par
}
\end{sourceblock}


\begin{explainblock}
这一段是原书论述链的一部分。读的时候把它当作桥梁：它通常负责把上一个定义、例子或定理，引到下一个计数方法。

\smallskip
\textbf{初学者检查点：}
\begin{itemize}[leftmargin=2em,itemsep=0.25em]
\item 找出这一段承接的上一个概念。
\item 找出它准备引出的下一个概念。
\item 把中间的关键等式或构造单独抄出来。
\end{itemize}

\smallskip
\textbf{本章读法提醒：} 第三章先定义坏事件，再做容斥；坏事件定义不清楚，公式没有意义。
\end{explainblock}


\begin{sourceblock}{§ 3.2 容斥原理的应用}
{\small\sloppy
\noindent § 3.2. 容斥原理的应用\par
}
\end{sourceblock}


\begin{explainblock}
这一节先不要急着背公式。先问三个问题：对象是什么，哪些限制条件不能丢，最后到底是在数所有对象、还是在数某个统计量的分布。把这三件事说清楚，后面的公式才会有位置。

\smallskip
\textbf{初学者检查点：}
\begin{itemize}[leftmargin=2em,itemsep=0.25em]
\item 把本节出现的新对象写成一句话定义。
\item 把每个公式左边和右边分别翻译成“在数什么”。
\item 找一个最小非平凡例子，手工列出所有对象。
\end{itemize}

\smallskip
\textbf{本章读法提醒：} 第三章先定义坏事件，再做容斥；坏事件定义不清楚，公式没有意义。
\end{explainblock}


\begin{sourceblock}{例 3.2.1 设集合 [n] 的错排数为 Dn , 用容斥原理证明:}
{\small\sloppy
\noindent 例 3.2.1 设集合 [n] 的错排数为 Dn , 用容斥原理证明:\par
\noindent X\par
\noindent n\par
\noindent (−1)k\par
\noindent Dn = n! .\par
\noindent k!\par
\noindent k=0\par
\noindent 证明: 令 Sn 表示集合 [n] 上所有排列的集合, 则 |Sn | = n!. 定义 Sn 上的性质\par
\noindent 集合:\par
\noindent P = \{P1 , P2 , . . . , Pn \},\par
\noindent 其中, Pi 表示排列中 i 被固定在它原来位置的, 即 i 是排列的固定点. 令 Ai 为 Sn\par
\noindent 中满足性质 Pi 的所有排列的集合. 考虑 Ai 中的任意排列:\par
\noindent j1 j2 · · · ji−1 iji+1 · · · jn ,\par
\noindent 显然地, j1 j2 . . . ji−1 ji+1 · · · jn 是 \{1, 2, · · · , i − 1, i + 1, . . . , n\} 的全排列, 所以:\par
\noindent |Ai | = (n − 1)!, 1 ≤ i ≤ n.\par
\noindent 同理:\par
\noindent |Ai ∩ Aj | = (n − 2)!, 1 ≤ i 6= j ≤ m,\par
\noindent 一般地,\par
}
\end{sourceblock}


\begin{explainblock}
例题是学习这本书最重要的部分。不要只看答案，要把每个限制条件变成一个计数动作：先安排谁、再插入谁、有没有重复、需不需要除以对称性。

\smallskip
\textbf{初学者检查点：}
\begin{itemize}[leftmargin=2em,itemsep=0.25em]
\item 先写出没有限制时的总数。
\item 逐条加入限制条件，观察总数怎样变化。
\item 最后检查有没有把同一种方案重复算了多次。
\end{itemize}

\smallskip
\textbf{本章读法提醒：} 第三章先定义坏事件，再做容斥；坏事件定义不清楚，公式没有意义。
\end{explainblock}


\begin{sourceblock}{第 3 章续读}
{\small\sloppy
\noindent |Ai1 ∩ Ai2 · · · ∩ Aik | = (n − k)!, 其中若 l 6= j, 则 il 6= ij ,\par
\noindent 显然 Dn 就是 Sn 中不满足性质 P1 , P2 , . . . , Pn 的所有排列, 由容斥原理可得:\par
\noindent Dn = |Ā1 ∩ Ā2 ∩ · · · ∩ Ān |\par
\noindent X\par
\noindent n X\par
\noindent = |Sn | + (−1)k |Ai1 ∩ Ai2 ∩ · · · ∩ Aik |\par
\noindent k=1 1≤i1 <i2 <···<ik ≤n\par
\noindent n n n n\par
\noindent = n! − (n − 1)! + (n − 2)! + · · · + (−1) 0!\par
\noindent Xn\par
\noindent (−1)k\par
\noindent = n! .\par
\noindent k!\par
\noindent k=0\par
}
\end{sourceblock}


\begin{explainblock}
这一段是原书论述链的一部分。读的时候把它当作桥梁：它通常负责把上一个定义、例子或定理，引到下一个计数方法。

\smallskip
\textbf{初学者检查点：}
\begin{itemize}[leftmargin=2em,itemsep=0.25em]
\item 找出这一段承接的上一个概念。
\item 找出它准备引出的下一个概念。
\item 把中间的关键等式或构造单独抄出来。
\end{itemize}

\smallskip
\textbf{本章读法提醒：} 第三章先定义坏事件，再做容斥；坏事件定义不清楚，公式没有意义。
\end{explainblock}


\begin{sourceblock}{例 3.2.2 设 n 为正整数, 欧拉数 φ(n) 定义为集合 [n] 中与 n 互素的数个数. 若}
{\small\sloppy
\noindent 例 3.2.2 设 n 为正整数, 欧拉数 φ(n) 定义为集合 [n] 中与 n 互素的数个数. 若\par
\noindent p1 , p2 , . . . , pk 是 n 的全部素因子, 则\par
\noindent φ(n) = n 1 − 1− ··· 1 − .\par
\noindent p1 p2 pk\par
\noindent 证明: 设集合 P = \{P1 , P2 , . . . , Pk \} 是与 [n] 中元素有关的一组性质集合. 设\par
\noindent 1 ≤ i ≤ k, 若 m ∈ [n] 且 m 能被 pi 整除, 则称 m 满足性质 Pi , 并将集合 [n] 中具有\par
\noindent Pi 性质的所有元素构成的集合记为 Ai , 则 Ai 显然是集合 [n] 中能被 pi 整除的数所\par
\noindent 组成的集合, 故 |Ai | = pni , 则至少具有性质集合 P 中任意 j 个性质 Pi1 , Pi2 , . . . , Pij\par
\noindent 元素的个数为\par
\noindent n\par
\noindent |Ai1 ∩ Ai2 ∩ · · · ∩ Aij | = .\par
\noindent pi1 pi2 · · · pij\par
\noindent 由容斥原理知在 [n] 与 n 互素的的元素的个数为\par
\noindent X\par
\noindent k X n\par
\noindent φ(n) = |Ā1 ∩ Ā2 ∩ · · · ∩ Āk | = n − (−1)j\par
\noindent p i1 p i2 · · · p ij\par
\noindent j=1 1≤i1 <i2 <···<ij ≤k\par
}
\end{sourceblock}


\begin{explainblock}
例题是学习这本书最重要的部分。不要只看答案，要把每个限制条件变成一个计数动作：先安排谁、再插入谁、有没有重复、需不需要除以对称性。

\smallskip
\textbf{初学者检查点：}
\begin{itemize}[leftmargin=2em,itemsep=0.25em]
\item 先写出没有限制时的总数。
\item 逐条加入限制条件，观察总数怎样变化。
\item 最后检查有没有把同一种方案重复算了多次。
\end{itemize}

\smallskip
\textbf{本章读法提醒：} 第三章先定义坏事件，再做容斥；坏事件定义不清楚，公式没有意义。
\end{explainblock}


\begin{sourceblock}{第 3 章续读}
{\small\sloppy
\noindent =n 1− 1− ··· 1 − .\par
\noindent p1 p2 pk\par
}
\end{sourceblock}


\begin{explainblock}
这一段是原书论述链的一部分。读的时候把它当作桥梁：它通常负责把上一个定义、例子或定理，引到下一个计数方法。

\smallskip
\textbf{初学者检查点：}
\begin{itemize}[leftmargin=2em,itemsep=0.25em]
\item 找出这一段承接的上一个概念。
\item 找出它准备引出的下一个概念。
\item 把中间的关键等式或构造单独抄出来。
\end{itemize}

\smallskip
\textbf{本章读法提醒：} 第三章先定义坏事件，再做容斥；坏事件定义不清楚，公式没有意义。
\end{explainblock}


\begin{sourceblock}{例 3.2.3 用容斥原理求集合 [n] 的 r–可重排列中每个元素都至少出现一次的排}
{\small\sloppy
\noindent 例 3.2.3 用容斥原理求集合 [n] 的 r–可重排列中每个元素都至少出现一次的排\par
\noindent 列数 qr .\par
\noindent 解: 设 [n] 的所有 r–可重排列的组成的集合为 S, 则 |S| = nr . 定义集合 S 的性\par
\noindent 质集合:\par
\noindent P = \{P1 , P2 , . . . , Pn \},\par
\noindent 其中 Pi 表示在 r-排列中不存在元素 i, 记 Ai 为集合 S 中满足性质 Pi 的元素组成\par
\noindent 的集合, 则\par
\noindent |Ai | = (n − 1)r .\par
\noindent 对任意 1 ≤ i < j ≤ n , 显然有\par
\noindent |Ai ∩ Aj | = (n − 2)r ,\par
\noindent 一般地, 对 1 ≤ i1 < i2 < . . . < ik ≤ n, 可得\par
\noindent |Ai1 ∩ Ai2 ∩ · · · ∩ Aik | = (n − k)r .\par
\noindent 由题意知 qr 即为 S 中不满足 P 中任意一个性质的排列的个数, 由容斥原理可得:\par
\noindent n n n r\par
\noindent qr = n r − (n − 1)r + (n − 2)r + · · · + (−1)n 0\par
\noindent X\par
\noindent n n 0,\par
\noindent n r > n,\par
}
\end{sourceblock}


\begin{explainblock}
例题是学习这本书最重要的部分。不要只看答案，要把每个限制条件变成一个计数动作：先安排谁、再插入谁、有没有重复、需不需要除以对称性。

\smallskip
\textbf{初学者检查点：}
\begin{itemize}[leftmargin=2em,itemsep=0.25em]
\item 先写出没有限制时的总数。
\item 逐条加入限制条件，观察总数怎样变化。
\item 最后检查有没有把同一种方案重复算了多次。
\end{itemize}

\smallskip
\textbf{本章读法提醒：} 第三章先定义坏事件，再做容斥；坏事件定义不清楚，公式没有意义。
\end{explainblock}


\begin{sourceblock}{第 3 章续读}
{\small\sloppy
\noindent = (−1)k\par
\noindent (n − k)r =\par
\noindent k n!, r = n.\par
\noindent k=0\par
}
\end{sourceblock}


\begin{explainblock}
这一段是原书论述链的一部分。读的时候把它当作桥梁：它通常负责把上一个定义、例子或定理，引到下一个计数方法。

\smallskip
\textbf{初学者检查点：}
\begin{itemize}[leftmargin=2em,itemsep=0.25em]
\item 找出这一段承接的上一个概念。
\item 找出它准备引出的下一个概念。
\item 把中间的关键等式或构造单独抄出来。
\end{itemize}

\smallskip
\textbf{本章读法提醒：} 第三章先定义坏事件，再做容斥；坏事件定义不清楚，公式没有意义。
\end{explainblock}


\begin{sourceblock}{例 3.2.4 求不定方程}
{\small\sloppy
\noindent 例 3.2.4 求不定方程\par
\noindent x1 + x2 + · · · + xn = r (3.12)\par
\noindent 满足:\par
\noindent xi < s, 1 ≤ i ≤ n,\par
\noindent 的非负整数解的个数 Ns .\par
\noindent 解: 设 X 为不定方程 (3.12) 非负整数解的集合, 则\par
\noindent n+r−1\par
\noindent |X| = .\par
\noindent r\par
\noindent 设 P = \{P1 , P2 , . . . , Pn \} 是与 X 中元素相关的性质, 设 x = (x1 , x2 , . . . , xn ) ∈ X, 若\par
\noindent xi ≥ s 则称 x 满足性质 Pi . 记 Ai 为 X 中满足性质 Pi 元素构成的集合, 则\par
\noindent n+r−s−1\par
\noindent |Ai | = ,\par
\noindent r−s\par
\noindent 一般地, X 中满足性质集 P 中任意 k 个性质 Pi1 , Pi2 , . . . , Pik 的元素构成的集合则\par
\noindent n + r − ks − 1\par
\noindent |Ai1 ∩ Ai2 ∩ · · · ∩ Aik | = .\par
\noindent r − ks\par
}
\end{sourceblock}


\begin{explainblock}
例题是学习这本书最重要的部分。不要只看答案，要把每个限制条件变成一个计数动作：先安排谁、再插入谁、有没有重复、需不需要除以对称性。

\smallskip
\textbf{初学者检查点：}
\begin{itemize}[leftmargin=2em,itemsep=0.25em]
\item 先写出没有限制时的总数。
\item 逐条加入限制条件，观察总数怎样变化。
\item 最后检查有没有把同一种方案重复算了多次。
\end{itemize}

\smallskip
\textbf{本章读法提醒：} 第三章先定义坏事件，再做容斥；坏事件定义不清楚，公式没有意义。
\end{explainblock}


\begin{sourceblock}{第 3 章续读}
{\small\sloppy
\noindent 故由容斥原理可得\par
\noindent Xn\par
\noindent n+r−1 k n n + r − ks − 1\par
\noindent Ns = − (−1)\par
\noindent r k r − ks\par
\noindent k=1\par
\noindent X\par
\noindent n\par
\noindent k n n + r − ks − 1\par
\noindent = (−1) .\par
\noindent k r − ks\par
\noindent k=0\par
\noindent 特别地, 取 s = 1 得\par
\noindent \par
\noindent X\par
\noindent n 1, r = 0,\par
\noindent n n + r − k − 1\par
\noindent N1 = (−1)k =\par
}
\end{sourceblock}


\begin{explainblock}
这一段是原书论述链的一部分。读的时候把它当作桥梁：它通常负责把上一个定义、例子或定理，引到下一个计数方法。

\smallskip
\textbf{初学者检查点：}
\begin{itemize}[leftmargin=2em,itemsep=0.25em]
\item 找出这一段承接的上一个概念。
\item 找出它准备引出的下一个概念。
\item 把中间的关键等式或构造单独抄出来。
\end{itemize}

\smallskip
\textbf{本章读法提醒：} 第三章先定义坏事件，再做容斥；坏事件定义不清楚，公式没有意义。
\end{explainblock}


\begin{sourceblock}{第 3 章续读}
{\small\sloppy
\noindent k r−k 0, r > 0.\par
\noindent k=0\par
}
\end{sourceblock}


\begin{explainblock}
这一段是原书论述链的一部分。读的时候把它当作桥梁：它通常负责把上一个定义、例子或定理，引到下一个计数方法。

\smallskip
\textbf{初学者检查点：}
\begin{itemize}[leftmargin=2em,itemsep=0.25em]
\item 找出这一段承接的上一个概念。
\item 找出它准备引出的下一个概念。
\item 把中间的关键等式或构造单独抄出来。
\end{itemize}

\smallskip
\textbf{本章读法提醒：} 第三章先定义坏事件，再做容斥；坏事件定义不清楚，公式没有意义。
\end{explainblock}


\begin{sourceblock}{例 3.2.5 (Eratosthenes 筛法) 求不大于 n 的所有素数的个数. 特别地, 求不大于}
{\small\sloppy
\noindent 例 3.2.5 (Eratosthenes 筛法) 求不大于 n 的所有素数的个数. 特别地, 求不大于\par
\noindent 解: 设 p1 (= 2), p2 (= 3), p3 (= 5), . . ., 是单调递增的素数序列. 对于 x > 0,\par
\noindent √\par
\noindent x ∈ R, 设 π(x) 表示不大于 x 的素数个数. 设 k = π( n), 对于整数 i(1 ≤ i ≤ k),\par
\noindent 设 Ai 是 \{2, 3, . . . , n\} 中 pi 的所有倍数组成的集合, 则 |Ai | = [n/pi ]; 一般地, 对于\par
\noindent 1 ≤ m1 < m2 < · · · < mi ≤ k, 我们显然有\par
\noindent n\par
\noindent |Am1 Am2 · · · Ami | = ,\par
\noindent pm1 pm2 · · · pmi\par
\noindent 其中 [x] 表示不大于 x 的整数.\par
\noindent 设 q ∈ \{2, 3, . . . , n\}, 若 q 不能被 p1 , p2 , . . . , pk 整除, 即 q ∈ Ā1 Ā2 · · · A¯k . 所以 q\par
\noindent √ √\par
\noindent 的素因子大于 n. 又因 q ≤ n, 故 q 是大于 n 而不大于 n 的素数 (否则 q 至少有\par
\noindent √\par
\noindent 两个大于 n 的素因子, 乘积显然大于 n, 与假设矛盾). 所以集合 Ā1 Ā2 · · · A¯k 由大\par
\noindent √\par
\noindent 于 n 但不大于 n 的所有素数组成, 由此得:\par
\noindent √\par
}
\end{sourceblock}


\begin{explainblock}
例题是学习这本书最重要的部分。不要只看答案，要把每个限制条件变成一个计数动作：先安排谁、再插入谁、有没有重复、需不需要除以对称性。

\smallskip
\textbf{初学者检查点：}
\begin{itemize}[leftmargin=2em,itemsep=0.25em]
\item 先写出没有限制时的总数。
\item 逐条加入限制条件，观察总数怎样变化。
\item 最后检查有没有把同一种方案重复算了多次。
\end{itemize}

\smallskip
\textbf{本章读法提醒：} 第三章先定义坏事件，再做容斥；坏事件定义不清楚，公式没有意义。
\end{explainblock}


\begin{sourceblock}{第 3 章续读}
{\small\sloppy
\noindent π(n) − π( n) = |Ā1 Ā2 · · · A¯k |.\par
\noindent 由容斥原理可得:\par
\noindent X\par
\noindent k X\par
\noindent |Ā1 Ā2 · · · A¯k | = n − 1 − (−1)i |Am1 Am2 · · · Ami |\par
\noindent i 1≤m1 <m2 <···<mi ≤k\par
\noindent X\par
\noindent k X\par
\noindent n\par
\noindent =n−1− (−1) i\par
\noindent . (3.13)\par
\noindent pm1 pm2 · · · pmi\par
\noindent i=1 1≤m1 <m2 <···<mi ≤k\par
\noindent 所以\par
\noindent X\par
\noindent k X\par
\noindent √ n\par
\noindent π(n) = π( n) + (n − 1) − (−1)i . (3.14)\par
}
\end{sourceblock}


\begin{explainblock}
这一段是原书论述链的一部分。读的时候把它当作桥梁：它通常负责把上一个定义、例子或定理，引到下一个计数方法。

\smallskip
\textbf{初学者检查点：}
\begin{itemize}[leftmargin=2em,itemsep=0.25em]
\item 找出这一段承接的上一个概念。
\item 找出它准备引出的下一个概念。
\item 把中间的关键等式或构造单独抄出来。
\end{itemize}

\smallskip
\textbf{本章读法提醒：} 第三章先定义坏事件，再做容斥；坏事件定义不清楚，公式没有意义。
\end{explainblock}


\begin{sourceblock}{第 3 章续读}
{\small\sloppy
\noindent pi1 pi2 · · · pim\par
\noindent i=1 1≤m1 <m2 <···<mi ≤k\par
\noindent √\par
\noindent 当 n = 100 时, 不大于 n = 10 的素数有 2, 3, 5, 7. 故 π(10) = 4. 由 (3.14) 可\par
\noindent 得:\par
\noindent π(100) = π(10)+99−(50+33+20+14)+(16+10+7+6+4+2)−(3+2+1+0) = 25\par
\noindent 容斥原理虽然能给出我们要求的精确公式, 但是冗长的项数却是它的致命缺陷.\par
\noindent 对固定的 n 来说, 计算次数达到 2n 次. 当 n 比较大的时候, 计算次数非常大, 耗时\par
\noindent 较长. 而很多时候我们可能只需要一个比较接近的值, 所以我们就思考能不能同过舍\par
\noindent 去一些项, 使得计算次数减少的同时也保证得到的值与精确值比较接近呢? 我们首\par
\noindent 先回答第一个问题.\par
}
\end{sourceblock}


\begin{explainblock}
这一段是原书论述链的一部分。读的时候把它当作桥梁：它通常负责把上一个定义、例子或定理，引到下一个计数方法。

\smallskip
\textbf{初学者检查点：}
\begin{itemize}[leftmargin=2em,itemsep=0.25em]
\item 找出这一段承接的上一个概念。
\item 找出它准备引出的下一个概念。
\item 把中间的关键等式或构造单独抄出来。
\end{itemize}

\smallskip
\textbf{本章读法提醒：} 第三章先定义坏事件，再做容斥；坏事件定义不清楚，公式没有意义。
\end{explainblock}


\begin{sourceblock}{定理 3.2.1 (Bonferroni 不等式) 令}
{\small\sloppy
\noindent 定理 3.2.1 (Bonferroni 不等式) 令\par
\noindent X\par
\noindent Sk = |Ai1 ∩ · · · ∩ Aik |, 1 ≤ k ≤ n.\par
\noindent 1≤i1 <i2 <···<ik ≤n\par
\noindent 设 0 ≤ l ≤ n, 则\par
\noindent X\par
\noindent l\par
\noindent |A1 ∪ A2 ∪ · · · ∪ An | ≤ (−1)k−1 Sk , l 为奇数, (3.15)\par
\noindent k=1\par
\noindent Xl\par
\noindent |A1 ∪ A2 ∪ · · · ∪ An | ≥ (−1)k−1 Sk , l 为偶数. (3.16)\par
\noindent k=1\par
\noindent 类似的,\par
\noindent X\par
\noindent l\par
\noindent |Ā1 ∩ Ā2 ∩ · · · ∩ Ān | ≥ |S| − (−1)k−1 Sk , l 为奇数, (3.17)\par
\noindent k=1\par
\noindent Xl\par
}
\end{sourceblock}


\begin{explainblock}
定理段不要只看结论，要看它把哪两个计数方式连在一起。组合学证明最常见的技巧是：左边按一种标准数，右边按另一种标准数，然后说明两边数的是同一批对象。

\smallskip
\textbf{初学者检查点：}
\begin{itemize}[leftmargin=2em,itemsep=0.25em]
\item 先复述定理的假设，不要把适用范围扩大。
\item 找出证明中的基本计数原则：乘法、加法、双射、递推、容斥或生成函数。
\item 检查最后的等式化简是否和前面的对象解释一致。
\end{itemize}

\smallskip
\textbf{本章读法提醒：} 第三章先定义坏事件，再做容斥；坏事件定义不清楚，公式没有意义。
\end{explainblock}


\begin{sourceblock}{第 3 章续读}
{\small\sloppy
\noindent |Ā1 ∩ Ā2 ∩ · · · ∩ Ān | ≤ |S| − (−1)k−1 Sk , l 为偶数. (3.18)\par
\noindent k=1\par
\noindent 特别地, 当 l = 1 时, 由 (3.15) 我们得到了 Boole 不等式\par
\noindent X\par
\noindent |A1 ∪ A2 ∪ · · · ∪ An | ≤ |Ai |. (3.19)\par
\noindent 1≤i≤n\par
\noindent 证明: 设 A = A1 ∪ A2 ∪ A3 ∪ · · · ∪ Am . 对任意 a ∈ A, 定义 I(a) = \{i : a ∈ Ai \},\par
\noindent 则 1 ≤ |I(a)| ≤ m. 注意到:\par
\noindent |I(a)|\par
\noindent X\par
\noindent |I(a)|\par
\noindent (−1)k−1 = 1, (3.20)\par
\noindent k\par
\noindent k=1\par
\noindent 我们有\par
\noindent X X |I(a)|\par
\noindent X\par
\noindent |I(a)|\par
}
\end{sourceblock}


\begin{explainblock}
这一段是原书论述链的一部分。读的时候把它当作桥梁：它通常负责把上一个定义、例子或定理，引到下一个计数方法。

\smallskip
\textbf{初学者检查点：}
\begin{itemize}[leftmargin=2em,itemsep=0.25em]
\item 找出这一段承接的上一个概念。
\item 找出它准备引出的下一个概念。
\item 把中间的关键等式或构造单独抄出来。
\end{itemize}

\smallskip
\textbf{本章读法提醒：} 第三章先定义坏事件，再做容斥；坏事件定义不清楚，公式没有意义。
\end{explainblock}


\begin{sourceblock}{第 3 章续读}
{\small\sloppy
\noindent |A| = 1= (−1) k−1\par
\noindent . (3.21)\par
\noindent k\par
\noindent a∈A a∈A k=1\par
\noindent 因为,\par
\noindent |I(a)|\par
\noindent X Xm\par
\noindent k−1 |I(a)| k−1 |I(a)|\par
\noindent (−1) = (−1) . (3.22)\par
\noindent k k\par
\noindent k=1 k=1\par
\noindent 所以\par
\noindent X\par
\noindent m X |I(a)|\par
\noindent |A| = (−1) k−1\par
\noindent k\par
\noindent k=1 a∈A\par
\noindent Xm X X\par
}
\end{sourceblock}


\begin{explainblock}
这一段是原书论述链的一部分。读的时候把它当作桥梁：它通常负责把上一个定义、例子或定理，引到下一个计数方法。

\smallskip
\textbf{初学者检查点：}
\begin{itemize}[leftmargin=2em,itemsep=0.25em]
\item 找出这一段承接的上一个概念。
\item 找出它准备引出的下一个概念。
\item 把中间的关键等式或构造单独抄出来。
\end{itemize}

\smallskip
\textbf{本章读法提醒：} 第三章先定义坏事件，再做容斥；坏事件定义不清楚，公式没有意义。
\end{explainblock}


\begin{sourceblock}{第 3 章续读}
{\small\sloppy
\noindent = (−1)k−1 χ∩i∈T Ai (a)\par
\noindent k=1 a∈A T ⊂[m]\par
\noindent |T |=k\par
\noindent X\par
\noindent m X \textbackslash{}\par
\noindent = Ai .\par
\noindent k=1 T ⊂[m] i∈T\par
\noindent |T |=k\par
\noindent 对于 1 ≤ ℓ ≤ m, 我们考虑\par
\noindent X\par
\noindent ℓ X \textbackslash{} XX\par
\noindent ℓ\par
\noindent k−1 |I(a)|\par
\noindent Ai = (−1) . (3.23)\par
\noindent k\par
\noindent k=1 T ⊂[m] i∈T a∈A k=1\par
\noindent |T |=k\par
\noindent 因为\par
}
\end{sourceblock}


\begin{explainblock}
这一段是原书论述链的一部分。读的时候把它当作桥梁：它通常负责把上一个定义、例子或定理，引到下一个计数方法。

\smallskip
\textbf{初学者检查点：}
\begin{itemize}[leftmargin=2em,itemsep=0.25em]
\item 找出这一段承接的上一个概念。
\item 找出它准备引出的下一个概念。
\item 把中间的关键等式或构造单独抄出来。
\end{itemize}

\smallskip
\textbf{本章读法提醒：} 第三章先定义坏事件，再做容斥；坏事件定义不清楚，公式没有意义。
\end{explainblock}


\begin{sourceblock}{第 3 章续读}
{\small\sloppy
\noindent XX ℓ XX ℓ\par
\noindent k−1 |I(a)| |I(a)| − 1 |I(a)| − 1\par
\noindent (−1) = +\par
\noindent k k−1 k\par
\noindent a∈A k=1 a∈A k=1\par
\noindent X\par
\noindent |I(a)| − 1\par
\noindent = 1 − (−1)ℓ\par
\noindent ℓ\par
\noindent a∈A\par
\noindent X |I(a)| − 1\par
\noindent = |A| − (−1) ℓ\par
\noindent .\par
\noindent ℓ\par
\noindent a∈A\par
\noindent 因为 |I(a)| ≥ 1 且 1 ≤ ℓ ≤ m, 所以\par
\noindent X |I(a)| − 1\par
\noindent ≥ 0,\par
}
\end{sourceblock}


\begin{explainblock}
这一段是原书论述链的一部分。读的时候把它当作桥梁：它通常负责把上一个定义、例子或定理，引到下一个计数方法。

\smallskip
\textbf{初学者检查点：}
\begin{itemize}[leftmargin=2em,itemsep=0.25em]
\item 找出这一段承接的上一个概念。
\item 找出它准备引出的下一个概念。
\item 把中间的关键等式或构造单独抄出来。
\end{itemize}

\smallskip
\textbf{本章读法提醒：} 第三章先定义坏事件，再做容斥；坏事件定义不清楚，公式没有意义。
\end{explainblock}


\begin{sourceblock}{第 3 章续读}
{\small\sloppy
\noindent ℓ\par
\noindent a∈A\par
\noindent 故 Bonferroni 不等式得证.\par
}
\end{sourceblock}


\begin{explainblock}
这一段是原书论述链的一部分。读的时候把它当作桥梁：它通常负责把上一个定义、例子或定理，引到下一个计数方法。

\smallskip
\textbf{初学者检查点：}
\begin{itemize}[leftmargin=2em,itemsep=0.25em]
\item 找出这一段承接的上一个概念。
\item 找出它准备引出的下一个概念。
\item 把中间的关键等式或构造单独抄出来。
\end{itemize}

\smallskip
\textbf{本章读法提醒：} 第三章先定义坏事件，再做容斥；坏事件定义不清楚，公式没有意义。
\end{explainblock}


\begin{sourceblock}{例 3.2.6 利用 Bonfferoni 不等式估计不大于 10000 的素数的个数.}
{\small\sloppy
\noindent 例 3.2.6 利用 Bonfferoni 不等式估计不大于 10000 的素数的个数.\par
\noindent 解: 由上个例题知, π(100) = 25. 所以由 (3.14) 可以求 π(10000). 但是计算量很\par
\noindent √\par
\noindent 大. 不妨截取 (3.14) 的前 b kc 项, 由 Bonferrnoi’s 不等式对不大于 n 的素数进行\par
\noindent 估计. 对 n = 10000 时, π(100) = k = 25. 通过截取和式的前几项由 Bonferroni’s 不\par
\noindent 等式可得: 所以\par
\noindent l 的值 1 2 3 4 5 6\par
\noindent Pl m−1 S\par
\noindent m=0 (−1) m -7992 5856 319 1262 1229 1229\par
\noindent 与 π(10000) 比较 ≤ ≥ ≤ ≥ ≤ ≥\par
\noindent π(10000) = 1229. (3.24)\par
\noindent P25\par
}
\end{sourceblock}


\begin{explainblock}
例题是学习这本书最重要的部分。不要只看答案，要把每个限制条件变成一个计数动作：先安排谁、再插入谁、有没有重复、需不需要除以对称性。

\smallskip
\textbf{初学者检查点：}
\begin{itemize}[leftmargin=2em,itemsep=0.25em]
\item 先写出没有限制时的总数。
\item 逐条加入限制条件，观察总数怎样变化。
\item 最后检查有没有把同一种方案重复算了多次。
\end{itemize}

\smallskip
\textbf{本章读法提醒：} 第三章先定义坏事件，再做容斥；坏事件定义不清楚，公式没有意义。
\end{explainblock}


\begin{sourceblock}{注记 3.2.2 就上题而言, 当 n = 10000 时, 仅仅截取和式 m=0 (−1)}
{\small\sloppy
\noindent 注记 3.2.2 就上题而言, 当 n = 10000 时, 仅仅截取和式 m=0 (−1)\par
\noindent m−1 S\par
\noindent m 前 5\par
\noindent 项求和得到的结果就与 π(n) 的实际值相等. 事实上, 对于有容斥原理所得到的含有\par
\noindent √\par
\noindent k 项和式的表达式中, 截取前 b kc 项求和所得到的估计值与实际值相比误差是比较\par
\noindent 小的.\par
}
\end{sourceblock}


\begin{explainblock}
注记通常是在补充术语、记号、历史或一个容易忽略的约定。它未必给新公式，但常常决定你后面会不会把符号读错。

\smallskip
\textbf{初学者检查点：}
\begin{itemize}[leftmargin=2em,itemsep=0.25em]
\item 记下新符号的读法。
\item 确认这个约定是否只在本章使用。
\item 把注记里的限制条件补到前面的定义旁边。
\end{itemize}

\smallskip
\textbf{本章读法提醒：} 第三章先定义坏事件，再做容斥；坏事件定义不清楚，公式没有意义。
\end{explainblock}


\begin{sourceblock}{习 题 三}
{\small\sloppy
\noindent 习 题 三\par
}
\end{sourceblock}


\begin{explainblock}
习题段不要先追求技巧。先给题目分类：它是在问排列、组合、递推、生成函数、容斥，还是结构计数。分类正确，工具通常就只剩一两个候选。

\smallskip
\textbf{初学者检查点：}
\begin{itemize}[leftmargin=2em,itemsep=0.25em]
\item 判断题目所属工具。
\item 列出变量和限制条件。
\item 先做小规模 n=2、3、4 的例子，确认直觉。
\end{itemize}

\smallskip
\textbf{本章读法提醒：} 第三章先定义坏事件，再做容斥；坏事件定义不清楚，公式没有意义。
\end{explainblock}


\begin{sourceblock}{习题 3.1 用容斥原理计算下列式子并用生成函数验算.}
{\small\sloppy
\noindent 习题 3.1 用容斥原理计算下列式子并用生成函数验算.\par
\noindent (1)\par
\noindent X\par
\noindent n\par
\noindent k n\par
\noindent (−1) k ,\par
\noindent k\par
\noindent k=1\par
\noindent (2)\par
\noindent X\par
\noindent n\par
\noindent n n−k\par
\noindent k\par
\noindent (−1) 2 .\par
\noindent k\par
\noindent k=0\par
}
\end{sourceblock}


\begin{explainblock}
习题段不要先追求技巧。先给题目分类：它是在问排列、组合、递推、生成函数、容斥，还是结构计数。分类正确，工具通常就只剩一两个候选。

\smallskip
\textbf{初学者检查点：}
\begin{itemize}[leftmargin=2em,itemsep=0.25em]
\item 判断题目所属工具。
\item 列出变量和限制条件。
\item 先做小规模 n=2、3、4 的例子，确认直觉。
\end{itemize}

\smallskip
\textbf{本章读法提醒：} 第三章先定义坏事件，再做容斥；坏事件定义不清楚，公式没有意义。
\end{explainblock}


\begin{sourceblock}{习题 3.2 在 [n] 上的所有排列中, 有多少个排列不含有集合 \{12, 23, 34, . . . , (n−}
{\small\sloppy
\noindent 习题 3.2 在 [n] 上的所有排列中, 有多少个排列不含有集合 \{12, 23, 34, . . . , (n−\par
\noindent 1)n\} 中任何一个二元子序列?\par
}
\end{sourceblock}


\begin{explainblock}
习题段不要先追求技巧。先给题目分类：它是在问排列、组合、递推、生成函数、容斥，还是结构计数。分类正确，工具通常就只剩一两个候选。

\smallskip
\textbf{初学者检查点：}
\begin{itemize}[leftmargin=2em,itemsep=0.25em]
\item 判断题目所属工具。
\item 列出变量和限制条件。
\item 先做小规模 n=2、3、4 的例子，确认直觉。
\end{itemize}

\smallskip
\textbf{本章读法提醒：} 第三章先定义坏事件，再做容斥；坏事件定义不清楚，公式没有意义。
\end{explainblock}


\begin{sourceblock}{习题 3.3 求由 a, b, c, d, e, f 这 6 个字符构成的全排列中不允许出现 ace 和 df}
{\small\sloppy
\noindent 习题 3.3 求由 a, b, c, d, e, f 这 6 个字符构成的全排列中不允许出现 ace 和 df\par
\noindent 图像的排列数.\par
}
\end{sourceblock}


\begin{explainblock}
习题段不要先追求技巧。先给题目分类：它是在问排列、组合、递推、生成函数、容斥，还是结构计数。分类正确，工具通常就只剩一两个候选。

\smallskip
\textbf{初学者检查点：}
\begin{itemize}[leftmargin=2em,itemsep=0.25em]
\item 判断题目所属工具。
\item 列出变量和限制条件。
\item 先做小规模 n=2、3、4 的例子，确认直觉。
\end{itemize}

\smallskip
\textbf{本章读法提醒：} 第三章先定义坏事件，再做容斥；坏事件定义不清楚，公式没有意义。
\end{explainblock}


\begin{sourceblock}{习题 3.4 用 26 个英文字母作不允许重复的全排列, 要求排除 dog, god, gum,}
{\small\sloppy
\noindent 习题 3.4 用 26 个英文字母作不允许重复的全排列, 要求排除 dog, god, gum,\par
\noindent depth, thing 字样出现, 求满足这些条件的排列数.\par
}
\end{sourceblock}


\begin{explainblock}
习题段不要先追求技巧。先给题目分类：它是在问排列、组合、递推、生成函数、容斥，还是结构计数。分类正确，工具通常就只剩一两个候选。

\smallskip
\textbf{初学者检查点：}
\begin{itemize}[leftmargin=2em,itemsep=0.25em]
\item 判断题目所属工具。
\item 列出变量和限制条件。
\item 先做小规模 n=2、3、4 的例子，确认直觉。
\end{itemize}

\smallskip
\textbf{本章读法提醒：} 第三章先定义坏事件，再做容斥；坏事件定义不清楚，公式没有意义。
\end{explainblock}


\begin{sourceblock}{习题 3.5 在 S = (2 · a1 , 2 · a2 , 2 · a3 , . . . , 2 · an ) 的全排列中, 若任何两个相同的}
{\small\sloppy
\noindent 习题 3.5 在 S = (2 · a1 , 2 · a2 , 2 · a3 , . . . , 2 · an ) 的全排列中, 若任何两个相同的\par
\noindent 字符都不相邻, 则叫做 “简单字”. 例如 a2 a1 a3 a2 a1 a3 是一个 n = 3 的简单字. 试用\par
\noindent 容斥原理, 求简单字的个数.\par
}
\end{sourceblock}


\begin{explainblock}
习题段不要先追求技巧。先给题目分类：它是在问排列、组合、递推、生成函数、容斥，还是结构计数。分类正确，工具通常就只剩一两个候选。

\smallskip
\textbf{初学者检查点：}
\begin{itemize}[leftmargin=2em,itemsep=0.25em]
\item 判断题目所属工具。
\item 列出变量和限制条件。
\item 先做小规模 n=2、3、4 的例子，确认直觉。
\end{itemize}

\smallskip
\textbf{本章读法提醒：} 第三章先定义坏事件，再做容斥；坏事件定义不清楚，公式没有意义。
\end{explainblock}


\begin{sourceblock}{习题 3.6 n 对夫妇 (n > 2) 围圆桌就座，要求男女相间，且每对夫妇不得相邻，}
{\small\sloppy
\noindent 习题 3.6 n 对夫妇 (n > 2) 围圆桌就座，要求男女相间，且每对夫妇不得相邻，\par
\noindent 问共有多少种就座方法?\par
}
\end{sourceblock}


\begin{explainblock}
习题段不要先追求技巧。先给题目分类：它是在问排列、组合、递推、生成函数、容斥，还是结构计数。分类正确，工具通常就只剩一两个候选。

\smallskip
\textbf{初学者检查点：}
\begin{itemize}[leftmargin=2em,itemsep=0.25em]
\item 判断题目所属工具。
\item 列出变量和限制条件。
\item 先做小规模 n=2、3、4 的例子，确认直觉。
\end{itemize}

\smallskip
\textbf{本章读法提醒：} 第三章先定义坏事件，再做容斥；坏事件定义不清楚，公式没有意义。
\end{explainblock}


\begin{sourceblock}{习题 3.7 有 A、B、C、D、E、F 共 6 个不同病的病人，又有 1、2、3、4、5、}
{\small\sloppy
\noindent 习题 3.7 有 A、B、C、D、E、F 共 6 个不同病的病人，又有 1、2、3、4、5、\par
\noindent 6 共 6 种不同的新药品，为了检验这些新药的治疗效果，需要对 6 个病人进行试验。\par
\noindent 已知 A、B 均不能服用 1、2 号药，C、D 均不能服用 3 号药，E 不能服用 4、6 号\par
\noindent 药，F 不能服用 5、6 号药。若要求每个病人试验一种他所能服用的药，问有多少种\par
\noindent 试验方案。\par
}
\end{sourceblock}


\begin{explainblock}
习题段不要先追求技巧。先给题目分类：它是在问排列、组合、递推、生成函数、容斥，还是结构计数。分类正确，工具通常就只剩一两个候选。

\smallskip
\textbf{初学者检查点：}
\begin{itemize}[leftmargin=2em,itemsep=0.25em]
\item 判断题目所属工具。
\item 列出变量和限制条件。
\item 先做小规模 n=2、3、4 的例子，确认直觉。
\end{itemize}

\smallskip
\textbf{本章读法提醒：} 第三章先定义坏事件，再做容斥；坏事件定义不清楚，公式没有意义。
\end{explainblock}


\begin{sourceblock}{习题 3.8 已知集合 S = \{1, 2, . . . , n\} 的 r–组合中含有子集 \{1, 2, . . . , m\} 的组}
{\small\sloppy
\noindent 习题 3.8 已知集合 S = \{1, 2, . . . , n\} 的 r–组合中含有子集 \{1, 2, . . . , m\} 的组\par
\noindent n−m\par
\noindent 合方案数为 r−m ，利用容斥原理证明恒等式：\par
\noindent X\par
\noindent m\par
\noindent k m n−k n−m\par
\noindent (−1) = .\par
\noindent k r r−m\par
\noindent k=0\par
}
\end{sourceblock}


\begin{explainblock}
习题段不要先追求技巧。先给题目分类：它是在问排列、组合、递推、生成函数、容斥，还是结构计数。分类正确，工具通常就只剩一两个候选。

\smallskip
\textbf{初学者检查点：}
\begin{itemize}[leftmargin=2em,itemsep=0.25em]
\item 判断题目所属工具。
\item 列出变量和限制条件。
\item 先做小规模 n=2、3、4 的例子，确认直觉。
\end{itemize}

\smallskip
\textbf{本章读法提醒：} 第三章先定义坏事件，再做容斥；坏事件定义不清楚，公式没有意义。
\end{explainblock}


\begin{sourceblock}{习题 3.9 把 2m 件相异物品 a1 , a2 , . . . , am , b1 , b2 , . . . , bm 分给 n 个人 (1 ≤ n ≤}
{\small\sloppy
\noindent 习题 3.9 把 2m 件相异物品 a1 , a2 , . . . , am , b1 , b2 , . . . , bm 分给 n 个人 (1 ≤ n ≤\par
\noindent 2m)，使得每人至少分得一个物品且 aj 与 bj 不分给同一个人 (j = 1, 2, . . . , m)，问\par
\noindent 有多少种不同的分法\par
}
\end{sourceblock}


\begin{explainblock}
习题段不要先追求技巧。先给题目分类：它是在问排列、组合、递推、生成函数、容斥，还是结构计数。分类正确，工具通常就只剩一两个候选。

\smallskip
\textbf{初学者检查点：}
\begin{itemize}[leftmargin=2em,itemsep=0.25em]
\item 判断题目所属工具。
\item 列出变量和限制条件。
\item 先做小规模 n=2、3、4 的例子，确认直觉。
\end{itemize}

\smallskip
\textbf{本章读法提醒：} 第三章先定义坏事件，再做容斥；坏事件定义不清楚，公式没有意义。
\end{explainblock}


\begin{sourceblock}{习题 3.10 以 f (r, n) 表示把 r 件相异物品分给 n 个人，且使得没有人恰好分}
{\small\sloppy
\noindent 习题 3.10 以 f (r, n) 表示把 r 件相异物品分给 n 个人，且使得没有人恰好分\par
\noindent 得一个物品的不同分法数，求 f (r, n) 的计数公式.\par
}
\end{sourceblock}


\begin{explainblock}
习题段不要先追求技巧。先给题目分类：它是在问排列、组合、递推、生成函数、容斥，还是结构计数。分类正确，工具通常就只剩一两个候选。

\smallskip
\textbf{初学者检查点：}
\begin{itemize}[leftmargin=2em,itemsep=0.25em]
\item 判断题目所属工具。
\item 列出变量和限制条件。
\item 先做小规模 n=2、3、4 的例子，确认直觉。
\end{itemize}

\smallskip
\textbf{本章读法提醒：} 第三章先定义坏事件，再做容斥；坏事件定义不清楚，公式没有意义。
\end{explainblock}


\begin{sourceblock}{习题 3.11 求不超过 200 的正整数中素数的个数.}
{\small\sloppy
\noindent 习题 3.11 求不超过 200 的正整数中素数的个数.\par
}
\end{sourceblock}


\begin{explainblock}
习题段不要先追求技巧。先给题目分类：它是在问排列、组合、递推、生成函数、容斥，还是结构计数。分类正确，工具通常就只剩一两个候选。

\smallskip
\textbf{初学者检查点：}
\begin{itemize}[leftmargin=2em,itemsep=0.25em]
\item 判断题目所属工具。
\item 列出变量和限制条件。
\item 先做小规模 n=2、3、4 的例子，确认直觉。
\end{itemize}

\smallskip
\textbf{本章读法提醒：} 第三章先定义坏事件，再做容斥；坏事件定义不清楚，公式没有意义。
\end{explainblock}


\begin{sourceblock}{习题 3.12 n 位的四进制数中，数字 1, 2, 3 各自至少出现一次的数有多少个?}
{\small\sloppy
\noindent 习题 3.12 n 位的四进制数中，数字 1, 2, 3 各自至少出现一次的数有多少个?\par
}
\end{sourceblock}


\begin{explainblock}
习题段不要先追求技巧。先给题目分类：它是在问排列、组合、递推、生成函数、容斥，还是结构计数。分类正确，工具通常就只剩一两个候选。

\smallskip
\textbf{初学者检查点：}
\begin{itemize}[leftmargin=2em,itemsep=0.25em]
\item 判断题目所属工具。
\item 列出变量和限制条件。
\item 先做小规模 n=2、3、4 的例子，确认直觉。
\end{itemize}

\smallskip
\textbf{本章读法提醒：} 第三章先定义坏事件，再做容斥；坏事件定义不清楚，公式没有意义。
\end{explainblock}


\begin{sourceblock}{习题 3.13 单位举行晚会，有个 6 部门各表演一个节目，上场次序编号为}
{\small\sloppy
\noindent 习题 3.13 单位举行晚会，有个 6 部门各表演一个节目，上场次序编号为\par
\noindent 1, 2, . . . , 6。现进行抽签，以决定上场次序。但其中有一个部门希望自己抽到的编号\par
\noindent 为偶数，另一个部门不希望抽到 4 或 6，还有一个部门不希望自己抽到的编号是 3\par
\noindent 的倍数。那么，抽签结果使大家都满意的概率是多少？\par
}
\end{sourceblock}


\begin{explainblock}
习题段不要先追求技巧。先给题目分类：它是在问排列、组合、递推、生成函数、容斥，还是结构计数。分类正确，工具通常就只剩一两个候选。

\smallskip
\textbf{初学者检查点：}
\begin{itemize}[leftmargin=2em,itemsep=0.25em]
\item 判断题目所属工具。
\item 列出变量和限制条件。
\item 先做小规模 n=2、3、4 的例子，确认直觉。
\end{itemize}

\smallskip
\textbf{本章读法提醒：} 第三章先定义坏事件，再做容斥；坏事件定义不清楚，公式没有意义。
\end{explainblock}


\begin{sourceblock}{习题 3.14 在 4 个 x, 3 个 y, 2 个 z 的排列中求不出现 x x x x, y y y, z z 图像}
{\small\sloppy
\noindent 习题 3.14 在 4 个 x, 3 个 y, 2 个 z 的排列中求不出现 x x x x, y y y, z z 图像\par
\noindent 的排列数.\par
}
\end{sourceblock}


\begin{explainblock}
习题段不要先追求技巧。先给题目分类：它是在问排列、组合、递推、生成函数、容斥，还是结构计数。分类正确，工具通常就只剩一两个候选。

\smallskip
\textbf{初学者检查点：}
\begin{itemize}[leftmargin=2em,itemsep=0.25em]
\item 判断题目所属工具。
\item 列出变量和限制条件。
\item 先做小规模 n=2、3、4 的例子，确认直觉。
\end{itemize}

\smallskip
\textbf{本章读法提醒：} 第三章先定义坏事件，再做容斥；坏事件定义不清楚，公式没有意义。
\end{explainblock}


\begin{sourceblock}{习题 3.15 用容斥原理求解方程}
{\small\sloppy
\noindent 习题 3.15 用容斥原理求解方程\par
\noindent x1 + x2 + x3 + x4 = 20\par
\noindent 的整数解的个数，其中\par
\noindent 1 ≤ x1 ≤ 6, 0 ≤ x2 ≤ 7, 4 ≤ x3 ≤ 8, 2 ≤ x4 ≤ 6.\par
}
\end{sourceblock}


\begin{explainblock}
习题段不要先追求技巧。先给题目分类：它是在问排列、组合、递推、生成函数、容斥，还是结构计数。分类正确，工具通常就只剩一两个候选。

\smallskip
\textbf{初学者检查点：}
\begin{itemize}[leftmargin=2em,itemsep=0.25em]
\item 判断题目所属工具。
\item 列出变量和限制条件。
\item 先做小规模 n=2、3、4 的例子，确认直觉。
\end{itemize}

\smallskip
\textbf{本章读法提醒：} 第三章先定义坏事件，再做容斥；坏事件定义不清楚，公式没有意义。
\end{explainblock}


\begin{sourceblock}{习题 3.16 求 [n] = \{1, 2, . . . , n\} 到 \{y1 , y2 , . . . , yk \} 的满射有多少个?}
{\small\sloppy
\noindent 习题 3.16 求 [n] = \{1, 2, . . . , n\} 到 \{y1 , y2 , . . . , yk \} 的满射有多少个?\par
}
\end{sourceblock}


\begin{explainblock}
习题段不要先追求技巧。先给题目分类：它是在问排列、组合、递推、生成函数、容斥，还是结构计数。分类正确，工具通常就只剩一两个候选。

\smallskip
\textbf{初学者检查点：}
\begin{itemize}[leftmargin=2em,itemsep=0.25em]
\item 判断题目所属工具。
\item 列出变量和限制条件。
\item 先做小规模 n=2、3、4 的例子，确认直觉。
\end{itemize}

\smallskip
\textbf{本章读法提醒：} 第三章先定义坏事件，再做容斥；坏事件定义不清楚，公式没有意义。
\end{explainblock}


\chapter{两类 Stirling 数}
\begin{roadmapblock}
本章保留原书正文的主要数学骨架，并在每个定义、定理、例题和论述片段后加入解读。
阅读时不要把目标放在“背下答案”上，而是放在“看懂这一步为什么这样数”上。

\smallskip
\textbf{本章最低目标：} 能说出每个公式左边在数什么、右边在数什么，以及为什么两边相等。
\end{roadmapblock}


\begin{sourceblock}{第 4 章正文}
{\small\sloppy
\noindent 两类 Stirling 数\par
\noindent James Stirling 是一位苏格兰籍的数学家. 他 1692 年 5 月出生于苏格兰斯特灵\par
\noindent 郡, 1770 年 10 月逝世于爱丁堡. Stirling 数以及 Stirling 渐进公式都是以他的名字\par
\noindent 命名的. James Stirling 18 岁进入牛津大学贝利奥尔学院 (Balliol College, Oxford)\par
\noindent 求学, 后被驱逐至威尼斯. 在威尼斯期间, James Stirling 在 Isaac Newton 的帮助下,\par
\noindent 与皇家科学院取得联系并邮寄了他的一篇论文 Methodus differentials Newtoniana\par
\noindent illustrata (Phil.Trans., 1718). 后又在 Newton 的帮助下于 1725 年回到伦敦. 在伦\par
\noindent 敦的十年时间里, 他一直致力于学术工作, 1730 年, 他最重要的工作 The methodus\par
\noindent differentialis, sive tractatus de summatione et interpolatione serierum infinitarum\par
\noindent (4to, London) 发表.\par
\noindent 这里我们主要介绍一下他的一个重要工作—Stirling 数. Stirling 数分为第一类和\par
\noindent 第二类.\par
}
\end{sourceblock}


\begin{explainblock}
这一段是原书论述链的一部分。读的时候把它当作桥梁：它通常负责把上一个定义、例子或定理，引到下一个计数方法。

\smallskip
\textbf{初学者检查点：}
\begin{itemize}[leftmargin=2em,itemsep=0.25em]
\item 找出这一段承接的上一个概念。
\item 找出它准备引出的下一个概念。
\item 把中间的关键等式或构造单独抄出来。
\end{itemize}

\smallskip
\textbf{本章读法提醒：} 第四章区分两类 Stirling 数：第一类对应排列的圈，第二类对应集合划分。
\end{explainblock}


\begin{sourceblock}{§ 4.1 排列的圈结构与第一类 Stirling 数}
{\small\sloppy
\noindent § 4.1. 排列的圈结构与第一类 Stirling 数\par
\noindent 设 π = i1 i2 . . . in 是集合 [n] 的任一排列, 我们可将其表示成下述形式:\par
\noindent !\par
\noindent π= .\par
\noindent i1 i2 · · · in\par
\noindent 它可看作 [n] 到自身的一个双射:\par
\noindent π ↓ ↓ ··· ↓ ,\par
\noindent i1 i2 · · · in\par
\noindent 即 π(k) = ik , k = 1, 2, . . . , n.\par
\noindent 对每一个 k ∈ \{1, 2, . . . , n\}, 考虑序列 k, π(k), π 2 (k), . . .. 因为 π 是双射, 所以\par
\noindent 肯定存在唯一的 ℓ ≥ 1 使得 π ℓ (k) = k, 而元素 k, π(k), . . . , π ℓ−1 (k) 互不相同, 称\par
\noindent (k, π(k), . . . , π ℓ−1 (k)) 为 π 的一个长为 ℓ 的圈. 我们认为圈 (k, π(k), . . . , π ℓ−1 (k)) 与\par
\noindent (π j (k), π j+1 (k), . . . , π ℓ−1 (k), k, . . . , π j−1 (k)) 是等价的. [n] 中每个元素都出现在 π\par
\noindent 的唯一圈中, 因此我们可把 π 看作一些不相交的圈的并或者不同圈 C1 , C2 , . . . , Ck\par
\noindent 66 第4章 两类 Stirling 数\par
\noindent 的乘积, 记为 π = C1 C2 · · · Ck .\par
}
\end{sourceblock}


\begin{explainblock}
这一节先不要急着背公式。先问三个问题：对象是什么，哪些限制条件不能丢，最后到底是在数所有对象、还是在数某个统计量的分布。把这三件事说清楚，后面的公式才会有位置。

\smallskip
\textbf{初学者检查点：}
\begin{itemize}[leftmargin=2em,itemsep=0.25em]
\item 把本节出现的新对象写成一句话定义。
\item 把每个公式左边和右边分别翻译成“在数什么”。
\item 找一个最小非平凡例子，手工列出所有对象。
\end{itemize}

\smallskip
\textbf{本章读法提醒：} 第四章区分两类 Stirling 数：第一类对应排列的圈，第二类对应集合划分。
\end{explainblock}


\begin{sourceblock}{例如 π = 4271365, 即 π(1) = 4, π(2) = 2, π(3) = 7, π(4) = 1, π(5) = 3, π(6) =}
{\small\sloppy
\noindent 例如 π = 4271365, 即 π(1) = 4, π(2) = 2, π(3) = 7, π(4) = 1, π(5) = 3, π(6) =\par
\noindent 6, π(7) = 1; 则它表示为圈结构为 π = (14)(2)(375)(6). 当然 π 表示成不交圈的记号\par
\noindent 是不唯一的, 比如它还可记为 π = (753)(41)(6)(2). 我们可以定义一种标准表示. 要\par
\noindent 求这满足:\par
\noindent (a) 每一个圈中将最大元素放在首位;\par
\noindent (b) 圈按照圈中最大元素从小到大排列.\par
\noindent 比如上面排列 π 的标准圈表示为 (2)(41)(6)(753).\par
}
\end{sourceblock}


\begin{explainblock}
例题是学习这本书最重要的部分。不要只看答案，要把每个限制条件变成一个计数动作：先安排谁、再插入谁、有没有重复、需不需要除以对称性。

\smallskip
\textbf{初学者检查点：}
\begin{itemize}[leftmargin=2em,itemsep=0.25em]
\item 先写出没有限制时的总数。
\item 逐条加入限制条件，观察总数怎样变化。
\item 最后检查有没有把同一种方案重复算了多次。
\end{itemize}

\smallskip
\textbf{本章读法提醒：} 第四章区分两类 Stirling 数：第一类对应排列的圈，第二类对应集合划分。
\end{explainblock}


\begin{sourceblock}{定义 4.1.1 n 元集合所有排列中恰好有 k 个圈的排列的个数, 称为无符号的第一}
{\small\sloppy
\noindent 定义 4.1.1 n 元集合所有排列中恰好有 k 个圈的排列的个数, 称为无符号的第一\par
\noindent 类 Stirling 数, 记为 c(n, k). 并称\par
\noindent s(n, k) = (−1)n−k c(n, k),\par
\noindent 为有符号的第一类 Stirling 数或第一类 Stirling 数.\par
}
\end{sourceblock}


\begin{explainblock}
定义段的任务是定边界。这里要特别注意参数范围、是否允许重复、是否考虑顺序、对象有没有标签。后面的每一道例题和证明，本质上都在反复使用这些边界。

\smallskip
\textbf{初学者检查点：}
\begin{itemize}[leftmargin=2em,itemsep=0.25em]
\item 圈出被定义的对象名称。
\item 圈出参数范围，例如 n、k 是否要求正整数。
\item 判断它是有序对象、无序对象、带标签对象还是结构对象。
\end{itemize}

\smallskip
\textbf{本章读法提醒：} 第四章区分两类 Stirling 数：第一类对应排列的圈，第二类对应集合划分。
\end{explainblock}


\begin{sourceblock}{注记 4.1.2 由定义可知当 k > n ≥ 0 时, c(n, k) = 0; 当 n ≥ 1 时,}
{\small\sloppy
\noindent 注记 4.1.2 由定义可知当 k > n ≥ 0 时, c(n, k) = 0; 当 n ≥ 1 时,\par
\noindent c(n, 0) = 0 c(n, n) = 1;\par
\noindent 并规定 c(0, 0) = 1.\par
}
\end{sourceblock}


\begin{explainblock}
注记通常是在补充术语、记号、历史或一个容易忽略的约定。它未必给新公式，但常常决定你后面会不会把符号读错。

\smallskip
\textbf{初学者检查点：}
\begin{itemize}[leftmargin=2em,itemsep=0.25em]
\item 记下新符号的读法。
\item 确认这个约定是否只在本章使用。
\item 把注记里的限制条件补到前面的定义旁边。
\end{itemize}

\smallskip
\textbf{本章读法提醒：} 第四章区分两类 Stirling 数：第一类对应排列的圈，第二类对应集合划分。
\end{explainblock}


\begin{sourceblock}{例 4.1.1 n = 3 时, 所有排列用圈表示如下:}
{\small\sloppy
\noindent 例 4.1.1 n = 3 时, 所有排列用圈表示如下:\par
\noindent (1)(2)(3), (12)(3), (13)(2), (1)(23), (123), (132),\par
\noindent 所以\par
\noindent c(3, 1) = 2 c(3, 2) = 3 c(3, 3) = 1.\par
\noindent n = 4 时也可以验证有:\par
\noindent c(4, 1) = 6 c(4, 2) = 11 c(4, 3) = 6 c(4, 4) = 1.\par
}
\end{sourceblock}


\begin{explainblock}
例题是学习这本书最重要的部分。不要只看答案，要把每个限制条件变成一个计数动作：先安排谁、再插入谁、有没有重复、需不需要除以对称性。

\smallskip
\textbf{初学者检查点：}
\begin{itemize}[leftmargin=2em,itemsep=0.25em]
\item 先写出没有限制时的总数。
\item 逐条加入限制条件，观察总数怎样变化。
\item 最后检查有没有把同一种方案重复算了多次。
\end{itemize}

\smallskip
\textbf{本章读法提醒：} 第四章区分两类 Stirling 数：第一类对应排列的圈，第二类对应集合划分。
\end{explainblock}


\begin{sourceblock}{定理 4.1.3 当 n ≥ 1, k ≥ 1 时,}
{\small\sloppy
\noindent 定理 4.1.3 当 n ≥ 1, k ≥ 1 时,\par
\noindent c(n, k) = (n − 1)c(n − 1, k) + c(n − 1, k − 1). (4.1)\par
\noindent 证明: 不妨考虑集合 [n] 上的排列. 对于集合 [n] 上的任意一个恰好含有 k 个\par
\noindent 圈的排列, 考虑元素 n, 它或是独立成一个圈, 或是与其他元素一起构成一个圈, 我们\par
\noindent 分为这两种情况讨论:\par
\noindent 当 n 独立成一个圈时, 则其它 k − 1 个圈就是集合 [n − 1] 的 k − 1 个圈的排列,\par
\noindent 此时共有 c(n − 1, k − 1) 种情形.\par
\noindent 当 n 与其他元素一起构成一个圈时, 则将 n 去掉得到集合 [n − 1] 的恰有 k 个\par
\noindent 圈的排列, 但是在将元素 n 插入到集合 [n − 1] 的恰好含有 k 个圈的排列时, 可得到\par
\noindent 4.1 排列的圈结构与第一类 Stirling 数 67\par
\noindent n − 1 个集合 [n] 的恰好含有 k 个圈的排列, 从而这种排列的个数等于集合 [n − 1]\par
\noindent 的恰好含有 k 个圈的排列个数的 n − 1 倍.\par
\noindent 综上,\par
\noindent c(n, k) = (n − 1)c(n − 1, k) + c(n − 1, k − 1).\par
\noindent k\par
\noindent 0 1 2 3 4 5 6 7 8 9\par
\noindent n\par
\noindent 6 0 120 274 225 85 15 1\par
}
\end{sourceblock}


\begin{explainblock}
定理段不要只看结论，要看它把哪两个计数方式连在一起。组合学证明最常见的技巧是：左边按一种标准数，右边按另一种标准数，然后说明两边数的是同一批对象。

\smallskip
\textbf{初学者检查点：}
\begin{itemize}[leftmargin=2em,itemsep=0.25em]
\item 先复述定理的假设，不要把适用范围扩大。
\item 找出证明中的基本计数原则：乘法、加法、双射、递推、容斥或生成函数。
\item 检查最后的等式化简是否和前面的对象解释一致。
\end{itemize}

\smallskip
\textbf{本章读法提醒：} 第四章区分两类 Stirling 数：第一类对应排列的圈，第二类对应集合划分。
\end{explainblock}


\begin{sourceblock}{第 4 章续读}
{\small\sloppy
\noindent 7 0 720 1764 1624 735 175 21 1\par
\noindent 8 0 5040 13068 13132 6769 1960 322 28 1\par
\noindent 9 0 40320 109584 118124 67284 22449 4536 546 36 1\par
\noindent 对于固定的整数 n(n ≥ 0), 下面定理给出序列 \{c(n, k)\}k≥0 的普通型生成函数.\par
}
\end{sourceblock}


\begin{explainblock}
这一段是原书论述链的一部分。读的时候把它当作桥梁：它通常负责把上一个定义、例子或定理，引到下一个计数方法。

\smallskip
\textbf{初学者检查点：}
\begin{itemize}[leftmargin=2em,itemsep=0.25em]
\item 找出这一段承接的上一个概念。
\item 找出它准备引出的下一个概念。
\item 把中间的关键等式或构造单独抄出来。
\end{itemize}

\smallskip
\textbf{本章读法提醒：} 第四章区分两类 Stirling 数：第一类对应排列的圈，第二类对应集合划分。
\end{explainblock}


\begin{sourceblock}{定理 4.1.4 设 n ≥ 0, 则}
{\small\sloppy
\noindent 定理 4.1.4 设 n ≥ 0, 则\par
\noindent X\par
\noindent n\par
\noindent c(n, k)xk = x(x + 1)(x + 2) · · · (x + n − 1). (4.2)\par
\noindent k=0\par
\noindent 证明: 不妨设\par
\noindent X\par
\noindent n\par
\noindent Fn (x) = c(n, k)xk .\par
\noindent k=0\par
\noindent 当 n = 0 时,\par
\noindent F0 (x) = c(0, 0) = 1,\par
\noindent 结论成立 (空的乘积等于 1).\par
\noindent 当 n ≥ 1 时, 因为 c(n, 0) = 0, 所以\par
\noindent X\par
\noindent n\par
\noindent Fn (x) = c(n, k)xk .\par
\noindent k=1\par
}
\end{sourceblock}


\begin{explainblock}
定理段不要只看结论，要看它把哪两个计数方式连在一起。组合学证明最常见的技巧是：左边按一种标准数，右边按另一种标准数，然后说明两边数的是同一批对象。

\smallskip
\textbf{初学者检查点：}
\begin{itemize}[leftmargin=2em,itemsep=0.25em]
\item 先复述定理的假设，不要把适用范围扩大。
\item 找出证明中的基本计数原则：乘法、加法、双射、递推、容斥或生成函数。
\item 检查最后的等式化简是否和前面的对象解释一致。
\end{itemize}

\smallskip
\textbf{本章读法提醒：} 第四章区分两类 Stirling 数：第一类对应排列的圈，第二类对应集合划分。
\end{explainblock}


\begin{sourceblock}{第 4 章续读}
{\small\sloppy
\noindent 由 c(n, k) 递推关系式 (4.1) 得:\par
\noindent X\par
\noindent n\par
\noindent Fn (x) = c(n, k)xk\par
\noindent k=1\par
\noindent Xn\par
\noindent = ((n − 1)c(n − 1, k) + c(n − 1, k − 1))xk\par
\noindent k=1\par
\noindent 68 第4章 两类 Stirling 数\par
\noindent X\par
\noindent n−1 X\par
\noindent n−1\par
\noindent = (n − 1) c(n − 1, k)xk + c(n − 1, k)xk+1\par
\noindent k=1 k=1\par
\noindent = (n − 1)Fn−1 (x) + xFn−1 (x)\par
\noindent = (n − 1 + x)Fn−1 (x),\par
\noindent 所以\par
\noindent Y\par
}
\end{sourceblock}


\begin{explainblock}
这一段是原书论述链的一部分。读的时候把它当作桥梁：它通常负责把上一个定义、例子或定理，引到下一个计数方法。

\smallskip
\textbf{初学者检查点：}
\begin{itemize}[leftmargin=2em,itemsep=0.25em]
\item 找出这一段承接的上一个概念。
\item 找出它准备引出的下一个概念。
\item 把中间的关键等式或构造单独抄出来。
\end{itemize}

\smallskip
\textbf{本章读法提醒：} 第四章区分两类 Stirling 数：第一类对应排列的圈，第二类对应集合划分。
\end{explainblock}


\begin{sourceblock}{第 4 章续读}
{\small\sloppy
\noindent n−1\par
\noindent Fn (x) = (n − k + x)F1 (x) = x(x + 1) · · · (x + n − 1).\par
\noindent k=1\par
}
\end{sourceblock}


\begin{explainblock}
这一段是原书论述链的一部分。读的时候把它当作桥梁：它通常负责把上一个定义、例子或定理，引到下一个计数方法。

\smallskip
\textbf{初学者检查点：}
\begin{itemize}[leftmargin=2em,itemsep=0.25em]
\item 找出这一段承接的上一个概念。
\item 找出它准备引出的下一个概念。
\item 把中间的关键等式或构造单独抄出来。
\end{itemize}

\smallskip
\textbf{本章读法提醒：} 第四章区分两类 Stirling 数：第一类对应排列的圈，第二类对应集合划分。
\end{explainblock}


\begin{sourceblock}{推论 4.1.5 设 n ≥ 0, 则}
{\small\sloppy
\noindent 推论 4.1.5 设 n ≥ 0, 则\par
\noindent X\par
\noindent n\par
\noindent s(n, k)xk = (x)n . (4.3)\par
\noindent k=0\par
\noindent 证明: 因为 c(n, k) = (−1)n−k s(n, k), 所以\par
\noindent X\par
\noindent n\par
\noindent s(n, k)(−1)n (−x)k = x(x + 1) · · · (x + n − 1),\par
\noindent k=0\par
\noindent 两边同时乘以 (−1)n 得\par
\noindent X\par
\noindent n\par
\noindent s(n, k)(−x)k = (−x)(−x − 1) · · · (−x − n + 1),\par
\noindent k=0\par
\noindent 将 −x 换成 x 立得\par
\noindent X\par
\noindent n\par
}
\end{sourceblock}


\begin{explainblock}
定理段不要只看结论，要看它把哪两个计数方式连在一起。组合学证明最常见的技巧是：左边按一种标准数，右边按另一种标准数，然后说明两边数的是同一批对象。

\smallskip
\textbf{初学者检查点：}
\begin{itemize}[leftmargin=2em,itemsep=0.25em]
\item 先复述定理的假设，不要把适用范围扩大。
\item 找出证明中的基本计数原则：乘法、加法、双射、递推、容斥或生成函数。
\item 检查最后的等式化简是否和前面的对象解释一致。
\end{itemize}

\smallskip
\textbf{本章读法提醒：} 第四章区分两类 Stirling 数：第一类对应排列的圈，第二类对应集合划分。
\end{explainblock}


\begin{sourceblock}{第 4 章续读}
{\small\sloppy
\noindent s(n, k)xk = x(x − 1) · · · (x − n + 1) = (x)n .\par
\noindent k=0\par
}
\end{sourceblock}


\begin{explainblock}
这一段是原书论述链的一部分。读的时候把它当作桥梁：它通常负责把上一个定义、例子或定理，引到下一个计数方法。

\smallskip
\textbf{初学者检查点：}
\begin{itemize}[leftmargin=2em,itemsep=0.25em]
\item 找出这一段承接的上一个概念。
\item 找出它准备引出的下一个概念。
\item 把中间的关键等式或构造单独抄出来。
\end{itemize}

\smallskip
\textbf{本章读法提醒：} 第四章区分两类 Stirling 数：第一类对应排列的圈，第二类对应集合划分。
\end{explainblock}


\begin{sourceblock}{§ 4.2 集合划分与第二类 Stirling 数}
{\small\sloppy
\noindent § 4.2. 集合划分与第二类 Stirling 数\par
}
\end{sourceblock}


\begin{explainblock}
这一节先不要急着背公式。先问三个问题：对象是什么，哪些限制条件不能丢，最后到底是在数所有对象、还是在数某个统计量的分布。把这三件事说清楚，后面的公式才会有位置。

\smallskip
\textbf{初学者检查点：}
\begin{itemize}[leftmargin=2em,itemsep=0.25em]
\item 把本节出现的新对象写成一句话定义。
\item 把每个公式左边和右边分别翻译成“在数什么”。
\item 找一个最小非平凡例子，手工列出所有对象。
\end{itemize}

\smallskip
\textbf{本章读法提醒：} 第四章区分两类 Stirling 数：第一类对应排列的圈，第二类对应集合划分。
\end{explainblock}


\begin{sourceblock}{定义 4.2.1 设 A 是一非空集合, π = \{A1 , A2 , . . . , Ak \} 是集合 A 的 k 个非空子集}
{\small\sloppy
\noindent 定义 4.2.1 设 A 是一非空集合, π = \{A1 , A2 , . . . , Ak \} 是集合 A 的 k 个非空子集\par
\noindent 组成的子集族. 若 A1 , A2 , . . . , Ak 满足:\par
\noindent (1) 当 i 6= j 时, Ai ∩ Aj = ϕ,\par
\noindent (2) A = A1 ∪ A2 ∪ A3 · · · ∪ Ak ,\par
\noindent 则称 π 是集合 A 的一个 k–部划分, 其中 A1 , A2 , . . . , Ak 称为划分 π 的块.\par
\noindent 换言之, 集合 A 的一个 k–部划分就是将 A 分成 k 个互不相交的集合或者按某\par
\noindent 种等价关系将集合 A 分成了 k 个等价类.\par
}
\end{sourceblock}


\begin{explainblock}
定义段的任务是定边界。这里要特别注意参数范围、是否允许重复、是否考虑顺序、对象有没有标签。后面的每一道例题和证明，本质上都在反复使用这些边界。

\smallskip
\textbf{初学者检查点：}
\begin{itemize}[leftmargin=2em,itemsep=0.25em]
\item 圈出被定义的对象名称。
\item 圈出参数范围，例如 n、k 是否要求正整数。
\item 判断它是有序对象、无序对象、带标签对象还是结构对象。
\end{itemize}

\smallskip
\textbf{本章读法提醒：} 第四章区分两类 Stirling 数：第一类对应排列的圈，第二类对应集合划分。
\end{explainblock}


\begin{sourceblock}{定义 4.2.2 一个 n 元集合的所有的 k–部划分个数称为第二类 Stirling 数, 记作}
{\small\sloppy
\noindent 定义 4.2.2 一个 n 元集合的所有的 k–部划分个数称为第二类 Stirling 数, 记作\par
\noindent S(n, k), 它所有的划分的个数称为 Bell 数 (Bell number or exponential number), 记\par
\noindent 作 B(n).\par
\noindent 4.2 集合划分与第二类 Stirling 数 69\par
\noindent 换言之, B(n) 计数了 n 元集合的所有的等价关系. 由定义可知:\par
\noindent X\par
\noindent n\par
\noindent B(n) = S(n, k).\par
\noindent k=0\par
\noindent 设 A 是一个有限集, 设\par
\noindent n o\par
\noindent π = \{a11 , a12 , . . . , a1i1 \}, \{a21 , a22 , . . . , a2i2 \}, . . . , \{ak1 , ak2 , . . . , akik \}\par
\noindent 是 A 的一个 k–部划分, 为了方便, 通常将 π 表示为\par
\noindent a11 a12 · · · a1i1 |a21 a22 · · · a2i2 | · · · |ak1 ak2 · · · akik .\par
}
\end{sourceblock}


\begin{explainblock}
定义段的任务是定边界。这里要特别注意参数范围、是否允许重复、是否考虑顺序、对象有没有标签。后面的每一道例题和证明，本质上都在反复使用这些边界。

\smallskip
\textbf{初学者检查点：}
\begin{itemize}[leftmargin=2em,itemsep=0.25em]
\item 圈出被定义的对象名称。
\item 圈出参数范围，例如 n、k 是否要求正整数。
\item 判断它是有序对象、无序对象、带标签对象还是结构对象。
\end{itemize}

\smallskip
\textbf{本章读法提醒：} 第四章区分两类 Stirling 数：第一类对应排列的圈，第二类对应集合划分。
\end{explainblock}


\begin{sourceblock}{例 4.2.1 集合 \{1, 2, 3, 4\} 的所有划分: 块数为 1 时, 只有 \{1234\}; 块数为 2 时,}
{\small\sloppy
\noindent 例 4.2.1 集合 \{1, 2, 3, 4\} 的所有划分: 块数为 1 时, 只有 \{1234\}; 块数为 2 时,\par
\noindent 有:\par
\noindent 1|234, 2|134, 3|124, 4|123, 12|34, 13|24, 14|23;\par
\noindent 块数为 3 块时, 有:\par
\noindent 1|2|34, 1|3|24. 1|4|23, 2|3|14, 2|4|13, 3|4|12;\par
\noindent 块数为 4 时, 只有 1|2|3|4.\par
\noindent 所以\par
\noindent S(4, 1) = S(4, 4) = 1, S(4, 2) = 7, S(4, 3) = 6,\par
\noindent 从而\par
\noindent B(4) = S(4, 1) + S(4, 2) + S(4, 3) + S(4, 4) = 15.\par
\noindent 第二类 Stirling 数是一类重要的计数函数. 对任意一对正整数 n, k, 由定义知:\par
\noindent S(n, n) = S(n, 1) = 1,\par
\noindent 且当 k > n 时, S(n, k) = 0. 规定 n ≥ 1 时,\par
\noindent S(n, 0) = S(0, n) = 0,\par
\noindent 并令 S(0, 0) = 1. 此时, 第二类 Stirling 数 S(n, k) 对任意一对非负整数 n, k 都有定\par
\noindent 义.\par
}
\end{sourceblock}


\begin{explainblock}
例题是学习这本书最重要的部分。不要只看答案，要把每个限制条件变成一个计数动作：先安排谁、再插入谁、有没有重复、需不需要除以对称性。

\smallskip
\textbf{初学者检查点：}
\begin{itemize}[leftmargin=2em,itemsep=0.25em]
\item 先写出没有限制时的总数。
\item 逐条加入限制条件，观察总数怎样变化。
\item 最后检查有没有把同一种方案重复算了多次。
\end{itemize}

\smallskip
\textbf{本章读法提醒：} 第四章区分两类 Stirling 数：第一类对应排列的圈，第二类对应集合划分。
\end{explainblock}


\begin{sourceblock}{定理 4.2.3 当 n ≥ 1, k ≥ 1,}
{\small\sloppy
\noindent 定理 4.2.3 当 n ≥ 1, k ≥ 1,\par
\noindent S(n, k) = S(n − 1, k − 1) + kS(n − 1, k). (4.4)\par
\noindent 证明: 由 S(n, k) 的定义可知当 n < k 及 n = k 时, 可验证 (4.4) 成立. 因此,\par
\noindent 只需证 n > k ≥ 1 时, (4.4) 成立即可.\par
\noindent 当 n > k ≥ 1 时, 设 π 是 [n] 的一个 k–部划分, 则\par
\noindent 70 第4章 两类 Stirling 数\par
\noindent (1) 若 \{n\} 是划分 π 中的一块, 则这类划分的个数与 [n − 1] 的 k − 1–部划分的个\par
\noindent 数相等, 即 S(n − 1, k − 1);\par
\noindent (2) 若 \{n\} 不是划分 π 中的一块, 则这类划分可通过将 n 插入到 [n − 1] 的 k–部\par
\noindent 划分中的任意一块得到, 所以它的个数为: kS(n − 1, k).\par
\noindent 所以当 n > k ≥ 1 时, (4.4) 也成立.\par
}
\end{sourceblock}


\begin{explainblock}
定理段不要只看结论，要看它把哪两个计数方式连在一起。组合学证明最常见的技巧是：左边按一种标准数，右边按另一种标准数，然后说明两边数的是同一批对象。

\smallskip
\textbf{初学者检查点：}
\begin{itemize}[leftmargin=2em,itemsep=0.25em]
\item 先复述定理的假设，不要把适用范围扩大。
\item 找出证明中的基本计数原则：乘法、加法、双射、递推、容斥或生成函数。
\item 检查最后的等式化简是否和前面的对象解释一致。
\end{itemize}

\smallskip
\textbf{本章读法提醒：} 第四章区分两类 Stirling 数：第一类对应排列的圈，第二类对应集合划分。
\end{explainblock}


\begin{sourceblock}{定理得证.}
{\small\sloppy
\noindent 定理得证.\par
\noindent k\par
\noindent 0 1 2 3 4 5 6 7 8 9\par
\noindent n\par
\noindent 6 0 1 31 90 65 15 1\par
\noindent 7 0 1 63 301 350 140 21 1\par
\noindent 8 0 1 127 966 1701 1050 266 28 1\par
\noindent 9 0 1 255 3025 7770 6951 2646 462 36 1\par
}
\end{sourceblock}


\begin{explainblock}
定理段不要只看结论，要看它把哪两个计数方式连在一起。组合学证明最常见的技巧是：左边按一种标准数，右边按另一种标准数，然后说明两边数的是同一批对象。

\smallskip
\textbf{初学者检查点：}
\begin{itemize}[leftmargin=2em,itemsep=0.25em]
\item 先复述定理的假设，不要把适用范围扩大。
\item 找出证明中的基本计数原则：乘法、加法、双射、递推、容斥或生成函数。
\item 检查最后的等式化简是否和前面的对象解释一致。
\end{itemize}

\smallskip
\textbf{本章读法提醒：} 第四章区分两类 Stirling 数：第一类对应排列的圈，第二类对应集合划分。
\end{explainblock}


\begin{sourceblock}{定理 4.2.4 设 k ≥ 0, 则}
{\small\sloppy
\noindent 定理 4.2.4 设 k ≥ 0, 则\par
\noindent ∞\par
\noindent X xk\par
\noindent S(n, k)xn = . (4.5)\par
\noindent (1 − x)(1 − 2x) · · · (1 − kx)\par
\noindent n=0\par
\noindent 证明: 不妨设\par
\noindent ∞\par
\noindent X ∞\par
\noindent X\par
\noindent Bk (x) = S(n, k)xn = S(n, k)xn . (4.6)\par
\noindent n=0 n=k\par
\noindent 当 k = 0 时, B0 = S(0, 0) = 1, 结论成立.\par
\noindent 当 k ≥ 1 时, 由等式 (4.4) 得\par
\noindent ∞\par
\noindent X\par
\noindent Bk (x) = (S(n − 1, k − 1) + kS(n − 1, k))xn\par
\noindent n=1\par
}
\end{sourceblock}


\begin{explainblock}
定理段不要只看结论，要看它把哪两个计数方式连在一起。组合学证明最常见的技巧是：左边按一种标准数，右边按另一种标准数，然后说明两边数的是同一批对象。

\smallskip
\textbf{初学者检查点：}
\begin{itemize}[leftmargin=2em,itemsep=0.25em]
\item 先复述定理的假设，不要把适用范围扩大。
\item 找出证明中的基本计数原则：乘法、加法、双射、递推、容斥或生成函数。
\item 检查最后的等式化简是否和前面的对象解释一致。
\end{itemize}

\smallskip
\textbf{本章读法提醒：} 第四章区分两类 Stirling 数：第一类对应排列的圈，第二类对应集合划分。
\end{explainblock}


\begin{sourceblock}{第 4 章续读}
{\small\sloppy
\noindent ∞\par
\noindent X ∞\par
\noindent X\par
\noindent = S(n − 1, k − 1)xn + k S(n − 1, k)xn\par
\noindent n=1 n=1\par
\noindent ∞\par
\noindent X ∞\par
\noindent X\par
\noindent =x S(n, k − 1)x + kx\par
\noindent n\par
\noindent S(n, k)xn\par
\noindent n=0 n=0\par
\noindent = xBk−1 (x) + kxBk (x)\par
\noindent 4.2 集合划分与第二类 Stirling 数 71\par
\noindent 所以\par
\noindent x\par
\noindent Bk (x) = Bk−1 (x).\par
\noindent 从而, 当 k ≥ 1 时,\par
}
\end{sourceblock}


\begin{explainblock}
这一段是原书论述链的一部分。读的时候把它当作桥梁：它通常负责把上一个定义、例子或定理，引到下一个计数方法。

\smallskip
\textbf{初学者检查点：}
\begin{itemize}[leftmargin=2em,itemsep=0.25em]
\item 找出这一段承接的上一个概念。
\item 找出它准备引出的下一个概念。
\item 把中间的关键等式或构造单独抄出来。
\end{itemize}

\smallskip
\textbf{本章读法提醒：} 第四章区分两类 Stirling 数：第一类对应排列的圈，第二类对应集合划分。
\end{explainblock}


\begin{sourceblock}{第 4 章续读}
{\small\sloppy
\noindent xk\par
\noindent Bk (x) = .\par
\noindent (1 − x)(1 − 2x) · · · (1 − kx)\par
\noindent 综上, 定理得证.\par
}
\end{sourceblock}


\begin{explainblock}
这一段是原书论述链的一部分。读的时候把它当作桥梁：它通常负责把上一个定义、例子或定理，引到下一个计数方法。

\smallskip
\textbf{初学者检查点：}
\begin{itemize}[leftmargin=2em,itemsep=0.25em]
\item 找出这一段承接的上一个概念。
\item 找出它准备引出的下一个概念。
\item 把中间的关键等式或构造单独抄出来。
\end{itemize}

\smallskip
\textbf{本章读法提醒：} 第四章区分两类 Stirling 数：第一类对应排列的圈，第二类对应集合划分。
\end{explainblock}


\begin{sourceblock}{定理 4.2.5 设 n ≥ 0, k ≥ 0, 则}
{\small\sloppy
\noindent 定理 4.2.5 设 n ≥ 0, k ≥ 0, 则\par
\noindent k\par
\noindent i k\par
\noindent S(n, k) = (−1) (k − i)n . (4.7)\par
\noindent k! i\par
\noindent i=0\par
\noindent 证明: (方法一) 由定理 4.2.4 知我们可通过将\par
\noindent xk\par
\noindent (1 − x)(1 − 2x) · · · (1 − kx)\par
\noindent 写成幂级数的形式来求 S(n, k) 的表达式.\par
\noindent 设\par
\noindent = .\par
\noindent (1 − x)(1 − 2x) · · · (1 − kx) 1 − jx\par
\noindent j=1\par
\noindent 对于固定的 j(1 ≤ j ≤ k), 两边同时乘以 (1 − jx) 后再令 x = 1j , 求得\par
\noindent αj =\par
\noindent (1 − 1j )(1 − 2j ) · · · (1 − j−1 j+1\par
\noindent j )(1 − j ) · · · (1 − j )\par
}
\end{sourceblock}


\begin{explainblock}
定理段不要只看结论，要看它把哪两个计数方式连在一起。组合学证明最常见的技巧是：左边按一种标准数，右边按另一种标准数，然后说明两边数的是同一批对象。

\smallskip
\textbf{初学者检查点：}
\begin{itemize}[leftmargin=2em,itemsep=0.25em]
\item 先复述定理的假设，不要把适用范围扩大。
\item 找出证明中的基本计数原则：乘法、加法、双射、递推、容斥或生成函数。
\item 检查最后的等式化简是否和前面的对象解释一致。
\end{itemize}

\smallskip
\textbf{本章读法提醒：} 第四章区分两类 Stirling 数：第一类对应排列的圈，第二类对应集合划分。
\end{explainblock}


\begin{sourceblock}{第 4 章续读}
{\small\sloppy
\noindent k\par
\noindent k\par
\noindent k−j k j\par
\noindent = (−1) ,\par
\noindent j k!\par
\noindent 所以\par
\noindent X\par
\noindent k k X\par
\noindent ∞\par
\noindent xk k k−j k j\par
\noindent =x (−1) j m xm\par
\noindent (1 − x)(1 − 2x) · · · (1 − kx) j k!\par
\noindent j=1 m=0\par
\noindent ∞ X\par
\noindent X k m+k\par
\noindent k j\par
\noindent = (−1)k−j xm+k ,\par
\noindent j k!\par
}
\end{sourceblock}


\begin{explainblock}
这一段是原书论述链的一部分。读的时候把它当作桥梁：它通常负责把上一个定义、例子或定理，引到下一个计数方法。

\smallskip
\textbf{初学者检查点：}
\begin{itemize}[leftmargin=2em,itemsep=0.25em]
\item 找出这一段承接的上一个概念。
\item 找出它准备引出的下一个概念。
\item 把中间的关键等式或构造单独抄出来。
\end{itemize}

\smallskip
\textbf{本章读法提醒：} 第四章区分两类 Stirling 数：第一类对应排列的圈，第二类对应集合划分。
\end{explainblock}


\begin{sourceblock}{第 4 章续读}
{\small\sloppy
\noindent m=0 j=1\par
\noindent ∞ X\par
\noindent X k n\par
\noindent k−j k j\par
\noindent = (−1) xn . (4.8)\par
\noindent j k!\par
\noindent n=k j=1\par
\noindent 比较 (4.6) 与 (4.8) 中 xn 的系数可得\par
\noindent X\par
\noindent k n\par
\noindent k−j k j\par
\noindent S(n, k) = (−1)\par
\noindent j k!\par
\noindent j=1\par
\noindent 72 第4章 两类 Stirling 数\par
\noindent X\par
\noindent k\par
\noindent k (k − i)n\par
}
\end{sourceblock}


\begin{explainblock}
这一段是原书论述链的一部分。读的时候把它当作桥梁：它通常负责把上一个定义、例子或定理，引到下一个计数方法。

\smallskip
\textbf{初学者检查点：}
\begin{itemize}[leftmargin=2em,itemsep=0.25em]
\item 找出这一段承接的上一个概念。
\item 找出它准备引出的下一个概念。
\item 把中间的关键等式或构造单独抄出来。
\end{itemize}

\smallskip
\textbf{本章读法提醒：} 第四章区分两类 Stirling 数：第一类对应排列的圈，第二类对应集合划分。
\end{explainblock}


\begin{sourceblock}{第 4 章续读}
{\small\sloppy
\noindent = (−1)i .\par
\noindent i k!\par
\noindent i=0\par
\noindent (方法二) 将 n 个不同的球放到标记了 1, 2, . . . , k 的 k 个盒子里, 允许盒子为空, 用\par
\noindent S 表示所有这样的放法组成的集合, 显然 |S| = k n . 对任意 1 ≤ i ≤ k, 定义 Pi 为性\par
\noindent 质 “第 i 个盒子是空的”, Ai 为 S 中满足性质 Pi 的放法组成的集合, 则由 S(n, k) 的\par
\noindent 组合意义可知:\par
\noindent k!S(n, k) = |Ā1 ∩ Ā2 ∩ · · · ∩ Āk |\par
\noindent 注意到对于 1 ≤ i1 < i2 < · · · < ij ≤ k, Ai1 ∩ Ai2 ∩ · · · ∩ Aij 表示的意义是 S 中满足\par
\noindent 性质 Pi1 , Pi2 , . . . , Pij 的放法组成的集合, 即标号为 i1 , i2 , . . . , ij 的盒子为空, 因此 n\par
\noindent 个球只能放进其它 k − j 个盒子中, 从而\par
\noindent |Ai1 ∩ Ai2 ∩ · · · ∩ Aij | = (k − j)n .\par
\noindent 由容斥原理得\par
\noindent X\par
\noindent k X\par
\noindent k!S(n, k) = |S| − (−1)j−1 |Ai1 ∩ Ai2 ∩ · · · ∩ Aij |\par
\noindent j=1 1≤i1 <i2 <···<ij ≤k\par
\noindent X\par
}
\end{sourceblock}


\begin{explainblock}
这一段是原书论述链的一部分。读的时候把它当作桥梁：它通常负责把上一个定义、例子或定理，引到下一个计数方法。

\smallskip
\textbf{初学者检查点：}
\begin{itemize}[leftmargin=2em,itemsep=0.25em]
\item 找出这一段承接的上一个概念。
\item 找出它准备引出的下一个概念。
\item 把中间的关键等式或构造单独抄出来。
\end{itemize}

\smallskip
\textbf{本章读法提醒：} 第四章区分两类 Stirling 数：第一类对应排列的圈，第二类对应集合划分。
\end{explainblock}


\begin{sourceblock}{第 4 章续读}
{\small\sloppy
\noindent k\par
\noindent j k\par
\noindent = (−1) (k − j)n .\par
\noindent j\par
\noindent j=0\par
}
\end{sourceblock}


\begin{explainblock}
这一段是原书论述链的一部分。读的时候把它当作桥梁：它通常负责把上一个定义、例子或定理，引到下一个计数方法。

\smallskip
\textbf{初学者检查点：}
\begin{itemize}[leftmargin=2em,itemsep=0.25em]
\item 找出这一段承接的上一个概念。
\item 找出它准备引出的下一个概念。
\item 把中间的关键等式或构造单独抄出来。
\end{itemize}

\smallskip
\textbf{本章读法提醒：} 第四章区分两类 Stirling 数：第一类对应排列的圈，第二类对应集合划分。
\end{explainblock}


\begin{sourceblock}{定理得证.}
{\small\sloppy
\noindent 定理得证.\par
}
\end{sourceblock}


\begin{explainblock}
定理段不要只看结论，要看它把哪两个计数方式连在一起。组合学证明最常见的技巧是：左边按一种标准数，右边按另一种标准数，然后说明两边数的是同一批对象。

\smallskip
\textbf{初学者检查点：}
\begin{itemize}[leftmargin=2em,itemsep=0.25em]
\item 先复述定理的假设，不要把适用范围扩大。
\item 找出证明中的基本计数原则：乘法、加法、双射、递推、容斥或生成函数。
\item 检查最后的等式化简是否和前面的对象解释一致。
\end{itemize}

\smallskip
\textbf{本章读法提醒：} 第四章区分两类 Stirling 数：第一类对应排列的圈，第二类对应集合划分。
\end{explainblock}


\begin{sourceblock}{定理 4.2.6 设 k ≥ 0, 则 \{S(n, k)\}n≥0 生成函数为}
{\small\sloppy
\noindent 定理 4.2.6 设 k ≥ 0, 则 \{S(n, k)\}n≥0 生成函数为\par
\noindent ∞\par
\noindent X xn (ex − 1)k\par
\noindent S(n, k) = . (4.9)\par
\noindent n! k!\par
\noindent n=0\par
\noindent 证明: 利用定理 4.2.5 得到的 S(n, k) 的显式表达式可得\par
\noindent ∞\par
\noindent X ∞ k\par
\noindent xn 1 XX j k xn\par
\noindent S(n, k) = (−1) (k − j)n\par
\noindent n! k! j n!\par
\noindent n=0 n=0 j=0\par
\noindent X\par
\noindent ∞\par
\noindent k\par
\noindent k xn\par
\noindent = (−1)j (k − j)n\par
}
\end{sourceblock}


\begin{explainblock}
定理段不要只看结论，要看它把哪两个计数方式连在一起。组合学证明最常见的技巧是：左边按一种标准数，右边按另一种标准数，然后说明两边数的是同一批对象。

\smallskip
\textbf{初学者检查点：}
\begin{itemize}[leftmargin=2em,itemsep=0.25em]
\item 先复述定理的假设，不要把适用范围扩大。
\item 找出证明中的基本计数原则：乘法、加法、双射、递推、容斥或生成函数。
\item 检查最后的等式化简是否和前面的对象解释一致。
\end{itemize}

\smallskip
\textbf{本章读法提醒：} 第四章区分两类 Stirling 数：第一类对应排列的圈，第二类对应集合划分。
\end{explainblock}


\begin{sourceblock}{第 4 章续读}
{\small\sloppy
\noindent k! j n!\par
\noindent j=0 n=0\par
\noindent k\par
\noindent j k\par
\noindent = (−1) e(k−j)x\par
\noindent k! j\par
\noindent j=0\par
\noindent (ex − 1)k\par
\noindent = .\par
\noindent k!\par
\noindent 4.2 集合划分与第二类 Stirling 数 73\par
}
\end{sourceblock}


\begin{explainblock}
这一段是原书论述链的一部分。读的时候把它当作桥梁：它通常负责把上一个定义、例子或定理，引到下一个计数方法。

\smallskip
\textbf{初学者检查点：}
\begin{itemize}[leftmargin=2em,itemsep=0.25em]
\item 找出这一段承接的上一个概念。
\item 找出它准备引出的下一个概念。
\item 把中间的关键等式或构造单独抄出来。
\end{itemize}

\smallskip
\textbf{本章读法提醒：} 第四章区分两类 Stirling 数：第一类对应排列的圈，第二类对应集合划分。
\end{explainblock}


\begin{sourceblock}{定理 4.2.7 设 B(n) 是第 n 个 Bell 数, 则}
{\small\sloppy
\noindent 定理 4.2.7 设 B(n) 是第 n 个 Bell 数, 则\par
\noindent ∞\par
\noindent X xn\par
\noindent = ee −1 .\par
\noindent x\par
\noindent B(n)\par
\noindent n!\par
\noindent n=0\par
\noindent 证明: 因为\par
\noindent X\par
\noindent n\par
\noindent B(n) = S(n, k),\par
\noindent k=0\par
\noindent 故\par
\noindent ∞\par
\noindent X ∞\par
\noindent xn X X\par
\noindent n\par
}
\end{sourceblock}


\begin{explainblock}
定理段不要只看结论，要看它把哪两个计数方式连在一起。组合学证明最常见的技巧是：左边按一种标准数，右边按另一种标准数，然后说明两边数的是同一批对象。

\smallskip
\textbf{初学者检查点：}
\begin{itemize}[leftmargin=2em,itemsep=0.25em]
\item 先复述定理的假设，不要把适用范围扩大。
\item 找出证明中的基本计数原则：乘法、加法、双射、递推、容斥或生成函数。
\item 检查最后的等式化简是否和前面的对象解释一致。
\end{itemize}

\smallskip
\textbf{本章读法提醒：} 第四章区分两类 Stirling 数：第一类对应排列的圈，第二类对应集合划分。
\end{explainblock}


\begin{sourceblock}{第 4 章续读}
{\small\sloppy
\noindent xn\par
\noindent B(n) = S(n, k)\par
\noindent n! n!\par
\noindent n=0 n=0 k=0\par
\noindent ∞\par
\noindent ∞ X\par
\noindent X xn\par
\noindent = S(n, k)\par
\noindent n!\par
\noindent k=0 n=k\par
\noindent ∞\par
\noindent X (ex − 1)k\par
\noindent =\par
\noindent k!\par
\noindent k=0\par
\noindent = ee −1 .\par
\noindent x\par
}
\end{sourceblock}


\begin{explainblock}
这一段是原书论述链的一部分。读的时候把它当作桥梁：它通常负责把上一个定义、例子或定理，引到下一个计数方法。

\smallskip
\textbf{初学者检查点：}
\begin{itemize}[leftmargin=2em,itemsep=0.25em]
\item 找出这一段承接的上一个概念。
\item 找出它准备引出的下一个概念。
\item 把中间的关键等式或构造单独抄出来。
\end{itemize}

\smallskip
\textbf{本章读法提醒：} 第四章区分两类 Stirling 数：第一类对应排列的圈，第二类对应集合划分。
\end{explainblock}


\begin{sourceblock}{定理 4.2.8 设 n ≥ 0, 则对任意实数 x,}
{\small\sloppy
\noindent 定理 4.2.8 设 n ≥ 0, 则对任意实数 x,\par
\noindent X\par
\noindent n\par
\noindent S(n, k)(x)k = xn . (4.10)\par
\noindent k=0\par
\noindent 证明: 当 n = 0 时, 显然成立. 当 n ≥ 1, 此时 S(n, 0) = 0. 所以只需证明对任\par
\noindent 意的实数 x,\par
\noindent X\par
\noindent n\par
\noindent S(n, k)(x)k = xn .\par
\noindent k=1\par
\noindent 不妨设 x 为正整数, S 为集合 [n] 到 x 元集合上的映射组成的集合, 则 |S| = xn .\par
\noindent 任取 φ ∈ S, 若对 i, j ∈ [n], 且 i 6= j, 我们 φ(i) = φ(j), 则称 i 与 j 等价; 从而 φ 确\par
\noindent 定了集合 [n] 的一个等价关系, 进而确定了 [n] 的一个划分. 因此 S 中的任一映射都\par
\noindent 对应集合 [n] 的一个划分. 反之, 集合 [n] 的任意一个 k–部划分显然对应 (x)k 个不\par
\noindent 同的映射. 故当 x 为正整数时,\par
\noindent X\par
\noindent n\par
}
\end{sourceblock}


\begin{explainblock}
定理段不要只看结论，要看它把哪两个计数方式连在一起。组合学证明最常见的技巧是：左边按一种标准数，右边按另一种标准数，然后说明两边数的是同一批对象。

\smallskip
\textbf{初学者检查点：}
\begin{itemize}[leftmargin=2em,itemsep=0.25em]
\item 先复述定理的假设，不要把适用范围扩大。
\item 找出证明中的基本计数原则：乘法、加法、双射、递推、容斥或生成函数。
\item 检查最后的等式化简是否和前面的对象解释一致。
\end{itemize}

\smallskip
\textbf{本章读法提醒：} 第四章区分两类 Stirling 数：第一类对应排列的圈，第二类对应集合划分。
\end{explainblock}


\begin{sourceblock}{第 4 章续读}
{\small\sloppy
\noindent xn = S(n, k)(x)k . (4.11)\par
\noindent k=1\par
\noindent 设\par
\noindent X\par
\noindent n\par
\noindent f (x) = xn − Sn,k (x)k .\par
\noindent k=1\par
\noindent 74 第4章 两类 Stirling 数\par
\noindent n\par
\noindent 因为 (x)k 是系数为 1 的 k 次多项式, 且 S(n, n) = 1, S(n, n − 1) = 6= 0, 所\par
\noindent 以 f (x) 是一个 n − 1 次的多项式. 由 (4.11) 可知 x = 1, 2, . . . , n 都是多项式 f (x)\par
\noindent 的根, 即 n − 1 次多项式 f (x) 有 n 个根, 故 f (x) ≡ 0, 所以对任意的实数 x, 等式\par
\noindent (4.10) 都成立.\par
}
\end{sourceblock}


\begin{explainblock}
这一段是原书论述链的一部分。读的时候把它当作桥梁：它通常负责把上一个定义、例子或定理，引到下一个计数方法。

\smallskip
\textbf{初学者检查点：}
\begin{itemize}[leftmargin=2em,itemsep=0.25em]
\item 找出这一段承接的上一个概念。
\item 找出它准备引出的下一个概念。
\item 把中间的关键等式或构造单独抄出来。
\end{itemize}

\smallskip
\textbf{本章读法提醒：} 第四章区分两类 Stirling 数：第一类对应排列的圈，第二类对应集合划分。
\end{explainblock}


\begin{sourceblock}{定理得证.}
{\small\sloppy
\noindent 定理得证.\par
}
\end{sourceblock}


\begin{explainblock}
定理段不要只看结论，要看它把哪两个计数方式连在一起。组合学证明最常见的技巧是：左边按一种标准数，右边按另一种标准数，然后说明两边数的是同一批对象。

\smallskip
\textbf{初学者检查点：}
\begin{itemize}[leftmargin=2em,itemsep=0.25em]
\item 先复述定理的假设，不要把适用范围扩大。
\item 找出证明中的基本计数原则：乘法、加法、双射、递推、容斥或生成函数。
\item 检查最后的等式化简是否和前面的对象解释一致。
\end{itemize}

\smallskip
\textbf{本章读法提醒：} 第四章区分两类 Stirling 数：第一类对应排列的圈，第二类对应集合划分。
\end{explainblock}


\begin{sourceblock}{定理 4.2.9 设 n, m 为非负整数, 则}
{\small\sloppy
\noindent 定理 4.2.9 设 n, m 为非负整数, 则\par
\noindent (\par
\noindent X\par
\noindent n\par
\noindent 1, m = n;\par
\noindent S(n, k)s(k, m) = δnm = (4.12)\par
\noindent k=1\par
\noindent 0, m 6= n.\par
\noindent 证明: (方法一) 由推论 4.1.5 和定理 4.2.8 知:\par
\noindent X\par
\noindent n X\par
\noindent n\par
\noindent (x)n = s(n, k)xk , xn = S(n, k)(x)k ,\par
\noindent k=0 k=0\par
\noindent 所以\par
\noindent X\par
\noindent n\par
\noindent n\par
}
\end{sourceblock}


\begin{explainblock}
定理段不要只看结论，要看它把哪两个计数方式连在一起。组合学证明最常见的技巧是：左边按一种标准数，右边按另一种标准数，然后说明两边数的是同一批对象。

\smallskip
\textbf{初学者检查点：}
\begin{itemize}[leftmargin=2em,itemsep=0.25em]
\item 先复述定理的假设，不要把适用范围扩大。
\item 找出证明中的基本计数原则：乘法、加法、双射、递推、容斥或生成函数。
\item 检查最后的等式化简是否和前面的对象解释一致。
\end{itemize}

\smallskip
\textbf{本章读法提醒：} 第四章区分两类 Stirling 数：第一类对应排列的圈，第二类对应集合划分。
\end{explainblock}


\begin{sourceblock}{第 4 章续读}
{\small\sloppy
\noindent x = S(n, k)(x)k\par
\noindent k=0\par
\noindent !\par
\noindent X\par
\noindent n X\par
\noindent k\par
\noindent = S(n, k) s(k, m)xm\par
\noindent k=0 m=0\par
\noindent n X\par
\noindent X n\par
\noindent = S(n, k)s(k, m)xm ,\par
\noindent m=0 k=m\par
\noindent 比较两边 xm (m = 0, 1, . . . , n) 的系数即可.\par
\noindent (方法二): 将标记有 1, 2, . . . , n 的 n 个不同球放入到 k(0 ≤ m ≤ k ≤ n) 个相同\par
\noindent 的盒子中后 (盒子不可为空), 再将这 k 个盒子摆放成 m 个圈, 每个圈中盒子的个数\par
\noindent 称为圈的长度 (单独一个盒子也被看作一个圈, 长度为 1). 将按上面规则得到的所有\par
\noindent 方案组成的集合记为 Am (n), 则显然有 |Am (n)| = S(n, k)c(k, m).\par
\noindent 对于 Am (n) 中的任意方案, 我们考虑圈中球的个数大于 1 的所有圈中标记最大\par
}
\end{sourceblock}


\begin{explainblock}
这一段是原书论述链的一部分。读的时候把它当作桥梁：它通常负责把上一个定义、例子或定理，引到下一个计数方法。

\smallskip
\textbf{初学者检查点：}
\begin{itemize}[leftmargin=2em,itemsep=0.25em]
\item 找出这一段承接的上一个概念。
\item 找出它准备引出的下一个概念。
\item 把中间的关键等式或构造单独抄出来。
\end{itemize}

\smallskip
\textbf{本章读法提醒：} 第四章区分两类 Stirling 数：第一类对应排列的圈，第二类对应集合划分。
\end{explainblock}


\begin{sourceblock}{第 4 章续读}
{\small\sloppy
\noindent 的球. 若它单独在一个盒子, 则将它拿出放到它后面的盒子里, 并把它之前所在的\par
\noindent 盒子去掉; 若除了它之外还有其他的球, 则把它取出放到一个新盒子里, 然后把新盒\par
\noindent 子放到它之前所在盒子的前面. 显然对于 m 6= n, 这个变换是一个对合变换, 它将\par
\noindent Am (n) 中含有 k 个盒子的方案映到 Am (n) 中含有 k + 1 或 k − 1 个盒子的方案. 本\par
\noindent 质上对于固定的 m 6= n, 这个对合将盒子数为奇数的方案与盒子数为偶数的方案一\par
\noindent 一对应, 而对于 m = n, 方案只有一种. 所以\par
\noindent X\par
\noindent n X\par
\noindent n\par
\noindent S(n, k)s(k, m) = (−1)k−m S(n, k)c(k, m) = δmn .\par
\noindent k=m k=m\par
\noindent 4.2 集合划分与第二类 Stirling 数 75\par
\noindent 记 Sn = (S(i, j))n×n , sn = (s(i, j))n×n 分别是由第一类 Stirling 数和第二类\par
\noindent Stirling 数得到的 n × n 矩阵. 由 (4.12) 可知:\par
\noindent Sn sn = I,\par
\noindent 即 Sn 与 sn 互为逆矩阵.\par
\noindent 设 V [x] 是由复系数多项式组成的向量空间. 设 B1 = (1, x, x2 , . . .)T , B2 =\par
\noindent (1, (x)1 , (x)2 , . . .)T . 则 B1 与 B2 都是 V [x] 的一组基. 由 (4.10) 与 (4.3) 可知,\par
}
\end{sourceblock}


\begin{explainblock}
这一段是原书论述链的一部分。读的时候把它当作桥梁：它通常负责把上一个定义、例子或定理，引到下一个计数方法。

\smallskip
\textbf{初学者检查点：}
\begin{itemize}[leftmargin=2em,itemsep=0.25em]
\item 找出这一段承接的上一个概念。
\item 找出它准备引出的下一个概念。
\item 把中间的关键等式或构造单独抄出来。
\end{itemize}

\smallskip
\textbf{本章读法提醒：} 第四章区分两类 Stirling 数：第一类对应排列的圈，第二类对应集合划分。
\end{explainblock}


\begin{sourceblock}{第 4 章续读}
{\small\sloppy
\noindent     \par
\noindent 1 S(0, 0) S(0, 1) · · · S(0, n) · · · 1\par
\noindent     \par
\noindent  x   S(1, 0) S(1, 1) · · · S(1, n) · · ·   (x)1 \par
\noindent     \par
\noindent  x2  =  S(2, 0) S(2, 1) · · · S(2, n) · · ·   (x)2  (4.13)\par
\noindent     \par
\noindent     \par
\noindent  ···   ··· ··· ··· ··· ···  ··· \par
\noindent ... ... ... ... ... ... ...\par
\noindent 和\par
\noindent     \par
\noindent 1 s(0, 0) s(0, 1) · · · s(0, n) · · · 1\par
\noindent     \par
\noindent  (x)1   s(1, 0) s(1, 1) · · · s(1, n) · · ·   x \par
\noindent     \par
\noindent  (x)2  =  s(2, 0) s(2, 1) · · · s(2, n) · · ·   x2  (4.14)\par
\noindent     \par
}
\end{sourceblock}


\begin{explainblock}
这一段是原书论述链的一部分。读的时候把它当作桥梁：它通常负责把上一个定义、例子或定理，引到下一个计数方法。

\smallskip
\textbf{初学者检查点：}
\begin{itemize}[leftmargin=2em,itemsep=0.25em]
\item 找出这一段承接的上一个概念。
\item 找出它准备引出的下一个概念。
\item 把中间的关键等式或构造单独抄出来。
\end{itemize}

\smallskip
\textbf{本章读法提醒：} 第四章区分两类 Stirling 数：第一类对应排列的圈，第二类对应集合划分。
\end{explainblock}


\begin{sourceblock}{第 4 章续读}
{\small\sloppy
\noindent     \par
\noindent  ···   ··· ··· ··· ··· ···  ··· \par
\noindent ... ... ... ... ... ... ...\par
\noindent 所以无穷矩阵 S 是基 B2 到 B1 之间的过度矩阵, 而无穷矩阵 s 是基 B1 到 B2\par
\noindent 的过渡矩阵, 因此 S 与 s 互逆. 所以从这个角度也能得到\par
\noindent X\par
\noindent n\par
\noindent S(n, k)s(k, m) = δmn .\par
\noindent k=m\par
}
\end{sourceblock}


\begin{explainblock}
这一段是原书论述链的一部分。读的时候把它当作桥梁：它通常负责把上一个定义、例子或定理，引到下一个计数方法。

\smallskip
\textbf{初学者检查点：}
\begin{itemize}[leftmargin=2em,itemsep=0.25em]
\item 找出这一段承接的上一个概念。
\item 找出它准备引出的下一个概念。
\item 把中间的关键等式或构造单独抄出来。
\end{itemize}

\smallskip
\textbf{本章读法提醒：} 第四章区分两类 Stirling 数：第一类对应排列的圈，第二类对应集合划分。
\end{explainblock}


\begin{sourceblock}{推论 4.2.10 设 \{an \}n≥0 和 \{bn \}n≥0 都是实数列. 则下面两个条件可以互相推出:}
{\small\sloppy
\noindent 推论 4.2.10 设 \{an \}n≥0 和 \{bn \}n≥0 都是实数列. 则下面两个条件可以互相推出:\par
\noindent (i) 对任意 n ≥ 0\par
\noindent X\par
\noindent n\par
\noindent bn = S(n, k)ak ;\par
\noindent k=0\par
\noindent (ii) 对任意 n ≥ 0\par
\noindent X\par
\noindent n\par
\noindent an = S(n, k)bk .\par
\noindent k=0\par
\noindent 证明: 不妨设 a = (a0 , a1 , a2 , . . .)T , b = (b0 , b1 , b2 , . . .)T , 则条件 (i) 说明 Sa = b,\par
\noindent 因此左乘 s 得到 a = sb, 也就是条件 (ii). 类似地, 由 (ii) 也能推出 (i).\par
\noindent 下面我们来介绍一下第二类 Stirling 数在求连续自然数的整数次幂和的应用.\par
}
\end{sourceblock}


\begin{explainblock}
定理段不要只看结论，要看它把哪两个计数方式连在一起。组合学证明最常见的技巧是：左边按一种标准数，右边按另一种标准数，然后说明两边数的是同一批对象。

\smallskip
\textbf{初学者检查点：}
\begin{itemize}[leftmargin=2em,itemsep=0.25em]
\item 先复述定理的假设，不要把适用范围扩大。
\item 找出证明中的基本计数原则：乘法、加法、双射、递推、容斥或生成函数。
\item 检查最后的等式化简是否和前面的对象解释一致。
\end{itemize}

\smallskip
\textbf{本章读法提醒：} 第四章区分两类 Stirling 数：第一类对应排列的圈，第二类对应集合划分。
\end{explainblock}


\begin{sourceblock}{定理 4.2.11}
{\small\sloppy
\noindent 定理 4.2.11\par
\noindent X\par
\noindent m\par
\noindent n+1\par
\noindent m m m\par
\noindent = S(m, k)k! .\par
\noindent k+1\par
\noindent k=0\par
\noindent 76 第4章 两类 Stirling 数\par
\noindent 证明: 由定理 4.2.8 知\par
\noindent X\par
\noindent m\par
\noindent m\par
\noindent x = S(m, k)(x)k ,\par
\noindent k=0\par
\noindent 所以\par
\noindent X\par
\noindent n n X\par
}
\end{sourceblock}


\begin{explainblock}
定理段不要只看结论，要看它把哪两个计数方式连在一起。组合学证明最常见的技巧是：左边按一种标准数，右边按另一种标准数，然后说明两边数的是同一批对象。

\smallskip
\textbf{初学者检查点：}
\begin{itemize}[leftmargin=2em,itemsep=0.25em]
\item 先复述定理的假设，不要把适用范围扩大。
\item 找出证明中的基本计数原则：乘法、加法、双射、递推、容斥或生成函数。
\item 检查最后的等式化简是否和前面的对象解释一致。
\end{itemize}

\smallskip
\textbf{本章读法提醒：} 第四章区分两类 Stirling 数：第一类对应排列的圈，第二类对应集合划分。
\end{explainblock}


\begin{sourceblock}{第 4 章续读}
{\small\sloppy
\noindent X m\par
\noindent m i\par
\noindent i = S(m, k)k!\par
\noindent k\par
\noindent i=1 i=1 k=0\par
\noindent Xm Xn\par
\noindent i\par
\noindent = S(m, k)k! .\par
\noindent k\par
\noindent k=0 i=0\par
\noindent 由帕斯卡公式 (1.2.5):\par
\noindent n n−1 n−1\par
\noindent = + ,\par
\noindent k k−1 k\par
\noindent 知\par
\noindent k k+1 n k+1 k+1 n\par
\noindent + + ··· + = + + ··· +\par
\noindent k k k k+1 k k\par
}
\end{sourceblock}


\begin{explainblock}
这一段是原书论述链的一部分。读的时候把它当作桥梁：它通常负责把上一个定义、例子或定理，引到下一个计数方法。

\smallskip
\textbf{初学者检查点：}
\begin{itemize}[leftmargin=2em,itemsep=0.25em]
\item 找出这一段承接的上一个概念。
\item 找出它准备引出的下一个概念。
\item 把中间的关键等式或构造单独抄出来。
\end{itemize}

\smallskip
\textbf{本章读法提醒：} 第四章区分两类 Stirling 数：第一类对应排列的圈，第二类对应集合划分。
\end{explainblock}


\begin{sourceblock}{第 4 章续读}
{\small\sloppy
\noindent k+2 k+2 n\par
\noindent = + + ··· +\par
\noindent k+1 k k\par
\noindent = ······\par
\noindent n n\par
\noindent = +\par
\noindent k+1 k\par
\noindent n+1\par
\noindent = .\par
\noindent k+1\par
\noindent 故而:\par
\noindent n\par
\noindent X n\par
\noindent X\par
\noindent i i n+1\par
\noindent = = . (4.15)\par
\noindent k k k+1\par
\noindent i=0 i=k\par
}
\end{sourceblock}


\begin{explainblock}
这一段是原书论述链的一部分。读的时候把它当作桥梁：它通常负责把上一个定义、例子或定理，引到下一个计数方法。

\smallskip
\textbf{初学者检查点：}
\begin{itemize}[leftmargin=2em,itemsep=0.25em]
\item 找出这一段承接的上一个概念。
\item 找出它准备引出的下一个概念。
\item 把中间的关键等式或构造单独抄出来。
\end{itemize}

\smallskip
\textbf{本章读法提醒：} 第四章区分两类 Stirling 数：第一类对应排列的圈，第二类对应集合划分。
\end{explainblock}


\begin{sourceblock}{第 4 章续读}
{\small\sloppy
\noindent 所以\par
\noindent X\par
\noindent n X\par
\noindent m n\par
\noindent X X\par
\noindent m\par
\noindent m i n+1\par
\noindent i = S(m, k)k! = S(m, k)k! .\par
\noindent k k+1\par
\noindent i=1 k=0 i=0 k=0\par
}
\end{sourceblock}


\begin{explainblock}
这一段是原书论述链的一部分。读的时候把它当作桥梁：它通常负责把上一个定义、例子或定理，引到下一个计数方法。

\smallskip
\textbf{初学者检查点：}
\begin{itemize}[leftmargin=2em,itemsep=0.25em]
\item 找出这一段承接的上一个概念。
\item 找出它准备引出的下一个概念。
\item 把中间的关键等式或构造单独抄出来。
\end{itemize}

\smallskip
\textbf{本章读法提醒：} 第四章区分两类 Stirling 数：第一类对应排列的圈，第二类对应集合划分。
\end{explainblock}


\begin{sourceblock}{例 4.2.2 计算}
{\small\sloppy
\noindent 例 4.2.2 计算\par
\noindent 14 + 24 + 34 + 44 + · · · + n4 .\par
\noindent 解:\par
\noindent 14 +24 + 34 + 44 + · · · + n4\par
\noindent X\par
\noindent n+1\par
\noindent = S(4, k)k!\par
\noindent k+1\par
\noindent k=0\par
\noindent n+1 n+1 n+1 n+1\par
\noindent = S(4, 1) + 2!S(4, 2) + 3!S(4, 3) + 4!S(4, 4)\par
\noindent n+1 n+1 n+1 n+1\par
\noindent = + 2! · 7 + 3! · 6 + 4!\par
\noindent = n(n + 1)(2n + 1) 3n2 + 3n − 1\par
\noindent 6Tn 2 − Tn (2n + 1)\par
\noindent = ,\par
\noindent 其中 Tn = n(n+1)\par
}
\end{sourceblock}


\begin{explainblock}
例题是学习这本书最重要的部分。不要只看答案，要把每个限制条件变成一个计数动作：先安排谁、再插入谁、有没有重复、需不需要除以对称性。

\smallskip
\textbf{初学者检查点：}
\begin{itemize}[leftmargin=2em,itemsep=0.25em]
\item 先写出没有限制时的总数。
\item 逐条加入限制条件，观察总数怎样变化。
\item 最后检查有没有把同一种方案重复算了多次。
\end{itemize}

\smallskip
\textbf{本章读法提醒：} 第四章区分两类 Stirling 数：第一类对应排列的圈，第二类对应集合划分。
\end{explainblock}


\begin{sourceblock}{注记 4.2.12 设 Sm (n) = 1m + 2m + 3m + · · · + nm , 德国数学家 Johann Faulhaber}
{\small\sloppy
\noindent 注记 4.2.12 设 Sm (n) = 1m + 2m + 3m + · · · + nm , 德国数学家 Johann Faulhaber\par
\noindent 发现并证明了:\par
\noindent (1) 当 m 是奇数时, Sm (n) 为 Tn 的多项式. 例如\par
\noindent 1 + 2 + 3 + · · · + n = Tn ,\par
\noindent 13 + 23 + 33 + · · · + n3 = Tn2 ,\par
\noindent 4Tn3 − Tn2\par
\noindent 15 + 25 + 35 + · · · + n5 = .\par
\noindent (2) 当 m 是偶数时, Sm (n) 是 Tn 的多项式乘以 2n + 1. 例如\par
\noindent Tn\par
\noindent 12 + 22 + 32 + · · · + n2 = (2n + 1) ,\par
\noindent 6Tn2 − Tn\par
\noindent 14 + 24 + 34 + · · · + n4 = (2n + 1) .\par
\noindent 有兴趣的读者可以参考 [24, 28].\par
\noindent 下面我们给出关于第二类 Stirling 数一个未解决的猜想.\par
\noindent 猜想 4.2.13 (Wilf’s Conjecture) 设 S(n, k) 是第二类 Stirling 数. 对于 n > 2, 当\par
\noindent 0 ≤ k ≤ n 时, S(n, k) 的 n 项交错和不为 0, 也即:\par
\noindent X\par
\noindent n\par
}
\end{sourceblock}


\begin{explainblock}
注记通常是在补充术语、记号、历史或一个容易忽略的约定。它未必给新公式，但常常决定你后面会不会把符号读错。

\smallskip
\textbf{初学者检查点：}
\begin{itemize}[leftmargin=2em,itemsep=0.25em]
\item 记下新符号的读法。
\item 确认这个约定是否只在本章使用。
\item 把注记里的限制条件补到前面的定义旁边。
\end{itemize}

\smallskip
\textbf{本章读法提醒：} 第四章区分两类 Stirling 数：第一类对应排列的圈，第二类对应集合划分。
\end{explainblock}


\begin{sourceblock}{第 4 章续读}
{\small\sloppy
\noindent (−1)k S(n, k) 6= 0. (4.16)\par
\noindent k=0\par
}
\end{sourceblock}


\begin{explainblock}
这一段是原书论述链的一部分。读的时候把它当作桥梁：它通常负责把上一个定义、例子或定理，引到下一个计数方法。

\smallskip
\textbf{初学者检查点：}
\begin{itemize}[leftmargin=2em,itemsep=0.25em]
\item 找出这一段承接的上一个概念。
\item 找出它准备引出的下一个概念。
\item 把中间的关键等式或构造单独抄出来。
\end{itemize}

\smallskip
\textbf{本章读法提醒：} 第四章区分两类 Stirling 数：第一类对应排列的圈，第二类对应集合划分。
\end{explainblock}


\begin{sourceblock}{§ 4.3 差分算子}
{\small\sloppy
\noindent § 4.3. 差分算子\par
}
\end{sourceblock}


\begin{explainblock}
这一节先不要急着背公式。先问三个问题：对象是什么，哪些限制条件不能丢，最后到底是在数所有对象、还是在数某个统计量的分布。把这三件事说清楚，后面的公式才会有位置。

\smallskip
\textbf{初学者检查点：}
\begin{itemize}[leftmargin=2em,itemsep=0.25em]
\item 把本节出现的新对象写成一句话定义。
\item 把每个公式左边和右边分别翻译成“在数什么”。
\item 找一个最小非平凡例子，手工列出所有对象。
\end{itemize}

\smallskip
\textbf{本章读法提醒：} 第四章区分两类 Stirling 数：第一类对应排列的圈，第二类对应集合划分。
\end{explainblock}


\begin{sourceblock}{定理 4.2.5 也可用差分算子来证明. 下面我们引入差分算子的概念.}
{\small\sloppy
\noindent 定理 4.2.5 也可用差分算子来证明. 下面我们引入差分算子的概念.\par
}
\end{sourceblock}


\begin{explainblock}
定理段不要只看结论，要看它把哪两个计数方式连在一起。组合学证明最常见的技巧是：左边按一种标准数，右边按另一种标准数，然后说明两边数的是同一批对象。

\smallskip
\textbf{初学者检查点：}
\begin{itemize}[leftmargin=2em,itemsep=0.25em]
\item 先复述定理的假设，不要把适用范围扩大。
\item 找出证明中的基本计数原则：乘法、加法、双射、递推、容斥或生成函数。
\item 检查最后的等式化简是否和前面的对象解释一致。
\end{itemize}

\smallskip
\textbf{本章读法提醒：} 第四章区分两类 Stirling 数：第一类对应排列的圈，第二类对应集合划分。
\end{explainblock}


\begin{sourceblock}{定义 4.3.1 设 f 是定义在非负整数上的函数. 当 n ≥ 0 时, 称}
{\small\sloppy
\noindent 定义 4.3.1 设 f 是定义在非负整数上的函数. 当 n ≥ 0 时, 称\par
\noindent ∆f (n) := f (n + 1) − f (n), (4.17)\par
\noindent 是 f 在 n 处的一阶差分, ∆ 称为差分算子. 对于任意正整数 k(k ≥ 2), 可递归定义\par
\noindent 78 第4章 两类 Stirling 数\par
\noindent f 在 n 处的 k– 阶差分为:\par
\noindent ∆k f (n) = ∆(∆k−1 f (n)). (4.18)\par
\noindent 一般地, 记 ∆0 f (n) = f (n).\par
}
\end{sourceblock}


\begin{explainblock}
定义段的任务是定边界。这里要特别注意参数范围、是否允许重复、是否考虑顺序、对象有没有标签。后面的每一道例题和证明，本质上都在反复使用这些边界。

\smallskip
\textbf{初学者检查点：}
\begin{itemize}[leftmargin=2em,itemsep=0.25em]
\item 圈出被定义的对象名称。
\item 圈出参数范围，例如 n、k 是否要求正整数。
\item 判断它是有序对象、无序对象、带标签对象还是结构对象。
\end{itemize}

\smallskip
\textbf{本章读法提醒：} 第四章区分两类 Stirling 数：第一类对应排列的圈，第二类对应集合划分。
\end{explainblock}


\begin{sourceblock}{例 4.3.1 设 f (n) = n2 , 计算 f (n) 的各阶差分.}
{\small\sloppy
\noindent 例 4.3.1 设 f (n) = n2 , 计算 f (n) 的各阶差分.\par
\noindent 解:\par
\noindent ∆n2 = (n + 1)2 − n2 = 2n + 1,\par
\noindent ∆2 n2 = ∆(2n + 1) = 2,\par
\noindent ∆3 n2 = ∆2 = 0,\par
\noindent 所以 \par
\noindent \par
\noindent 2n + 1, k = 1,\par
\noindent \par
\noindent ∆k n2 = 2, k = 2,\par
\noindent \par
\noindent \par
\noindent \par
\noindent 0, k ≥ 3.\par
}
\end{sourceblock}


\begin{explainblock}
例题是学习这本书最重要的部分。不要只看答案，要把每个限制条件变成一个计数动作：先安排谁、再插入谁、有没有重复、需不需要除以对称性。

\smallskip
\textbf{初学者检查点：}
\begin{itemize}[leftmargin=2em,itemsep=0.25em]
\item 先写出没有限制时的总数。
\item 逐条加入限制条件，观察总数怎样变化。
\item 最后检查有没有把同一种方案重复算了多次。
\end{itemize}

\smallskip
\textbf{本章读法提醒：} 第四章区分两类 Stirling 数：第一类对应排列的圈，第二类对应集合划分。
\end{explainblock}


\begin{sourceblock}{例 4.3.2 求 n4 各阶差分分别在 n 取 0, 1, 2 时的值.}
{\small\sloppy
\noindent 例 4.3.2 求 n4 各阶差分分别在 n 取 0, 1, 2 时的值.\par
\noindent 解: 我们当然可以先求 n4 的差分再计算具体的值, 但是这样的计算比较复杂. 不\par
\noindent 妨利用差分的定义 ∆f (n) = f (n + 1) − f (n) 知 f (n) 的一阶差分在 n 处的取值就是\par
\noindent f (n + 1) 与 f (n) 的差值, 再由差分的递归定义可知只要知道 f (n) 的在任意整数时\par
\noindent 的取值, 则各阶差分在任意整数上的取值也就能被唯一确定. 所以我们通过计算 n4\par
\noindent 在 n = 0, 1, 2, 3, 4, 5, 6 的值并列表如下:\par
\noindent 0 1 16 81 256 625 1296 ···\par
\noindent 1 15 65 175 369 671 ···\par
\noindent 14 50 110 194 302 ···\par
\noindent 因此 n4 各阶差分的在 n = 1, 2, 3 处的取值便可查表得出. 而且从上表显然知 n4\par
\noindent 的 5 阶及以上的差分的值均为 0.\par
}
\end{sourceblock}


\begin{explainblock}
例题是学习这本书最重要的部分。不要只看答案，要把每个限制条件变成一个计数动作：先安排谁、再插入谁、有没有重复、需不需要除以对称性。

\smallskip
\textbf{初学者检查点：}
\begin{itemize}[leftmargin=2em,itemsep=0.25em]
\item 先写出没有限制时的总数。
\item 逐条加入限制条件，观察总数怎样变化。
\item 最后检查有没有把同一种方案重复算了多次。
\end{itemize}

\smallskip
\textbf{本章读法提醒：} 第四章区分两类 Stirling 数：第一类对应排列的圈，第二类对应集合划分。
\end{explainblock}


\begin{sourceblock}{注记 4.3.2 如例 4.3.2 中绘制的表 4.3, 一般地, 对给定的函数 f , 如果把数值}
{\small\sloppy
\noindent 注记 4.3.2 如例 4.3.2 中绘制的表 4.3, 一般地, 对给定的函数 f , 如果把数值\par
\noindent f (0), f (1), f (2), . . . 排列成一行, 再把 ∆f (0), ∆f (1), ∆f (2), . . . 排在下面一行, 这样\par
\noindent 一直作下去, 也即把 ∆k f (0), ∆k f (1), ∆k f (2), . . . 排在第 k + 1 行. 得到的这个数值\par
\noindent 表就称作 f 的差分表. 差分表在数值计算中有着重要意义.\par
\noindent 注意到在前面两个例子中 n2 的大于 2 阶差分与 n4 的大于 5 阶的差分的值均为\par
\noindent 0, 但这并不是偶然的现象. 对于任意多项式的差分, 我们有下面的定理.\par
}
\end{sourceblock}


\begin{explainblock}
注记通常是在补充术语、记号、历史或一个容易忽略的约定。它未必给新公式，但常常决定你后面会不会把符号读错。

\smallskip
\textbf{初学者检查点：}
\begin{itemize}[leftmargin=2em,itemsep=0.25em]
\item 记下新符号的读法。
\item 确认这个约定是否只在本章使用。
\item 把注记里的限制条件补到前面的定义旁边。
\end{itemize}

\smallskip
\textbf{本章读法提醒：} 第四章区分两类 Stirling 数：第一类对应排列的圈，第二类对应集合划分。
\end{explainblock}


\begin{sourceblock}{定理 4.3.3 设 f (x) 是 x 的 m(m ≥ 1) 次多项式:}
{\small\sloppy
\noindent 定理 4.3.3 设 f (x) 是 x 的 m(m ≥ 1) 次多项式:\par
\noindent f (x) = am xm + am−1 xm−1 + · · · + a0 , am 6= 0.\par
\noindent 若 n 是大于 m 的正整数, 则对任意 x,\par
\noindent ∆n f (x) = 0.\par
\noindent 证明: 事实上, 只需证明当 n = m + 1 命题成立即可. 对 m 用数学归纳法来证\par
\noindent 明. 当 m = 0 时, f (x) = a0 , 则 ∆f (x) = f (x + 1) − f (x) = 0, 命题成立. 假设命题\par
\noindent 对于 x 的 m − 1 次的多项式都成立. 显然地, x 的 m 次的多项式 f (x) 的一阶差分\par
\noindent 是:\par
\noindent ∆f (x) = f (x + 1) − f (x) = (am (x + 1)m + am−1 (x + 1)m−1 + · · · + a1 (x + 1) + a0 )\par
\noindent − (am xm + am−1 xm−1 + · · · + a1 x + a0 ).\par
\noindent 由二项式定理可知,\par
\noindent m m−1\par
\noindent am (x + 1) − am x\par
\noindent m m\par
\noindent = am x + m\par
\noindent x + · · · + 1 − am xm\par
\noindent m m−1\par
\noindent = am x + · · · + am .\par
}
\end{sourceblock}


\begin{explainblock}
定理段不要只看结论，要看它把哪两个计数方式连在一起。组合学证明最常见的技巧是：左边按一种标准数，右边按另一种标准数，然后说明两边数的是同一批对象。

\smallskip
\textbf{初学者检查点：}
\begin{itemize}[leftmargin=2em,itemsep=0.25em]
\item 先复述定理的假设，不要把适用范围扩大。
\item 找出证明中的基本计数原则：乘法、加法、双射、递推、容斥或生成函数。
\item 检查最后的等式化简是否和前面的对象解释一致。
\end{itemize}

\smallskip
\textbf{本章读法提醒：} 第四章区分两类 Stirling 数：第一类对应排列的圈，第二类对应集合划分。
\end{explainblock}


\begin{sourceblock}{第 4 章续读}
{\small\sloppy
\noindent 这说明, ∆f (x) 至多是 x 的 m − 1 次多项式. 由假设可知对任意的实数 x,\par
\noindent ∆m+1 f (x) = ∆m (∆f (x)) = 0.\par
}
\end{sourceblock}


\begin{explainblock}
这一段是原书论述链的一部分。读的时候把它当作桥梁：它通常负责把上一个定义、例子或定理，引到下一个计数方法。

\smallskip
\textbf{初学者检查点：}
\begin{itemize}[leftmargin=2em,itemsep=0.25em]
\item 找出这一段承接的上一个概念。
\item 找出它准备引出的下一个概念。
\item 把中间的关键等式或构造单独抄出来。
\end{itemize}

\smallskip
\textbf{本章读法提醒：} 第四章区分两类 Stirling 数：第一类对应排列的圈，第二类对应集合划分。
\end{explainblock}


\begin{sourceblock}{例 4.3.3 设 n 为正整数, 计算 (x)n 关于 x 的各阶差分.}
{\small\sloppy
\noindent 例 4.3.3 设 n 为正整数, 计算 (x)n 关于 x 的各阶差分.\par
\noindent 解: 因为 (x)n 是 x 的 n 次多项式, 所以当 k > n 时,\par
\noindent ∆k (x)n = 0.\par
\noindent 又\par
\noindent ∆(x)n = (x + 1)n − (x)n = n(x)n−1 ,\par
\noindent 所以\par
\noindent ∆2 (x)n = n∆(x)n−1 = n(n − 1)(x)n−2 ,\par
\noindent 故对任意 1 ≤ k ≤ n, 可得\par
\noindent ∆k (x)n = (n)k (x)n−k .\par
\noindent 特别地, 故当 k = n 时,\par
\noindent ∆n (x)n = n!.\par
\noindent 80 第4章 两类 Stirling 数\par
\noindent 我们引进 (向后) 移位算子 E : Ef (n) = f (n + 1), 与恒等算子 I : If (n) = f (n),\par
\noindent 显然 EI = IE 可以交换顺序, 且有 ∆ = E − I.\par
}
\end{sourceblock}


\begin{explainblock}
例题是学习这本书最重要的部分。不要只看答案，要把每个限制条件变成一个计数动作：先安排谁、再插入谁、有没有重复、需不需要除以对称性。

\smallskip
\textbf{初学者检查点：}
\begin{itemize}[leftmargin=2em,itemsep=0.25em]
\item 先写出没有限制时的总数。
\item 逐条加入限制条件，观察总数怎样变化。
\item 最后检查有没有把同一种方案重复算了多次。
\end{itemize}

\smallskip
\textbf{本章读法提醒：} 第四章区分两类 Stirling 数：第一类对应排列的圈，第二类对应集合划分。
\end{explainblock}


\begin{sourceblock}{定理 4.3.4 设 k 为任意非负整数, f (n) 是定义在非负整数上的函数, 则}
{\small\sloppy
\noindent 定理 4.3.4 设 k 为任意非负整数, f (n) 是定义在非负整数上的函数, 则\par
\noindent X\par
\noindent k\par
\noindent k\par
\noindent ∆k f (n) = (−1)k−i f (n + i). (4.19)\par
\noindent i\par
\noindent i=0\par
\noindent 证明: 因为 E 与 I 可交换, 应用二项式定理得:\par
\noindent X\par
\noindent k\par
\noindent k−i k\par
\noindent ∆ = (E − I) =\par
\noindent k\par
\noindent (−1) k\par
\noindent Ei,\par
\noindent i\par
\noindent i=0\par
\noindent 故\par
}
\end{sourceblock}


\begin{explainblock}
定理段不要只看结论，要看它把哪两个计数方式连在一起。组合学证明最常见的技巧是：左边按一种标准数，右边按另一种标准数，然后说明两边数的是同一批对象。

\smallskip
\textbf{初学者检查点：}
\begin{itemize}[leftmargin=2em,itemsep=0.25em]
\item 先复述定理的假设，不要把适用范围扩大。
\item 找出证明中的基本计数原则：乘法、加法、双射、递推、容斥或生成函数。
\item 检查最后的等式化简是否和前面的对象解释一致。
\end{itemize}

\smallskip
\textbf{本章读法提醒：} 第四章区分两类 Stirling 数：第一类对应排列的圈，第二类对应集合划分。
\end{explainblock}


\begin{sourceblock}{第 4 章续读}
{\small\sloppy
\noindent X\par
\noindent k X\par
\noindent k\par
\noindent k k−i k i k−i k\par
\noindent ∆ f (n) = (−1) E f (n) = (−1) f (n + i).\par
\noindent i i\par
\noindent i=0 i=0\par
}
\end{sourceblock}


\begin{explainblock}
这一段是原书论述链的一部分。读的时候把它当作桥梁：它通常负责把上一个定义、例子或定理，引到下一个计数方法。

\smallskip
\textbf{初学者检查点：}
\begin{itemize}[leftmargin=2em,itemsep=0.25em]
\item 找出这一段承接的上一个概念。
\item 找出它准备引出的下一个概念。
\item 把中间的关键等式或构造单独抄出来。
\end{itemize}

\smallskip
\textbf{本章读法提醒：} 第四章区分两类 Stirling 数：第一类对应排列的圈，第二类对应集合划分。
\end{explainblock}


\begin{sourceblock}{注记 4.3.5 在定理 4.3.4 中取 n = 0, 则}
{\small\sloppy
\noindent 注记 4.3.5 在定理 4.3.4 中取 n = 0, 则\par
\noindent X\par
\noindent k\par
\noindent k k−i k\par
\noindent ∆ f (0) = (−1) f (i), (4.20)\par
\noindent i\par
\noindent i=0\par
\noindent 这表示 f (n) 的 k 阶差分在 n = 0 处的取值可由 f (0), f (1), · · · , f (k) 的值表示.\par
}
\end{sourceblock}


\begin{explainblock}
注记通常是在补充术语、记号、历史或一个容易忽略的约定。它未必给新公式，但常常决定你后面会不会把符号读错。

\smallskip
\textbf{初学者检查点：}
\begin{itemize}[leftmargin=2em,itemsep=0.25em]
\item 记下新符号的读法。
\item 确认这个约定是否只在本章使用。
\item 把注记里的限制条件补到前面的定义旁边。
\end{itemize}

\smallskip
\textbf{本章读法提醒：} 第四章区分两类 Stirling 数：第一类对应排列的圈，第二类对应集合划分。
\end{explainblock}


\begin{sourceblock}{例 4.3.4 利用 (4.20) 计算 n4 的各阶差分在 n = 0 的取值.}
{\small\sloppy
\noindent 例 4.3.4 利用 (4.20) 计算 n4 的各阶差分在 n = 0 的取值.\par
\noindent 解: 由定理 4.3.4 知\par
\noindent X\par
\noindent k\par
\noindent k 4 k−i k 4\par
\noindent ∆ 0 = (−1) i ,\par
\noindent i\par
\noindent i=0\par
\noindent 所以\par
\noindent ∆04 = 1, ∆2 04 = −2 · 14 + 24 = 14,\par
\noindent ∆3 04 = 3 · 14 − 3 · 24 + 34 = 36, ∆4 = −4 · 14 + 6 · 24 − 4 · 34 + 44 = 24,\par
\noindent 再由 4.3.3 知 k ≥ 5 时,\par
\noindent X\par
\noindent k\par
\noindent k 4\par
\noindent ∆k 04 = (−1)k−i i = 0.\par
\noindent i\par
\noindent i=0\par
}
\end{sourceblock}


\begin{explainblock}
例题是学习这本书最重要的部分。不要只看答案，要把每个限制条件变成一个计数动作：先安排谁、再插入谁、有没有重复、需不需要除以对称性。

\smallskip
\textbf{初学者检查点：}
\begin{itemize}[leftmargin=2em,itemsep=0.25em]
\item 先写出没有限制时的总数。
\item 逐条加入限制条件，观察总数怎样变化。
\item 最后检查有没有把同一种方案重复算了多次。
\end{itemize}

\smallskip
\textbf{本章读法提醒：} 第四章区分两类 Stirling 数：第一类对应排列的圈，第二类对应集合划分。
\end{explainblock}


\begin{sourceblock}{例 4.3.5 用差分算子证明:}
{\small\sloppy
\noindent 例 4.3.5 用差分算子证明:\par
\noindent X\par
\noindent k\par
\noindent k−j k n\par
\noindent k!S(n, k) = (−1) j .\par
\noindent j\par
\noindent j=0\par
\noindent 证明: 由定理 4.2.8 知\par
\noindent X\par
\noindent n\par
\noindent S(n, j)(x)j = xn .\par
\noindent j=0\par
\noindent 对上式两边关于 x 作 k(0 ≤ k ≤ n) 次差分, 由定理 4.3.3可得\par
\noindent X\par
\noindent n\par
\noindent k!S(n, k) + S(n, j)∆k (x)j = ∆k xn . (4.21)\par
\noindent j=k+1\par
\noindent 因为\par
}
\end{sourceblock}


\begin{explainblock}
例题是学习这本书最重要的部分。不要只看答案，要把每个限制条件变成一个计数动作：先安排谁、再插入谁、有没有重复、需不需要除以对称性。

\smallskip
\textbf{初学者检查点：}
\begin{itemize}[leftmargin=2em,itemsep=0.25em]
\item 先写出没有限制时的总数。
\item 逐条加入限制条件，观察总数怎样变化。
\item 最后检查有没有把同一种方案重复算了多次。
\end{itemize}

\smallskip
\textbf{本章读法提醒：} 第四章区分两类 Stirling 数：第一类对应排列的圈，第二类对应集合划分。
\end{explainblock}


\begin{sourceblock}{第 4 章续读}
{\small\sloppy
\noindent X\par
\noindent n X\par
\noindent n\par
\noindent S(n, j)∆k (x)j = S(n, j)(j)k (x)j−k ,\par
\noindent j=k+1 j=k+1\par
\noindent 是不含常数项的多项式, 所以在 (4.21) 中令 x = 0 得\par
\noindent k!S(n, k) = ∆k 0n .\par
\noindent 在 (4.20) 中取 f (n) = nm 立得\par
\noindent X\par
\noindent k\par
\noindent k m k−i k\par
\noindent ∆ 0 = (−1) im .\par
\noindent i\par
\noindent i=0\par
\noindent 所以\par
\noindent X\par
\noindent k\par
\noindent k n k−i k\par
}
\end{sourceblock}


\begin{explainblock}
这一段是原书论述链的一部分。读的时候把它当作桥梁：它通常负责把上一个定义、例子或定理，引到下一个计数方法。

\smallskip
\textbf{初学者检查点：}
\begin{itemize}[leftmargin=2em,itemsep=0.25em]
\item 找出这一段承接的上一个概念。
\item 找出它准备引出的下一个概念。
\item 把中间的关键等式或构造单独抄出来。
\end{itemize}

\smallskip
\textbf{本章读法提醒：} 第四章区分两类 Stirling 数：第一类对应排列的圈，第二类对应集合划分。
\end{explainblock}


\begin{sourceblock}{第 4 章续读}
{\small\sloppy
\noindent k!S(n, k) = ∆ 0 = (−1) in .\par
\noindent i\par
\noindent i=0\par
\noindent 在前面定理 4.3.4 中我们知任意定义在非负整数上的函数 f (n), 它的任意 k 阶差\par
\noindent 分在 0 处的取值都可由 f (n) 前 k 项的函数值唯一确定. 当然反过来, f (n) 在 k 处\par
\noindent 的取值也由其各阶差分在 0 处的取值确定.\par
}
\end{sourceblock}


\begin{explainblock}
这一段是原书论述链的一部分。读的时候把它当作桥梁：它通常负责把上一个定义、例子或定理，引到下一个计数方法。

\smallskip
\textbf{初学者检查点：}
\begin{itemize}[leftmargin=2em,itemsep=0.25em]
\item 找出这一段承接的上一个概念。
\item 找出它准备引出的下一个概念。
\item 把中间的关键等式或构造单独抄出来。
\end{itemize}

\smallskip
\textbf{本章读法提醒：} 第四章区分两类 Stirling 数：第一类对应排列的圈，第二类对应集合划分。
\end{explainblock}


\begin{sourceblock}{定理 4.3.6 设 f (n) 是定义在非负整数上的函数, 则}
{\small\sloppy
\noindent 定理 4.3.6 设 f (n) 是定义在非负整数上的函数, 则\par
\noindent n\par
\noindent X n\par
\noindent f (n) = ∆k f (0).\par
\noindent k\par
\noindent k=0\par
\noindent 证明: 因为 f (n) = E n f (0), 再由 E = I + ∆ 且 I 与 ∆ 可交换, 应用二项式定\par
\noindent 理可得\par
\noindent f (n) = E n f (0) = (I + ∆)n f (0)\par
\noindent n\par
\noindent X n\par
\noindent = ∆k f (0).\par
\noindent k\par
\noindent k=0\par
\noindent 82 第4章 两类 Stirling 数\par
}
\end{sourceblock}


\begin{explainblock}
定理段不要只看结论，要看它把哪两个计数方式连在一起。组合学证明最常见的技巧是：左边按一种标准数，右边按另一种标准数，然后说明两边数的是同一批对象。

\smallskip
\textbf{初学者检查点：}
\begin{itemize}[leftmargin=2em,itemsep=0.25em]
\item 先复述定理的假设，不要把适用范围扩大。
\item 找出证明中的基本计数原则：乘法、加法、双射、递推、容斥或生成函数。
\item 检查最后的等式化简是否和前面的对象解释一致。
\end{itemize}

\smallskip
\textbf{本章读法提醒：} 第四章区分两类 Stirling 数：第一类对应排列的圈，第二类对应集合划分。
\end{explainblock}


\begin{sourceblock}{例 4.3.6 设 f (x) = x3 + 3x2 − 2x + 1 是一个实多项式, 证明:}
{\small\sloppy
\noindent 例 4.3.6 设 f (x) = x3 + 3x2 − 2x + 1 是一个实多项式, 证明:\par
\noindent x x x x\par
\noindent f (x) = +2 + 12 +6 .\par
\noindent 证明: 不妨设 x 是非负整数, 计算 f (x) 的各阶差分在 x = 0 处的取值得:\par
\noindent ∆0 f (0) = 1, ∆f (0) = 2, ∆2 f (0) = 12, ∆3 f (0) = 6,\par
\noindent 当 k ≥ 4 时, ∆k f (0) = 0, 则由定理 4.3.6 得\par
\noindent x x x x\par
\noindent f (x) = +2 + 12 +6 ,\par
\noindent 即上式对任意非负整数均成立, 所以对任意实数 x 也均成立.\par
}
\end{sourceblock}


\begin{explainblock}
例题是学习这本书最重要的部分。不要只看答案，要把每个限制条件变成一个计数动作：先安排谁、再插入谁、有没有重复、需不需要除以对称性。

\smallskip
\textbf{初学者检查点：}
\begin{itemize}[leftmargin=2em,itemsep=0.25em]
\item 先写出没有限制时的总数。
\item 逐条加入限制条件，观察总数怎样变化。
\item 最后检查有没有把同一种方案重复算了多次。
\end{itemize}

\smallskip
\textbf{本章读法提醒：} 第四章区分两类 Stirling 数：第一类对应排列的圈，第二类对应集合划分。
\end{explainblock}


\begin{sourceblock}{例 4.3.7 计算前 n 个自然数的 4 次幂的和.}
{\small\sloppy
\noindent 例 4.3.7 计算前 n 个自然数的 4 次幂的和.\par
\noindent 解: 由例 4.3.2 知 n4 的各阶差分在 0 处的取值分别为\par
\noindent 0, 1, 14, 36, 24, 0, 0, . . . ,\par
\noindent 所以由定理 4.3.6 知\par
\noindent n n n n n\par
\noindent n =0·\par
\noindent +1· + 14 · + 36 · + 24 · ,\par
\noindent 故\par
\noindent X\par
\noindent n n\par
\noindent X n\par
\noindent X n\par
\noindent X n\par
\noindent X\par
\noindent k = + 14 + 36 + 24\par
\noindent k=1 k=1 k=1 k=1 k=1\par
\noindent n+1 n+1 n+1 n+1\par
\noindent = + 14 + 36 + 24\par
}
\end{sourceblock}


\begin{explainblock}
例题是学习这本书最重要的部分。不要只看答案，要把每个限制条件变成一个计数动作：先安排谁、再插入谁、有没有重复、需不需要除以对称性。

\smallskip
\textbf{初学者检查点：}
\begin{itemize}[leftmargin=2em,itemsep=0.25em]
\item 先写出没有限制时的总数。
\item 逐条加入限制条件，观察总数怎样变化。
\item 最后检查有没有把同一种方案重复算了多次。
\end{itemize}

\smallskip
\textbf{本章读法提醒：} 第四章区分两类 Stirling 数：第一类对应排列的圈，第二类对应集合划分。
\end{explainblock}


\begin{sourceblock}{第 4 章续读}
{\small\sloppy
\noindent (6Tn2 − Tn )(2n + 1)\par
\noindent = ,\par
\noindent 其中 Tn = n(n+1)\par
\noindent 2 . 这与例 4.2.2 中得到的结果一样.\par
}
\end{sourceblock}


\begin{explainblock}
这一段是原书论述链的一部分。读的时候把它当作桥梁：它通常负责把上一个定义、例子或定理，引到下一个计数方法。

\smallskip
\textbf{初学者检查点：}
\begin{itemize}[leftmargin=2em,itemsep=0.25em]
\item 找出这一段承接的上一个概念。
\item 找出它准备引出的下一个概念。
\item 把中间的关键等式或构造单独抄出来。
\end{itemize}

\smallskip
\textbf{本章读法提醒：} 第四章区分两类 Stirling 数：第一类对应排列的圈，第二类对应集合划分。
\end{explainblock}


\begin{sourceblock}{§ 4.4 文法}
{\small\sloppy
\noindent § 4.4. 文法\par
\noindent 形式文法定义了一种由序列和符号组成的语言. 形式文法的基本思想是从一些初\par
\noindent 始元素开始, 通过规则产生新的元素.\par
}
\end{sourceblock}


\begin{explainblock}
这一节先不要急着背公式。先问三个问题：对象是什么，哪些限制条件不能丢，最后到底是在数所有对象、还是在数某个统计量的分布。把这三件事说清楚，后面的公式才会有位置。

\smallskip
\textbf{初学者检查点：}
\begin{itemize}[leftmargin=2em,itemsep=0.25em]
\item 把本节出现的新对象写成一句话定义。
\item 把每个公式左边和右边分别翻译成“在数什么”。
\item 找一个最小非平凡例子，手工列出所有对象。
\end{itemize}

\smallskip
\textbf{本章读法提醒：} 第四章区分两类 Stirling 数：第一类对应排列的圈，第二类对应集合划分。
\end{explainblock}


\begin{sourceblock}{定义 4.4.1 (上下文无关文法) 假设 A 是一个字母表, 其中字母均为互相独立且可}
{\small\sloppy
\noindent 定义 4.4.1 (上下文无关文法) 假设 A 是一个字母表, 其中字母均为互相独立且可\par
\noindent 交换的不定元, A 上的形式函数 (formal function) 可定义如下:\par
\noindent (a) A 中每个字母都是一个形式函数;\par
\noindent (b) 如果 u 和 v 是形式函数, 那么 u + v, uv 均为形式函数;\par
\noindent (c) 若 f 是解析函数, u 为形式函数, 那么 f (u) 为形式函数;\par
\noindent (d) 每一个形式函数都可以由有限个上述过程得到.\par
}
\end{sourceblock}


\begin{explainblock}
定义段的任务是定边界。这里要特别注意参数范围、是否允许重复、是否考虑顺序、对象有没有标签。后面的每一道例题和证明，本质上都在反复使用这些边界。

\smallskip
\textbf{初学者检查点：}
\begin{itemize}[leftmargin=2em,itemsep=0.25em]
\item 圈出被定义的对象名称。
\item 圈出参数范围，例如 n、k 是否要求正整数。
\item 判断它是有序对象、无序对象、带标签对象还是结构对象。
\end{itemize}

\smallskip
\textbf{本章读法提醒：} 第四章区分两类 Stirling 数：第一类对应排列的圈，第二类对应集合划分。
\end{explainblock}


\begin{sourceblock}{定义在字母集 A 上的上下文无关文法 G, 就是将 A 中字母替换为 A 中形式函数的}
{\small\sloppy
\noindent 定义在字母集 A 上的上下文无关文法 G, 就是将 A 中字母替换为 A 中形式函数的\par
\noindent 替换规则 (substitution rules) 的集合.\par
}
\end{sourceblock}


\begin{explainblock}
定义段的任务是定边界。这里要特别注意参数范围、是否允许重复、是否考虑顺序、对象有没有标签。后面的每一道例题和证明，本质上都在反复使用这些边界。

\smallskip
\textbf{初学者检查点：}
\begin{itemize}[leftmargin=2em,itemsep=0.25em]
\item 圈出被定义的对象名称。
\item 圈出参数范围，例如 n、k 是否要求正整数。
\item 判断它是有序对象、无序对象、带标签对象还是结构对象。
\end{itemize}

\smallskip
\textbf{本章读法提醒：} 第四章区分两类 Stirling 数：第一类对应排列的圈，第二类对应集合划分。
\end{explainblock}


\begin{sourceblock}{定义 4.4.2 如下定义的微分算子 D 是良性的：}
{\small\sloppy
\noindent 定义 4.4.2 如下定义的微分算子 D 是良性的：\par
\noindent (a) 对于形式函数 u 和 v,\par
\noindent D(u + v) = D(u) + D(v),\par
\noindent D(uv) = D(u)v + uD(v).\par
\noindent (b) 对解析函数 f (x) 和形式函数 u,\par
\noindent ∂f (u)\par
\noindent D(f (u)) = D(u).\par
\noindent ∂u\par
\noindent (c) 对于任意 u ∈ A, 若文法 G 中包含替换规则 u → v, 则 D(u) = v; 否则,\par
\noindent D(u) = 0, 称这样的 u 为终点 (terminal).\par
\noindent 我们称满足如上条件的 D 为文法 G 对应的形式微分.\par
}
\end{sourceblock}


\begin{explainblock}
定义段的任务是定边界。这里要特别注意参数范围、是否允许重复、是否考虑顺序、对象有没有标签。后面的每一道例题和证明，本质上都在反复使用这些边界。

\smallskip
\textbf{初学者检查点：}
\begin{itemize}[leftmargin=2em,itemsep=0.25em]
\item 圈出被定义的对象名称。
\item 圈出参数范围，例如 n、k 是否要求正整数。
\item 判断它是有序对象、无序对象、带标签对象还是结构对象。
\end{itemize}

\smallskip
\textbf{本章读法提醒：} 第四章区分两类 Stirling 数：第一类对应排列的圈，第二类对应集合划分。
\end{explainblock}


\begin{sourceblock}{例如, 若}
{\small\sloppy
\noindent 例如, 若\par
\noindent G = \{x → y, y → xy\},\par
\noindent 那么 D(x) = y, D(y) = xy, D(z) = 0.\par
\noindent 值得注意的是, Leibniz 公式对于微分算子依然成立：\par
\noindent X n\par
\noindent n\par
\noindent D (xy) = Dk (x)Dn−k (y).\par
\noindent k\par
\noindent k=0\par
}
\end{sourceblock}


\begin{explainblock}
例题是学习这本书最重要的部分。不要只看答案，要把每个限制条件变成一个计数动作：先安排谁、再插入谁、有没有重复、需不需要除以对称性。

\smallskip
\textbf{初学者检查点：}
\begin{itemize}[leftmargin=2em,itemsep=0.25em]
\item 先写出没有限制时的总数。
\item 逐条加入限制条件，观察总数怎样变化。
\item 最后检查有没有把同一种方案重复算了多次。
\end{itemize}

\smallskip
\textbf{本章读法提醒：} 第四章区分两类 Stirling 数：第一类对应排列的圈，第二类对应集合划分。
\end{explainblock}


\begin{sourceblock}{例 4.4.1 设序列 \{an \}n ≥ 0, \{bn \}n≥0 满足:}
{\small\sloppy
\noindent 例 4.4.1 设序列 \{an \}n ≥ 0, \{bn \}n≥0 满足:\par
\noindent n\par
\noindent X n\par
\noindent an = bk .\par
\noindent k\par
\noindent k=0\par
\noindent 用文法证明:\par
\noindent n\par
\noindent X n\par
\noindent bn = (−1)n−k ak .\par
\noindent k\par
\noindent k=0\par
\noindent 证明: 设\par
\noindent G = \{f → f, ci → ci+1 \},\par
\noindent 其中 i = 0, 1, . . . 且 ci f = bi . 设 D 是文法 G 对应的形式微分, 则 bn = f Dn (c0 ). 由\par
\noindent Leibniz 公式得:\par
\noindent bn = f · Dn (c0 )\par
\noindent 84 第4章 两类 Stirling 数\par
}
\end{sourceblock}


\begin{explainblock}
例题是学习这本书最重要的部分。不要只看答案，要把每个限制条件变成一个计数动作：先安排谁、再插入谁、有没有重复、需不需要除以对称性。

\smallskip
\textbf{初学者检查点：}
\begin{itemize}[leftmargin=2em,itemsep=0.25em]
\item 先写出没有限制时的总数。
\item 逐条加入限制条件，观察总数怎样变化。
\item 最后检查有没有把同一种方案重复算了多次。
\end{itemize}

\smallskip
\textbf{本章读法提醒：} 第四章区分两类 Stirling 数：第一类对应排列的圈，第二类对应集合划分。
\end{explainblock}


\begin{sourceblock}{第 4 章续读}
{\small\sloppy
\noindent = f · Dn (c0 f f −1 )\par
\noindent X n\par
\noindent n\par
\noindent =f Dk (c0 f )Dn−k (f −1 ).\par
\noindent k\par
\noindent k=0\par
\noindent 因为\par
\noindent D(c0 f ) = c0 f + f c1 ;\par
\noindent D2 (c0 f ) = c0 f + 2f c1 + f c2 ;\par
\noindent ..\par
\noindent .\par
\noindent Xn\par
\noindent n\par
\noindent Dn (c0 f ) = ck f\par
\noindent k\par
\noindent k=0\par
\noindent 所以\par
\noindent n\par
}
\end{sourceblock}


\begin{explainblock}
这一段是原书论述链的一部分。读的时候把它当作桥梁：它通常负责把上一个定义、例子或定理，引到下一个计数方法。

\smallskip
\textbf{初学者检查点：}
\begin{itemize}[leftmargin=2em,itemsep=0.25em]
\item 找出这一段承接的上一个概念。
\item 找出它准备引出的下一个概念。
\item 把中间的关键等式或构造单独抄出来。
\end{itemize}

\smallskip
\textbf{本章读法提醒：} 第四章区分两类 Stirling 数：第一类对应排列的圈，第二类对应集合划分。
\end{explainblock}


\begin{sourceblock}{第 4 章续读}
{\small\sloppy
\noindent X\par
\noindent n n\par
\noindent D (b0 ) = bk = an .\par
\noindent k\par
\noindent k=0\par
\noindent 又因为\par
\noindent D(f −1 ) = − D(f ) = −f −1 ,\par
\noindent f2\par
\noindent 所以\par
\noindent Dn (f −1 ) = (−1)n f −1 .\par
\noindent 从而\par
\noindent Xn\par
\noindent n\par
\noindent bn = f Dk (c0 f )Dn−k (f −1 ) (4.22)\par
\noindent k\par
\noindent k=0\par
\noindent Xn\par
\noindent n\par
}
\end{sourceblock}


\begin{explainblock}
这一段是原书论述链的一部分。读的时候把它当作桥梁：它通常负责把上一个定义、例子或定理，引到下一个计数方法。

\smallskip
\textbf{初学者检查点：}
\begin{itemize}[leftmargin=2em,itemsep=0.25em]
\item 找出这一段承接的上一个概念。
\item 找出它准备引出的下一个概念。
\item 把中间的关键等式或构造单独抄出来。
\end{itemize}

\smallskip
\textbf{本章读法提醒：} 第四章区分两类 Stirling 数：第一类对应排列的圈，第二类对应集合划分。
\end{explainblock}


\begin{sourceblock}{第 4 章续读}
{\small\sloppy
\noindent =f ak · (−1)k f −1 (4.23)\par
\noindent k\par
\noindent k=0\par
\noindent Xn\par
\noindent k n\par
\noindent = (−1) ak . (4.24)\par
\noindent k\par
\noindent k=0\par
}
\end{sourceblock}


\begin{explainblock}
这一段是原书论述链的一部分。读的时候把它当作桥梁：它通常负责把上一个定义、例子或定理，引到下一个计数方法。

\smallskip
\textbf{初学者检查点：}
\begin{itemize}[leftmargin=2em,itemsep=0.25em]
\item 找出这一段承接的上一个概念。
\item 找出它准备引出的下一个概念。
\item 把中间的关键等式或构造单独抄出来。
\end{itemize}

\smallskip
\textbf{本章读法提醒：} 第四章区分两类 Stirling 数：第一类对应排列的圈，第二类对应集合划分。
\end{explainblock}


\begin{sourceblock}{定理 4.4.3 设 G = \{x → xy, , y → yz, z → z 2 \}, D 为 G 对应的形式微分, 则}
{\small\sloppy
\noindent 定理 4.4.3 设 G = \{x → xy, , y → yz, z → z 2 \}, D 为 G 对应的形式微分, 则\par
\noindent Dn (x) = xy(y + z)(y + 2z) · · · (y + (n − 1)z). (4.25)\par
\noindent 设 c(n, k) 为第一类生成函数, 则它还可写作\par
\noindent X\par
\noindent n\par
\noindent Dn (x) = x c(n, k)y k z n−k . (4.26)\par
\noindent k=0\par
\noindent 证明: 我们不妨先进行几次计算来寻找规律,\par
\noindent D(x) = xy;\par
\noindent D2 (x) = xy 2 + xyz = xy(y + z);\par
\noindent D3 (x) = xy 3 + 3xy 2 z + 2xyz 2 = xy(y + z)(y + 2z).\par
\noindent 根据以上的计算, 不妨假设\par
\noindent Dn+1 = (y + nz)Dn (x).\par
\noindent 并用归纳法进行证明.\par
\noindent 假设对 n 情形, 满足\par
\noindent Dn = (y + (n − 1)z)Dn−1 (x).\par
\noindent 那么,\par
\noindent Dn+1 (x) = D(Dn (x))\par
}
\end{sourceblock}


\begin{explainblock}
定理段不要只看结论，要看它把哪两个计数方式连在一起。组合学证明最常见的技巧是：左边按一种标准数，右边按另一种标准数，然后说明两边数的是同一批对象。

\smallskip
\textbf{初学者检查点：}
\begin{itemize}[leftmargin=2em,itemsep=0.25em]
\item 先复述定理的假设，不要把适用范围扩大。
\item 找出证明中的基本计数原则：乘法、加法、双射、递推、容斥或生成函数。
\item 检查最后的等式化简是否和前面的对象解释一致。
\end{itemize}

\smallskip
\textbf{本章读法提醒：} 第四章区分两类 Stirling 数：第一类对应排列的圈，第二类对应集合划分。
\end{explainblock}


\begin{sourceblock}{第 4 章续读}
{\small\sloppy
\noindent = D(y + (n − 1)z)Dn−1 (x) + (y + (n − 1)z)Dn (x)\par
\noindent = (yz + (n − 1)z 2 )Dn−1 (x) + (y + (n − 1)z)Dn (x)\par
\noindent = zDn (x) + (y + (n − 1)z)Dn (x)\par
\noindent = (y + nz)Dn (x).\par
\noindent 所以 (4.25) 成立.\par
\noindent 对于 (4.26), 我们依然用归纳法证明.\par
\noindent 当 n = 1 时, (4.26) 显然成立:\par
\noindent xc(1, 0)z + xc(1, 1)y = xy = D(x).\par
\noindent 假设对于 n, (4.26) 成立. 则对于 n + 1,\par
\noindent Dn+1 (x) = D(Dn (x))\par
\noindent X\par
\noindent n\par
\noindent = D(x c(n, k)y k z n−k )\par
\noindent k=0\par
\noindent X\par
\noindent n\par
\noindent =x c(n, k) ky k−1 yz · z n−k + y k (n − k)z n−k−1 z 2\par
\noindent k=0\par
}
\end{sourceblock}


\begin{explainblock}
这一段是原书论述链的一部分。读的时候把它当作桥梁：它通常负责把上一个定义、例子或定理，引到下一个计数方法。

\smallskip
\textbf{初学者检查点：}
\begin{itemize}[leftmargin=2em,itemsep=0.25em]
\item 找出这一段承接的上一个概念。
\item 找出它准备引出的下一个概念。
\item 把中间的关键等式或构造单独抄出来。
\end{itemize}

\smallskip
\textbf{本章读法提醒：} 第四章区分两类 Stirling 数：第一类对应排列的圈，第二类对应集合划分。
\end{explainblock}


\begin{sourceblock}{第 4 章续读}
{\small\sloppy
\noindent X\par
\noindent n\par
\noindent + xy c(n, k)y k z n−k\par
\noindent k=0\par
\noindent X\par
\noindent n X\par
\noindent n\par
\noindent =x c(n, k)y k+1 z n−k + x nc(n, k)y k z n−k+1\par
\noindent k=0 k=0\par
\noindent 86 第4章 两类 Stirling 数\par
\noindent X\par
\noindent n+1 X\par
\noindent n\par
\noindent =x c(n, k − 1)y k z n−k+1 + x nc(n, k)y k z n−k+1\par
\noindent k=1 k=0\par
\noindent X\par
\noindent n\par
\noindent =x (c(n, k − 1) + nc(n, k))y k z n−k+1 + c(n, n)xy n+1\par
}
\end{sourceblock}


\begin{explainblock}
这一段是原书论述链的一部分。读的时候把它当作桥梁：它通常负责把上一个定义、例子或定理，引到下一个计数方法。

\smallskip
\textbf{初学者检查点：}
\begin{itemize}[leftmargin=2em,itemsep=0.25em]
\item 找出这一段承接的上一个概念。
\item 找出它准备引出的下一个概念。
\item 把中间的关键等式或构造单独抄出来。
\end{itemize}

\smallskip
\textbf{本章读法提醒：} 第四章区分两类 Stirling 数：第一类对应排列的圈，第二类对应集合划分。
\end{explainblock}


\begin{sourceblock}{第 4 章续读}
{\small\sloppy
\noindent k=1\par
\noindent X\par
\noindent n+1\par
\noindent =x c(n + 1, k)y k z n−k+1 .\par
\noindent k=0\par
\noindent 即对于 n + 1 情形, (4.26) 也成立.\par
\noindent 综上, 定理得证.\par
\noindent 下面给出这一文法的组合意义. 对于一个用标准圈结构表示的排列 ω ∈ Sn , 我们\par
\noindent 将每个排列标记为一个 x, 将排列中每个圈的首位标记为 y, 其余字母标记为 z, 这\par
\noindent 样我们就可以用文法中的形式函数将排列标记. 例如, ω = (2)(41)(85763) 被标记为\par
\noindent xyyzyzzzz = xy 3 z 5 . 当 n = 1 时, 排列 ω = 1 被标记为 xy = D(x); 当 n = 2 时, 排\par
\noindent 列 ω1 = (1)(2), ω2 = (21) 分别被标记为 xy 2 , xyz. 此时有 xy 2 + xyz = D2 (x).\par
\noindent 因为集合 [n + 1] 上任意一 k 个圈排列, 可以将 n + 1 插入到一个有 k 或 k − 1\par
\noindent 个圈的 [n] 上的排列中得到. 若插入到有 k 个圈 [n] 的排列中得到\par
\noindent 现将 n + 1 插入到集合 [n] 的任一有 k 个圈的排列中, 为得到一个仍然是 k 个圈\par
\noindent 的排列, 则这时将增加一个新的标记 z, 且若 n + 1 插入的圈原本只有一个字母, 则\par
\noindent 标记 y 将变为 yz, 若插入长度大于等于两个字母的圈, 则标记 z 变为 z 2 ; 若将 n + 1\par
\noindent 插入到集合 [n] 的任一有 k − 1 个圈的排列中, 为得到一个是 k 个圈的排列, 则 n + 1\par
}
\end{sourceblock}


\begin{explainblock}
这一段是原书论述链的一部分。读的时候把它当作桥梁：它通常负责把上一个定义、例子或定理，引到下一个计数方法。

\smallskip
\textbf{初学者检查点：}
\begin{itemize}[leftmargin=2em,itemsep=0.25em]
\item 找出这一段承接的上一个概念。
\item 找出它准备引出的下一个概念。
\item 把中间的关键等式或构造单独抄出来。
\end{itemize}

\smallskip
\textbf{本章读法提醒：} 第四章区分两类 Stirling 数：第一类对应排列的圈，第二类对应集合划分。
\end{explainblock}


\begin{sourceblock}{第 4 章续读}
{\small\sloppy
\noindent 只能作为一个新圈, 这时增加一个新标记 y, 可看作将 x 变为 xy. 因此, 用文法中形\par
\noindent 式函数表示: 即为 G = \{x → xy, y → yz, z → z 2 \}.\par
\noindent 由定理 4.4.3 知\par
\noindent X\par
\noindent n\par
\noindent xy(y + z)(y + 2z) · · · (y + (n − 1)z). = x c(n, k)y k z n−k ,\par
\noindent k=0\par
\noindent 取 x = z = 1, 立得定理 4.3.4:\par
\noindent X\par
\noindent n\par
\noindent c(n, k)y k = y(y + 1) · · · (y + n − 1).\par
\noindent k=0\par
}
\end{sourceblock}


\begin{explainblock}
这一段是原书论述链的一部分。读的时候把它当作桥梁：它通常负责把上一个定义、例子或定理，引到下一个计数方法。

\smallskip
\textbf{初学者检查点：}
\begin{itemize}[leftmargin=2em,itemsep=0.25em]
\item 找出这一段承接的上一个概念。
\item 找出它准备引出的下一个概念。
\item 把中间的关键等式或构造单独抄出来。
\end{itemize}

\smallskip
\textbf{本章读法提醒：} 第四章区分两类 Stirling 数：第一类对应排列的圈，第二类对应集合划分。
\end{explainblock}


\begin{sourceblock}{定理 4.4.4 设 G = \{x → xy, y → y\}, D 是文法 G 对应的形式微分. 设 S(n, k) 是}
{\small\sloppy
\noindent 定理 4.4.4 设 G = \{x → xy, y → y\}, D 是文法 G 对应的形式微分. 设 S(n, k) 是\par
\noindent 第二类 Stirling 数, 则\par
\noindent X\par
\noindent n\par
\noindent Dn (x) = S(n, k)xy k . (4.27)\par
\noindent k=0\par
\noindent 证明: 对集合的一个划分进行标号, 这个划分的整体标为 x, 对这个划分的每个\par
\noindent 部分标为 y, 然后把这些标号相乘作为这个划分的权重. 例如: 集合 \{1, 2, 3\} 的一个\par
\noindent 2–部划分 1|23 的权重为 xy 2 .\par
\noindent 为得到集合 \{1, 2, · · · , n\} 的所有 k–部划分, 我们可以在 \{1, 2, · · · , n − 1\} 的每个\par
\noindent (k − 1)–划分中添加一个块 \{n\}; 也可以在 \{1, 2, · · · , n − 1\} 的每个 k–划分的子块内\par
\noindent 加入元素 n 得到 \{1, 2, · · · , n\} 上 k 个不同的 k–部划分.\par
\noindent 对于前者, 对应于文法 G 的 x → xy 替换规则: 对于后者, 对应于文法的 y → y\par
\noindent 替换规则.\par
\noindent 由微分的性质\par
\noindent Xn X\par
\noindent n\par
\noindent n+1 n k\par
}
\end{sourceblock}


\begin{explainblock}
定理段不要只看结论，要看它把哪两个计数方式连在一起。组合学证明最常见的技巧是：左边按一种标准数，右边按另一种标准数，然后说明两边数的是同一批对象。

\smallskip
\textbf{初学者检查点：}
\begin{itemize}[leftmargin=2em,itemsep=0.25em]
\item 先复述定理的假设，不要把适用范围扩大。
\item 找出证明中的基本计数原则：乘法、加法、双射、递推、容斥或生成函数。
\item 检查最后的等式化简是否和前面的对象解释一致。
\end{itemize}

\smallskip
\textbf{本章读法提醒：} 第四章区分两类 Stirling 数：第一类对应排列的圈，第二类对应集合划分。
\end{explainblock}


\begin{sourceblock}{第 4 章续读}
{\small\sloppy
\noindent D (x) = D(D (x)) = D( S(n, k)xy ) = S(n, k)(xy k+1 + kxy k ),\par
\noindent k=0 k=0\par
\noindent 而由 (4.27) 知\par
\noindent X\par
\noindent n+1\par
\noindent Dn+1 (x) = S(n + 1, k)xy k ,\par
\noindent k=0\par
\noindent 所以\par
\noindent X\par
\noindent n X\par
\noindent n+1\par
\noindent k+1 k\par
\noindent S(n, k)(xy + kxy ) = S(n + 1, k)xy k ,\par
\noindent k=0 k=0\par
\noindent 比较 xy k 系数即可得到 S(n, k) 的递推关系:\par
\noindent S(n + 1, k) = S(n, k − 1) + kS(n, k).\par
\noindent 从另一方面考虑:\par
\noindent n\par
}
\end{sourceblock}


\begin{explainblock}
这一段是原书论述链的一部分。读的时候把它当作桥梁：它通常负责把上一个定义、例子或定理，引到下一个计数方法。

\smallskip
\textbf{初学者检查点：}
\begin{itemize}[leftmargin=2em,itemsep=0.25em]
\item 找出这一段承接的上一个概念。
\item 找出它准备引出的下一个概念。
\item 把中间的关键等式或构造单独抄出来。
\end{itemize}

\smallskip
\textbf{本章读法提醒：} 第四章区分两类 Stirling 数：第一类对应排列的圈，第二类对应集合划分。
\end{explainblock}


\begin{sourceblock}{第 4 章续读}
{\small\sloppy
\noindent X\par
\noindent n+1 n nn\par
\noindent D (x) = D (D(x)) = D (xy) = Dj (x)Dn−j (y)\par
\noindent j\par
\noindent j=0\par
\noindent Xn X j\par
\noindent n\par
\noindent = S(j, k)xy k+1\par
\noindent j\par
\noindent j=0 k=0\par
\noindent Xn Xn\par
\noindent n\par
\noindent = S(j, k)xy k+1\par
\noindent j\par
\noindent k=0 j=k\par
\noindent 所以\par
\noindent X\par
\noindent n+1 X n\par
}
\end{sourceblock}


\begin{explainblock}
这一段是原书论述链的一部分。读的时候把它当作桥梁：它通常负责把上一个定义、例子或定理，引到下一个计数方法。

\smallskip
\textbf{初学者检查点：}
\begin{itemize}[leftmargin=2em,itemsep=0.25em]
\item 找出这一段承接的上一个概念。
\item 找出它准备引出的下一个概念。
\item 把中间的关键等式或构造单独抄出来。
\end{itemize}

\smallskip
\textbf{本章读法提醒：} 第四章区分两类 Stirling 数：第一类对应排列的圈，第二类对应集合划分。
\end{explainblock}


\begin{sourceblock}{第 4 章续读}
{\small\sloppy
\noindent n X\par
\noindent k n\par
\noindent S(n + 1, k)xy = S(j, k)xy k+1\par
\noindent j\par
\noindent k=0 k=0 j=k\par
\noindent 比较 y k 的系数得:\par
\noindent n\par
\noindent X n\par
\noindent S(n + 1, k + 1) = S(j, k). (4.28)\par
\noindent j\par
\noindent j=k\par
}
\end{sourceblock}


\begin{explainblock}
这一段是原书论述链的一部分。读的时候把它当作桥梁：它通常负责把上一个定义、例子或定理，引到下一个计数方法。

\smallskip
\textbf{初学者检查点：}
\begin{itemize}[leftmargin=2em,itemsep=0.25em]
\item 找出这一段承接的上一个概念。
\item 找出它准备引出的下一个概念。
\item 把中间的关键等式或构造单独抄出来。
\end{itemize}

\smallskip
\textbf{本章读法提醒：} 第四章区分两类 Stirling 数：第一类对应排列的圈，第二类对应集合划分。
\end{explainblock}


\begin{sourceblock}{例 4.4.2 利用文法证明:}
{\small\sloppy
\noindent 例 4.4.2 利用文法证明:\par
\noindent n\par
\noindent X\par
\noindent i+j n\par
\noindent S(n, i + j) = S(k, i)S(n − k, j). (4.29)\par
\noindent i k\par
\noindent k=0\par
\noindent 证明: 记文法\par
\noindent Gx = \{f → f x, x → x\},\par
\noindent 88 第4章 两类 Stirling 数\par
\noindent Gy = \{f → f y, y → y\}.\par
\noindent 我们来构造一个新的文法:\par
\noindent Gx+y = \{f → f (x + y), x → x, y → y\}.\par
\noindent 不难验证, 在字母表 A = \{f, x, y\} 上, 有 Dx+y = Dx + Dy , 并且在 A 上, 还有\par
\noindent Dx Dy = Dy Dx .\par
\noindent 利用之前结论, 我们有\par
\noindent X\par
\noindent n X\par
}
\end{sourceblock}


\begin{explainblock}
例题是学习这本书最重要的部分。不要只看答案，要把每个限制条件变成一个计数动作：先安排谁、再插入谁、有没有重复、需不需要除以对称性。

\smallskip
\textbf{初学者检查点：}
\begin{itemize}[leftmargin=2em,itemsep=0.25em]
\item 先写出没有限制时的总数。
\item 逐条加入限制条件，观察总数怎样变化。
\item 最后检查有没有把同一种方案重复算了多次。
\end{itemize}

\smallskip
\textbf{本章读法提醒：} 第四章区分两类 Stirling 数：第一类对应排列的圈，第二类对应集合划分。
\end{explainblock}


\begin{sourceblock}{第 4 章续读}
{\small\sloppy
\noindent n k\par
\noindent X\par
\noindent n k k\par
\noindent Dx+y (f ) = f S(n, k)(x + y) = f S(n, k) xi y k−i ,\par
\noindent i\par
\noindent k=1 k=1 i=0\par
\noindent 另一方面,\par
\noindent n\par
\noindent X n\par
\noindent n\par
\noindent Dx+y (f ) = (Dx + Dy )n (f ) = Dxk Dyn−k (f )\par
\noindent k\par
\noindent k=0\par
\noindent n !\par
\noindent X n X\par
\noindent n−k\par
\noindent = k\par
\noindent Dx f S(n − k, j)y j\par
}
\end{sourceblock}


\begin{explainblock}
这一段是原书论述链的一部分。读的时候把它当作桥梁：它通常负责把上一个定义、例子或定理，引到下一个计数方法。

\smallskip
\textbf{初学者检查点：}
\begin{itemize}[leftmargin=2em,itemsep=0.25em]
\item 找出这一段承接的上一个概念。
\item 找出它准备引出的下一个概念。
\item 把中间的关键等式或构造单独抄出来。
\end{itemize}

\smallskip
\textbf{本章读法提醒：} 第四章区分两类 Stirling 数：第一类对应排列的圈，第二类对应集合划分。
\end{explainblock}


\begin{sourceblock}{第 4 章续读}
{\small\sloppy
\noindent k\par
\noindent k=0 j=1\par
\noindent n\par
\noindent X Xk X\par
\noindent n−k\par
\noindent n\par
\noindent = f S(k, i)xi S(n − k, j)y j .\par
\noindent k\par
\noindent k=0 i=1 j=1\par
\noindent 比较 f xi y j 前的系数, 等式得证.\par
}
\end{sourceblock}


\begin{explainblock}
这一段是原书论述链的一部分。读的时候把它当作桥梁：它通常负责把上一个定义、例子或定理，引到下一个计数方法。

\smallskip
\textbf{初学者检查点：}
\begin{itemize}[leftmargin=2em,itemsep=0.25em]
\item 找出这一段承接的上一个概念。
\item 找出它准备引出的下一个概念。
\item 把中间的关键等式或构造单独抄出来。
\end{itemize}

\smallskip
\textbf{本章读法提醒：} 第四章区分两类 Stirling 数：第一类对应排列的圈，第二类对应集合划分。
\end{explainblock}


\begin{sourceblock}{例 4.4.3 利用文法证明:}
{\small\sloppy
\noindent 例 4.4.3 利用文法证明:\par
\noindent X m\par
\noindent S(m + n, k) = im−j S(n, i)S(j, k − i). (4.30)\par
\noindent j\par
\noindent i+j≥k\par
\noindent 证明: 用 Stirling Grammar :\par
\noindent G = \{f → f g, g → g\},\par
\noindent 首先得到\par
\noindent X\par
\noindent m+n\par
\noindent m+n\par
\noindent D (f ) = f S(m + n, k)g k .\par
\noindent k=1\par
\noindent 同时我们有\par
\noindent Dm+n (f ) = Dm Dn (f )\par
\noindent !\par
\noindent X\par
\noindent n\par
}
\end{sourceblock}


\begin{explainblock}
例题是学习这本书最重要的部分。不要只看答案，要把每个限制条件变成一个计数动作：先安排谁、再插入谁、有没有重复、需不需要除以对称性。

\smallskip
\textbf{初学者检查点：}
\begin{itemize}[leftmargin=2em,itemsep=0.25em]
\item 先写出没有限制时的总数。
\item 逐条加入限制条件，观察总数怎样变化。
\item 最后检查有没有把同一种方案重复算了多次。
\end{itemize}

\smallskip
\textbf{本章读法提醒：} 第四章区分两类 Stirling 数：第一类对应排列的圈，第二类对应集合划分。
\end{explainblock}


\begin{sourceblock}{第 4 章续读}
{\small\sloppy
\noindent m i\par
\noindent =D f S(n, i)g\par
\noindent i=1\par
\noindent X\par
\noindent n\par
\noindent = S(n, i)Dm f g i\par
\noindent i=1\par
\noindent X\par
\noindent n m\par
\noindent X m\par
\noindent = S(n, i) Dj (f )im−j g i .\par
\noindent j\par
\noindent i=1 j=0\par
\noindent 比较 f g k 前的系数, 等式得证.\par
\noindent 90 第4章 两类 Stirling 数\par
}
\end{sourceblock}


\begin{explainblock}
这一段是原书论述链的一部分。读的时候把它当作桥梁：它通常负责把上一个定义、例子或定理，引到下一个计数方法。

\smallskip
\textbf{初学者检查点：}
\begin{itemize}[leftmargin=2em,itemsep=0.25em]
\item 找出这一段承接的上一个概念。
\item 找出它准备引出的下一个概念。
\item 把中间的关键等式或构造单独抄出来。
\end{itemize}

\smallskip
\textbf{本章读法提醒：} 第四章区分两类 Stirling 数：第一类对应排列的圈，第二类对应集合划分。
\end{explainblock}


\begin{sourceblock}{习 题 四}
{\small\sloppy
\noindent 习 题 四\par
}
\end{sourceblock}


\begin{explainblock}
习题段不要先追求技巧。先给题目分类：它是在问排列、组合、递推、生成函数、容斥，还是结构计数。分类正确，工具通常就只剩一两个候选。

\smallskip
\textbf{初学者检查点：}
\begin{itemize}[leftmargin=2em,itemsep=0.25em]
\item 判断题目所属工具。
\item 列出变量和限制条件。
\item 先做小规模 n=2、3、4 的例子，确认直觉。
\end{itemize}

\smallskip
\textbf{本章读法提醒：} 第四章区分两类 Stirling 数：第一类对应排列的圈，第二类对应集合划分。
\end{explainblock}


\begin{sourceblock}{习题 4.1 证明无符号的第一类 Stirling 数满足}
{\small\sloppy
\noindent 习题 4.1 证明无符号的第一类 Stirling 数满足\par
\noindent i) c(n, 1) = (n − 1)! (n ≥ 1);\par
\noindent n\par
\noindent ii) c(n, n − 1) = (n ≥ 1).\par
\noindent n n\par
\noindent iii) )s(n, n − 2) = 2 +3 .\par
}
\end{sourceblock}


\begin{explainblock}
习题段不要先追求技巧。先给题目分类：它是在问排列、组合、递推、生成函数、容斥，还是结构计数。分类正确，工具通常就只剩一两个候选。

\smallskip
\textbf{初学者检查点：}
\begin{itemize}[leftmargin=2em,itemsep=0.25em]
\item 判断题目所属工具。
\item 列出变量和限制条件。
\item 先做小规模 n=2、3、4 的例子，确认直觉。
\end{itemize}

\smallskip
\textbf{本章读法提醒：} 第四章区分两类 Stirling 数：第一类对应排列的圈，第二类对应集合划分。
\end{explainblock}


\begin{sourceblock}{习题 4.2 设 c(n, m) 表示无符号的第一类 Stirling 数. 证明}
{\small\sloppy
\noindent 习题 4.2 设 c(n, m) 表示无符号的第一类 Stirling 数. 证明\par
\noindent m m+1 n\par
\noindent c(n, m) + c(n, m + 1) + · · · + c(n, n) = c(n + 1, m + 1).\par
\noindent m m m\par
}
\end{sourceblock}


\begin{explainblock}
习题段不要先追求技巧。先给题目分类：它是在问排列、组合、递推、生成函数、容斥，还是结构计数。分类正确，工具通常就只剩一两个候选。

\smallskip
\textbf{初学者检查点：}
\begin{itemize}[leftmargin=2em,itemsep=0.25em]
\item 判断题目所属工具。
\item 列出变量和限制条件。
\item 先做小规模 n=2、3、4 的例子，确认直觉。
\end{itemize}

\smallskip
\textbf{本章读法提醒：} 第四章区分两类 Stirling 数：第一类对应排列的圈，第二类对应集合划分。
\end{explainblock}


\begin{sourceblock}{习题 4.3 计算第二类 Stirling 数 S(8, k) (k = 0, 1, . . . , 8), 并求 Bell 数 B8 .}
{\small\sloppy
\noindent 习题 4.3 计算第二类 Stirling 数 S(8, k) (k = 0, 1, . . . , 8), 并求 Bell 数 B8 .\par
}
\end{sourceblock}


\begin{explainblock}
习题段不要先追求技巧。先给题目分类：它是在问排列、组合、递推、生成函数、容斥，还是结构计数。分类正确，工具通常就只剩一两个候选。

\smallskip
\textbf{初学者检查点：}
\begin{itemize}[leftmargin=2em,itemsep=0.25em]
\item 判断题目所属工具。
\item 列出变量和限制条件。
\item 先做小规模 n=2、3、4 的例子，确认直觉。
\end{itemize}

\smallskip
\textbf{本章读法提醒：} 第四章区分两类 Stirling 数：第一类对应排列的圈，第二类对应集合划分。
\end{explainblock}


\begin{sourceblock}{习题 4.4 证明第二类 Stirling 数满足关系}
{\small\sloppy
\noindent 习题 4.4 证明第二类 Stirling 数满足关系\par
\noindent i) S(n, 1) = 1 (n ≥ 1);\par
\noindent ii) S(n, 2) = 2n−1 − 1 (n ≥ 2);\par
\noindent n\par
\noindent iii) S(n, n − 1) = (n ≥ 1);\par
\noindent n n\par
\noindent iv)S(n, n − 2) = +3 (n ≥ 2).\par
}
\end{sourceblock}


\begin{explainblock}
习题段不要先追求技巧。先给题目分类：它是在问排列、组合、递推、生成函数、容斥，还是结构计数。分类正确，工具通常就只剩一两个候选。

\smallskip
\textbf{初学者检查点：}
\begin{itemize}[leftmargin=2em,itemsep=0.25em]
\item 判断题目所属工具。
\item 列出变量和限制条件。
\item 先做小规模 n=2、3、4 的例子，确认直觉。
\end{itemize}

\smallskip
\textbf{本章读法提醒：} 第四章区分两类 Stirling 数：第一类对应排列的圈，第二类对应集合划分。
\end{explainblock}


\begin{sourceblock}{习题 4.5 证明}
{\small\sloppy
\noindent 习题 4.5 证明\par
\noindent n n n\par
\noindent S(n + 1, k) = S(0, k − 1) + S(1, k − 1) + · · · + S(n, k − 1)\par
}
\end{sourceblock}


\begin{explainblock}
习题段不要先追求技巧。先给题目分类：它是在问排列、组合、递推、生成函数、容斥，还是结构计数。分类正确，工具通常就只剩一两个候选。

\smallskip
\textbf{初学者检查点：}
\begin{itemize}[leftmargin=2em,itemsep=0.25em]
\item 判断题目所属工具。
\item 列出变量和限制条件。
\item 先做小规模 n=2、3、4 的例子，确认直觉。
\end{itemize}

\smallskip
\textbf{本章读法提醒：} 第四章区分两类 Stirling 数：第一类对应排列的圈，第二类对应集合划分。
\end{explainblock}


\begin{sourceblock}{习题 4.6 定义 Si (n, k) 是 n 元集合的每块至少有 i 个元的 k 部划分的个数}
{\small\sloppy
\noindent 习题 4.6 定义 Si (n, k) 是 n 元集合的每块至少有 i 个元的 k 部划分的个数\par
\noindent (ki ≤ n). 证明\par
\noindent n−1\par
\noindent Si (n, k) = Si (n − i, k − 1) + kSi (n − 1, k).\par
\noindent i−1\par
}
\end{sourceblock}


\begin{explainblock}
习题段不要先追求技巧。先给题目分类：它是在问排列、组合、递推、生成函数、容斥，还是结构计数。分类正确，工具通常就只剩一两个候选。

\smallskip
\textbf{初学者检查点：}
\begin{itemize}[leftmargin=2em,itemsep=0.25em]
\item 判断题目所属工具。
\item 列出变量和限制条件。
\item 先做小规模 n=2、3、4 的例子，确认直觉。
\end{itemize}

\smallskip
\textbf{本章读法提醒：} 第四章区分两类 Stirling 数：第一类对应排列的圈，第二类对应集合划分。
\end{explainblock}


\begin{sourceblock}{习题 4.7 求 15 + 25 + 35 + · · · + n5 的和, 并验证它是 Tn = n(n+1)}
{\small\sloppy
\noindent 习题 4.7 求 15 + 25 + 35 + · · · + n5 的和, 并验证它是 Tn = n(n+1)\par
\noindent Pb\par
}
\end{sourceblock}


\begin{explainblock}
习题段不要先追求技巧。先给题目分类：它是在问排列、组合、递推、生成函数、容斥，还是结构计数。分类正确，工具通常就只剩一两个候选。

\smallskip
\textbf{初学者检查点：}
\begin{itemize}[leftmargin=2em,itemsep=0.25em]
\item 判断题目所属工具。
\item 列出变量和限制条件。
\item 先做小规模 n=2、3、4 的例子，确认直觉。
\end{itemize}

\smallskip
\textbf{本章读法提醒：} 第四章区分两类 Stirling 数：第一类对应排列的圈，第二类对应集合划分。
\end{explainblock}


\begin{sourceblock}{习题 4.8 i}
{\small\sloppy
\noindent 习题 4.8 i\par
\noindent 计算 i=a n , 其中 b > a ≥ n.\par
}
\end{sourceblock}


\begin{explainblock}
习题段不要先追求技巧。先给题目分类：它是在问排列、组合、递推、生成函数、容斥，还是结构计数。分类正确，工具通常就只剩一两个候选。

\smallskip
\textbf{初学者检查点：}
\begin{itemize}[leftmargin=2em,itemsep=0.25em]
\item 判断题目所属工具。
\item 列出变量和限制条件。
\item 先做小规模 n=2、3、4 的例子，确认直觉。
\end{itemize}

\smallskip
\textbf{本章读法提醒：} 第四章区分两类 Stirling 数：第一类对应排列的圈，第二类对应集合划分。
\end{explainblock}


\begin{sourceblock}{习题 4.9 设 0 ≤ k ≤ n , 利用 S(n, k) 的递推关系证明}
{\small\sloppy
\noindent 习题 4.9 设 0 ≤ k ≤ n , 利用 S(n, k) 的递推关系证明\par
\noindent X\par
\noindent n\par
\noindent xn = S(n, k)(x)k .\par
\noindent k=0\par
}
\end{sourceblock}


\begin{explainblock}
习题段不要先追求技巧。先给题目分类：它是在问排列、组合、递推、生成函数、容斥，还是结构计数。分类正确，工具通常就只剩一两个候选。

\smallskip
\textbf{初学者检查点：}
\begin{itemize}[leftmargin=2em,itemsep=0.25em]
\item 判断题目所属工具。
\item 列出变量和限制条件。
\item 先做小规模 n=2、3、4 的例子，确认直觉。
\end{itemize}

\smallskip
\textbf{本章读法提醒：} 第四章区分两类 Stirling 数：第一类对应排列的圈，第二类对应集合划分。
\end{explainblock}


\begin{sourceblock}{习题 4.10 证明当 n ≥ 0 时, Bell 数 Bn 满足下面恒等式}
{\small\sloppy
\noindent 习题 4.10 证明当 n ≥ 0 时, Bell 数 Bn 满足下面恒等式\par
\noindent n\par
\noindent X n\par
\noindent Bn+1 = Bk .\par
\noindent k\par
\noindent k=0\par
\noindent 92 第4章 两类 Stirling 数\par
}
\end{sourceblock}


\begin{explainblock}
习题段不要先追求技巧。先给题目分类：它是在问排列、组合、递推、生成函数、容斥，还是结构计数。分类正确，工具通常就只剩一两个候选。

\smallskip
\textbf{初学者检查点：}
\begin{itemize}[leftmargin=2em,itemsep=0.25em]
\item 判断题目所属工具。
\item 列出变量和限制条件。
\item 先做小规模 n=2、3、4 的例子，确认直觉。
\end{itemize}

\smallskip
\textbf{本章读法提醒：} 第四章区分两类 Stirling 数：第一类对应排列的圈，第二类对应集合划分。
\end{explainblock}


\begin{sourceblock}{习题 4.11 设 Gn (x) 是序列 \{ S(m,n)}
{\small\sloppy
\noindent 习题 4.11 设 Gn (x) 是序列 \{ S(m,n)\par
\noindent m! \}m≥n 的普通型生成函数, 即\par
\noindent ∞\par
\noindent X S(m, n)\par
\noindent Gn (x) = xm .\par
\noindent m=n\par
\noindent m!\par
\noindent 证明:\par
\noindent (1)\par
\noindent d\par
\noindent Gn (x) − nGn (x) = Gn−1 (x), Gn (0) = 0.\par
\noindent dx\par
\noindent (2)\par
\noindent G1 (x) = ex − 1.\par
}
\end{sourceblock}


\begin{explainblock}
习题段不要先追求技巧。先给题目分类：它是在问排列、组合、递推、生成函数、容斥，还是结构计数。分类正确，工具通常就只剩一两个候选。

\smallskip
\textbf{初学者检查点：}
\begin{itemize}[leftmargin=2em,itemsep=0.25em]
\item 判断题目所属工具。
\item 列出变量和限制条件。
\item 先做小规模 n=2、3、4 的例子，确认直觉。
\end{itemize}

\smallskip
\textbf{本章读法提醒：} 第四章区分两类 Stirling 数：第一类对应排列的圈，第二类对应集合划分。
\end{explainblock}


\begin{sourceblock}{习题 4.12 设 Gn (x) 是序列 \{ c(m,n)}
{\small\sloppy
\noindent 习题 4.12 设 Gn (x) 是序列 \{ c(m,n)\par
\noindent m! \}m≥n 的普通型生成函数, 即\par
\noindent ∞\par
\noindent X c(m, n)\par
\noindent Gn (x) = xm .\par
\noindent m=n\par
\noindent m!\par
\noindent 证明:\par
\noindent (1)\par
\noindent d\par
\noindent (1 − x) Gn (x) = Gn−1 (x), Gn (0) = 0.\par
\noindent dx\par
\noindent (2)\par
\noindent G1 (x) = − log(1 − x).\par
\noindent Pn\par
}
\end{sourceblock}


\begin{explainblock}
习题段不要先追求技巧。先给题目分类：它是在问排列、组合、递推、生成函数、容斥，还是结构计数。分类正确，工具通常就只剩一两个候选。

\smallskip
\textbf{初学者检查点：}
\begin{itemize}[leftmargin=2em,itemsep=0.25em]
\item 判断题目所属工具。
\item 列出变量和限制条件。
\item 先做小规模 n=2、3、4 的例子，确认直觉。
\end{itemize}

\smallskip
\textbf{本章读法提醒：} 第四章区分两类 Stirling 数：第一类对应排列的圈，第二类对应集合划分。
\end{explainblock}


\begin{sourceblock}{习题 4.13 设 f (n) = 2n2 − n + 3, 确定其差分表, 并求出 k=0 f (k) 的表达式.}
{\small\sloppy
\noindent 习题 4.13 设 f (n) = 2n2 − n + 3, 确定其差分表, 并求出 k=0 f (k) 的表达式.\par
}
\end{sourceblock}


\begin{explainblock}
习题段不要先追求技巧。先给题目分类：它是在问排列、组合、递推、生成函数、容斥，还是结构计数。分类正确，工具通常就只剩一两个候选。

\smallskip
\textbf{初学者检查点：}
\begin{itemize}[leftmargin=2em,itemsep=0.25em]
\item 判断题目所属工具。
\item 列出变量和限制条件。
\item 先做小规模 n=2、3、4 的例子，确认直觉。
\end{itemize}

\smallskip
\textbf{本章读法提醒：} 第四章区分两类 Stirling 数：第一类对应排列的圈，第二类对应集合划分。
\end{explainblock}


\begin{sourceblock}{习题 4.14 设 f (n) 是 n 的三次多项式, 如果其差分表的第 1 行的前 4 个数是}
{\small\sloppy
\noindent 习题 4.14 设 f (n) 是 n 的三次多项式, 如果其差分表的第 1 行的前 4 个数是\par
\noindent P\par
\noindent 1, −1, 1, 10, 试确定 f (n) 并计算 nk=0 f (k).\par
}
\end{sourceblock}


\begin{explainblock}
习题段不要先追求技巧。先给题目分类：它是在问排列、组合、递推、生成函数、容斥，还是结构计数。分类正确，工具通常就只剩一两个候选。

\smallskip
\textbf{初学者检查点：}
\begin{itemize}[leftmargin=2em,itemsep=0.25em]
\item 判断题目所属工具。
\item 列出变量和限制条件。
\item 先做小规模 n=2、3、4 的例子，确认直觉。
\end{itemize}

\smallskip
\textbf{本章读法提醒：} 第四章区分两类 Stirling 数：第一类对应排列的圈，第二类对应集合划分。
\end{explainblock}


\begin{sourceblock}{习题 4.15 如果 f (n) 是 n 的 m 次多项式, 证明使得}
{\small\sloppy
\noindent 习题 4.15 如果 f (n) 是 n 的 m 次多项式, 证明使得\par
\noindent n n n\par
\noindent f (n) = c0 + c1 + · · · + cm ,\par
\noindent 的常数 c0 , c1 , . . . , cm 是唯一确定的.\par
}
\end{sourceblock}


\begin{explainblock}
习题段不要先追求技巧。先给题目分类：它是在问排列、组合、递推、生成函数、容斥，还是结构计数。分类正确，工具通常就只剩一两个候选。

\smallskip
\textbf{初学者检查点：}
\begin{itemize}[leftmargin=2em,itemsep=0.25em]
\item 判断题目所属工具。
\item 列出变量和限制条件。
\item 先做小规模 n=2、3、4 的例子，确认直觉。
\end{itemize}

\smallskip
\textbf{本章读法提醒：} 第四章区分两类 Stirling 数：第一类对应排列的圈，第二类对应集合划分。
\end{explainblock}


\begin{sourceblock}{习题 4.16 证明: 设 f (n) 是 n 的 m 次多项式, 则}
{\small\sloppy
\noindent 习题 4.16 证明: 设 f (n) 是 n 的 m 次多项式, 则\par
\noindent X\par
\noindent n m\par
\noindent X\par
\noindent n+1\par
\noindent f (k) = ∆j f (0).\par
\noindent j+1\par
\noindent k=0 j=0\par
}
\end{sourceblock}


\begin{explainblock}
习题段不要先追求技巧。先给题目分类：它是在问排列、组合、递推、生成函数、容斥，还是结构计数。分类正确，工具通常就只剩一两个候选。

\smallskip
\textbf{初学者检查点：}
\begin{itemize}[leftmargin=2em,itemsep=0.25em]
\item 判断题目所属工具。
\item 列出变量和限制条件。
\item 先做小规模 n=2、3、4 的例子，确认直觉。
\end{itemize}

\smallskip
\textbf{本章读法提醒：} 第四章区分两类 Stirling 数：第一类对应排列的圈，第二类对应集合划分。
\end{explainblock}


\begin{sourceblock}{习题 4.17 利用差分的方法求和}
{\small\sloppy
\noindent 习题 4.17 利用差分的方法求和\par
\noindent X\par
\noindent n\par
\noindent k(k − 1)(k + 3). (4.31)\par
\noindent k=0\par
}
\end{sourceblock}


\begin{explainblock}
习题段不要先追求技巧。先给题目分类：它是在问排列、组合、递推、生成函数、容斥，还是结构计数。分类正确，工具通常就只剩一两个候选。

\smallskip
\textbf{初学者检查点：}
\begin{itemize}[leftmargin=2em,itemsep=0.25em]
\item 判断题目所属工具。
\item 列出变量和限制条件。
\item 先做小规模 n=2、3、4 的例子，确认直觉。
\end{itemize}

\smallskip
\textbf{本章读法提醒：} 第四章区分两类 Stirling 数：第一类对应排列的圈，第二类对应集合划分。
\end{explainblock}


\chapter{排列与两类欧拉数}
\begin{roadmapblock}
本章保留原书正文的主要数学骨架，并在每个定义、定理、例题和论述片段后加入解读。
阅读时不要把目标放在“背下答案”上，而是放在“看懂这一步为什么这样数”上。

\smallskip
\textbf{本章最低目标：} 能说出每个公式左边在数什么、右边在数什么，以及为什么两边相等。
\end{roadmapblock}


\begin{sourceblock}{第 5 章正文}
{\small\sloppy
\noindent 排列与两类欧拉数\par
\noindent 在这一章, 我们将介绍由欧拉定义在排列上的两类组合数. 为了区分, 我们称定\par
\noindent 义在排列上下降位的组合数为 Eulerian 数, 与交错排列有关的组合数为 Euler 数.\par
}
\end{sourceblock}


\begin{explainblock}
这一段是原书论述链的一部分。读的时候把它当作桥梁：它通常负责把上一个定义、例子或定理，引到下一个计数方法。

\smallskip
\textbf{初学者检查点：}
\begin{itemize}[leftmargin=2em,itemsep=0.25em]
\item 找出这一段承接的上一个概念。
\item 找出它准备引出的下一个概念。
\item 把中间的关键等式或构造单独抄出来。
\end{itemize}

\smallskip
\textbf{本章读法提醒：} 第五章注意下降数、上升数、逆序数不是同一个统计量。
\end{explainblock}


\begin{sourceblock}{§ 5.1 Eulerian 数的定义与性质}
{\small\sloppy
\noindent § 5.1. Eulerian 数的定义与性质\par
\noindent P\par
\noindent Eulerian 数是欧拉在研究交错的和式 (−1)k k n 时首先发现的. 近二百年来,\par
\noindent Eulerian 数不断的从不同方向被发现, 在数论, 组合数学, 及概率与统计等方向都有\par
\noindent 所应用. 其重要性不言而喻.\par
\noindent 我们首先引入排列的一个重要的统计量: 下降数 (上升数). 此外排列还有其它各\par
\noindent 样统计量, 对排列的统计量的研究是组合数学中一个重要课题, 有兴趣的读者可参\par
\noindent 考 [59].\par
}
\end{sourceblock}


\begin{explainblock}
这一节先不要急着背公式。先问三个问题：对象是什么，哪些限制条件不能丢，最后到底是在数所有对象、还是在数某个统计量的分布。把这三件事说清楚，后面的公式才会有位置。

\smallskip
\textbf{初学者检查点：}
\begin{itemize}[leftmargin=2em,itemsep=0.25em]
\item 把本节出现的新对象写成一句话定义。
\item 把每个公式左边和右边分别翻译成“在数什么”。
\item 找一个最小非平凡例子，手工列出所有对象。
\end{itemize}

\smallskip
\textbf{本章读法提醒：} 第五章注意下降数、上升数、逆序数不是同一个统计量。
\end{explainblock}


\begin{sourceblock}{定义 5.1.1 (下降数、上升数) 对于一个排列 π = π1 π2 · · · πn , 若 πi > πi+1 , 那么}
{\small\sloppy
\noindent 定义 5.1.1 (下降数、上升数) 对于一个排列 π = π1 π2 · · · πn , 若 πi > πi+1 , 那么\par
\noindent 我们称 i 为一个下降位 (desent), 用 des(π) 来表示 π 下降位的个数; 若 πi < πi+1 ,\par
\noindent 则称 i 是一个上升位 (ascent), 用 asc(π) 来表示 π 上升位的个数.\par
\noindent 注意下降位与上升位指得是排列 π 中元素的位置而不是 π 中的元素.\par
\noindent 对任意排列 π, 由定义显然知\par
\noindent des(π) + asc(π) = n − 1. (5.1)\par
}
\end{sourceblock}


\begin{explainblock}
定义段的任务是定边界。这里要特别注意参数范围、是否允许重复、是否考虑顺序、对象有没有标签。后面的每一道例题和证明，本质上都在反复使用这些边界。

\smallskip
\textbf{初学者检查点：}
\begin{itemize}[leftmargin=2em,itemsep=0.25em]
\item 圈出被定义的对象名称。
\item 圈出参数范围，例如 n、k 是否要求正整数。
\item 判断它是有序对象、无序对象、带标签对象还是结构对象。
\end{itemize}

\smallskip
\textbf{本章读法提醒：} 第五章注意下降数、上升数、逆序数不是同一个统计量。
\end{explainblock}


\begin{sourceblock}{例 5.1.1 π = 43521 是集合 [5] 上的一个排列, 它的下降位为 1, 3, 4, 上升位为 2,}
{\small\sloppy
\noindent 例 5.1.1 π = 43521 是集合 [5] 上的一个排列, 它的下降位为 1, 3, 4, 上升位为 2,\par
\noindent 所以 des(π) = 3, asc(π) = 1.\par
}
\end{sourceblock}


\begin{explainblock}
例题是学习这本书最重要的部分。不要只看答案，要把每个限制条件变成一个计数动作：先安排谁、再插入谁、有没有重复、需不需要除以对称性。

\smallskip
\textbf{初学者检查点：}
\begin{itemize}[leftmargin=2em,itemsep=0.25em]
\item 先写出没有限制时的总数。
\item 逐条加入限制条件，观察总数怎样变化。
\item 最后检查有没有把同一种方案重复算了多次。
\end{itemize}

\smallskip
\textbf{本章读法提醒：} 第五章注意下降数、上升数、逆序数不是同一个统计量。
\end{explainblock}


\begin{sourceblock}{定义 5.1.2 集合 n 的所有排列中下降位个数为 k − 1 的排列的个数, 称为第 k 个}
{\small\sloppy
\noindent 定义 5.1.2 集合 n 的所有排列中下降位个数为 k − 1 的排列的个数, 称为第 k 个\par
\noindent Eulerian 数, 记为 A(n, k). 由 (5.1) 知 A(n, k) 也表示集合 [n] 的全部排列中具有\par
\noindent n − k 个上升位的排列的个数.\par
}
\end{sourceblock}


\begin{explainblock}
定义段的任务是定边界。这里要特别注意参数范围、是否允许重复、是否考虑顺序、对象有没有标签。后面的每一道例题和证明，本质上都在反复使用这些边界。

\smallskip
\textbf{初学者检查点：}
\begin{itemize}[leftmargin=2em,itemsep=0.25em]
\item 圈出被定义的对象名称。
\item 圈出参数范围，例如 n、k 是否要求正整数。
\item 判断它是有序对象、无序对象、带标签对象还是结构对象。
\end{itemize}

\smallskip
\textbf{本章读法提醒：} 第五章注意下降数、上升数、逆序数不是同一个统计量。
\end{explainblock}


\begin{sourceblock}{注记 5.1.3 由定义可知 A(n, 1) = A(n, n) = 1, 且有}
{\small\sloppy
\noindent 注记 5.1.3 由定义可知 A(n, 1) = A(n, n) = 1, 且有\par
\noindent X\par
\noindent n\par
\noindent A(n, k) = n!.\par
\noindent k=1\par
\noindent 并规定当 n ≥ 1 时, A(n, 0) = 0 且 A(0, 0) = 1.\par
}
\end{sourceblock}


\begin{explainblock}
注记通常是在补充术语、记号、历史或一个容易忽略的约定。它未必给新公式，但常常决定你后面会不会把符号读错。

\smallskip
\textbf{初学者检查点：}
\begin{itemize}[leftmargin=2em,itemsep=0.25em]
\item 记下新符号的读法。
\item 确认这个约定是否只在本章使用。
\item 把注记里的限制条件补到前面的定义旁边。
\end{itemize}

\smallskip
\textbf{本章读法提醒：} 第五章注意下降数、上升数、逆序数不是同一个统计量。
\end{explainblock}


\begin{sourceblock}{例 5.1.2 集合 \{1, 2, 3\} 中的所有排列有}
{\small\sloppy
\noindent 例 5.1.2 集合 \{1, 2, 3\} 中的所有排列有\par
\noindent 123, 132, 213, 231, 312, 321,\par
\noindent 其中有 0 个下降位的有 123, 有 1 个下降位的有\par
\noindent 132, 213, 231, 312,\par
\noindent 有两个下降位的有 321, 所以\par
\noindent A(3, 1) = 1, A(3, 2) = 4, A(3, 3) = 1.\par
\noindent k\par
\noindent 1 2 3 4 5 6 7 8 9 ···\par
\noindent n\par
\noindent 6 1 57 302 302 57 1\par
\noindent 7 1 120 1191 2416 1191 120 1\par
\noindent 8 1 247 4293 15619 15619 4293 247 1\par
\noindent 9 1 502 14608 88234 156190 88234 14608 502 1\par
\noindent .. .. .. .. .. .. .. .. .. .. ..\par
\noindent . . . . . . . . . . .\par
\noindent 性质 5.1.4 (对称性) 设 n ≥ k ≥ 0, 则\par
\noindent A(n, k) = A(n, n − k + 1).\par
\noindent 证明: 设 π = π1 π2 · · · πn 是集合 [n] 上任意一个排列, 记 π r = πn πn−1 · · · π1 ,\par
}
\end{sourceblock}


\begin{explainblock}
例题是学习这本书最重要的部分。不要只看答案，要把每个限制条件变成一个计数动作：先安排谁、再插入谁、有没有重复、需不需要除以对称性。

\smallskip
\textbf{初学者检查点：}
\begin{itemize}[leftmargin=2em,itemsep=0.25em]
\item 先写出没有限制时的总数。
\item 逐条加入限制条件，观察总数怎样变化。
\item 最后检查有没有把同一种方案重复算了多次。
\end{itemize}

\smallskip
\textbf{本章读法提醒：} 第五章注意下降数、上升数、逆序数不是同一个统计量。
\end{explainblock}


\begin{sourceblock}{第 5 章续读}
{\small\sloppy
\noindent 称为 π 的转置, 显然转置是集合 [n] 上所有排列组成的集合上的对合映射. 由转置的\par
}
\end{sourceblock}


\begin{explainblock}
这一段是原书论述链的一部分。读的时候把它当作桥梁：它通常负责把上一个定义、例子或定理，引到下一个计数方法。

\smallskip
\textbf{初学者检查点：}
\begin{itemize}[leftmargin=2em,itemsep=0.25em]
\item 找出这一段承接的上一个概念。
\item 找出它准备引出的下一个概念。
\item 把中间的关键等式或构造单独抄出来。
\end{itemize}

\smallskip
\textbf{本章读法提醒：} 第五章注意下降数、上升数、逆序数不是同一个统计量。
\end{explainblock}


\begin{sourceblock}{定义可知, 若 π 具有 k − 1 个下降位, 则 π r 有 n − k 个下降位, 又 π 与 π r 一一对}
{\small\sloppy
\noindent 定义可知, 若 π 具有 k − 1 个下降位, 则 π r 有 n − k 个下降位, 又 π 与 π r 一一对\par
\noindent 应, 命题得证.\par
}
\end{sourceblock}


\begin{explainblock}
定义段的任务是定边界。这里要特别注意参数范围、是否允许重复、是否考虑顺序、对象有没有标签。后面的每一道例题和证明，本质上都在反复使用这些边界。

\smallskip
\textbf{初学者检查点：}
\begin{itemize}[leftmargin=2em,itemsep=0.25em]
\item 圈出被定义的对象名称。
\item 圈出参数范围，例如 n、k 是否要求正整数。
\item 判断它是有序对象、无序对象、带标签对象还是结构对象。
\end{itemize}

\smallskip
\textbf{本章读法提醒：} 第五章注意下降数、上升数、逆序数不是同一个统计量。
\end{explainblock}


\begin{sourceblock}{定理 5.1.5 (递推关系) 设 n ≥ 1, k ≥ 1, 则}
{\small\sloppy
\noindent 定理 5.1.5 (递推关系) 设 n ≥ 1, k ≥ 1, 则\par
\noindent A(n, k) = kA(n − 1, k) + (n − k + 1)A(n − 1, k − 1). (5.2)\par
\noindent 证明: 因为任意一个 n 元排列都可通过将 n 插入到 n − 1 元排列中得到, 所以\par
\noindent 不妨考虑将 n 插入到 n − 1 元排列中时下降位的变化. 设 π 是集合 [n − 1] 上的具\par
\noindent 有 r − 1 个下降位的排列, 则通过将 n 插入到 π 中任意位置显然可得到下降位个数\par
\noindent 不同两类新排列:\par
\noindent (1) 若将 n 插入到 π 的 r − 1 个下降位的后面或最后一个位置 (共计 r 个位置),\par
\noindent 5.2 Eulerian 多项式 95\par
\noindent 则得到的 n 元排列也有 r − 1 个下降位;\par
\noindent (2) 若将 n 插入到非下降位且不能插到最后一位 (共计 n − r 个位置), 则得到的 n\par
\noindent 元排列有 r 个下降位.\par
\noindent 所以具有 k − 1 个下降位的 n 元排列只能通过将 n 插入到具有 k − 1 或 k − 2 个下\par
\noindent 降位的 n − 1 元排列中得到, 且具有 k − 1 个下降位的 n − 1 元的排列只有 k 个可\par
\noindent 插入的位置能保证下降位个数不变, 而具有 k − 2 个下降位的 n − 1 元的排列只有\par
\noindent n − (k − 1) 个可插入的位置能使下降位个数加 1, 故式 (5.2) 成立. 定理得证.\par
}
\end{sourceblock}


\begin{explainblock}
定理段不要只看结论，要看它把哪两个计数方式连在一起。组合学证明最常见的技巧是：左边按一种标准数，右边按另一种标准数，然后说明两边数的是同一批对象。

\smallskip
\textbf{初学者检查点：}
\begin{itemize}[leftmargin=2em,itemsep=0.25em]
\item 先复述定理的假设，不要把适用范围扩大。
\item 找出证明中的基本计数原则：乘法、加法、双射、递推、容斥或生成函数。
\item 检查最后的等式化简是否和前面的对象解释一致。
\end{itemize}

\smallskip
\textbf{本章读法提醒：} 第五章注意下降数、上升数、逆序数不是同一个统计量。
\end{explainblock}


\begin{sourceblock}{注记 5.1.6 由 A(n, k) 的递推关系, 很容易验证 A(n, k) 的对称性.}
{\small\sloppy
\noindent 注记 5.1.6 由 A(n, k) 的递推关系, 很容易验证 A(n, k) 的对称性.\par
}
\end{sourceblock}


\begin{explainblock}
注记通常是在补充术语、记号、历史或一个容易忽略的约定。它未必给新公式，但常常决定你后面会不会把符号读错。

\smallskip
\textbf{初学者检查点：}
\begin{itemize}[leftmargin=2em,itemsep=0.25em]
\item 记下新符号的读法。
\item 确认这个约定是否只在本章使用。
\item 把注记里的限制条件补到前面的定义旁边。
\end{itemize}

\smallskip
\textbf{本章读法提醒：} 第五章注意下降数、上升数、逆序数不是同一个统计量。
\end{explainblock}


\begin{sourceblock}{§ 5.2 Eulerian 多项式}
{\small\sloppy
\noindent § 5.2. Eulerian 多项式\par
\noindent 设 An (x) 为 Eulerian 数的普通型生成函数, 即:\par
\noindent X\par
\noindent n\par
\noindent An (x) = A(n, k)xk .\par
\noindent k=0\par
\noindent 显然 An (x) 是关于 x 的 n 次多项式, 称为 Eulerian 多项式. 由 Eulerian 数的组合\par
\noindent 意义可知:\par
\noindent X\par
\noindent An (x) = xdes(π)+1 ,\par
\noindent π∈Sn\par
\noindent 其中 Sn 为集合 [n] 的所有排列.\par
}
\end{sourceblock}


\begin{explainblock}
这一节先不要急着背公式。先问三个问题：对象是什么，哪些限制条件不能丢，最后到底是在数所有对象、还是在数某个统计量的分布。把这三件事说清楚，后面的公式才会有位置。

\smallskip
\textbf{初学者检查点：}
\begin{itemize}[leftmargin=2em,itemsep=0.25em]
\item 把本节出现的新对象写成一句话定义。
\item 把每个公式左边和右边分别翻译成“在数什么”。
\item 找一个最小非平凡例子，手工列出所有对象。
\end{itemize}

\smallskip
\textbf{本章读法提醒：} 第五章注意下降数、上升数、逆序数不是同一个统计量。
\end{explainblock}


\begin{sourceblock}{例 5.2.1 前 5 个 Eulerian 多项式:}
{\small\sloppy
\noindent 例 5.2.1 前 5 个 Eulerian 多项式:\par
\noindent A0 (x) =1;\par
\noindent A1 (x) =x;\par
\noindent A2 (x) =x + x2 ;\par
\noindent A3 (x) =x + 4x2 + x3 ;\par
\noindent A4 (x) =x + 11x2 + 11x3 + x4 ;\par
\noindent A5 (x) =x + 26x2 + 66x3 + 26x4 + x5 .\par
}
\end{sourceblock}


\begin{explainblock}
例题是学习这本书最重要的部分。不要只看答案，要把每个限制条件变成一个计数动作：先安排谁、再插入谁、有没有重复、需不需要除以对称性。

\smallskip
\textbf{初学者检查点：}
\begin{itemize}[leftmargin=2em,itemsep=0.25em]
\item 先写出没有限制时的总数。
\item 逐条加入限制条件，观察总数怎样变化。
\item 最后检查有没有把同一种方案重复算了多次。
\end{itemize}

\smallskip
\textbf{本章读法提醒：} 第五章注意下降数、上升数、逆序数不是同一个统计量。
\end{explainblock}


\begin{sourceblock}{定理 5.2.1 设 n ≥ 1, 则}
{\small\sloppy
\noindent 定理 5.2.1 设 n ≥ 1, 则\par
\noindent An (x) = x(1 − x)A′n−1 (x) + nxAn−1 (x). (5.3)\par
\noindent 证明: 由 A(n, k) 的递推关系 (5.2) 可得\par
\noindent X\par
\noindent n X\par
\noindent n\par
\noindent A(n, k)xk = ((kA(n − 1, k) + (n − k + 1)A(n − 1, k − 1))xk\par
\noindent k=1 k=1\par
\noindent X\par
\noindent n−1 X\par
\noindent n\par
\noindent = kA(n − 1, k)xk + (n − k + 1)A(n − 1, k − 1)xk\par
\noindent k=1 k=2\par
\noindent X\par
\noindent n−1\par
\noindent = xA′n−1 (x) + (n − k)A(n − 1, k)xk+1\par
\noindent k=1\par
\noindent X\par
}
\end{sourceblock}


\begin{explainblock}
定理段不要只看结论，要看它把哪两个计数方式连在一起。组合学证明最常见的技巧是：左边按一种标准数，右边按另一种标准数，然后说明两边数的是同一批对象。

\smallskip
\textbf{初学者检查点：}
\begin{itemize}[leftmargin=2em,itemsep=0.25em]
\item 先复述定理的假设，不要把适用范围扩大。
\item 找出证明中的基本计数原则：乘法、加法、双射、递推、容斥或生成函数。
\item 检查最后的等式化简是否和前面的对象解释一致。
\end{itemize}

\smallskip
\textbf{本章读法提醒：} 第五章注意下降数、上升数、逆序数不是同一个统计量。
\end{explainblock}


\begin{sourceblock}{第 5 章续读}
{\small\sloppy
\noindent n−1 X\par
\noindent n−1\par
\noindent = xA′n−1 (x) + nA(n − 1, k)xk+1 − x2 kA(n − 1, k)xk−1\par
\noindent k=1 k=1\par
\noindent = x(1 − x)A′n−1 (x) + nxAn−1 (x),\par
}
\end{sourceblock}


\begin{explainblock}
这一段是原书论述链的一部分。读的时候把它当作桥梁：它通常负责把上一个定义、例子或定理，引到下一个计数方法。

\smallskip
\textbf{初学者检查点：}
\begin{itemize}[leftmargin=2em,itemsep=0.25em]
\item 找出这一段承接的上一个概念。
\item 找出它准备引出的下一个概念。
\item 把中间的关键等式或构造单独抄出来。
\end{itemize}

\smallskip
\textbf{本章读法提醒：} 第五章注意下降数、上升数、逆序数不是同一个统计量。
\end{explainblock}


\begin{sourceblock}{定理得证.}
{\small\sloppy
\noindent 定理得证.\par
}
\end{sourceblock}


\begin{explainblock}
定理段不要只看结论，要看它把哪两个计数方式连在一起。组合学证明最常见的技巧是：左边按一种标准数，右边按另一种标准数，然后说明两边数的是同一批对象。

\smallskip
\textbf{初学者检查点：}
\begin{itemize}[leftmargin=2em,itemsep=0.25em]
\item 先复述定理的假设，不要把适用范围扩大。
\item 找出证明中的基本计数原则：乘法、加法、双射、递推、容斥或生成函数。
\item 检查最后的等式化简是否和前面的对象解释一致。
\end{itemize}

\smallskip
\textbf{本章读法提醒：} 第五章注意下降数、上升数、逆序数不是同一个统计量。
\end{explainblock}


\begin{sourceblock}{定理 5.2.2}
{\small\sloppy
\noindent 定理 5.2.2\par
\noindent ∞\par
\noindent X An (x)\par
\noindent k n xk = . (5.4)\par
\noindent (1 − x)n+1\par
\noindent k=1\par
\noindent 证明: 下面我们用归纳法证明该定理. 当 n = 1 时, (5.4) 左边为:\par
\noindent ∞\par
\noindent X ′\par
\noindent kxk = x = ,\par
\noindent 1−x (1 − x)2\par
\noindent k=1\par
\noindent 再由 A1 (x) = x 知等式成立.\par
\noindent 假设式 (5.4) 对比 n(n ≥ 2) 小的正整数都成立, 则\par
\noindent ∞\par
\noindent X An−1 (x)\par
\noindent k n−1 xk = .\par
\noindent (1 − x)n\par
}
\end{sourceblock}


\begin{explainblock}
定理段不要只看结论，要看它把哪两个计数方式连在一起。组合学证明最常见的技巧是：左边按一种标准数，右边按另一种标准数，然后说明两边数的是同一批对象。

\smallskip
\textbf{初学者检查点：}
\begin{itemize}[leftmargin=2em,itemsep=0.25em]
\item 先复述定理的假设，不要把适用范围扩大。
\item 找出证明中的基本计数原则：乘法、加法、双射、递推、容斥或生成函数。
\item 检查最后的等式化简是否和前面的对象解释一致。
\end{itemize}

\smallskip
\textbf{本章读法提醒：} 第五章注意下降数、上升数、逆序数不是同一个统计量。
\end{explainblock}


\begin{sourceblock}{第 5 章续读}
{\small\sloppy
\noindent k=1\par
\noindent 对上式两边关于 x 求导后再乘上 x 可得\par
\noindent ∞\par
\noindent X x(1 − x)A′n−1 (x) + nxAn−1 (x)\par
\noindent k n xk = ,\par
\noindent (1 − x)n+1\par
\noindent k=1\par
\noindent 再由等式 (5.3) 知 (5.12) 成立.\par
}
\end{sourceblock}


\begin{explainblock}
这一段是原书论述链的一部分。读的时候把它当作桥梁：它通常负责把上一个定义、例子或定理，引到下一个计数方法。

\smallskip
\textbf{初学者检查点：}
\begin{itemize}[leftmargin=2em,itemsep=0.25em]
\item 找出这一段承接的上一个概念。
\item 找出它准备引出的下一个概念。
\item 把中间的关键等式或构造单独抄出来。
\end{itemize}

\smallskip
\textbf{本章读法提醒：} 第五章注意下降数、上升数、逆序数不是同一个统计量。
\end{explainblock}


\begin{sourceblock}{定理得证.}
{\small\sloppy
\noindent 定理得证.\par
\noindent 可用公式 (5.4) 快速计算 An (x).\par
}
\end{sourceblock}


\begin{explainblock}
定理段不要只看结论，要看它把哪两个计数方式连在一起。组合学证明最常见的技巧是：左边按一种标准数，右边按另一种标准数，然后说明两边数的是同一批对象。

\smallskip
\textbf{初学者检查点：}
\begin{itemize}[leftmargin=2em,itemsep=0.25em]
\item 先复述定理的假设，不要把适用范围扩大。
\item 找出证明中的基本计数原则：乘法、加法、双射、递推、容斥或生成函数。
\item 检查最后的等式化简是否和前面的对象解释一致。
\end{itemize}

\smallskip
\textbf{本章读法提醒：} 第五章注意下降数、上升数、逆序数不是同一个统计量。
\end{explainblock}


\begin{sourceblock}{例 5.2.2 求 A4 (x).}
{\small\sloppy
\noindent 例 5.2.2 求 A4 (x).\par
\noindent 解: 由(5.4) 知\par
\noindent ∞\par
\noindent X\par
\noindent A4 (x) =(1 − x) 5\par
\noindent m4 xm\par
\noindent m=0\par
\noindent ∞\par
\noindent X\par
\noindent =(1 − x) (1 − x)\par
\noindent m 4 xm\par
\noindent m=0\par
\noindent ∞\par
\noindent X\par
\noindent =(1 − x)4 (m4 − (m − 1)4 )xm ,\par
\noindent m=0\par
\noindent 5.2 Eulerian 多项式 97\par
\noindent 令 a1 (m) = m4 − (m − 1)4 , 则可依上面的计算得\par
}
\end{sourceblock}


\begin{explainblock}
例题是学习这本书最重要的部分。不要只看答案，要把每个限制条件变成一个计数动作：先安排谁、再插入谁、有没有重复、需不需要除以对称性。

\smallskip
\textbf{初学者检查点：}
\begin{itemize}[leftmargin=2em,itemsep=0.25em]
\item 先写出没有限制时的总数。
\item 逐条加入限制条件，观察总数怎样变化。
\item 最后检查有没有把同一种方案重复算了多次。
\end{itemize}

\smallskip
\textbf{本章读法提醒：} 第五章注意下降数、上升数、逆序数不是同一个统计量。
\end{explainblock}


\begin{sourceblock}{第 5 章续读}
{\small\sloppy
\noindent ∞\par
\noindent X\par
\noindent A4 (x) =(1 − x) 4\par
\noindent a1 (m)xm\par
\noindent m=0\par
\noindent X∞\par
\noindent =(1 − x)3 a2 (m)xm\par
\noindent m=0\par
\noindent =··· .\par
\noindent 我们将 ai (m) 的值绘制成下表: 所以\par
\noindent m4 0 1 24 34 44 54\par
\noindent a1 (m) 0 1 15 65 175 369\par
\noindent a2 (m) 0 1 14 50 110 194\par
\noindent a3 (m) 0 1 13 36 60 84\par
\noindent a4 (m) 0 1 12 23 24 24\par
\noindent a5 (m) 0 1 11 11 1 0\par
\noindent A4 (x) = x + 11x2 + 11x3 + x4 .\par
}
\end{sourceblock}


\begin{explainblock}
这一段是原书论述链的一部分。读的时候把它当作桥梁：它通常负责把上一个定义、例子或定理，引到下一个计数方法。

\smallskip
\textbf{初学者检查点：}
\begin{itemize}[leftmargin=2em,itemsep=0.25em]
\item 找出这一段承接的上一个概念。
\item 找出它准备引出的下一个概念。
\item 把中间的关键等式或构造单独抄出来。
\end{itemize}

\smallskip
\textbf{本章读法提醒：} 第五章注意下降数、上升数、逆序数不是同一个统计量。
\end{explainblock}


\begin{sourceblock}{定理 5.2.3 设 A(x, t) 为 Eulerian 多项式 An (x) 的指数型生成函数, 即}
{\small\sloppy
\noindent 定理 5.2.3 设 A(x, t) 为 Eulerian 多项式 An (x) 的指数型生成函数, 即\par
\noindent X tn\par
\noindent A(x, t) := An (x) ,\par
\noindent n!\par
\noindent n≥0\par
\noindent 则\par
\noindent 1−x\par
\noindent A(x, t) = . (5.5)\par
\noindent 证明: 由定理 5.2.2 知\par
\noindent ∞\par
\noindent X X An (x) tn ∞\par
\noindent An (x) tn 1\par
\noindent = +\par
\noindent (1 − x) n+1 n! 1−x (1 − x)n+1 n!\par
\noindent n=0 n=1\par
\noindent ∞\par
\noindent X ∞\par
\noindent ∞ X\par
}
\end{sourceblock}


\begin{explainblock}
定理段不要只看结论，要看它把哪两个计数方式连在一起。组合学证明最常见的技巧是：左边按一种标准数，右边按另一种标准数，然后说明两边数的是同一批对象。

\smallskip
\textbf{初学者检查点：}
\begin{itemize}[leftmargin=2em,itemsep=0.25em]
\item 先复述定理的假设，不要把适用范围扩大。
\item 找出证明中的基本计数原则：乘法、加法、双射、递推、容斥或生成函数。
\item 检查最后的等式化简是否和前面的对象解释一致。
\end{itemize}

\smallskip
\textbf{本章读法提醒：} 第五章注意下降数、上升数、逆序数不是同一个统计量。
\end{explainblock}


\begin{sourceblock}{第 5 章续读}
{\small\sloppy
\noindent X\par
\noindent k tn\par
\noindent = x + k n xk\par
\noindent n!\par
\noindent k=0 n=1 k=1\par
\noindent ∞\par
\noindent X ∞\par
\noindent X ∞\par
\noindent X\par
\noindent k k (kt)n\par
\noindent = x + x\par
\noindent n!\par
\noindent k=0 k=1 n=1\par
\noindent ∞\par
\noindent X ∞\par
\noindent X (kt)n\par
\noindent =1+ xk\par
\noindent n!\par
}
\end{sourceblock}


\begin{explainblock}
这一段是原书论述链的一部分。读的时候把它当作桥梁：它通常负责把上一个定义、例子或定理，引到下一个计数方法。

\smallskip
\textbf{初学者检查点：}
\begin{itemize}[leftmargin=2em,itemsep=0.25em]
\item 找出这一段承接的上一个概念。
\item 找出它准备引出的下一个概念。
\item 把中间的关键等式或构造单独抄出来。
\end{itemize}

\smallskip
\textbf{本章读法提醒：} 第五章注意下降数、上升数、逆序数不是同一个统计量。
\end{explainblock}


\begin{sourceblock}{第 5 章续读}
{\small\sloppy
\noindent k=1 n=0\par
\noindent ∞\par
\noindent X 1\par
\noindent = xk ekt = .\par
\noindent k=0\par
\noindent 两边同时乘上 1 − x, 然后将 t 换成 t(1 − x) 立得 (5.5).\par
}
\end{sourceblock}


\begin{explainblock}
这一段是原书论述链的一部分。读的时候把它当作桥梁：它通常负责把上一个定义、例子或定理，引到下一个计数方法。

\smallskip
\textbf{初学者检查点：}
\begin{itemize}[leftmargin=2em,itemsep=0.25em]
\item 找出这一段承接的上一个概念。
\item 找出它准备引出的下一个概念。
\item 把中间的关键等式或构造单独抄出来。
\end{itemize}

\smallskip
\textbf{本章读法提醒：} 第五章注意下降数、上升数、逆序数不是同一个统计量。
\end{explainblock}


\begin{sourceblock}{§ 5.3 含 Eulerian 数的组合恒等式}
{\small\sloppy
\noindent § 5.3. 含 Eulerian 数的组合恒等式\par
\noindent 这节我们将介绍一些与 Eulerian 数相关的恒等式.\par
}
\end{sourceblock}


\begin{explainblock}
这一节先不要急着背公式。先问三个问题：对象是什么，哪些限制条件不能丢，最后到底是在数所有对象、还是在数某个统计量的分布。把这三件事说清楚，后面的公式才会有位置。

\smallskip
\textbf{初学者检查点：}
\begin{itemize}[leftmargin=2em,itemsep=0.25em]
\item 把本节出现的新对象写成一句话定义。
\item 把每个公式左边和右边分别翻译成“在数什么”。
\item 找一个最小非平凡例子，手工列出所有对象。
\end{itemize}

\smallskip
\textbf{本章读法提醒：} 第五章注意下降数、上升数、逆序数不是同一个统计量。
\end{explainblock}


\begin{sourceblock}{定理 5.3.1 (Eulerian 数的显示表达式) 设 n 与 k 非负整数且 n ≥ k, 则}
{\small\sloppy
\noindent 定理 5.3.1 (Eulerian 数的显示表达式) 设 n 与 k 非负整数且 n ≥ k, 则\par
\noindent X\par
\noindent k\par
\noindent n+1\par
\noindent A(n, k) = (−1)i (k − i)n . (5.6)\par
\noindent i\par
\noindent i=0\par
\noindent 证明: 利用容斥原理证明. 考虑将集合 [n] 中的元素放到 k − 1 个竖线产生的\par
\noindent k − 1 个隔间中, 并要求隔间中元素是按递增顺序排列, 这样得到的所有带竖线排列\par
\noindent 组成的集合记为 S, 则显然 |S| = k n .\par
\noindent 我们定义集合 S 上的性质集合 P = \{P1 , P2 , . . . , Pk−1 \}, 其中 Pi (1 ≤ i ≤ k − 1)\par
\noindent 表示集合 S 中带竖线的排列第 i 条竖线不产生下降位, 并令 Ai 为集合 S 中所有具\par
\noindent 有性质 Pi 的所有带竖线排列组成的集合. 因此集合 S 中所有不具有集合 P 中任何\par
\noindent 一个性质的排列就是集合 [n] 中具有 k − 1 个下降位的排列. 故\par
\noindent A(n, k) = |Ā1 ∩ Ā2 ∩ · · · ∩ Āk−1 |.\par
\noindent 又由容斥原理可得\par
\noindent X X\par
\noindent |Ā1 ∩ Ā2 | ∩ · · · ∩ Āk−1 | = |S| − |Aj | + |Aj1 ∩ Aj2 | + · · ·\par
}
\end{sourceblock}


\begin{explainblock}
定理段不要只看结论，要看它把哪两个计数方式连在一起。组合学证明最常见的技巧是：左边按一种标准数，右边按另一种标准数，然后说明两边数的是同一批对象。

\smallskip
\textbf{初学者检查点：}
\begin{itemize}[leftmargin=2em,itemsep=0.25em]
\item 先复述定理的假设，不要把适用范围扩大。
\item 找出证明中的基本计数原则：乘法、加法、双射、递推、容斥或生成函数。
\item 检查最后的等式化简是否和前面的对象解释一致。
\end{itemize}

\smallskip
\textbf{本章读法提醒：} 第五章注意下降数、上升数、逆序数不是同一个统计量。
\end{explainblock}


\begin{sourceblock}{第 5 章续读}
{\small\sloppy
\noindent 1≤j≤k−1 1≤j1 <j2 ≤k−1\par
\noindent X\par
\noindent + (−1)i |Aj1 ∩ · · · ∩ Aji | + · · ·\par
\noindent 1≤j1 <···<ji ≤k−1\par
\noindent + (−1)k−1 |A1 ∩ A2 ∩ · · · ∩ Ak−1 |\par
\noindent X\par
\noindent k−1 X\par
\noindent = |S| − (−1)i |Aj1 ∩ Aj2 ∩ · · · ∩ Aji |.\par
\noindent i=1 1≤j1 <···<ji ≤k−1\par
\noindent (5.7)\par
\noindent 下面我们考虑计算后面的和式. 设 1 ≤ i ≤ k − 1 现将集合 [n] 元素放到由 k − 1 − i\par
\noindent 个竖线产生的隔间中且也要求它们在隔间的顺序是递增的, 方法数为 (k − i)n . 对于\par
\noindent 这样得到的任意排列, 我们再将剩余的 i 条竖线插入到其中, 要求不能插入到原有竖\par
\noindent 线位置的后面, 所以共有 n + 1 个位置可供选择, 则方法数有 n+1\par
\noindent i , 显然插入的这 i\par
\noindent 条竖线都不产生下降位. 由于插入位置的任意性可知:\par
\noindent X\par
\noindent n+1\par
}
\end{sourceblock}


\begin{explainblock}
这一段是原书论述链的一部分。读的时候把它当作桥梁：它通常负责把上一个定义、例子或定理，引到下一个计数方法。

\smallskip
\textbf{初学者检查点：}
\begin{itemize}[leftmargin=2em,itemsep=0.25em]
\item 找出这一段承接的上一个概念。
\item 找出它准备引出的下一个概念。
\item 把中间的关键等式或构造单独抄出来。
\end{itemize}

\smallskip
\textbf{本章读法提醒：} 第五章注意下降数、上升数、逆序数不是同一个统计量。
\end{explainblock}


\begin{sourceblock}{第 5 章续读}
{\small\sloppy
\noindent |Aj1 ∩ Aj2 ∩ · · · ∩ Aji | = (k − i)n ,\par
\noindent i\par
\noindent 1≤j1 <···<ji ≤k−1\par
\noindent 代入 (5.7) 立得 (5.6).\par
\noindent 5.3 含 Eulerian 数的组合恒等式 99\par
}
\end{sourceblock}


\begin{explainblock}
这一段是原书论述链的一部分。读的时候把它当作桥梁：它通常负责把上一个定义、例子或定理，引到下一个计数方法。

\smallskip
\textbf{初学者检查点：}
\begin{itemize}[leftmargin=2em,itemsep=0.25em]
\item 找出这一段承接的上一个概念。
\item 找出它准备引出的下一个概念。
\item 把中间的关键等式或构造单独抄出来。
\end{itemize}

\smallskip
\textbf{本章读法提醒：} 第五章注意下降数、上升数、逆序数不是同一个统计量。
\end{explainblock}


\begin{sourceblock}{注记 5.3.2 定理 5.3.1 也可由定理 5.2.2 给出的证明. 由定理 5.2.2 知}
{\small\sloppy
\noindent 注记 5.3.2 定理 5.3.1 也可由定理 5.2.2 给出的证明. 由定理 5.2.2 知\par
\noindent ∞\par
\noindent X\par
\noindent An (x) = (1 − x) n+1\par
\noindent j n xj , (5.8)\par
\noindent j=1\par
\noindent 将 (1 − x)n+1 由二项式定理展开代入得:\par
\noindent X\par
\noindent n+1\par
\noindent n+1\par
\noindent ∞\par
\noindent X\par
\noindent i\par
\noindent An (x) = (−x) j n xj\par
\noindent i\par
\noindent i=0 j=1\par
\noindent ∞\par
\noindent X k\par
}
\end{sourceblock}


\begin{explainblock}
注记通常是在补充术语、记号、历史或一个容易忽略的约定。它未必给新公式，但常常决定你后面会不会把符号读错。

\smallskip
\textbf{初学者检查点：}
\begin{itemize}[leftmargin=2em,itemsep=0.25em]
\item 记下新符号的读法。
\item 确认这个约定是否只在本章使用。
\item 把注记里的限制条件补到前面的定义旁边。
\end{itemize}

\smallskip
\textbf{本章读法提醒：} 第五章注意下降数、上升数、逆序数不是同一个统计量。
\end{explainblock}


\begin{sourceblock}{第 5 章续读}
{\small\sloppy
\noindent X\par
\noindent n+1\par
\noindent = xk\par
\noindent (−1)i (k − i)n ,\par
\noindent i\par
\noindent k=1 i=0\par
\noindent 比较 xk 的系数立得 (5.6).\par
}
\end{sourceblock}


\begin{explainblock}
这一段是原书论述链的一部分。读的时候把它当作桥梁：它通常负责把上一个定义、例子或定理，引到下一个计数方法。

\smallskip
\textbf{初学者检查点：}
\begin{itemize}[leftmargin=2em,itemsep=0.25em]
\item 找出这一段承接的上一个概念。
\item 找出它准备引出的下一个概念。
\item 把中间的关键等式或构造单独抄出来。
\end{itemize}

\smallskip
\textbf{本章读法提醒：} 第五章注意下降数、上升数、逆序数不是同一个统计量。
\end{explainblock}


\begin{sourceblock}{定理 5.3.3 设 n ≥ 1, 则对于任意实数 x,}
{\small\sloppy
\noindent 定理 5.3.3 设 n ≥ 1, 则对于任意实数 x,\par
\noindent X\par
\noindent n\par
\noindent x+n−k\par
\noindent xn = A(n, k) (5.9)\par
\noindent n\par
\noindent k=1\par
\noindent 证明: (方法一) 由对称性 A(n, k) = A(n, n − k + 1) 可得\par
\noindent X\par
\noindent n Xn Xn\par
\noindent x+n−k x+n−k x+j−1\par
\noindent A(n, k) = A(n, n−k+1) = A(n, j) .\par
\noindent n n n\par
\noindent k=1 k=1 j=1\par
\noindent 下面可用归纳法证明\par
\noindent X\par
\noindent n\par
\noindent n x+j−1\par
}
\end{sourceblock}


\begin{explainblock}
定理段不要只看结论，要看它把哪两个计数方式连在一起。组合学证明最常见的技巧是：左边按一种标准数，右边按另一种标准数，然后说明两边数的是同一批对象。

\smallskip
\textbf{初学者检查点：}
\begin{itemize}[leftmargin=2em,itemsep=0.25em]
\item 先复述定理的假设，不要把适用范围扩大。
\item 找出证明中的基本计数原则：乘法、加法、双射、递推、容斥或生成函数。
\item 检查最后的等式化简是否和前面的对象解释一致。
\end{itemize}

\smallskip
\textbf{本章读法提醒：} 第五章注意下降数、上升数、逆序数不是同一个统计量。
\end{explainblock}


\begin{sourceblock}{第 5 章续读}
{\small\sloppy
\noindent x = A(n, j) . (5.10)\par
\noindent n\par
\noindent j=1\par
\noindent 显然当 n = 1 时等式 (5.10) 左边与右边都等于 x. 假设对于小于 n(n ≥ 2) 的情形等\par
\noindent 式 (5.10) 均成立, 再由 A(n, k) 的递推关系式 (5.2) 可得\par
\noindent X\par
\noindent n\par
\noindent x+k−1\par
\noindent A(n, k)\par
\noindent n\par
\noindent k=1\par
\noindent n\par
\noindent X\par
\noindent x+k−1\par
\noindent = kA(n − 1, k) + (n − k + 1)A(n − 1, k − 1)\par
\noindent n\par
\noindent k=1\par
\noindent X\par
}
\end{sourceblock}


\begin{explainblock}
这一段是原书论述链的一部分。读的时候把它当作桥梁：它通常负责把上一个定义、例子或定理，引到下一个计数方法。

\smallskip
\textbf{初学者检查点：}
\begin{itemize}[leftmargin=2em,itemsep=0.25em]
\item 找出这一段承接的上一个概念。
\item 找出它准备引出的下一个概念。
\item 把中间的关键等式或构造单独抄出来。
\end{itemize}

\smallskip
\textbf{本章读法提醒：} 第五章注意下降数、上升数、逆序数不是同一个统计量。
\end{explainblock}


\begin{sourceblock}{第 5 章续读}
{\small\sloppy
\noindent n Xn\par
\noindent x+k−1 x+k−1\par
\noindent = kA(n − 1, k) + (n − k + 1)A(n − 1, k − 1)\par
\noindent n n\par
\noindent k=1 k=1\par
\noindent X\par
\noindent n Xn\par
\noindent x+k−1 x+k−1\par
\noindent = kA(n − 1, k) + (n − k + 1)A(n − 1, k − 1)\par
\noindent n n\par
\noindent k=1 k=2\par
\noindent X\par
\noindent n n−1\par
\noindent X\par
\noindent x+k−1 x+k\par
\noindent = kA(n − 1, k) + (n − k)A(n − 1, k)\par
\noindent n n\par
\noindent k=1 k=1\par
}
\end{sourceblock}


\begin{explainblock}
这一段是原书论述链的一部分。读的时候把它当作桥梁：它通常负责把上一个定义、例子或定理，引到下一个计数方法。

\smallskip
\textbf{初学者检查点：}
\begin{itemize}[leftmargin=2em,itemsep=0.25em]
\item 找出这一段承接的上一个概念。
\item 找出它准备引出的下一个概念。
\item 把中间的关键等式或构造单独抄出来。
\end{itemize}

\smallskip
\textbf{本章读法提醒：} 第五章注意下降数、上升数、逆序数不是同一个统计量。
\end{explainblock}


\begin{sourceblock}{第 5 章续读}
{\small\sloppy
\noindent X\par
\noindent n−1\par
\noindent x+k−1 x+k\par
\noindent = A(n − 1, k) k + (n − k)\par
\noindent n n\par
\noindent k=1\par
\noindent X\par
\noindent n−1\par
\noindent x+k−1\par
\noindent =x · A(n − 1, k)\par
\noindent n−1\par
\noindent k=1\par
\noindent =x · xn−1\par
\noindent =xn .\par
\noindent 所以 (5.10) 成立, 定理得证.\par
\noindent (方法二) 由代数学知识, 我们只需证明 (5.9) 对 x 是任意正整数时成立即可. 设\par
\noindent x = m, 其中 m 是任意一个正整数, 即证:\par
\noindent X\par
}
\end{sourceblock}


\begin{explainblock}
这一段是原书论述链的一部分。读的时候把它当作桥梁：它通常负责把上一个定义、例子或定理，引到下一个计数方法。

\smallskip
\textbf{初学者检查点：}
\begin{itemize}[leftmargin=2em,itemsep=0.25em]
\item 找出这一段承接的上一个概念。
\item 找出它准备引出的下一个概念。
\item 把中间的关键等式或构造单独抄出来。
\end{itemize}

\smallskip
\textbf{本章读法提醒：} 第五章注意下降数、上升数、逆序数不是同一个统计量。
\end{explainblock}


\begin{sourceblock}{第 5 章续读}
{\small\sloppy
\noindent n\par
\noindent n m+n−k\par
\noindent m = A(n, k) . (5.11)\par
\noindent n\par
\noindent k=1\par
\noindent 显然等式 (5.11) 的左边计数了长为 n 的可重排列 a1 a2 · · · an 的个数, 其中 1 ≤ ai ≤\par
\noindent m(1 ≤ i ≤ n); 该排列也可写成两行的形式:\par
\noindent !\par
\noindent A= .\par
\noindent a1 a2 · · · an\par
\noindent 将 A 中的每列按 ai 从小到大的顺序重新排列, 且当 ai = aj 时, 按 i, j 从小到大顺\par
\noindent 序排列; 这样重新排列后记为 A′ : 阵记为:\par
\noindent !\par
\noindent π\par
\noindent A′ =\par
\noindent ω\par
\noindent 其中 π = i1 i2 · · · in , 它是集合 [n] 上的排列, ω = ai1 ai2 · · · ain 它是集合 [m] 上的 n\par
\noindent 长可重排列, 且满足当 r 是排列 i1 i2 · · · in 中的下降位, 我们有 air < air+1 . 显然 A\par
}
\end{sourceblock}


\begin{explainblock}
这一段是原书论述链的一部分。读的时候把它当作桥梁：它通常负责把上一个定义、例子或定理，引到下一个计数方法。

\smallskip
\textbf{初学者检查点：}
\begin{itemize}[leftmargin=2em,itemsep=0.25em]
\item 找出这一段承接的上一个概念。
\item 找出它准备引出的下一个概念。
\item 把中间的关键等式或构造单独抄出来。
\end{itemize}

\smallskip
\textbf{本章读法提醒：} 第五章注意下降数、上升数、逆序数不是同一个统计量。
\end{explainblock}


\begin{sourceblock}{第 5 章续读}
{\small\sloppy
\noindent 与 A′ 一一对应.\par
\noindent 比如: !\par
\noindent 与 !\par
\noindent 对应.\par
\noindent 下面我们求形如 A′ 的二元组的个数. 对于 A′ 的第一行元素 π, 我们知它是集合\par
\noindent [n] 上的一个排列, 则当这个排列具有 k − 1 个下降位时, 它有 A(n, k) 种情形; 对于\par
\noindent 5.3 含 Eulerian 数的组合恒等式 101\par
\noindent A′ 中第二行元素 ω, 它满足\par
\noindent 1 ≤ ai1 ≤ ai2 ≤ · · · ≤ ain ≤ m,\par
\noindent 且当 r 是排列 i1 i2 · · · in 的下降位时, air 与 air+1 之间的小于等于号变为严格小于\par
\noindent 号. 设 i1 i2 · · · in 中的 k − 1 个下降位分别为 j1 , j2 , . . . jk−1 , 则\par
\noindent 1 ≤ ai1 ≤ · · · ≤ aij1 < aij1 +1 ≤ · · · ≤ aij2 < aij2 +1 ≤ · · · ≤ ain ≤ m.\par
\noindent 当 ℓ ≤ j1 , 令 biℓ = aiℓ + ℓ − 1; 则\par
\noindent 1 ≤ bi1 < bi2 < · · · < bij1 ;\par
\noindent 一般地, 我们可设 1 < r ≤ k − 1, 当 jr−1 < ℓ ≤ jr 时, 令 biℓ = aiℓ + ℓ − r, 所以\par
\noindent bijr−1 +1 < bijr−1 +2 < · · · < bijr .\par
\noindent 当 jk−1 < ℓ ≤ n, 令 biℓ = aiℓ + ℓ − k, 则\par
\noindent bijk−1 +1 < bijk−1 +1 < · · · < bin ≤ m + n − k.\par
}
\end{sourceblock}


\begin{explainblock}
这一段是原书论述链的一部分。读的时候把它当作桥梁：它通常负责把上一个定义、例子或定理，引到下一个计数方法。

\smallskip
\textbf{初学者检查点：}
\begin{itemize}[leftmargin=2em,itemsep=0.25em]
\item 找出这一段承接的上一个概念。
\item 找出它准备引出的下一个概念。
\item 把中间的关键等式或构造单独抄出来。
\end{itemize}

\smallskip
\textbf{本章读法提醒：} 第五章注意下降数、上升数、逆序数不是同一个统计量。
\end{explainblock}


\begin{sourceblock}{第 5 章续读}
{\small\sloppy
\noindent 设 1 < r ≤ k − 1, 由上面的变换可知\par
\noindent bijr−1 = aijr−1 + jr−1 − (r − 1) < bijr−1 +1 = aijr−1 +1 + jr−1 + 1 − r.\par
\noindent 且有\par
\noindent bijk−1 = aijk−1 + jk−1 − (k − 1) < bijk−1 +1 = aijk−1 +1 + jk−1 + 1 − k.\par
\noindent 故我们有:\par
\noindent 1 ≤ bi1 < bi2 < · · · < bin ≤ m + n − k. (5.12)\par
\noindent 显 然 ai1 ai2 · · · ain 与 排 列 bi1 bi2 · · · bin 一 一 对 应, 而 满 足 (5.12) 的 序 列 显 然 有\par
\noindent m+n−k\par
\noindent 个. 所以当 A′ 第一行确定的排列 i1 i2 · · · in 有 k − 1 个下降位时, A′\par
\noindent n\par
\noindent 第二行确定的 ai1 ai2 · · · ain 有 m+n−k n 种情形; 由加法原理和乘法原理可知形如 A′\par
\noindent 的二元组有:\par
\noindent X\par
\noindent n\par
\noindent m+n−k\par
\noindent A(n, k) .\par
\noindent n\par
\noindent k=1\par
}
\end{sourceblock}


\begin{explainblock}
这一段是原书论述链的一部分。读的时候把它当作桥梁：它通常负责把上一个定义、例子或定理，引到下一个计数方法。

\smallskip
\textbf{初学者检查点：}
\begin{itemize}[leftmargin=2em,itemsep=0.25em]
\item 找出这一段承接的上一个概念。
\item 找出它准备引出的下一个概念。
\item 把中间的关键等式或构造单独抄出来。
\end{itemize}

\smallskip
\textbf{本章读法提醒：} 第五章注意下降数、上升数、逆序数不是同一个统计量。
\end{explainblock}


\begin{sourceblock}{第 5 章续读}
{\small\sloppy
\noindent 所以等式 (5.11) 成立; 若 x 不是一个正整数, 由于等式 (5.9) 两边都可以看作是\par
\noindent 关于变量 x 的多项式, 而它们在无穷多个数值上取值相同, 所以它们本身必须是相等\par
\noindent 的.\par
\noindent 下面我们应用等式 (5.10) 给出前 m 个自然数正数次幂求和的一个公式.\par
}
\end{sourceblock}


\begin{explainblock}
这一段是原书论述链的一部分。读的时候把它当作桥梁：它通常负责把上一个定义、例子或定理，引到下一个计数方法。

\smallskip
\textbf{初学者检查点：}
\begin{itemize}[leftmargin=2em,itemsep=0.25em]
\item 找出这一段承接的上一个概念。
\item 找出它准备引出的下一个概念。
\item 把中间的关键等式或构造单独抄出来。
\end{itemize}

\smallskip
\textbf{本章读法提醒：} 第五章注意下降数、上升数、逆序数不是同一个统计量。
\end{explainblock}


\begin{sourceblock}{定理 5.3.4}
{\small\sloppy
\noindent 定理 5.3.4\par
\noindent X\par
\noindent n X\par
\noindent m\par
\noindent m k+n\par
\noindent i = A(m, k) . (5.13)\par
\noindent m+1\par
\noindent i=1 k=1\par
\noindent 证明: 由 (5.10) 知当 m ≥ 1 时,\par
\noindent X\par
\noindent m\par
\noindent m i+k−1\par
\noindent i = A(m, k) ,\par
\noindent m\par
\noindent k=1\par
\noindent 求和得\par
\noindent X\par
\noindent n X\par
}
\end{sourceblock}


\begin{explainblock}
定理段不要只看结论，要看它把哪两个计数方式连在一起。组合学证明最常见的技巧是：左边按一种标准数，右边按另一种标准数，然后说明两边数的是同一批对象。

\smallskip
\textbf{初学者检查点：}
\begin{itemize}[leftmargin=2em,itemsep=0.25em]
\item 先复述定理的假设，不要把适用范围扩大。
\item 找出证明中的基本计数原则：乘法、加法、双射、递推、容斥或生成函数。
\item 检查最后的等式化简是否和前面的对象解释一致。
\end{itemize}

\smallskip
\textbf{本章读法提醒：} 第五章注意下降数、上升数、逆序数不是同一个统计量。
\end{explainblock}


\begin{sourceblock}{第 5 章续读}
{\small\sloppy
\noindent n X\par
\noindent m\par
\noindent i+k−1\par
\noindent im = A(m, k)\par
\noindent m\par
\noindent i=1 i=1 k=1\par
\noindent Xm Xn\par
\noindent i+k−1\par
\noindent = A(m, k)\par
\noindent m\par
\noindent k=1 i=1\par
\noindent Xm Xn\par
\noindent i+k i+k−1\par
\noindent = A(m, k) −\par
\noindent m+1 m+1\par
\noindent k=1 i=1\par
\noindent Xm\par
\noindent n+k k\par
}
\end{sourceblock}


\begin{explainblock}
这一段是原书论述链的一部分。读的时候把它当作桥梁：它通常负责把上一个定义、例子或定理，引到下一个计数方法。

\smallskip
\textbf{初学者检查点：}
\begin{itemize}[leftmargin=2em,itemsep=0.25em]
\item 找出这一段承接的上一个概念。
\item 找出它准备引出的下一个概念。
\item 把中间的关键等式或构造单独抄出来。
\end{itemize}

\smallskip
\textbf{本章读法提醒：} 第五章注意下降数、上升数、逆序数不是同一个统计量。
\end{explainblock}


\begin{sourceblock}{第 5 章续读}
{\small\sloppy
\noindent = A(m, k) −\par
\noindent m+1 m+1\par
\noindent k=1\par
\noindent Xm\par
\noindent n+k\par
\noindent = A(m, k) .\par
\noindent m+1\par
\noindent k=1\par
}
\end{sourceblock}


\begin{explainblock}
这一段是原书论述链的一部分。读的时候把它当作桥梁：它通常负责把上一个定义、例子或定理，引到下一个计数方法。

\smallskip
\textbf{初学者检查点：}
\begin{itemize}[leftmargin=2em,itemsep=0.25em]
\item 找出这一段承接的上一个概念。
\item 找出它准备引出的下一个概念。
\item 把中间的关键等式或构造单独抄出来。
\end{itemize}

\smallskip
\textbf{本章读法提醒：} 第五章注意下降数、上升数、逆序数不是同一个统计量。
\end{explainblock}


\begin{sourceblock}{例 5.3.1 求 13 + 23 + 33 + · · · + n3 .}
{\small\sloppy
\noindent 例 5.3.1 求 13 + 23 + 33 + · · · + n3 .\par
\noindent 解: 由定理 5.3.4 可得\par
\noindent X\par
\noindent k+n\par
\noindent 1 + 2 + 3 + ··· + n =\par
\noindent A(3, k)\par
\noindent k=1\par
\noindent n+1 n+2 n+3\par
\noindent = +4 +\par
\noindent n2 (n + 1)2\par
\noindent = .\par
\noindent 下面我们用组合方法建立了 Eulerian 数 A(n, k) 与第二类 Stirling 数 S(n, k) 之\par
\noindent 间的联系.\par
}
\end{sourceblock}


\begin{explainblock}
例题是学习这本书最重要的部分。不要只看答案，要把每个限制条件变成一个计数动作：先安排谁、再插入谁、有没有重复、需不需要除以对称性。

\smallskip
\textbf{初学者检查点：}
\begin{itemize}[leftmargin=2em,itemsep=0.25em]
\item 先写出没有限制时的总数。
\item 逐条加入限制条件，观察总数怎样变化。
\item 最后检查有没有把同一种方案重复算了多次。
\end{itemize}

\smallskip
\textbf{本章读法提醒：} 第五章注意下降数、上升数、逆序数不是同一个统计量。
\end{explainblock}


\begin{sourceblock}{定理 5.3.5 设 n, r 为正整数, 则}
{\small\sloppy
\noindent 定理 5.3.5 设 n, r 为正整数, 则\par
\noindent X\par
\noindent r\par
\noindent n−k\par
\noindent r!S(n, r) = A(n, k) . (5.14)\par
\noindent r−k\par
\noindent k=1\par
\noindent 证明: 等式 (5.14) 左边计数了非空集合 [n] 的有序 r 块划分的个数. 不妨\par
\noindent 设 B = B1 |B2 | · · · |Br 是集合 [n] 的任意一个 r 有序划分. 对任意 1 ≤ i ≤ r, 将\par
\noindent 5.3 含 Eulerian 数的组合恒等式 103\par
\noindent Br 中的元素按递增的顺序排列后 (单个元素也认为是递增的) 仍用 Bi 表示. 则\par
\noindent 排列 B1 B2 · · · Br 是集合 [n] 上的一个排列, 为了方便叙述我们用带竖线的排列\par
\noindent B1 |B2 | · · · |Br 来表示它, 并记为 π. 显然这种带竖线的排列也可按照如下的方式构\par
\noindent 造: 先给定一个具有 k − 1(1 ≤ k ≤ r) 个下降位的排列 π, 然后将 r − 1 条竖线插入\par
\noindent 到 π 中, 要求 π 中 k − 1 个下降位的位置都要插入 1 条竖线, 其余 r − k 条竖线插\par
\noindent 入到其它 n − k 位置中的 r − k 个位置. 故这种带竖线的排列共有 A(n, k) n−k\par
\noindent r−k 个,\par
\noindent 所以\par
}
\end{sourceblock}


\begin{explainblock}
定理段不要只看结论，要看它把哪两个计数方式连在一起。组合学证明最常见的技巧是：左边按一种标准数，右边按另一种标准数，然后说明两边数的是同一批对象。

\smallskip
\textbf{初学者检查点：}
\begin{itemize}[leftmargin=2em,itemsep=0.25em]
\item 先复述定理的假设，不要把适用范围扩大。
\item 找出证明中的基本计数原则：乘法、加法、双射、递推、容斥或生成函数。
\item 检查最后的等式化简是否和前面的对象解释一致。
\end{itemize}

\smallskip
\textbf{本章读法提醒：} 第五章注意下降数、上升数、逆序数不是同一个统计量。
\end{explainblock}


\begin{sourceblock}{第 5 章续读}
{\small\sloppy
\noindent X\par
\noindent r\par
\noindent n−k\par
\noindent r!S(n, r) = A(n, k) .\par
\noindent r−k\par
\noindent k=1\par
}
\end{sourceblock}


\begin{explainblock}
这一段是原书论述链的一部分。读的时候把它当作桥梁：它通常负责把上一个定义、例子或定理，引到下一个计数方法。

\smallskip
\textbf{初学者检查点：}
\begin{itemize}[leftmargin=2em,itemsep=0.25em]
\item 找出这一段承接的上一个概念。
\item 找出它准备引出的下一个概念。
\item 把中间的关键等式或构造单独抄出来。
\end{itemize}

\smallskip
\textbf{本章读法提醒：} 第五章注意下降数、上升数、逆序数不是同一个统计量。
\end{explainblock}


\begin{sourceblock}{定理 5.3.6 设 n, k 为任意的正整数且 n ≥ k, 则}
{\small\sloppy
\noindent 定理 5.3.6 设 n, k 为任意的正整数且 n ≥ k, 则\par
\noindent X\par
\noindent k\par
\noindent n−r\par
\noindent A(n, k) = S(n, r)r! (−1)k−r . (5.15)\par
\noindent k−r\par
\noindent r=1\par
\noindent 证明: 由定理 5.3.5 可把\par
\noindent X\par
\noindent r\par
\noindent n−k\par
\noindent r!S(n, r) = A(n, k) ,\par
\noindent r−k\par
\noindent k=1\par
\noindent 代入 (5.15) 的右边可得\par
\noindent X\par
\noindent k X\par
\noindent k X\par
}
\end{sourceblock}


\begin{explainblock}
定理段不要只看结论，要看它把哪两个计数方式连在一起。组合学证明最常见的技巧是：左边按一种标准数，右边按另一种标准数，然后说明两边数的是同一批对象。

\smallskip
\textbf{初学者检查点：}
\begin{itemize}[leftmargin=2em,itemsep=0.25em]
\item 先复述定理的假设，不要把适用范围扩大。
\item 找出证明中的基本计数原则：乘法、加法、双射、递推、容斥或生成函数。
\item 检查最后的等式化简是否和前面的对象解释一致。
\end{itemize}

\smallskip
\textbf{本章读法提醒：} 第五章注意下降数、上升数、逆序数不是同一个统计量。
\end{explainblock}


\begin{sourceblock}{第 5 章续读}
{\small\sloppy
\noindent r\par
\noindent n−r k−r k−r n − r n−i\par
\noindent S(n, r)r! (−1) = (−1) A(n, i)\par
\noindent k−r k−r r−i\par
\noindent r=1 r=1 i=1\par
\noindent X\par
\noindent k X\par
\noindent k\par
\noindent k−r n − i n−r\par
\noindent = A(n, i) (−1) .\par
\noindent r−i k−r\par
\noindent i=1 r=i\par
\noindent 注意到当 1 ≤ i < k 时,\par
\noindent k\par
\noindent X k\par
\noindent X\par
\noindent n−i n−r n−i k−n−1 k−i−1\par
\noindent (−1)k−r = = = 0,\par
}
\end{sourceblock}


\begin{explainblock}
这一段是原书论述链的一部分。读的时候把它当作桥梁：它通常负责把上一个定义、例子或定理，引到下一个计数方法。

\smallskip
\textbf{初学者检查点：}
\begin{itemize}[leftmargin=2em,itemsep=0.25em]
\item 找出这一段承接的上一个概念。
\item 找出它准备引出的下一个概念。
\item 把中间的关键等式或构造单独抄出来。
\end{itemize}

\smallskip
\textbf{本章读法提醒：} 第五章注意下降数、上升数、逆序数不是同一个统计量。
\end{explainblock}


\begin{sourceblock}{第 5 章续读}
{\small\sloppy
\noindent r−i k−r r−i k−r k−i\par
\noindent r=i r=i\par
\noindent 其中最后第二个等号是由 Chu-Vandermonde 恒等式 (见习题 1.16) 得到的. 所以\par
\noindent X\par
\noindent k X\par
\noindent k X\par
\noindent k\par
\noindent n−r k−r k−r n − i n−r\par
\noindent S(n, r)r! (−1) = A(n, i) (−1)\par
\noindent k−r r−i k−r\par
\noindent r=1 i=1 r=i\par
\noindent = A(n, k).\par
}
\end{sourceblock}


\begin{explainblock}
这一段是原书论述链的一部分。读的时候把它当作桥梁：它通常负责把上一个定义、例子或定理，引到下一个计数方法。

\smallskip
\textbf{初学者检查点：}
\begin{itemize}[leftmargin=2em,itemsep=0.25em]
\item 找出这一段承接的上一个概念。
\item 找出它准备引出的下一个概念。
\item 把中间的关键等式或构造单独抄出来。
\end{itemize}

\smallskip
\textbf{本章读法提醒：} 第五章注意下降数、上升数、逆序数不是同一个统计量。
\end{explainblock}


\begin{sourceblock}{§ 5.4 交错排列与 Euler 数}
{\small\sloppy
\noindent § 5.4. 交错排列与 Euler 数\par
\noindent 对于排列我们已经建立很多概念, 如错排, 下降位, 上升位等等, 对应的有错排数\par
\noindent Dn , Eulerian 数, Eulerian 多项式等等. 这节我们将讨论下降位与上升位交替出现的\par
\noindent 排列.\par
}
\end{sourceblock}


\begin{explainblock}
这一节先不要急着背公式。先问三个问题：对象是什么，哪些限制条件不能丢，最后到底是在数所有对象、还是在数某个统计量的分布。把这三件事说清楚，后面的公式才会有位置。

\smallskip
\textbf{初学者检查点：}
\begin{itemize}[leftmargin=2em,itemsep=0.25em]
\item 把本节出现的新对象写成一句话定义。
\item 把每个公式左边和右边分别翻译成“在数什么”。
\item 找一个最小非平凡例子，手工列出所有对象。
\end{itemize}

\smallskip
\textbf{本章读法提醒：} 第五章注意下降数、上升数、逆序数不是同一个统计量。
\end{explainblock}


\begin{sourceblock}{例 5.4.1 集合 \{1, 2, 3, 4\} 所有排列中, 上升位和下降位交替出现的排列有:}
{\small\sloppy
\noindent 例 5.4.1 集合 \{1, 2, 3, 4\} 所有排列中, 上升位和下降位交替出现的排列有:\par
\noindent 1324, 2413, 1423, 3412, 2314,\par
\noindent 4231, 3142, 3241, 2143, 4132.\par
\noindent 观察发现, 上面五个排列是先出现上升位再出现下降位, 后面五个则先出现下降位再\par
\noindent 出现上升位.\par
}
\end{sourceblock}


\begin{explainblock}
例题是学习这本书最重要的部分。不要只看答案，要把每个限制条件变成一个计数动作：先安排谁、再插入谁、有没有重复、需不需要除以对称性。

\smallskip
\textbf{初学者检查点：}
\begin{itemize}[leftmargin=2em,itemsep=0.25em]
\item 先写出没有限制时的总数。
\item 逐条加入限制条件，观察总数怎样变化。
\item 最后检查有没有把同一种方案重复算了多次。
\end{itemize}

\smallskip
\textbf{本章读法提醒：} 第五章注意下降数、上升数、逆序数不是同一个统计量。
\end{explainblock}


\begin{sourceblock}{定义 5.4.1 (交错排列) 设集合 [n] 的所有排列组成的集合为 Sn . 设排列 π =}
{\small\sloppy
\noindent 定义 5.4.1 (交错排列) 设集合 [n] 的所有排列组成的集合为 Sn . 设排列 π =\par
\noindent π1 π2 · · · πn ∈ Sn , 若 π 满 足 π1 > π2 < π3 > π4 < · · · , 则 称 π 是 交 错 的; 若\par
\noindent π1 < π2 > π3 < π4 > · · · , 则称 π 是反向交错的. Sn 中所有交错排列的个数记为\par
\noindent En , 称为 Euler 数, 规定 E0 = 1.\par
\noindent 由上例可知 E4 = 5. 我们这里列出了前几项 Euler 数的值:\par
\noindent n 1 2 3 4 5 6 7 8\par
\noindent En 1 1 2 5 16 61 272 1385\par
\noindent 性质 5.4.2 相同长度的交错排列和反向交错排列的个数相等.\par
\noindent 证明: 设 π = π1 π2 · · · πn ∈ Sn 为交错排列, 令\par
\noindent φ(πi ) = n + 1 − πi .\par
\noindent 则 π ′ = φ(π1 )φ(π2 ) · · · φ(πn ) ∈ Sn 为反向交错排列, 显然这个变换是集合 Sn 上到自\par
\noindent 身的对合映射. 所以长度为 n 的交错排列与反向交错排列之间一一对应, 它们的个\par
\noindent 数因此相等.\par
}
\end{sourceblock}


\begin{explainblock}
定义段的任务是定边界。这里要特别注意参数范围、是否允许重复、是否考虑顺序、对象有没有标签。后面的每一道例题和证明，本质上都在反复使用这些边界。

\smallskip
\textbf{初学者检查点：}
\begin{itemize}[leftmargin=2em,itemsep=0.25em]
\item 圈出被定义的对象名称。
\item 圈出参数范围，例如 n、k 是否要求正整数。
\item 判断它是有序对象、无序对象、带标签对象还是结构对象。
\end{itemize}

\smallskip
\textbf{本章读法提醒：} 第五章注意下降数、上升数、逆序数不是同一个统计量。
\end{explainblock}


\begin{sourceblock}{定理 5.4.3 Euler 数满足下列递推关系:}
{\small\sloppy
\noindent 定理 5.4.3 Euler 数满足下列递推关系:\par
\noindent n\par
\noindent X n\par
\noindent 2En+1 = Ei En−i , n ≥ 1. (5.16)\par
\noindent i\par
\noindent i=0\par
\noindent 证明: 我们考虑 Sn+1 中下降位与上升位交替出现的所有排列记为集合 Πn+1 ,\par
\noindent 由性质 5.4.2 知 |Πn+1 | = 2En+1 . 设 π ∈ Πn+1 , πi+1 = n + 1(0 ≤ i ≤ n). 记\par
\noindent u = π1 π2 · · · πi , v = πi+2 πi+3 · · · πn+1 . 则当 i 是偶数时, 知 u 是集合 [n] 的 i 元子集\par
\noindent 的一个交错排列, 而 v 是剩下元素的一个反向交错排列; 当 i 是奇数时, 知 u 是集合\par
\noindent [n] 的 i 元子集的一个反向交错排列, 而 v 依旧是剩下元素的一个反向交错排列. 对\par
\noindent i 从 0 到 n 进行求和得:\par
\noindent n\par
\noindent X n\par
\noindent 2En+1 = Ei En−i .\par
\noindent i\par
\noindent i=0\par
\noindent 5.4 交错排列与 Euler 数 105\par
}
\end{sourceblock}


\begin{explainblock}
定理段不要只看结论，要看它把哪两个计数方式连在一起。组合学证明最常见的技巧是：左边按一种标准数，右边按另一种标准数，然后说明两边数的是同一批对象。

\smallskip
\textbf{初学者检查点：}
\begin{itemize}[leftmargin=2em,itemsep=0.25em]
\item 先复述定理的假设，不要把适用范围扩大。
\item 找出证明中的基本计数原则：乘法、加法、双射、递推、容斥或生成函数。
\item 检查最后的等式化简是否和前面的对象解释一致。
\end{itemize}

\smallskip
\textbf{本章读法提醒：} 第五章注意下降数、上升数、逆序数不是同一个统计量。
\end{explainblock}


\begin{sourceblock}{第 5 章续读}
{\small\sloppy
\noindent 下面我们考虑 Euler 数的生成函数.\par
}
\end{sourceblock}


\begin{explainblock}
这一段是原书论述链的一部分。读的时候把它当作桥梁：它通常负责把上一个定义、例子或定理，引到下一个计数方法。

\smallskip
\textbf{初学者检查点：}
\begin{itemize}[leftmargin=2em,itemsep=0.25em]
\item 找出这一段承接的上一个概念。
\item 找出它准备引出的下一个概念。
\item 把中间的关键等式或构造单独抄出来。
\end{itemize}

\smallskip
\textbf{本章读法提醒：} 第五章注意下降数、上升数、逆序数不是同一个统计量。
\end{explainblock}


\begin{sourceblock}{定理 5.4.4 设 n 为非负整数, En 为 Euler 数, 则}
{\small\sloppy
\noindent 定理 5.4.4 设 n 为非负整数, En 为 Euler 数, 则\par
\noindent ∞\par
\noindent X xn\par
\noindent En = sec x + tan x. (5.17)\par
\noindent n!\par
\noindent n=0\par
\noindent 证明: 不妨设\par
\noindent ∞\par
\noindent X xn\par
\noindent y= En ,\par
\noindent n!\par
\noindent n=0\par
\noindent 则 y(0) = 1 且\par
\noindent ∞\par
\noindent X\par
\noindent ′ xn\par
\noindent y = En+1 ,\par
\noindent n!\par
}
\end{sourceblock}


\begin{explainblock}
定理段不要只看结论，要看它把哪两个计数方式连在一起。组合学证明最常见的技巧是：左边按一种标准数，右边按另一种标准数，然后说明两边数的是同一批对象。

\smallskip
\textbf{初学者检查点：}
\begin{itemize}[leftmargin=2em,itemsep=0.25em]
\item 先复述定理的假设，不要把适用范围扩大。
\item 找出证明中的基本计数原则：乘法、加法、双射、递推、容斥或生成函数。
\item 检查最后的等式化简是否和前面的对象解释一致。
\end{itemize}

\smallskip
\textbf{本章读法提醒：} 第五章注意下降数、上升数、逆序数不是同一个统计量。
\end{explainblock}


\begin{sourceblock}{第 5 章续读}
{\small\sloppy
\noindent n=0\par
\noindent 所以\par
\noindent ∞\par
\noindent X X ∞\par
\noindent ′ xn xn\par
\noindent 2y = 2En+1 =2+ 2En+1 . (5.18)\par
\noindent n! n!\par
\noindent n=0 n=1\par
\noindent 将 (5.16) 代入 (5.18) 可得:\par
\noindent n\par
\noindent ∞ X\par
\noindent X n xn\par
\noindent 2y ′ = 2 + Ek En−k\par
\noindent k n!\par
\noindent n=1 k=0\par
\noindent ∞\par
\noindent X ∞ n\par
\noindent n xn X X n xn\par
}
\end{sourceblock}


\begin{explainblock}
这一段是原书论述链的一部分。读的时候把它当作桥梁：它通常负责把上一个定义、例子或定理，引到下一个计数方法。

\smallskip
\textbf{初学者检查点：}
\begin{itemize}[leftmargin=2em,itemsep=0.25em]
\item 找出这一段承接的上一个概念。
\item 找出它准备引出的下一个概念。
\item 把中间的关键等式或构造单独抄出来。
\end{itemize}

\smallskip
\textbf{本章读法提醒：} 第五章注意下降数、上升数、逆序数不是同一个统计量。
\end{explainblock}


\begin{sourceblock}{第 5 章续读}
{\small\sloppy
\noindent =2+ E0 En + Ek En−k\par
\noindent n=1 n=1 k=1\par
\noindent ∞\par
\noindent X ∞ ∞\par
\noindent xn X xk X xn−k\par
\noindent =1+ En + Ek En−k\par
\noindent n! k! (n − k)!\par
\noindent n=0 k=1 n=k\par
\noindent = 1 + y + (y − 1)y = y 2 + 1.\par
\noindent 因此得到下面的微分方程:\par
\noindent 2y ′ = y 2 + 1, y(0) = 1.\par
\noindent 解之可得\par
\noindent dy\par
\noindent y +1\par
\noindent 对上式两边积分得\par
\noindent 2 arctan y = x + C.\par
\noindent 由 y(0) = 1 知 C = π2 , 所以有\par
\noindent x + π/2\par
}
\end{sourceblock}


\begin{explainblock}
这一段是原书论述链的一部分。读的时候把它当作桥梁：它通常负责把上一个定义、例子或定理，引到下一个计数方法。

\smallskip
\textbf{初学者检查点：}
\begin{itemize}[leftmargin=2em,itemsep=0.25em]
\item 找出这一段承接的上一个概念。
\item 找出它准备引出的下一个概念。
\item 把中间的关键等式或构造单独抄出来。
\end{itemize}

\smallskip
\textbf{本章读法提醒：} 第五章注意下降数、上升数、逆序数不是同一个统计量。
\end{explainblock}


\begin{sourceblock}{第 5 章续读}
{\small\sloppy
\noindent arctan y = .\par
\noindent 故而\par
\noindent x + π/2\par
\noindent y = tan\par
\noindent sin x+π/2\par
\noindent =\par
\noindent cos x+π/2\par
\noindent =\par
\noindent =\par
\noindent sin(x + π2 )\par
\noindent =\par
\noindent cos x\par
\noindent = sec x + tan x.\par
\noindent 在 (5.17) 将 x 替换为 −x 时, 得到\par
\noindent ∞\par
\noindent X xn\par
\noindent (−1)n En = sec x − tan x. (5.19)\par
\noindent n!\par
}
\end{sourceblock}


\begin{explainblock}
这一段是原书论述链的一部分。读的时候把它当作桥梁：它通常负责把上一个定义、例子或定理，引到下一个计数方法。

\smallskip
\textbf{初学者检查点：}
\begin{itemize}[leftmargin=2em,itemsep=0.25em]
\item 找出这一段承接的上一个概念。
\item 找出它准备引出的下一个概念。
\item 把中间的关键等式或构造单独抄出来。
\end{itemize}

\smallskip
\textbf{本章读法提醒：} 第五章注意下降数、上升数、逆序数不是同一个统计量。
\end{explainblock}


\begin{sourceblock}{第 5 章续读}
{\small\sloppy
\noindent n=0\par
\noindent 通过将 (5.17) 与 (5.17) 左右两边分别求和或分别求差即可得到下面的推论.\par
}
\end{sourceblock}


\begin{explainblock}
这一段是原书论述链的一部分。读的时候把它当作桥梁：它通常负责把上一个定义、例子或定理，引到下一个计数方法。

\smallskip
\textbf{初学者检查点：}
\begin{itemize}[leftmargin=2em,itemsep=0.25em]
\item 找出这一段承接的上一个概念。
\item 找出它准备引出的下一个概念。
\item 把中间的关键等式或构造单独抄出来。
\end{itemize}

\smallskip
\textbf{本章读法提醒：} 第五章注意下降数、上升数、逆序数不是同一个统计量。
\end{explainblock}


\begin{sourceblock}{推论 5.4.5 设 n 为非负整数, En 为 Euler 数, 则}
{\small\sloppy
\noindent 推论 5.4.5 设 n 为非负整数, En 为 Euler 数, 则\par
\noindent ∞\par
\noindent X x2n\par
\noindent E2n = sec x, (5.20)\par
\noindent (2n)!\par
\noindent n=0\par
\noindent ∞\par
\noindent X x2n+1\par
\noindent E2n+1 = tan x. (5.21)\par
\noindent (2n + 1)!\par
\noindent n=0\par
\noindent 这里我们给出它的组合证明.\par
\noindent 证明: 首先我们来证明 (5.20). 因为\par
\noindent X ∞\par
\noindent = cos x = (−1)j ,\par
\noindent sec x (2j)!\par
\noindent j=0\par
\noindent 所以我们只需证明\par
}
\end{sourceblock}


\begin{explainblock}
定理段不要只看结论，要看它把哪两个计数方式连在一起。组合学证明最常见的技巧是：左边按一种标准数，右边按另一种标准数，然后说明两边数的是同一批对象。

\smallskip
\textbf{初学者检查点：}
\begin{itemize}[leftmargin=2em,itemsep=0.25em]
\item 先复述定理的假设，不要把适用范围扩大。
\item 找出证明中的基本计数原则：乘法、加法、双射、递推、容斥或生成函数。
\item 检查最后的等式化简是否和前面的对象解释一致。
\end{itemize}

\smallskip
\textbf{本章读法提醒：} 第五章注意下降数、上升数、逆序数不是同一个统计量。
\end{explainblock}


\begin{sourceblock}{第 5 章续读}
{\small\sloppy
\noindent ∞\par
\noindent X X ∞ ∞\par
\noindent x2n 1 x2k X x2j\par
\noindent E2n × = E2k (−1)j\par
\noindent (2n)! sec x (2k)! (2j)!\par
\noindent n=0 k=0 j=0\par
\noindent ∞ X\par
\noindent X n\par
\noindent n−k 2n x2n\par
\noindent = (−1) E2k = 1,\par
\noindent 2k (2n)!\par
\noindent n=0 k=0\par
\noindent 5.4 交错排列与 Euler 数 107\par
\noindent 即证明 (\par
\noindent X\par
\noindent n\par
\noindent n−k 2n 1, n = 0,\par
\noindent (−1) E2k = (5.22)\par
}
\end{sourceblock}


\begin{explainblock}
这一段是原书论述链的一部分。读的时候把它当作桥梁：它通常负责把上一个定义、例子或定理，引到下一个计数方法。

\smallskip
\textbf{初学者检查点：}
\begin{itemize}[leftmargin=2em,itemsep=0.25em]
\item 找出这一段承接的上一个概念。
\item 找出它准备引出的下一个概念。
\item 把中间的关键等式或构造单独抄出来。
\end{itemize}

\smallskip
\textbf{本章读法提醒：} 第五章注意下降数、上升数、逆序数不是同一个统计量。
\end{explainblock}


\begin{sourceblock}{第 5 章续读}
{\small\sloppy
\noindent k=0\par
\noindent 2k 0, n ≥ 1.\par
\noindent 考虑集合 [2n], 设 S ⊂ [2n], S C 是 S 的补集. 设 |S| = 2k, 构造有序对 (α, β), 其\par
\noindent 中 α = a1 a2 · · · a2k 是集合 S 中元素的一个反向交错排列, 即 a1 < a2 > a3 < · · · >\par
\noindent a2k−1 < a2k ; β = b1 b2 · · · b2n−2k 是集合 S C 中元素按递增的顺序排列. 设有序对\par
\noindent (α, β) 的权重为 (−1)n−k .\par
\noindent 1. 当 a2k > b1 或 k = 0 时, 对有序对 (α, β) 作下面的变换:\par
\noindent (α, β) 7→ (α′ , β ′ )\par
\noindent 其中 α′ = a1 a2 · · · a2k b1 b2 , 是长为 2k + 2 的反向交错排列, β ′ = b3 b4 · · · b2n−2k 是长\par
\noindent 为 2n − 2k − 2 的递增排列; (α′ , β ′ ) 对应的权重为 (−1)n−k−1 ;\par
\noindent 2. 当 a2k < b1 或 k = n 时, 对有序对 (α, β) 作下面的变换:\par
\noindent (α, β) 7→ (α′′ , β ′′ )\par
\noindent 其中 α′′ = a1 a2 · · · a2k−2 是 2k − 2 长反向交错排列, β ′′ = a2k−1 a2k b1 · · · b2n−2k\par
\noindent 是递增排列. (α′′ , β ′′ ) 对应的权重为 (−1)n−k+1 ; 显然这种变换是可逆的. 例如\par
\noindent (2516, 347) 对应 (251634, 7), 3612|457 对应 36|12457. 因为 2k 个不同元素组成的反\par
\noindent 向交错排列的个数为 E2k , 由上述变换知等式 (5.22) 成立, 所以等式 (5.20) 成立.\par
\noindent 下面证明 (5.21). 我们只需证明\par
\noindent ∞\par
}
\end{sourceblock}


\begin{explainblock}
这一段是原书论述链的一部分。读的时候把它当作桥梁：它通常负责把上一个定义、例子或定理，引到下一个计数方法。

\smallskip
\textbf{初学者检查点：}
\begin{itemize}[leftmargin=2em,itemsep=0.25em]
\item 找出这一段承接的上一个概念。
\item 找出它准备引出的下一个概念。
\item 把中间的关键等式或构造单独抄出来。
\end{itemize}

\smallskip
\textbf{本章读法提醒：} 第五章注意下降数、上升数、逆序数不是同一个统计量。
\end{explainblock}


\begin{sourceblock}{第 5 章续读}
{\small\sloppy
\noindent X X ∞ ∞\par
\noindent x2k+1 x2k+1 X x2j\par
\noindent E2k+1 × cos x = E2k+1 (−1)j\par
\noindent (2k + 1)! (2k + 1)! (2j)!\par
\noindent k=0 k=0 j=0\par
\noindent ∞ X\par
\noindent X n\par
\noindent n−k 2n + 1 x2n+1\par
\noindent = (−1) E2k+1 = sin x.\par
\noindent 2k + 1 (2n + 1)!\par
\noindent n=0 k=0\par
\noindent 而\par
\noindent ∞\par
\noindent X x2n+1\par
\noindent sin x = (−1)n ,\par
\noindent (2n + 1)!\par
\noindent n=0\par
\noindent 所以只需证明\par
}
\end{sourceblock}


\begin{explainblock}
这一段是原书论述链的一部分。读的时候把它当作桥梁：它通常负责把上一个定义、例子或定理，引到下一个计数方法。

\smallskip
\textbf{初学者检查点：}
\begin{itemize}[leftmargin=2em,itemsep=0.25em]
\item 找出这一段承接的上一个概念。
\item 找出它准备引出的下一个概念。
\item 把中间的关键等式或构造单独抄出来。
\end{itemize}

\smallskip
\textbf{本章读法提醒：} 第五章注意下降数、上升数、逆序数不是同一个统计量。
\end{explainblock}


\begin{sourceblock}{第 5 章续读}
{\small\sloppy
\noindent X\par
\noindent n\par
\noindent n−k 2n + 1\par
\noindent (−1) E2k+1 = (−1)n . (5.23)\par
\noindent 2k + 1\par
\noindent k=0\par
\noindent 设 n ≥ 1, 考虑集合 [2n + 1], 设 S 是 [2n + 1] 的 2k 元子集, S C 为集合 S 的补\par
\noindent 集. 定义由集合 S 中元素反向交错排列和集合 S C 中元素递降的排列组成的有序对\par
\noindent (α, β), 并设其权重为 (−1)n−k . 不妨设 α = a1 a2 · · · a2k+1 , β = b1 b2 · · · b2n−2k .\par
\noindent 1. 当 k 6= 0 时, 若 k = n 或者 a2k+1 > b1 , 对有序对 (α, β) 作如下变换:\par
\noindent (α, β) 7→ (α1 , β1 ),\par
\noindent 其中 α1 = a1 a2 · · · a2k−1 是长为 2k − 1 的反向交错排列, β1 = a2k a2k+1 b1 · · · b2n−2k\par
\noindent 是长为 2n − 2k + 2 的递减排列; (α1 , β1 ) 的权重为 (−1)n−k+1 ;\par
\noindent 2. 当 a2k+1 < b1 时, 对有序对 (α, β) 作如下变换:\par
\noindent (α, β) 7→ (α2 , β2 ),\par
\noindent 其中 α2 = a1 a2 · · · a2k+1 b1 b2 是长为 2k + 3 的反向交错排列, β2 = b3 b4 · · · b2n−2k 是\par
\noindent 长为 2n − 2k − 2 的递减排列; (α2 , β2 ) 的权重为 (−1)n−k−1 ;\par
\noindent 当 k = 0 且 a1 > b1 时, 有序对只有 (2n + 1, 2n(2n − 1) · · · 1), 其权重为 (−1)n ,\par
}
\end{sourceblock}


\begin{explainblock}
这一段是原书论述链的一部分。读的时候把它当作桥梁：它通常负责把上一个定义、例子或定理，引到下一个计数方法。

\smallskip
\textbf{初学者检查点：}
\begin{itemize}[leftmargin=2em,itemsep=0.25em]
\item 找出这一段承接的上一个概念。
\item 找出它准备引出的下一个概念。
\item 把中间的关键等式或构造单独抄出来。
\end{itemize}

\smallskip
\textbf{本章读法提醒：} 第五章注意下降数、上升数、逆序数不是同一个统计量。
\end{explainblock}


\begin{sourceblock}{第 5 章续读}
{\small\sloppy
\noindent 我们对其不做任何变换.\par
\noindent 由以上讨论可知等式 (5.23) 成立, 所以 (5.21) 也成立.\par
}
\end{sourceblock}


\begin{explainblock}
这一段是原书论述链的一部分。读的时候把它当作桥梁：它通常负责把上一个定义、例子或定理，引到下一个计数方法。

\smallskip
\textbf{初学者检查点：}
\begin{itemize}[leftmargin=2em,itemsep=0.25em]
\item 找出这一段承接的上一个概念。
\item 找出它准备引出的下一个概念。
\item 把中间的关键等式或构造单独抄出来。
\end{itemize}

\smallskip
\textbf{本章读法提醒：} 第五章注意下降数、上升数、逆序数不是同一个统计量。
\end{explainblock}


\begin{sourceblock}{注记 5.4.6 E2n 有时也被称作正割 (secant) 数, E2n+1 被称作正切 (tangent) 数.}
{\small\sloppy
\noindent 注记 5.4.6 E2n 有时也被称作正割 (secant) 数, E2n+1 被称作正切 (tangent) 数.\par
\noindent 5.4 交错排列与 Euler 数 109\par
}
\end{sourceblock}


\begin{explainblock}
注记通常是在补充术语、记号、历史或一个容易忽略的约定。它未必给新公式，但常常决定你后面会不会把符号读错。

\smallskip
\textbf{初学者检查点：}
\begin{itemize}[leftmargin=2em,itemsep=0.25em]
\item 记下新符号的读法。
\item 确认这个约定是否只在本章使用。
\item 把注记里的限制条件补到前面的定义旁边。
\end{itemize}

\smallskip
\textbf{本章读法提醒：} 第五章注意下降数、上升数、逆序数不是同一个统计量。
\end{explainblock}


\begin{sourceblock}{习 题 五}
{\small\sloppy
\noindent 习 题 五\par
}
\end{sourceblock}


\begin{explainblock}
习题段不要先追求技巧。先给题目分类：它是在问排列、组合、递推、生成函数、容斥，还是结构计数。分类正确，工具通常就只剩一两个候选。

\smallskip
\textbf{初学者检查点：}
\begin{itemize}[leftmargin=2em,itemsep=0.25em]
\item 判断题目所属工具。
\item 列出变量和限制条件。
\item 先做小规模 n=2、3、4 的例子，确认直觉。
\end{itemize}

\smallskip
\textbf{本章读法提醒：} 第五章注意下降数、上升数、逆序数不是同一个统计量。
\end{explainblock}


\begin{sourceblock}{习题 5.1 给出 Eulerian 数对称性, 即下面等式的另一种证明:}
{\small\sloppy
\noindent 习题 5.1 给出 Eulerian 数对称性, 即下面等式的另一种证明:\par
\noindent A(n, k) = A(n, n − k + 1).\par
}
\end{sourceblock}


\begin{explainblock}
习题段不要先追求技巧。先给题目分类：它是在问排列、组合、递推、生成函数、容斥，还是结构计数。分类正确，工具通常就只剩一两个候选。

\smallskip
\textbf{初学者检查点：}
\begin{itemize}[leftmargin=2em,itemsep=0.25em]
\item 判断题目所属工具。
\item 列出变量和限制条件。
\item 先做小规模 n=2、3、4 的例子，确认直觉。
\end{itemize}

\smallskip
\textbf{本章读法提醒：} 第五章注意下降数、上升数、逆序数不是同一个统计量。
\end{explainblock}


\begin{sourceblock}{习题 5.2 计算 Eulerian 多项式 A6 (x).}
{\small\sloppy
\noindent 习题 5.2 计算 Eulerian 多项式 A6 (x).\par
}
\end{sourceblock}


\begin{explainblock}
习题段不要先追求技巧。先给题目分类：它是在问排列、组合、递推、生成函数、容斥，还是结构计数。分类正确，工具通常就只剩一两个候选。

\smallskip
\textbf{初学者检查点：}
\begin{itemize}[leftmargin=2em,itemsep=0.25em]
\item 判断题目所属工具。
\item 列出变量和限制条件。
\item 先做小规模 n=2、3、4 的例子，确认直觉。
\end{itemize}

\smallskip
\textbf{本章读法提醒：} 第五章注意下降数、上升数、逆序数不是同一个统计量。
\end{explainblock}


\begin{sourceblock}{习题 5.3 计算 16 + 26 + 3+ · · · + n6 .}
{\small\sloppy
\noindent 习题 5.3 计算 16 + 26 + 3+ · · · + n6 .\par
}
\end{sourceblock}


\begin{explainblock}
习题段不要先追求技巧。先给题目分类：它是在问排列、组合、递推、生成函数、容斥，还是结构计数。分类正确，工具通常就只剩一两个候选。

\smallskip
\textbf{初学者检查点：}
\begin{itemize}[leftmargin=2em,itemsep=0.25em]
\item 判断题目所属工具。
\item 列出变量和限制条件。
\item 先做小规模 n=2、3、4 的例子，确认直觉。
\end{itemize}

\smallskip
\textbf{本章读法提醒：} 第五章注意下降数、上升数、逆序数不是同一个统计量。
\end{explainblock}


\begin{sourceblock}{习题 5.4 设有 n 个由 1, 2, . . . , n 标记的盒子, 我们做一个 n 步的实验. 在第 i}
{\small\sloppy
\noindent 习题 5.4 设有 n 个由 1, 2, . . . , n 标记的盒子, 我们做一个 n 步的实验. 在第 i\par
\noindent 步将一个球放入到标记有 1, 2, . . . , i 的盒子中. 记 B(n, k) 表示 n 步实验完成时有\par
\noindent k − 1 个盒子为空的不同的实验方法数. 证明:\par
\noindent B(n, k) = A(n, k).\par
}
\end{sourceblock}


\begin{explainblock}
习题段不要先追求技巧。先给题目分类：它是在问排列、组合、递推、生成函数、容斥，还是结构计数。分类正确，工具通常就只剩一两个候选。

\smallskip
\textbf{初学者检查点：}
\begin{itemize}[leftmargin=2em,itemsep=0.25em]
\item 判断题目所属工具。
\item 列出变量和限制条件。
\item 先做小规模 n=2、3、4 的例子，确认直觉。
\end{itemize}

\smallskip
\textbf{本章读法提醒：} 第五章注意下降数、上升数、逆序数不是同一个统计量。
\end{explainblock}


\begin{sourceblock}{习题 5.5 令 M (n, k) 表示集合 [n] 上排列中至多有 k 个上升位的排列个数, 即:}
{\small\sloppy
\noindent 习题 5.5 令 M (n, k) 表示集合 [n] 上排列中至多有 k 个上升位的排列个数, 即:\par
\noindent X\par
\noindent k\par
\noindent M (n, k) = A(n, n − k + i).\par
\noindent i=0\par
\noindent 试构造一个恰当的满射, 证明\par
\noindent M (n, k) ≤ k!S(n, k).\par
}
\end{sourceblock}


\begin{explainblock}
习题段不要先追求技巧。先给题目分类：它是在问排列、组合、递推、生成函数、容斥，还是结构计数。分类正确，工具通常就只剩一两个候选。

\smallskip
\textbf{初学者检查点：}
\begin{itemize}[leftmargin=2em,itemsep=0.25em]
\item 判断题目所属工具。
\item 列出变量和限制条件。
\item 先做小规模 n=2、3、4 的例子，确认直觉。
\end{itemize}

\smallskip
\textbf{本章读法提醒：} 第五章注意下降数、上升数、逆序数不是同一个统计量。
\end{explainblock}


\begin{sourceblock}{习题 5.6 设 π = π1 π2 · · · πn 是集合 n 上的排列, 我们称 π 在位置 i 处改变了}
{\small\sloppy
\noindent 习题 5.6 设 π = π1 π2 · · · πn 是集合 n 上的排列, 我们称 π 在位置 i 处改变了\par
\noindent 方向, 如果 πi−1 < πi > πi+1 或者 πi−1 > πi < πi+1 . 比如: 排列 3561247 中有 2 个\par
\noindent 改变方向的位置分别为 3, 4. 设 G(n, k) 表示集合 [n] 上排列中有 k − 1 个位置改变\par
\noindent 了方向的排列个数. 设 n, k 都是整正数, 证明\par
\noindent G(n, k) = kG(n − 1, k) + 2G(n − 1, k − 1) + (n − k)G(n − 1, k − 2).\par
}
\end{sourceblock}


\begin{explainblock}
习题段不要先追求技巧。先给题目分类：它是在问排列、组合、递推、生成函数、容斥，还是结构计数。分类正确，工具通常就只剩一两个候选。

\smallskip
\textbf{初学者检查点：}
\begin{itemize}[leftmargin=2em,itemsep=0.25em]
\item 判断题目所属工具。
\item 列出变量和限制条件。
\item 先做小规模 n=2、3、4 的例子，确认直觉。
\end{itemize}

\smallskip
\textbf{本章读法提醒：} 第五章注意下降数、上升数、逆序数不是同一个统计量。
\end{explainblock}


\begin{sourceblock}{习题 5.7 利用容斥原理证明定理 5.3.6, 即对任意 n ≥ k ≥ 1,}
{\small\sloppy
\noindent 习题 5.7 利用容斥原理证明定理 5.3.6, 即对任意 n ≥ k ≥ 1,\par
\noindent X\par
\noindent k\par
\noindent n−r\par
\noindent A(n, k) = S(n, r)r! (−1)k−r .\par
\noindent k−r\par
\noindent r=1\par
}
\end{sourceblock}


\begin{explainblock}
习题段不要先追求技巧。先给题目分类：它是在问排列、组合、递推、生成函数、容斥，还是结构计数。分类正确，工具通常就只剩一两个候选。

\smallskip
\textbf{初学者检查点：}
\begin{itemize}[leftmargin=2em,itemsep=0.25em]
\item 判断题目所属工具。
\item 列出变量和限制条件。
\item 先做小规模 n=2、3、4 的例子，确认直觉。
\end{itemize}

\smallskip
\textbf{本章读法提醒：} 第五章注意下降数、上升数、逆序数不是同一个统计量。
\end{explainblock}


\begin{sourceblock}{习题 5.8 证明:(1) 设 S = \{s1 , s2 , . . . , sk \} ⊆ [n − 1], 记 α(S) 为集合 [n] 上排列}
{\small\sloppy
\noindent 习题 5.8 证明:(1) 设 S = \{s1 , s2 , . . . , sk \} ⊆ [n − 1], 记 α(S) 为集合 [n] 上排列\par
\noindent 的下降位集是 S 的子集的个数, 则\par
\noindent n n − s1 n − sk−1\par
\noindent α(S) = ··· .\par
\noindent s1 s2 − s1 sk − sk−1\par
\noindent (2) 记 β(S) 为集合 [n] 的下降位集是 S 的排列的个数, 则\par
\noindent X\par
\noindent β(S) = (−1)|S|−|T | α(T ),\par
\noindent T ⊆S\par
\noindent 其中 |S|, |T | 分别表示集合 S 和 T 中元素的个数.\par
}
\end{sourceblock}


\begin{explainblock}
习题段不要先追求技巧。先给题目分类：它是在问排列、组合、递推、生成函数、容斥，还是结构计数。分类正确，工具通常就只剩一两个候选。

\smallskip
\textbf{初学者检查点：}
\begin{itemize}[leftmargin=2em,itemsep=0.25em]
\item 判断题目所属工具。
\item 列出变量和限制条件。
\item 先做小规模 n=2、3、4 的例子，确认直觉。
\end{itemize}

\smallskip
\textbf{本章读法提醒：} 第五章注意下降数、上升数、逆序数不是同一个统计量。
\end{explainblock}


\begin{sourceblock}{习题 5.9 证明}
{\small\sloppy
\noindent 习题 5.9 证明\par
\noindent X\par
\noindent n\par
\noindent An (x) = x k!S(n, k)(x − 1)n−k .\par
\noindent k=1\par
}
\end{sourceblock}


\begin{explainblock}
习题段不要先追求技巧。先给题目分类：它是在问排列、组合、递推、生成函数、容斥，还是结构计数。分类正确，工具通常就只剩一两个候选。

\smallskip
\textbf{初学者检查点：}
\begin{itemize}[leftmargin=2em,itemsep=0.25em]
\item 判断题目所属工具。
\item 列出变量和限制条件。
\item 先做小规模 n=2、3、4 的例子，确认直觉。
\end{itemize}

\smallskip
\textbf{本章读法提醒：} 第五章注意下降数、上升数、逆序数不是同一个统计量。
\end{explainblock}


\begin{sourceblock}{习 题 5.10 设 S(n, k) 是 第 二 类 Stirling 数, A(n, k) 是 Eulerian 数. 利 用}
{\small\sloppy
\noindent 习 题 5.10 设 S(n, k) 是 第 二 类 Stirling 数, A(n, k) 是 Eulerian 数. 利 用\par
\noindent P P\par
\noindent S(n, k) = ki=0 (−1)i ki (k − i)n /k! 及 mn = nk=0 A(n, k) m+n−k\par
\noindent n 证明:\par
\noindent X\par
\noindent r\par
\noindent n−k\par
\noindent r!S(n, r) = A(n, k) .\par
\noindent r−k\par
\noindent k=0\par
}
\end{sourceblock}


\begin{explainblock}
习题段不要先追求技巧。先给题目分类：它是在问排列、组合、递推、生成函数、容斥，还是结构计数。分类正确，工具通常就只剩一两个候选。

\smallskip
\textbf{初学者检查点：}
\begin{itemize}[leftmargin=2em,itemsep=0.25em]
\item 判断题目所属工具。
\item 列出变量和限制条件。
\item 先做小规模 n=2、3、4 的例子，确认直觉。
\end{itemize}

\smallskip
\textbf{本章读法提醒：} 第五章注意下降数、上升数、逆序数不是同一个统计量。
\end{explainblock}


\begin{sourceblock}{习题 5.11 设 2-重集合 \{1, 1, 2, 2, . . . , n, n\} 上排列中满足对任意 1 ≤ m ≤ n}
{\small\sloppy
\noindent 习题 5.11 设 2-重集合 \{1, 1, 2, 2, . . . , n, n\} 上排列中满足对任意 1 ≤ m ≤ n\par
\noindent 在排列中 2 个 m 之间的元素均大于 m 的排列组成的集合记为 An . 比如 2-重集合\par
\noindent \{1, 1, 2, 2\} 上这种排列有\par
\noindent A2 = \{1122, 2211, 1221\}.\par
\noindent (1) 设 En,k 表示集合 An 中具有 k 个上升位的排列的个数. 比如 n = 2 时,\par
\noindent A2,0 = 1, A2,1 = 2.\par
\noindent 证明:\par
\noindent En,k = (k + 1)En−1,k + (2n − 1 − k)En−1,k−1 .\par
\noindent (2) 证明:\par
\noindent X\par
\noindent n−1\par
\noindent En,k = (2n − 1)!!.\par
\noindent k=0\par
}
\end{sourceblock}


\begin{explainblock}
习题段不要先追求技巧。先给题目分类：它是在问排列、组合、递推、生成函数、容斥，还是结构计数。分类正确，工具通常就只剩一两个候选。

\smallskip
\textbf{初学者检查点：}
\begin{itemize}[leftmargin=2em,itemsep=0.25em]
\item 判断题目所属工具。
\item 列出变量和限制条件。
\item 先做小规模 n=2、3、4 的例子，确认直觉。
\end{itemize}

\smallskip
\textbf{本章读法提醒：} 第五章注意下降数、上升数、逆序数不是同一个统计量。
\end{explainblock}


\begin{sourceblock}{习题 5.12 试给出三角恒等式 sin2 θ + cos2 θ = 1 的一个组合证明.}
{\small\sloppy
\noindent 习题 5.12 试给出三角恒等式 sin2 θ + cos2 θ = 1 的一个组合证明.\par
}
\end{sourceblock}


\begin{explainblock}
习题段不要先追求技巧。先给题目分类：它是在问排列、组合、递推、生成函数、容斥，还是结构计数。分类正确，工具通常就只剩一两个候选。

\smallskip
\textbf{初学者检查点：}
\begin{itemize}[leftmargin=2em,itemsep=0.25em]
\item 判断题目所属工具。
\item 列出变量和限制条件。
\item 先做小规模 n=2、3、4 的例子，确认直觉。
\end{itemize}

\smallskip
\textbf{本章读法提醒：} 第五章注意下降数、上升数、逆序数不是同一个统计量。
\end{explainblock}


\begin{sourceblock}{习题 5.13 试出下面三角恒等式的组合证明:}
{\small\sloppy
\noindent 习题 5.13 试出下面三角恒等式的组合证明:\par
\noindent tan x + tan y\par
\noindent tan(x + y) = .\par
}
\end{sourceblock}


\begin{explainblock}
习题段不要先追求技巧。先给题目分类：它是在问排列、组合、递推、生成函数、容斥，还是结构计数。分类正确，工具通常就只剩一两个候选。

\smallskip
\textbf{初学者检查点：}
\begin{itemize}[leftmargin=2em,itemsep=0.25em]
\item 判断题目所属工具。
\item 列出变量和限制条件。
\item 先做小规模 n=2、3、4 的例子，确认直觉。
\end{itemize}

\smallskip
\textbf{本章读法提醒：} 第五章注意下降数、上升数、逆序数不是同一个统计量。
\end{explainblock}


\begin{sourceblock}{习题 5.14 例用定理 5.3.3 证明定理 5.3.1.}
{\small\sloppy
\noindent 习题 5.14 例用定理 5.3.3 证明定理 5.3.1.\par
}
\end{sourceblock}


\begin{explainblock}
习题段不要先追求技巧。先给题目分类：它是在问排列、组合、递推、生成函数、容斥，还是结构计数。分类正确，工具通常就只剩一两个候选。

\smallskip
\textbf{初学者检查点：}
\begin{itemize}[leftmargin=2em,itemsep=0.25em]
\item 判断题目所属工具。
\item 列出变量和限制条件。
\item 先做小规模 n=2、3、4 的例子，确认直觉。
\end{itemize}

\smallskip
\textbf{本章读法提醒：} 第五章注意下降数、上升数、逆序数不是同一个统计量。
\end{explainblock}


\chapter{Bernoulli 数与 Bernoulli 多项式}
\begin{roadmapblock}
本章保留原书正文的主要数学骨架，并在每个定义、定理、例题和论述片段后加入解读。
阅读时不要把目标放在“背下答案”上，而是放在“看懂这一步为什么这样数”上。

\smallskip
\textbf{本章最低目标：} 能说出每个公式左边在数什么、右边在数什么，以及为什么两边相等。
\end{roadmapblock}


\begin{sourceblock}{第 6 章正文}
{\small\sloppy
\noindent Bernoulli 数与 Bernoulli 多项式\par
\noindent 连续自然数方幂求和的问题由来已久. 为了解决这一问题, 雅各布 · 伯努利对\par
\noindent Bernoulli 数进行研究并利用 Bernoulli 数的相关性质得到了一种快速计算出连续自\par
\noindent 然数的方幂和的方法, 并宣称用了不到七分半钟就能前 1000 个自然数的十次幂的\par
\noindent 和:\par
\noindent 110 + 210 + · · · + 100010 = 91409924241424243424241924242500.\par
}
\end{sourceblock}


\begin{explainblock}
这一段是原书论述链的一部分。读的时候把它当作桥梁：它通常负责把上一个定义、例子或定理，引到下一个计数方法。

\smallskip
\textbf{初学者检查点：}
\begin{itemize}[leftmargin=2em,itemsep=0.25em]
\item 找出这一段承接的上一个概念。
\item 找出它准备引出的下一个概念。
\item 把中间的关键等式或构造单独抄出来。
\end{itemize}

\smallskip
\textbf{本章读法提醒：} 第六章把 Bernoulli 数放回生成函数里看，不要把它们当成孤立常数。
\end{explainblock}


\begin{sourceblock}{§ 6.1 Bernoulli 数的基本性质}
{\small\sloppy
\noindent § 6.1. Bernoulli 数的基本性质\par
\noindent 下面我们先给出 Bernoulli 数的定义.\par
}
\end{sourceblock}


\begin{explainblock}
这一节先不要急着背公式。先问三个问题：对象是什么，哪些限制条件不能丢，最后到底是在数所有对象、还是在数某个统计量的分布。把这三件事说清楚，后面的公式才会有位置。

\smallskip
\textbf{初学者检查点：}
\begin{itemize}[leftmargin=2em,itemsep=0.25em]
\item 把本节出现的新对象写成一句话定义。
\item 把每个公式左边和右边分别翻译成“在数什么”。
\item 找一个最小非平凡例子，手工列出所有对象。
\end{itemize}

\smallskip
\textbf{本章读法提醒：} 第六章把 Bernoulli 数放回生成函数里看，不要把它们当成孤立常数。
\end{explainblock}


\begin{sourceblock}{定义 6.1.1 设 B0 = 1, 则由下面的递推关系定义:}
{\small\sloppy
\noindent 定义 6.1.1 设 B0 = 1, 则由下面的递推关系定义:\par
\noindent n−1\par
\noindent Bn = − Bk , (6.1)\par
\noindent n+1 k\par
\noindent k=0\par
\noindent 序列称为 Bernoulli 数.\par
}
\end{sourceblock}


\begin{explainblock}
定义段的任务是定边界。这里要特别注意参数范围、是否允许重复、是否考虑顺序、对象有没有标签。后面的每一道例题和证明，本质上都在反复使用这些边界。

\smallskip
\textbf{初学者检查点：}
\begin{itemize}[leftmargin=2em,itemsep=0.25em]
\item 圈出被定义的对象名称。
\item 圈出参数范围，例如 n、k 是否要求正整数。
\item 判断它是有序对象、无序对象、带标签对象还是结构对象。
\end{itemize}

\smallskip
\textbf{本章读法提醒：} 第六章把 Bernoulli 数放回生成函数里看，不要把它们当成孤立常数。
\end{explainblock}


\begin{sourceblock}{例 6.1.1 由 Bernoulli 数的定义可直接计算得:}
{\small\sloppy
\noindent 例 6.1.1 由 Bernoulli 数的定义可直接计算得:\par
\noindent B1 = − , B2 = , B3 = 0,\par
\noindent B4 = − , B5 = 0, B6 = .\par
}
\end{sourceblock}


\begin{explainblock}
例题是学习这本书最重要的部分。不要只看答案，要把每个限制条件变成一个计数动作：先安排谁、再插入谁、有没有重复、需不需要除以对称性。

\smallskip
\textbf{初学者检查点：}
\begin{itemize}[leftmargin=2em,itemsep=0.25em]
\item 先写出没有限制时的总数。
\item 逐条加入限制条件，观察总数怎样变化。
\item 最后检查有没有把同一种方案重复算了多次。
\end{itemize}

\smallskip
\textbf{本章读法提醒：} 第六章把 Bernoulli 数放回生成函数里看，不要把它们当成孤立常数。
\end{explainblock}


\begin{sourceblock}{注记 6.1.2 由 (6.1) 可知当 n ≥ 1 时,}
{\small\sloppy
\noindent 注记 6.1.2 由 (6.1) 可知当 n ≥ 1 时,\par
\noindent n\par
\noindent X\par
\noindent n+1\par
\noindent Bk = 0, (6.2)\par
\noindent k\par
\noindent k=0\par
\noindent 所以\par
\noindent X\par
\noindent n+1\par
\noindent n+1\par
\noindent Bk = Bn+1 , n ≥ 1, B0 = 1. (6.3)\par
\noindent k\par
\noindent k=0\par
\noindent 112 第 6 章 Bernoulli 数与 Bernoulli 多项式\par
\noindent 上式也可作为 Bernoulli 数的定义. 另外, 由式 (6.2) 可知\par
\noindent n\par
\noindent X n\par
}
\end{sourceblock}


\begin{explainblock}
注记通常是在补充术语、记号、历史或一个容易忽略的约定。它未必给新公式，但常常决定你后面会不会把符号读错。

\smallskip
\textbf{初学者检查点：}
\begin{itemize}[leftmargin=2em,itemsep=0.25em]
\item 记下新符号的读法。
\item 确认这个约定是否只在本章使用。
\item 把注记里的限制条件补到前面的定义旁边。
\end{itemize}

\smallskip
\textbf{本章读法提醒：} 第六章把 Bernoulli 数放回生成函数里看，不要把它们当成孤立常数。
\end{explainblock}


\begin{sourceblock}{第 6 章续读}
{\small\sloppy
\noindent n+1 n+1 X n+1\par
\noindent Bk = 1 − + Bk = 0,\par
\noindent k 2 k\par
\noindent k=0 k=2\par
\noindent 所以\par
\noindent n\par
\noindent n+1 X n+1\par
\noindent 1+ + Bk = n + 1. (6.4)\par
\noindent k=2\par
\noindent 故若取 B1 = 12 , (6.4) 则可写为:\par
\noindent n\par
\noindent X\par
\noindent n+1\par
\noindent Bk = n + 1,\par
\noindent k\par
\noindent k=0\par
\noindent 因此, 也有书中以上式作为 Bernoulli 数的定义. 本书采用第一种定义方式.\par
}
\end{sourceblock}


\begin{explainblock}
这一段是原书论述链的一部分。读的时候把它当作桥梁：它通常负责把上一个定义、例子或定理，引到下一个计数方法。

\smallskip
\textbf{初学者检查点：}
\begin{itemize}[leftmargin=2em,itemsep=0.25em]
\item 找出这一段承接的上一个概念。
\item 找出它准备引出的下一个概念。
\item 把中间的关键等式或构造单独抄出来。
\end{itemize}

\smallskip
\textbf{本章读法提醒：} 第六章把 Bernoulli 数放回生成函数里看，不要把它们当成孤立常数。
\end{explainblock}


\begin{sourceblock}{定理 6.1.3 设 B(x) 为 Bernoulli 数的指数型生成函数, 则}
{\small\sloppy
\noindent 定理 6.1.3 设 B(x) 为 Bernoulli 数的指数型生成函数, 则\par
\noindent ∞\par
\noindent X xn x\par
\noindent B(x) = Bn = x .\par
\noindent n! e −1\par
\noindent n=0\par
\noindent 证明: 因为\par
\noindent ∞\par
\noindent X xn\par
\noindent B(x) = Bn\par
\noindent n!\par
\noindent n=0\par
\noindent ∞\par
\noindent x X xn\par
\noindent =1− + Bn ,\par
\noindent n=2\par
\noindent 再由 (6.3) 知当 n ≥ 2 时\par
\noindent n\par
}
\end{sourceblock}


\begin{explainblock}
定理段不要只看结论，要看它把哪两个计数方式连在一起。组合学证明最常见的技巧是：左边按一种标准数，右边按另一种标准数，然后说明两边数的是同一批对象。

\smallskip
\textbf{初学者检查点：}
\begin{itemize}[leftmargin=2em,itemsep=0.25em]
\item 先复述定理的假设，不要把适用范围扩大。
\item 找出证明中的基本计数原则：乘法、加法、双射、递推、容斥或生成函数。
\item 检查最后的等式化简是否和前面的对象解释一致。
\end{itemize}

\smallskip
\textbf{本章读法提醒：} 第六章把 Bernoulli 数放回生成函数里看，不要把它们当成孤立常数。
\end{explainblock}


\begin{sourceblock}{第 6 章续读}
{\small\sloppy
\noindent X n\par
\noindent Bk = Bn ,\par
\noindent k\par
\noindent k=0\par
\noindent 所以\par
\noindent ∞ n\par
\noindent x XX n xn\par
\noindent B(x) = 1 − + Bk\par
\noindent n=2 k=0\par
\noindent ∞\par
\noindent x X X Bk xk\par
\noindent n\par
\noindent xn−k\par
\noindent =1− + × .\par
\noindent n=2 k=0\par
\noindent 当 n = 1 时,\par
\noindent n\par
\noindent X\par
}
\end{sourceblock}


\begin{explainblock}
这一段是原书论述链的一部分。读的时候把它当作桥梁：它通常负责把上一个定义、例子或定理，引到下一个计数方法。

\smallskip
\textbf{初学者检查点：}
\begin{itemize}[leftmargin=2em,itemsep=0.25em]
\item 找出这一段承接的上一个概念。
\item 找出它准备引出的下一个概念。
\item 把中间的关键等式或构造单独抄出来。
\end{itemize}

\smallskip
\textbf{本章读法提醒：} 第六章把 Bernoulli 数放回生成函数里看，不要把它们当成孤立常数。
\end{explainblock}


\begin{sourceblock}{第 6 章续读}
{\small\sloppy
\noindent n 1 1 1 1\par
\noindent Bk = B0 + B1 = 1 − = ,\par
\noindent k 0 1 2 2\par
\noindent k=0\par
\noindent 当 n = 0 时,\par
\noindent n\par
\noindent X n\par
\noindent Bk = 1,\par
\noindent k\par
\noindent k=0\par
\noindent 6.1 Bernoulli 数的基本性质 113\par
\noindent 所以\par
\noindent ∞\par
\noindent x X X Bk xk\par
\noindent n\par
\noindent xn−k\par
\noindent B(x) = 1 − + ×\par
\noindent n=2 k=0\par
}
\end{sourceblock}


\begin{explainblock}
这一段是原书论述链的一部分。读的时候把它当作桥梁：它通常负责把上一个定义、例子或定理，引到下一个计数方法。

\smallskip
\textbf{初学者检查点：}
\begin{itemize}[leftmargin=2em,itemsep=0.25em]
\item 找出这一段承接的上一个概念。
\item 找出它准备引出的下一个概念。
\item 把中间的关键等式或构造单独抄出来。
\end{itemize}

\smallskip
\textbf{本章读法提醒：} 第六章把 Bernoulli 数放回生成函数里看，不要把它们当成孤立常数。
\end{explainblock}


\begin{sourceblock}{第 6 章续读}
{\small\sloppy
\noindent ∞ X\par
\noindent X n\par
\noindent xk xn−k\par
\noindent = −x + Bk ×\par
\noindent k! (n − k)!\par
\noindent n=0 k=0\par
\noindent ∞\par
\noindent X ∞\par
\noindent Bk xk X xn−k\par
\noindent = −x +\par
\noindent k! (n − k)!\par
\noindent k=0 n=k\par
\noindent = −x + B(x)ex ,\par
\noindent 故\par
\noindent x\par
\noindent B(x) = .\par
\noindent ex − 1\par
}
\end{sourceblock}


\begin{explainblock}
这一段是原书论述链的一部分。读的时候把它当作桥梁：它通常负责把上一个定义、例子或定理，引到下一个计数方法。

\smallskip
\textbf{初学者检查点：}
\begin{itemize}[leftmargin=2em,itemsep=0.25em]
\item 找出这一段承接的上一个概念。
\item 找出它准备引出的下一个概念。
\item 把中间的关键等式或构造单独抄出来。
\end{itemize}

\smallskip
\textbf{本章读法提醒：} 第六章把 Bernoulli 数放回生成函数里看，不要把它们当成孤立常数。
\end{explainblock}


\begin{sourceblock}{注记 6.1.4 该定理也可作为 Bernoulli 数定义.}
{\small\sloppy
\noindent 注记 6.1.4 该定理也可作为 Bernoulli 数定义.\par
\noindent 由 Bernoulli 数的指数型生成函数很容易得到下面的性质.\par
\noindent 性质 6.1.5 设 n 为任意正整数, 则\par
\noindent B2n+1 = 0.\par
\noindent 证明: 因为\par
\noindent ∞\par
\noindent X x2n+1\par
\noindent B(x) − B(−x) = 2 B2n+1\par
\noindent (2n + 1)!\par
\noindent n=0\par
\noindent x −x\par
\noindent = −\par
\noindent ex − 1 e−x − 1\par
\noindent = −x,\par
\noindent 比较系数即可.\par
\noindent 下面我们再给出 Bernoulli 数的另一个性质, 证明将在 6.2 节给出.\par
\noindent 性质 6.1.6 设 n 为正整数, 则\par
\noindent (−1)n−1 B2n > 0.\par
}
\end{sourceblock}


\begin{explainblock}
注记通常是在补充术语、记号、历史或一个容易忽略的约定。它未必给新公式，但常常决定你后面会不会把符号读错。

\smallskip
\textbf{初学者检查点：}
\begin{itemize}[leftmargin=2em,itemsep=0.25em]
\item 记下新符号的读法。
\item 确认这个约定是否只在本章使用。
\item 把注记里的限制条件补到前面的定义旁边。
\end{itemize}

\smallskip
\textbf{本章读法提醒：} 第六章把 Bernoulli 数放回生成函数里看，不要把它们当成孤立常数。
\end{explainblock}


\begin{sourceblock}{第 6 章续读}
{\small\sloppy
\noindent 下面我们给出由 Bernoulli 数计算前 n 个自然数的正整数次方的和的公式.\par
}
\end{sourceblock}


\begin{explainblock}
这一段是原书论述链的一部分。读的时候把它当作桥梁：它通常负责把上一个定义、例子或定理，引到下一个计数方法。

\smallskip
\textbf{初学者检查点：}
\begin{itemize}[leftmargin=2em,itemsep=0.25em]
\item 找出这一段承接的上一个概念。
\item 找出它准备引出的下一个概念。
\item 把中间的关键等式或构造单独抄出来。
\end{itemize}

\smallskip
\textbf{本章读法提醒：} 第六章把 Bernoulli 数放回生成函数里看，不要把它们当成孤立常数。
\end{explainblock}


\begin{sourceblock}{定理 6.1.7 设 Sm (n) = 1m + 2m + 3m + · · · + nm , 则}
{\small\sloppy
\noindent 定理 6.1.7 设 Sm (n) = 1m + 2m + 3m + · · · + nm , 则\par
\noindent m\par
\noindent X\par
\noindent m + 1 (n + 1)m−k+1\par
\noindent Sm (n) = Bk .\par
\noindent k m+1\par
\noindent k=0\par
\noindent 114 第 6 章 Bernoulli 数与 Bernoulli 多项式\par
\noindent 证明: (方法一) 因为\par
\noindent ∞\par
\noindent X\par
\noindent jx xm\par
\noindent e = jm ,\par
\noindent m!\par
\noindent m=0\par
\noindent 所以\par
\noindent ∞\par
\noindent X xm\par
}
\end{sourceblock}


\begin{explainblock}
定理段不要只看结论，要看它把哪两个计数方式连在一起。组合学证明最常见的技巧是：左边按一种标准数，右边按另一种标准数，然后说明两边数的是同一批对象。

\smallskip
\textbf{初学者检查点：}
\begin{itemize}[leftmargin=2em,itemsep=0.25em]
\item 先复述定理的假设，不要把适用范围扩大。
\item 找出证明中的基本计数原则：乘法、加法、双射、递推、容斥或生成函数。
\item 检查最后的等式化简是否和前面的对象解释一致。
\end{itemize}

\smallskip
\textbf{本章读法提醒：} 第六章把 Bernoulli 数放回生成函数里看，不要把它们当成孤立常数。
\end{explainblock}


\begin{sourceblock}{第 6 章续读}
{\small\sloppy
\noindent x\par
\noindent e +e 2x\par
\noindent + ··· + e nx\par
\noindent = Sm (n) . (6.5)\par
\noindent m!\par
\noindent m=0\par
\noindent 另一方面, 由等比数列求和公式可知:\par
\noindent ex (enx − 1)\par
\noindent ex + e2x + · · · + enx = .\par
\noindent ex − 1\par
\noindent 而\par
\noindent ex (enx − 1) e(n+1)x − ex x\par
\noindent = × x\par
\noindent e −1\par
\noindent x x e −1\par
\noindent ∞ ∞\par
\noindent 1X xj X xk\par
\noindent = ((n + 1)j − 1) Bk\par
}
\end{sourceblock}


\begin{explainblock}
这一段是原书论述链的一部分。读的时候把它当作桥梁：它通常负责把上一个定义、例子或定理，引到下一个计数方法。

\smallskip
\textbf{初学者检查点：}
\begin{itemize}[leftmargin=2em,itemsep=0.25em]
\item 找出这一段承接的上一个概念。
\item 找出它准备引出的下一个概念。
\item 把中间的关键等式或构造单独抄出来。
\end{itemize}

\smallskip
\textbf{本章读法提醒：} 第六章把 Bernoulli 数放回生成函数里看，不要把它们当成孤立常数。
\end{explainblock}


\begin{sourceblock}{第 6 章续读}
{\small\sloppy
\noindent x j! k!\par
\noindent j=0 k=0\par
\noindent ∞\par
\noindent X ∞\par
\noindent (n + 1)j+1 − 1 xj X xk\par
\noindent = Bk\par
\noindent j+1 j! k!\par
\noindent j=0 k=0\par
\noindent ∞\par
\noindent X\par
\noindent xm X m (n + 1)j+1 − 1\par
\noindent = Bk\par
\noindent m! k+j=m k j+1\par
\noindent m=0\par
\noindent 0≤k≤m\par
\noindent ∞\par
\noindent X m\par
\noindent xm X m (n + 1)m−k+1 − 1\par
}
\end{sourceblock}


\begin{explainblock}
这一段是原书论述链的一部分。读的时候把它当作桥梁：它通常负责把上一个定义、例子或定理，引到下一个计数方法。

\smallskip
\textbf{初学者检查点：}
\begin{itemize}[leftmargin=2em,itemsep=0.25em]
\item 找出这一段承接的上一个概念。
\item 找出它准备引出的下一个概念。
\item 把中间的关键等式或构造单独抄出来。
\end{itemize}

\smallskip
\textbf{本章读法提醒：} 第六章把 Bernoulli 数放回生成函数里看，不要把它们当成孤立常数。
\end{explainblock}


\begin{sourceblock}{第 6 章续读}
{\small\sloppy
\noindent = Bk\par
\noindent m! k m−k+1\par
\noindent m=0 k=0\par
\noindent ∞\par
\noindent X m\par
\noindent xm X m + 1 (n + 1)m−k+1 − 1\par
\noindent = Bk . (6.6)\par
\noindent m! k m+1\par
\noindent m=0 k=0\par
\noindent 比较 (6.5) 和 (6.6) 两个式子的系数可得:\par
\noindent m\par
\noindent X\par
\noindent m + 1 (n + 1)m−k+1 − 1\par
\noindent Sm (n) = Bk\par
\noindent k m+1\par
\noindent k=0\par
\noindent m\par
\noindent X m\par
}
\end{sourceblock}


\begin{explainblock}
这一段是原书论述链的一部分。读的时候把它当作桥梁：它通常负责把上一个定义、例子或定理，引到下一个计数方法。

\smallskip
\textbf{初学者检查点：}
\begin{itemize}[leftmargin=2em,itemsep=0.25em]
\item 找出这一段承接的上一个概念。
\item 找出它准备引出的下一个概念。
\item 把中间的关键等式或构造单独抄出来。
\end{itemize}

\smallskip
\textbf{本章读法提醒：} 第六章把 Bernoulli 数放回生成函数里看，不要把它们当成孤立常数。
\end{explainblock}


\begin{sourceblock}{第 6 章续读}
{\small\sloppy
\noindent m + 1 (n + 1)m−k+1 1 X m+1\par
\noindent = Bk − Bk\par
\noindent k m+1 m+1 k\par
\noindent k=0 k=0\par
\noindent m\par
\noindent X\par
\noindent m + 1 (n + 1)m−k+1\par
\noindent = Bk . (6.7)\par
\noindent k m+1\par
\noindent k=0\par
\noindent 下面我们用 Telescoping 的方法给出定理 6.1.7 的另一种证明.\par
\noindent 6.1 Bernoulli 数的基本性质 115\par
\noindent (方法二) 设 tn = nm , 我们先求满足\par
\noindent zn+1 − zn = tn (6.8)\par
\noindent 的超几何项 zn . 所谓超几何项是指一个序列 \{zn \}n≥0 的相邻两项的比 zn+1\par
\noindent zn 是关于\par
\noindent n 的有理函数. 等式 (6.8)两边同时除以 zn , 可得\par
\noindent zn+1 tn\par
}
\end{sourceblock}


\begin{explainblock}
这一段是原书论述链的一部分。读的时候把它当作桥梁：它通常负责把上一个定义、例子或定理，引到下一个计数方法。

\smallskip
\textbf{初学者检查点：}
\begin{itemize}[leftmargin=2em,itemsep=0.25em]
\item 找出这一段承接的上一个概念。
\item 找出它准备引出的下一个概念。
\item 把中间的关键等式或构造单独抄出来。
\end{itemize}

\smallskip
\textbf{本章读法提醒：} 第六章把 Bernoulli 数放回生成函数里看，不要把它们当成孤立常数。
\end{explainblock}


\begin{sourceblock}{第 6 章续读}
{\small\sloppy
\noindent −1= ,\par
\noindent zn zn\par
\noindent 故 tn /zn 是关于 n 的有理函数, 又因为 tn = nm 是 n 的多项式, 所以 zn 也是关于\par
\noindent n 的有理函数. 由 zn 满足等式 (6.8) 得, zn 是关于 n 的多项式, 且 zn 的最高次数为\par
\noindent m + 1 [16]. 设\par
\noindent X\par
\noindent m+1\par
\noindent zn = ai · ni ,\par
\noindent i=0\par
\noindent 代入上式得,\par
\noindent X\par
\noindent m+1 X\par
\noindent m+1\par
\noindent ai · (n + 1) −\par
\noindent i\par
\noindent ai · ni = nm . (6.9)\par
\noindent i=0 i=0\par
\noindent 显然 a0 的值不影响结果, 不妨取 a0 = 0. 比较等式 (6.9) 两边相同次幂的系数得,\par
}
\end{sourceblock}


\begin{explainblock}
这一段是原书论述链的一部分。读的时候把它当作桥梁：它通常负责把上一个定义、例子或定理，引到下一个计数方法。

\smallskip
\textbf{初学者检查点：}
\begin{itemize}[leftmargin=2em,itemsep=0.25em]
\item 找出这一段承接的上一个概念。
\item 找出它准备引出的下一个概念。
\item 把中间的关键等式或构造单独抄出来。
\end{itemize}

\smallskip
\textbf{本章读法提醒：} 第六章把 Bernoulli 数放回生成函数里看，不要把它们当成孤立常数。
\end{explainblock}


\begin{sourceblock}{第 6 章续读}
{\small\sloppy
\noindent am+1 = m+1 , 且对于任意 0 ≤ i ≤ m − 1, 有\par
\noindent X\par
\noindent m+1\par
\noindent j\par
\noindent aj = 0.\par
\noindent i\par
\noindent j=i+1\par
\noindent 令 j → m + 1 − j, 则可得\par
\noindent X\par
\noindent m−i\par
\noindent m+1−i\par
\noindent · j! · (m + 1 − j)! · am+1−j = 0. (6.10)\par
\noindent j\par
\noindent j=0\par
\noindent 令\par
\noindent Aj,m = j!(m + 1 − j)!am+1−j /m!,\par
\noindent 则 A0,m = (m + 1)am+1 = 1. 又由 Bernoulli 数的定义知 B0 = 1 且 n ≥ 1 时满足\par
\noindent (6.2), 即:\par
}
\end{sourceblock}


\begin{explainblock}
这一段是原书论述链的一部分。读的时候把它当作桥梁：它通常负责把上一个定义、例子或定理，引到下一个计数方法。

\smallskip
\textbf{初学者检查点：}
\begin{itemize}[leftmargin=2em,itemsep=0.25em]
\item 找出这一段承接的上一个概念。
\item 找出它准备引出的下一个概念。
\item 把中间的关键等式或构造单独抄出来。
\end{itemize}

\smallskip
\textbf{本章读法提醒：} 第六章把 Bernoulli 数放回生成函数里看，不要把它们当成孤立常数。
\end{explainblock}


\begin{sourceblock}{第 6 章续读}
{\small\sloppy
\noindent n\par
\noindent X\par
\noindent n+1\par
\noindent Bk = 0.\par
\noindent k\par
\noindent k=0\par
\noindent 所以对于固定的常数 m(m ≥ 1) 可知 Aj,m = Bj , 所以\par
\noindent Bj Bj m+1\par
\noindent am+1−j = m! = ,\par
\noindent j!(m + 1 − j)! m+1 j\par
\noindent 即对于 1 ≤ i ≤ m + 1\par
\noindent Bm+1−i m+1\par
\noindent ai = .\par
\noindent m+1 m+1−i\par
\noindent 116 第 6 章 Bernoulli 数与 Bernoulli 多项式\par
\noindent 故\par
\noindent X\par
\noindent m+1\par
}
\end{sourceblock}


\begin{explainblock}
这一段是原书论述链的一部分。读的时候把它当作桥梁：它通常负责把上一个定义、例子或定理，引到下一个计数方法。

\smallskip
\textbf{初学者检查点：}
\begin{itemize}[leftmargin=2em,itemsep=0.25em]
\item 找出这一段承接的上一个概念。
\item 找出它准备引出的下一个概念。
\item 把中间的关键等式或构造单独抄出来。
\end{itemize}

\smallskip
\textbf{本章读法提醒：} 第六章把 Bernoulli 数放回生成函数里看，不要把它们当成孤立常数。
\end{explainblock}


\begin{sourceblock}{第 6 章续读}
{\small\sloppy
\noindent zn = Bm+1−i ni .\par
\noindent m+1 m+1−i\par
\noindent i=1\par
\noindent 因此,\par
\noindent Sm (n) = 1m + 2m + 3m + · · · + (n − 1)m + nm\par
\noindent = (z2 − z1 ) + (z3 − z2 ) + · · · + (zn − zn−1 ) − (zn+1 − zn )\par
\noindent = zn+1 − zn\par
\noindent X\par
\noindent m+1 X 1 m+1\par
\noindent m+1\par
\noindent = Bm+1−k · (n + 1) −\par
\noindent k\par
\noindent Bm+1−k\par
\noindent m+1 m+1−k m+1 m+1−k\par
\noindent k=1 k=1\par
\noindent X\par
\noindent m X\par
\noindent m\par
}
\end{sourceblock}


\begin{explainblock}
这一段是原书论述链的一部分。读的时候把它当作桥梁：它通常负责把上一个定义、例子或定理，引到下一个计数方法。

\smallskip
\textbf{初学者检查点：}
\begin{itemize}[leftmargin=2em,itemsep=0.25em]
\item 找出这一段承接的上一个概念。
\item 找出它准备引出的下一个概念。
\item 把中间的关键等式或构造单独抄出来。
\end{itemize}

\smallskip
\textbf{本章读法提醒：} 第六章把 Bernoulli 数放回生成函数里看，不要把它们当成孤立常数。
\end{explainblock}


\begin{sourceblock}{第 6 章续读}
{\small\sloppy
\noindent = Bk · (n + 1) m+1−k\par
\noindent − Bk\par
\noindent m+1 k m+1 k\par
\noindent k=0 k=0\par
\noindent X\par
\noindent m\par
\noindent (n + 1)m+1−k m + 1\par
\noindent = Bk .\par
\noindent m+1 k\par
\noindent k=0\par
}
\end{sourceblock}


\begin{explainblock}
这一段是原书论述链的一部分。读的时候把它当作桥梁：它通常负责把上一个定义、例子或定理，引到下一个计数方法。

\smallskip
\textbf{初学者检查点：}
\begin{itemize}[leftmargin=2em,itemsep=0.25em]
\item 找出这一段承接的上一个概念。
\item 找出它准备引出的下一个概念。
\item 把中间的关键等式或构造单独抄出来。
\end{itemize}

\smallskip
\textbf{本章读法提醒：} 第六章把 Bernoulli 数放回生成函数里看，不要把它们当成孤立常数。
\end{explainblock}


\begin{sourceblock}{例 6.1.2 计算 S4 (n) = 14 + 24 + 34 + · · · + n4 .}
{\small\sloppy
\noindent 例 6.1.2 计算 S4 (n) = 14 + 24 + 34 + · · · + n4 .\par
\noindent 解: 由\par
\noindent B0 = 1, B1 = − , B2 = , B3 = 0, B4 = − ,\par
\noindent 代入 (6.7) 知\par
\noindent (n + 1)5 (n + 1)4 (n + 1)3 n + 1 (6Tn2 − Tn )(1 + 2n)\par
\noindent S4 (n) = − + − = ,\par
\noindent 其中 Tn = n(n+1)\par
\noindent 下面我们引进 Bernoulli 多项式的定义.\par
}
\end{sourceblock}


\begin{explainblock}
例题是学习这本书最重要的部分。不要只看答案，要把每个限制条件变成一个计数动作：先安排谁、再插入谁、有没有重复、需不需要除以对称性。

\smallskip
\textbf{初学者检查点：}
\begin{itemize}[leftmargin=2em,itemsep=0.25em]
\item 先写出没有限制时的总数。
\item 逐条加入限制条件，观察总数怎样变化。
\item 最后检查有没有把同一种方案重复算了多次。
\end{itemize}

\smallskip
\textbf{本章读法提醒：} 第六章把 Bernoulli 数放回生成函数里看，不要把它们当成孤立常数。
\end{explainblock}


\begin{sourceblock}{定义 6.1.8 设 n ≥ 0, 多项式 Bn (x) 按如下方式定义:}
{\small\sloppy
\noindent 定义 6.1.8 设 n ≥ 0, 多项式 Bn (x) 按如下方式定义:\par
\noindent n\par
\noindent X n\par
\noindent Bn (x) = Bk xn−k , (6.11)\par
\noindent k\par
\noindent k=0\par
\noindent 我们则称 Bn (x) 为 Bernoulli 多项式.\par
\noindent 由定义可知当 n = 0 时\par
\noindent B0 (x) = 1,\par
\noindent 当 n ≥ 2 时,\par
\noindent n\par
\noindent X n\par
\noindent Bn (0) = Bn , Bn (1) = Bk = Bn ,\par
\noindent k\par
\noindent k=0\par
\noindent 6.1 Bernoulli 数的基本性质 117\par
\noindent 即当 n ≥ 2 时,\par
\noindent Bn (0) = Bn (1) = Bn . (6.12)\par
}
\end{sourceblock}


\begin{explainblock}
定义段的任务是定边界。这里要特别注意参数范围、是否允许重复、是否考虑顺序、对象有没有标签。后面的每一道例题和证明，本质上都在反复使用这些边界。

\smallskip
\textbf{初学者检查点：}
\begin{itemize}[leftmargin=2em,itemsep=0.25em]
\item 圈出被定义的对象名称。
\item 圈出参数范围，例如 n、k 是否要求正整数。
\item 判断它是有序对象、无序对象、带标签对象还是结构对象。
\end{itemize}

\smallskip
\textbf{本章读法提醒：} 第六章把 Bernoulli 数放回生成函数里看，不要把它们当成孤立常数。
\end{explainblock}


\begin{sourceblock}{例 6.1.3 由}
{\small\sloppy
\noindent 例 6.1.3 由\par
\noindent B0 = 1, B1 = − , B2 = , B3 = 0, B4 = − , B5 = 0, B6 = ,\par
\noindent 可计算出前 6 个 Bernoulli 多项式:\par
\noindent B0 (x) = 1,\par
\noindent B1 (x) = x − ,\par
\noindent B2 (x) = x2 − x + ,\par
\noindent B3 (x) = x3 − x2 + ,\par
\noindent B4 (x) = x4 − 2x3 + x2 − ,\par
\noindent B5 (x) = x5 − x4 + x3 − x,\par
\noindent B6 (x) = x6 − 3x5 + x4 − x2 + .\par
\noindent 自然数的整数幂次和公式也可由含有 Bernoulli 多项式的等式来表示.\par
}
\end{sourceblock}


\begin{explainblock}
例题是学习这本书最重要的部分。不要只看答案，要把每个限制条件变成一个计数动作：先安排谁、再插入谁、有没有重复、需不需要除以对称性。

\smallskip
\textbf{初学者检查点：}
\begin{itemize}[leftmargin=2em,itemsep=0.25em]
\item 先写出没有限制时的总数。
\item 逐条加入限制条件，观察总数怎样变化。
\item 最后检查有没有把同一种方案重复算了多次。
\end{itemize}

\smallskip
\textbf{本章读法提醒：} 第六章把 Bernoulli 数放回生成函数里看，不要把它们当成孤立常数。
\end{explainblock}


\begin{sourceblock}{定理 6.1.9 设 n, m 均为正整数, Sm (n) = 1m + 2m + 3m + · · · + nm , 则}
{\small\sloppy
\noindent 定理 6.1.9 设 n, m 均为正整数, Sm (n) = 1m + 2m + 3m + · · · + nm , 则\par
\noindent Bm+1 (n + 1) − Bm+1 (0) Bm+1 (n + 1) − Bm+1 (1)\par
\noindent Sm (n) = = . (6.13)\par
\noindent m+1 m+1\par
\noindent 证明: 因为\par
\noindent m !\par
\noindent X m + 1 (n + 1)m−k+1 1 X\par
\noindent m+1\par
\noindent m+1\par
\noindent Bk = Bk (n + 1) m+1−k\par
\noindent − Bm+1\par
\noindent k m+1 m+1 k\par
\noindent k=0 k=0\par
\noindent Bm+1 (n + 1) − Bm+1\par
\noindent = ,\par
\noindent m+1\par
\noindent 故由公式 (6.7) 及 (6.12) 立得 (6.13).\par
\noindent 下面我们考虑 Bn (x) 关于 x 的导数. 由 (6.11) 可知当 n ≥ 1 时,\par
}
\end{sourceblock}


\begin{explainblock}
定理段不要只看结论，要看它把哪两个计数方式连在一起。组合学证明最常见的技巧是：左边按一种标准数，右边按另一种标准数，然后说明两边数的是同一批对象。

\smallskip
\textbf{初学者检查点：}
\begin{itemize}[leftmargin=2em,itemsep=0.25em]
\item 先复述定理的假设，不要把适用范围扩大。
\item 找出证明中的基本计数原则：乘法、加法、双射、递推、容斥或生成函数。
\item 检查最后的等式化简是否和前面的对象解释一致。
\end{itemize}

\smallskip
\textbf{本章读法提醒：} 第六章把 Bernoulli 数放回生成函数里看，不要把它们当成孤立常数。
\end{explainblock}


\begin{sourceblock}{第 6 章续读}
{\small\sloppy
\noindent X\par
\noindent n−1\par
\noindent n\par
\noindent Bn′ (x) = (n − k)Bk xn−1−k = nBn−1 (x). (6.14)\par
\noindent k\par
\noindent k=0\par
\noindent 因此, 对 Bn (x) 连续求导可得下面的性质.\par
\noindent 118 第 6 章 Bernoulli 数与 Bernoulli 多项式\par
\noindent 性质 6.1.10 设 n ≥ k ≥ 1, 则\par
\noindent dk n!\par
\noindent Bn (x) = Bn(k) (x) = Bn−k (x). (6.15)\par
\noindent dxk (n − k)!\par
\noindent 特别地, 当 n = k ≥ 1 时,\par
\noindent Bn(n) (x) = n!. (6.16)\par
\noindent 当 k = n − 1 时, 由 (6.15) 知\par
\noindent Bnn−1 (x) = n!B1 (x) = n! x− ,\par
\noindent 将 x = 1 和 x = 0 分别代入上式可得\par
\noindent n!\par
}
\end{sourceblock}


\begin{explainblock}
这一段是原书论述链的一部分。读的时候把它当作桥梁：它通常负责把上一个定义、例子或定理，引到下一个计数方法。

\smallskip
\textbf{初学者检查点：}
\begin{itemize}[leftmargin=2em,itemsep=0.25em]
\item 找出这一段承接的上一个概念。
\item 找出它准备引出的下一个概念。
\item 把中间的关键等式或构造单独抄出来。
\end{itemize}

\smallskip
\textbf{本章读法提醒：} 第六章把 Bernoulli 数放回生成函数里看，不要把它们当成孤立常数。
\end{explainblock}


\begin{sourceblock}{第 6 章续读}
{\small\sloppy
\noindent Bn(n−1) (1) = −Bn(n−1) (0) = . (6.17)\par
\noindent 当 n − k ≥ 2 时, 在 (6.15) 中分别取 x = 1 和 x = 0 可得\par
\noindent n!\par
\noindent Bn(k) (1) = Bn−k (1),\par
\noindent (n − k)!\par
\noindent 和\par
\noindent n!\par
\noindent Bn(k) (0) = Bn−k (0).\par
\noindent (n − k)!\par
\noindent 再由 (6.12) 可得当 n − k ≥ 2 时,\par
\noindent n!\par
\noindent Bn(k) (1) = Bn(k) (0) = Bn−k . (6.18)\par
\noindent (n − k)!\par
\noindent 下面我们给出 Bernoulli 多项式的指数型生成函数.\par
}
\end{sourceblock}


\begin{explainblock}
这一段是原书论述链的一部分。读的时候把它当作桥梁：它通常负责把上一个定义、例子或定理，引到下一个计数方法。

\smallskip
\textbf{初学者检查点：}
\begin{itemize}[leftmargin=2em,itemsep=0.25em]
\item 找出这一段承接的上一个概念。
\item 找出它准备引出的下一个概念。
\item 把中间的关键等式或构造单独抄出来。
\end{itemize}

\smallskip
\textbf{本章读法提醒：} 第六章把 Bernoulli 数放回生成函数里看，不要把它们当成孤立常数。
\end{explainblock}


\begin{sourceblock}{定理 6.1.11 设 Bn (x) 为 Bernoulli 多项式, 则}
{\small\sloppy
\noindent 定理 6.1.11 设 Bn (x) 为 Bernoulli 多项式, 则\par
\noindent ∞\par
\noindent X tn tetx\par
\noindent Bn (x) = t . (6.19)\par
\noindent n! e −1\par
\noindent n=0\par
\noindent 证明: 因为\par
\noindent ∞\par
\noindent X xm tm X ∞\par
\noindent t tk\par
\noindent e · t\par
\noindent tx\par
\noindent = Bk\par
\noindent e −1 m! k!\par
\noindent m=0 k=0\par
\noindent X ∞\par
\noindent ∞ X\par
\noindent m+k tm+k\par
}
\end{sourceblock}


\begin{explainblock}
定理段不要只看结论，要看它把哪两个计数方式连在一起。组合学证明最常见的技巧是：左边按一种标准数，右边按另一种标准数，然后说明两边数的是同一批对象。

\smallskip
\textbf{初学者检查点：}
\begin{itemize}[leftmargin=2em,itemsep=0.25em]
\item 先复述定理的假设，不要把适用范围扩大。
\item 找出证明中的基本计数原则：乘法、加法、双射、递推、容斥或生成函数。
\item 检查最后的等式化简是否和前面的对象解释一致。
\end{itemize}

\smallskip
\textbf{本章读法提醒：} 第六章把 Bernoulli 数放回生成函数里看，不要把它们当成孤立常数。
\end{explainblock}


\begin{sourceblock}{第 6 章续读}
{\small\sloppy
\noindent = Bk xm\par
\noindent k (m + k)!\par
\noindent m=0 k=0\par
\noindent n\par
\noindent ∞ X\par
\noindent X n tn\par
\noindent = Bk xn−k ,\par
\noindent k n!\par
\noindent n=0 k=0\par
\noindent 所以等式 (6.19) 成立.\par
\noindent 6.2 Bernoulli 数与黎曼 ζ–函数 119\par
\noindent 性质 6.1.12 设 n ≥ 0, 则\par
\noindent (\par
\noindent 0, n = 0,\par
\noindent Bn (x + 1) − Bn (x) = (6.20)\par
\noindent nxn−1 , n ≥ 1.\par
\noindent 证明: 由 Bernoulli 多项式的指数型生成函数可知\par
\noindent ∞\par
}
\end{sourceblock}


\begin{explainblock}
这一段是原书论述链的一部分。读的时候把它当作桥梁：它通常负责把上一个定义、例子或定理，引到下一个计数方法。

\smallskip
\textbf{初学者检查点：}
\begin{itemize}[leftmargin=2em,itemsep=0.25em]
\item 找出这一段承接的上一个概念。
\item 找出它准备引出的下一个概念。
\item 把中间的关键等式或构造单独抄出来。
\end{itemize}

\smallskip
\textbf{本章读法提醒：} 第六章把 Bernoulli 数放回生成函数里看，不要把它们当成孤立常数。
\end{explainblock}


\begin{sourceblock}{第 6 章续读}
{\small\sloppy
\noindent X tn te(x+1)t text\par
\noindent Bn (x + 1) = t = t + text\par
\noindent n! e −1 e −1\par
\noindent n=0\par
\noindent ∞\par
\noindent X ∞\par
\noindent tn X n tn+1\par
\noindent = Bn (x) + x ,\par
\noindent n! n!\par
\noindent n=0 n=0\par
\noindent 比较两边 tn 的系数立得 (6.20).\par
\noindent 由 (6.20) 知当 m ≥ 0 时,\par
\noindent Bm+1 (x + 1) − Bm+1 (x)\par
\noindent xm = ,\par
\noindent m+1\par
\noindent 在上式中分别取 x = 1, 2, . . . , n 再求和又一次得到了 (6.13).\par
}
\end{sourceblock}


\begin{explainblock}
这一段是原书论述链的一部分。读的时候把它当作桥梁：它通常负责把上一个定义、例子或定理，引到下一个计数方法。

\smallskip
\textbf{初学者检查点：}
\begin{itemize}[leftmargin=2em,itemsep=0.25em]
\item 找出这一段承接的上一个概念。
\item 找出它准备引出的下一个概念。
\item 把中间的关键等式或构造单独抄出来。
\end{itemize}

\smallskip
\textbf{本章读法提醒：} 第六章把 Bernoulli 数放回生成函数里看，不要把它们当成孤立常数。
\end{explainblock}


\begin{sourceblock}{§ 6.2 Bernoulli 数与黎曼 ζ–函数}
{\small\sloppy
\noindent § 6.2. Bernoulli 数与黎曼 ζ–函数\par
\noindent 设 s 为复数且 Re(s) > 1, 黎曼 ζ–函数定义如下:\par
\noindent ∞\par
\noindent X 1\par
\noindent ζ(s) = .\par
\noindent ns\par
\noindent n=1\par
\noindent 它由欧拉引进的, 当时欧拉考虑的是 s 为实数时的情形, 并给出了 ζ–函数的另外一\par
\noindent 种表示方式, 揭示了它与素数的分布有密切的联系. 黎曼在 1859 年找到了 ζ–函数的\par
\noindent 解析延拓: Z ∞\par
\noindent Γ(1 − s) (−z)s\par
\noindent ζ(s) = , (6.21)\par
\noindent 2πi ∞ ez − 1 z\par
\noindent 其中 Γ 是 Γ-函数, 它是阶乘函数在复平面上的解析延拓:\par
\noindent ns n!\par
\noindent Γ(s) = lim .\par
\noindent n→∞ s(s + 1) · · · (s + n − 1)(s + n)\par
\noindent 式 (6.21) 中积分是在一个绕正实轴 (即从 +∞ 出发, 沿实轴上方积分至原点附近,\par
}
\end{sourceblock}


\begin{explainblock}
这一节先不要急着背公式。先问三个问题：对象是什么，哪些限制条件不能丢，最后到底是在数所有对象、还是在数某个统计量的分布。把这三件事说清楚，后面的公式才会有位置。

\smallskip
\textbf{初学者检查点：}
\begin{itemize}[leftmargin=2em,itemsep=0.25em]
\item 把本节出现的新对象写成一句话定义。
\item 把每个公式左边和右边分别翻译成“在数什么”。
\item 找一个最小非平凡例子，手工列出所有对象。
\end{itemize}

\smallskip
\textbf{本章读法提醒：} 第六章把 Bernoulli 数放回生成函数里看，不要把它们当成孤立常数。
\end{explainblock}


\begin{sourceblock}{第 6 章续读}
{\small\sloppy
\noindent 环绕原点积分至实轴下方, 再沿实轴下方积分至 +∞, 积分路径离实轴的距离及环绕\par
\noindent 原点的半径均趋于 0) 进行的围道积分. 可以证明 ζ–函数积分表达式除了 s = 1 处\par
\noindent 有一个单个极点外, 在整个复平面上处处解析. 运用 ζ–函数积分表达式可以得到下\par
\noindent 面的函数方程\par
\noindent ζ(s) = 2Γ(1 − s)(2π)s−1 sin(πs/2)ζ(1 − s),\par
\noindent 从这个关系式中不难发现, ζ–函数在 s = −2n(n ∈ Z) 处取值为零—因为 sin(πs/2)\par
\noindent 为 0. 复平面上的这种使 ζ-函数取值为零的点被称为 ζ–函数的零点. 因此 s =\par
\noindent 120 第 6 章 Bernoulli 数与 Bernoulli 多项式\par
\noindent −2n(n ∈ Z) 是 ζ–函数的零点. 这些零点分布有序, 性质简单, 被称为 ζ–函数的平凡\par
\noindent 零点. 除了这些平凡零点外, ζ–函数还有许多其它零点, 它们的性质远比那些平凡零\par
\noindent 点来得复杂, 被恰如其分地称为非平凡零点. 对 ζ–函数非平凡零点的研究构成了现\par
\noindent 代数学中最艰深的课题之一. 数学界迄今最重要的猜想之一—黎曼猜想就与之有关,\par
\noindent 叙述如下:\par
\noindent 猜想 6.2.1 (黎曼) ζ-函数的所有非平凡零点都位于复平面上 Re(s) = 1/2 的直线\par
\noindent 上.\par
\noindent 有兴趣的读者可以参考 [21, 51].\par
\noindent 利用算术基本定理, 欧拉将 ζ–函数表示成与素数有关的无穷乘积.\par
}
\end{sourceblock}


\begin{explainblock}
这一段是原书论述链的一部分。读的时候把它当作桥梁：它通常负责把上一个定义、例子或定理，引到下一个计数方法。

\smallskip
\textbf{初学者检查点：}
\begin{itemize}[leftmargin=2em,itemsep=0.25em]
\item 找出这一段承接的上一个概念。
\item 找出它准备引出的下一个概念。
\item 把中间的关键等式或构造单独抄出来。
\end{itemize}

\smallskip
\textbf{本章读法提醒：} 第六章把 Bernoulli 数放回生成函数里看，不要把它们当成孤立常数。
\end{explainblock}


\begin{sourceblock}{定理 6.2.2 (欧拉乘积公式) 设 P 是由所有素数组成的集合, 则对于任意 s > 1,}
{\small\sloppy
\noindent 定理 6.2.2 (欧拉乘积公式) 设 P 是由所有素数组成的集合, 则对于任意 s > 1,\par
\noindent X 1 Y 1\par
\noindent −1\par
\noindent = 1− s . (6.22)\par
\noindent ns p\par
\noindent n≥1 p∈P\par
\noindent 欧拉乘积公式实际上与下面的算术基本定理等价.\par
}
\end{sourceblock}


\begin{explainblock}
定理段不要只看结论，要看它把哪两个计数方式连在一起。组合学证明最常见的技巧是：左边按一种标准数，右边按另一种标准数，然后说明两边数的是同一批对象。

\smallskip
\textbf{初学者检查点：}
\begin{itemize}[leftmargin=2em,itemsep=0.25em]
\item 先复述定理的假设，不要把适用范围扩大。
\item 找出证明中的基本计数原则：乘法、加法、双射、递推、容斥或生成函数。
\item 检查最后的等式化简是否和前面的对象解释一致。
\end{itemize}

\smallskip
\textbf{本章读法提醒：} 第六章把 Bernoulli 数放回生成函数里看，不要把它们当成孤立常数。
\end{explainblock}


\begin{sourceblock}{定理 6.2.3 设整数 n > 1, 那么必有}
{\small\sloppy
\noindent 定理 6.2.3 设整数 n > 1, 那么必有\par
\noindent n = pα1 1 pα2 2 · · · pαk k , (6.23)\par
\noindent 其中 p1 , p2 , . . . , pk 是不同的素数, αi (1 ≤ i ≤ k) 是自然数, 而且在不计次序的意义\par
\noindent 下, (6.23) 式唯一.\par
\noindent 下面我们利用算术基本定理给出欧拉乘积公式的证明.\par
\noindent 证明: 我们首先说明当 s > 1 时, 无穷乘积\par
\noindent Y\par
\noindent 1− s (6.24)\par
\noindent p\par
\noindent p∈P\par
\noindent 收敛且大于 1.\par
\noindent 设 s > 1, p 为素数. 因为\par
\noindent ln 1 − s = ln 1 + s < s ,\par
\noindent p p −1 p −1\par
\noindent 所以有\par
\noindent X\par
\noindent ln 1 − s < < 2 < 2 .\par
\noindent p ps − 1 ps ns\par
}
\end{sourceblock}


\begin{explainblock}
定理段不要只看结论，要看它把哪两个计数方式连在一起。组合学证明最常见的技巧是：左边按一种标准数，右边按另一种标准数，然后说明两边数的是同一批对象。

\smallskip
\textbf{初学者检查点：}
\begin{itemize}[leftmargin=2em,itemsep=0.25em]
\item 先复述定理的假设，不要把适用范围扩大。
\item 找出证明中的基本计数原则：乘法、加法、双射、递推、容斥或生成函数。
\item 检查最后的等式化简是否和前面的对象解释一致。
\end{itemize}

\smallskip
\textbf{本章读法提醒：} 第六章把 Bernoulli 数放回生成函数里看，不要把它们当成孤立常数。
\end{explainblock}


\begin{sourceblock}{第 6 章续读}
{\small\sloppy
\noindent p∈P p∈P p∈P n≥1\par
\noindent P 1\par
\noindent 因为级数 n≥1 ns 在 s > 1 时收敛, 所以正项级数\par
\noindent X\par
\noindent ln 1 − s\par
\noindent p\par
\noindent p∈P\par
\noindent 6.2 Bernoulli 数与黎曼 ζ–函数 121\par
\noindent 亦收敛. 由此可推出无穷乘积 (6.24) 收敛, 而且值显然大于 1.\par
\noindent PN\par
\noindent 给定正整数 N , 取正整数 k, 使得 2k−1 ≤ N < 2k . 下面我们考虑和式 1\par
\noindent n=1 ns .\par
\noindent 由算术基本定理知任意正整数 n ≥ 2 可以表示为某些素数的乘积, 而且表示方式唯\par
\noindent 一. 所以我们有下面的不等式:\par
\noindent X\par
\noindent N\par
\noindent Y 1 1\par
\noindent ≤ 1 + s + 2s + · · · + ks < 1 + s + 2s + · · ·\par
}
\end{sourceblock}


\begin{explainblock}
这一段是原书论述链的一部分。读的时候把它当作桥梁：它通常负责把上一个定义、例子或定理，引到下一个计数方法。

\smallskip
\textbf{初学者检查点：}
\begin{itemize}[leftmargin=2em,itemsep=0.25em]
\item 找出这一段承接的上一个概念。
\item 找出它准备引出的下一个概念。
\item 把中间的关键等式或构造单独抄出来。
\end{itemize}

\smallskip
\textbf{本章读法提醒：} 第六章把 Bernoulli 数放回生成函数里看，不要把它们当成孤立常数。
\end{explainblock}


\begin{sourceblock}{第 6 章续读}
{\small\sloppy
\noindent ns p p p p p\par
\noindent n=1 p∈P p∈P\par
\noindent p≤N p≤N\par
\noindent Y\par
\noindent < 1− s < 1− s .\par
\noindent p p\par
\noindent p∈P p∈P\par
\noindent p≤N\par
\noindent P 1\par
\noindent 所以级数 n≥1 ns 收敛, 且\par
\noindent X 1 Y\par
\noindent ≤ 1− s . (6.25)\par
\noindent ns p\par
\noindent n≥1 p∈P\par
\noindent 反过来, 对任意正整数 M 及 h, 取\par
\noindent Y\par
\noindent N1 = ph ,\par
\noindent p∈P\par
}
\end{sourceblock}


\begin{explainblock}
这一段是原书论述链的一部分。读的时候把它当作桥梁：它通常负责把上一个定义、例子或定理，引到下一个计数方法。

\smallskip
\textbf{初学者检查点：}
\begin{itemize}[leftmargin=2em,itemsep=0.25em]
\item 找出这一段承接的上一个概念。
\item 找出它准备引出的下一个概念。
\item 把中间的关键等式或构造单独抄出来。
\end{itemize}

\smallskip
\textbf{本章读法提醒：} 第六章把 Bernoulli 数放回生成函数里看，不要把它们当成孤立常数。
\end{explainblock}


\begin{sourceblock}{第 6 章续读}
{\small\sloppy
\noindent p≤M\par
\noindent 则由算术基本定理可知下式成立:\par
\noindent Y 1 1\par
\noindent XN1\par
\noindent 1 + s + · · · hs ≤ ≤ .\par
\noindent p p ns ns\par
\noindent p∈P n=1 n≥1\par
\noindent p≤M\par
\noindent 令 h → ∞ 得:\par
\noindent Y\par
\noindent 1− s ≤ .\par
\noindent p ns\par
\noindent p∈P n≥1\par
\noindent p≤M\par
\noindent 再令 M → ∞ 得:\par
\noindent Y 1\par
\noindent −1 X 1\par
\noindent 1− s ≤ . (6.26)\par
}
\end{sourceblock}


\begin{explainblock}
这一段是原书论述链的一部分。读的时候把它当作桥梁：它通常负责把上一个定义、例子或定理，引到下一个计数方法。

\smallskip
\textbf{初学者检查点：}
\begin{itemize}[leftmargin=2em,itemsep=0.25em]
\item 找出这一段承接的上一个概念。
\item 找出它准备引出的下一个概念。
\item 把中间的关键等式或构造单独抄出来。
\end{itemize}

\smallskip
\textbf{本章读法提醒：} 第六章把 Bernoulli 数放回生成函数里看，不要把它们当成孤立常数。
\end{explainblock}


\begin{sourceblock}{第 6 章续读}
{\small\sloppy
\noindent p ns\par
\noindent p∈P n≥1\par
\noindent 比较 (6.25) 与 (6.26) 立得 (6.22).\par
\noindent 由欧拉乘积公式可以很容易推出数论中的一个定理:\par
}
\end{sourceblock}


\begin{explainblock}
这一段是原书论述链的一部分。读的时候把它当作桥梁：它通常负责把上一个定义、例子或定理，引到下一个计数方法。

\smallskip
\textbf{初学者检查点：}
\begin{itemize}[leftmargin=2em,itemsep=0.25em]
\item 找出这一段承接的上一个概念。
\item 找出它准备引出的下一个概念。
\item 把中间的关键等式或构造单独抄出来。
\end{itemize}

\smallskip
\textbf{本章读法提醒：} 第六章把 Bernoulli 数放回生成函数里看，不要把它们当成孤立常数。
\end{explainblock}


\begin{sourceblock}{定理 6.2.4 素数的个数有无穷多个.}
{\small\sloppy
\noindent 定理 6.2.4 素数的个数有无穷多个.\par
\noindent 证明: 假设素数个数有有限个, 则极限\par
\noindent Y\par
\noindent lim (1 − p−s )−1\par
\noindent s→1+\par
\noindent p∈P\par
\noindent 122 第 6 章 Bernoulli 数与 Bernoulli 多项式\par
\noindent 存在且有限. 由欧拉乘积公式知\par
\noindent X 1 Y\par
\noindent lim s\par
\noindent = lim (1 − p−s )−1 ,\par
\noindent s→1+ n s→1+\par
\noindent n≥1 p∈P\par
\noindent P∞ 1\par
\noindent 而正项级数 n≥1 n 是发散, 推出矛盾, 故假设不成立, 也即素数个数有无穷多\par
\noindent 个.\par
\noindent 下面我们介绍 ζ–函数与 Bernoulli 数之间的关系. 欧拉当时考虑正弦函数 sin(πz)\par
\noindent 的零点: z = 0, ±1, ±2, . . ., 并猜测它能写成下面无穷乘积的形式:\par
}
\end{sourceblock}


\begin{explainblock}
定理段不要只看结论，要看它把哪两个计数方式连在一起。组合学证明最常见的技巧是：左边按一种标准数，右边按另一种标准数，然后说明两边数的是同一批对象。

\smallskip
\textbf{初学者检查点：}
\begin{itemize}[leftmargin=2em,itemsep=0.25em]
\item 先复述定理的假设，不要把适用范围扩大。
\item 找出证明中的基本计数原则：乘法、加法、双射、递推、容斥或生成函数。
\item 检查最后的等式化简是否和前面的对象解释一致。
\end{itemize}

\smallskip
\textbf{本章读法提醒：} 第六章把 Bernoulli 数放回生成函数里看，不要把它们当成孤立常数。
\end{explainblock}


\begin{sourceblock}{第 6 章续读}
{\small\sloppy
\noindent z z z z\par
\noindent sin (πz) =πz 1 − 1+ 1− 1+ ···\par
\noindent z2 z2 z2\par
\noindent =πz 1 − 2 1− 2 1 − 2 ··· . (6.27)\par
\noindent 这个猜想由魏尔斯特拉斯∗ [22, p.137–138] 给出严格的证明.\par
\noindent 下面我们就应用 sin (πz) 的乘积表达式 (6.27) 来证明 Bernoulli 数与 ζ–函数之\par
\noindent 间的关系.\par
}
\end{sourceblock}


\begin{explainblock}
这一段是原书论述链的一部分。读的时候把它当作桥梁：它通常负责把上一个定义、例子或定理，引到下一个计数方法。

\smallskip
\textbf{初学者检查点：}
\begin{itemize}[leftmargin=2em,itemsep=0.25em]
\item 找出这一段承接的上一个概念。
\item 找出它准备引出的下一个概念。
\item 把中间的关键等式或构造单独抄出来。
\end{itemize}

\smallskip
\textbf{本章读法提醒：} 第六章把 Bernoulli 数放回生成函数里看，不要把它们当成孤立常数。
\end{explainblock}


\begin{sourceblock}{定理 6.2.5 设 k 为正整数, 则}
{\small\sloppy
\noindent 定理 6.2.5 设 k 为正整数, 则\par
\noindent 22k−1 2k\par
\noindent ζ(2k) = (−1)k−1 B2k π . (6.28)\par
\noindent (2k)!\par
\noindent 证明: 设 z 为任意复数, 对 (6.27), 即:\par
\noindent Y z2\par
\noindent sin z = z 1− 2 2 ,\par
\noindent n π\par
\noindent n≥1\par
\noindent 先取对数后再求导可得\par
\noindent X z2\par
\noindent z cot z = 1 + 2 .\par
\noindent z 2 − n2 π 2\par
\noindent n≥1\par
\noindent 因为\par
\noindent z2 X 1 z 2k\par
\noindent = ,\par
\noindent z 2 − n2 π 2 n2k π 2k\par
}
\end{sourceblock}


\begin{explainblock}
定理段不要只看结论，要看它把哪两个计数方式连在一起。组合学证明最常见的技巧是：左边按一种标准数，右边按另一种标准数，然后说明两边数的是同一批对象。

\smallskip
\textbf{初学者检查点：}
\begin{itemize}[leftmargin=2em,itemsep=0.25em]
\item 先复述定理的假设，不要把适用范围扩大。
\item 找出证明中的基本计数原则：乘法、加法、双射、递推、容斥或生成函数。
\item 检查最后的等式化简是否和前面的对象解释一致。
\end{itemize}

\smallskip
\textbf{本章读法提醒：} 第六章把 Bernoulli 数放回生成函数里看，不要把它们当成孤立常数。
\end{explainblock}


\begin{sourceblock}{第 6 章续读}
{\small\sloppy
\noindent k≥1\par
\noindent 所以\par
\noindent X X 1 z 2k\par
\noindent z cot z = 1 − 2 . (6.29)\par
\noindent n2k π 2k\par
\noindent n≥1 k≥1\par
\noindent 又因为\par
\noindent eiz − e−iz eiz + e−iz\par
\noindent sin(z) = , cos(z) = ,\par
\noindent 2i 2\par
\noindent ∗\par
\noindent 卡尔 · 特奥多尔 · 威廉 · 魏尔施特拉斯 (Karl Theodor Wilhelm Weierstrass,1815–1897), 德国数\par
\noindent 学家, 被誉为 “现代分析之父”.\par
\noindent 6.2 Bernoulli 数与黎曼 ζ–函数 123\par
\noindent 所以\par
\noindent eiz + e−iz 2zi\par
\noindent z cot z = zi iz −iz\par
\noindent = 2iz + zi. (6.30)\par
}
\end{sourceblock}


\begin{explainblock}
这一段是原书论述链的一部分。读的时候把它当作桥梁：它通常负责把上一个定义、例子或定理，引到下一个计数方法。

\smallskip
\textbf{初学者检查点：}
\begin{itemize}[leftmargin=2em,itemsep=0.25em]
\item 找出这一段承接的上一个概念。
\item 找出它准备引出的下一个概念。
\item 把中间的关键等式或构造单独抄出来。
\end{itemize}

\smallskip
\textbf{本章读法提醒：} 第六章把 Bernoulli 数放回生成函数里看，不要把它们当成孤立常数。
\end{explainblock}


\begin{sourceblock}{第 6 章续读}
{\small\sloppy
\noindent e −e e −1\par
\noindent 另一方面, 当 k ≥ 1 时, B2k+1 = 0, 所以 Bernoulli 数的生成函数:\par
\noindent ∞\par
\noindent X ∞\par
\noindent xn x X x2k x\par
\noindent Bn =1− + B2k = x ,\par
\noindent n! 2 (2k)! e −1\par
\noindent n=0 k=1\par
\noindent 令 x = 2iz 可得\par
\noindent ∞\par
\noindent X 22k z 2k 2iz\par
\noindent 1 − iz + (−1)k B2k = 2iz . (6.31)\par
\noindent (2k)! e −1\par
\noindent k=1\par
\noindent 由 (6.30) 和 (6.31) 可得:\par
\noindent ∞\par
\noindent X 22k z 2k\par
\noindent z cot z = 1 − (−1)k−1 B2k . (6.32)\par
}
\end{sourceblock}


\begin{explainblock}
这一段是原书论述链的一部分。读的时候把它当作桥梁：它通常负责把上一个定义、例子或定理，引到下一个计数方法。

\smallskip
\textbf{初学者检查点：}
\begin{itemize}[leftmargin=2em,itemsep=0.25em]
\item 找出这一段承接的上一个概念。
\item 找出它准备引出的下一个概念。
\item 把中间的关键等式或构造单独抄出来。
\end{itemize}

\smallskip
\textbf{本章读法提醒：} 第六章把 Bernoulli 数放回生成函数里看，不要把它们当成孤立常数。
\end{explainblock}


\begin{sourceblock}{第 6 章续读}
{\small\sloppy
\noindent (2k)!\par
\noindent k=1\par
\noindent 比较式 (6.29) 和 (6.32) 中的 z 2k 的系数立得\par
\noindent 22k−1 2k\par
\noindent ζ(2k) = (−1)k−1 B2k π .\par
\noindent (2k)!\par
}
\end{sourceblock}


\begin{explainblock}
这一段是原书论述链的一部分。读的时候把它当作桥梁：它通常负责把上一个定义、例子或定理，引到下一个计数方法。

\smallskip
\textbf{初学者检查点：}
\begin{itemize}[leftmargin=2em,itemsep=0.25em]
\item 找出这一段承接的上一个概念。
\item 找出它准备引出的下一个概念。
\item 把中间的关键等式或构造单独抄出来。
\end{itemize}

\smallskip
\textbf{本章读法提醒：} 第六章把 Bernoulli 数放回生成函数里看，不要把它们当成孤立常数。
\end{explainblock}


\begin{sourceblock}{例 6.2.1 由 (6.28) 计算 ζ(2), ζ(4) 和 ζ(6) 的值.}
{\small\sloppy
\noindent 例 6.2.1 由 (6.28) 计算 ζ(2), ζ(4) 和 ζ(6) 的值.\par
\noindent 解: 因为\par
\noindent B2 = , B4 = − , B6 = ,\par
\noindent 由定理 6.2.5 可得\par
\noindent (2π)2 π2\par
\noindent ζ(2) = B2 = ,\par
\noindent (2π)4 π4\par
\noindent ζ(4) = − B4 = ,\par
\noindent (2π)6 π6\par
\noindent ζ(6) = B6 = .\par
}
\end{sourceblock}


\begin{explainblock}
例题是学习这本书最重要的部分。不要只看答案，要把每个限制条件变成一个计数动作：先安排谁、再插入谁、有没有重复、需不需要除以对称性。

\smallskip
\textbf{初学者检查点：}
\begin{itemize}[leftmargin=2em,itemsep=0.25em]
\item 先写出没有限制时的总数。
\item 逐条加入限制条件，观察总数怎样变化。
\item 最后检查有没有把同一种方案重复算了多次。
\end{itemize}

\smallskip
\textbf{本章读法提醒：} 第六章把 Bernoulli 数放回生成函数里看，不要把它们当成孤立常数。
\end{explainblock}


\begin{sourceblock}{注记 6.2.6 欧拉在猜测正弦函数 sin(πz) 有上面的乘积表达形式后, 即:}
{\small\sloppy
\noindent 注记 6.2.6 欧拉在猜测正弦函数 sin(πz) 有上面的乘积表达形式后, 即:\par
\noindent Y z2\par
\noindent sin(πz) = πz 1− 2 .\par
\noindent n\par
\noindent n≥1\par
\noindent 随后又考虑 sin(πz) 的 Taylor 展开式:\par
\noindent ∞\par
\noindent X π 2n−1 z 2n−1\par
\noindent sin (πz) = (−1)n−1 . (6.33)\par
\noindent (2n − 1)!\par
\noindent n=1\par
\noindent 124 第 6 章 Bernoulli 数与 Bernoulli 多项式\par
\noindent 并由之得到\par
\noindent Y z2\par
\noindent ∞\par
\noindent X π 2n−1 z 2n−1\par
\noindent πz 1− 2 = (−1)n−1 .\par
\noindent n (2n − 1)!\par
}
\end{sourceblock}


\begin{explainblock}
注记通常是在补充术语、记号、历史或一个容易忽略的约定。它未必给新公式，但常常决定你后面会不会把符号读错。

\smallskip
\textbf{初学者检查点：}
\begin{itemize}[leftmargin=2em,itemsep=0.25em]
\item 记下新符号的读法。
\item 确认这个约定是否只在本章使用。
\item 把注记里的限制条件补到前面的定义旁边。
\end{itemize}

\smallskip
\textbf{本章读法提醒：} 第六章把 Bernoulli 数放回生成函数里看，不要把它们当成孤立常数。
\end{explainblock}


\begin{sourceblock}{第 6 章续读}
{\small\sloppy
\noindent n≥1 n=1\par
\noindent 比较两边 z 3 项的系数可得:\par
\noindent z2 z2 z2 1 1 π3\par
\noindent [z ]πz 1 − 2\par
\noindent 1− 2 1 − 2 · · · = −π + + · · · = − .\par
\noindent 所以\par
\noindent ∞\par
\noindent X 1 π2\par
\noindent ζ(2) = 2\par
\noindent = .\par
\noindent n 6\par
\noindent n=1\par
\noindent 当 k ≥ 1 时, 由定理 6.2.5 可得\par
\noindent 2(2k)!\par
\noindent B2k = (−1)k−1 ζ(2k), (6.34)\par
\noindent (2π)2k\par
\noindent 因为 ζ(2k) 在 k ≥ 1 总是正的, 因此我们就给出了性质 6.1.6 的证明.\par
}
\end{sourceblock}


\begin{explainblock}
这一段是原书论述链的一部分。读的时候把它当作桥梁：它通常负责把上一个定义、例子或定理，引到下一个计数方法。

\smallskip
\textbf{初学者检查点：}
\begin{itemize}[leftmargin=2em,itemsep=0.25em]
\item 找出这一段承接的上一个概念。
\item 找出它准备引出的下一个概念。
\item 把中间的关键等式或构造单独抄出来。
\end{itemize}

\smallskip
\textbf{本章读法提醒：} 第六章把 Bernoulli 数放回生成函数里看，不要把它们当成孤立常数。
\end{explainblock}


\begin{sourceblock}{推论 6.2.7 当 k → ∞ 时,}
{\small\sloppy
\noindent 推论 6.2.7 当 k → ∞ 时,\par
\noindent 2k\par
\noindent √ k\par
\noindent |B2k | ∼ 4 πk .\par
\noindent πe\par
\noindent 证明: 当 k 趋于无穷大时, 显然 ζ(2k) 趋于 1. 因此, 对式子 (6.34) 应用 Stirling\par
\noindent 公式: k\par
\noindent √ k\par
\noindent k! ∼ 2πk ,\par
\noindent e\par
\noindent 可得\par
\noindent 2k 2k\par
\noindent |B2k | ∼ 4πk = 4 πk .\par
\noindent (2π)2k e πe\par
}
\end{sourceblock}


\begin{explainblock}
定理段不要只看结论，要看它把哪两个计数方式连在一起。组合学证明最常见的技巧是：左边按一种标准数，右边按另一种标准数，然后说明两边数的是同一批对象。

\smallskip
\textbf{初学者检查点：}
\begin{itemize}[leftmargin=2em,itemsep=0.25em]
\item 先复述定理的假设，不要把适用范围扩大。
\item 找出证明中的基本计数原则：乘法、加法、双射、递推、容斥或生成函数。
\item 检查最后的等式化简是否和前面的对象解释一致。
\end{itemize}

\smallskip
\textbf{本章读法提醒：} 第六章把 Bernoulli 数放回生成函数里看，不要把它们当成孤立常数。
\end{explainblock}


\begin{sourceblock}{§ 6.3 Euler–Maclaurin 公式及应用}
{\small\sloppy
\noindent § 6.3. Euler–Maclaurin 公式及应用\par
\noindent Euler–Maclaurin 公式 (也简称为 Euler 公式) 是数值积分, 渐近展开和求和问题\par
\noindent 中的重要公式.\par
}
\end{sourceblock}


\begin{explainblock}
这一节先不要急着背公式。先问三个问题：对象是什么，哪些限制条件不能丢，最后到底是在数所有对象、还是在数某个统计量的分布。把这三件事说清楚，后面的公式才会有位置。

\smallskip
\textbf{初学者检查点：}
\begin{itemize}[leftmargin=2em,itemsep=0.25em]
\item 把本节出现的新对象写成一句话定义。
\item 把每个公式左边和右边分别翻译成“在数什么”。
\item 找一个最小非平凡例子，手工列出所有对象。
\end{itemize}

\smallskip
\textbf{本章读法提醒：} 第六章把 Bernoulli 数放回生成函数里看，不要把它们当成孤立常数。
\end{explainblock}


\begin{sourceblock}{定理 6.3.1 (Euler–Maclaurin) 设整数 a ≤ b, M 为大于 1 的自然数. 若 f (x) 在}
{\small\sloppy
\noindent 定理 6.3.1 (Euler–Maclaurin) 设整数 a ≤ b, M 为大于 1 的自然数. 若 f (x) 在\par
\noindent [a, b] 上存在 M 阶导数, 则\par
\noindent X\par
\noindent b Z b X Bk+1 M −1\par
\noindent f (i) = f (x)dx + (f (a) + f (b)) + f (k) (b) − f (k) (a)\par
\noindent a 2 (k + 1)!\par
\noindent i=a k=1\par
\noindent Z b\par
\noindent (−1)M\par
\noindent − BM (x − [x])f (M ) (x)dx, (6.35)\par
\noindent M! a\par
\noindent 其中 Bk+1 是 Bernoulli 数, BM (x) 是 Bernoulli 多项式.\par
\noindent 6.3 Euler–Maclaurin 公式及应用 125\par
\noindent 证明: 设 g(x) 在 [0, 1] 上存在 M 阶导数. 由\par
\noindent B1 (x) = x −\par
\noindent 可知 B1′ (x) = 1, 所以\par
\noindent Z 1 Z 1\par
\noindent g(x)dx = g(x)dB1 (x).\par
}
\end{sourceblock}


\begin{explainblock}
定理段不要只看结论，要看它把哪两个计数方式连在一起。组合学证明最常见的技巧是：左边按一种标准数，右边按另一种标准数，然后说明两边数的是同一批对象。

\smallskip
\textbf{初学者检查点：}
\begin{itemize}[leftmargin=2em,itemsep=0.25em]
\item 先复述定理的假设，不要把适用范围扩大。
\item 找出证明中的基本计数原则：乘法、加法、双射、递推、容斥或生成函数。
\item 检查最后的等式化简是否和前面的对象解释一致。
\end{itemize}

\smallskip
\textbf{本章读法提醒：} 第六章把 Bernoulli 数放回生成函数里看，不要把它们当成孤立常数。
\end{explainblock}


\begin{sourceblock}{第 6 章续读}
{\small\sloppy
\noindent 由分部积分可得\par
\noindent Z 1 Z 1\par
\noindent g(x)dB1 (x) = g(x)B1 (x) − B1 (x)g ′ (x)dx\par
\noindent Z 1\par
\noindent = (g(1) + g(0)) − B1 (x)g ′ (x)dx.\par
\noindent 由 (6.14) 知当 k ≥ 0 时,\par
\noindent ′\par
\noindent Bk+1 (x)\par
\noindent Bk (x) = ,\par
\noindent k+1\par
\noindent 所以\par
\noindent Z 1 Z 1\par
\noindent g(x)dx = (g(1) + g(0)) − B2 (x)g ′ (x) + B2 (x)g ′′ (x)dx\par
\noindent = (g(1) + g(0)) − B2 (x)g ′ (x)\par
\noindent Z 1\par
\noindent + B3 (x)g ′′ (x) − B3 (x)g ′′′ (x)dx\par
\noindent 2·3 0 2·3 0 (6.36)\par
\noindent M −1\par
}
\end{sourceblock}


\begin{explainblock}
这一段是原书论述链的一部分。读的时候把它当作桥梁：它通常负责把上一个定义、例子或定理，引到下一个计数方法。

\smallskip
\textbf{初学者检查点：}
\begin{itemize}[leftmargin=2em,itemsep=0.25em]
\item 找出这一段承接的上一个概念。
\item 找出它准备引出的下一个概念。
\item 把中间的关键等式或构造单独抄出来。
\end{itemize}

\smallskip
\textbf{本章读法提醒：} 第六章把 Bernoulli 数放回生成函数里看，不要把它们当成孤立常数。
\end{explainblock}


\begin{sourceblock}{第 6 章续读}
{\small\sloppy
\noindent X (−1)k\par
\noindent = (g(1) + g(0)) + Bk+1 (x)g (k) (x)\par
\noindent k=1\par
\noindent Z 1\par
\noindent (−1)M\par
\noindent + BM (x)g (M ) (x)dx\par
\noindent M! 0\par
\noindent 再由式 (6.12) 及 (−1)k Bk = Bk 知\par
\noindent X\par
\noindent M −1\par
\noindent (−1)k 1 X (−1)k\par
\noindent M −1\par
\noindent Bk+1 (x)g (k) (x) = Bk+1 (1)g (k) (1) − Bk+1 (0)g (k) (0)\par
\noindent (k + 1)! 0 (k + 1)!\par
\noindent k=1 k=1\par
\noindent Bk+1 (k)\par
\noindent X\par
\noindent M −1\par
}
\end{sourceblock}


\begin{explainblock}
这一段是原书论述链的一部分。读的时候把它当作桥梁：它通常负责把上一个定义、例子或定理，引到下一个计数方法。

\smallskip
\textbf{初学者检查点：}
\begin{itemize}[leftmargin=2em,itemsep=0.25em]
\item 找出这一段承接的上一个概念。
\item 找出它准备引出的下一个概念。
\item 把中间的关键等式或构造单独抄出来。
\end{itemize}

\smallskip
\textbf{本章读法提醒：} 第六章把 Bernoulli 数放回生成函数里看，不要把它们当成孤立常数。
\end{explainblock}


\begin{sourceblock}{第 6 章续读}
{\small\sloppy
\noindent =− g (1) − g (k) (0) ,\par
\noindent (k + 1)!\par
\noindent k=1\par
\noindent 所以\par
\noindent Z 1\par
\noindent Bk+1 (k)\par
\noindent X\par
\noindent M −1\par
\noindent g(x)dx = (g(1) + g(0)) − g (1) − g (k) (0)\par
\noindent k=1\par
\noindent 126 第 6 章 Bernoulli 数与 Bernoulli 多项式\par
\noindent Z 1\par
\noindent (−1)M\par
\noindent + BM (x)g (M ) (x)dx,\par
\noindent M! 0\par
\noindent 故\par
\noindent Z 1\par
\noindent Bk+1 (k)\par
}
\end{sourceblock}


\begin{explainblock}
这一段是原书论述链的一部分。读的时候把它当作桥梁：它通常负责把上一个定义、例子或定理，引到下一个计数方法。

\smallskip
\textbf{初学者检查点：}
\begin{itemize}[leftmargin=2em,itemsep=0.25em]
\item 找出这一段承接的上一个概念。
\item 找出它准备引出的下一个概念。
\item 把中间的关键等式或构造单独抄出来。
\end{itemize}

\smallskip
\textbf{本章读法提醒：} 第六章把 Bernoulli 数放回生成函数里看，不要把它们当成孤立常数。
\end{explainblock}


\begin{sourceblock}{第 6 章续读}
{\small\sloppy
\noindent X\par
\noindent M −1\par
\noindent (g(1) + g(0)) = g(x)dx + g (1) − g (k) (0)\par
\noindent k=1\par
\noindent Z\par
\noindent (−1)M 1\par
\noindent − BM (x)g (M ) (x)dx.\par
\noindent M! 0\par
\noindent 现不妨设 f (x + i) = g(x), 其中 a ≤ i ≤ b − 1, 则\par
\noindent Z i+1\par
\noindent Bk+1 (k)\par
\noindent X\par
\noindent M −1\par
\noindent (f (i + 1) + f (i)) = f (x)dx + f (i + 1) − f (k) (i)\par
\noindent k=1\par
\noindent Z\par
\noindent (−1)M 1\par
\noindent − BM (x)f (M ) (i + x)dx.\par
}
\end{sourceblock}


\begin{explainblock}
这一段是原书论述链的一部分。读的时候把它当作桥梁：它通常负责把上一个定义、例子或定理，引到下一个计数方法。

\smallskip
\textbf{初学者检查点：}
\begin{itemize}[leftmargin=2em,itemsep=0.25em]
\item 找出这一段承接的上一个概念。
\item 找出它准备引出的下一个概念。
\item 把中间的关键等式或构造单独抄出来。
\end{itemize}

\smallskip
\textbf{本章读法提醒：} 第六章把 Bernoulli 数放回生成函数里看，不要把它们当成孤立常数。
\end{explainblock}


\begin{sourceblock}{第 6 章续读}
{\small\sloppy
\noindent M! 0\par
\noindent 因为\par
\noindent Z 1 Z i+1 Z i+1\par
\noindent BM (x)f (M ) (x + i)dx = BM (x − i)f (M ) (x)dx = BM (x − [x])f (M ) (x)dx,\par
\noindent 所以\par
\noindent Z i+1\par
\noindent Bk+1 (k)\par
\noindent X\par
\noindent M −1\par
\noindent (f (i + 1) + f (i)) = f (x)dx + f (i + 1) − f (k) (i)\par
\noindent k=1\par
\noindent Z\par
\noindent (−1)M i+1\par
\noindent − BM (x − [x])f (M ) (x)dx, (6.37)\par
\noindent M! i\par
\noindent 其中 [x] 表示不超过 x 的最大正整数.\par
\noindent 若在 (6.37) 中分别取 i = a, a + 1, . . . , b − 1, 然后相加得\par
\noindent Z b\par
}
\end{sourceblock}


\begin{explainblock}
这一段是原书论述链的一部分。读的时候把它当作桥梁：它通常负责把上一个定义、例子或定理，引到下一个计数方法。

\smallskip
\textbf{初学者检查点：}
\begin{itemize}[leftmargin=2em,itemsep=0.25em]
\item 找出这一段承接的上一个概念。
\item 找出它准备引出的下一个概念。
\item 把中间的关键等式或构造单独抄出来。
\end{itemize}

\smallskip
\textbf{本章读法提醒：} 第六章把 Bernoulli 数放回生成函数里看，不要把它们当成孤立常数。
\end{explainblock}


\begin{sourceblock}{第 6 章续读}
{\small\sloppy
\noindent Bk+1 (k)\par
\noindent X\par
\noindent b X\par
\noindent M −1\par
\noindent f (i) − (f (a) + f (b)) = f (x)dx + f (b) − f (k) (a)\par
\noindent i=a k=1\par
\noindent Z\par
\noindent (−1)M b\par
\noindent − BM (x − [x])f (M ) (x)dx, (6.38)\par
\noindent M! a\par
\noindent 移项立得 (6.35).\par
\noindent 当 k ≥ 1 时, B2k+1 = 0. 在 (6.35) 中令 M = 2m, 我们得到了下面的定理.\par
}
\end{sourceblock}


\begin{explainblock}
这一段是原书论述链的一部分。读的时候把它当作桥梁：它通常负责把上一个定义、例子或定理，引到下一个计数方法。

\smallskip
\textbf{初学者检查点：}
\begin{itemize}[leftmargin=2em,itemsep=0.25em]
\item 找出这一段承接的上一个概念。
\item 找出它准备引出的下一个概念。
\item 把中间的关键等式或构造单独抄出来。
\end{itemize}

\smallskip
\textbf{本章读法提醒：} 第六章把 Bernoulli 数放回生成函数里看，不要把它们当成孤立常数。
\end{explainblock}


\begin{sourceblock}{定理 6.3.2 设整数 a ≥ b, m 为大于 1 的自然数. 若 f (x) 在 [a, b] 上存在 2m 阶}
{\small\sloppy
\noindent 定理 6.3.2 设整数 a ≥ b, m 为大于 1 的自然数. 若 f (x) 在 [a, b] 上存在 2m 阶\par
\noindent 6.3 Euler–Maclaurin 公式及应用 127\par
\noindent 导数, 则\par
\noindent Z b\par
\noindent X X B2k (2k−1)\par
\noindent b m\par
\noindent f (i) = f (x)dx + (f (a) + f (b)) + f (b) − f (2k−1) (a)\par
\noindent a 2 (2k)!\par
\noindent i=a k=1\par
\noindent Z b\par
\noindent − B2m (x − [x])f (2m) (x)dx.\par
\noindent (2m)! a\par
\noindent 特别地, 取 a = 1, b = n, 我们得到了下面的推论.\par
}
\end{sourceblock}


\begin{explainblock}
定理段不要只看结论，要看它把哪两个计数方式连在一起。组合学证明最常见的技巧是：左边按一种标准数，右边按另一种标准数，然后说明两边数的是同一批对象。

\smallskip
\textbf{初学者检查点：}
\begin{itemize}[leftmargin=2em,itemsep=0.25em]
\item 先复述定理的假设，不要把适用范围扩大。
\item 找出证明中的基本计数原则：乘法、加法、双射、递推、容斥或生成函数。
\item 检查最后的等式化简是否和前面的对象解释一致。
\end{itemize}

\smallskip
\textbf{本章读法提醒：} 第六章把 Bernoulli 数放回生成函数里看，不要把它们当成孤立常数。
\end{explainblock}


\begin{sourceblock}{推论 6.3.3 设 m 为大于 1 的自然数. 若 f (x) 在 [a, b] 上存在 2m 阶导数, 则}
{\small\sloppy
\noindent 推论 6.3.3 设 m 为大于 1 的自然数. 若 f (x) 在 [a, b] 上存在 2m 阶导数, 则\par
\noindent Z n\par
\noindent X X B2k (2k−1)\par
\noindent n m\par
\noindent f (i) = f (x)dx + (f (1) + f (n)) + f (n) − f (2k−1) (1)\par
\noindent i=1 k=1\par
\noindent Z n\par
\noindent − B2m (x − [x])f (2m) (x)dx, (6.39)\par
\noindent (2m)! 1\par
\noindent 下面作为 Euler–Maclaurin 公式的一个应用, 我们给出 Stirling 公式的一个证明.\par
}
\end{sourceblock}


\begin{explainblock}
定理段不要只看结论，要看它把哪两个计数方式连在一起。组合学证明最常见的技巧是：左边按一种标准数，右边按另一种标准数，然后说明两边数的是同一批对象。

\smallskip
\textbf{初学者检查点：}
\begin{itemize}[leftmargin=2em,itemsep=0.25em]
\item 先复述定理的假设，不要把适用范围扩大。
\item 找出证明中的基本计数原则：乘法、加法、双射、递推、容斥或生成函数。
\item 检查最后的等式化简是否和前面的对象解释一致。
\end{itemize}

\smallskip
\textbf{本章读法提醒：} 第六章把 Bernoulli 数放回生成函数里看，不要把它们当成孤立常数。
\end{explainblock}


\begin{sourceblock}{定理 6.3.4 (Striling 公式) 当 n 趋于无穷时,}
{\small\sloppy
\noindent 定理 6.3.4 (Striling 公式) 当 n 趋于无穷时,\par
\noindent √ n n\par
\noindent n! ∼ 2πn . (6.40)\par
\noindent e\par
\noindent 证明: 取 f (x) = log x, 代入到 (6.39) 中可得\par
\noindent X\par
\noindent n Z n X B2k m\par
\noindent log i = log n! = log xdx + log n + · (2k − 2)! −1\par
\noindent i=1 k=1\par
\noindent Z n\par
\noindent + dx,\par
\noindent 2m 1 x2m\par
\noindent 所以\par
\noindent X\par
\noindent m−1\par
\noindent log n! = n + log n − n + 1 + · 2k−1 + R2m + Km , (6.41)\par
\noindent k=1\par
\noindent 其中\par
}
\end{sourceblock}


\begin{explainblock}
定理段不要只看结论，要看它把哪两个计数方式连在一起。组合学证明最常见的技巧是：左边按一种标准数，右边按另一种标准数，然后说明两边数的是同一批对象。

\smallskip
\textbf{初学者检查点：}
\begin{itemize}[leftmargin=2em,itemsep=0.25em]
\item 先复述定理的假设，不要把适用范围扩大。
\item 找出证明中的基本计数原则：乘法、加法、双射、递推、容斥或生成函数。
\item 检查最后的等式化简是否和前面的对象解释一致。
\end{itemize}

\smallskip
\textbf{本章读法提醒：} 第六章把 Bernoulli 数放回生成函数里看，不要把它们当成孤立常数。
\end{explainblock}


\begin{sourceblock}{第 6 章续读}
{\small\sloppy
\noindent B2m\par
\noindent R2m = n−2m+1 − Ω2m (n),\par
\noindent 2m(2m − 1)\par
\noindent Z ∞\par
\noindent Ω2m (n) = dx,\par
\noindent 2m n x2m\par
\noindent 及\par
\noindent Z ∞ X m\par
\noindent 1 B2m (x − [x]) B2k\par
\noindent Km = dx − .\par
\noindent 2m 1 x2m 2k(2k − 1)\par
\noindent k=1\par
\noindent 128 第 6 章 Bernoulli 数与 Bernoulli 多项式\par
\noindent 显然 Z ∞\par
\noindent |B2m | |B2m |\par
\noindent |Ω2m (n)| ≤ x−2m dx = n−2m+1 ,\par
\noindent 2m n 2m(2m − 1)\par
\noindent 故 R2m 与 B2m 同号, 不妨设\par
}
\end{sourceblock}


\begin{explainblock}
这一段是原书论述链的一部分。读的时候把它当作桥梁：它通常负责把上一个定义、例子或定理，引到下一个计数方法。

\smallskip
\textbf{初学者检查点：}
\begin{itemize}[leftmargin=2em,itemsep=0.25em]
\item 找出这一段承接的上一个概念。
\item 找出它准备引出的下一个概念。
\item 把中间的关键等式或构造单独抄出来。
\end{itemize}

\smallskip
\textbf{本章读法提醒：} 第六章把 Bernoulli 数放回生成函数里看，不要把它们当成孤立常数。
\end{explainblock}


\begin{sourceblock}{第 6 章续读}
{\small\sloppy
\noindent |B2m |\par
\noindent R2m = θm n−2m+1 , (6.42)\par
\noindent 2m(2m − 1)\par
\noindent 则 θm > 0. 又因为 Km 与 n 的值无关, 因此,\par
\noindent lim log n! − n + log n + n − 1 = Km ,\par
\noindent n→∞ 2\par
\noindent 所以 Km 是个常量, 不妨设 Km = K.\par
\noindent 现将 (6.41) 中 m 替换为 m + 1 可得\par
\noindent X\par
\noindent m\par
\noindent log n! = n+ log n − n + 1 + + R2m+2 + K,\par
\noindent k=1\par
\noindent 对比 (6.41) 可得\par
\noindent B2m\par
\noindent R2m = n−2m+1 + R2m+2 ,\par
\noindent 2m(2m − 1)\par
\noindent 由 (6.42) 及\par
\noindent B2m+2\par
}
\end{sourceblock}


\begin{explainblock}
这一段是原书论述链的一部分。读的时候把它当作桥梁：它通常负责把上一个定义、例子或定理，引到下一个计数方法。

\smallskip
\textbf{初学者检查点：}
\begin{itemize}[leftmargin=2em,itemsep=0.25em]
\item 找出这一段承接的上一个概念。
\item 找出它准备引出的下一个概念。
\item 把中间的关键等式或构造单独抄出来。
\end{itemize}

\smallskip
\textbf{本章读法提醒：} 第六章把 Bernoulli 数放回生成函数里看，不要把它们当成孤立常数。
\end{explainblock}


\begin{sourceblock}{第 6 章续读}
{\small\sloppy
\noindent R2m+2 = θm+1 n−2m−1 , θm+1 > 0,\par
\noindent 2(m + 1)(2m + 1)\par
\noindent 可得\par
\noindent B2m+2 2m(2m − 1)\par
\noindent θm = 1 − θm+1 2\par
\noindent ,\par
\noindent B2m n (2m + 2)(2m + 1)\par
\noindent 因此对任意 m ≥ 1,\par
\noindent 现不妨取 m = 1, 则\par
\noindent log n! = K + n + log n − n + 1 + .\par
\noindent 故\par
\noindent θ1\par
\noindent n! = eK nn+ 2 e−n+1 e 12n .\par
\noindent 由 Wallis 公式:\par
\noindent 24n (n!)4 π\par
\noindent lim = ,\par
\noindent n→∞ (2n + 1)((2n)!)2 2\par
\noindent 知\par
}
\end{sourceblock}


\begin{explainblock}
这一段是原书论述链的一部分。读的时候把它当作桥梁：它通常负责把上一个定义、例子或定理，引到下一个计数方法。

\smallskip
\textbf{初学者检查点：}
\begin{itemize}[leftmargin=2em,itemsep=0.25em]
\item 找出这一段承接的上一个概念。
\item 找出它准备引出的下一个概念。
\item 把中间的关键等式或构造单独抄出来。
\end{itemize}

\smallskip
\textbf{本章读法提醒：} 第六章把 Bernoulli 数放回生成函数里看，不要把它们当成孤立常数。
\end{explainblock}


\begin{sourceblock}{第 6 章续读}
{\small\sloppy
\noindent 24n (eK nn+ 2 e−n+1 e 12n )4 n2 π\par
\noindent lim ′ = lim e2K+2 = ,\par
\noindent n→∞ 1 θ1 n→∞ 2n(2n + 1) 2\par
\noindent (2n + 1)(eK (2n)2n+ 2 e−2n+1 e 24n )2\par
\noindent 6.3 Euler–Maclaurin 公式及应用 129\par
\noindent 所以\par
\noindent √\par
\noindent K = log 2π − 1.\par
\noindent 故\par
\noindent √ θ1\par
\noindent 2πnn+ 2 e−n e 12n ,\par
\noindent n! =\par
\noindent 所以当 n 趋于无穷时, n n\par
\noindent √\par
\noindent n! ∼ 2πn .\par
\noindent e\par
\noindent 下面我们再给个例子进行说明.\par
}
\end{sourceblock}


\begin{explainblock}
这一段是原书论述链的一部分。读的时候把它当作桥梁：它通常负责把上一个定义、例子或定理，引到下一个计数方法。

\smallskip
\textbf{初学者检查点：}
\begin{itemize}[leftmargin=2em,itemsep=0.25em]
\item 找出这一段承接的上一个概念。
\item 找出它准备引出的下一个概念。
\item 把中间的关键等式或构造单独抄出来。
\end{itemize}

\smallskip
\textbf{本章读法提醒：} 第六章把 Bernoulli 数放回生成函数里看，不要把它们当成孤立常数。
\end{explainblock}


\begin{sourceblock}{例 6.3.1 利用 Euler-Maclaurin 公式计算}
{\small\sloppy
\noindent 例 6.3.1 利用 Euler-Maclaurin 公式计算\par
\noindent S2m (n) = 12m + 22m + 32m + · · · + n2m .\par
\noindent 解: 不妨考虑 f (x) = x2m , 则\par
\noindent (\par
\noindent (2m)k x2m−k , 0 ≤ k ≤ 2m,\par
\noindent f (k) (x) =\par
\noindent 0, k ≥ 2m + 1.\par
\noindent 应用 Euler–Maclaurin 公式可得\par
\noindent X\par
\noindent n Z n X B2k m\par
\noindent S2m (n) = f (i) = f (x)dx + (f (0) + f (n)) + (f (2k−1) (n) − f (2k−1) (0))\par
\noindent i=0 k=1\par
\noindent m\par
\noindent X\par
\noindent n2m+1 n2m 1 2m + 1\par
\noindent = + + B2k n2m−2k+1\par
\noindent 2m + 1 2 2m + 1 2k\par
\noindent k=1\par
}
\end{sourceblock}


\begin{explainblock}
例题是学习这本书最重要的部分。不要只看答案，要把每个限制条件变成一个计数动作：先安排谁、再插入谁、有没有重复、需不需要除以对称性。

\smallskip
\textbf{初学者检查点：}
\begin{itemize}[leftmargin=2em,itemsep=0.25em]
\item 先写出没有限制时的总数。
\item 逐条加入限制条件，观察总数怎样变化。
\item 最后检查有没有把同一种方案重复算了多次。
\end{itemize}

\smallskip
\textbf{本章读法提醒：} 第六章把 Bernoulli 数放回生成函数里看，不要把它们当成孤立常数。
\end{explainblock}


\begin{sourceblock}{第 6 章续读}
{\small\sloppy
\noindent 1 2m + 1 2m+1 1 2m + 1\par
\noindent = B0 n + B1 n2m\par
\noindent 2m + 1 0 2m + 1 1\par
\noindent X 2m+1\par
\noindent + Bk n2m−k+1 − n2m\par
\noindent 2m + 1 k\par
\noindent k=2\par
\noindent X 2m+1\par
\noindent = Bk n2m−k+1 − n2m\par
\noindent 2m + 1 k\par
\noindent k=0\par
\noindent = B2m+1 (n) − n2m ,\par
\noindent 2m + 1\par
\noindent 所以\par
\noindent B2m+1 (n)\par
\noindent S2m (n − 1) = ,\par
\noindent 2m + 1\par
\noindent 这与 (6.13) 一致 (因为当 m ≥ 1 时, B2m+1 (1) = B2m+1 = 0).\par
}
\end{sourceblock}


\begin{explainblock}
这一段是原书论述链的一部分。读的时候把它当作桥梁：它通常负责把上一个定义、例子或定理，引到下一个计数方法。

\smallskip
\textbf{初学者检查点：}
\begin{itemize}[leftmargin=2em,itemsep=0.25em]
\item 找出这一段承接的上一个概念。
\item 找出它准备引出的下一个概念。
\item 把中间的关键等式或构造单独抄出来。
\end{itemize}

\smallskip
\textbf{本章读法提醒：} 第六章把 Bernoulli 数放回生成函数里看，不要把它们当成孤立常数。
\end{explainblock}


\begin{sourceblock}{第 6 章续读}
{\small\sloppy
\noindent 130 第 6 章 Bernoulli 数与 Bernoulli 多项式\par
\noindent 6.3 Euler–Maclaurin 公式及应用 131\par
}
\end{sourceblock}


\begin{explainblock}
这一段是原书论述链的一部分。读的时候把它当作桥梁：它通常负责把上一个定义、例子或定理，引到下一个计数方法。

\smallskip
\textbf{初学者检查点：}
\begin{itemize}[leftmargin=2em,itemsep=0.25em]
\item 找出这一段承接的上一个概念。
\item 找出它准备引出的下一个概念。
\item 把中间的关键等式或构造单独抄出来。
\end{itemize}

\smallskip
\textbf{本章读法提醒：} 第六章把 Bernoulli 数放回生成函数里看，不要把它们当成孤立常数。
\end{explainblock}


\begin{sourceblock}{习 题 六}
{\small\sloppy
\noindent 习 题 六\par
}
\end{sourceblock}


\begin{explainblock}
习题段不要先追求技巧。先给题目分类：它是在问排列、组合、递推、生成函数、容斥，还是结构计数。分类正确，工具通常就只剩一两个候选。

\smallskip
\textbf{初学者检查点：}
\begin{itemize}[leftmargin=2em,itemsep=0.25em]
\item 判断题目所属工具。
\item 列出变量和限制条件。
\item 先做小规模 n=2、3、4 的例子，确认直觉。
\end{itemize}

\smallskip
\textbf{本章读法提醒：} 第六章把 Bernoulli 数放回生成函数里看，不要把它们当成孤立常数。
\end{explainblock}


\begin{sourceblock}{习题 6.1 计算前 8 个 Bernoulli 数并求}
{\small\sloppy
\noindent 习题 6.1 计算前 8 个 Bernoulli 数并求\par
\noindent (1)\par
\noindent ζ(8) = 1 + 8\par
\noindent + 8 + ··· + 8 + ··· ,\par
\noindent (2)\par
\noindent 18 + 2 8 + 3 8 + · · · + n8 .\par
}
\end{sourceblock}


\begin{explainblock}
习题段不要先追求技巧。先给题目分类：它是在问排列、组合、递推、生成函数、容斥，还是结构计数。分类正确，工具通常就只剩一两个候选。

\smallskip
\textbf{初学者检查点：}
\begin{itemize}[leftmargin=2em,itemsep=0.25em]
\item 判断题目所属工具。
\item 列出变量和限制条件。
\item 先做小规模 n=2、3、4 的例子，确认直觉。
\end{itemize}

\smallskip
\textbf{本章读法提醒：} 第六章把 Bernoulli 数放回生成函数里看，不要把它们当成孤立常数。
\end{explainblock}


\begin{sourceblock}{习题 6.2 证明当 n ≥ 2 时, Bernoulli 数满足}
{\small\sloppy
\noindent 习题 6.2 证明当 n ≥ 2 时, Bernoulli 数满足\par
\noindent n\par
\noindent X\par
\noindent 2n\par
\noindent (1 − 2n)B2n = B2m B2n−2m .\par
\noindent 2m\par
\noindent m=0\par
}
\end{sourceblock}


\begin{explainblock}
习题段不要先追求技巧。先给题目分类：它是在问排列、组合、递推、生成函数、容斥，还是结构计数。分类正确，工具通常就只剩一两个候选。

\smallskip
\textbf{初学者检查点：}
\begin{itemize}[leftmargin=2em,itemsep=0.25em]
\item 判断题目所属工具。
\item 列出变量和限制条件。
\item 先做小规模 n=2、3、4 的例子，确认直觉。
\end{itemize}

\smallskip
\textbf{本章读法提醒：} 第六章把 Bernoulli 数放回生成函数里看，不要把它们当成孤立常数。
\end{explainblock}


\begin{sourceblock}{习题 6.3 证明当 n ≥ 0 时, Bernoulli 数与第二类 Striling 数满足下面的等式:}
{\small\sloppy
\noindent 习题 6.3 证明当 n ≥ 0 时, Bernoulli 数与第二类 Striling 数满足下面的等式:\par
\noindent X\par
\noindent n\par
\noindent (−1)k\par
\noindent Bn = k!S(n, k).\par
\noindent k+1\par
\noindent k=0\par
}
\end{sourceblock}


\begin{explainblock}
习题段不要先追求技巧。先给题目分类：它是在问排列、组合、递推、生成函数、容斥，还是结构计数。分类正确，工具通常就只剩一两个候选。

\smallskip
\textbf{初学者检查点：}
\begin{itemize}[leftmargin=2em,itemsep=0.25em]
\item 判断题目所属工具。
\item 列出变量和限制条件。
\item 先做小规模 n=2、3、4 的例子，确认直觉。
\end{itemize}

\smallskip
\textbf{本章读法提醒：} 第六章把 Bernoulli 数放回生成函数里看，不要把它们当成孤立常数。
\end{explainblock}


\begin{sourceblock}{习题 6.4 证明当 n ≥ 0 时, Bernoulli 数与第一类 Striling 数满足下面的等式:}
{\small\sloppy
\noindent 习题 6.4 证明当 n ≥ 0 时, Bernoulli 数与第一类 Striling 数满足下面的等式:\par
\noindent X\par
\noindent n\par
\noindent (−1)n n!\par
\noindent s(n, k)Bk = .\par
\noindent n+1\par
\noindent k=0\par
}
\end{sourceblock}


\begin{explainblock}
习题段不要先追求技巧。先给题目分类：它是在问排列、组合、递推、生成函数、容斥，还是结构计数。分类正确，工具通常就只剩一两个候选。

\smallskip
\textbf{初学者检查点：}
\begin{itemize}[leftmargin=2em,itemsep=0.25em]
\item 判断题目所属工具。
\item 列出变量和限制条件。
\item 先做小规模 n=2、3、4 的例子，确认直觉。
\end{itemize}

\smallskip
\textbf{本章读法提醒：} 第六章把 Bernoulli 数放回生成函数里看，不要把它们当成孤立常数。
\end{explainblock}


\begin{sourceblock}{习题 6.5 证明 Bernoulli 数 B2n 的分母是所有满足 p − 1 整除 2n 的素数 p 的}
{\small\sloppy
\noindent 习题 6.5 证明 Bernoulli 数 B2n 的分母是所有满足 p − 1 整除 2n 的素数 p 的\par
\noindent 乘积.\par
}
\end{sourceblock}


\begin{explainblock}
习题段不要先追求技巧。先给题目分类：它是在问排列、组合、递推、生成函数、容斥，还是结构计数。分类正确，工具通常就只剩一两个候选。

\smallskip
\textbf{初学者检查点：}
\begin{itemize}[leftmargin=2em,itemsep=0.25em]
\item 判断题目所属工具。
\item 列出变量和限制条件。
\item 先做小规模 n=2、3、4 的例子，确认直觉。
\end{itemize}

\smallskip
\textbf{本章读法提醒：} 第六章把 Bernoulli 数放回生成函数里看，不要把它们当成孤立常数。
\end{explainblock}


\begin{sourceblock}{习题 6.6 证明当 n ≥ 0 时,}
{\small\sloppy
\noindent 习题 6.6 证明当 n ≥ 0 时,\par
\noindent Bn (x) = (−1)n Bn (1 − x).\par
}
\end{sourceblock}


\begin{explainblock}
习题段不要先追求技巧。先给题目分类：它是在问排列、组合、递推、生成函数、容斥，还是结构计数。分类正确，工具通常就只剩一两个候选。

\smallskip
\textbf{初学者检查点：}
\begin{itemize}[leftmargin=2em,itemsep=0.25em]
\item 判断题目所属工具。
\item 列出变量和限制条件。
\item 先做小规模 n=2、3、4 的例子，确认直觉。
\end{itemize}

\smallskip
\textbf{本章读法提醒：} 第六章把 Bernoulli 数放回生成函数里看，不要把它们当成孤立常数。
\end{explainblock}


\begin{sourceblock}{习题 6.7 证明当 n ≥ 2 时, Bernoulli 多项式满足下列的关系式:}
{\small\sloppy
\noindent 习题 6.7 证明当 n ≥ 2 时, Bernoulli 多项式满足下列的关系式:\par
\noindent n\par
\noindent X n\par
\noindent (1) Bk (x) = Bn (x + 1),\par
\noindent k\par
\noindent k=0\par
\noindent X\par
\noindent n−1\par
\noindent n\par
\noindent (2) Bk (x) = nxn−1 ,\par
\noindent k\par
\noindent k=0\par
\noindent Z x+1\par
\noindent (3) Bn (t)dt = xn .\par
\noindent x\par
}
\end{sourceblock}


\begin{explainblock}
习题段不要先追求技巧。先给题目分类：它是在问排列、组合、递推、生成函数、容斥，还是结构计数。分类正确，工具通常就只剩一两个候选。

\smallskip
\textbf{初学者检查点：}
\begin{itemize}[leftmargin=2em,itemsep=0.25em]
\item 判断题目所属工具。
\item 列出变量和限制条件。
\item 先做小规模 n=2、3、4 的例子，确认直觉。
\end{itemize}

\smallskip
\textbf{本章读法提醒：} 第六章把 Bernoulli 数放回生成函数里看，不要把它们当成孤立常数。
\end{explainblock}


\begin{sourceblock}{习题 6.8 证明}
{\small\sloppy
\noindent 习题 6.8 证明\par
\noindent n\par
\noindent X n\par
\noindent Bn (x + y) = Bk (y)xn−k .\par
\noindent k\par
\noindent k=0\par
\noindent 132 第 6 章 Bernoulli 数与 Bernoulli 多项式\par
}
\end{sourceblock}


\begin{explainblock}
习题段不要先追求技巧。先给题目分类：它是在问排列、组合、递推、生成函数、容斥，还是结构计数。分类正确，工具通常就只剩一两个候选。

\smallskip
\textbf{初学者检查点：}
\begin{itemize}[leftmargin=2em,itemsep=0.25em]
\item 判断题目所属工具。
\item 列出变量和限制条件。
\item 先做小规模 n=2、3、4 的例子，确认直觉。
\end{itemize}

\smallskip
\textbf{本章读法提醒：} 第六章把 Bernoulli 数放回生成函数里看，不要把它们当成孤立常数。
\end{explainblock}


\begin{sourceblock}{习题 6.9 证明 Bernoulli 多项式的乘法公式}
{\small\sloppy
\noindent 习题 6.9 证明 Bernoulli 多项式的乘法公式\par
\noindent X\par
\noindent m−1\par
\noindent n−1 k\par
\noindent Bn (mx) = m Bn x+ ,\par
\noindent m\par
\noindent k=0\par
\noindent 其中 m 是任意正整数. 又由此导出下列公式\par
\noindent X\par
\noindent m−1\par
\noindent k 1\par
\noindent Bn =− 1− Bn ,\par
\noindent m mn−1\par
\noindent k=1\par
\noindent 其特殊情形为\par
\noindent Bn = − 1 − n−1 Bn .\par
}
\end{sourceblock}


\begin{explainblock}
习题段不要先追求技巧。先给题目分类：它是在问排列、组合、递推、生成函数、容斥，还是结构计数。分类正确，工具通常就只剩一两个候选。

\smallskip
\textbf{初学者检查点：}
\begin{itemize}[leftmargin=2em,itemsep=0.25em]
\item 判断题目所属工具。
\item 列出变量和限制条件。
\item 先做小规模 n=2、3、4 的例子，确认直觉。
\end{itemize}

\smallskip
\textbf{本章读法提醒：} 第六章把 Bernoulli 数放回生成函数里看，不要把它们当成孤立常数。
\end{explainblock}


\begin{sourceblock}{习题 6.10 设 Bn (x) 为 Bernoulli 数, 证明多项式 B(x − [x]) 是周期为 1 的函}
{\small\sloppy
\noindent 习题 6.10 设 Bn (x) 为 Bernoulli 数, 证明多项式 B(x − [x]) 是周期为 1 的函\par
\noindent 数, 并证明:\par
\noindent ∞\par
\noindent X cos 2πkx\par
\noindent B2n (x − [x]) = 2(−1)n+1 (2n)!\par
\noindent (2πk)2n\par
\noindent k=1\par
\noindent X∞\par
\noindent sin 2πkx\par
\noindent B2n−1 (x − [x]) = 2(−1)n (2n − 1)! .\par
\noindent (2πk)2n−1\par
\noindent k=1\par
}
\end{sourceblock}


\begin{explainblock}
习题段不要先追求技巧。先给题目分类：它是在问排列、组合、递推、生成函数、容斥，还是结构计数。分类正确，工具通常就只剩一两个候选。

\smallskip
\textbf{初学者检查点：}
\begin{itemize}[leftmargin=2em,itemsep=0.25em]
\item 判断题目所属工具。
\item 列出变量和限制条件。
\item 先做小规模 n=2、3、4 的例子，确认直觉。
\end{itemize}

\smallskip
\textbf{本章读法提醒：} 第六章把 Bernoulli 数放回生成函数里看，不要把它们当成孤立常数。
\end{explainblock}


\chapter{指数型生成函数的乘积公式和复合公式}
\begin{roadmapblock}
本章保留原书正文的主要数学骨架，并在每个定义、定理、例题和论述片段后加入解读。
阅读时不要把目标放在“背下答案”上，而是放在“看懂这一步为什么这样数”上。

\smallskip
\textbf{本章最低目标：} 能说出每个公式左边在数什么、右边在数什么，以及为什么两边相等。
\end{roadmapblock}


\begin{sourceblock}{第 7 章正文}
{\small\sloppy
\noindent 指数型生成函数的乘积公式和复合公\par
\noindent 式\par
}
\end{sourceblock}


\begin{explainblock}
这一段是原书论述链的一部分。读的时候把它当作桥梁：它通常负责把上一个定义、例子或定理，引到下一个计数方法。

\smallskip
\textbf{初学者检查点：}
\begin{itemize}[leftmargin=2em,itemsep=0.25em]
\item 找出这一段承接的上一个概念。
\item 找出它准备引出的下一个概念。
\item 把中间的关键等式或构造单独抄出来。
\end{itemize}

\smallskip
\textbf{本章读法提醒：} 第七章指数型生成函数的标签分配是核心，别和普通生成函数混用。
\end{explainblock}


\begin{sourceblock}{§ 7.1 乘积公式}
{\small\sloppy
\noindent § 7.1. 乘积公式\par
\noindent 为方便叙述, 我们只考虑复数域上形式幂级数构成的整环 C⟦x⟧, 其中 x 为不定\par
\noindent 元. 注意形式幂级数构成的整环与不定元的选取无关, 即我们总认为 C⟦x⟧ 与 C⟦y⟧\par
\noindent 是相同的. 在本章中我们始终用 Ef (x) 表示序列 \{fn \}n≥0 或 \{f (n)\}n≥0 的指数型生\par
\noindent 成函数, 在不混淆的情形下, 简记为 Ef .\par
\noindent 现考虑复数域 C 上的序列 \{an \}n≥0 和 \{bn \}n≥0 , 设 \{cn \}n≥0 定义如下:\par
\noindent n\par
\noindent X n\par
\noindent cn = ak bn−k , n ≥ 0,\par
\noindent k\par
\noindent k=0\par
\noindent 则由形式幂级数的乘法可知\par
\noindent ∞\par
\noindent X ∞\par
\noindent xn X xn\par
\noindent Ea Eb = an bn\par
\noindent n! n!\par
\noindent n=0 n=0\par
}
\end{sourceblock}


\begin{explainblock}
这一节先不要急着背公式。先问三个问题：对象是什么，哪些限制条件不能丢，最后到底是在数所有对象、还是在数某个统计量的分布。把这三件事说清楚，后面的公式才会有位置。

\smallskip
\textbf{初学者检查点：}
\begin{itemize}[leftmargin=2em,itemsep=0.25em]
\item 把本节出现的新对象写成一句话定义。
\item 把每个公式左边和右边分别翻译成“在数什么”。
\item 找一个最小非平凡例子，手工列出所有对象。
\end{itemize}

\smallskip
\textbf{本章读法提醒：} 第七章指数型生成函数的标签分配是核心，别和普通生成函数混用。
\end{explainblock}


\begin{sourceblock}{第 7 章续读}
{\small\sloppy
\noindent ∞ X\par
\noindent X n\par
\noindent ak bn−k n\par
\noindent = x\par
\noindent k! (n − k)!\par
\noindent n=0 k=0\par
\noindent n\par
\noindent ∞ X\par
\noindent X n xn\par
\noindent = ak bn−k\par
\noindent k n!\par
\noindent n=0 k=0\par
\noindent ∞\par
\noindent X xn\par
\noindent = cn ,\par
\noindent n!\par
\noindent n=0\par
\noindent 即:\par
}
\end{sourceblock}


\begin{explainblock}
这一段是原书论述链的一部分。读的时候把它当作桥梁：它通常负责把上一个定义、例子或定理，引到下一个计数方法。

\smallskip
\textbf{初学者检查点：}
\begin{itemize}[leftmargin=2em,itemsep=0.25em]
\item 找出这一段承接的上一个概念。
\item 找出它准备引出的下一个概念。
\item 把中间的关键等式或构造单独抄出来。
\end{itemize}

\smallskip
\textbf{本章读法提醒：} 第七章指数型生成函数的标签分配是核心，别和普通生成函数混用。
\end{explainblock}


\begin{sourceblock}{第 7 章续读}
{\small\sloppy
\noindent Ec = Ea · Eb .\par
\noindent 这个结果虽然简单, 但是若从组合的角度来描述它, 就能充分展现它在组合计数上的\par
\noindent 134 第 7 章 指数型生成函数的乘积公式和复合公式\par
\noindent 威力. 为此我们还需要再作一些说明.\par
\noindent 为了叙述方便, 我们先给出集合 S 弱划分的定义并采用 \#S 表示集合 S 中元素\par
\noindent 的个数.\par
}
\end{sourceblock}


\begin{explainblock}
这一段是原书论述链的一部分。读的时候把它当作桥梁：它通常负责把上一个定义、例子或定理，引到下一个计数方法。

\smallskip
\textbf{初学者检查点：}
\begin{itemize}[leftmargin=2em,itemsep=0.25em]
\item 找出这一段承接的上一个概念。
\item 找出它准备引出的下一个概念。
\item 把中间的关键等式或构造单独抄出来。
\end{itemize}

\smallskip
\textbf{本章读法提醒：} 第七章指数型生成函数的标签分配是核心，别和普通生成函数混用。
\end{explainblock}


\begin{sourceblock}{定义 7.1.1 设 T1 , T2 , . . . , Tk 是集合 X 的 k 个子集, 若它们两两不交且满足}
{\small\sloppy
\noindent 定义 7.1.1 设 T1 , T2 , . . . , Tk 是集合 X 的 k 个子集, 若它们两两不交且满足\par
\noindent T1 ∪ T2 ∪ · · · ∪ Tk = X,\par
\noindent 则称 \{T1 , T2 , . . . , Tk \} 是集合 X 的 k 部弱划分. 若将集合 X 的任意一个弱划分:\par
\noindent \{T1 , T2 , . . . , Tk \} 的不同排序视为集合 X 的不同的弱划分, 并称 (T1 , T2 , . . . , Tk ) 是集\par
\noindent 合 X 的 k 部有序弱划分.\par
}
\end{sourceblock}


\begin{explainblock}
定义段的任务是定边界。这里要特别注意参数范围、是否允许重复、是否考虑顺序、对象有没有标签。后面的每一道例题和证明，本质上都在反复使用这些边界。

\smallskip
\textbf{初学者检查点：}
\begin{itemize}[leftmargin=2em,itemsep=0.25em]
\item 圈出被定义的对象名称。
\item 圈出参数范围，例如 n、k 是否要求正整数。
\item 判断它是有序对象、无序对象、带标签对象还是结构对象。
\end{itemize}

\smallskip
\textbf{本章读法提醒：} 第七章指数型生成函数的标签分配是核心，别和普通生成函数混用。
\end{explainblock}


\begin{sourceblock}{例 7.1.1 集合 \{1, 2\}, \{3, 4\} 是 \{1, 2, 3, 4\} 的一个 2 部 (弱) 划分, 集合 \{1, 2, 3, 4\}, ∅, ∅}
{\small\sloppy
\noindent 例 7.1.1 集合 \{1, 2\}, \{3, 4\} 是 \{1, 2, 3, 4\} 的一个 2 部 (弱) 划分, 集合 \{1, 2, 3, 4\}, ∅, ∅\par
\noindent 是集合 \{1, 2, 3, 4\} 的 3 部弱划分.\par
}
\end{sourceblock}


\begin{explainblock}
例题是学习这本书最重要的部分。不要只看答案，要把每个限制条件变成一个计数动作：先安排谁、再插入谁、有没有重复、需不需要除以对称性。

\smallskip
\textbf{初学者检查点：}
\begin{itemize}[leftmargin=2em,itemsep=0.25em]
\item 先写出没有限制时的总数。
\item 逐条加入限制条件，观察总数怎样变化。
\item 最后检查有没有把同一种方案重复算了多次。
\end{itemize}

\smallskip
\textbf{本章读法提醒：} 第七章指数型生成函数的标签分配是核心，别和普通生成函数混用。
\end{explainblock}


\begin{sourceblock}{例 7.1.2 集合 \{1, 2, 3\} 的所有 2 部弱划分有:}
{\small\sloppy
\noindent 例 7.1.2 集合 \{1, 2, 3\} 的所有 2 部弱划分有:\par
\noindent ∅|123, 1|23, 2|13, 3|12,\par
\noindent 调换划分的顺序可得到它的全部 2 部有序弱划分, 即除上面外还应包括\par
\noindent 123|∅, 23|1, 13|2, 12|3.\par
\noindent 有了 (有序) 弱划分的概念, 我们便可从组合角度将两个形式幂级数的乘积结果\par
\noindent 叙述如下.\par
}
\end{sourceblock}


\begin{explainblock}
例题是学习这本书最重要的部分。不要只看答案，要把每个限制条件变成一个计数动作：先安排谁、再插入谁、有没有重复、需不需要除以对称性。

\smallskip
\textbf{初学者检查点：}
\begin{itemize}[leftmargin=2em,itemsep=0.25em]
\item 先写出没有限制时的总数。
\item 逐条加入限制条件，观察总数怎样变化。
\item 最后检查有没有把同一种方案重复算了多次。
\end{itemize}

\smallskip
\textbf{本章读法提醒：} 第七章指数型生成函数的标签分配是核心，别和普通生成函数混用。
\end{explainblock}


\begin{sourceblock}{定理 7.1.2 设 S 为任意可数集合, 函数 f, g 与 h 均是定义在非负整数上的函数,}
{\small\sloppy
\noindent 定理 7.1.2 设 S 为任意可数集合, 函数 f, g 与 h 均是定义在非负整数上的函数,\par
\noindent 且满足\par
\noindent X\par
\noindent h(\#S) = f (\#T1 )g(\#T2 ),\par
\noindent (T1 ,T2 )\par
\noindent 其中 (T1 , T2 ) 为集合 X 上的 2 部有序弱划分, 则\par
\noindent Eh = Ef Eg . (7.1)\par
\noindent n\par
\noindent 证明: 不妨设 \#S = n, 则显然有 k 个子集满足 \#T1 = k, \#T2 = n − k, 故\par
\noindent X n\par
\noindent X n\par
\noindent h(n) = f (\#T1 )g(\#T2 ) = f (k)g(n − k),\par
\noindent k\par
\noindent (T1 ,T2 ) k=0\par
\noindent 因此 (7.1) 成立.\par
\noindent 这个定理虽然是两个形式幂级数乘积的直接结果, 但却被赋予了组合意义. 假设\par
\noindent 有两种结构附加到集合 S 上, 例如结构 α 与结构 β. 若这两种结构是只与集合中元\par
\noindent 素有关的一种新的 “结合” 类型的结构, 并用 α ∪ β 表示, 并用 h(\#S) 表示集合 S\par
}
\end{sourceblock}


\begin{explainblock}
定理段不要只看结论，要看它把哪两个计数方式连在一起。组合学证明最常见的技巧是：左边按一种标准数，右边按另一种标准数，然后说明两边数的是同一批对象。

\smallskip
\textbf{初学者检查点：}
\begin{itemize}[leftmargin=2em,itemsep=0.25em]
\item 先复述定理的假设，不要把适用范围扩大。
\item 找出证明中的基本计数原则：乘法、加法、双射、递推、容斥或生成函数。
\item 检查最后的等式化简是否和前面的对象解释一致。
\end{itemize}

\smallskip
\textbf{本章读法提醒：} 第七章指数型生成函数的标签分配是核心，别和普通生成函数混用。
\end{explainblock}


\begin{sourceblock}{第 7 章续读}
{\small\sloppy
\noindent 附加这种结构的个数. 若这种结构可通过分别在 T1 和 T2 附加结构 α 和结构 β 得\par
\noindent 到, 且分别作 f (\#T1 ) 与 g(\#T2 ) 表示在集合 T1 附加结构 α 和在集合 T2 附加结构\par
\noindent β 的个数, 则显然就有\par
\noindent X\par
\noindent h(\#S) = f (\#T1 )g(\#T2 ).\par
\noindent (T1 ,T2 )\par
}
\end{sourceblock}


\begin{explainblock}
这一段是原书论述链的一部分。读的时候把它当作桥梁：它通常负责把上一个定义、例子或定理，引到下一个计数方法。

\smallskip
\textbf{初学者检查点：}
\begin{itemize}[leftmargin=2em,itemsep=0.25em]
\item 找出这一段承接的上一个概念。
\item 找出它准备引出的下一个概念。
\item 把中间的关键等式或构造单独抄出来。
\end{itemize}

\smallskip
\textbf{本章读法提醒：} 第七章指数型生成函数的标签分配是核心，别和普通生成函数混用。
\end{explainblock}


\begin{sourceblock}{例 7.1.3 给定 n 元集合 S, 设 h(\#S) 计数的是将 S 依次分成两个不交的集合}
{\small\sloppy
\noindent 例 7.1.3 给定 n 元集合 S, 设 h(\#S) 计数的是将 S 依次分成两个不交的集合\par
\noindent T1 , T2 (允许 T1 , T2 为空) 后, 将 T1 中的元素进行排列同时选取集合 T2 的一个子集\par
\noindent 的方法数, 求 Eh .\par
\noindent 解: 因为集合 T1 中元素的排序的个数为 (\#T1 )!, 集合 T2 中选取一个子集的个数\par
\noindent 为 2\#T2 , 所以不妨对任意非负整数 n, 设 f (n) = n!, g(n) = 2n . 故而\par
\noindent X\par
\noindent h(\#S) = f (\#T1 )g(\#T2 ),\par
\noindent (T1 ,T2 )\par
\noindent 所以由定理 7.1.2 可知\par
\noindent ∞\par
\noindent ! ∞\par
\noindent !\par
\noindent X xn X xn e2x\par
\noindent Eh = Ef Eg = n! · 2n = .\par
\noindent n! n! 1−x\par
\noindent n=0 n=0\par
}
\end{sourceblock}


\begin{explainblock}
例题是学习这本书最重要的部分。不要只看答案，要把每个限制条件变成一个计数动作：先安排谁、再插入谁、有没有重复、需不需要除以对称性。

\smallskip
\textbf{初学者检查点：}
\begin{itemize}[leftmargin=2em,itemsep=0.25em]
\item 先写出没有限制时的总数。
\item 逐条加入限制条件，观察总数怎样变化。
\item 最后检查有没有把同一种方案重复算了多次。
\end{itemize}

\smallskip
\textbf{本章读法提醒：} 第七章指数型生成函数的标签分配是核心，别和普通生成函数混用。
\end{explainblock}


\begin{sourceblock}{例 7.1.4 设 n 为任意正整数, Dn 表示集合 [n] 上所有错排的个数, 并规定}
{\small\sloppy
\noindent 例 7.1.4 设 n 为任意正整数, Dn 表示集合 [n] 上所有错排的个数, 并规定\par
\noindent D0 = 1, 求 ED .\par
\noindent 解: 对任意正整数 n, 设 h(n) 表示集合 [n] 上所有排列的个数, 设 f (n) = 1 及\par
\noindent g(n) = Dn , 并规定 h(0) = f (0) = g(0) = 1. 因为任意一个排列都可看作将集合 [n]\par
\noindent 分成 T1 , T2 两部分 (允许 T1 , T2 为空), 然后将 T1 中的元素视为排列的固定点, 将集\par
\noindent 合 T2 中的元素进行错位排列, 则\par
\noindent X\par
\noindent h(n) = f (\#T1 )g(\#T2 ),\par
\noindent (T1 ,T2 )\par
\noindent 因此,\par
\noindent Eh = Ef · Eg .\par
\noindent 又\par
\noindent ∞\par
\noindent X xn 1\par
\noindent Eh = n! = ,\par
\noindent n! 1−x\par
\noindent n=0\par
\noindent ∞\par
}
\end{sourceblock}


\begin{explainblock}
例题是学习这本书最重要的部分。不要只看答案，要把每个限制条件变成一个计数动作：先安排谁、再插入谁、有没有重复、需不需要除以对称性。

\smallskip
\textbf{初学者检查点：}
\begin{itemize}[leftmargin=2em,itemsep=0.25em]
\item 先写出没有限制时的总数。
\item 逐条加入限制条件，观察总数怎样变化。
\item 最后检查有没有把同一种方案重复算了多次。
\end{itemize}

\smallskip
\textbf{本章读法提醒：} 第七章指数型生成函数的标签分配是核心，别和普通生成函数混用。
\end{explainblock}


\begin{sourceblock}{第 7 章续读}
{\small\sloppy
\noindent X xn\par
\noindent Ef = = ex ,\par
\noindent n!\par
\noindent n=0\par
\noindent 故\par
\noindent Eh e−x\par
\noindent ED = Eg = = .\par
\noindent Ef 1−x\par
}
\end{sourceblock}


\begin{explainblock}
这一段是原书论述链的一部分。读的时候把它当作桥梁：它通常负责把上一个定义、例子或定理，引到下一个计数方法。

\smallskip
\textbf{初学者检查点：}
\begin{itemize}[leftmargin=2em,itemsep=0.25em]
\item 找出这一段承接的上一个概念。
\item 找出它准备引出的下一个概念。
\item 把中间的关键等式或构造单独抄出来。
\end{itemize}

\smallskip
\textbf{本章读法提醒：} 第七章指数型生成函数的标签分配是核心，别和普通生成函数混用。
\end{explainblock}


\begin{sourceblock}{例 7.1.5 求 n 元集合二叉全划分的个数 h(n).}
{\small\sloppy
\noindent 例 7.1.5 求 n 元集合二叉全划分的个数 h(n).\par
}
\end{sourceblock}


\begin{explainblock}
例题是学习这本书最重要的部分。不要只看答案，要把每个限制条件变成一个计数动作：先安排谁、再插入谁、有没有重复、需不需要除以对称性。

\smallskip
\textbf{初学者检查点：}
\begin{itemize}[leftmargin=2em,itemsep=0.25em]
\item 先写出没有限制时的总数。
\item 逐条加入限制条件，观察总数怎样变化。
\item 最后检查有没有把同一种方案重复算了多次。
\end{itemize}

\smallskip
\textbf{本章读法提醒：} 第七章指数型生成函数的标签分配是核心，别和普通生成函数混用。
\end{explainblock}


\begin{sourceblock}{注记 7.1.3 集合 [n] 的二叉划分是指集合 [n] 的二部划分, 因此集合 [n] 的二叉划}
{\small\sloppy
\noindent 注记 7.1.3 集合 [n] 的二叉划分是指集合 [n] 的二部划分, 因此集合 [n] 的二叉划\par
\noindent 136 第 7 章 指数型生成函数的乘积公式和复合公式\par
\noindent 分的个数为 S(n, 2) = 2 2−2 = 2n−1 − 1. 二叉全划分是指从集合 [n] 开始将其划分成\par
\noindent n\par
\noindent 两个部分, 然后再将每一个不是孤立元素的部分继续分成两个部分, 以此类推, 直到\par
\noindent 分到剩下的都是只含一个元素的部分.\par
\noindent 解: 不妨规定 h(0) = 0, h(1) = 1, 由题意知当 n ≥ 2 时,\par
\noindent X\par
\noindent h(n) = h(\#T1 )h(\#T2 ),\par
\noindent \{T1 ,T2 \}\par
\noindent 其中 \{T1 , T2 \} 是集合 [n] 的弱划分, 因此当 n ≥ 2 时\par
\noindent h(n) = h(\#T1 )h(\#T2 ),\par
\noindent (T1 ,T2 )\par
\noindent 因为 h(0) = 0, h(1) = 1, 所以\par
\noindent Eh − x = Eh2 .\par
\noindent 解得: √\par
\noindent 2± 4 − 8x √\par
\noindent Eh = = 1 ± 1 − 2x.\par
}
\end{sourceblock}


\begin{explainblock}
注记通常是在补充术语、记号、历史或一个容易忽略的约定。它未必给新公式，但常常决定你后面会不会把符号读错。

\smallskip
\textbf{初学者检查点：}
\begin{itemize}[leftmargin=2em,itemsep=0.25em]
\item 记下新符号的读法。
\item 确认这个约定是否只在本章使用。
\item 把注记里的限制条件补到前面的定义旁边。
\end{itemize}

\smallskip
\textbf{本章读法提醒：} 第七章指数型生成函数的标签分配是核心，别和普通生成函数混用。
\end{explainblock}


\begin{sourceblock}{第 7 章续读}
{\small\sloppy
\noindent 再由 Eh (0) = h(0) = 0 知\par
\noindent √\par
\noindent Eh = 1 − 1 − 2x\par
\noindent ∞ 1\par
\noindent X\par
\noindent =1− 2 (−2x)k\par
\noindent n\par
\noindent n=0\par
\noindent ∞\par
\noindent !\par
\noindent X (2k − 3)!!\par
\noindent =1− 1−x− k\par
\noindent x\par
\noindent k!\par
\noindent n=2\par
\noindent ∞\par
\noindent X (2k − 3)!! k\par
\noindent =x+ x ,\par
}
\end{sourceblock}


\begin{explainblock}
这一段是原书论述链的一部分。读的时候把它当作桥梁：它通常负责把上一个定义、例子或定理，引到下一个计数方法。

\smallskip
\textbf{初学者检查点：}
\begin{itemize}[leftmargin=2em,itemsep=0.25em]
\item 找出这一段承接的上一个概念。
\item 找出它准备引出的下一个概念。
\item 把中间的关键等式或构造单独抄出来。
\end{itemize}

\smallskip
\textbf{本章读法提醒：} 第七章指数型生成函数的标签分配是核心，别和普通生成函数混用。
\end{explainblock}


\begin{sourceblock}{第 7 章续读}
{\small\sloppy
\noindent k!\par
\noindent k=2\par
\noindent 所以 \par
\noindent  1, n = 1,\par
\noindent h(n) =\par
\noindent \par
\noindent (2n − 3)!!, n ≥ 2.\par
\noindent 将乘积公式迭代若干次, 我们得到了下面的结论.\par
}
\end{sourceblock}


\begin{explainblock}
这一段是原书论述链的一部分。读的时候把它当作桥梁：它通常负责把上一个定义、例子或定理，引到下一个计数方法。

\smallskip
\textbf{初学者检查点：}
\begin{itemize}[leftmargin=2em,itemsep=0.25em]
\item 找出这一段承接的上一个概念。
\item 找出它准备引出的下一个概念。
\item 把中间的关键等式或构造单独抄出来。
\end{itemize}

\smallskip
\textbf{本章读法提醒：} 第七章指数型生成函数的标签分配是核心，别和普通生成函数混用。
\end{explainblock}


\begin{sourceblock}{定理 7.1.4 对任意正整数 k, 设 f1 , f2 , . . . , fk 及 h 是定义在非负整数上的函数, 且}
{\small\sloppy
\noindent 定理 7.1.4 对任意正整数 k, 设 f1 , f2 , . . . , fk 及 h 是定义在非负整数上的函数, 且\par
\noindent 满足\par
\noindent X\par
\noindent h(\#S) = f1 (\#T1 ) · · · fk (\#Tk ),\par
\noindent (T1 ,T2 ,...,Tk )\par
\noindent 其中 S 是任意可数集, (T1 , T2 , . . . , Tk ) 为集合 X 的 k 部有序弱划分, 则\par
\noindent Y\par
\noindent k\par
\noindent Eh = Efi .\par
\noindent i=1\par
\noindent 特别地, 当 f1 (n) = f2 (n) = · · · = fk (n) 时, 我们有下面的推论.\par
}
\end{sourceblock}


\begin{explainblock}
定理段不要只看结论，要看它把哪两个计数方式连在一起。组合学证明最常见的技巧是：左边按一种标准数，右边按另一种标准数，然后说明两边数的是同一批对象。

\smallskip
\textbf{初学者检查点：}
\begin{itemize}[leftmargin=2em,itemsep=0.25em]
\item 先复述定理的假设，不要把适用范围扩大。
\item 找出证明中的基本计数原则：乘法、加法、双射、递推、容斥或生成函数。
\item 检查最后的等式化简是否和前面的对象解释一致。
\end{itemize}

\smallskip
\textbf{本章读法提醒：} 第七章指数型生成函数的标签分配是核心，别和普通生成函数混用。
\end{explainblock}


\begin{sourceblock}{推论 7.1.5 对任意正整数 k, 设 f (n) 是定义在非负整数上的函数且 f (0) = 0,}
{\small\sloppy
\noindent 推论 7.1.5 对任意正整数 k, 设 f (n) 是定义在非负整数上的函数且 f (0) = 0,\par
\noindent h(n) 由 f (n) 按下面关系定义:\par
\noindent X\par
\noindent h(\#S) = f (\#B1 ) · · · f (\#Bk ), (7.2)\par
\noindent \{B1 ,B2 ,...,Bk \}\par
\noindent 其中 S 是任意可数集, \{B1 , B2 , . . . , Bk \} 是集合 S 的 k 部划分, 则\par
\noindent Efk\par
\noindent Eh = .\par
\noindent k!\par
\noindent 证明: 设\par
\noindent X\par
\noindent H(\#S) = f (\#T1 )f (\#T2 ) · · · f (\#Tk ),\par
\noindent (T1 ,T2 ,...,Tk )\par
\noindent 其中 (T1 , T2 , . . . Tk ) 为集合 X 上的所有 k 部有序弱划分, 则由定理 7.1.4 可得\par
\noindent EH = Efk .\par
\noindent 因为 f (0) = 0, 所以\par
\noindent X\par
\noindent H(\#S) = f (\#T1 )f (\#T2 ) · · · f (\#Tk )\par
}
\end{sourceblock}


\begin{explainblock}
定理段不要只看结论，要看它把哪两个计数方式连在一起。组合学证明最常见的技巧是：左边按一种标准数，右边按另一种标准数，然后说明两边数的是同一批对象。

\smallskip
\textbf{初学者检查点：}
\begin{itemize}[leftmargin=2em,itemsep=0.25em]
\item 先复述定理的假设，不要把适用范围扩大。
\item 找出证明中的基本计数原则：乘法、加法、双射、递推、容斥或生成函数。
\item 检查最后的等式化简是否和前面的对象解释一致。
\end{itemize}

\smallskip
\textbf{本章读法提醒：} 第七章指数型生成函数的标签分配是核心，别和普通生成函数混用。
\end{explainblock}


\begin{sourceblock}{第 7 章续读}
{\small\sloppy
\noindent (T1 ,T2 ,...,Tk )\par
\noindent X\par
\noindent = k! f (\#B1 )f (\#B2 ) · · · f (\#Bk ),\par
\noindent \{B1 ,B2 ,··· ,Bk \}\par
\noindent 其中 \{B1 , B2 , . . . , Bk \} 是集合 S 的 k 部划分, 所以\par
\noindent H(\#S)\par
\noindent h(\#S) = ,\par
\noindent k!\par
\noindent 故\par
\noindent EH Efk\par
\noindent Eh = = .\par
\noindent k! k!\par
}
\end{sourceblock}


\begin{explainblock}
这一段是原书论述链的一部分。读的时候把它当作桥梁：它通常负责把上一个定义、例子或定理，引到下一个计数方法。

\smallskip
\textbf{初学者检查点：}
\begin{itemize}[leftmargin=2em,itemsep=0.25em]
\item 找出这一段承接的上一个概念。
\item 找出它准备引出的下一个概念。
\item 把中间的关键等式或构造单独抄出来。
\end{itemize}

\smallskip
\textbf{本章读法提醒：} 第七章指数型生成函数的标签分配是核心，别和普通生成函数混用。
\end{explainblock}


\begin{sourceblock}{例 7.1.6 设 S(n, k) 为第二类 Stirling 数, 求 \{S(n, k)\}n≥0 的生成函数 ES .}
{\small\sloppy
\noindent 例 7.1.6 设 S(n, k) 为第二类 Stirling 数, 求 \{S(n, k)\}n≥0 的生成函数 ES .\par
\noindent 解: 设集合 S 为 n 元集合, 由第二类 Stirling 数的组合定义知\par
\noindent X\par
\noindent S(n, k) = 1,\par
\noindent \{B1 ,B2 ,...,Bk \}\par
\noindent 其中 \{B1 , B2 , . . . , Bk \} 为集合 S 的 k 部划分. 因此, 不妨在 (7.2) 中取 f (n) = 1,\par
\noindent 138 第 7 章 指数型生成函数的乘积公式和复合公式\par
\noindent h(n) = S(n, k). 所以由推论 7.1.5 得\par
\noindent Efk\par
\noindent ES = Eh = ,\par
\noindent k!\par
\noindent 又因为\par
\noindent ∞\par
\noindent X ∞\par
\noindent xn X xn\par
\noindent Ef = f (n) = = ex − 1,\par
\noindent n! n!\par
\noindent n=1 n=1\par
}
\end{sourceblock}


\begin{explainblock}
例题是学习这本书最重要的部分。不要只看答案，要把每个限制条件变成一个计数动作：先安排谁、再插入谁、有没有重复、需不需要除以对称性。

\smallskip
\textbf{初学者检查点：}
\begin{itemize}[leftmargin=2em,itemsep=0.25em]
\item 先写出没有限制时的总数。
\item 逐条加入限制条件，观察总数怎样变化。
\item 最后检查有没有把同一种方案重复算了多次。
\end{itemize}

\smallskip
\textbf{本章读法提醒：} 第七章指数型生成函数的标签分配是核心，别和普通生成函数混用。
\end{explainblock}


\begin{sourceblock}{第 7 章续读}
{\small\sloppy
\noindent 所以\par
\noindent (ex − 1)k\par
\noindent ES = .\par
\noindent k!\par
}
\end{sourceblock}


\begin{explainblock}
这一段是原书论述链的一部分。读的时候把它当作桥梁：它通常负责把上一个定义、例子或定理，引到下一个计数方法。

\smallskip
\textbf{初学者检查点：}
\begin{itemize}[leftmargin=2em,itemsep=0.25em]
\item 找出这一段承接的上一个概念。
\item 找出它准备引出的下一个概念。
\item 把中间的关键等式或构造单独抄出来。
\end{itemize}

\smallskip
\textbf{本章读法提醒：} 第七章指数型生成函数的标签分配是核心，别和普通生成函数混用。
\end{explainblock}


\begin{sourceblock}{§ 7.2 复合公式}
{\small\sloppy
\noindent § 7.2. 复合公式\par
\noindent 有了前面的结论我们就可以给出组合计数上的一个十分有用的公式—指数型形\par
\noindent 式幂级数的复合公式. 首先给出两个形式幂级数的复合的定义如下.\par
}
\end{sourceblock}


\begin{explainblock}
这一节先不要急着背公式。先问三个问题：对象是什么，哪些限制条件不能丢，最后到底是在数所有对象、还是在数某个统计量的分布。把这三件事说清楚，后面的公式才会有位置。

\smallskip
\textbf{初学者检查点：}
\begin{itemize}[leftmargin=2em,itemsep=0.25em]
\item 把本节出现的新对象写成一句话定义。
\item 把每个公式左边和右边分别翻译成“在数什么”。
\item 找一个最小非平凡例子，手工列出所有对象。
\end{itemize}

\smallskip
\textbf{本章读法提醒：} 第七章指数型生成函数的标签分配是核心，别和普通生成函数混用。
\end{explainblock}


\begin{sourceblock}{定义 7.2.1 设}
{\small\sloppy
\noindent 定义 7.2.1 设\par
\noindent ∞\par
\noindent X ∞\par
\noindent X\par
\noindent f (x) = an xn , g(x) = bn xn ,\par
\noindent n=0 n=0\par
\noindent 若 b0 = 0, 则称\par
\noindent ∞\par
\noindent X\par
\noindent f (g(x)) = an g n (x)\par
\noindent n=0\par
\noindent 为 f 与 g 的复合, 简记为 f ◦ g.\par
\noindent 有了复合的概念, 我们现在可以给出下面的复合公式.\par
}
\end{sourceblock}


\begin{explainblock}
定义段的任务是定边界。这里要特别注意参数范围、是否允许重复、是否考虑顺序、对象有没有标签。后面的每一道例题和证明，本质上都在反复使用这些边界。

\smallskip
\textbf{初学者检查点：}
\begin{itemize}[leftmargin=2em,itemsep=0.25em]
\item 圈出被定义的对象名称。
\item 圈出参数范围，例如 n、k 是否要求正整数。
\item 判断它是有序对象、无序对象、带标签对象还是结构对象。
\end{itemize}

\smallskip
\textbf{本章读法提醒：} 第七章指数型生成函数的标签分配是核心，别和普通生成函数混用。
\end{explainblock}


\begin{sourceblock}{定理 7.2.2 (复合公式) 设 f, g 与 h 是定 义在 非负 整数 上的 函数, 若 f (0) =}
{\small\sloppy
\noindent 定理 7.2.2 (复合公式) 设 f, g 与 h 是定 义在 非负 整数 上的 函数, 若 f (0) =\par
\noindent 0, g(0) = h(0) = 1, 且当 n ≥ 1 时, 满足下面的关系:\par
\noindent X\par
\noindent n X\par
\noindent h(\#S) = f (\#B1 ) · · · f (\#Bk )g(k), (7.3)\par
\noindent k=1 \{B1 ,B2 ,...,Bk \}\par
\noindent 其中 S 为任意可数集, \{B1 , B2 , . . . , Bk \} 为 S 的 k 部划分, 则\par
\noindent Eh = Eg (Ef ).\par
\noindent 证明: 对任意正整数 k, 设\par
\noindent X\par
\noindent tk (\#S) = g(k) f (\#B1 ) · · · f (\#Bk ),\par
\noindent \{B1 ,B2 ,...,Bk \}\par
\noindent 其中 \{B1 , B2 , . . . , Bk \} 是集合 S 的 k 部划分, 并规定 tk (0) = 0, 显然当 n < k 时,\par
\noindent tk (n) = 0. 由推论 7.1.5 可得\par
\noindent ∞\par
\noindent X X ∞\par
\noindent xn xn g(k) k\par
\noindent Etk = tk (n) = tk (n) = E (7.4)\par
}
\end{sourceblock}


\begin{explainblock}
定理段不要只看结论，要看它把哪两个计数方式连在一起。组合学证明最常见的技巧是：左边按一种标准数，右边按另一种标准数，然后说明两边数的是同一批对象。

\smallskip
\textbf{初学者检查点：}
\begin{itemize}[leftmargin=2em,itemsep=0.25em]
\item 先复述定理的假设，不要把适用范围扩大。
\item 找出证明中的基本计数原则：乘法、加法、双射、递推、容斥或生成函数。
\item 检查最后的等式化简是否和前面的对象解释一致。
\end{itemize}

\smallskip
\textbf{本章读法提醒：} 第七章指数型生成函数的标签分配是核心，别和普通生成函数混用。
\end{explainblock}


\begin{sourceblock}{第 7 章续读}
{\small\sloppy
\noindent n! n! k! f\par
\noindent n=1 n=k\par
\noindent 现不妨设 \#S = n(n ≥ 1), 则由 h(n) 定义可得:\par
\noindent X\par
\noindent n X\par
\noindent n\par
\noindent h(\#S) = tk (\#S) = tk (n),\par
\noindent k=1 k=1\par
\noindent 所以\par
\noindent ∞\par
\noindent X X ∞\par
\noindent xn xn\par
\noindent Eh (x) = h(n) =1+ h(n)\par
\noindent n! n!\par
\noindent n=0 n=1\par
\noindent ∞ X\par
\noindent X n\par
\noindent xn\par
}
\end{sourceblock}


\begin{explainblock}
这一段是原书论述链的一部分。读的时候把它当作桥梁：它通常负责把上一个定义、例子或定理，引到下一个计数方法。

\smallskip
\textbf{初学者检查点：}
\begin{itemize}[leftmargin=2em,itemsep=0.25em]
\item 找出这一段承接的上一个概念。
\item 找出它准备引出的下一个概念。
\item 把中间的关键等式或构造单独抄出来。
\end{itemize}

\smallskip
\textbf{本章读法提醒：} 第七章指数型生成函数的标签分配是核心，别和普通生成函数混用。
\end{explainblock}


\begin{sourceblock}{第 7 章续读}
{\small\sloppy
\noindent =1+ tk (n)\par
\noindent n!\par
\noindent n=1 k=1\par
\noindent ∞\par
\noindent ∞ X\par
\noindent X xn\par
\noindent =1+ tk (n) . (7.5)\par
\noindent n!\par
\noindent k=1 n=k\par
\noindent 将 (7.4) 代入 (7.5) 得\par
\noindent ∞ X\par
\noindent X ∞\par
\noindent xn\par
\noindent Eh = 1 + tk (n)\par
\noindent n!\par
\noindent k=1 n=k\par
\noindent ∞\par
\noindent X Efk\par
}
\end{sourceblock}


\begin{explainblock}
这一段是原书论述链的一部分。读的时候把它当作桥梁：它通常负责把上一个定义、例子或定理，引到下一个计数方法。

\smallskip
\textbf{初学者检查点：}
\begin{itemize}[leftmargin=2em,itemsep=0.25em]
\item 找出这一段承接的上一个概念。
\item 找出它准备引出的下一个概念。
\item 把中间的关键等式或构造单独抄出来。
\end{itemize}

\smallskip
\textbf{本章读法提醒：} 第七章指数型生成函数的标签分配是核心，别和普通生成函数混用。
\end{explainblock}


\begin{sourceblock}{第 7 章续读}
{\small\sloppy
\noindent =1+ g(k)\par
\noindent k!\par
\noindent k=1\par
\noindent ∞\par
\noindent X Efk\par
\noindent = g(k) = Eg (Ef ).\par
\noindent k!\par
\noindent k=0\par
}
\end{sourceblock}


\begin{explainblock}
这一段是原书论述链的一部分。读的时候把它当作桥梁：它通常负责把上一个定义、例子或定理，引到下一个计数方法。

\smallskip
\textbf{初学者检查点：}
\begin{itemize}[leftmargin=2em,itemsep=0.25em]
\item 找出这一段承接的上一个概念。
\item 找出它准备引出的下一个概念。
\item 把中间的关键等式或构造单独抄出来。
\end{itemize}

\smallskip
\textbf{本章读法提醒：} 第七章指数型生成函数的标签分配是核心，别和普通生成函数混用。
\end{explainblock}


\begin{sourceblock}{例 7.2.1 求将由 n 个人组成的集合 S 分成若干小组, 每组都排成一个小队, 小}
{\small\sloppy
\noindent 例 7.2.1 求将由 n 个人组成的集合 S 分成若干小组, 每组都排成一个小队, 小\par
\noindent 队的第一个人有序的站在同一圆上的方法数 h(n).\par
\noindent 解: 当 n 为正整数时, 设 f (n) = n!, g(n) = (n − 1)!, 规定 f (0) = 0. 由题意可知\par
\noindent X\par
\noindent n X\par
\noindent h(\#S) = f (\#B1 )f (\#B2 ) · · · f (\#Bk )g(k),\par
\noindent k=1 \{B1 ,B2 ,...,Bk \}\par
\noindent 其中 \{B1 , B2 , . . . , Bk \} 是集合 S 的 k 部划分, 若设 h(0) = g(0) = 1, 则由定理 7.2.2\par
\noindent 得\par
\noindent Eh = Eg (Ef ).\par
\noindent 140 第 7 章 指数型生成函数的乘积公式和复合公式\par
\noindent 因为\par
\noindent ∞\par
\noindent X xn\par
\noindent Eg = 1 + (n − 1)! = 1 − log(1 − x),\par
\noindent n!\par
\noindent n=1\par
\noindent ∞\par
}
\end{sourceblock}


\begin{explainblock}
例题是学习这本书最重要的部分。不要只看答案，要把每个限制条件变成一个计数动作：先安排谁、再插入谁、有没有重复、需不需要除以对称性。

\smallskip
\textbf{初学者检查点：}
\begin{itemize}[leftmargin=2em,itemsep=0.25em]
\item 先写出没有限制时的总数。
\item 逐条加入限制条件，观察总数怎样变化。
\item 最后检查有没有把同一种方案重复算了多次。
\end{itemize}

\smallskip
\textbf{本章读法提醒：} 第七章指数型生成函数的标签分配是核心，别和普通生成函数混用。
\end{explainblock}


\begin{sourceblock}{第 7 章续读}
{\small\sloppy
\noindent X x\par
\noindent Ef = xn = ,\par
\noindent 1−x\par
\noindent n=1\par
\noindent 所以\par
\noindent x\par
\noindent Eh = 1 − log(1 − ) = 1 − log(1 − 2x) + log(1 − x)\par
\noindent 1−x\par
\noindent ∞\par
\noindent X ∞\par
\noindent xn X xn\par
\noindent =1+ 2n (n − 1)! − (n − 1)!\par
\noindent n! n!\par
\noindent n=1 n=1\par
\noindent ∞\par
\noindent X xn\par
\noindent =1+ (2n − 1)(n − 1)! ,\par
\noindent n!\par
}
\end{sourceblock}


\begin{explainblock}
这一段是原书论述链的一部分。读的时候把它当作桥梁：它通常负责把上一个定义、例子或定理，引到下一个计数方法。

\smallskip
\textbf{初学者检查点：}
\begin{itemize}[leftmargin=2em,itemsep=0.25em]
\item 找出这一段承接的上一个概念。
\item 找出它准备引出的下一个概念。
\item 把中间的关键等式或构造单独抄出来。
\end{itemize}

\smallskip
\textbf{本章读法提醒：} 第七章指数型生成函数的标签分配是核心，别和普通生成函数混用。
\end{explainblock}


\begin{sourceblock}{第 7 章续读}
{\small\sloppy
\noindent n=1\par
\noindent 故\par
\noindent h(n) = (2n − 1)(n − 1)!, n ≥ 1.\par
\noindent 复合公式与乘积公式的区别在于, 对乘积公式来说, 是对集合 \#S 的所有的 k 部\par
\noindent 有序弱划分, k 为固定的正整数, 而对复合公式而言, 是对集合 S 的全部划分. 而且,\par
\noindent 对乘积公式而言, h(\#S) 计数的是将 k 个结构分别附加到集合 S 的 k 个子集上得到\par
\noindent 的新结构的个数, 但是对复合公式而言, 集合 S 上的某结构可以看作是连通分支的\par
\noindent 并, 而且在每一个连通分支上都存在一个结构, 同时也有一个额外结构附加在这个连\par
\noindent 通集上, 并用 g(k) 表示具有 k 个连通分支的连通集附加一个额外结构的方法数, 此\par
\noindent 时 h(n) 计数的则是集合 S 附加了这些结构的所有方法数.\par
\noindent 特别地, 对任意非负整数 n, 取 g(n) = 1, 则得到了下面推论.\par
}
\end{sourceblock}


\begin{explainblock}
这一段是原书论述链的一部分。读的时候把它当作桥梁：它通常负责把上一个定义、例子或定理，引到下一个计数方法。

\smallskip
\textbf{初学者检查点：}
\begin{itemize}[leftmargin=2em,itemsep=0.25em]
\item 找出这一段承接的上一个概念。
\item 找出它准备引出的下一个概念。
\item 把中间的关键等式或构造单独抄出来。
\end{itemize}

\smallskip
\textbf{本章读法提醒：} 第七章指数型生成函数的标签分配是核心，别和普通生成函数混用。
\end{explainblock}


\begin{sourceblock}{推论 7.2.3 设 f, h 是定义在非负整数上的函数, 若 f (0) = 0, h(0) = 1, 且当 n ≥ 1}
{\small\sloppy
\noindent 推论 7.2.3 设 f, h 是定义在非负整数上的函数, 若 f (0) = 0, h(0) = 1, 且当 n ≥ 1\par
\noindent 时, 满足下面的关系:\par
\noindent X\par
\noindent n X\par
\noindent h(\#S) = f (\#B1 ) · · · f (\#Bk ), (7.6)\par
\noindent k=1 \{B1 ,B2 ,...,Bk \}\par
\noindent 其中 S 为任意可数集, \{B1 , B2 , . . . , Bk \} 为 S 的 k 部划分, 则\par
\noindent Eh = exp(Ef ).\par
\noindent 从组合角度来看这个推论, 可以理解为一个附加在集合上的结构是由若干连通分\par
\noindent 支结合起来的, 对于整个分支不再赋予额外的结构.\par
}
\end{sourceblock}


\begin{explainblock}
定理段不要只看结论，要看它把哪两个计数方式连在一起。组合学证明最常见的技巧是：左边按一种标准数，右边按另一种标准数，然后说明两边数的是同一批对象。

\smallskip
\textbf{初学者检查点：}
\begin{itemize}[leftmargin=2em,itemsep=0.25em]
\item 先复述定理的假设，不要把适用范围扩大。
\item 找出证明中的基本计数原则：乘法、加法、双射、递推、容斥或生成函数。
\item 检查最后的等式化简是否和前面的对象解释一致。
\end{itemize}

\smallskip
\textbf{本章读法提醒：} 第七章指数型生成函数的标签分配是核心，别和普通生成函数混用。
\end{explainblock}


\begin{sourceblock}{例 7.2.2 设 n ≥ 1, 求将集合 [n] 分成若干块后, 再对每块中的元素线性排序的}
{\small\sloppy
\noindent 例 7.2.2 设 n ≥ 1, 求将集合 [n] 分成若干块后, 再对每块中的元素线性排序的\par
\noindent 方法数 h(n).\par
\noindent 解: 对任意正整数 n, 设 f (n) = n!, 并规定 f (0) = 0. 由题意知, 当 n ≥ 1 时,\par
\noindent X\par
\noindent n X\par
\noindent h(n) = f (\#B1 )f (\#B2 ) · · · f (\#Bk ),\par
\noindent k=1 \{B1 ,B2 ,...,Bk \}\par
\noindent 所以由推论 7.2.3 知\par
\noindent ∞\par
\noindent !\par
\noindent X x\par
\noindent Eh = exp (Ef ) = exp xn = e 1−x .\par
\noindent n=1\par
}
\end{sourceblock}


\begin{explainblock}
例题是学习这本书最重要的部分。不要只看答案，要把每个限制条件变成一个计数动作：先安排谁、再插入谁、有没有重复、需不需要除以对称性。

\smallskip
\textbf{初学者检查点：}
\begin{itemize}[leftmargin=2em,itemsep=0.25em]
\item 先写出没有限制时的总数。
\item 逐条加入限制条件，观察总数怎样变化。
\item 最后检查有没有把同一种方案重复算了多次。
\end{itemize}

\smallskip
\textbf{本章读法提醒：} 第七章指数型生成函数的标签分配是核心，别和普通生成函数混用。
\end{explainblock}


\begin{sourceblock}{例 7.2.3 求 Bell 数序列 \{B(n)\}n≥0 的生成函数 EB .}
{\small\sloppy
\noindent 例 7.2.3 求 Bell 数序列 \{B(n)\}n≥0 的生成函数 EB .\par
\noindent 解: 设集合 X 为非空可数集合, 由 Bell 数 B(\#X) 表示集合 X 的划分个数, 即\par
\noindent X\par
\noindent \#X\par
\noindent B(\#X) = S(\#X, k),\par
\noindent k=1\par
\noindent 其中 S(\#X, k) 是第二类 Stirling 数, 并规定 B(0) = 1. 不妨设 \#X = n(n ≥ 1), 由\par
}
\end{sourceblock}


\begin{explainblock}
例题是学习这本书最重要的部分。不要只看答案，要把每个限制条件变成一个计数动作：先安排谁、再插入谁、有没有重复、需不需要除以对称性。

\smallskip
\textbf{初学者检查点：}
\begin{itemize}[leftmargin=2em,itemsep=0.25em]
\item 先写出没有限制时的总数。
\item 逐条加入限制条件，观察总数怎样变化。
\item 最后检查有没有把同一种方案重复算了多次。
\end{itemize}

\smallskip
\textbf{本章读法提醒：} 第七章指数型生成函数的标签分配是核心，别和普通生成函数混用。
\end{explainblock}


\begin{sourceblock}{例 7.1.6 知}
{\small\sloppy
\noindent 例 7.1.6 知\par
\noindent X\par
\noindent S(n, k) = 1,\par
\noindent \{B1 ,B2 ,...,Bk \}\par
\noindent 其中 \{B1 , B2 , . . . , Bk \} 为集合 X 的 k 部划分, 所以\par
\noindent X\par
\noindent n X\par
\noindent B(n) = B(\#X) = 1, n ≥ 1. (7.7)\par
\noindent k=1 \{B1 ,B2 ,...,Bk \}\par
\noindent 对任意正整数 n, 设 h(n) = B(n), f (n) = 1, 并规定 h(0) = B(0) = 1, f (0) = 0, 则\par
\noindent 由推论 7.2.3 得\par
\noindent ∞\par
\noindent !\par
\noindent X xn\par
\noindent = ee −1 .\par
\noindent x\par
\noindent EB = Eh = exp (Ef ) = exp\par
\noindent n!\par
}
\end{sourceblock}


\begin{explainblock}
例题是学习这本书最重要的部分。不要只看答案，要把每个限制条件变成一个计数动作：先安排谁、再插入谁、有没有重复、需不需要除以对称性。

\smallskip
\textbf{初学者检查点：}
\begin{itemize}[leftmargin=2em,itemsep=0.25em]
\item 先写出没有限制时的总数。
\item 逐条加入限制条件，观察总数怎样变化。
\item 最后检查有没有把同一种方案重复算了多次。
\end{itemize}

\smallskip
\textbf{本章读法提醒：} 第七章指数型生成函数的标签分配是核心，别和普通生成函数混用。
\end{explainblock}


\begin{sourceblock}{第 7 章续读}
{\small\sloppy
\noindent n=1\par
}
\end{sourceblock}


\begin{explainblock}
这一段是原书论述链的一部分。读的时候把它当作桥梁：它通常负责把上一个定义、例子或定理，引到下一个计数方法。

\smallskip
\textbf{初学者检查点：}
\begin{itemize}[leftmargin=2em,itemsep=0.25em]
\item 找出这一段承接的上一个概念。
\item 找出它准备引出的下一个概念。
\item 把中间的关键等式或构造单独抄出来。
\end{itemize}

\smallskip
\textbf{本章读法提醒：} 第七章指数型生成函数的标签分配是核心，别和普通生成函数混用。
\end{explainblock}


\begin{sourceblock}{例 7.2.4 设 S 为非空 n(n ≥ 1) 元集合, 求对集合 S 划分后, 再对每一块进行划}
{\small\sloppy
\noindent 例 7.2.4 设 S 为非空 n(n ≥ 1) 元集合, 求对集合 S 划分后, 再对每一块进行划\par
\noindent 分的方法数 B2 (\#S) 的指数型生成函数.\par
\noindent 解: 由题意可知\par
\noindent X\par
\noindent B2 (\#S) = B(\#T1 )B(\#T2 ) · · · B(\#Tk ),\par
\noindent \{T1 ,T2 ,...,Tk \}\par
\noindent 其中 \{T1 , T2 , . . . , Tk \} 是集合 S 的 k 部划分. 若设 B2 (0) = 1, 则由推论 7.2.3 知\par
\noindent ∞\par
\noindent X xn\par
\noindent B2 (n) = exp(EB − 1) = exp(exp(ex − 1) − 1).\par
\noindent n!\par
\noindent n=0\par
}
\end{sourceblock}


\begin{explainblock}
例题是学习这本书最重要的部分。不要只看答案，要把每个限制条件变成一个计数动作：先安排谁、再插入谁、有没有重复、需不需要除以对称性。

\smallskip
\textbf{初学者检查点：}
\begin{itemize}[leftmargin=2em,itemsep=0.25em]
\item 先写出没有限制时的总数。
\item 逐条加入限制条件，观察总数怎样变化。
\item 最后检查有没有把同一种方案重复算了多次。
\end{itemize}

\smallskip
\textbf{本章读法提醒：} 第七章指数型生成函数的标签分配是核心，别和普通生成函数混用。
\end{explainblock}


\begin{sourceblock}{例 7.2.5 求顶点集为 n(n ≥ 1) 元集合 S 的连通图的个数 c(n) 的指数型生成函}
{\small\sloppy
\noindent 例 7.2.5 求顶点集为 n(n ≥ 1) 元集合 S 的连通图的个数 c(n) 的指数型生成函\par
\noindent 142 第 7 章 指数型生成函数的乘积公式和复合公式\par
\noindent 数 Ec . (注: 连通图是指图上任意两顶点均有一条路.)\par
\noindent 解: 对任意正整数 n, 设 h(\#S) 表示顶点集为 n 元集合 S 的图的个数, 显然\par
\noindent n\par
\noindent h(n) = 2( 2 ) , 并规定 h(0) = 1. 因为任意一个 n 个顶点的图都可以看作由任意\par
\noindent k(1 ≤ k ≤ n) 个连通图构成, 因此对任意 n ≥ 1,\par
\noindent X\par
\noindent n X\par
\noindent h(\#S) = c(\#B1 )c(\#B2 ) · · · c(\#Bk ),\par
\noindent k=1 \{B1 ,B2 ,...,Bk \}\par
\noindent 若规定 c(0) = 0, 则由推论 7.2.3 知:\par
\noindent Eh = exp (Ec ),\par
\noindent 所以  \par
\noindent ∞\par
\noindent X n\par
\noindent x\par
\noindent Ec = log(Eh ) = log  2( ) .\par
}
\end{sourceblock}


\begin{explainblock}
例题是学习这本书最重要的部分。不要只看答案，要把每个限制条件变成一个计数动作：先安排谁、再插入谁、有没有重复、需不需要除以对称性。

\smallskip
\textbf{初学者检查点：}
\begin{itemize}[leftmargin=2em,itemsep=0.25em]
\item 先写出没有限制时的总数。
\item 逐条加入限制条件，观察总数怎样变化。
\item 最后检查有没有把同一种方案重复算了多次。
\end{itemize}

\smallskip
\textbf{本章读法提醒：} 第七章指数型生成函数的标签分配是核心，别和普通生成函数混用。
\end{explainblock}


\begin{sourceblock}{第 7 章续读}
{\small\sloppy
\noindent n\par
\noindent n!\par
\noindent n≥0\par
}
\end{sourceblock}


\begin{explainblock}
这一段是原书论述链的一部分。读的时候把它当作桥梁：它通常负责把上一个定义、例子或定理，引到下一个计数方法。

\smallskip
\textbf{初学者检查点：}
\begin{itemize}[leftmargin=2em,itemsep=0.25em]
\item 找出这一段承接的上一个概念。
\item 找出它准备引出的下一个概念。
\item 把中间的关键等式或构造单独抄出来。
\end{itemize}

\smallskip
\textbf{本章读法提醒：} 第七章指数型生成函数的标签分配是核心，别和普通生成函数混用。
\end{explainblock}


\begin{sourceblock}{例 7.2.6 求顶点集为 n(n ≥ 1) 元集 S 上具有 k 个连通分支的图的个数 ck (n)}
{\small\sloppy
\noindent 例 7.2.6 求顶点集为 n(n ≥ 1) 元集 S 上具有 k 个连通分支的图的个数 ck (n)\par
\noindent 的指数型生成函数 Eck .\par
\noindent 解: 由题意可知\par
\noindent X\par
\noindent ck (\#S) = c(\#B1 )c(\#B2 ) · · · c(\#Bk ),\par
\noindent \{B1 ,B2 ,...,Bk \}\par
\noindent 其中 \{B1 , B2 , . . . , Bk \} 是集合 S 的 k 部划分, 则由推论 7.1.5 知\par
\noindent Eck\par
\noindent Eck = ,\par
\noindent k!\par
\noindent 再由例 7.2.5 得 !\par
\noindent ∞\par
\noindent X\par
\noindent Eck = logk 2( )\par
\noindent k! n!\par
\noindent n=0\par
}
\end{sourceblock}


\begin{explainblock}
例题是学习这本书最重要的部分。不要只看答案，要把每个限制条件变成一个计数动作：先安排谁、再插入谁、有没有重复、需不需要除以对称性。

\smallskip
\textbf{初学者检查点：}
\begin{itemize}[leftmargin=2em,itemsep=0.25em]
\item 先写出没有限制时的总数。
\item 逐条加入限制条件，观察总数怎样变化。
\item 最后检查有没有把同一种方案重复算了多次。
\end{itemize}

\smallskip
\textbf{本章读法提醒：} 第七章指数型生成函数的标签分配是核心，别和普通生成函数混用。
\end{explainblock}


\begin{sourceblock}{注记 7.2.4 一般地, 对任意正整数 n, 若 h(n) 计数的 n 元集合具某种结构 α 的个}
{\small\sloppy
\noindent 注记 7.2.4 一般地, 对任意正整数 n, 若 h(n) 计数的 n 元集合具某种结构 α 的个\par
\noindent 数, 且规定 h(0) = 1. 若这种结构可以通过将 n 分成 k(1 ≤ k ≤ n) 个小块, 而每一个\par
\noindent 小块都具有同一种其它的结构 β, 若对任意 n 元集合, 它上面具有结构 β 的个数只与\par
\noindent n 有关, 不妨记为 f (n), 且规定 f (0) = 0. 若 fk (n) 表示集合 n 上恰有 k(1 ≤ k ≤ n)\par
\noindent 个分支具有结构 β 的个数 (显然, f1 (n) = f (n)), 则由推论 7.2.3 及推论 7.1.5 知:\par
\noindent Efk\par
\noindent Eh = exp (Ef ), E fk = ,\par
\noindent k!\par
\noindent 故而\par
\noindent logk (Eh )\par
\noindent Efk = .\par
\noindent k!\par
\noindent 下面我们再举个例子说明.\par
}
\end{sourceblock}


\begin{explainblock}
注记通常是在补充术语、记号、历史或一个容易忽略的约定。它未必给新公式，但常常决定你后面会不会把符号读错。

\smallskip
\textbf{初学者检查点：}
\begin{itemize}[leftmargin=2em,itemsep=0.25em]
\item 记下新符号的读法。
\item 确认这个约定是否只在本章使用。
\item 把注记里的限制条件补到前面的定义旁边。
\end{itemize}

\smallskip
\textbf{本章读法提醒：} 第七章指数型生成函数的标签分配是核心，别和普通生成函数混用。
\end{explainblock}


\begin{sourceblock}{例 7.2.7 求无符号第一类 Stirling 数 c(n, k) 的关于变量 n 的生成函数.}
{\small\sloppy
\noindent 例 7.2.7 求无符号第一类 Stirling 数 c(n, k) 的关于变量 n 的生成函数.\par
\noindent 解: 设 S 为 n 元非空集合, 记 h(\#S) 为集合 S 的排列个数, 因此 h(\#S) = n!,\par
\noindent 并规定 h(0) = 1, 所以\par
\noindent ∞\par
\noindent X xn 1\par
\noindent Eh = n! = .\par
\noindent n! 1−x\par
\noindent n=0\par
\noindent 由于集合 S 上任意一个排列都有唯一的圈结构与其对应, 换言之, 集合 S 上的排列\par
\noindent 可以看作将 S 分成互不相交的 k 部分, 然后每一部分形成一个圈, 因此由注记 7.2.4\par
\noindent 知\par
\noindent ∞\par
\noindent X\par
\noindent xn logk (Eh ) 1 1\par
\noindent c(n, k) = = logk .\par
\noindent n! k! k! 1−x\par
\noindent n=0\par
\noindent 下面我们给出复合公式在排列的圈表示上的形式.\par
}
\end{sourceblock}


\begin{explainblock}
例题是学习这本书最重要的部分。不要只看答案，要把每个限制条件变成一个计数动作：先安排谁、再插入谁、有没有重复、需不需要除以对称性。

\smallskip
\textbf{初学者检查点：}
\begin{itemize}[leftmargin=2em,itemsep=0.25em]
\item 先写出没有限制时的总数。
\item 逐条加入限制条件，观察总数怎样变化。
\item 最后检查有没有把同一种方案重复算了多次。
\end{itemize}

\smallskip
\textbf{本章读法提醒：} 第七章指数型生成函数的标签分配是核心，别和普通生成函数混用。
\end{explainblock}


\begin{sourceblock}{定理 7.2.5 设 f, g 与 h 是定义在非负整数上的函数, 且 f (0) = 0, h(0) = g(0). 若}
{\small\sloppy
\noindent 定理 7.2.5 设 f, g 与 h 是定义在非负整数上的函数, 且 f (0) = 0, h(0) = g(0). 若\par
\noindent f, g, h 满足下面的关系:\par
\noindent X\par
\noindent h(\#X) = f (\#C1 ) · · · f (\#Ck )g(k), (7.8)\par
\noindent π∈S(X)\par
\noindent 其中集合 X 为任意非空可数集, S(X) 表示集合 X 上的所有排列构成的集合,\par
\noindent C1 , C2 , . . . , Ck 是排列 π 表示成圈结构时对应的圈, 则\par
\noindent ∞\par
\noindent !\par
\noindent X xn\par
\noindent Eh (x) = Eg f (n) .\par
\noindent n\par
\noindent n=1\par
\noindent 证明: 因为对于任意 j 元集合的圈排列的个数为 (j − 1)!, 所以当 n ≥ 1 时,\par
\noindent (7.8) 也等价于\par
\noindent X\par
\noindent n X\par
\noindent h(\#X) = h(n) = (\#B1 − 1)!f (\#B1 ) · · · (\#Bk − 1)!f (\#Bk )g(k),\par
}
\end{sourceblock}


\begin{explainblock}
定理段不要只看结论，要看它把哪两个计数方式连在一起。组合学证明最常见的技巧是：左边按一种标准数，右边按另一种标准数，然后说明两边数的是同一批对象。

\smallskip
\textbf{初学者检查点：}
\begin{itemize}[leftmargin=2em,itemsep=0.25em]
\item 先复述定理的假设，不要把适用范围扩大。
\item 找出证明中的基本计数原则：乘法、加法、双射、递推、容斥或生成函数。
\item 检查最后的等式化简是否和前面的对象解释一致。
\end{itemize}

\smallskip
\textbf{本章读法提醒：} 第七章指数型生成函数的标签分配是核心，别和普通生成函数混用。
\end{explainblock}


\begin{sourceblock}{第 7 章续读}
{\small\sloppy
\noindent k=1 \{B1 ,B2 ,...,Bk \}\par
\noindent (7.9)\par
\noindent 其中 \{B1 , B2 , . . . , Bk \} 为非空集合 X 的 k 部划分, 所以由复合公式 7.2.2 可得\par
\noindent ∞\par
\noindent !\par
\noindent X zn\par
\noindent Eh (x) = Eg (n − 1)!f (n) ,\par
\noindent n!\par
\noindent n=1\par
\noindent ∞\par
\noindent !\par
\noindent X xn\par
\noindent = Eg f (n) .\par
\noindent n\par
\noindent n=1\par
\noindent 特别地, 取 g(n) = 1, 我们得到了下面的结论.\par
}
\end{sourceblock}


\begin{explainblock}
这一段是原书论述链的一部分。读的时候把它当作桥梁：它通常负责把上一个定义、例子或定理，引到下一个计数方法。

\smallskip
\textbf{初学者检查点：}
\begin{itemize}[leftmargin=2em,itemsep=0.25em]
\item 找出这一段承接的上一个概念。
\item 找出它准备引出的下一个概念。
\item 把中间的关键等式或构造单独抄出来。
\end{itemize}

\smallskip
\textbf{本章读法提醒：} 第七章指数型生成函数的标签分配是核心，别和普通生成函数混用。
\end{explainblock}


\begin{sourceblock}{推论 7.2.6 设 f, h 是定义在非负整数上的函数, 且 f (0) = 0, h(0) = 1. 若 f, h 满}
{\small\sloppy
\noindent 推论 7.2.6 设 f, h 是定义在非负整数上的函数, 且 f (0) = 0, h(0) = 1. 若 f, h 满\par
\noindent 足下面的关系:\par
\noindent X\par
\noindent h(\#X) = f (\#C1 ) · · · f (\#Ck ), (7.10)\par
\noindent π∈S(X)\par
\noindent 其中集合 X 为任意非空可数集, S(X) 表示集合 X 上的所有排列构成的集合,\par
\noindent 144 第 7 章 指数型生成函数的乘积公式和复合公式\par
\noindent C1 , C2 , . . . , Ck 是排列 π 表示成圈结构时对应的圈, 则\par
\noindent ∞\par
\noindent !\par
\noindent X xn\par
\noindent Eh (x) = exp f (n) .\par
\noindent n\par
\noindent n=1\par
}
\end{sourceblock}


\begin{explainblock}
定理段不要只看结论，要看它把哪两个计数方式连在一起。组合学证明最常见的技巧是：左边按一种标准数，右边按另一种标准数，然后说明两边数的是同一批对象。

\smallskip
\textbf{初学者检查点：}
\begin{itemize}[leftmargin=2em,itemsep=0.25em]
\item 先复述定理的假设，不要把适用范围扩大。
\item 找出证明中的基本计数原则：乘法、加法、双射、递推、容斥或生成函数。
\item 检查最后的等式化简是否和前面的对象解释一致。
\end{itemize}

\smallskip
\textbf{本章读法提醒：} 第七章指数型生成函数的标签分配是核心，别和普通生成函数混用。
\end{explainblock}


\begin{sourceblock}{例 7.2.8 现将集合 [n] 上所有排列用圈表示, 设排列中每个圈的长度都小于等于}
{\small\sloppy
\noindent 例 7.2.8 现将集合 [n] 上所有排列用圈表示, 设排列中每个圈的长度都小于等于\par
\noindent 2 的排列的个数为 h(n), 求 Eh .\par
\noindent 解: 设 S(n) 为集合 [n] 上所有排列用圈表示的集合. 设\par
\noindent \par
\noindent \par
\noindent 0, n = 0,\par
\noindent \par
\noindent \par
\noindent \par
\noindent f (n) = 1, n = 1, 2,\par
\noindent \par
\noindent \par
\noindent \par
\noindent \par
\noindent 0, n > 2.\par
\noindent 则由题意可知:\par
\noindent X\par
\noindent h(n) = f (\#C1 )f (\#C2 ) · · · f (\#Ck ),\par
}
\end{sourceblock}


\begin{explainblock}
例题是学习这本书最重要的部分。不要只看答案，要把每个限制条件变成一个计数动作：先安排谁、再插入谁、有没有重复、需不需要除以对称性。

\smallskip
\textbf{初学者检查点：}
\begin{itemize}[leftmargin=2em,itemsep=0.25em]
\item 先写出没有限制时的总数。
\item 逐条加入限制条件，观察总数怎样变化。
\item 最后检查有没有把同一种方案重复算了多次。
\end{itemize}

\smallskip
\textbf{本章读法提醒：} 第七章指数型生成函数的标签分配是核心，别和普通生成函数混用。
\end{explainblock}


\begin{sourceblock}{第 7 章续读}
{\small\sloppy
\noindent π∈S(n)\par
\noindent 其中 C1 , C2 , . . . , Ck 是 π 对应的圈. 故由推论 7.2.6 可得\par
\noindent ∞\par
\noindent !\par
\noindent X xn x2\par
\noindent Eh (x) = exp f (n) = exp x + .\par
\noindent n 2\par
\noindent n=1\par
}
\end{sourceblock}


\begin{explainblock}
这一段是原书论述链的一部分。读的时候把它当作桥梁：它通常负责把上一个定义、例子或定理，引到下一个计数方法。

\smallskip
\textbf{初学者检查点：}
\begin{itemize}[leftmargin=2em,itemsep=0.25em]
\item 找出这一段承接的上一个概念。
\item 找出它准备引出的下一个概念。
\item 把中间的关键等式或构造单独抄出来。
\end{itemize}

\smallskip
\textbf{本章读法提醒：} 第七章指数型生成函数的标签分配是核心，别和普通生成函数混用。
\end{explainblock}


\begin{sourceblock}{习 题 七}
{\small\sloppy
\noindent 习 题 七\par
}
\end{sourceblock}


\begin{explainblock}
习题段不要先追求技巧。先给题目分类：它是在问排列、组合、递推、生成函数、容斥，还是结构计数。分类正确，工具通常就只剩一两个候选。

\smallskip
\textbf{初学者检查点：}
\begin{itemize}[leftmargin=2em,itemsep=0.25em]
\item 判断题目所属工具。
\item 列出变量和限制条件。
\item 先做小规模 n=2、3、4 的例子，确认直觉。
\end{itemize}

\smallskip
\textbf{本章读法提醒：} 第七章指数型生成函数的标签分配是核心，别和普通生成函数混用。
\end{explainblock}


\begin{sourceblock}{习题 7.1 某足球队有 n 名成员, 现教练将球员分成两队. 两队都各自排成一列,}
{\small\sloppy
\noindent 习题 7.1 某足球队有 n 名成员, 现教练将球员分成两队. 两队都各自排成一列,\par
\noindent 第一队要求这队球员或穿蓝色球衣或穿白色球衣或穿黄色球衣, 第二队只能穿红色\par
\noindent 球衣. 请问有多少种情形?\par
}
\end{sourceblock}


\begin{explainblock}
习题段不要先追求技巧。先给题目分类：它是在问排列、组合、递推、生成函数、容斥，还是结构计数。分类正确，工具通常就只剩一两个候选。

\smallskip
\textbf{初学者检查点：}
\begin{itemize}[leftmargin=2em,itemsep=0.25em]
\item 判断题目所属工具。
\item 列出变量和限制条件。
\item 先做小规模 n=2、3、4 的例子，确认直觉。
\end{itemize}

\smallskip
\textbf{本章读法提醒：} 第七章指数型生成函数的标签分配是核心，别和普通生成函数混用。
\end{explainblock}


\begin{sourceblock}{习题 7.2 将 n 个人分成若干组, 每一组绕圆桌就坐, 问有多少种情形?}
{\small\sloppy
\noindent 习题 7.2 将 n 个人分成若干组, 每一组绕圆桌就坐, 问有多少种情形?\par
}
\end{sourceblock}


\begin{explainblock}
习题段不要先追求技巧。先给题目分类：它是在问排列、组合、递推、生成函数、容斥，还是结构计数。分类正确，工具通常就只剩一两个候选。

\smallskip
\textbf{初学者检查点：}
\begin{itemize}[leftmargin=2em,itemsep=0.25em]
\item 判断题目所属工具。
\item 列出变量和限制条件。
\item 先做小规模 n=2、3、4 的例子，确认直觉。
\end{itemize}

\smallskip
\textbf{本章读法提醒：} 第七章指数型生成函数的标签分配是核心，别和普通生成函数混用。
\end{explainblock}


\begin{sourceblock}{习题 7.3 设 fn 表示将集合 [n] 划分成若干块, 每块中元素个数只能是 3 或 4}
{\small\sloppy
\noindent 习题 7.3 设 fn 表示将集合 [n] 划分成若干块, 每块中元素个数只能是 3 或 4\par
\noindent 或 9 的划分的个数, 求 \{fn \} 的指数型生成函数.\par
}
\end{sourceblock}


\begin{explainblock}
习题段不要先追求技巧。先给题目分类：它是在问排列、组合、递推、生成函数、容斥，还是结构计数。分类正确，工具通常就只剩一两个候选。

\smallskip
\textbf{初学者检查点：}
\begin{itemize}[leftmargin=2em,itemsep=0.25em]
\item 判断题目所属工具。
\item 列出变量和限制条件。
\item 先做小规模 n=2、3、4 的例子，确认直觉。
\end{itemize}

\smallskip
\textbf{本章读法提醒：} 第七章指数型生成函数的标签分配是核心，别和普通生成函数混用。
\end{explainblock}


\begin{sourceblock}{习题 7.4 设 tn 是集合 [n] 上只有长为 a1 , a1 , . . . , ak 的圈的排列的个数, 求}
{\small\sloppy
\noindent 习题 7.4 设 tn 是集合 [n] 上只有长为 a1 , a1 , . . . , ak 的圈的排列的个数, 求\par
\noindent \{tn \} 的指数型生成函数.\par
}
\end{sourceblock}


\begin{explainblock}
习题段不要先追求技巧。先给题目分类：它是在问排列、组合、递推、生成函数、容斥，还是结构计数。分类正确，工具通常就只剩一两个候选。

\smallskip
\textbf{初学者检查点：}
\begin{itemize}[leftmargin=2em,itemsep=0.25em]
\item 判断题目所属工具。
\item 列出变量和限制条件。
\item 先做小规模 n=2、3、4 的例子，确认直觉。
\end{itemize}

\smallskip
\textbf{本章读法提醒：} 第七章指数型生成函数的标签分配是核心，别和普通生成函数混用。
\end{explainblock}


\begin{sourceblock}{习题 7.5 设 H2,3 (n) 表示集合 [n] 上只有长为 2 或 3 的圈的排列, 利用习题}
{\small\sloppy
\noindent 习题 7.5 设 H2,3 (n) 表示集合 [n] 上只有长为 2 或 3 的圈的排列, 利用习题\par
\noindent 7.4 求 H2,3 (n) 递推关系.\par
}
\end{sourceblock}


\begin{explainblock}
习题段不要先追求技巧。先给题目分类：它是在问排列、组合、递推、生成函数、容斥，还是结构计数。分类正确，工具通常就只剩一两个候选。

\smallskip
\textbf{初学者检查点：}
\begin{itemize}[leftmargin=2em,itemsep=0.25em]
\item 判断题目所属工具。
\item 列出变量和限制条件。
\item 先做小规模 n=2、3、4 的例子，确认直觉。
\end{itemize}

\smallskip
\textbf{本章读法提醒：} 第七章指数型生成函数的标签分配是核心，别和普通生成函数混用。
\end{explainblock}


\begin{sourceblock}{习题 7.6 设有 n 张不同的卡片, 将这些卡片分成非空的若干组, 且要求每组都}
{\small\sloppy
\noindent 习题 7.6 设有 n 张不同的卡片, 将这些卡片分成非空的若干组, 且要求每组都\par
\noindent 有偶数张卡片. 再将每组的卡片进行排序, 之后将排列好的各组卡片进行线性排开,\par
\noindent 求有多少种情形发生?\par
}
\end{sourceblock}


\begin{explainblock}
习题段不要先追求技巧。先给题目分类：它是在问排列、组合、递推、生成函数、容斥，还是结构计数。分类正确，工具通常就只剩一两个候选。

\smallskip
\textbf{初学者检查点：}
\begin{itemize}[leftmargin=2em,itemsep=0.25em]
\item 判断题目所属工具。
\item 列出变量和限制条件。
\item 先做小规模 n=2、3、4 的例子，确认直觉。
\end{itemize}

\smallskip
\textbf{本章读法提醒：} 第七章指数型生成函数的标签分配是核心，别和普通生成函数混用。
\end{explainblock}


\begin{sourceblock}{习题 7.7 设 hn 为集合 [n] 的有序划分的个数, 计算 hn 的指数型生成函数, 并}
{\small\sloppy
\noindent 习题 7.7 设 hn 为集合 [n] 的有序划分的个数, 计算 hn 的指数型生成函数, 并\par
\noindent 给出 hn 的表达式. (注: 所谓的有序划分, 是指将集合划分后再对所有的块进行排\par
\noindent 列.)\par
}
\end{sourceblock}


\begin{explainblock}
习题段不要先追求技巧。先给题目分类：它是在问排列、组合、递推、生成函数、容斥，还是结构计数。分类正确，工具通常就只剩一两个候选。

\smallskip
\textbf{初学者检查点：}
\begin{itemize}[leftmargin=2em,itemsep=0.25em]
\item 判断题目所属工具。
\item 列出变量和限制条件。
\item 先做小规模 n=2、3、4 的例子，确认直觉。
\end{itemize}

\smallskip
\textbf{本章读法提醒：} 第七章指数型生成函数的标签分配是核心，别和普通生成函数混用。
\end{explainblock}


\begin{sourceblock}{习题 7.8 设 kn 为将集合 [n] 的有序划分中的每一块中的元素再排列的方法数,}
{\small\sloppy
\noindent 习题 7.8 设 kn 为将集合 [n] 的有序划分中的每一块中的元素再排列的方法数,\par
\noindent 求 kn 的生成函数.\par
}
\end{sourceblock}


\begin{explainblock}
习题段不要先追求技巧。先给题目分类：它是在问排列、组合、递推、生成函数、容斥，还是结构计数。分类正确，工具通常就只剩一两个候选。

\smallskip
\textbf{初学者检查点：}
\begin{itemize}[leftmargin=2em,itemsep=0.25em]
\item 判断题目所属工具。
\item 列出变量和限制条件。
\item 先做小规模 n=2、3、4 的例子，确认直觉。
\end{itemize}

\smallskip
\textbf{本章读法提醒：} 第七章指数型生成函数的标签分配是核心，别和普通生成函数混用。
\end{explainblock}


\begin{sourceblock}{习题 7.9 设 k 为任意正整数, hk (n) 为集合 [n] 上的排列表示成圈结构后, 每}
{\small\sloppy
\noindent 习题 7.9 设 k 为任意正整数, hk (n) 为集合 [n] 上的排列表示成圈结构后, 每\par
\noindent 个圈的长度是 k 的倍数的排列的个数,\par
\noindent (1) 求 \{hk (n)\}n≥0 的生成函数.\par
\noindent (2) 验证:\par
\noindent X\par
\noindent n\par
\noindent hk (n) = (−1)k S(n, k)2n−k k!,\par
\noindent k=1\par
\noindent 其中 S(n, k) 为第二类 Stirling 数.\par
}
\end{sourceblock}


\begin{explainblock}
习题段不要先追求技巧。先给题目分类：它是在问排列、组合、递推、生成函数、容斥，还是结构计数。分类正确，工具通常就只剩一两个候选。

\smallskip
\textbf{初学者检查点：}
\begin{itemize}[leftmargin=2em,itemsep=0.25em]
\item 判断题目所属工具。
\item 列出变量和限制条件。
\item 先做小规模 n=2、3、4 的例子，确认直觉。
\end{itemize}

\smallskip
\textbf{本章读法提醒：} 第七章指数型生成函数的标签分配是核心，别和普通生成函数混用。
\end{explainblock}


\begin{sourceblock}{习题 7.10 设 a(n) 为集合 [n] 上圈长为奇数排列个数, b(n) 为集合 [n] 上圈长}
{\small\sloppy
\noindent 习题 7.10 设 a(n) 为集合 [n] 上圈长为奇数排列个数, b(n) 为集合 [n] 上圈长\par
\noindent 为偶数排列的个数.\par
\noindent (1) 证明: r\par
\noindent ∞\par
\noindent X xn 1+x\par
\noindent a(n) = ,\par
\noindent n! 1−x\par
\noindent n=0\par
\noindent ∞\par
\noindent X xn 1\par
\noindent b(n) =√ .\par
\noindent n=0\par
\noindent n! 1 − x2\par
\noindent (2) 求 a(n) 与 b(n) 之间的关系式.\par
\noindent 146 第 7 章 指数型生成函数的乘积公式和复合公式\par
}
\end{sourceblock}


\begin{explainblock}
习题段不要先追求技巧。先给题目分类：它是在问排列、组合、递推、生成函数、容斥，还是结构计数。分类正确，工具通常就只剩一两个候选。

\smallskip
\textbf{初学者检查点：}
\begin{itemize}[leftmargin=2em,itemsep=0.25em]
\item 判断题目所属工具。
\item 列出变量和限制条件。
\item 先做小规模 n=2、3、4 的例子，确认直觉。
\end{itemize}

\smallskip
\textbf{本章读法提醒：} 第七章指数型生成函数的标签分配是核心，别和普通生成函数混用。
\end{explainblock}


\chapter{Catalan 数与 Narayana 数}
\begin{roadmapblock}
本章保留原书正文的主要数学骨架，并在每个定义、定理、例题和论述片段后加入解读。
阅读时不要把目标放在“背下答案”上，而是放在“看懂这一步为什么这样数”上。

\smallskip
\textbf{本章最低目标：} 能说出每个公式左边在数什么、右边在数什么，以及为什么两边相等。
\end{roadmapblock}


\begin{sourceblock}{第 8 章正文}
{\small\sloppy
\noindent Catalan 数与 Narayana 数\par
}
\end{sourceblock}


\begin{explainblock}
这一段是原书论述链的一部分。读的时候把它当作桥梁：它通常负责把上一个定义、例子或定理，引到下一个计数方法。

\smallskip
\textbf{初学者检查点：}
\begin{itemize}[leftmargin=2em,itemsep=0.25em]
\item 找出这一段承接的上一个概念。
\item 找出它准备引出的下一个概念。
\item 把中间的关键等式或构造单独抄出来。
\end{itemize}

\smallskip
\textbf{本章读法提醒：} 第八章 Catalan 问题的核心常常是前缀不越界或递归分解。
\end{explainblock}


\begin{sourceblock}{§ 8.1 Catalan 数与 Dyck 路}
{\small\sloppy
\noindent § 8.1. Catalan 数与 Dyck 路\par
\noindent Catalan 数首次出现在 Catalan 的著作中, 它是在研究凸多边形剖分问题时被发\par
\noindent 现的. 后来出现在很多组合问题的研究中, 它由法裔比利时数学家 Eugène Charles\par
\noindent Catalan (1814–1894) 所命名. 他的主要工作是关于连分式, 画法几何学, 数论以及组\par
\noindent 合学. 他在 1855 年发现了 unique surface (periodic minimal surface in the space )\par
\noindent 并给它命名. 1844 年, 他已经发表了了著名的 Catalan 猜想, 并把 Catalan 数带入了\par
\noindent 解决组合问题的领域. 这个猜想在 2002 年被罗马尼亚数学家 Preda Mihǎilescu 证\par
\noindent 明.\par
\noindent Catalan 数是一类经典的组合数, 也是当代计数组合学研究的热点问题之一. 当\par
\noindent 代组合数学大师 Richard P. Stanley 在他的著作 Catalan Numbers [47] 中, 列举了\par
\noindent 214 种可以用 Catalan 数计数的事物. 并且, Stanley 在他的个人主页上持续记录着\par
\noindent 最新被发现的可用 Catalan 数计数的事物.\par
\noindent 1751 年 9 月 4 日, 欧拉 (Leonhard Euler) 在写给哥德巴赫 (Christian Goldbach)\par
\noindent 的信中提到凸多边形剖分问题. 即通过连接凸多边形的顶点 (使得对角线内部不交),\par
\noindent 将凸多边形划分成三角形的集合. 在信中, 欧拉给出了凸五边形的所有划分, 如下图:\par
\noindent a a a a a\par
\noindent e b e b e b e b e b\par
\noindent d c d c d c d c d c\par
}
\end{sourceblock}


\begin{explainblock}
这一节先不要急着背公式。先问三个问题：对象是什么，哪些限制条件不能丢，最后到底是在数所有对象、还是在数某个统计量的分布。把这三件事说清楚，后面的公式才会有位置。

\smallskip
\textbf{初学者检查点：}
\begin{itemize}[leftmargin=2em,itemsep=0.25em]
\item 把本节出现的新对象写成一句话定义。
\item 把每个公式左边和右边分别翻译成“在数什么”。
\item 找一个最小非平凡例子，手工列出所有对象。
\end{itemize}

\smallskip
\textbf{本章读法提醒：} 第八章 Catalan 问题的核心常常是前缀不越界或递归分解。
\end{explainblock}


\begin{sourceblock}{第 8 章续读}
{\small\sloppy
\noindent 另外, 欧拉还计算出了凸六边形, 凸七边形, 等结果. 如下表 (n 表示凸多边形的\par
\noindent 边数, Tn 表示剖分的个数):\par
\noindent n 3 4 5 6 7 8 9 10\par
\noindent Tn 1 2 5 14 42 152 429 1430\par
\noindent 148 第 8 章 Catalan 数与 Narayana 数\par
\noindent 并猜测对于一般的 n(≥ 3),\par
\noindent 2 · 6 · 10 · 14 · 18 · 22 · · · (4n − 10)\par
\noindent Tn = ,\par
\noindent 2 · 3 · 4 · 5 · 6 · 7 · · · (n − 1)\par
\noindent 这里的 Tn 就是著名的 Catalan 数, 写成现在的记号:\par
\noindent Tn = Cn−2 = .\par
\noindent n−1 n−2\par
}
\end{sourceblock}


\begin{explainblock}
这一段是原书论述链的一部分。读的时候把它当作桥梁：它通常负责把上一个定义、例子或定理，引到下一个计数方法。

\smallskip
\textbf{初学者检查点：}
\begin{itemize}[leftmargin=2em,itemsep=0.25em]
\item 找出这一段承接的上一个概念。
\item 找出它准备引出的下一个概念。
\item 把中间的关键等式或构造单独抄出来。
\end{itemize}

\smallskip
\textbf{本章读法提醒：} 第八章 Catalan 问题的核心常常是前缀不越界或递归分解。
\end{explainblock}


\begin{sourceblock}{注记 8.1.1 清朝数学家明安图在其数学著作《割圜密率捷法》中将 sin(2α) 展开}
{\small\sloppy
\noindent 注记 8.1.1 清朝数学家明安图在其数学著作《割圜密率捷法》中将 sin(2α) 展开\par
\noindent 为:\par
\noindent ∞\par
\noindent X Cn−1\par
\noindent sin(2α) = 2 sin(α) − sin2n+1 (α).\par
\noindent 4n−1\par
\noindent n=1\par
\noindent 这应该是 Catalan 数最早的记载. 但明安图当时并没有意识到它的重要性.\par
\noindent 关于 Catalan 数的定义有很多种, 比如上面关于凸多边形剖分三角形的个数就可\par
\noindent 以定义为 Catalan 数, 这里选用递归关系来定义 Catalan 数.\par
}
\end{sourceblock}


\begin{explainblock}
注记通常是在补充术语、记号、历史或一个容易忽略的约定。它未必给新公式，但常常决定你后面会不会把符号读错。

\smallskip
\textbf{初学者检查点：}
\begin{itemize}[leftmargin=2em,itemsep=0.25em]
\item 记下新符号的读法。
\item 确认这个约定是否只在本章使用。
\item 把注记里的限制条件补到前面的定义旁边。
\end{itemize}

\smallskip
\textbf{本章读法提醒：} 第八章 Catalan 问题的核心常常是前缀不越界或递归分解。
\end{explainblock}


\begin{sourceblock}{定义 8.1.2 规定 C0 = 1, 当 n ≥ 1 时, 满足下面递归关系的一类数称为 Catalan}
{\small\sloppy
\noindent 定义 8.1.2 规定 C0 = 1, 当 n ≥ 1 时, 满足下面递归关系的一类数称为 Catalan\par
\noindent 数:\par
\noindent Cn = C0 Cn−1 + C1 Cn−2 + · · · + Cn−1 C0 ,\par
\noindent 也即:\par
\noindent X\par
\noindent n−1 X\par
\noindent n\par
\noindent Cn = Ck Cn−1−k = Ck−1 Cn−k .\par
\noindent k=0 k=1\par
\noindent 下面我们来严格地说明凸 n 边形三角剖分个数满足 Catalan 数的定义. 我们考\par
\noindent 虑 (n + 2) 条边的凸多边形. 顶点按顺时针从 1 到 n + 2 进行标号 (如图 8.1), 设它\par
\noindent 的三角剖分的个数为 Tn+2 .\par
\noindent n+2 1\par
\noindent D1\par
\noindent D2\par
\noindent k+2 k−1\par
\noindent k+1 k\par
\noindent 显然地, 对任意 k(2 ≤ k ≤ n − 1), 由标号分别为 1, k 和 n + 2 将 (n + 2) 条边的\par
}
\end{sourceblock}


\begin{explainblock}
定义段的任务是定边界。这里要特别注意参数范围、是否允许重复、是否考虑顺序、对象有没有标签。后面的每一道例题和证明，本质上都在反复使用这些边界。

\smallskip
\textbf{初学者检查点：}
\begin{itemize}[leftmargin=2em,itemsep=0.25em]
\item 圈出被定义的对象名称。
\item 圈出参数范围，例如 n、k 是否要求正整数。
\item 判断它是有序对象、无序对象、带标签对象还是结构对象。
\end{itemize}

\smallskip
\textbf{本章读法提醒：} 第八章 Catalan 问题的核心常常是前缀不越界或递归分解。
\end{explainblock}


\begin{sourceblock}{第 8 章续读}
{\small\sloppy
\noindent 凸多边形分成了有 k 条边凸多边形 D1 与 ((n + 2) − k + 1) 条边的凸多边形 D2 . 则\par
\noindent D1 与 D2 对应的三角剖分数分别为 Tk 与 Tn−k+3 . 说明: 因为 k 取 2 或 k 取 n − 1\par
\noindent 时, D1 或 D2 不存在, 为计数需要这里设 T2 = 1.\par
\noindent 由乘法原理知, 当确定一个顶点为 1, k, n + 2 的三角形时, (n + 2) 条边的凸多边\par
\noindent 8.1 Catalan 数与 Dyck 路 149\par
\noindent 形存在 Tk · Tn−k+3 种三角剖分. 对 k 从 2 到 n − 1 求和可得:\par
\noindent X\par
\noindent n+1\par
\noindent Tn+2 = Tk · Tn−k+3 , T2 = 1.\par
\noindent k=2\par
\noindent 所以序列 \{Tn+2 \}n≥0 与 Catalan 数的初值及递推关系均相同, 故而 Tn+2 = Cn .\par
}
\end{sourceblock}


\begin{explainblock}
这一段是原书论述链的一部分。读的时候把它当作桥梁：它通常负责把上一个定义、例子或定理，引到下一个计数方法。

\smallskip
\textbf{初学者检查点：}
\begin{itemize}[leftmargin=2em,itemsep=0.25em]
\item 找出这一段承接的上一个概念。
\item 找出它准备引出的下一个概念。
\item 把中间的关键等式或构造单独抄出来。
\end{itemize}

\smallskip
\textbf{本章读法提醒：} 第八章 Catalan 问题的核心常常是前缀不越界或递归分解。
\end{explainblock}


\begin{sourceblock}{定理 8.1.3 设 c(x) 为 Catalan 数的生成函数, 即}
{\small\sloppy
\noindent 定理 8.1.3 设 c(x) 为 Catalan 数的生成函数, 即\par
\noindent ∞\par
\noindent X\par
\noindent c(x) = Cn x n ,\par
\noindent n=0\par
\noindent 则 √\par
\noindent 1− 1 − 4x\par
\noindent c(x) = .\par
\noindent 2x\par
\noindent 证明: 由\par
\noindent ∞\par
\noindent X\par
\noindent c(x) = Cn x n\par
\noindent n=0\par
\noindent ∞ X\par
\noindent X n\par
\noindent =1+ Ck Cn−1−k xn\par
\noindent n=1 k=0\par
}
\end{sourceblock}


\begin{explainblock}
定理段不要只看结论，要看它把哪两个计数方式连在一起。组合学证明最常见的技巧是：左边按一种标准数，右边按另一种标准数，然后说明两边数的是同一批对象。

\smallskip
\textbf{初学者检查点：}
\begin{itemize}[leftmargin=2em,itemsep=0.25em]
\item 先复述定理的假设，不要把适用范围扩大。
\item 找出证明中的基本计数原则：乘法、加法、双射、递推、容斥或生成函数。
\item 检查最后的等式化简是否和前面的对象解释一致。
\end{itemize}

\smallskip
\textbf{本章读法提醒：} 第八章 Catalan 问题的核心常常是前缀不越界或递归分解。
\end{explainblock}


\begin{sourceblock}{第 8 章续读}
{\small\sloppy
\noindent ∞\par
\noindent ∞ X\par
\noindent X\par
\noindent =1+ Ck Cn−1−k xn\par
\noindent k=0 n=k+1\par
\noindent ∞ X\par
\noindent X ∞\par
\noindent =1+ Ck Cm xm+k+1\par
\noindent k=0 m=0\par
\noindent = 1 + xc(x)2\par
\noindent 即 c(x) 满足:\par
\noindent c(x) = 1 + xc2 (x).\par
\noindent 解得: √\par
\noindent 1− 1 − 4x\par
\noindent c(x) = .\par
\noindent 2x\par
\noindent 由 Catalan 数的生成函数, 我们可得到 Cn 的显式表达式.\par
}
\end{sourceblock}


\begin{explainblock}
这一段是原书论述链的一部分。读的时候把它当作桥梁：它通常负责把上一个定义、例子或定理，引到下一个计数方法。

\smallskip
\textbf{初学者检查点：}
\begin{itemize}[leftmargin=2em,itemsep=0.25em]
\item 找出这一段承接的上一个概念。
\item 找出它准备引出的下一个概念。
\item 把中间的关键等式或构造单独抄出来。
\end{itemize}

\smallskip
\textbf{本章读法提醒：} 第八章 Catalan 问题的核心常常是前缀不越界或递归分解。
\end{explainblock}


\begin{sourceblock}{定理 8.1.4 设 Cn 为第 n 个 Catalan 数, 则}
{\small\sloppy
\noindent 定理 8.1.4 设 Cn 为第 n 个 Catalan 数, 则\par
\noindent Cn = = .\par
\noindent n+1 n n!(n + 1)!\par
\noindent 150 第 8 章 Catalan 数与 Narayana 数\par
\noindent 证明: 由牛顿二项式定理可知:\par
\noindent √ ∞ 1\par
\noindent !\par
\noindent 1− 1 − 4x 1 X\par
\noindent c(x) = = 1− 2 (−1)n 4n xn\par
\noindent 2x 2x n\par
\noindent n=0\par
\noindent ∞ 1\par
\noindent !\par
\noindent =− 2 (−1)n 4n xn ,\par
\noindent 2x n\par
\noindent n=1\par
\noindent 注意到:\par
\noindent ∞ 1\par
}
\end{sourceblock}


\begin{explainblock}
定理段不要只看结论，要看它把哪两个计数方式连在一起。组合学证明最常见的技巧是：左边按一种标准数，右边按另一种标准数，然后说明两边数的是同一批对象。

\smallskip
\textbf{初学者检查点：}
\begin{itemize}[leftmargin=2em,itemsep=0.25em]
\item 先复述定理的假设，不要把适用范围扩大。
\item 找出证明中的基本计数原则：乘法、加法、双射、递推、容斥或生成函数。
\item 检查最后的等式化简是否和前面的对象解释一致。
\end{itemize}

\smallskip
\textbf{本章读法提醒：} 第八章 Catalan 问题的核心常常是前缀不越界或递归分解。
\end{explainblock}


\begin{sourceblock}{第 8 章续读}
{\small\sloppy
\noindent X X\par
\noindent 2 (−1)n 4n xn = (−1) · × 4n · xn\par
\noindent n n!\par
\noindent n=1 n≥1\par
\noindent X 1 × 2 × 3 × · · · × (2n − 3) × (2n − 2)\par
\noindent =− × 2n · xn\par
\noindent 2 × 4 × 6 × · · · (2n − 2) × n!\par
\noindent n≥1\par
\noindent X 2 · (2(n − 1))!\par
\noindent =− xn\par
\noindent (n − 1)!n!\par
\noindent n≥1\par
\noindent X 1 2n\par
\noindent = −2x xn ,\par
\noindent n+1 n\par
\noindent n≥0\par
\noindent 所以\par
\noindent ∞\par
}
\end{sourceblock}


\begin{explainblock}
这一段是原书论述链的一部分。读的时候把它当作桥梁：它通常负责把上一个定义、例子或定理，引到下一个计数方法。

\smallskip
\textbf{初学者检查点：}
\begin{itemize}[leftmargin=2em,itemsep=0.25em]
\item 找出这一段承接的上一个概念。
\item 找出它准备引出的下一个概念。
\item 把中间的关键等式或构造单独抄出来。
\end{itemize}

\smallskip
\textbf{本章读法提醒：} 第八章 Catalan 问题的核心常常是前缀不越界或递归分解。
\end{explainblock}


\begin{sourceblock}{第 8 章续读}
{\small\sloppy
\noindent X\par
\noindent c(x) = x ,\par
\noindent n+1 n\par
\noindent n=0\par
\noindent 故 1\par
\noindent Cn = − 2 (−4)n+1 = = .\par
\noindent 2 n+1 n+1 n n!(n + 1)!\par
}
\end{sourceblock}


\begin{explainblock}
这一段是原书论述链的一部分。读的时候把它当作桥梁：它通常负责把上一个定义、例子或定理，引到下一个计数方法。

\smallskip
\textbf{初学者检查点：}
\begin{itemize}[leftmargin=2em,itemsep=0.25em]
\item 找出这一段承接的上一个概念。
\item 找出它准备引出的下一个概念。
\item 把中间的关键等式或构造单独抄出来。
\end{itemize}

\smallskip
\textbf{本章读法提醒：} 第八章 Catalan 问题的核心常常是前缀不越界或递归分解。
\end{explainblock}


\begin{sourceblock}{注记 8.1.5 这里 Cn = n+1 n 一般认为是 Cn 的标准写法, 但有时为了将 Cn 写}
{\small\sloppy
\noindent 注记 8.1.5 这里 Cn = n+1 n 一般认为是 Cn 的标准写法, 但有时为了将 Cn 写\par
\noindent 成\par
\noindent Cn = ,\par
\noindent 2n + 1 n\par
\noindent 或\par
\noindent Cn = , n ≥ 1, C0 = 1,\par
\noindent n n−1\par
\noindent 可能会更方便.\par
}
\end{sourceblock}


\begin{explainblock}
注记通常是在补充术语、记号、历史或一个容易忽略的约定。它未必给新公式，但常常决定你后面会不会把符号读错。

\smallskip
\textbf{初学者检查点：}
\begin{itemize}[leftmargin=2em,itemsep=0.25em]
\item 记下新符号的读法。
\item 确认这个约定是否只在本章使用。
\item 把注记里的限制条件补到前面的定义旁边。
\end{itemize}

\smallskip
\textbf{本章读法提醒：} 第八章 Catalan 问题的核心常常是前缀不越界或递归分解。
\end{explainblock}


\begin{sourceblock}{例 8.1.1 集合 [n] 上的划分 π = \{B1 , B2 , . . . , Bk \}, 若对任意 1 ≤ a < b < c <}
{\small\sloppy
\noindent 例 8.1.1 集合 [n] 上的划分 π = \{B1 , B2 , . . . , Bk \}, 若对任意 1 ≤ a < b < c <\par
\noindent d ≤ n, 当 \{a, c\} ⊂ Bi , \{b, d\} ⊂ Bj 时推得 i = j, 则称该划分是不交划分. 如图 8.2 是\par
\noindent 集合 [9] 上的一个不交划分. 设 N C(n) 计数集合 [n] 的不交划分的个数, 求 N C(n).\par
\noindent 1̇ 2̇ 3̇ 4̇ 5̇ 6̇ 7̇ 8̇ 9̇\par
\noindent 8.1 Catalan 数与 Dyck 路 151\par
\noindent 解: 设 π 是集合 [n] 上的一个不交划分, 设 1 所在块是 B1 , 若 1 单独一块, 则剩\par
\noindent 余的划分可以看作集合 \{2, 3, . . . , n\} 的不交划分, 有 N C(n − 1) 种情形; 若 1 不是\par
\noindent 单独一块的, 设在 B1 中除 1 外最小的数记为 k; 去掉元素 1, 则该不交划分分成了两\par
\noindent 部分; 它可以看作集合 \{2, 3, . . . , k − 1\} 的不交划分和集合 \{k, k + 1, · · · , n\} 的不交\par
\noindent 划分的并; 所以共计有 N C(k − 2)N C(n − k + 1) 种情形; 所以\par
\noindent X\par
\noindent n\par
\noindent N C(n) = N C(n − 1) + N C(k − 2)N C(n − k + 1)\par
\noindent k=2\par
\noindent X\par
\noindent n−1\par
\noindent = N C(i)N C(n − 1 − i).\par
\noindent i=0\par
}
\end{sourceblock}


\begin{explainblock}
例题是学习这本书最重要的部分。不要只看答案，要把每个限制条件变成一个计数动作：先安排谁、再插入谁、有没有重复、需不需要除以对称性。

\smallskip
\textbf{初学者检查点：}
\begin{itemize}[leftmargin=2em,itemsep=0.25em]
\item 先写出没有限制时的总数。
\item 逐条加入限制条件，观察总数怎样变化。
\item 最后检查有没有把同一种方案重复算了多次。
\end{itemize}

\smallskip
\textbf{本章读法提醒：} 第八章 Catalan 问题的核心常常是前缀不越界或递归分解。
\end{explainblock}


\begin{sourceblock}{第 8 章续读}
{\small\sloppy
\noindent 取 N C(0) = 1, 则 N C(n) 与 Catalan 数 Cn 的递推关系及初值一样, 故 N C(n) = Cn .\par
}
\end{sourceblock}


\begin{explainblock}
这一段是原书论述链的一部分。读的时候把它当作桥梁：它通常负责把上一个定义、例子或定理，引到下一个计数方法。

\smallskip
\textbf{初学者检查点：}
\begin{itemize}[leftmargin=2em,itemsep=0.25em]
\item 找出这一段承接的上一个概念。
\item 找出它准备引出的下一个概念。
\item 把中间的关键等式或构造单独抄出来。
\end{itemize}

\smallskip
\textbf{本章读法提醒：} 第八章 Catalan 问题的核心常常是前缀不越界或递归分解。
\end{explainblock}


\begin{sourceblock}{例 8.1.2 n 个非结合的二元加法运算由 n 对括号决定顺序, 问可以有多少种加}
{\small\sloppy
\noindent 例 8.1.2 n 个非结合的二元加法运算由 n 对括号决定顺序, 问可以有多少种加\par
\noindent 括号的方式? 例如, n = 3 时共有 5 种加括号的方式:\par
\noindent (((a + b) + c + d), ((a + b) + (c + d)), ((a + (b + c)) + d),\par
\noindent a + ((b + c) + d), (a + (b + (c + d))).\par
\noindent 解: 我们不妨考虑只有左括号 “(” 和右括号 “)” 构成且满足\par
\noindent 1) 总长度为 2n , 且恰好包含有 n 个左括号和 n 个右括号,\par
\noindent 2) 从左往右数, 左括号个数永远不小于右括号个数, 的长为 2n 的合法括号对序列,\par
\noindent 其个数设为 f (n), 并令 f (0)=1. 我们称一个 2n 长的括号对序列是本原的, 如果去掉\par
\noindent 它的第一个左括号和最后一个右括号后的括号序列依旧是合法括号对序列. 换言之,\par
\noindent 本原的合法括号序列是指从左到右在达到尽头之前任何位置左括号的个数总是大于\par
\noindent 右括号的个数. 长为 2n 的本原的合法括号序列的个数设为 g(n). 例如, 当 n = 3 时,\par
\noindent 合法括号对序列有:\par
\noindent ()()(); ()(()); (())(); (()()); ((())).\par
\noindent 其中后 2 个是本原的, 也即 f (3) = 5, g(3) = 2. 当 n ≥ 1 时, 因为任意一个 2(n + 1)\par
\noindent 长的对本原的合法括号序列, 去掉后前后左右括号就得到了一个 2n 长的合法括号序\par
\noindent 列, 而任意一个 2n 长的合法括号序列在前面加上左括号, 在后面加上右括号也得到\par
\noindent 了一个 2(n + 1) 长的对本原的合法括号序列, 故而有\par
\noindent f (n) = g(n + 1). (8.1)\par
}
\end{sourceblock}


\begin{explainblock}
例题是学习这本书最重要的部分。不要只看答案，要把每个限制条件变成一个计数动作：先安排谁、再插入谁、有没有重复、需不需要除以对称性。

\smallskip
\textbf{初学者检查点：}
\begin{itemize}[leftmargin=2em,itemsep=0.25em]
\item 先写出没有限制时的总数。
\item 逐条加入限制条件，观察总数怎样变化。
\item 最后检查有没有把同一种方案重复算了多次。
\end{itemize}

\smallskip
\textbf{本章读法提醒：} 第八章 Catalan 问题的核心常常是前缀不越界或递归分解。
\end{explainblock}


\begin{sourceblock}{第 8 章续读}
{\small\sloppy
\noindent 设 k(k ≥ 1) 为一合法括号序列中从左开始第一次达到左, 右括号数目相等时括号对\par
\noindent 的个数, 则这 2k 长的括号序列显然是本原的合法序列, 而剩下的是 2n − 2k 长的合\par
\noindent 法括号序列, 故当 n ≥ 1 时, f (n) 有下面的递推关系:\par
\noindent X\par
\noindent n X\par
\noindent n\par
\noindent f (n) = g(k)f (n − k) = f (k − 1)f (n − k).\par
\noindent k=1 k=1\par
\noindent 152 第 8 章 Catalan 数与 Narayana 数\par
\noindent 故 f (n) 与 Catalan 数 Cn 具有相同的初值及递推关系, 所以 f (n) = Cn = 2n n .\par
}
\end{sourceblock}


\begin{explainblock}
这一段是原书论述链的一部分。读的时候把它当作桥梁：它通常负责把上一个定义、例子或定理，引到下一个计数方法。

\smallskip
\textbf{初学者检查点：}
\begin{itemize}[leftmargin=2em,itemsep=0.25em]
\item 找出这一段承接的上一个概念。
\item 找出它准备引出的下一个概念。
\item 把中间的关键等式或构造单独抄出来。
\end{itemize}

\smallskip
\textbf{本章读法提醒：} 第八章 Catalan 问题的核心常常是前缀不越界或递归分解。
\end{explainblock}


\begin{sourceblock}{注记 8.1.6 在例 8.1.2 中, 如果将左括号都换成 1, 右括号都换成 −1, 就得到了由}
{\small\sloppy
\noindent 注记 8.1.6 在例 8.1.2 中, 如果将左括号都换成 1, 右括号都换成 −1, 就得到了由\par
\noindent n 个 1 和 n 个 −1 组成的长为 2n 的 1, −1 序列, 且从左到右 1 的个数始终不小于\par
\noindent −1 的个数. 换言之, 若定义这个 2n 长序列为 a1 a2 · · · a2n , 则对任意 k(1 ≤ k ≤ 2n),\par
\noindent 都有\par
\noindent a1 + a2 + a3 + · · · + ak−1 + ak ≥ 0. (8.2)\par
\noindent 因此, 我们也可将 Catalan 数 Cn 定义为计数由 −1 与 1 组成的长为 2n 且对任意\par
\noindent k(1 ≤ k ≤ 2n) 都满足 (8.2) 的序列的个数.\par
}
\end{sourceblock}


\begin{explainblock}
注记通常是在补充术语、记号、历史或一个容易忽略的约定。它未必给新公式，但常常决定你后面会不会把符号读错。

\smallskip
\textbf{初学者检查点：}
\begin{itemize}[leftmargin=2em,itemsep=0.25em]
\item 记下新符号的读法。
\item 确认这个约定是否只在本章使用。
\item 把注记里的限制条件补到前面的定义旁边。
\end{itemize}

\smallskip
\textbf{本章读法提醒：} 第八章 Catalan 问题的核心常常是前缀不越界或递归分解。
\end{explainblock}


\begin{sourceblock}{定义 8.1.7 在 (x, y) 平面上定义两种移动方式, 称沿 (1, 1) 方向 ( 东北方向) 的移}
{\small\sloppy
\noindent 定义 8.1.7 在 (x, y) 平面上定义两种移动方式, 称沿 (1, 1) 方向 ( 东北方向) 的移\par
\noindent 动为一个上升, 记为 u. 沿 (1, −1) 方向 (东南方向) 的移动为一个下降, 记为 d. 经\par
\noindent 这两种移动方式, 从坐标原点 (0, 0) 处到 (2n, 0) 的移动路径称为长为 2n 的自由的\par
\noindent Dyck 路. 对于任一个长为 2n 的自由 Dyck 路在 x 下方的下降称为它的裂缝 (flaw),\par
\noindent 具有 k(0 ≤ k ≤ n) 个裂缝的自由 n Dyck 路的个数记为 Dn,k . 特别地, 裂缝个数为\par
\noindent 0 的 2n 长的自由 Dyck 路称为长为 2n 的 Dyck 路, 换言之, 2n 长的 Dyck 路是指\par
\noindent 从 (0, 0) 到 (2n, 0) 经上升与下降移动且不经过 x 轴的下方的路径. 长为 2n 的 Dyck\par
\noindent 的个数简记为 Dn . 对任意一条长为 2n 的 (自由) Dyck 路, 若将其移动方式用 u 与\par
\noindent d 来表示, 得到的只有 u 与 d 的序列称为 2n 长的 (自由) Dyck 字.\par
}
\end{sourceblock}


\begin{explainblock}
定义段的任务是定边界。这里要特别注意参数范围、是否允许重复、是否考虑顺序、对象有没有标签。后面的每一道例题和证明，本质上都在反复使用这些边界。

\smallskip
\textbf{初学者检查点：}
\begin{itemize}[leftmargin=2em,itemsep=0.25em]
\item 圈出被定义的对象名称。
\item 圈出参数范围，例如 n、k 是否要求正整数。
\item 判断它是有序对象、无序对象、带标签对象还是结构对象。
\end{itemize}

\smallskip
\textbf{本章读法提醒：} 第八章 Catalan 问题的核心常常是前缀不越界或递归分解。
\end{explainblock}


\begin{sourceblock}{例 8.1.3 如图 8.3, n = 2 时, 所有自由的 Dyck 路.}
{\small\sloppy
\noindent 例 8.1.3 如图 8.3, n = 2 时, 所有自由的 Dyck 路.\par
\noindent y y y\par
\noindent x x x\par
\noindent 字: udud; 裂缝个数为 0 字: uudd; 裂缝个数为 0 字: uddu; 裂缝个数为 1\par
\noindent y y y\par
\noindent x x x\par
\noindent 字: duud; 裂缝个数为 1 字: dduu; 裂缝个数为 2 字: dudu; 裂缝个数为 2\par
}
\end{sourceblock}


\begin{explainblock}
例题是学习这本书最重要的部分。不要只看答案，要把每个限制条件变成一个计数动作：先安排谁、再插入谁、有没有重复、需不需要除以对称性。

\smallskip
\textbf{初学者检查点：}
\begin{itemize}[leftmargin=2em,itemsep=0.25em]
\item 先写出没有限制时的总数。
\item 逐条加入限制条件，观察总数怎样变化。
\item 最后检查有没有把同一种方案重复算了多次。
\end{itemize}

\smallskip
\textbf{本章读法提醒：} 第八章 Catalan 问题的核心常常是前缀不越界或递归分解。
\end{explainblock}


\begin{sourceblock}{定理 8.1.8 设长为 2n 的 Dyck 路的个数为 Dn , Cn 为 Catalan 数, 则}
{\small\sloppy
\noindent 定理 8.1.8 设长为 2n 的 Dyck 路的个数为 Dn , Cn 为 Catalan 数, 则\par
\noindent D n = Cn = . (8.3)\par
\noindent n+1 n\par
\noindent 证明: (方法一) 对于任意 2n 长的 Dyck 路, 若它第一次到达 x 轴的坐标为\par
\noindent (k, 0), 其中 1 ≤ k ≤ n. 则这条 2n 长 Dyck 路可依此为界将它分成了一条从 (0, 0)\par
\noindent 到 (2k, 0) 的路径和从 (2k, 0) 到 (2n, 0) 的路径; 其中从 (0, 0) 到 (2k, 0) 的路径, 去\par
\noindent 掉第一步和最后一步, 就可以看作一条 2(k − 1) 长的 Dyck 路; 而从 (2k, 0) 到 (2n, 0)\par
\noindent 8.1 Catalan 数与 Dyck 路 153\par
\noindent 的路径也可以看作一条 2(n − k) 长的路径; 因此有下面的递推关系:\par
\noindent X\par
\noindent n\par
\noindent Dn = Dk−1 Dn−k .\par
\noindent k=1\par
\noindent 所以 Dn 与 Catalan 数的初值及递推关系一样, 因此,\par
\noindent Dn = .\par
\noindent n+1 n\par
\noindent (方法二) 由 Dyck 路的定义可知在长为 2n 的 Dyck 中恰有 n 个 u 和 n 个 d, 且\par
\noindent 从左往右读, d 的个数不超过 u 的个数. 若不然, 则它对应的路径必然要越过 x 轴,\par
}
\end{sourceblock}


\begin{explainblock}
定理段不要只看结论，要看它把哪两个计数方式连在一起。组合学证明最常见的技巧是：左边按一种标准数，右边按另一种标准数，然后说明两边数的是同一批对象。

\smallskip
\textbf{初学者检查点：}
\begin{itemize}[leftmargin=2em,itemsep=0.25em]
\item 先复述定理的假设，不要把适用范围扩大。
\item 找出证明中的基本计数原则：乘法、加法、双射、递推、容斥或生成函数。
\item 检查最后的等式化简是否和前面的对象解释一致。
\end{itemize}

\smallskip
\textbf{本章读法提醒：} 第八章 Catalan 问题的核心常常是前缀不越界或递归分解。
\end{explainblock}


\begin{sourceblock}{第 8 章续读}
{\small\sloppy
\noindent 不合 Dyck 路的定义. 现将 Dyck 字中的 u, d 分别用 1, −1 替换得到一个长为 2n 的\par
\noindent 1, −1 序列, 同时从左到右 −1 的个数始终不大于 1 的个数, 所以由注记 8.1.6 知长为\par
\noindent 2n 的 Dyck 字的个数对应第 n 个 Catalan 数, 而 Dyck 字与 Dyck 路一一对应, 所\par
\noindent 以定理成立.\par
\noindent 2n\par
\noindent (方法三) 考虑所有从 (0, 0) 到 (2n, 0) 的自由 Dyck 路, 显然共计有 n 条. 记\par
\noindent P 为任意一条从 (0, 0) 到 (2n, 0) 的非 Dyck 路的自由 Dyck 路, 则它必然越过 x 轴\par
\noindent 与 y = −1 的直线相交. 下面我们建立从 (0, 0) 到 (2n, 0) 的非 Dyck 路的自由 Dyck\par
\noindent 路到从 (0, −2) 到 (2n, 0) 经 (1, 1), (1, −1) 两种移动方式的路径的 1 − 1 映射. 如图\par
\noindent y\par
\noindent (0, 0)\par
\noindent x\par
\noindent −1\par
\noindent A\par
\noindent 1) 找出路径与 y = −1 相交的第一个点, 记为 A.\par
\noindent 2) 将 A 点之前的路径沿 y = −1 翻折, 得到 (0, −2) 到 (2n, 0) 的一条路径.\par
\noindent 显然变换过程是可逆的. 又因为从 (0, −2) 到 (2n, 0) 的移动路径恰好有 2n\par
\noindent n−1 个, 故\par
}
\end{sourceblock}


\begin{explainblock}
这一段是原书论述链的一部分。读的时候把它当作桥梁：它通常负责把上一个定义、例子或定理，引到下一个计数方法。

\smallskip
\textbf{初学者检查点：}
\begin{itemize}[leftmargin=2em,itemsep=0.25em]
\item 找出这一段承接的上一个概念。
\item 找出它准备引出的下一个概念。
\item 把中间的关键等式或构造单独抄出来。
\end{itemize}

\smallskip
\textbf{本章读法提醒：} 第八章 Catalan 问题的核心常常是前缀不越界或递归分解。
\end{explainblock}


\begin{sourceblock}{第 8 章续读}
{\small\sloppy
\noindent 经 (1, 1), (1, −1) 两种移动方式从原点 (0, 0) 到 (2n, 0) 所有路径中经过 x 轴下方的\par
\noindent 2n\par
\noindent 路径个数为 n−1 . 所以 2n 长的 Dyck 路个数为:\par
\noindent 2n 2n 1 2n\par
\noindent |Dn | = − = .\par
\noindent n n−1 n+1 n\par
\noindent 这个方法用到技巧称为反射原理 (Reflection Principle), 该技巧由 Désiré André∗\par
\noindent ∗\par
\noindent Désiré André 法国数学家.\par
\noindent 154 第 8 章 Catalan 数与 Narayana 数\par
\noindent 首先使用, 并在组合证明中有很好的应用, 有兴趣的读者可以参考 [10].\par
}
\end{sourceblock}


\begin{explainblock}
这一段是原书论述链的一部分。读的时候把它当作桥梁：它通常负责把上一个定义、例子或定理，引到下一个计数方法。

\smallskip
\textbf{初学者检查点：}
\begin{itemize}[leftmargin=2em,itemsep=0.25em]
\item 找出这一段承接的上一个概念。
\item 找出它准备引出的下一个概念。
\item 把中间的关键等式或构造单独抄出来。
\end{itemize}

\smallskip
\textbf{本章读法提醒：} 第八章 Catalan 问题的核心常常是前缀不越界或递归分解。
\end{explainblock}


\begin{sourceblock}{例 8.1.4 在 2 × n 的表格中填入 1, 2, . . . , 2n, 如果表格中的数字按行与列递增,}
{\small\sloppy
\noindent 例 8.1.4 在 2 × n 的表格中填入 1, 2, . . . , 2n, 如果表格中的数字按行与列递增,\par
\noindent 我们称该表格为形如 2 × n 的标准 Young 表, 记它的个数为 Yn . 证明:\par
\noindent Yn = Dn = .\par
\noindent n+1 n\par
\noindent 证明: 我们通过建立 2 × n 形的标准 Young 表与 2n 长的 Dyck 路径之间的双\par
\noindent 射来证明 Yn = Dn .\par
\noindent 对于任意的一个 2 × n 形的标准 Young 表 λ, 通过在表中第 1 行和第 2 行所填\par
\noindent 数字的后面分别添加字母 u 和 d 得到一个新表 µ. 按字母前数字从小到大的顺序将\par
\noindent 字母写成一个长为 2n 的字 w, 显然字 w 首字母为 u 尾字母为 d, 且含有 n 个字母\par
\noindent u 和 n 个字母 d. 若设在字 w 中第一次出现字母 d 个数大于字母 u 的个数位置为\par
\noindent 第 k 个位置, 则 k 显然为奇数, 则数字 k 填在标准 Young 表的第 2 行, 而在数字 k\par
\noindent 的方格中所填的数字都将比 k 要大, 这与标准 Young 表的定义矛盾. 所以字 w 满足\par
\noindent 从左至右字母 u 的个数不比字母 d 的个数少. 分别将 u 和 d 对应 Dyck 路径中的上\par
\noindent 升和下降移动, 则字 w 为 Dyck 字. 而每一个 2n 长的 Dyck 字与 2n 长的 Dyck 路\par
\noindent 径一一对应.\par
}
\end{sourceblock}


\begin{explainblock}
例题是学习这本书最重要的部分。不要只看答案，要把每个限制条件变成一个计数动作：先安排谁、再插入谁、有没有重复、需不需要除以对称性。

\smallskip
\textbf{初学者检查点：}
\begin{itemize}[leftmargin=2em,itemsep=0.25em]
\item 先写出没有限制时的总数。
\item 逐条加入限制条件，观察总数怎样变化。
\item 最后检查有没有把同一种方案重复算了多次。
\end{itemize}

\smallskip
\textbf{本章读法提醒：} 第八章 Catalan 问题的核心常常是前缀不越界或递归分解。
\end{explainblock}


\begin{sourceblock}{例如 n = 3 时的 2 × 3 形的标准 Young 表 5 个:}
{\small\sloppy
\noindent 例如 n = 3 时的 2 × 3 形的标准 Young 表 5 个:\par
\noindent 1u 2u 3u 1u 2u 5u 1u 2u 4u\par
\noindent , , ,\par
\noindent 4d 5d 6d 3d 4d 6d 3d 5d 6d\par
\noindent 1u 3u 4u 1u 3u 5u\par
\noindent , .\par
\noindent 2d 5d 6d 2d 4d 6d\par
\noindent 所以对应的 Dyck 字和 Dyck 路如图 8.4.\par
\noindent y y y\par
\noindent x x x\par
\noindent 字:uuuddd 字:uuddud 字:uududd\par
\noindent y y\par
\noindent x x\par
\noindent 字:uduudd 字:ududud\par
\noindent 可验证上述变换是可逆的. 所以 2 × n 形的标准 Young 表的个数等于 2n 长的\par
\noindent Dyck 路径的个数, 均为 Catalan 数.\par
}
\end{sourceblock}


\begin{explainblock}
例题是学习这本书最重要的部分。不要只看答案，要把每个限制条件变成一个计数动作：先安排谁、再插入谁、有没有重复、需不需要除以对称性。

\smallskip
\textbf{初学者检查点：}
\begin{itemize}[leftmargin=2em,itemsep=0.25em]
\item 先写出没有限制时的总数。
\item 逐条加入限制条件，观察总数怎样变化。
\item 最后检查有没有把同一种方案重复算了多次。
\end{itemize}

\smallskip
\textbf{本章读法提醒：} 第八章 Catalan 问题的核心常常是前缀不越界或递归分解。
\end{explainblock}


\begin{sourceblock}{注记 8.1.9 设 2 × n 形标准 Young 表的第 2 行中第一次出现偶数的列为第 k 列,}
{\small\sloppy
\noindent 注记 8.1.9 设 2 × n 形标准 Young 表的第 2 行中第一次出现偶数的列为第 k 列,\par
\noindent 若将后面 n − k 列中的数字都减小 2k 则后面形成了 2 × (n − k) 形的标准 Young\par
\noindent 8.1 Catalan 数与 Dyck 路 155\par
\noindent 表; 而前 k 列如果去掉 1 和 2k, 然后将第 1 行的数字均向前移动一格, 再将格子中\par
\noindent 数字都减 1 则得到了一个 2 × (k − 1) 的标准 Young 表; 设 2 × n 的标准 Young 表\par
\noindent 的个数为 Y (n), 规定 Y (0) = 1, 则满足下面递推关系:\par
\noindent X\par
\noindent n\par
\noindent Y (n + 1) = Y (k)Y (n − k),\par
\noindent k=0\par
\noindent 可推知 Y (n) 为 Catalan 数 C(n).\par
\noindent 下面的定理是我国著名概率学家钟开莱与 W. Feller 于 1949 年发表的论文 [34]\par
\noindent 中提出的 Chung-Feller 定理的一个特殊情况. 它表明自由 Dyck 路中裂缝的个数的\par
\noindent 分布是均匀的. 在 [34] 中, 作者是通过研究一个特殊的抛硬币问题来得到的这个结\par
\noindent 论, 在文中作者证明 Chung-Feller 定理的方法有浓重的分析特色. 这里我们给出它\par
\noindent 的一个组合证明, 它还有其它的证明, 有兴趣的同学可以参考 [6].\par
}
\end{sourceblock}


\begin{explainblock}
注记通常是在补充术语、记号、历史或一个容易忽略的约定。它未必给新公式，但常常决定你后面会不会把符号读错。

\smallskip
\textbf{初学者检查点：}
\begin{itemize}[leftmargin=2em,itemsep=0.25em]
\item 记下新符号的读法。
\item 确认这个约定是否只在本章使用。
\item 把注记里的限制条件补到前面的定义旁边。
\end{itemize}

\smallskip
\textbf{本章读法提醒：} 第八章 Catalan 问题的核心常常是前缀不越界或递归分解。
\end{explainblock}


\begin{sourceblock}{定理 8.1.10 (Chung-Feller 定理) 设 2n 长的自由 Dyck 路径中有 k 个裂缝的个数}
{\small\sloppy
\noindent 定理 8.1.10 (Chung-Feller 定理) 设 2n 长的自由 Dyck 路径中有 k 个裂缝的个数\par
\noindent 为 Dn,k 个, 其中 0 ≤ k ≤ n, Cn 为第 n 个 Catalan 数, 则\par
\noindent Dn,0 = Dn,1 = · · · = Dn,n = Cn = .\par
\noindent n+1 n\par
\noindent 证明: 依旧用 u, d 分别表示上升移动 (1, 1) 与下降移动 (1, −1). 首先, 我们考\par
\noindent 虑裂缝个数为 0 时 (Dyck 路) 与裂缝个数为 1 之间的双射. 对任意一个 Dyck 路,\par
\noindent 我们考虑它从开始到第一次下降到 x 轴上的路径对应的 Dyck 字, 记为 uAdα, 其\par
\noindent 中 A, α 是 Dyck 字或空字. 现将 uAdα 变换为 Aduα, 就得到了一个裂缝为 1 的自\par
\noindent 由 Dyck 字, 对应的 Dyck 路到裂缝为 1 的自由 Dyck 路间的变换如图 (8.5) 所示;\par
\noindent 反之, 对任意一个裂缝为 1 的自由的 Dyck 路对应的字 Aduα (其中 A, α 为空字或\par
\noindent Dyck 字) 可将 u 调到 Ad 前面变换为 uAdα, 显然这是一个 Dyck 字 (如图 8.5). 也\par
\noindent 即是说, 上面的过程建立了从 Dyck 路到裂缝为 1 的自由 Dyck 路之间的双射.\par
\noindent 因为对任意的一个有 k(1 ≤ k < n) 个裂缝的自由 Dyck 路都可以表示为 BuAdα,\par
\noindent 其中 A 是 Dyck 路 (字) 或空路 (字), B 是在 x 轴下方的路径或空路径, α 为空路\par
\noindent (字) 或自由 Dyck 路 (字); 也即是说, u 是从 x 轴第一次的上升移动, d 是第一次下\par
\noindent 降到 x 轴如图 8.6. 所以可以类似的建立具有 k 个裂缝的自由 Dyck 路与 k + 1 个\par
\noindent 裂缝的自由路之间的双射. 具体来说, 通过将 Bu 与 Ad 调换位置得到 AdBuα(如图\par
\noindent 8.6, 此时 d 成为 x 轴下方新的下降移动, 也即 AdBuα 中裂缝的个数比 BuAdα 裂\par
}
\end{sourceblock}


\begin{explainblock}
定理段不要只看结论，要看它把哪两个计数方式连在一起。组合学证明最常见的技巧是：左边按一种标准数，右边按另一种标准数，然后说明两边数的是同一批对象。

\smallskip
\textbf{初学者检查点：}
\begin{itemize}[leftmargin=2em,itemsep=0.25em]
\item 先复述定理的假设，不要把适用范围扩大。
\item 找出证明中的基本计数原则：乘法、加法、双射、递推、容斥或生成函数。
\item 检查最后的等式化简是否和前面的对象解释一致。
\end{itemize}

\smallskip
\textbf{本章读法提醒：} 第八章 Catalan 问题的核心常常是前缀不越界或递归分解。
\end{explainblock}


\begin{sourceblock}{第 8 章续读}
{\small\sloppy
\noindent 缝的个数多 1. 这里需要说明的是对于自由 Dyck 路 AdBuα 中的 d 表示在 x 轴下\par
\noindent 方的第一次下降移动, u 表示从 x 轴下方第一次到达 x 轴的上升移动. 因此, 只需找\par
\noindent 到具有 k + 1 个裂缝的自由 Dyck 路中在 x 轴下方的第一次下降移动的位置与从 x\par
\noindent 轴下方第一次到达 x 轴的上升移动的位置, 将其表示成 AdBuα 的形式, 通过调换\par
\noindent Ad 与 Bu 的即可得到具有 k 个裂缝的自由路 BuAdα, 过程可逆的. 从而建立了从\par
\noindent 具有 k 个裂缝的自由 Dyck 路到 k + 1 个裂缝的自由 Dyck 路的双射. 这样我们证\par
\noindent 明了有 k(0 ≤ k ≤ n) 个裂缝的 2n 长自由 Dyck 路的个数与 k 无关. 而易知从原点\par
\noindent 到 (2n, 0) 的 2n 长自由 Dyck 路有 2n\par
\noindent n 个. 从而有 k 个裂缝的 2n 长自由 Dyck 路\par
\noindent 的个数为 Cn .\par
\noindent 156 第 8 章 Catalan 数与 Narayana 数\par
\noindent y\par
\noindent A\par
\noindent u d α\par
\noindent x\par
\noindent y\par
\noindent A α\par
\noindent d u x\par
}
\end{sourceblock}


\begin{explainblock}
这一段是原书论述链的一部分。读的时候把它当作桥梁：它通常负责把上一个定义、例子或定理，引到下一个计数方法。

\smallskip
\textbf{初学者检查点：}
\begin{itemize}[leftmargin=2em,itemsep=0.25em]
\item 找出这一段承接的上一个概念。
\item 找出它准备引出的下一个概念。
\item 把中间的关键等式或构造单独抄出来。
\end{itemize}

\smallskip
\textbf{本章读法提醒：} 第八章 Catalan 问题的核心常常是前缀不越界或递归分解。
\end{explainblock}


\begin{sourceblock}{§ 8.2 Narayana 数与 Narayana 多项式}
{\small\sloppy
\noindent § 8.2. Narayana 数与 Narayana 多项式\par
\noindent 这节我们考虑与 Dyck 路有关的另一个组合数—Narayana 数. 首先给出 Dyck 路\par
\noindent 峰与谷的定义.\par
}
\end{sourceblock}


\begin{explainblock}
这一节先不要急着背公式。先问三个问题：对象是什么，哪些限制条件不能丢，最后到底是在数所有对象、还是在数某个统计量的分布。把这三件事说清楚，后面的公式才会有位置。

\smallskip
\textbf{初学者检查点：}
\begin{itemize}[leftmargin=2em,itemsep=0.25em]
\item 把本节出现的新对象写成一句话定义。
\item 把每个公式左边和右边分别翻译成“在数什么”。
\item 找一个最小非平凡例子，手工列出所有对象。
\end{itemize}

\smallskip
\textbf{本章读法提醒：} 第八章 Catalan 问题的核心常常是前缀不越界或递归分解。
\end{explainblock}


\begin{sourceblock}{定义 8.2.1 对任意一条 Dyck 路, 其移动方式由上升变为下降称为峰; 反之, 移动}
{\small\sloppy
\noindent 定义 8.2.1 对任意一条 Dyck 路, 其移动方式由上升变为下降称为峰; 反之, 移动\par
\noindent 方式由下降变为上升称为谷.\par
\noindent 若用 u, d 分别表示生成 Dyck 路所需的上升和下降移动, 则任意一条 Dyck 路都\par
\noindent 对应一个由字母 u 和 d 组成的 Dyck 字. Dyck 路的峰和谷的位置则恰是 Dyck 字\par
\noindent 中 u 和 d 相邻的位置, 且 ud 形成峰, du 形成谷.\par
\noindent 如图 8.7, 3−Dyck 路中的峰与谷的位置与对应的 Dyck 字中 ud 与 du 出现的位\par
\noindent 置:\par
\noindent 可以验证任意一条 Dyck 路的峰的个数比谷的个数多 1. 例如图 8.7 中, 第一条\par
\noindent Dyck 路有一个峰; 中间三条 Dyck 路各有两个峰, 一个谷; 最后一条 Dyck 路有三个\par
\noindent 峰, 两个谷. 由峰与谷的定义可知对任意一条 2n 长的 Dyck 路, 其形状完全由其峰\par
\noindent 与谷的 (个数与) 位置完全确定.\par
}
\end{sourceblock}


\begin{explainblock}
定义段的任务是定边界。这里要特别注意参数范围、是否允许重复、是否考虑顺序、对象有没有标签。后面的每一道例题和证明，本质上都在反复使用这些边界。

\smallskip
\textbf{初学者检查点：}
\begin{itemize}[leftmargin=2em,itemsep=0.25em]
\item 圈出被定义的对象名称。
\item 圈出参数范围，例如 n、k 是否要求正整数。
\item 判断它是有序对象、无序对象、带标签对象还是结构对象。
\end{itemize}

\smallskip
\textbf{本章读法提醒：} 第八章 Catalan 问题的核心常常是前缀不越界或递归分解。
\end{explainblock}


\begin{sourceblock}{定义 8.2.2 我们称 2n 长的 Dyck 路中有 k(1 ≤ k ≤ n) 个峰 (或 k − 1 个谷) 的路}
{\small\sloppy
\noindent 定义 8.2.2 我们称 2n 长的 Dyck 路中有 k(1 ≤ k ≤ n) 个峰 (或 k − 1 个谷) 的路\par
\noindent 径个数为 Narayana 数, 记为 Nn,k .\par
\noindent 由图 8.7 可知 N3,1 = 1, N3,2 = 3, N3,3 = 1. 对于任意 2n 长的 Dyck 路, 具有一\par
\noindent 个峰的 Dyck 路和 n 个峰的 2n 长的 Dyck 路显然都只有一条, 即 Nn,1 = Nn,n = 1,\par
\noindent 分别对应 Dyck 字为 uuu · · · ud · · · ddd 与 udud · · · udud. 当然, 我们可通过规定\par
\noindent N0,0 = 1 及当 k > n ≥ 1 时, 令 Nn,k = Nn,0 = 0 使得对任意非负整数 n, k, Nn,k 都\par
\noindent 有定义.\par
\noindent 由 Narayana 数 Nn,k 定义及定理 8.1.8, 我们显然有下面的结论.\par
\noindent 8.2 Narayana 数与 Narayana 多项式 157\par
\noindent y\par
\noindent A\par
\noindent u d\par
\noindent α\par
\noindent x\par
\noindent B\par
\noindent y\par
\noindent α\par
\noindent A\par
}
\end{sourceblock}


\begin{explainblock}
定义段的任务是定边界。这里要特别注意参数范围、是否允许重复、是否考虑顺序、对象有没有标签。后面的每一道例题和证明，本质上都在反复使用这些边界。

\smallskip
\textbf{初学者检查点：}
\begin{itemize}[leftmargin=2em,itemsep=0.25em]
\item 圈出被定义的对象名称。
\item 圈出参数范围，例如 n、k 是否要求正整数。
\item 判断它是有序对象、无序对象、带标签对象还是结构对象。
\end{itemize}

\smallskip
\textbf{本章读法提醒：} 第八章 Catalan 问题的核心常常是前缀不越界或递归分解。
\end{explainblock}


\begin{sourceblock}{第 8 章续读}
{\small\sloppy
\noindent d u\par
\noindent x\par
\noindent B\par
\noindent y y y\par
\noindent 字:uuu ↑ ddd x 字:uu ↑ dd ↓ u ↑ d x 字:uu ↑ d ↓ u ↑ dd x\par
\noindent y y\par
\noindent 字:u ↑ d ↓ uu ↑ dd x 字:u ↑ d ↓ u ↑ d ↓ u ↑ dx\par
\noindent 性质 8.2.3 设 Dn 为 2n 长的 Dyck 路的个数, Cn 为第 n 个 Catalan 数, 则\par
\noindent X\par
\noindent n\par
\noindent Cn = D n = Nn,k .\par
\noindent k=1\par
\noindent 对固定的 n ≥ 0, 我们考虑 Nn,k 的生成函数, 记为 Hn (x), 则\par
\noindent X\par
\noindent n\par
\noindent Hn (x) = Nn,k xk ,\par
\noindent k=0\par
\noindent 称 Hn (x) 为 Narayana 多项式. 显然 H0 (x) = 1, H1 (x) = x. 一般的我们有下面的关\par
}
\end{sourceblock}


\begin{explainblock}
这一段是原书论述链的一部分。读的时候把它当作桥梁：它通常负责把上一个定义、例子或定理，引到下一个计数方法。

\smallskip
\textbf{初学者检查点：}
\begin{itemize}[leftmargin=2em,itemsep=0.25em]
\item 找出这一段承接的上一个概念。
\item 找出它准备引出的下一个概念。
\item 把中间的关键等式或构造单独抄出来。
\end{itemize}

\smallskip
\textbf{本章读法提醒：} 第八章 Catalan 问题的核心常常是前缀不越界或递归分解。
\end{explainblock}


\begin{sourceblock}{第 8 章续读}
{\small\sloppy
\noindent 系式.\par
\noindent 158 第 8 章 Catalan 数与 Narayana 数\par
}
\end{sourceblock}


\begin{explainblock}
这一段是原书论述链的一部分。读的时候把它当作桥梁：它通常负责把上一个定义、例子或定理，引到下一个计数方法。

\smallskip
\textbf{初学者检查点：}
\begin{itemize}[leftmargin=2em,itemsep=0.25em]
\item 找出这一段承接的上一个概念。
\item 找出它准备引出的下一个概念。
\item 把中间的关键等式或构造单独抄出来。
\end{itemize}

\smallskip
\textbf{本章读法提醒：} 第八章 Catalan 问题的核心常常是前缀不越界或递归分解。
\end{explainblock}


\begin{sourceblock}{定理 8.2.4 当 n ≥ 2 时,}
{\small\sloppy
\noindent 定理 8.2.4 当 n ≥ 2 时,\par
\noindent X\par
\noindent n−1\par
\noindent Hn (x) = xHn−1 (x) + Hi (x)Hn−1−i (x). (8.4)\par
\noindent i=1\par
\noindent 证明: 从 Dyck 路的类型出发分析.\par
}
\end{sourceblock}


\begin{explainblock}
定理段不要只看结论，要看它把哪两个计数方式连在一起。组合学证明最常见的技巧是：左边按一种标准数，右边按另一种标准数，然后说明两边数的是同一批对象。

\smallskip
\textbf{初学者检查点：}
\begin{itemize}[leftmargin=2em,itemsep=0.25em]
\item 先复述定理的假设，不要把适用范围扩大。
\item 找出证明中的基本计数原则：乘法、加法、双射、递推、容斥或生成函数。
\item 检查最后的等式化简是否和前面的对象解释一致。
\end{itemize}

\smallskip
\textbf{本章读法提醒：} 第八章 Catalan 问题的核心常常是前缀不越界或递归分解。
\end{explainblock}


\begin{sourceblock}{定理 8.2.5 设 H(x, t) 是 Hn (x) 的生成函数, 即}
{\small\sloppy
\noindent 定理 8.2.5 设 H(x, t) 是 Hn (x) 的生成函数, 即\par
\noindent ∞\par
\noindent X\par
\noindent H(x, t) = Hn (x)tn ,\par
\noindent n=0\par
\noindent 则 p\par
\noindent 1 + t − xt − (xt − t − 1)2 − 4t\par
\noindent H(x, t) = . (8.5)\par
\noindent 2t\par
\noindent 证明: 由定理 8.2.4 可知\par
\noindent ∞\par
\noindent X\par
\noindent H(x, t) = Hn (x)tn\par
\noindent n=0\par
\noindent ∞\par
\noindent X\par
\noindent = 1 + xt + Hn (x)tn\par
\noindent n=2\par
}
\end{sourceblock}


\begin{explainblock}
定理段不要只看结论，要看它把哪两个计数方式连在一起。组合学证明最常见的技巧是：左边按一种标准数，右边按另一种标准数，然后说明两边数的是同一批对象。

\smallskip
\textbf{初学者检查点：}
\begin{itemize}[leftmargin=2em,itemsep=0.25em]
\item 先复述定理的假设，不要把适用范围扩大。
\item 找出证明中的基本计数原则：乘法、加法、双射、递推、容斥或生成函数。
\item 检查最后的等式化简是否和前面的对象解释一致。
\end{itemize}

\smallskip
\textbf{本章读法提醒：} 第八章 Catalan 问题的核心常常是前缀不越界或递归分解。
\end{explainblock}


\begin{sourceblock}{第 8 章续读}
{\small\sloppy
\noindent ∞\par
\noindent !\par
\noindent X X\par
\noindent n−1\par
\noindent = 1 + xt + xHn−1 (x) + Hi (x)Hn−1−i (x) tn\par
\noindent n=2 i=1\par
\noindent X∞ ∞ n−1\par
\noindent X X\par
\noindent n−1\par
\noindent = 1 + xt + xt Hn−1 (x)t + Hi (x)Hn−1−i (x)tn\par
\noindent n=2 n=2 i=1\par
\noindent ∞\par
\noindent ! ∞ X\par
\noindent X X n\par
\noindent = 1 + xt 1 + Hn (x)tn +t Hi (x)Hn−i (x)tn\par
\noindent n=1 n=1 i=1\par
\noindent ∞ X\par
\noindent ∞ ∞\par
}
\end{sourceblock}


\begin{explainblock}
这一段是原书论述链的一部分。读的时候把它当作桥梁：它通常负责把上一个定义、例子或定理，引到下一个计数方法。

\smallskip
\textbf{初学者检查点：}
\begin{itemize}[leftmargin=2em,itemsep=0.25em]
\item 找出这一段承接的上一个概念。
\item 找出它准备引出的下一个概念。
\item 把中间的关键等式或构造单独抄出来。
\end{itemize}

\smallskip
\textbf{本章读法提醒：} 第八章 Catalan 问题的核心常常是前缀不越界或递归分解。
\end{explainblock}


\begin{sourceblock}{第 8 章续读}
{\small\sloppy
\noindent !\par
\noindent X X\par
\noindent = 1 + xtH(x, t) + t Hi (x)Hn−i (x)t − n\par
\noindent H0 (x)Hn (x)t n\par
\noindent n=0 i=0 n=0\par
\noindent = 1 + xtH(x, t) + tH(x, t) − tH(x, t),\par
\noindent 所以\par
\noindent tH(x, t)2 + (xt − t − 1)H(x, t) + 1 = 0.\par
\noindent 故 p\par
\noindent 1 + t − xt ± (xt − t − 1)2 − 4t\par
\noindent H(x, t) = .\par
\noindent 2t\par
\noindent 又 H(x, 0) = H0 (x) = 1, 所以\par
\noindent p\par
\noindent 1 + t − xt − (xt − t − 1)2 − 4t\par
\noindent H(x, t) = .\par
\noindent 2t\par
\noindent 下面我们给出 Narayana 数的计算公式.\par
}
\end{sourceblock}


\begin{explainblock}
这一段是原书论述链的一部分。读的时候把它当作桥梁：它通常负责把上一个定义、例子或定理，引到下一个计数方法。

\smallskip
\textbf{初学者检查点：}
\begin{itemize}[leftmargin=2em,itemsep=0.25em]
\item 找出这一段承接的上一个概念。
\item 找出它准备引出的下一个概念。
\item 把中间的关键等式或构造单独抄出来。
\end{itemize}

\smallskip
\textbf{本章读法提醒：} 第八章 Catalan 问题的核心常常是前缀不越界或递归分解。
\end{explainblock}


\begin{sourceblock}{第 8 章续读}
{\small\sloppy
\noindent 8.2 Narayana 数与 Narayana 多项式 159\par
}
\end{sourceblock}


\begin{explainblock}
这一段是原书论述链的一部分。读的时候把它当作桥梁：它通常负责把上一个定义、例子或定理，引到下一个计数方法。

\smallskip
\textbf{初学者检查点：}
\begin{itemize}[leftmargin=2em,itemsep=0.25em]
\item 找出这一段承接的上一个概念。
\item 找出它准备引出的下一个概念。
\item 把中间的关键等式或构造单独抄出来。
\end{itemize}

\smallskip
\textbf{本章读法提醒：} 第八章 Catalan 问题的核心常常是前缀不越界或递归分解。
\end{explainblock}


\begin{sourceblock}{定理 8.2.6 设 Nn,k (n ≥ k ≥ 1), 为 Narayana 数, 则}
{\small\sloppy
\noindent 定理 8.2.6 设 Nn,k (n ≥ k ≥ 1), 为 Narayana 数, 则\par
\noindent Nn,k = = . (8.6)\par
\noindent n k k−1 k k−1 k−1\par
\noindent 证明: 下面通过展开 H(x, t) 的表达式 (8.5) 来求解 Nn,k 或 Nn,k . 由 (8.5) 可\par
\noindent 得\par
\noindent p\par
\noindent 1 + t(1 − x) − (1 + t(1 − x))2 − 4t\par
\noindent H(x, t) =\par
\noindent 2t\par
\noindent q\par
\noindent 1 1 − 1 − (1+t(1−x))\par
\noindent 4t\par
\noindent = .\par
\noindent (1+t(1−x))2\par
\noindent 注意到 √ ∞\par
\noindent 1− 1 − 4z X 1 2k k\par
\noindent = z ,\par
\noindent 2z k+1 k\par
}
\end{sourceblock}


\begin{explainblock}
定理段不要只看结论，要看它把哪两个计数方式连在一起。组合学证明最常见的技巧是：左边按一种标准数，右边按另一种标准数，然后说明两边数的是同一批对象。

\smallskip
\textbf{初学者检查点：}
\begin{itemize}[leftmargin=2em,itemsep=0.25em]
\item 先复述定理的假设，不要把适用范围扩大。
\item 找出证明中的基本计数原则：乘法、加法、双射、递推、容斥或生成函数。
\item 检查最后的等式化简是否和前面的对象解释一致。
\end{itemize}

\smallskip
\textbf{本章读法提醒：} 第八章 Catalan 问题的核心常常是前缀不越界或递归分解。
\end{explainblock}


\begin{sourceblock}{第 8 章续读}
{\small\sloppy
\noindent k=0\par
\noindent 可得\par
\noindent X∞ k\par
\noindent H(x, t) =\par
\noindent 1 + t(1 − x) k+1 k (1 + t(1 − x))2\par
\noindent k=0\par
\noindent ∞\par
\noindent X\par
\noindent = .\par
\noindent k + 1 k (1 + t(1 − x))2k+1\par
\noindent k=0\par
\noindent 又因为\par
\noindent ∞\par
\noindent X\par
\noindent tk j 2k + j\par
\noindent = (−1) tj+k (1 − x)j (6a)\par
\noindent (1 + t(1 − x))2k+1 j\par
\noindent j=0\par
}
\end{sourceblock}


\begin{explainblock}
这一段是原书论述链的一部分。读的时候把它当作桥梁：它通常负责把上一个定义、例子或定理，引到下一个计数方法。

\smallskip
\textbf{初学者检查点：}
\begin{itemize}[leftmargin=2em,itemsep=0.25em]
\item 找出这一段承接的上一个概念。
\item 找出它准备引出的下一个概念。
\item 把中间的关键等式或构造单独抄出来。
\end{itemize}

\smallskip
\textbf{本章读法提醒：} 第八章 Catalan 问题的核心常常是前缀不越界或递归分解。
\end{explainblock}


\begin{sourceblock}{第 8 章续读}
{\small\sloppy
\noindent X∞ X j\par
\noindent i+j 2k + j j j+k i\par
\noindent = (−1) t x\par
\noindent 2k i\par
\noindent j=0 i=0\par
\noindent X∞ X ∞\par
\noindent i+j 2k + j j j+k i\par
\noindent = (−1) t x\par
\noindent 2k i\par
\noindent i=0 j=i\par
\noindent X∞ X ∞\par
\noindent j 2k + j + i j + i i+j+k i\par
\noindent = (−1) t x,\par
\noindent 2k i\par
\noindent i=0 j=0\par
\noindent 所以\par
\noindent ∞ X\par
\noindent X n\par
}
\end{sourceblock}


\begin{explainblock}
这一段是原书论述链的一部分。读的时候把它当作桥梁：它通常负责把上一个定义、例子或定理，引到下一个计数方法。

\smallskip
\textbf{初学者检查点：}
\begin{itemize}[leftmargin=2em,itemsep=0.25em]
\item 找出这一段承接的上一个概念。
\item 找出它准备引出的下一个概念。
\item 把中间的关键等式或构造单独抄出来。
\end{itemize}

\smallskip
\textbf{本章读法提醒：} 第八章 Catalan 问题的核心常常是前缀不越界或递归分解。
\end{explainblock}


\begin{sourceblock}{第 8 章续读}
{\small\sloppy
\noindent H(x, t) = 1 + Nn,k xk tn\par
\noindent n=1 k=1\par
\noindent ∞ X\par
\noindent X ∞\par
\noindent ∞ X\par
\noindent (−1)j 2k 2k + j + i j + i i+j+k i\par
\noindent = t x\par
\noindent k+1 k 2k i\par
\noindent k=0 i=0 j=0\par
\noindent 160 第 8 章 Catalan 数与 Narayana 数\par
\noindent ∞ X\par
\noindent X ∞\par
\noindent ∞ X\par
\noindent (−1)j+k 2k i+j+k i + j − k i+j i\par
\noindent = t x\par
\noindent k+1 k 2k i\par
\noindent k=0 i=0 j=k\par
\noindent ∞ X\par
}
\end{sourceblock}


\begin{explainblock}
这一段是原书论述链的一部分。读的时候把它当作桥梁：它通常负责把上一个定义、例子或定理，引到下一个计数方法。

\smallskip
\textbf{初学者检查点：}
\begin{itemize}[leftmargin=2em,itemsep=0.25em]
\item 找出这一段承接的上一个概念。
\item 找出它准备引出的下一个概念。
\item 把中间的关键等式或构造单独抄出来。
\end{itemize}

\smallskip
\textbf{本章读法提醒：} 第八章 Catalan 问题的核心常常是前缀不越界或递归分解。
\end{explainblock}


\begin{sourceblock}{第 8 章续读}
{\small\sloppy
\noindent X ∞ X\par
\noindent j\par
\noindent (−1)j+k 2k i+j+k i + j − k i+j i\par
\noindent = t x\par
\noindent k+1 k 2k i\par
\noindent i=0 j=0 k=0\par
\noindent X ∞ X\par
\noindent ∞ X j\par
\noindent (−1)j+k i + j i+j+k j i+j i\par
\noindent = t x\par
\noindent k+1 i k k\par
\noindent i=0 j=0 k=0\par
\noindent ∞ X\par
\noindent ∞ X\par
\noindent X j\par
\noindent (−1)j i + j −i − j − 1 j+1\par
\noindent = ti+j xi\par
\noindent j+1 i k k+1\par
}
\end{sourceblock}


\begin{explainblock}
这一段是原书论述链的一部分。读的时候把它当作桥梁：它通常负责把上一个定义、例子或定理，引到下一个计数方法。

\smallskip
\textbf{初学者检查点：}
\begin{itemize}[leftmargin=2em,itemsep=0.25em]
\item 找出这一段承接的上一个概念。
\item 找出它准备引出的下一个概念。
\item 把中间的关键等式或构造单独抄出来。
\end{itemize}

\smallskip
\textbf{本章读法提醒：} 第八章 Catalan 问题的核心常常是前缀不越界或递归分解。
\end{explainblock}


\begin{sourceblock}{第 8 章续读}
{\small\sloppy
\noindent i=0 j=0 k=0\par
\noindent X∞ X ∞ X j\par
\noindent (−1)j i + j −i − j − 1 j + 1 i+j i\par
\noindent = t x\par
\noindent j+1 i k j−k\par
\noindent i=0 j=0 k=0\par
\noindent X∞ X ∞\par
\noindent (−1)j i + j −i i+j i\par
\noindent = t x\par
\noindent j+1 i j\par
\noindent i=0 j=0\par
\noindent X∞ X ∞\par
\noindent 1 i+j i + j − 1 i+j i\par
\noindent = t x\par
\noindent j+1 i j\par
\noindent i=0 j=0\par
\noindent X∞ X ∞\par
\noindent = x t\par
}
\end{sourceblock}


\begin{explainblock}
这一段是原书论述链的一部分。读的时候把它当作桥梁：它通常负责把上一个定义、例子或定理，引到下一个计数方法。

\smallskip
\textbf{初学者检查点：}
\begin{itemize}[leftmargin=2em,itemsep=0.25em]
\item 找出这一段承接的上一个概念。
\item 找出它准备引出的下一个概念。
\item 把中间的关键等式或构造单独抄出来。
\end{itemize}

\smallskip
\textbf{本章读法提醒：} 第八章 Catalan 问题的核心常常是前缀不越界或递归分解。
\end{explainblock}


\begin{sourceblock}{第 8 章续读}
{\small\sloppy
\noindent n−k+1 k n−k\par
\noindent k=0 n=k\par
\noindent X∞ X n\par
\noindent = x t\par
\noindent n−k+1 k n−k\par
\noindent n=0 k=0\par
\noindent X∞ X n\par
\noindent =1+ x t\par
\noindent n−k+1 k k−1\par
\noindent n=1 k=1\par
\noindent X∞ X n\par
\noindent =1+ xk tn ,\par
\noindent n k k−1\par
\noindent n=1 k=1\par
\noindent 当 n ≥ k ≥ 1 时, 比较系数可得,\par
\noindent Nn,k = .\par
\noindent n k k−1\par
}
\end{sourceblock}


\begin{explainblock}
这一段是原书论述链的一部分。读的时候把它当作桥梁：它通常负责把上一个定义、例子或定理，引到下一个计数方法。

\smallskip
\textbf{初学者检查点：}
\begin{itemize}[leftmargin=2em,itemsep=0.25em]
\item 找出这一段承接的上一个概念。
\item 找出它准备引出的下一个概念。
\item 把中间的关键等式或构造单独抄出来。
\end{itemize}

\smallskip
\textbf{本章读法提醒：} 第八章 Catalan 问题的核心常常是前缀不越界或递归分解。
\end{explainblock}


\begin{sourceblock}{注记 8.2.7 这个定理的证明我们将在学习过圈引理后再给出其它证明, 详见第 8.3}
{\small\sloppy
\noindent 注记 8.2.7 这个定理的证明我们将在学习过圈引理后再给出其它证明, 详见第 8.3\par
\noindent 节.\par
\noindent 由定理 8.2.6 我们很容易发现 Nn,k 的对称性:\par
\noindent Nn,k = = = Nn,n−k+1 . (8.7)\par
\noindent n k k−1 n n−k n−k+1\par
}
\end{sourceblock}


\begin{explainblock}
注记通常是在补充术语、记号、历史或一个容易忽略的约定。它未必给新公式，但常常决定你后面会不会把符号读错。

\smallskip
\textbf{初学者检查点：}
\begin{itemize}[leftmargin=2em,itemsep=0.25em]
\item 记下新符号的读法。
\item 确认这个约定是否只在本章使用。
\item 把注记里的限制条件补到前面的定义旁边。
\end{itemize}

\smallskip
\textbf{本章读法提醒：} 第八章 Catalan 问题的核心常常是前缀不越界或递归分解。
\end{explainblock}


\begin{sourceblock}{§ 8.3 圈引理}
{\small\sloppy
\noindent § 8.3. 圈引理\par
\noindent 圈引理是组合计数上一个很有用的方法. 首先我们给出由两个字母组成的字\par
\noindent k-dominating 字的定义.\par
}
\end{sourceblock}


\begin{explainblock}
这一节先不要急着背公式。先问三个问题：对象是什么，哪些限制条件不能丢，最后到底是在数所有对象、还是在数某个统计量的分布。把这三件事说清楚，后面的公式才会有位置。

\smallskip
\textbf{初学者检查点：}
\begin{itemize}[leftmargin=2em,itemsep=0.25em]
\item 把本节出现的新对象写成一句话定义。
\item 把每个公式左边和右边分别翻译成“在数什么”。
\item 找一个最小非平凡例子，手工列出所有对象。
\end{itemize}

\smallskip
\textbf{本章读法提醒：} 第八章 Catalan 问题的核心常常是前缀不越界或递归分解。
\end{explainblock}


\begin{sourceblock}{定义 8.3.1 设由字母 a, b 组成的字 W = w1 w2 · · · wn , 若对任意 i(1 ≤ i ≤ n), 字}
{\small\sloppy
\noindent 定义 8.3.1 设由字母 a, b 组成的字 W = w1 w2 · · · wn , 若对任意 i(1 ≤ i ≤ n), 字\par
\noindent W 的因子 w1 w2 · · · wi 中字母 a 个数总是大于字母 b 个数的 k 倍, 则称字 W 是\par
\noindent k-dominating 的.\par
\noindent 对任意 n 长字 W , 我们用 (W ) 表示由这个字形成的 n 长圈形字. 显然, 对任意\par
\noindent 一个 n 长的圈形字从任意两个字母的位置剪开就得到了一个 n 长字.\par
\noindent 有了上面的概念, 圈引理可叙述如下.\par
}
\end{sourceblock}


\begin{explainblock}
定义段的任务是定边界。这里要特别注意参数范围、是否允许重复、是否考虑顺序、对象有没有标签。后面的每一道例题和证明，本质上都在反复使用这些边界。

\smallskip
\textbf{初学者检查点：}
\begin{itemize}[leftmargin=2em,itemsep=0.25em]
\item 圈出被定义的对象名称。
\item 圈出参数范围，例如 n、k 是否要求正整数。
\item 判断它是有序对象、无序对象、带标签对象还是结构对象。
\end{itemize}

\smallskip
\textbf{本章读法提醒：} 第八章 Catalan 问题的核心常常是前缀不越界或递归分解。
\end{explainblock}


\begin{sourceblock}{定理 8.3.2 (Cycle Lemma) 设 (W ) = (w1 w2 · · · wm+n ) 是由 m 个字母 a 与 n 个}
{\small\sloppy
\noindent 定理 8.3.2 (Cycle Lemma) 设 (W ) = (w1 w2 · · · wm+n ) 是由 m 个字母 a 与 n 个\par
\noindent 字母 b 组成的圈形字, 设 k 为正整数, 若 m > kn, 则从 (W ) 任意位置剪开得到的\par
\noindent m + n 个字中 (可能有相同的), 恰有 m − kn 个字是 k-dominating 的.\par
\noindent 证明: 我们首先叙述下面的事实: 在 (W ) 选定初始位置, 并按顺时针方向读,\par
\noindent 若出现 k 个字母 a 后出现字母 b, 即\par
\noindent · · · aa · · · a\} b · · · ,\par
\noindent | \{z\par
\noindent k个a\par
\noindent 则从这 k + 1 个字母形成的 n 个位置处剪开得到的字显然不是 k-dominating 的, 所\par
\noindent 以删去这 k + 1 个字母并不影响在 (W ) 任意位置剪开得到的 m + n 个字中 (可能有\par
\noindent 相同的) k-dominating 的字的个数.\par
\noindent 下面我们说明 (W ) 中存在这样的字母列. 因为对于由 m 个字母 a 和 n 个字母\par
\noindent b 组成的任意圈形字, 都可以看作先将 n 个字母 b 放到圆上, 形成 n 个空隙, 再将 m\par
\noindent 个字母分成 n 份 (允许为空) 放入空隙中得到, 因为 m > kn, 所以由鸽巢原理可知\par
\noindent 存在两个字母 b 中间字母 a 的个数大于 k, 故而按顺时针方向存在字母 b 前有 k 个\par
\noindent 字母 a 后的字母列. 所以我们可以将其从 (W ) 中删去得到 (W ′ ), 这不改变所求结\par
\noindent 果. 而此时 (W ′ ) 显然也存在这样的序列, 所以我们可以一直作下去, 直到 n 次后得\par
\noindent 到的圈形字中不再含有字母 b, 换言之, 最后得到的圈形字中只有 m − kn 个字母 a,\par
}
\end{sourceblock}


\begin{explainblock}
定理段不要只看结论，要看它把哪两个计数方式连在一起。组合学证明最常见的技巧是：左边按一种标准数，右边按另一种标准数，然后说明两边数的是同一批对象。

\smallskip
\textbf{初学者检查点：}
\begin{itemize}[leftmargin=2em,itemsep=0.25em]
\item 先复述定理的假设，不要把适用范围扩大。
\item 找出证明中的基本计数原则：乘法、加法、双射、递推、容斥或生成函数。
\item 检查最后的等式化简是否和前面的对象解释一致。
\end{itemize}

\smallskip
\textbf{本章读法提醒：} 第八章 Catalan 问题的核心常常是前缀不越界或递归分解。
\end{explainblock}


\begin{sourceblock}{第 8 章续读}
{\small\sloppy
\noindent 从这 m − kn 个位置剪开得到的字显然都是 k-dominating 的字. 定理得证.\par
\noindent 我们现在考虑由字母 a 和 b 组成的字. 设由 n + 1 个字母 a 和 n 个字母 b 组成\par
\noindent 的任意字的集合为 Sw , 则显然有 |Sw | = 2n+1\par
\noindent n .\par
\noindent 设 W = w1 w2 · · · w2n+1 ∈ Sw . 用 (W ) 表示由 W 生成的圈形字, 则我们有以下\par
\noindent 结论.\par
}
\end{sourceblock}


\begin{explainblock}
这一段是原书论述链的一部分。读的时候把它当作桥梁：它通常负责把上一个定义、例子或定理，引到下一个计数方法。

\smallskip
\textbf{初学者检查点：}
\begin{itemize}[leftmargin=2em,itemsep=0.25em]
\item 找出这一段承接的上一个概念。
\item 找出它准备引出的下一个概念。
\item 把中间的关键等式或构造单独抄出来。
\end{itemize}

\smallskip
\textbf{本章读法提醒：} 第八章 Catalan 问题的核心常常是前缀不越界或递归分解。
\end{explainblock}


\begin{sourceblock}{定理 8.3.3 集合 Sw 中的任意字都是本原字.}
{\small\sloppy
\noindent 定理 8.3.3 集合 Sw 中的任意字都是本原字.\par
}
\end{sourceblock}


\begin{explainblock}
定理段不要只看结论，要看它把哪两个计数方式连在一起。组合学证明最常见的技巧是：左边按一种标准数，右边按另一种标准数，然后说明两边数的是同一批对象。

\smallskip
\textbf{初学者检查点：}
\begin{itemize}[leftmargin=2em,itemsep=0.25em]
\item 先复述定理的假设，不要把适用范围扩大。
\item 找出证明中的基本计数原则：乘法、加法、双射、递推、容斥或生成函数。
\item 检查最后的等式化简是否和前面的对象解释一致。
\end{itemize}

\smallskip
\textbf{本章读法提醒：} 第八章 Catalan 问题的核心常常是前缀不越界或递归分解。
\end{explainblock}


\begin{sourceblock}{注记 8.3.4 w 是本原字是指不存在自然数 n(≥ 2) 及字 v 使得 w = v n .}
{\small\sloppy
\noindent 注记 8.3.4 w 是本原字是指不存在自然数 n(≥ 2) 及字 v 使得 w = v n .\par
\noindent 证明: 假设不然, 则存在 W = w1 w2 · · · w2n+1 ∈ Sw , 使得对某一个自然数\par
\noindent k(k ≥ 2) 及 j(≥ 1) 长字 v, 满足 W = v k , 所以 v = w1 w2 · · · wj 且 jk = 2n + 1, 但\par
\noindent 这是不可能的, 因为若设 v 中有 x 个字母 a 和 y 个字母 b, 则有:\par
\noindent (\par
\noindent kx = n + 1,\par
\noindent ky = n.\par
\noindent 即可推得 k 是 n + 1 与 n 的公因子, 而由 n + 1 与 n 互素推出 k = 1, 这与假设\par
\noindent k ≥ 2 矛盾. 定理得证.\par
\noindent 162 第 8 章 Catalan 数与 Narayana 数\par
}
\end{sourceblock}


\begin{explainblock}
注记通常是在补充术语、记号、历史或一个容易忽略的约定。它未必给新公式，但常常决定你后面会不会把符号读错。

\smallskip
\textbf{初学者检查点：}
\begin{itemize}[leftmargin=2em,itemsep=0.25em]
\item 记下新符号的读法。
\item 确认这个约定是否只在本章使用。
\item 把注记里的限制条件补到前面的定义旁边。
\end{itemize}

\smallskip
\textbf{本章读法提醒：} 第八章 Catalan 问题的核心常常是前缀不越界或递归分解。
\end{explainblock}


\begin{sourceblock}{定理 8.3.5 由集合 Sw 中的字生成的所有圈形字记为集合 Scw , 则 |Scw | = Cn , 其}
{\small\sloppy
\noindent 定理 8.3.5 由集合 Sw 中的字生成的所有圈形字记为集合 Scw , 则 |Scw | = Cn , 其\par
\noindent 中 Cn 为 Catalan 数.\par
\noindent 证明: 由定理 8.3.3 知, 下面 2n + 1 个字各不相同\par
\noindent w1 w2 · · · w2n+1 , w2 w3 · · · w2n+1 w1 , ··· , w2n+1 w1 · · · w2n .\par
\noindent 所以每一个圈形字从不同位置剪开可以得到 2n + 1 个不同的字, 故而:\par
\noindent |Sw | 1 2n + 1 1 2n\par
\noindent |Scw | = = = = Cn .\par
\noindent 2n + 1 2n + 1 n n+1 n\par
\noindent 下面我们应用圈引理给出定理 8.2.6(叙述如下) 的证明.\par
}
\end{sourceblock}


\begin{explainblock}
定理段不要只看结论，要看它把哪两个计数方式连在一起。组合学证明最常见的技巧是：左边按一种标准数，右边按另一种标准数，然后说明两边数的是同一批对象。

\smallskip
\textbf{初学者检查点：}
\begin{itemize}[leftmargin=2em,itemsep=0.25em]
\item 先复述定理的假设，不要把适用范围扩大。
\item 找出证明中的基本计数原则：乘法、加法、双射、递推、容斥或生成函数。
\item 检查最后的等式化简是否和前面的对象解释一致。
\end{itemize}

\smallskip
\textbf{本章读法提醒：} 第八章 Catalan 问题的核心常常是前缀不越界或递归分解。
\end{explainblock}


\begin{sourceblock}{定理 8.3.6 设 Nn,k 为 Narayana 数, 则}
{\small\sloppy
\noindent 定理 8.3.6 设 Nn,k 为 Narayana 数, 则\par
\noindent Nn,k = . (8.8)\par
\noindent k+1 k k\par
\noindent 证明: 由 Dyck 路峰与谷的定义可知峰的个数等于 Dyck 字中出现 ud 字段的\par
\noindent 次数, 谷的个数等于 Dyck 字出现 du 字段的次数. 比如, ududud 出现 ud 三次对\par
\noindent 应该 Dyck 路的三个峰, 而 du 字段则出现两次, 对应 Dyck 路中的两个谷. 为证明\par
\noindent Nn,k 的计算公式, 我们考虑由 n + 1 个字母 u 和 n 个字母 d 组成的一类字. 这类字\par
\noindent 满足下列两个条件:\par
\noindent 1. 字的首字母为 u, 尾字母为 d;\par
\noindent 2. 字中仅出现 k(0 ≤ k ≤ n − 1) 个 du 的字母段.\par
\noindent 对于给定 n 与 k, 满足上述两个条件的字可以通过下面三步得到:\par
\noindent 1) 将 n + 1 个字母 u 排成一排, 从它们产生的 n 个空隙中找到 k 个位置, 共有\par
\noindent n\par
\noindent k 中方法;\par
\noindent 2) 将 n 个字母 d 排成一排, 从它们产生的 n − 1 个空隙中找到 k 个位置, 将其分\par
\noindent 成不空的 k + 1 部分, 共有 n−1\par
\noindent k ;\par
\noindent 3) 将 2) 中分成的 k + 1 部分依次放到 1) 中选取的 k 个位置和最后一个字母 u\par
}
\end{sourceblock}


\begin{explainblock}
定理段不要只看结论，要看它把哪两个计数方式连在一起。组合学证明最常见的技巧是：左边按一种标准数，右边按另一种标准数，然后说明两边数的是同一批对象。

\smallskip
\textbf{初学者检查点：}
\begin{itemize}[leftmargin=2em,itemsep=0.25em]
\item 先复述定理的假设，不要把适用范围扩大。
\item 找出证明中的基本计数原则：乘法、加法、双射、递推、容斥或生成函数。
\item 检查最后的等式化简是否和前面的对象解释一致。
\end{itemize}

\smallskip
\textbf{本章读法提醒：} 第八章 Catalan 问题的核心常常是前缀不越界或递归分解。
\end{explainblock}


\begin{sourceblock}{第 8 章续读}
{\small\sloppy
\noindent 的后面.\par
\noindent 不妨将这些字组成的集合为 Πn,k , 则\par
\noindent n−1 n\par
\noindent |Πn,k | = .\par
\noindent k k\par
\noindent 3−1\par
}
\end{sourceblock}


\begin{explainblock}
这一段是原书论述链的一部分。读的时候把它当作桥梁：它通常负责把上一个定义、例子或定理，引到下一个计数方法。

\smallskip
\textbf{初学者检查点：}
\begin{itemize}[leftmargin=2em,itemsep=0.25em]
\item 找出这一段承接的上一个概念。
\item 找出它准备引出的下一个概念。
\item 把中间的关键等式或构造单独抄出来。
\end{itemize}

\smallskip
\textbf{本章读法提醒：} 第八章 Catalan 问题的核心常常是前缀不越界或递归分解。
\end{explainblock}


\begin{sourceblock}{例如: 取 n = 3, 当 k = 0 时, 满足条件的字有 0 0 = 1 个:}
{\small\sloppy
\noindent 例如: 取 n = 3, 当 k = 0 时, 满足条件的字有 0 0 = 1 个:\par
\noindent uuuuddd;\par
\noindent 3−1\par
\noindent 当 k = 1 时, 满足条件的字有 1 1 = 6 个:\par
\noindent uduuudd, udduuud, uuduudd,\par
\noindent uudduud, uuududd, uuuddud;\par
\noindent 当 k = 2 时, 满足条件的字有 32 3−1\par
\noindent ududuud, uduudud, uududud.\par
\noindent 因为这些字以 u 为首且 du 的个数是 k, 所以去掉这些字的首字母不影响 du 片\par
\noindent 段的个数. 不妨设 W ∈ Πn,k . 若去掉 W 的首字母 u, 得到的字是具有 k 个 du 片段\par
\noindent 的 Dyck 字, 因此 W 对应一条具有 k 个谷的 Dyck 路; 反之, 在任意一个 Dyck 字\par
\noindent 前面加上字母 u 得到的字显然都属于 Πn,k , 且因为对任意一个 Dyck 字从左到右读\par
\noindent u 的个数始终不小于 d 的个数, 所以 Dyck 字加上首字母 u 时得到的字肯定满足从\par
\noindent 左到右读 u 的个数始终大于 d 的个数. 因此 Nn,k 也表示集合 Πn,k 中满足: 字母 u\par
\noindent 的个数从左到右总比字母 d 个数大的字的个数. 因为 Πn,k 中有 n 个 d 和 n + 1 个\par
\noindent u, 所以也可说 Nn,k+1 表示集合 Πn,k 中 1-dominating 字的个数. 下面我们证明集合\par
\noindent Πn,k 中 1-dominating 字的个数 k+1 k k ,\par
\noindent 由条件 1 和条件 2 知 w ∈ Πn,k 当且仅当 w = α1 α2 · · · αk+1 , 其中 αi 均为\par
}
\end{sourceblock}


\begin{explainblock}
例题是学习这本书最重要的部分。不要只看答案，要把每个限制条件变成一个计数动作：先安排谁、再插入谁、有没有重复、需不需要除以对称性。

\smallskip
\textbf{初学者检查点：}
\begin{itemize}[leftmargin=2em,itemsep=0.25em]
\item 先写出没有限制时的总数。
\item 逐条加入限制条件，观察总数怎样变化。
\item 最后检查有没有把同一种方案重复算了多次。
\end{itemize}

\smallskip
\textbf{本章读法提醒：} 第八章 Catalan 问题的核心常常是前缀不越界或递归分解。
\end{explainblock}


\begin{sourceblock}{第 8 章续读}
{\small\sloppy
\noindent u · · · ud · · · d 的形式.\par
}
\end{sourceblock}


\begin{explainblock}
这一段是原书论述链的一部分。读的时候把它当作桥梁：它通常负责把上一个定义、例子或定理，引到下一个计数方法。

\smallskip
\textbf{初学者检查点：}
\begin{itemize}[leftmargin=2em,itemsep=0.25em]
\item 找出这一段承接的上一个概念。
\item 找出它准备引出的下一个概念。
\item 把中间的关键等式或构造单独抄出来。
\end{itemize}

\smallskip
\textbf{本章读法提醒：} 第八章 Catalan 问题的核心常常是前缀不越界或递归分解。
\end{explainblock}


\begin{sourceblock}{例如 n = 3, k = 1 时,}
{\small\sloppy
\noindent 例如 n = 3, k = 1 时,\par
\noindent ud|uuudd, udd|uuud, uud|uudd,\par
\noindent uudd|uud, uuud|udd, uuudd|ud;\par
\noindent n = 3, k = 2 时,\par
\noindent ud|ud|uud, ud|uud|ud, uud|ud|ud.\par
\noindent 对于任意的 w ∈ Πn,k , 不妨设为 w = α1 α2 · · · αk+1 . 对任意 i(0 ≤ i ≤ k), 记\par
\noindent i+1 · · · αk+1 α1 · · · αi (规定 α0 = αk+1 ), 则 w\par
\noindent w(i) = α (i) ∈ Π\par
\noindent n,k .\par
\noindent 设 (w) 表示由 w 生成的圈型字, 则对于 0 ≤ i ≤ k, w(i) 都可以看作是将圈型字\par
\noindent (w) 从 αi 与 αi+1 相邻的位置线性展开得到的, 这些位置显然都在片段 du 中间, 反\par
\noindent 之, 若不是在片段 du 之间展开, 则得到的字显然不是 Πn,k 中的字. 还需说明的是将\par
\noindent 圈形字 (w) 从不同位置展开得到的字是不同的. 这是因为 w 中含有 n + 1 个字母 u\par
\noindent 和 n 个字母 d, 由定理 8.3.3 知这些字都是本原字, 所以由定理 8.3.5 知这些字从形\par
\noindent 成的圈形字从不同位置剪开得到的字是不同的. 故而若对 Πn,k 中的元素按它们是否\par
\noindent 该能表示成同一个圈型字来进行分类, 则每一类恰好有 k + 1 个不同字. 因此, Πn,k\par
\noindent 可以分成 k+1 k k 类.\par
}
\end{sourceblock}


\begin{explainblock}
例题是学习这本书最重要的部分。不要只看答案，要把每个限制条件变成一个计数动作：先安排谁、再插入谁、有没有重复、需不需要除以对称性。

\smallskip
\textbf{初学者检查点：}
\begin{itemize}[leftmargin=2em,itemsep=0.25em]
\item 先写出没有限制时的总数。
\item 逐条加入限制条件，观察总数怎样变化。
\item 最后检查有没有把同一种方案重复算了多次。
\end{itemize}

\smallskip
\textbf{本章读法提醒：} 第八章 Catalan 问题的核心常常是前缀不越界或递归分解。
\end{explainblock}


\begin{sourceblock}{例如 n = 3, k = 1 时, 分成的三类. 第 1 类:}
{\small\sloppy
\noindent 例如 n = 3, k = 1 时, 分成的三类. 第 1 类:\par
\noindent ud|uuudd, uuudd|ud;\par
\noindent 164 第 8 章 Catalan 数与 Narayana 数\par
\noindent 这一类中只有 uuudd|ud 满足条件; 第 2 类:\par
\noindent udd|uuud, uuud|udd;\par
\noindent 只有 uuud|udd 满足条件; 第 3 类:\par
\noindent uud|uudd, uudd|uud.\par
\noindent 只有 uud|uudd 满足条件.\par
\noindent n = 3, k = 2 时, 只有一类\par
\noindent ud|ud|uud, ud|uud|ud, uud|ud|ud,\par
\noindent 其中 uud|ud|ud 满足条件.\par
\noindent 换言之, Πn,k 中只能形成 k+1 k k 个圈形字, 应用圈引理 8.3.2, 取 m =\par
\noindent n+1, k = 1, 知每一个圈形字从不同位置展开的所有字中只有 1 个字是 1-dominating\par
\noindent 的, 而这个 1-dominating 字显然从 u 开始, 到 d 结束, 所以这个字是在片段 du 中间\par
\noindent 位置展开的, 在上面我们已经说明, 在这些位置展开的字肯定是 Πn,k 中的字, 所以\par
\noindent Πn,k 满足 1-dominating 的字的个数为 k+1 k k . 故\par
\noindent Nn,k+1 = .\par
\noindent k+1 k k\par
}
\end{sourceblock}


\begin{explainblock}
例题是学习这本书最重要的部分。不要只看答案，要把每个限制条件变成一个计数动作：先安排谁、再插入谁、有没有重复、需不需要除以对称性。

\smallskip
\textbf{初学者检查点：}
\begin{itemize}[leftmargin=2em,itemsep=0.25em]
\item 先写出没有限制时的总数。
\item 逐条加入限制条件，观察总数怎样变化。
\item 最后检查有没有把同一种方案重复算了多次。
\end{itemize}

\smallskip
\textbf{本章读法提醒：} 第八章 Catalan 问题的核心常常是前缀不越界或递归分解。
\end{explainblock}


\begin{sourceblock}{§ 8.4 拉格朗日反演公式}
{\small\sloppy
\noindent § 8.4. 拉格朗日反演公式\par
\noindent 组合数学中, 尤其在计数组合学中我们经常会遇到形如\par
\noindent f (x) = xG(f (x)) (8.9)\par
\noindent 的 函 数 方 程. 其 中 f (x) ∈ xC⟦x⟧, G(x) ∈ C⟦x⟧. 比 如 Catalan 数 的 生 成 函 数\par
\noindent P\par
\noindent c(x) = n≥0 Cn xn 满足\par
\noindent c(x) = 1 + xc(x)2 .\par
\noindent 若令\par
\noindent f (x) = c(x) − 1, G(x) = (x + 1)2 ,\par
\noindent 则 f (x) 与 G(x) 满足函数方程 (8.9).\par
\noindent 在前面我们已经看到在利用生成函数求 Catalan 数时计算是比较复杂的. 事实\par
\noindent 上, 通常情形下, 对于任意正整数 n, f (x) 中 xn 的系数是很难计算的. 但是 G(x) 中\par
\noindent 各项系数通常比较容易确定, 这节我们所要介绍的拉格朗日反演公式正是利用这一\par
\noindent 特性, 通过 G(x) 的系数公式求解 f (x) 中系数的表达式. 这在计数组合学中有着重\par
\noindent 要的理论价值.\par
\noindent 在叙述这个重要的定理之前, 我们首先给出形式幂级数复合的一个基本却十分重\par
\noindent 要的结果.\par
\noindent P∞\par
}
\end{sourceblock}


\begin{explainblock}
这一节先不要急着背公式。先问三个问题：对象是什么，哪些限制条件不能丢，最后到底是在数所有对象、还是在数某个统计量的分布。把这三件事说清楚，后面的公式才会有位置。

\smallskip
\textbf{初学者检查点：}
\begin{itemize}[leftmargin=2em,itemsep=0.25em]
\item 把本节出现的新对象写成一句话定义。
\item 把每个公式左边和右边分别翻译成“在数什么”。
\item 找一个最小非平凡例子，手工列出所有对象。
\end{itemize}

\smallskip
\textbf{本章读法提醒：} 第八章 Catalan 问题的核心常常是前缀不越界或递归分解。
\end{explainblock}


\begin{sourceblock}{定理 8.4.1 设 f (x) = n=0 an x}
{\small\sloppy
\noindent 定理 8.4.1 设 f (x) = n=0 an x\par
\noindent n 是任意给定的形式幂级数, 则存在形式幂级数\par
\noindent P\par
\noindent g(x) = ∞ n\par
\noindent n=0 bn x 使得\par
\noindent g(0) = b0 = 0, f (g(x)) = x, (8.10)\par
\noindent 的充要条件是\par
\noindent f (0) = a0 = 0, f ′ (0) = a1 6= 0. (8.11)\par
\noindent 若 g 存在, 那么 g 是唯一的, 而且 g(f (x)) = x.\par
\noindent 证明: 我们先证明必要性. 若条件 (8.10) 成立, 则\par
\noindent f (g(x)) = a0 + a1 (b1 x + b2 x2 + · · · ) + a2 (b1 x + b2 x2 + · · · )2 + · · · = x,\par
\noindent 比较两边系数可知\par
\noindent a0 = 0, a1 b1 = 1.\par
\noindent 因此条件 (8.11) 是必要的.\par
\noindent P∞ n\par
\noindent 下面我们证明充分性. 若条件 (8.11) 成立, 假设存在这样的 g(x) = n=1 bn x\par
\noindent 使得\par
\noindent f (g(x)) = a1 (b1 x + b2 x2 + · · · ) + a2 (b1 x + b2 x2 + · · · )2 + · · · = 0,\par
}
\end{sourceblock}


\begin{explainblock}
定理段不要只看结论，要看它把哪两个计数方式连在一起。组合学证明最常见的技巧是：左边按一种标准数，右边按另一种标准数，然后说明两边数的是同一批对象。

\smallskip
\textbf{初学者检查点：}
\begin{itemize}[leftmargin=2em,itemsep=0.25em]
\item 先复述定理的假设，不要把适用范围扩大。
\item 找出证明中的基本计数原则：乘法、加法、双射、递推、容斥或生成函数。
\item 检查最后的等式化简是否和前面的对象解释一致。
\end{itemize}

\smallskip
\textbf{本章读法提醒：} 第八章 Catalan 问题的核心常常是前缀不越界或递归分解。
\end{explainblock}


\begin{sourceblock}{第 8 章续读}
{\small\sloppy
\noindent 即得到了下面的方程组:\par
\noindent \par
\noindent \par
\noindent  a1 b1 = 1,\par
\noindent \par
\noindent \par
\noindent \par
\noindent \par
\noindent \par
\noindent \par
\noindent  a1 b2 + a2 b21 = 0,\par
\noindent \par
\noindent \par
\noindent \par
\noindent a1 b3 + 2a2 b1 b2 + a3 b31 = 0,\par
\noindent \par
\noindent \par
\noindent \par
}
\end{sourceblock}


\begin{explainblock}
这一段是原书论述链的一部分。读的时候把它当作桥梁：它通常负责把上一个定义、例子或定理，引到下一个计数方法。

\smallskip
\textbf{初学者检查点：}
\begin{itemize}[leftmargin=2em,itemsep=0.25em]
\item 找出这一段承接的上一个概念。
\item 找出它准备引出的下一个概念。
\item 把中间的关键等式或构造单独抄出来。
\end{itemize}

\smallskip
\textbf{本章读法提醒：} 第八章 Catalan 问题的核心常常是前缀不越界或递归分解。
\end{explainblock}


\begin{sourceblock}{第 8 章续读}
{\small\sloppy
\noindent \par
\noindent \par
\noindent ··· ···\par
\noindent 显然当 a1 6= 0 时, b1 , b2 , b2 , . . . , 存在且唯一, 因此, g(x) 存在且唯一.\par
\noindent 上面求得的级数 g(x) 显然满足 g(0) = 0, g ′ (0) = b1 6= 0, 因此对 g(x) 应用上面\par
\noindent 刚对 f (x) 证明的结果可知存在一个形式幂级数 h(x), 满足\par
\noindent h(0) = 0, g(h(x)) = x,\par
\noindent 所以\par
\noindent h(x) = x ◦ h = (f ◦ g) ◦ h = f ◦ (g ◦ h) = f ◦ x = f (x),\par
\noindent 所以 h(x) 就是 f (x), 因此 g ◦ f = x.\par
\noindent 从定理 8.4.1 显然能看出若 f ◦ g = x, 则 f (x) 与 g(x) 都不含常数项. 因此, 不妨\par
\noindent 考虑复数域 C 上所有不含有常数项的形式幂级数, 并记为 xC⟦x⟧. 显然形式幂级数\par
\noindent 上的复合运算在集合上是 xC⟦x⟧ 是封闭的, 且满足结合律, 显然对于任意形式幂级\par
\noindent 166 第 8 章 Catalan 数与 Narayana 数\par
\noindent 数 f ∈ xC⟦x⟧ 都有 f ◦ x = f = x ◦ f , 所以 xC⟦x⟧ 与复合运算 ◦ 构成幺半群, 其中 x\par
\noindent 为该幺半群的单位元. 从这个角度来说, 若 f, g ∈ xC⟦x⟧, 且 f ◦ g = x, 则称 f 与 g\par
\noindent 互为复合逆 (Compositional inverse), 并记 g = f <−1> . 因此, 定理 8.4.1 可以重新叙\par
\noindent 述如下:\par
}
\end{sourceblock}


\begin{explainblock}
这一段是原书论述链的一部分。读的时候把它当作桥梁：它通常负责把上一个定义、例子或定理，引到下一个计数方法。

\smallskip
\textbf{初学者检查点：}
\begin{itemize}[leftmargin=2em,itemsep=0.25em]
\item 找出这一段承接的上一个概念。
\item 找出它准备引出的下一个概念。
\item 把中间的关键等式或构造单独抄出来。
\end{itemize}

\smallskip
\textbf{本章读法提醒：} 第八章 Catalan 问题的核心常常是前缀不越界或递归分解。
\end{explainblock}


\begin{sourceblock}{定理 8.4.2 设 f, g ∈ xC⟦x⟧, 则 f 具有复合逆当且仅当 f (0) = 0 与 f ′ (0) 6= 0 同时}
{\small\sloppy
\noindent 定理 8.4.2 设 f, g ∈ xC⟦x⟧, 则 f 具有复合逆当且仅当 f (0) = 0 与 f ′ (0) 6= 0 同时\par
\noindent 成立, 且若存在复合逆 f <−1> 则唯一. 换言之, 若存在 g 使得 f ◦ g = x 或 g ◦ f = x,\par
\noindent 则 g = f <−1> .\par
}
\end{sourceblock}


\begin{explainblock}
定理段不要只看结论，要看它把哪两个计数方式连在一起。组合学证明最常见的技巧是：左边按一种标准数，右边按另一种标准数，然后说明两边数的是同一批对象。

\smallskip
\textbf{初学者检查点：}
\begin{itemize}[leftmargin=2em,itemsep=0.25em]
\item 先复述定理的假设，不要把适用范围扩大。
\item 找出证明中的基本计数原则：乘法、加法、双射、递推、容斥或生成函数。
\item 检查最后的等式化简是否和前面的对象解释一致。
\end{itemize}

\smallskip
\textbf{本章读法提醒：} 第八章 Catalan 问题的核心常常是前缀不越界或递归分解。
\end{explainblock}


\begin{sourceblock}{例 8.4.1 求 ex − 1 的复合逆.}
{\small\sloppy
\noindent 例 8.4.1 求 ex − 1 的复合逆.\par
\noindent 解: 因为\par
\noindent ∞\par
\noindent X xn\par
\noindent ex − 1 = ∈ xC⟦x⟧,\par
\noindent n!\par
\noindent n=1\par
\noindent 且一次项系数不为 0, 所以存在复合逆. 又 elog(1+x) − 1 = x, 所以 (ex − 1)<−1> =\par
\noindent log(1 + x).\par
\noindent 下面我们给出这节重要的定理: 拉格朗日反演公式.\par
}
\end{sourceblock}


\begin{explainblock}
例题是学习这本书最重要的部分。不要只看答案，要把每个限制条件变成一个计数动作：先安排谁、再插入谁、有没有重复、需不需要除以对称性。

\smallskip
\textbf{初学者检查点：}
\begin{itemize}[leftmargin=2em,itemsep=0.25em]
\item 先写出没有限制时的总数。
\item 逐条加入限制条件，观察总数怎样变化。
\item 最后检查有没有把同一种方案重复算了多次。
\end{itemize}

\smallskip
\textbf{本章读法提醒：} 第八章 Catalan 问题的核心常常是前缀不越界或递归分解。
\end{explainblock}


\begin{sourceblock}{定理 8.4.3 (The Lagrange inversion formula) 设 F (x) = a1 x + a2 x2 + · · · ∈ xC⟦x⟧,}
{\small\sloppy
\noindent 定理 8.4.3 (The Lagrange inversion formula) 设 F (x) = a1 x + a2 x2 + · · · ∈ xC⟦x⟧,\par
\noindent 其中 a1 6= 0, 则\par
\noindent n\par
\noindent x\par
\noindent n\par
\noindent n[x ]F <−1> k\par
\noindent (x) = k[x n−k\par
\noindent ] = k[x−k ]F (x)−n , (8.12)\par
\noindent F (x)\par
\noindent 其中 k, n 均为整数.\par
\noindent 证明: 由定理 8.4.2 知若 F <−1> (x) 存在则必然 F <−1> (x) ∈ xC⟦x⟧, 所以不妨\par
\noindent 设\par
\noindent X\par
\noindent F <−1> (x)k = pi xi . (8.13)\par
\noindent i≥k\par
\noindent 则\par
\noindent n[xn ]F <−1> (x)k = npn . (8.14)\par
\noindent 将式 (8.13) 中 x 换成 F (x) 立得:\par
}
\end{sourceblock}


\begin{explainblock}
定理段不要只看结论，要看它把哪两个计数方式连在一起。组合学证明最常见的技巧是：左边按一种标准数，右边按另一种标准数，然后说明两边数的是同一批对象。

\smallskip
\textbf{初学者检查点：}
\begin{itemize}[leftmargin=2em,itemsep=0.25em]
\item 先复述定理的假设，不要把适用范围扩大。
\item 找出证明中的基本计数原则：乘法、加法、双射、递推、容斥或生成函数。
\item 检查最后的等式化简是否和前面的对象解释一致。
\end{itemize}

\smallskip
\textbf{本章读法提醒：} 第八章 Catalan 问题的核心常常是前缀不越界或递归分解。
\end{explainblock}


\begin{sourceblock}{第 8 章续读}
{\small\sloppy
\noindent X\par
\noindent xk = pi F (x)i .\par
\noindent i≥k\par
\noindent 对两边求导可得:\par
\noindent X\par
\noindent kxk−1 = ipi F (x)i−1 F ′ (x),\par
\noindent i≥k\par
\noindent 所以\par
\noindent kxk−1 X\par
\noindent = ipi F (x)i−1−n F ′ (x).\par
\noindent F (x)n\par
\noindent i≥k\par
\noindent 故而\par
\noindent kxk−1 X\par
\noindent [x−1 ] n\par
\noindent = [x−1 ] ipi F (x)i−1−n F ′ (x)\par
\noindent F (x)\par
\noindent i≥k\par
}
\end{sourceblock}


\begin{explainblock}
这一段是原书论述链的一部分。读的时候把它当作桥梁：它通常负责把上一个定义、例子或定理，引到下一个计数方法。

\smallskip
\textbf{初学者检查点：}
\begin{itemize}[leftmargin=2em,itemsep=0.25em]
\item 找出这一段承接的上一个概念。
\item 找出它准备引出的下一个概念。
\item 把中间的关键等式或构造单独抄出来。
\end{itemize}

\smallskip
\textbf{本章读法提醒：} 第八章 Catalan 问题的核心常常是前缀不越界或递归分解。
\end{explainblock}


\begin{sourceblock}{第 8 章续读}
{\small\sloppy
\noindent X 1 d\par
\noindent = [x−1 ] F (x)i−n + [x−1 ]npn F (x)−1 F ′ (x).\par
\noindent i≥k,\par
\noindent i − n dx\par
\noindent i̸=n\par
\noindent P\par
\noindent n∈Z cn x , Z 为全体整数, L(x) 的留数记为 Res(L),\par
\noindent 又因为对任意洛朗级数 L(x) = n\par
\noindent 则 Res(L) = c−1 . 设 L(x) 的导数为 L′ (x), 则 Res(L′ ) = 0, 也即 L′ (x) 中不含 x−1\par
\noindent 项, 所以\par
\noindent X 1 d\par
\noindent [x−1 ] F (x)i−n = 0,\par
\noindent i≥k,\par
\noindent i − n dx\par
\noindent i̸=n\par
\noindent 故:\par
\noindent kxk−1\par
\noindent [x−1 ] = [x−1 ]npn F (x)−1 F ′ (x)\par
}
\end{sourceblock}


\begin{explainblock}
这一段是原书论述链的一部分。读的时候把它当作桥梁：它通常负责把上一个定义、例子或定理，引到下一个计数方法。

\smallskip
\textbf{初学者检查点：}
\begin{itemize}[leftmargin=2em,itemsep=0.25em]
\item 找出这一段承接的上一个概念。
\item 找出它准备引出的下一个概念。
\item 把中间的关键等式或构造单独抄出来。
\end{itemize}

\smallskip
\textbf{本章读法提醒：} 第八章 Catalan 问题的核心常常是前缀不越界或递归分解。
\end{explainblock}


\begin{sourceblock}{第 8 章续读}
{\small\sloppy
\noindent F (x)n\par
\noindent −1 a1 + 2a2 x + · · ·\par
\noindent = [x ]npn\par
\noindent a1 x + a2 x2 + · · ·\par
\noindent −1 1\par
\noindent = [x ]npn + ···\par
\noindent x\par
\noindent = npn .\par
\noindent 所以 n\par
\noindent x kxk−1\par
\noindent k[x n−k\par
\noindent ] = k[x−k ]F (x)−n = [x−1 ] = npn .\par
\noindent F (x) F (x)n\par
\noindent 与 (8.14) 相比较, 立得 (8.12) 式, 即:\par
\noindent n\par
\noindent x\par
\noindent n\par
\noindent n[x ]F <−1> k\par
}
\end{sourceblock}


\begin{explainblock}
这一段是原书论述链的一部分。读的时候把它当作桥梁：它通常负责把上一个定义、例子或定理，引到下一个计数方法。

\smallskip
\textbf{初学者检查点：}
\begin{itemize}[leftmargin=2em,itemsep=0.25em]
\item 找出这一段承接的上一个概念。
\item 找出它准备引出的下一个概念。
\item 把中间的关键等式或构造单独抄出来。
\end{itemize}

\smallskip
\textbf{本章读法提醒：} 第八章 Catalan 问题的核心常常是前缀不越界或递归分解。
\end{explainblock}


\begin{sourceblock}{第 8 章续读}
{\small\sloppy
\noindent (x) = k[x n−k\par
\noindent ] = k[x−k ]F (x)−n .\par
\noindent F (x)\par
}
\end{sourceblock}


\begin{explainblock}
这一段是原书论述链的一部分。读的时候把它当作桥梁：它通常负责把上一个定义、例子或定理，引到下一个计数方法。

\smallskip
\textbf{初学者检查点：}
\begin{itemize}[leftmargin=2em,itemsep=0.25em]
\item 找出这一段承接的上一个概念。
\item 找出它准备引出的下一个概念。
\item 把中间的关键等式或构造单独抄出来。
\end{itemize}

\smallskip
\textbf{本章读法提醒：} 第八章 Catalan 问题的核心常常是前缀不越界或递归分解。
\end{explainblock}


\begin{sourceblock}{例 8.4.2 求 F (x) = x(1 − x) 的复合逆.}
{\small\sloppy
\noindent 例 8.4.2 求 F (x) = x(1 − x) 的复合逆.\par
\noindent P∞\par
\noindent 解: 设 f (x) = F <−1> (x) = n=1 an x\par
\noindent n , 则由 (8.12) 得\par
\noindent nan = [x−1 ]F (x)−n = [x−1 ](x − x2 )−n\par
\noindent X\par
\noindent −1 k −n\par
\noindent = [x ] (−1) xk−n\par
\noindent k\par
\noindent k≥0\par
\noindent n−1 −n 2n − 2\par
\noindent = (−1) = ,\par
\noindent n−1 n−1\par
\noindent 168 第 8 章 Catalan 数与 Narayana 数\par
\noindent 所以\par
\noindent an = = Cn−1 , n ≥ 1,\par
\noindent n n−1\par
\noindent 故 √\par
}
\end{sourceblock}


\begin{explainblock}
例题是学习这本书最重要的部分。不要只看答案，要把每个限制条件变成一个计数动作：先安排谁、再插入谁、有没有重复、需不需要除以对称性。

\smallskip
\textbf{初学者检查点：}
\begin{itemize}[leftmargin=2em,itemsep=0.25em]
\item 先写出没有限制时的总数。
\item 逐条加入限制条件，观察总数怎样变化。
\item 最后检查有没有把同一种方案重复算了多次。
\end{itemize}

\smallskip
\textbf{本章读法提醒：} 第八章 Catalan 问题的核心常常是前缀不越界或递归分解。
\end{explainblock}


\begin{sourceblock}{第 8 章续读}
{\small\sloppy
\noindent 1− 1 − 4x\par
\noindent f (x) = xc(x) = .\par
\noindent P\par
}
\end{sourceblock}


\begin{explainblock}
这一段是原书论述链的一部分。读的时候把它当作桥梁：它通常负责把上一个定义、例子或定理，引到下一个计数方法。

\smallskip
\textbf{初学者检查点：}
\begin{itemize}[leftmargin=2em,itemsep=0.25em]
\item 找出这一段承接的上一个概念。
\item 找出它准备引出的下一个概念。
\item 把中间的关键等式或构造单独抄出来。
\end{itemize}

\smallskip
\textbf{本章读法提醒：} 第八章 Catalan 问题的核心常常是前缀不越界或递归分解。
\end{explainblock}


\begin{sourceblock}{例 8.4.3 证明: xe−x 的复合逆为 n≥1 n}
{\small\sloppy
\noindent 例 8.4.3 证明: xe−x 的复合逆为 n≥1 n\par
\noindent n−1 xn .\par
\noindent n!\par
\noindent 证明: 令 F (x) = xe−x , 由拉格朗日反演公式直接得:\par
\noindent [xn ](xe−x )<−1> = [x ]e\par
\noindent n\par
\noindent = = .\par
\noindent n (n − 1)! n!\par
}
\end{sourceblock}


\begin{explainblock}
例题是学习这本书最重要的部分。不要只看答案，要把每个限制条件变成一个计数动作：先安排谁、再插入谁、有没有重复、需不需要除以对称性。

\smallskip
\textbf{初学者检查点：}
\begin{itemize}[leftmargin=2em,itemsep=0.25em]
\item 先写出没有限制时的总数。
\item 逐条加入限制条件，观察总数怎样变化。
\item 最后检查有没有把同一种方案重复算了多次。
\end{itemize}

\smallskip
\textbf{本章读法提醒：} 第八章 Catalan 问题的核心常常是前缀不越界或递归分解。
\end{explainblock}


\begin{sourceblock}{注记 8.4.4 拉格朗日反演公式实际上是复分析上的一个重要定理, 有兴趣的读者}
{\small\sloppy
\noindent 注记 8.4.4 拉格朗日反演公式实际上是复分析上的一个重要定理, 有兴趣的读者\par
\noindent 可以参考 [22], 纯组合性证明可以参考 [3].\par
}
\end{sourceblock}


\begin{explainblock}
注记通常是在补充术语、记号、历史或一个容易忽略的约定。它未必给新公式，但常常决定你后面会不会把符号读错。

\smallskip
\textbf{初学者检查点：}
\begin{itemize}[leftmargin=2em,itemsep=0.25em]
\item 记下新符号的读法。
\item 确认这个约定是否只在本章使用。
\item 把注记里的限制条件补到前面的定义旁边。
\end{itemize}

\smallskip
\textbf{本章读法提醒：} 第八章 Catalan 问题的核心常常是前缀不越界或递归分解。
\end{explainblock}


\begin{sourceblock}{推论 8.4.5 设 G(x) ∈ C⟦x⟧, 且 G(0) 6= 0. 若}
{\small\sloppy
\noindent 推论 8.4.5 设 G(x) ∈ C⟦x⟧, 且 G(0) 6= 0. 若\par
\noindent f (x) = xG(f (x)), (8.15)\par
\noindent 则\par
\noindent n[xn ]f (x)k = k[xn−k ]G(x)n . (8.16)\par
\noindent x\par
\noindent 证明: 令 G(x) = F (x) , 由 8.15 得\par
\noindent f (x)\par
\noindent f (x) = x ,\par
\noindent F (f (x))\par
\noindent 即\par
\noindent F (f (x)) = x,\par
\noindent 故而 f (x) = F <−1> (x), 应用定理 8.4.3 即可得 (8.16).\par
\noindent 特别地, 在式 (8.16) 中取 k = 1, 则有\par
\noindent n[xn ]f (x) = [xn−1 ]G(x)n . (8.17)\par
\noindent 若设 f (x) = α1 x + α2 x2 + · · · , 则 n ≥ 1 时, 有\par
\noindent [xn−1 ]G(x)n\par
\noindent αn = ; (8.18)\par
\noindent n\par
}
\end{sourceblock}


\begin{explainblock}
定理段不要只看结论，要看它把哪两个计数方式连在一起。组合学证明最常见的技巧是：左边按一种标准数，右边按另一种标准数，然后说明两边数的是同一批对象。

\smallskip
\textbf{初学者检查点：}
\begin{itemize}[leftmargin=2em,itemsep=0.25em]
\item 先复述定理的假设，不要把适用范围扩大。
\item 找出证明中的基本计数原则：乘法、加法、双射、递推、容斥或生成函数。
\item 检查最后的等式化简是否和前面的对象解释一致。
\end{itemize}

\smallskip
\textbf{本章读法提醒：} 第八章 Catalan 问题的核心常常是前缀不越界或递归分解。
\end{explainblock}


\begin{sourceblock}{第 8 章续读}
{\small\sloppy
\noindent 若设 f (x) = γ1 x + γ2 x2! + γ3 x3! + · · · , 则 n ≥ 1 时, 有\par
\noindent γn = (n − 1)![xn−1 ]G(x)n . (8.19)\par
\noindent P\par
}
\end{sourceblock}


\begin{explainblock}
这一段是原书论述链的一部分。读的时候把它当作桥梁：它通常负责把上一个定义、例子或定理，引到下一个计数方法。

\smallskip
\textbf{初学者检查点：}
\begin{itemize}[leftmargin=2em,itemsep=0.25em]
\item 找出这一段承接的上一个概念。
\item 找出它准备引出的下一个概念。
\item 把中间的关键等式或构造单独抄出来。
\end{itemize}

\smallskip
\textbf{本章读法提醒：} 第八章 Catalan 问题的核心常常是前缀不越界或递归分解。
\end{explainblock}


\begin{sourceblock}{例 8.4.4 设 c(x) = n}
{\small\sloppy
\noindent 例 8.4.4 设 c(x) = n\par
\noindent n≥0 Cn x 是 Catalan 数 Cn 的生成函数, 已知 c(x) 满足\par
\noindent 下面的方程:\par
\noindent c(x) = 1 + xc(x)2 .\par
\noindent 利用拉格朗日反演公式求 Cn .\par
\noindent P\par
\noindent 解: 令 f (x) = c(x) − 1 = n 2\par
\noindent n≥1 Cn x , G(x) = (1 + x) , 则\par
\noindent f (x) = xG(f (x)).\par
\noindent n ≥ 1 时, 由拉格朗日反演公式可知\par
\noindent [xn−1 ](1 + x)2n 1 2n 1 2n\par
\noindent Cn = = = .\par
\noindent n n n−1 n+1 n\par
}
\end{sourceblock}


\begin{explainblock}
例题是学习这本书最重要的部分。不要只看答案，要把每个限制条件变成一个计数动作：先安排谁、再插入谁、有没有重复、需不需要除以对称性。

\smallskip
\textbf{初学者检查点：}
\begin{itemize}[leftmargin=2em,itemsep=0.25em]
\item 先写出没有限制时的总数。
\item 逐条加入限制条件，观察总数怎样变化。
\item 最后检查有没有把同一种方案重复算了多次。
\end{itemize}

\smallskip
\textbf{本章读法提醒：} 第八章 Catalan 问题的核心常常是前缀不越界或递归分解。
\end{explainblock}


\begin{sourceblock}{推论 8.4.6 对任意 H(x) ∈ C⟦x⟧ (或 H(x) 是形式洛朗级数), 则}
{\small\sloppy
\noindent 推论 8.4.6 对任意 H(x) ∈ C⟦x⟧ (或 H(x) 是形式洛朗级数), 则\par
\noindent n\par
\noindent n −1 n−1 ′ x\par
\noindent n[x ]H(F (x)) = [x ]H (x) . (8.20)\par
\noindent F (x)\par
\noindent 等价地, 若 f (x) = xG(f (x)), 则有\par
\noindent n[xn ]H(f (x)) = [xn−1 ]H ′ (x)G(x)n . (8.21)\par
\noindent P∞\par
\noindent 证明: 设 N1 为任意整数, 设 H(x) = k\par
\noindent k=N1 hk x . 则由拉格朗日反演公式知\par
\noindent X\par
\noindent n+1\par
\noindent n <−1>\par
\noindent n[x ]H(F (x)) = n[xn ]F <−1> (x)k\par
\noindent k≥N1\par
\noindent X\par
\noindent n+1 n\par
\noindent n−1 k−1 x\par
}
\end{sourceblock}


\begin{explainblock}
定理段不要只看结论，要看它把哪两个计数方式连在一起。组合学证明最常见的技巧是：左边按一种标准数，右边按另一种标准数，然后说明两边数的是同一批对象。

\smallskip
\textbf{初学者检查点：}
\begin{itemize}[leftmargin=2em,itemsep=0.25em]
\item 先复述定理的假设，不要把适用范围扩大。
\item 找出证明中的基本计数原则：乘法、加法、双射、递推、容斥或生成函数。
\item 检查最后的等式化简是否和前面的对象解释一致。
\end{itemize}

\smallskip
\textbf{本章读法提醒：} 第八章 Catalan 问题的核心常常是前缀不越界或递归分解。
\end{explainblock}


\begin{sourceblock}{第 8 章续读}
{\small\sloppy
\noindent = k[x ]x .\par
\noindent F (x)\par
\noindent k≥N1\par
\noindent 显然\par
\noindent X\par
\noindent n+1 n X\par
\noindent n+1 n\par
\noindent n−1 k−1 x n−1 k−1 x\par
\noindent k[x ]x = k[x ]x\par
\noindent F (x) F (x)\par
\noindent k≥N1 k≥N1\par
\noindent n\par
\noindent n−1 ′ x\par
\noindent = k[x ]H (x) ,\par
\noindent F (x)\par
\noindent 所以 n\par
\noindent n −1 n−1 ′ x\par
\noindent n[x ]H(F (x)) = [x ]H (x) .\par
}
\end{sourceblock}


\begin{explainblock}
这一段是原书论述链的一部分。读的时候把它当作桥梁：它通常负责把上一个定义、例子或定理，引到下一个计数方法。

\smallskip
\textbf{初学者检查点：}
\begin{itemize}[leftmargin=2em,itemsep=0.25em]
\item 找出这一段承接的上一个概念。
\item 找出它准备引出的下一个概念。
\item 把中间的关键等式或构造单独抄出来。
\end{itemize}

\smallskip
\textbf{本章读法提醒：} 第八章 Catalan 问题的核心常常是前缀不越界或递归分解。
\end{explainblock}


\begin{sourceblock}{第 8 章续读}
{\small\sloppy
\noindent F (x)\par
}
\end{sourceblock}


\begin{explainblock}
这一段是原书论述链的一部分。读的时候把它当作桥梁：它通常负责把上一个定义、例子或定理，引到下一个计数方法。

\smallskip
\textbf{初学者检查点：}
\begin{itemize}[leftmargin=2em,itemsep=0.25em]
\item 找出这一段承接的上一个概念。
\item 找出它准备引出的下一个概念。
\item 把中间的关键等式或构造单独抄出来。
\end{itemize}

\smallskip
\textbf{本章读法提醒：} 第八章 Catalan 问题的核心常常是前缀不越界或递归分解。
\end{explainblock}


\begin{sourceblock}{例 8.4.5 由 Lagrange 反演公式求 Nayarana 数 Nn,k 表达式:}
{\small\sloppy
\noindent 例 8.4.5 由 Lagrange 反演公式求 Nayarana 数 Nn,k 表达式:\par
\noindent Nn,k = .\par
\noindent n k k−1\par
\noindent 170 第 8 章 Catalan 数与 Narayana 数\par
\noindent 证明: 因为 Nayarana 多项式的生成函数为\par
\noindent ∞\par
\noindent X\par
\noindent H(x, t) = Hn (x)tn\par
\noindent n=0\par
\noindent ∞ X\par
\noindent X n\par
\noindent =1+ Nn,k xk tn ,\par
\noindent n=1 k=1\par
\noindent 所以\par
\noindent x ∞ X\par
\noindent X n\par
\noindent H ,t = 1 + Nn,k xk tn−k .\par
\noindent t\par
}
\end{sourceblock}


\begin{explainblock}
例题是学习这本书最重要的部分。不要只看答案，要把每个限制条件变成一个计数动作：先安排谁、再插入谁、有没有重复、需不需要除以对称性。

\smallskip
\textbf{初学者检查点：}
\begin{itemize}[leftmargin=2em,itemsep=0.25em]
\item 先写出没有限制时的总数。
\item 逐条加入限制条件，观察总数怎样变化。
\item 最后检查有没有把同一种方案重复算了多次。
\end{itemize}

\smallskip
\textbf{本章读法提醒：} 第八章 Catalan 问题的核心常常是前缀不越界或递归分解。
\end{explainblock}


\begin{sourceblock}{第 8 章续读}
{\small\sloppy
\noindent n=1 k=1\par
\noindent 再由 H(x, t) 满足的方程可得:\par
\noindent x 2 x\par
\noindent tH ,t + H , t (x − t − 1) + 1 = 0,\par
\noindent t t\par
\noindent 所以\par
\noindent x x t x\par
\noindent H , t − 1 = (H , t − 1) + 1 H , t − 1 + 1 x,\par
\noindent t t x t\par
\noindent x\par
\noindent 令H , t − 1 = x · f (x, t), 则得\par
\noindent t\par
\noindent f (x, t) = (1 + xf (x, t))(1 + tf (x, t))\par
\noindent P∞ n\par
\noindent 令 G(z) = (1+xz) (1 + tz), 则 f (x, t) = G(f (x, t)). 不妨设 g(x, t; z) = n=0 an (x, t)z\par
\noindent 满足方程\par
\noindent g(x, t; z) = zG(g(x, t; z)), (8.22)\par
\noindent 显然 g(x, t; z) 满足 g(x, t; 0) = 0. 不妨设\par
}
\end{sourceblock}


\begin{explainblock}
这一段是原书论述链的一部分。读的时候把它当作桥梁：它通常负责把上一个定义、例子或定理，引到下一个计数方法。

\smallskip
\textbf{初学者检查点：}
\begin{itemize}[leftmargin=2em,itemsep=0.25em]
\item 找出这一段承接的上一个概念。
\item 找出它准备引出的下一个概念。
\item 把中间的关键等式或构造单独抄出来。
\end{itemize}

\smallskip
\textbf{本章读法提醒：} 第八章 Catalan 问题的核心常常是前缀不越界或递归分解。
\end{explainblock}


\begin{sourceblock}{第 8 章续读}
{\small\sloppy
\noindent ∞\par
\noindent X\par
\noindent f (x, t) = g(x, t; 1) = an (x, t),\par
\noindent n=1\par
\noindent 这显然是可行的, 因为 z = 1 时, 我们有\par
\noindent f (x, t) = G(f (x, t)).\par
\noindent 设 n ≥ 1, 对等式 (8.22) 应用 Lagrange 反演公式得\par
\noindent an (x, t) = [z ]G(z)n\par
\noindent n\par
\noindent = [z n−1 ](1 + xz)n (1 + yz)n\par
\noindent n\par
\noindent n\par
\noindent 1 n−1 X X n n i j i+j\par
\noindent n\par
\noindent = [z ] xt z\par
\noindent n i j\par
\noindent i=0 j=0\par
\noindent = xt ,\par
}
\end{sourceblock}


\begin{explainblock}
这一段是原书论述链的一部分。读的时候把它当作桥梁：它通常负责把上一个定义、例子或定理，引到下一个计数方法。

\smallskip
\textbf{初学者检查点：}
\begin{itemize}[leftmargin=2em,itemsep=0.25em]
\item 找出这一段承接的上一个概念。
\item 找出它准备引出的下一个概念。
\item 把中间的关键等式或构造单独抄出来。
\end{itemize}

\smallskip
\textbf{本章读法提醒：} 第八章 Catalan 问题的核心常常是前缀不越界或递归分解。
\end{explainblock}


\begin{sourceblock}{第 8 章续读}
{\small\sloppy
\noindent n i+j=n−1 i j\par
\noindent 0≤i,j≤n\par
\noindent 所以\par
\noindent ∞\par
\noindent X\par
\noindent f (x, t) = an (x, t)\par
\noindent n=1\par
\noindent X∞\par
\noindent = xt\par
\noindent n i+j=n−1 i j\par
\noindent n=1\par
\noindent 0≤i,j≤n\par
\noindent n\par
\noindent ∞ X\par
\noindent X\par
\noindent n n\par
\noindent = xk−1 tn−k .\par
\noindent k−1 n−k\par
}
\end{sourceblock}


\begin{explainblock}
这一段是原书论述链的一部分。读的时候把它当作桥梁：它通常负责把上一个定义、例子或定理，引到下一个计数方法。

\smallskip
\textbf{初学者检查点：}
\begin{itemize}[leftmargin=2em,itemsep=0.25em]
\item 找出这一段承接的上一个概念。
\item 找出它准备引出的下一个概念。
\item 把中间的关键等式或构造单独抄出来。
\end{itemize}

\smallskip
\textbf{本章读法提醒：} 第八章 Catalan 问题的核心常常是前缀不越界或递归分解。
\end{explainblock}


\begin{sourceblock}{第 8 章续读}
{\small\sloppy
\noindent n=1 k=1\par
\noindent 所以\par
\noindent x ∞ X\par
\noindent X n\par
\noindent H , t − 1 = xf (x, t) = xk tn−k\par
\noindent t n k−1 n−k\par
\noindent n=1 k=1\par
\noindent 又\par
\noindent x ∞ X\par
\noindent X n\par
\noindent H ,t − 1 = Nn,k xk tn−k .\par
\noindent t\par
\noindent n=1 k=1\par
\noindent 所以比较系数可得:\par
\noindent Nn,k = = .\par
\noindent n k−1 n−k n k−1 k\par
\noindent 172 第 8 章 Catalan 数与 Narayana 数\par
}
\end{sourceblock}


\begin{explainblock}
这一段是原书论述链的一部分。读的时候把它当作桥梁：它通常负责把上一个定义、例子或定理，引到下一个计数方法。

\smallskip
\textbf{初学者检查点：}
\begin{itemize}[leftmargin=2em,itemsep=0.25em]
\item 找出这一段承接的上一个概念。
\item 找出它准备引出的下一个概念。
\item 把中间的关键等式或构造单独抄出来。
\end{itemize}

\smallskip
\textbf{本章读法提醒：} 第八章 Catalan 问题的核心常常是前缀不越界或递归分解。
\end{explainblock}


\begin{sourceblock}{习 题 八}
{\small\sloppy
\noindent 习 题 八\par
}
\end{sourceblock}


\begin{explainblock}
习题段不要先追求技巧。先给题目分类：它是在问排列、组合、递推、生成函数、容斥，还是结构计数。分类正确，工具通常就只剩一两个候选。

\smallskip
\textbf{初学者检查点：}
\begin{itemize}[leftmargin=2em,itemsep=0.25em]
\item 判断题目所属工具。
\item 列出变量和限制条件。
\item 先做小规模 n=2、3、4 的例子，确认直觉。
\end{itemize}

\smallskip
\textbf{本章读法提醒：} 第八章 Catalan 问题的核心常常是前缀不越界或递归分解。
\end{explainblock}


\begin{sourceblock}{习题 8.1 证明 Catalan numbers 满足以下关系:}
{\small\sloppy
\noindent 习题 8.1 证明 Catalan numbers 满足以下关系:\par
\noindent (1)C2 = 2C1 ;\par
\noindent (2)C3 = 3C2 − C1 ;\par
\noindent (3)C4 = 4C3 − 3C2 ;\par
\noindent (4)Cn = n1 Cn−1 − n−1\par
\noindent n−2\par
\noindent 证明从 (0, 0) 到 (2n − 2, 0) 的所有 Dyck 路它们与 x 轴的交点的个\par
}
\end{sourceblock}


\begin{explainblock}
习题段不要先追求技巧。先给题目分类：它是在问排列、组合、递推、生成函数、容斥，还是结构计数。分类正确，工具通常就只剩一两个候选。

\smallskip
\textbf{初学者检查点：}
\begin{itemize}[leftmargin=2em,itemsep=0.25em]
\item 判断题目所属工具。
\item 列出变量和限制条件。
\item 先做小规模 n=2、3、4 的例子，确认直觉。
\end{itemize}

\smallskip
\textbf{本章读法提醒：} 第八章 Catalan 问题的核心常常是前缀不越界或递归分解。
\end{explainblock}


\begin{sourceblock}{习题 8.2}
{\small\sloppy
\noindent 习题 8.2\par
\noindent 数的和为第 n 个 Catalan 数 Cn = n+1 n . 如图 8.8, 4 长 Dyck 路与 x 轴的交点个\par
\noindent 数为 5.\par
}
\end{sourceblock}


\begin{explainblock}
习题段不要先追求技巧。先给题目分类：它是在问排列、组合、递推、生成函数、容斥，还是结构计数。分类正确，工具通常就只剩一两个候选。

\smallskip
\textbf{初学者检查点：}
\begin{itemize}[leftmargin=2em,itemsep=0.25em]
\item 判断题目所属工具。
\item 列出变量和限制条件。
\item 先做小规模 n=2、3、4 的例子，确认直觉。
\end{itemize}

\smallskip
\textbf{本章读法提醒：} 第八章 Catalan 问题的核心常常是前缀不越界或递归分解。
\end{explainblock}


\begin{sourceblock}{习题 8.3 证明从 (0, 0) 到 (2n,}
{\small\sloppy
\noindent 习题 8.3 证明从 (0, 0) 到 (2n,\par
\noindent 0) 的 Dyck 路的所有高度为 1 的峰的个数的和\par
\noindent 为第 n 个 Catalan 数 Cn = n+1 . 如图 8.9, n = 3 时有 5 个高度为 1 的峰.\par
\noindent n\par
}
\end{sourceblock}


\begin{explainblock}
习题段不要先追求技巧。先给题目分类：它是在问排列、组合、递推、生成函数、容斥，还是结构计数。分类正确，工具通常就只剩一两个候选。

\smallskip
\textbf{初学者检查点：}
\begin{itemize}[leftmargin=2em,itemsep=0.25em]
\item 判断题目所属工具。
\item 列出变量和限制条件。
\item 先做小规模 n=2、3、4 的例子，确认直觉。
\end{itemize}

\smallskip
\textbf{本章读法提醒：} 第八章 Catalan 问题的核心常常是前缀不越界或递归分解。
\end{explainblock}


\begin{sourceblock}{习题 8.4 在一次选举中有两名候选人 P 和 Q, 有 2n 个人投票, 候选人 P 和}
{\small\sloppy
\noindent 习题 8.4 在一次选举中有两名候选人 P 和 Q, 有 2n 个人投票, 候选人 P 和\par
\noindent Q 分别收到 p 和 q 张票.\par
\noindent (1) 若 p = q = n, 求计票过程中支持候选人 P 的选票数始终不低于候选人 Q 的\par
\noindent 方法数;\par
\noindent (2) 若 p > q, 求计票过程中支持候选人 P 的选票数始终领先候选人 Q 的概率.\par
}
\end{sourceblock}


\begin{explainblock}
习题段不要先追求技巧。先给题目分类：它是在问排列、组合、递推、生成函数、容斥，还是结构计数。分类正确，工具通常就只剩一两个候选。

\smallskip
\textbf{初学者检查点：}
\begin{itemize}[leftmargin=2em,itemsep=0.25em]
\item 判断题目所属工具。
\item 列出变量和限制条件。
\item 先做小规模 n=2、3、4 的例子，确认直觉。
\end{itemize}

\smallskip
\textbf{本章读法提醒：} 第八章 Catalan 问题的核心常常是前缀不越界或递归分解。
\end{explainblock}


\begin{sourceblock}{习题 8.5 在 (x, y) 平面上定义三种移动方式:(1, 1), (1, −1), (1, 0). 从 (0, 0) 到}
{\small\sloppy
\noindent 习题 8.5 在 (x, y) 平面上定义三种移动方式:(1, 1), (1, −1), (1, 0). 从 (0, 0) 到\par
\noindent (n, 0) 且不经过 x 轴下方的路径称为 Motzkin 路. 如图 8.10 当 n = 8 时的一条\par
\noindent Motzkin 路. 证明把 (1, 0) 涂成红色或蓝色的从\par
\noindent (0, 0) 到 (n − 1, 0) 的 Motzkin 路的\par
\noindent 个数为第 n 个 Catalan 数 Cn = n+1 .\par
\noindent n\par
}
\end{sourceblock}


\begin{explainblock}
习题段不要先追求技巧。先给题目分类：它是在问排列、组合、递推、生成函数、容斥，还是结构计数。分类正确，工具通常就只剩一两个候选。

\smallskip
\textbf{初学者检查点：}
\begin{itemize}[leftmargin=2em,itemsep=0.25em]
\item 判断题目所属工具。
\item 列出变量和限制条件。
\item 先做小规模 n=2、3、4 的例子，确认直觉。
\end{itemize}

\smallskip
\textbf{本章读法提醒：} 第八章 Catalan 问题的核心常常是前缀不越界或递归分解。
\end{explainblock}


\begin{sourceblock}{习题 8.6 设有 2n 个不同的点分布在圆周上, 证明以这 2n 个点为端点的不相}
{\small\sloppy
\noindent 习题 8.6 设有 2n 个不同的点分布在圆周上, 证明以这 2n 个点为端点的不相\par
\noindent 交的连弦的方法数为第 n 个 Catalan 数 Cn = n+1 n . 如图 8.11 n = 3 时, 有 5 种\par
\noindent 方式.\par
\noindent 174 第 8 章 Catalan 数与 Narayana 数\par
\noindent A A\par
\noindent A A A\par
}
\end{sourceblock}


\begin{explainblock}
习题段不要先追求技巧。先给题目分类：它是在问排列、组合、递推、生成函数、容斥，还是结构计数。分类正确，工具通常就只剩一两个候选。

\smallskip
\textbf{初学者检查点：}
\begin{itemize}[leftmargin=2em,itemsep=0.25em]
\item 判断题目所属工具。
\item 列出变量和限制条件。
\item 先做小规模 n=2、3、4 的例子，确认直觉。
\end{itemize}

\smallskip
\textbf{本章读法提醒：} 第八章 Catalan 问题的核心常常是前缀不越界或递归分解。
\end{explainblock}


\begin{sourceblock}{习题 8.7 证明: 当 n > 3 时,Catalan 数 Cn 都不是素数.}
{\small\sloppy
\noindent 习题 8.7 证明: 当 n > 3 时,Catalan 数 Cn 都不是素数.\par
}
\end{sourceblock}


\begin{explainblock}
习题段不要先追求技巧。先给题目分类：它是在问排列、组合、递推、生成函数、容斥，还是结构计数。分类正确，工具通常就只剩一两个候选。

\smallskip
\textbf{初学者检查点：}
\begin{itemize}[leftmargin=2em,itemsep=0.25em]
\item 判断题目所属工具。
\item 列出变量和限制条件。
\item 先做小规模 n=2、3、4 的例子，确认直觉。
\end{itemize}

\smallskip
\textbf{本章读法提醒：} 第八章 Catalan 问题的核心常常是前缀不越界或递归分解。
\end{explainblock}


\begin{sourceblock}{习题 8.8 设 Tn 是将凸 n(≥ 3) 多边形通过连接它的顶点且使内部对角线不交}
{\small\sloppy
\noindent 习题 8.8 设 Tn 是将凸 n(≥ 3) 多边形通过连接它的顶点且使内部对角线不交\par
\noindent 将其分成若干三角形的方法数. 当 n ≥ 4 时证明:\par
\noindent n\par
\noindent (n − 3)Tn = (T3 Tn−1 + T4 Tn−2 + · · · + Tn−1 T3 ).\par
}
\end{sourceblock}


\begin{explainblock}
习题段不要先追求技巧。先给题目分类：它是在问排列、组合、递推、生成函数、容斥，还是结构计数。分类正确，工具通常就只剩一两个候选。

\smallskip
\textbf{初学者检查点：}
\begin{itemize}[leftmargin=2em,itemsep=0.25em]
\item 判断题目所属工具。
\item 列出变量和限制条件。
\item 先做小规模 n=2、3、4 的例子，确认直觉。
\end{itemize}

\smallskip
\textbf{本章读法提醒：} 第八章 Catalan 问题的核心常常是前缀不越界或递归分解。
\end{explainblock}


\begin{sourceblock}{习题 8.9 设两个序列 a1 , a2 , . . . , an 和 b1 , b2 , . . . , bn 中的所有的数都不相同. 将}
{\small\sloppy
\noindent 习题 8.9 设两个序列 a1 , a2 , . . . , an 和 b1 , b2 , . . . , bn 中的所有的数都不相同. 将\par
\noindent 序列 a1 , a2 , . . . , an 和 b1 , b2 , . . . , bn 合并成一个序列, 若它们在合并的序列中, 每个序\par
\noindent 列都保持它们在原来序列中的相对位置称为两序列的融合. 比如 n = 2 时,\par
\noindent a1 a2 b1 b2 ; a1 b1 a2 b2 ;\par
\noindent a1 b1 b2 a2 ; b1 a1 a2 b2 ;\par
\noindent b1 a1 b2 a2 ; b1 b2 a1 a2 ;\par
\noindent (1) 求将 a1 , a2 , . . . , an 和 b1 , b2 , . . . , bn 融合会产生多少种不同的序列?\par
\noindent (2) 若融合后的序列中 ai 不在 bi 的前面则称它有一个倒置; 例如\par
\noindent a1 b1 b2 a2 b3 a3\par
\noindent 中有 2 个倒置; 在 n = 2 时, 具有 0, 1 或 2 个倒置的序列个数分别有 2 个. 求\par
\noindent 将 a1 , a2 , . . . , an 和 b1 , b2 , . . . , bn 融合后, 有多少个序列恰具有 k 个倒置, 其中\par
\noindent k = 0, 1, 2, . . . , n.\par
}
\end{sourceblock}


\begin{explainblock}
习题段不要先追求技巧。先给题目分类：它是在问排列、组合、递推、生成函数、容斥，还是结构计数。分类正确，工具通常就只剩一两个候选。

\smallskip
\textbf{初学者检查点：}
\begin{itemize}[leftmargin=2em,itemsep=0.25em]
\item 判断题目所属工具。
\item 列出变量和限制条件。
\item 先做小规模 n=2、3、4 的例子，确认直觉。
\end{itemize}

\smallskip
\textbf{本章读法提醒：} 第八章 Catalan 问题的核心常常是前缀不越界或递归分解。
\end{explainblock}


\begin{sourceblock}{习题 8.10 设 c(x) 为 Catalan 数的生成函数. 证明:}
{\small\sloppy
\noindent 习题 8.10 设 c(x) 为 Catalan 数的生成函数. 证明:\par
\noindent (1)\par
\noindent ∞\par
\noindent X\par
\noindent 2n + 1 n c(x)\par
\noindent x =√ ,\par
\noindent n=0\par
\noindent n 1 − 4x\par
\noindent (2)\par
\noindent c2 (x)\par
\noindent c′ (x) = √ ,\par
\noindent (3)\par
\noindent n\par
\noindent X 2k\par
\noindent (n + 1)Cn+1 = Cn−k+1 .\par
\noindent k\par
\noindent k=0\par
}
\end{sourceblock}


\begin{explainblock}
习题段不要先追求技巧。先给题目分类：它是在问排列、组合、递推、生成函数、容斥，还是结构计数。分类正确，工具通常就只剩一两个候选。

\smallskip
\textbf{初学者检查点：}
\begin{itemize}[leftmargin=2em,itemsep=0.25em]
\item 判断题目所属工具。
\item 列出变量和限制条件。
\item 先做小规模 n=2、3、4 的例子，确认直觉。
\end{itemize}

\smallskip
\textbf{本章读法提醒：} 第八章 Catalan 问题的核心常常是前缀不越界或递归分解。
\end{explainblock}


\begin{sourceblock}{习题 8.11 证明: 由数 1, 2, . . . , 2n 构成且同时满足下面三组条件:}
{\small\sloppy
\noindent 习题 8.11 证明: 由数 1, 2, . . . , 2n 构成且同时满足下面三组条件:\par
\noindent x11 < x12 < · · · < x1n ,\par
\noindent x21 < x22 < · · · < x2n ,\par
\noindent x1i < x2i , ∀1 ≤ i ≤ n\par
\noindent 的二行列阵 " \#\par
\noindent x11 x12 · · · x1n\par
\noindent x21 x22 · · · x2n\par
\noindent 共有 Cn 个.\par
}
\end{sourceblock}


\begin{explainblock}
习题段不要先追求技巧。先给题目分类：它是在问排列、组合、递推、生成函数、容斥，还是结构计数。分类正确，工具通常就只剩一两个候选。

\smallskip
\textbf{初学者检查点：}
\begin{itemize}[leftmargin=2em,itemsep=0.25em]
\item 判断题目所属工具。
\item 列出变量和限制条件。
\item 先做小规模 n=2、3、4 的例子，确认直觉。
\end{itemize}

\smallskip
\textbf{本章读法提醒：} 第八章 Catalan 问题的核心常常是前缀不越界或递归分解。
\end{explainblock}


\begin{sourceblock}{习题 8.12 在数学与计算机科学中的堆栈可排序列 (Stack-sortable permuta-}
{\small\sloppy
\noindent 习题 8.12 在数学与计算机科学中的堆栈可排序列 (Stack-sortable permuta-\par
\noindent tion) 的结构 (或称为树排列 (tree permutation)) 是一个排列; 排列中的元素可以通\par
\noindent 过一种算法进行排序，而该算法的内部存储限于单个堆栈数据结构. 具体而言, 将标\par
\noindent 有 1 到 n 的木块放到一个不透明的盒子中, 堆栈可排序列可看作按下面算法得到的\par
\noindent 排列:\par
\noindent (a) 初始化一个空栈;\par
\noindent (b) 若盒子非空, 则从盒子中任取一个木块, 标号设为 x; 若 x 比栈中最上方的数\par
\noindent 值大, 则将栈中数据依次输出; 否则将标有 x 的木块放到堆栈中最上方;\par
\noindent (c) 若盒子为空且堆栈非空则将栈中数据依次输出;\par
\noindent 显然上述算法输出的数据是集合 [n] 上的排列. 比如 n = 3 时, 得到的排列有:\par
\noindent 123, 132, 213, 不含有排列 231, 312, 321. 证明:\par
\noindent (1) 由栈可排得到的排列是避免 231-模式的. (所谓避免 231-模式的排列是指排\par
\noindent 列 π = π1 π2 · · · πn 不存在 i, j, k(i < j < k) 使得 πk < πi < πj .)\par
\noindent (2) 由栈可排得到的排列的个数等于 Cn , 其中 Cn 是第 n 个 Catalan 数.\par
}
\end{sourceblock}


\begin{explainblock}
习题段不要先追求技巧。先给题目分类：它是在问排列、组合、递推、生成函数、容斥，还是结构计数。分类正确，工具通常就只剩一两个候选。

\smallskip
\textbf{初学者检查点：}
\begin{itemize}[leftmargin=2em,itemsep=0.25em]
\item 判断题目所属工具。
\item 列出变量和限制条件。
\item 先做小规模 n=2、3、4 的例子，确认直觉。
\end{itemize}

\smallskip
\textbf{本章读法提醒：} 第八章 Catalan 问题的核心常常是前缀不越界或递归分解。
\end{explainblock}


\begin{sourceblock}{习题 8.13 设 n 长避免 231-模式的排列中具有 k 个下降位的序列个数为}
{\small\sloppy
\noindent 习题 8.13 设 n 长避免 231-模式的排列中具有 k 个下降位的序列个数为\par
\noindent C(n, k), 则由上题知:\par
\noindent X\par
\noindent n\par
\noindent Cn = C(n, k).\par
\noindent k=0\par
\noindent 176 第 8 章 Catalan 数与 Narayana 数\par
\noindent 设 Fn (t) 是 \{Cn,k \}nk=0 的生成函数, 即\par
\noindent X X\par
\noindent n\par
\noindent Fn (t) = tdes(ω) = Cn,k tk ,\par
\noindent ω∈Sn (231) k=0\par
\noindent 其中 Sn (231) 表示 n 长避免 231-模式的排列组成的集合, des(ω) 表示 ω 中下降位的\par
\noindent 个数.\par
\noindent (1) 证明当 n ≥ 1 时,\par
\noindent X\par
\noindent n−2\par
\noindent Fn (t) = Fn−1 (t) + t Fi (t)Fn−1−i (t).\par
}
\end{sourceblock}


\begin{explainblock}
习题段不要先追求技巧。先给题目分类：它是在问排列、组合、递推、生成函数、容斥，还是结构计数。分类正确，工具通常就只剩一两个候选。

\smallskip
\textbf{初学者检查点：}
\begin{itemize}[leftmargin=2em,itemsep=0.25em]
\item 判断题目所属工具。
\item 列出变量和限制条件。
\item 先做小规模 n=2、3、4 的例子，确认直觉。
\end{itemize}

\smallskip
\textbf{本章读法提醒：} 第八章 Catalan 问题的核心常常是前缀不越界或递归分解。
\end{explainblock}


\begin{sourceblock}{第 8 章续读}
{\small\sloppy
\noindent i=0\par
\noindent (2) 设\par
\noindent X\par
\noindent F (t, z) = Fn (t)z n .\par
\noindent n≥0\par
\noindent 证明: p\par
\noindent 1 + z(t − 1) − 1 − 2z(t + 1) + z 2 (t − 1)2\par
\noindent F (t, z) = .\par
\noindent 2tz\par
}
\end{sourceblock}


\begin{explainblock}
这一段是原书论述链的一部分。读的时候把它当作桥梁：它通常负责把上一个定义、例子或定理，引到下一个计数方法。

\smallskip
\textbf{初学者检查点：}
\begin{itemize}[leftmargin=2em,itemsep=0.25em]
\item 找出这一段承接的上一个概念。
\item 找出它准备引出的下一个概念。
\item 把中间的关键等式或构造单独抄出来。
\end{itemize}

\smallskip
\textbf{本章读法提醒：} 第八章 Catalan 问题的核心常常是前缀不越界或递归分解。
\end{explainblock}


\begin{sourceblock}{习题 8.14 设 C(x) 是 Catalan 数的生成函数, 证明:}
{\small\sloppy
\noindent 习题 8.14 设 C(x) 是 Catalan 数的生成函数, 证明:\par
\noindent X X\par
\noindent k(C(x) − 1)k = 4n−1 xn .\par
\noindent k≥1 n≥1\par
\noindent (提示: 参考文献 [7])\par
}
\end{sourceblock}


\begin{explainblock}
习题段不要先追求技巧。先给题目分类：它是在问排列、组合、递推、生成函数、容斥，还是结构计数。分类正确，工具通常就只剩一两个候选。

\smallskip
\textbf{初学者检查点：}
\begin{itemize}[leftmargin=2em,itemsep=0.25em]
\item 判断题目所属工具。
\item 列出变量和限制条件。
\item 先做小规模 n=2、3、4 的例子，确认直觉。
\end{itemize}

\smallskip
\textbf{本章读法提醒：} 第八章 Catalan 问题的核心常常是前缀不越界或递归分解。
\end{explainblock}


\begin{sourceblock}{习题 8.15 设 Cn 是 Catalan 数, 证明: 当且仅当 n 是 Mersenne† 数, 即存在}
{\small\sloppy
\noindent 习题 8.15 设 Cn 是 Catalan 数, 证明: 当且仅当 n 是 Mersenne† 数, 即存在\par
\noindent m ≥ 1, 使得 n = 2m − 1 时, Cn 是奇数.\par
\noindent (提示: 当 n ≥ 1 时, 由 Catalan 数 Cn 的递推关系:\par
\noindent Cn = C0 Cn−1 + C1 Cn−2 + · · · Cn−1 C0\par
\noindent 可得:\par
\noindent \par
\noindent  2(C0 Cn−1 + C1 Cn−2 + · · · + C n2 −1 C n2 ), n 是偶数,\par
\noindent Cn =\par
\noindent 2(C0 Cn−1 + C1 Cn−2 + · · · + C n−3 C n+1 ) + C 2n−1 , n 是奇数\par
\noindent 因此, 对于 n > 0, Cn 是奇数当且仅当 n 和 C n−1 都是奇数. 根据同样的理论我\par
\noindent 们知 Cn 是奇数当且仅当 n−1 n−1\par
\noindent 2 和 C n−3 都是奇数或 2 = 0. 这样递推下去, 我们可\par
\noindent 得结论 (1).)\par
}
\end{sourceblock}


\begin{explainblock}
习题段不要先追求技巧。先给题目分类：它是在问排列、组合、递推、生成函数、容斥，还是结构计数。分类正确，工具通常就只剩一两个候选。

\smallskip
\textbf{初学者检查点：}
\begin{itemize}[leftmargin=2em,itemsep=0.25em]
\item 判断题目所属工具。
\item 列出变量和限制条件。
\item 先做小规模 n=2、3、4 的例子，确认直觉。
\end{itemize}

\smallskip
\textbf{本章读法提醒：} 第八章 Catalan 问题的核心常常是前缀不越界或递归分解。
\end{explainblock}


\begin{sourceblock}{习题 8.16 设 τ 集合 [m] 上的一个排列, 我们称集合 [n] (n ≥ m) 上的排列 π}
{\small\sloppy
\noindent 习题 8.16 设 τ 集合 [m] 上的一个排列, 我们称集合 [n] (n ≥ m) 上的排列 π\par
\noindent 是避免 τ -模式的, 如果在 π 中任意删除 n − m 个元素后得到的 m 长排列与排列 τ\par
\noindent 中的大小排序都不一样. 比如排列 4321 是避免 231-模式的. 记集合 [n] 上所有排列\par
\noindent †\par
\noindent Marin Mersenne(1588–1648): 法国数学家.\par
\noindent 为 Sn , Sn 中避免模式 τ 的排列组成的集合记为 Sn (τ ). 设 π = π1 π2 · · · πn ∈ Sn , 映\par
\noindent 射 C 和 R 分别定义为\par
\noindent C(π) = (n + 1 − π1 )(n + 1 − π2 ) . . . (n + 1 − π1 ),\par
\noindent R(π) = πn πn−1 . . . π1 .\par
\noindent 显然映射 C 与 R 是集合 Sn 到自身的对合. 证明:\par
\noindent (1) 对于任意的模式 τ , 对合 C 是集合 Sn (τ ) 到集合 Sn (C(τ )) 之间的一一映射.\par
\noindent (2) 对于任意的模式 τ , 对合 R 是集合 Sn (τ ) 到集合 Sn (R(τ )) 之间的一一映射.\par
\noindent (3) 对于任意的排列 τ ∈ S3 , Sn (τ ) 与长为 2n 的 Dyck 路一一对应.\par
\noindent 提示: 由对合 C 与 R 的定义可知结论 (1) 与 (2) 是显然的. 集合 [3] 中所有的排列\par
\noindent 有 123, 132, 213, 231, 312, 321, 且有\par
\noindent C(123) = 321, C(132) = 312, C(213) = 231, C(231) = 213, C(312) = 132, C(321) = 123,\par
\noindent R(123) = 321, R(132) = 231, R(213) = 312, R(231) = 132, R(312) = 213, R(321) = 123.\par
\noindent 再由 (1) 与 (2) 可得避免 123-模式的与避免 321-模式的个数相等, 避免 132-模式的\par
}
\end{sourceblock}


\begin{explainblock}
习题段不要先追求技巧。先给题目分类：它是在问排列、组合、递推、生成函数、容斥，还是结构计数。分类正确，工具通常就只剩一两个候选。

\smallskip
\textbf{初学者检查点：}
\begin{itemize}[leftmargin=2em,itemsep=0.25em]
\item 判断题目所属工具。
\item 列出变量和限制条件。
\item 先做小规模 n=2、3、4 的例子，确认直觉。
\end{itemize}

\smallskip
\textbf{本章读法提醒：} 第八章 Catalan 问题的核心常常是前缀不越界或递归分解。
\end{explainblock}


\begin{sourceblock}{第 8 章续读}
{\small\sloppy
\noindent 排列的个数, 避免 312-模式的排列的个数, 避免 231-模式的排列的个数, 避免 213-模\par
\noindent 式的排列的个数都相等. 所以我们只需要分别构造集合 Sn (132) 和集合 Sn (123) 与\par
\noindent 长为 2n 的 Dyck 路所组成的集合之间的双射.\par
\noindent 设 π = π1 π2 . . . πn ∈ Sn . 若 πi 的左边元素均大于 πi ，即对于任意的 j < i 都有\par
\noindent πi < πj 成立，则称 πi 为从左到右极小元. 比如排列 74352681 的从左到右极小元为\par
\noindent 7, 4, 3, 2, 1.\par
\noindent 记 u 为上升移动 (1, 1), d 为下降移动 (1, −1), Dyck 路可用 u 和 d 组成的字表示,\par
\noindent 在这个意义下, 我们认为 Dyck 字等同于 Dyck 路, 并记 2n 长的 Dyck 路组成的集\par
\noindent 合为 Dn . 下面按照如下的规则定义一个映射: Φ : Sn → Dn . 给定一个排列 π ∈ Sn ,\par
\noindent 设 m1 , m2 , . . . , mk (m1 > m2 > · · · > mk ) 为 π 的从左到右极小元, 则 mk = 1 且 π\par
\noindent 可以唯一地分解成如下形式：\par
\noindent π = m1 w1 m2 w2 . . . mk wk (8.23)\par
\noindent 对于 1 ≤ i < k, wi 是排列 π 中位于 mi 与 mi+1 之间的子序列, wk 为 mk 后的子序\par
\noindent 列 (可为空序列). 设 m0 = n + 1, 从左往右读取分解形式 (8.23), 每个从左到右极小\par
\noindent 元 mi 对应 mi−1 − mi 个连续向上升移动 u. 每个子序列 wi 对应 |wi | + 1 个连续下\par
\noindent 降移动 d，其中 |wi | 是 wi 的长度, 这样得到的格路记为 Φ(π), 即\par
\noindent Φ(π) = uu · · · u\}\par
\noindent | \{z dd · · · d\} · · ·\par
}
\end{sourceblock}


\begin{explainblock}
这一段是原书论述链的一部分。读的时候把它当作桥梁：它通常负责把上一个定义、例子或定理，引到下一个计数方法。

\smallskip
\textbf{初学者检查点：}
\begin{itemize}[leftmargin=2em,itemsep=0.25em]
\item 找出这一段承接的上一个概念。
\item 找出它准备引出的下一个概念。
\item 把中间的关键等式或构造单独抄出来。
\end{itemize}

\smallskip
\textbf{本章读法提醒：} 第八章 Catalan 问题的核心常常是前缀不越界或递归分解。
\end{explainblock}


\begin{sourceblock}{第 8 章续读}
{\small\sloppy
\noindent | \{z uu · · · u\}\par
\noindent | \{z dd · · · d\} .\par
\noindent | \{z\par
\noindent m0 − m1 个 u |w1 | + 1 个 d mk−1 − mk 个 u |wk | + 1 个 d\par
\noindent 178 第 8 章 Catalan 数与 Narayana 数\par
\noindent 其中 u 的个数为\par
\noindent (m0 − m1 ) + (m1 − m2 ) + · · · + (mk−1 − mk ) = n + 1 − mk = n,\par
\noindent d 的个数为\par
\noindent (|w1 | + 1) + (|w2 | + 1) + · · · + (|wk | + 1) = n − k + k = n,\par
\noindent 所以 Φ(π) 是长为 2n. 故 Φ(π) 是一条从 (0, 0) 到 (2n, 0) 的格路. 下面为了说明\par
\noindent Φ(π) 是 Dyck 路, 只需说明从左到右读 Φ(π) 中 u 的个数总不小于 d 的个数, 为此,\par
\noindent 我们只需证明对任意 1 ≤ i ≤ k\par
\noindent (m0 − m1 ) + (m1 − m2 ) + · · · + (mi−1 − mi ) ≥ (|w1 | + 1) + (|w2 | + 1) + · · · + (|wi | + 1),\par
\noindent (8.24)\par
\noindent 显然 i = k 时成立. 下面我们说明 1 ≤ i < k 时(8.24)也成立. 因为\par
\noindent (m0 −m1 )+(m1 −m2 )+· · ·+(mi−1 −mi ) = n+1−mi ≥ n+1−(mi+1 +1) = n−mi+1 ,\par
\noindent 又 wi 是位于极小元 mi 与 mi+1 之间, 所以 wi 最大长度为 mi − mi+1 − 1, 故\par
\noindent (|w1 | + 1) + (|w2 | + 1) + · · · + (|wi | + 1) ≤ m1 − mi+1 ≤ n − mi+1 .\par
}
\end{sourceblock}


\begin{explainblock}
这一段是原书论述链的一部分。读的时候把它当作桥梁：它通常负责把上一个定义、例子或定理，引到下一个计数方法。

\smallskip
\textbf{初学者检查点：}
\begin{itemize}[leftmargin=2em,itemsep=0.25em]
\item 找出这一段承接的上一个概念。
\item 找出它准备引出的下一个概念。
\item 把中间的关键等式或构造单独抄出来。
\end{itemize}

\smallskip
\textbf{本章读法提醒：} 第八章 Catalan 问题的核心常常是前缀不越界或递归分解。
\end{explainblock}


\begin{sourceblock}{第 8 章续读}
{\small\sloppy
\noindent 或者从另一个角度, (|w1 |+1)+(|w2 |+1)+· · ·+(|wi |+1) 可以表示序列 m1 w1 · · · mi wi\par
\noindent 的长度, 它显然是小于等于 n − mi+1 . 所以 (8.24) 成立. 故 Φ(π) ∈ Dn , 即映射 Φ 是\par
\noindent 集合 Sn 到集合 Dn 的一个映射.\par
}
\end{sourceblock}


\begin{explainblock}
这一段是原书论述链的一部分。读的时候把它当作桥梁：它通常负责把上一个定义、例子或定理，引到下一个计数方法。

\smallskip
\textbf{初学者检查点：}
\begin{itemize}[leftmargin=2em,itemsep=0.25em]
\item 找出这一段承接的上一个概念。
\item 找出它准备引出的下一个概念。
\item 把中间的关键等式或构造单独抄出来。
\end{itemize}

\smallskip
\textbf{本章读法提醒：} 第八章 Catalan 问题的核心常常是前缀不越界或递归分解。
\end{explainblock}


\begin{sourceblock}{例如当 π = 74352681 时, π 可以唯一地分解成 7w1 4w2 3w3 2w4 1w5 , 其中 w1 , w2 ,}
{\small\sloppy
\noindent 例如当 π = 74352681 时, π 可以唯一地分解成 7w1 4w2 3w3 2w4 1w5 , 其中 w1 , w2 ,\par
\noindent w5 为空序列, w3 = 5，w4 = 68. 则\par
\noindent Φ(74352681) = uuduuududdudddud,\par
\noindent 如图 8.12.\par
\noindent 我们设 Φ132 和 Φ123 是映射 Φ 分别限制到集合 Sn (132) 和集合 Sn 123 上得到的\par
\noindent 映射. 下面我们通过分别构造它们的逆映射来说明 Φ132 和 Φ123 是双射.\par
\noindent 为了证明 Φ132 是一个双射, 我们构造一个映射 Ψ132 : Dn → Sn (132). 给定一条\par
\noindent Dyck 路 P ∈ Dn , 我们用 1, 2 . . . n 从右向左依次给上升移动 u 标号，然后从左到右\par
\noindent 按下面的规则给下降移动用 1, 2, . . . n 标号: 对于一个未标号的下降移动 d，我们用\par
\noindent 其左边向上升移动的标号中最小的且未被用于对下降移动标号的最小标号给 d 标\par
\noindent 号. 具体地说, 设从左向右第 i(1 ≤ i ≤ n) 个下降位置之前的上升移动标号集合记为\par
\noindent Ai , 下降移动标号集合为 Bi , 则第 i 个下降位置标号为 ai , 即\par
\noindent ai = min Ai \textbackslash{} Bi .\par
\noindent 我们令\par
\noindent Ψ132 (P ) = a1 a2 · · · an .\par
\noindent 由 Ψ132 的构造可知 Ψ132 (P ) ∈ Sn (132), 并且 Φ132 与 Ψ132 互为逆映射. 由此便证\par
\noindent 明了 Φ132 是双射.\par
\noindent 比如, P 如图 8.13, 我们得到 Ψ132 (P ) = 74352681.\par
}
\end{sourceblock}


\begin{explainblock}
例题是学习这本书最重要的部分。不要只看答案，要把每个限制条件变成一个计数动作：先安排谁、再插入谁、有没有重复、需不需要除以对称性。

\smallskip
\textbf{初学者检查点：}
\begin{itemize}[leftmargin=2em,itemsep=0.25em]
\item 先写出没有限制时的总数。
\item 逐条加入限制条件，观察总数怎样变化。
\item 最后检查有没有把同一种方案重复算了多次。
\end{itemize}

\smallskip
\textbf{本章读法提醒：} 第八章 Catalan 问题的核心常常是前缀不越界或递归分解。
\end{explainblock}


\begin{sourceblock}{第 8 章续读}
{\small\sloppy
\noindent 下面我们构造映射 Ψ123 : Dn → Sn (123). 给定一条 Dyck 路 P ∈ Dn , 我们用\par
\noindent 1, 2 . . . n 从右向左依次给上升移动 u 标号，然后从左到右地按下面的规则给下降移\par
\noindent 动 d 标号. 对于一个未标号的下降移动 d，若 d 的前一步是上升移动，我们用其左\par
\noindent 边对上升移动 u 的标号中未被用过标记下降移动 d 的最小标号给 d 标号；否则, 我\par
\noindent 们用其左边上升移动 u 的标号中未被用过标记下降移动 d 的最大标号给 d 标号. 具\par
\noindent 体而言, 设第 i 个下降移动 d 前的上升移动标号数值的集合和下降移动标号数值集\par
\noindent 合分别记为 Ai , Bi . 从左向右对下降移动 d 标号, 第 i 个未标号的下降移动 d 若它\par
\noindent 前一步是上升移动则对该下降位置用 bi = min Ai \textbackslash{} Bi , 否则用 bi = max Ai \textbackslash{} Bi . 我\par
\noindent 们令\par
\noindent Ψ123 (P ) = b1 b2 · · · bn .\par
\noindent 按照 Ψ123 的构造，我们容易验证 Ψ123 (P ) ∈ Sn (123), 并且 Φ123 与 Ψ123 互为逆映\par
\noindent 射, 所以 Φ123 是 Sn (123) 与 Dn 的双射.\par
\noindent 比如设 π = 74382651。π 可以唯一地分解成 7w1 4w2 3w3 2w4 1w5 , 其中 w1 ，w2 ，\par
\noindent w5 为空序列，w3 = 8，w4 = 65. 图 8.14 给出了排列 π 所对应的格路 Φ123 (π). 当\par
\noindent P 为如图8.15 所给出的 Dyck 路时, 我们得到 Ψ123 (P ) = 74382651.\par
}
\end{sourceblock}


\begin{explainblock}
这一段是原书论述链的一部分。读的时候把它当作桥梁：它通常负责把上一个定义、例子或定理，引到下一个计数方法。

\smallskip
\textbf{初学者检查点：}
\begin{itemize}[leftmargin=2em,itemsep=0.25em]
\item 找出这一段承接的上一个概念。
\item 找出它准备引出的下一个概念。
\item 把中间的关键等式或构造单独抄出来。
\end{itemize}

\smallskip
\textbf{本章读法提醒：} 第八章 Catalan 问题的核心常常是前缀不越界或递归分解。
\end{explainblock}


\begin{sourceblock}{注记 8.4.7 映射 Φ 是由 Krattenthaler 在文 [13] 中给出, 但他并没有给出 Φ 的逆}
{\small\sloppy
\noindent 注记 8.4.7 映射 Φ 是由 Krattenthaler 在文 [13] 中给出, 但他并没有给出 Φ 的逆\par
\noindent 映射的显性表示. 映射 Ψ123 与 Ψ132 是由 Chen, Deng, Du [9] 所给出的. Bandlow\par
\noindent 与 Killpatrick [4] 构造了 Sn (312) 与 Dn 之间的双射，并且证明了该双射将排列中\par
\noindent 180 第 8 章 Catalan 数与 Narayana 数\par
\noindent 的逆序数转换成对应 Dyck 路的面积. 事实上，他们的双射等价于双射 Φ ◦ C.\par
\noindent ‘\par
}
\end{sourceblock}


\begin{explainblock}
注记通常是在补充术语、记号、历史或一个容易忽略的约定。它未必给新公式，但常常决定你后面会不会把符号读错。

\smallskip
\textbf{初学者检查点：}
\begin{itemize}[leftmargin=2em,itemsep=0.25em]
\item 记下新符号的读法。
\item 确认这个约定是否只在本章使用。
\item 把注记里的限制条件补到前面的定义旁边。
\end{itemize}

\smallskip
\textbf{本章读法提醒：} 第八章 Catalan 问题的核心常常是前缀不越界或递归分解。
\end{explainblock}


\chapter{整数分拆}
\begin{roadmapblock}
本章保留原书正文的主要数学骨架，并在每个定义、定理、例题和论述片段后加入解读。
阅读时不要把目标放在“背下答案”上，而是放在“看懂这一步为什么这样数”上。

\smallskip
\textbf{本章最低目标：} 能说出每个公式左边在数什么、右边在数什么，以及为什么两边相等。
\end{roadmapblock}


\begin{sourceblock}{第 9 章正文}
{\small\sloppy
\noindent 整数分拆\par
}
\end{sourceblock}


\begin{explainblock}
这一段是原书论述链的一部分。读的时候把它当作桥梁：它通常负责把上一个定义、例子或定理，引到下一个计数方法。

\smallskip
\textbf{初学者检查点：}
\begin{itemize}[leftmargin=2em,itemsep=0.25em]
\item 找出这一段承接的上一个概念。
\item 找出它准备引出的下一个概念。
\item 把中间的关键等式或构造单独抄出来。
\end{itemize}

\smallskip
\textbf{本章读法提醒：} 第九章分拆题先判断部件是否允许重复、是否要求不同、是否限制奇偶。
\end{explainblock}


\begin{sourceblock}{§ 9.1 整数分拆的定义与图表示}
{\small\sloppy
\noindent § 9.1. 整数分拆的定义与图表示\par
\noindent 整数分拆 (Integer Partition) 问题最早出现在莱布尼兹在 1699 年写给约翰 · 伯\par
\noindent 努利的信中. 顾名思义, 正整数 n 的分拆实际上就是指将 n 表示成正整数的和的形\par
\noindent 式. 显然正整数 n 本身就是它自身的一个分拆.\par
\noindent 欧拉对整数问题的研究起源于 1740 年诺地在给他信中询问欧拉是否知道将正整\par
\noindent 数 n 表示若干不同或可相同的正整数和的方法数. 欧拉从此开始了对整数分拆的系\par
\noindent 统研究, 并首次将幂级数应用到对整数分拆的研究中, 由此给出了很多漂亮的公式.\par
\noindent 而这种方法正是组合数学中常用的生成函数方法, 因此欧拉也被认为是生成函数理\par
\noindent 论的创始人.\par
}
\end{sourceblock}


\begin{explainblock}
这一节先不要急着背公式。先问三个问题：对象是什么，哪些限制条件不能丢，最后到底是在数所有对象、还是在数某个统计量的分布。把这三件事说清楚，后面的公式才会有位置。

\smallskip
\textbf{初学者检查点：}
\begin{itemize}[leftmargin=2em,itemsep=0.25em]
\item 把本节出现的新对象写成一句话定义。
\item 把每个公式左边和右边分别翻译成“在数什么”。
\item 找一个最小非平凡例子，手工列出所有对象。
\end{itemize}

\smallskip
\textbf{本章读法提醒：} 第九章分拆题先判断部件是否允许重复、是否要求不同、是否限制奇偶。
\end{explainblock}


\begin{sourceblock}{定义 9.1.1 设 λ1 , λ2 , . . . , λm 都是正整数, 若它们满足 λ1 ≥ λ2 ≥ · · · ≥ λm , 且其}
{\small\sloppy
\noindent 定义 9.1.1 设 λ1 , λ2 , . . . , λm 都是正整数, 若它们满足 λ1 ≥ λ2 ≥ · · · ≥ λm , 且其\par
\noindent 和为 n, 即:\par
\noindent λ1 + λ2 + · · · + λm = n,\par
\noindent 若和式均不论次序, 则称 λ1 , λ2 , . . . , λm 是正整数 n 的一个无序分拆简称为分拆. 记\par
\noindent 为 λ = (λ1 , λ2 , . . . , λm ). 称 λi (1 ≤ i ≤ m) 为分拆 λ 的部分, λ1 称为分拆 λ 的最大\par
\noindent 部分; 称 m 为分拆 λ 的长, 记为 ℓ(λ).\par
}
\end{sourceblock}


\begin{explainblock}
定义段的任务是定边界。这里要特别注意参数范围、是否允许重复、是否考虑顺序、对象有没有标签。后面的每一道例题和证明，本质上都在反复使用这些边界。

\smallskip
\textbf{初学者检查点：}
\begin{itemize}[leftmargin=2em,itemsep=0.25em]
\item 圈出被定义的对象名称。
\item 圈出参数范围，例如 n、k 是否要求正整数。
\item 判断它是有序对象、无序对象、带标签对象还是结构对象。
\end{itemize}

\smallskip
\textbf{本章读法提醒：} 第九章分拆题先判断部件是否允许重复、是否要求不同、是否限制奇偶。
\end{explainblock}


\begin{sourceblock}{例如, λ = (6, 2, 1) 是 9 一个分拆, 其中最大部分为 6, ℓ(λ) = 3.}
{\small\sloppy
\noindent 例如, λ = (6, 2, 1) 是 9 一个分拆, 其中最大部分为 6, ℓ(λ) = 3.\par
}
\end{sourceblock}


\begin{explainblock}
例题是学习这本书最重要的部分。不要只看答案，要把每个限制条件变成一个计数动作：先安排谁、再插入谁、有没有重复、需不需要除以对称性。

\smallskip
\textbf{初学者检查点：}
\begin{itemize}[leftmargin=2em,itemsep=0.25em]
\item 先写出没有限制时的总数。
\item 逐条加入限制条件，观察总数怎样变化。
\item 最后检查有没有把同一种方案重复算了多次。
\end{itemize}

\smallskip
\textbf{本章读法提醒：} 第九章分拆题先判断部件是否允许重复、是否要求不同、是否限制奇偶。
\end{explainblock}


\begin{sourceblock}{定义 9.1.2 若对 n 的一个分拆 λ 的部分及分拆的长度没有任何限制, 则称 λ 为一}
{\small\sloppy
\noindent 定义 9.1.2 若对 n 的一个分拆 λ 的部分及分拆的长度没有任何限制, 则称 λ 为一\par
\noindent 个无限制的分拆.\par
}
\end{sourceblock}


\begin{explainblock}
定义段的任务是定边界。这里要特别注意参数范围、是否允许重复、是否考虑顺序、对象有没有标签。后面的每一道例题和证明，本质上都在反复使用这些边界。

\smallskip
\textbf{初学者检查点：}
\begin{itemize}[leftmargin=2em,itemsep=0.25em]
\item 圈出被定义的对象名称。
\item 圈出参数范围，例如 n、k 是否要求正整数。
\item 判断它是有序对象、无序对象、带标签对象还是结构对象。
\end{itemize}

\smallskip
\textbf{本章读法提醒：} 第九章分拆题先判断部件是否允许重复、是否要求不同、是否限制奇偶。
\end{explainblock}


\begin{sourceblock}{定义 9.1.3 正整数 n 的所有分拆的个数称为 n 的分拆数, 记为 p(n), 并规定}
{\small\sloppy
\noindent 定义 9.1.3 正整数 n 的所有分拆的个数称为 n 的分拆数, 记为 p(n), 并规定\par
\noindent p(0) = 1. 有时也称 p(n) 为整数分拆函数.\par
}
\end{sourceblock}


\begin{explainblock}
定义段的任务是定边界。这里要特别注意参数范围、是否允许重复、是否考虑顺序、对象有没有标签。后面的每一道例题和证明，本质上都在反复使用这些边界。

\smallskip
\textbf{初学者检查点：}
\begin{itemize}[leftmargin=2em,itemsep=0.25em]
\item 圈出被定义的对象名称。
\item 圈出参数范围，例如 n、k 是否要求正整数。
\item 判断它是有序对象、无序对象、带标签对象还是结构对象。
\end{itemize}

\smallskip
\textbf{本章读法提醒：} 第九章分拆题先判断部件是否允许重复、是否要求不同、是否限制奇偶。
\end{explainblock}


\begin{sourceblock}{例如 5 的分拆有:}
{\small\sloppy
\noindent 例如 5 的分拆有:\par
\noindent 5=4+1\par
\noindent =3+2\par
\noindent =3+1+1\par
\noindent =2+2+1\par
\noindent =2+1+1+1\par
\noindent = 1 + 1 + 1 + 1 + 1.\par
\noindent 所以 p(5) = 7.\par
\noindent 有上例可知分拆 λ 的部分可以相同. 若设 λ 中有 mi 个部分为 i, 则可将 λ 表示\par
\noindent 成\par
\noindent λ = 1m1 , 2m2 , · · · .\par
\noindent 除了这些表示方法, 分拆还可以由图来表示. 例如 8 的一个分拆 3, 2, 2, 1 可以表示\par
\noindent 成 λ = (3, 2, 2, 1), 也可以把每一个部分都用点表示, 点的个数由分拆对应的部分来\par
\noindent 确定, 将这些点横向列出, 不同部分按行分开, 并要求行与行之间左对齐. 在这种规\par
\noindent 则下, λ = (3, 2, 2, 1) 就可以表示为:\par
\noindent • • •\par
\noindent • •\par
\noindent • •\par
}
\end{sourceblock}


\begin{explainblock}
例题是学习这本书最重要的部分。不要只看答案，要把每个限制条件变成一个计数动作：先安排谁、再插入谁、有没有重复、需不需要除以对称性。

\smallskip
\textbf{初学者检查点：}
\begin{itemize}[leftmargin=2em,itemsep=0.25em]
\item 先写出没有限制时的总数。
\item 逐条加入限制条件，观察总数怎样变化。
\item 最后检查有没有把同一种方案重复算了多次。
\end{itemize}

\smallskip
\textbf{本章读法提醒：} 第九章分拆题先判断部件是否允许重复、是否要求不同、是否限制奇偶。
\end{explainblock}


\begin{sourceblock}{第 9 章续读}
{\small\sloppy
\noindent •\par
\noindent 显然这种表示与分拆 λ 是一一对应的, 不同分拆对应不同的图表示. 这种图表示\par
\noindent 的方法由 James Joseph Sylvester 和他的学生 Fabian Franklin 首次引入, 现在称为\par
\noindent Ferrers 图. 如果我们用方块来代替点, 则得到了下面的图:\par
\noindent 这也是分拆的一种图表示, 称为 Young 图. 如果在 Young 图中的方块中填写数字,\par
\noindent 则称为 Young 表.\par
}
\end{sourceblock}


\begin{explainblock}
这一段是原书论述链的一部分。读的时候把它当作桥梁：它通常负责把上一个定义、例子或定理，引到下一个计数方法。

\smallskip
\textbf{初学者检查点：}
\begin{itemize}[leftmargin=2em,itemsep=0.25em]
\item 找出这一段承接的上一个概念。
\item 找出它准备引出的下一个概念。
\item 把中间的关键等式或构造单独抄出来。
\end{itemize}

\smallskip
\textbf{本章读法提醒：} 第九章分拆题先判断部件是否允许重复、是否要求不同、是否限制奇偶。
\end{explainblock}


\begin{sourceblock}{定理 9.1.4 设 λ = (λ1 , λ2 , · · · , λl ), 则}
{\small\sloppy
\noindent 定理 9.1.4 设 λ = (λ1 , λ2 , · · · , λl ), 则\par
\noindent X X λ′\par
\noindent (i − 1)λi = i\par
\noindent .\par
\noindent i≥1 i≥1\par
\noindent 证明: 我们可以把 λ 表示为 Young 图的形式, 并在第 i 行的每一个格子内都\par
\noindent 填入数字 i − 1, 这样等式左端就可以看作按行指标进行求和, 而等式右端则可以看\par
\noindent 作按列指标进行求和, 显然二者相等.\par
}
\end{sourceblock}


\begin{explainblock}
定理段不要只看结论，要看它把哪两个计数方式连在一起。组合学证明最常见的技巧是：左边按一种标准数，右边按另一种标准数，然后说明两边数的是同一批对象。

\smallskip
\textbf{初学者检查点：}
\begin{itemize}[leftmargin=2em,itemsep=0.25em]
\item 先复述定理的假设，不要把适用范围扩大。
\item 找出证明中的基本计数原则：乘法、加法、双射、递推、容斥或生成函数。
\item 检查最后的等式化简是否和前面的对象解释一致。
\end{itemize}

\smallskip
\textbf{本章读法提醒：} 第九章分拆题先判断部件是否允许重复、是否要求不同、是否限制奇偶。
\end{explainblock}


\begin{sourceblock}{注记 9.1.5 上个定理是 Young 表的一个简单应用. 关于 Young 表的理论这里不}
{\small\sloppy
\noindent 注记 9.1.5 上个定理是 Young 表的一个简单应用. 关于 Young 表的理论这里不\par
\noindent 再详细介绍, 但是 Young 表是组合数学中很重要的研究对象之一.\par
}
\end{sourceblock}


\begin{explainblock}
注记通常是在补充术语、记号、历史或一个容易忽略的约定。它未必给新公式，但常常决定你后面会不会把符号读错。

\smallskip
\textbf{初学者检查点：}
\begin{itemize}[leftmargin=2em,itemsep=0.25em]
\item 记下新符号的读法。
\item 确认这个约定是否只在本章使用。
\item 把注记里的限制条件补到前面的定义旁边。
\end{itemize}

\smallskip
\textbf{本章读法提醒：} 第九章分拆题先判断部件是否允许重复、是否要求不同、是否限制奇偶。
\end{explainblock}


\begin{sourceblock}{定义 9.1.6 将分拆 λ 对应的 Ferrers 图或 Young 图作转置, 即把此示图的第 1 列,}
{\small\sloppy
\noindent 定义 9.1.6 将分拆 λ 对应的 Ferrers 图或 Young 图作转置, 即把此示图的第 1 列,\par
\noindent 第 2 列, · · · 分别作为另一示图的第 1 行, 第 2 行, · · · . 则后一示图代表另一个分拆,\par
\noindent 称为分拆 λ 的共轭, 记为 λ′ . 若 λ = λ′ , 称分拆 λ 是自共轭分拆.\par
}
\end{sourceblock}


\begin{explainblock}
定义段的任务是定边界。这里要特别注意参数范围、是否允许重复、是否考虑顺序、对象有没有标签。后面的每一道例题和证明，本质上都在反复使用这些边界。

\smallskip
\textbf{初学者检查点：}
\begin{itemize}[leftmargin=2em,itemsep=0.25em]
\item 圈出被定义的对象名称。
\item 圈出参数范围，例如 n、k 是否要求正整数。
\item 判断它是有序对象、无序对象、带标签对象还是结构对象。
\end{itemize}

\smallskip
\textbf{本章读法提醒：} 第九章分拆题先判断部件是否允许重复、是否要求不同、是否限制奇偶。
\end{explainblock}


\begin{sourceblock}{例如 λ = (3, 2, 2, 1) 的共轭分拆是 λ′ = (4, 3, 1).}
{\small\sloppy
\noindent 例如 λ = (3, 2, 2, 1) 的共轭分拆是 λ′ = (4, 3, 1).\par
\noindent 显然分拆 λ 的共轭分拆 λ′ 唯一由 λ 确定, 并且通过这种共轭关系, 我们可以直\par
\noindent 接得到一下定理:\par
}
\end{sourceblock}


\begin{explainblock}
例题是学习这本书最重要的部分。不要只看答案，要把每个限制条件变成一个计数动作：先安排谁、再插入谁、有没有重复、需不需要除以对称性。

\smallskip
\textbf{初学者检查点：}
\begin{itemize}[leftmargin=2em,itemsep=0.25em]
\item 先写出没有限制时的总数。
\item 逐条加入限制条件，观察总数怎样变化。
\item 最后检查有没有把同一种方案重复算了多次。
\end{itemize}

\smallskip
\textbf{本章读法提醒：} 第九章分拆题先判断部件是否允许重复、是否要求不同、是否限制奇偶。
\end{explainblock}


\begin{sourceblock}{定理 9.1.7 设 n 为正整数, 对于任意的 1 ≤ m ≤ n, n 的具有 m 个部分的分拆个}
{\small\sloppy
\noindent 定理 9.1.7 设 n 为正整数, 对于任意的 1 ≤ m ≤ n, n 的具有 m 个部分的分拆个\par
\noindent 数等于 n 的最大部分为 m 的分拆个数.\par
}
\end{sourceblock}


\begin{explainblock}
定理段不要只看结论，要看它把哪两个计数方式连在一起。组合学证明最常见的技巧是：左边按一种标准数，右边按另一种标准数，然后说明两边数的是同一批对象。

\smallskip
\textbf{初学者检查点：}
\begin{itemize}[leftmargin=2em,itemsep=0.25em]
\item 先复述定理的假设，不要把适用范围扩大。
\item 找出证明中的基本计数原则：乘法、加法、双射、递推、容斥或生成函数。
\item 检查最后的等式化简是否和前面的对象解释一致。
\end{itemize}

\smallskip
\textbf{本章读法提醒：} 第九章分拆题先判断部件是否允许重复、是否要求不同、是否限制奇偶。
\end{explainblock}


\begin{sourceblock}{§ 9.2 整数分拆的生成函数}
{\small\sloppy
\noindent § 9.2. 整数分拆的生成函数\par
}
\end{sourceblock}


\begin{explainblock}
这一节先不要急着背公式。先问三个问题：对象是什么，哪些限制条件不能丢，最后到底是在数所有对象、还是在数某个统计量的分布。把这三件事说清楚，后面的公式才会有位置。

\smallskip
\textbf{初学者检查点：}
\begin{itemize}[leftmargin=2em,itemsep=0.25em]
\item 把本节出现的新对象写成一句话定义。
\item 把每个公式左边和右边分别翻译成“在数什么”。
\item 找一个最小非平凡例子，手工列出所有对象。
\end{itemize}

\smallskip
\textbf{本章读法提醒：} 第九章分拆题先判断部件是否允许重复、是否要求不同、是否限制奇偶。
\end{explainblock}


\begin{sourceblock}{定理 9.2.1 (欧拉) 设 n 为正整数, p(n) 为分拆函数, 则}
{\small\sloppy
\noindent 定理 9.2.1 (欧拉) 设 n 为正整数, p(n) 为分拆函数, 则\par
\noindent ∞\par
\noindent X ∞\par
\noindent Y 1\par
\noindent p(n)q n = .\par
\noindent n=0 k=1\par
\noindent 证明: 首先, 我们可以重记原式左端的形式幂级数. 它的每一项都可以分别看\par
\noindent 作是 p(n) 个 q n 作和, 那么每一个 q n 我们就都可以一一地对应到一个权重为 n 的\par
\noindent 分拆.\par
\noindent 另一方面, 我们可以把右面的每一项乘积都分别写成一个形式幂级数的形式, 例\par
\noindent 如对一个固定的 k0 , 其对应的那项乘积可以写成\par
\noindent = 1 + q k0 + q k0 +k0 + q k0 +k0 +k0 + · · · .\par
\noindent 我们知道, 如果一个分拆 λ 中有 mi 个部分为 i, 则可将 λ 表示成\par
\noindent λ = 1m1 , 2m2 , · · · .\par
\noindent 注意到, 这种表示是一一对应的.\par
\noindent 所以, 如果一个分拆 λ 中有 mi 个部分为 i, 我们就可以对应到无穷乘积里的第 i\par
\noindent 项中的, 次数为 imi 的单项式. 如果我们把每个第 i 项中的, 次数为 imi 的这些单项\par
\noindent 式都乘起来, 就可以对应到分拆 λ. 这样我们就得到了从一个分拆, 到一些单项式的\par
}
\end{sourceblock}


\begin{explainblock}
定理段不要只看结论，要看它把哪两个计数方式连在一起。组合学证明最常见的技巧是：左边按一种标准数，右边按另一种标准数，然后说明两边数的是同一批对象。

\smallskip
\textbf{初学者检查点：}
\begin{itemize}[leftmargin=2em,itemsep=0.25em]
\item 先复述定理的假设，不要把适用范围扩大。
\item 找出证明中的基本计数原则：乘法、加法、双射、递推、容斥或生成函数。
\item 检查最后的等式化简是否和前面的对象解释一致。
\end{itemize}

\smallskip
\textbf{本章读法提醒：} 第九章分拆题先判断部件是否允许重复、是否要求不同、是否限制奇偶。
\end{explainblock}


\begin{sourceblock}{第 9 章续读}
{\small\sloppy
\noindent 乘积的一个一一对应.\par
\noindent 因此, 每一个权重为 n 的分拆 λ, 既可以对应到原式左端的某一项 q n , 也可以对\par
\noindent 应到一个那种带尖括号的表示, 从而对应到原式右端的一些单项式的乘积, 并且这些\par
\noindent 项的次数之和必然为 n. 所以, 如果我们把原式右端的乘积全部乘开, 重新写成一个\par
\noindent 形式幂级数的形式, 那么它每一项的次数必然是原式左端的对应次数前的系数. 不难\par
\noindent 验证, 我们之前用过的所有的加法和乘法都可以在一个包含整数集的, 以 q 为未定\par
\noindent 元的形式幂级数环中进行运算, 所以原式的左右两端在形式幂级数的意义下是相等\par
\noindent 的.\par
\noindent 对于有限制的分拆我们有以下定理.\par
}
\end{sourceblock}


\begin{explainblock}
这一段是原书论述链的一部分。读的时候把它当作桥梁：它通常负责把上一个定义、例子或定理，引到下一个计数方法。

\smallskip
\textbf{初学者检查点：}
\begin{itemize}[leftmargin=2em,itemsep=0.25em]
\item 找出这一段承接的上一个概念。
\item 找出它准备引出的下一个概念。
\item 把中间的关键等式或构造单独抄出来。
\end{itemize}

\smallskip
\textbf{本章读法提醒：} 第九章分拆题先判断部件是否允许重复、是否要求不同、是否限制奇偶。
\end{explainblock}


\begin{sourceblock}{定理 9.2.2 设 n 为正整数, 集合 S 是由正整数组成的集合. 集合 S 既可以为有限}
{\small\sloppy
\noindent 定理 9.2.2 设 n 为正整数, 集合 S 是由正整数组成的集合. 集合 S 既可以为有限\par
\noindent 集也定义为无穷集合. 设 P (S) 为 n 的所有分拆中每个部分都在 S 中的分拆所构成\par
\noindent 的集合, 则\par
\noindent X Y 1\par
\noindent xl(λ) q |λ| = .\par
\noindent λ∈P (S) s∈S\par
\noindent 证明: 证明的方法完全类似于前面定理的证明. 需要注意的是, 这里每一个乘\par
\noindent 积中对应的单项式变成了 xmi q imi , 其中 q imi 提供了 mi 个 i, 而 xmi 提供的就是,\par
\noindent 在对应的分拆中, 这些为 i 的部分对该分拆的长 ℓλ 的贡献 mi .\par
}
\end{sourceblock}


\begin{explainblock}
定理段不要只看结论，要看它把哪两个计数方式连在一起。组合学证明最常见的技巧是：左边按一种标准数，右边按另一种标准数，然后说明两边数的是同一批对象。

\smallskip
\textbf{初学者检查点：}
\begin{itemize}[leftmargin=2em,itemsep=0.25em]
\item 先复述定理的假设，不要把适用范围扩大。
\item 找出证明中的基本计数原则：乘法、加法、双射、递推、容斥或生成函数。
\item 检查最后的等式化简是否和前面的对象解释一致。
\end{itemize}

\smallskip
\textbf{本章读法提醒：} 第九章分拆题先判断部件是否允许重复、是否要求不同、是否限制奇偶。
\end{explainblock}


\begin{sourceblock}{§ 9.3 欧拉分拆定理}
{\small\sloppy
\noindent § 9.3. 欧拉分拆定理\par
}
\end{sourceblock}


\begin{explainblock}
这一节先不要急着背公式。先问三个问题：对象是什么，哪些限制条件不能丢，最后到底是在数所有对象、还是在数某个统计量的分布。把这三件事说清楚，后面的公式才会有位置。

\smallskip
\textbf{初学者检查点：}
\begin{itemize}[leftmargin=2em,itemsep=0.25em]
\item 把本节出现的新对象写成一句话定义。
\item 把每个公式左边和右边分别翻译成“在数什么”。
\item 找一个最小非平凡例子，手工列出所有对象。
\end{itemize}

\smallskip
\textbf{本章读法提醒：} 第九章分拆题先判断部件是否允许重复、是否要求不同、是否限制奇偶。
\end{explainblock}


\begin{sourceblock}{定理 9.3.1 (欧拉分拆定理) 设 n 为正整数, n 的各部分都不同的分拆组成的集合}
{\small\sloppy
\noindent 定理 9.3.1 (欧拉分拆定理) 设 n 为正整数, n 的各部分都不同的分拆组成的集合\par
\noindent 记为 Hd (n), 其个数记为 pd (n); n 的各部分均为奇数的分拆组成的集合记为 Ho (n),\par
\noindent 其个数记为 po (n); 则 pd (n) = po (n).\par
\noindent 证明: 显然 pd (n) 与 po (n) 的分别生成函数为:\par
\noindent X\par
\noindent pd (n)q n = (1 + q)(1 + q 2 ) · · · ,\par
\noindent n≥0\par
\noindent X 1\par
\noindent po (n)q n = ,\par
\noindent (1 − q 1 )(1 − q 3 ) · · ·\par
\noindent n≥0\par
\noindent 而\par
\noindent (1 + q)(1 − q)(1 + q 2 )(1 − q 2 ) · · · 1\par
\noindent (1 + q)(1 + q 2 ) · · · = = ,\par
\noindent (1 − q)(1 − q 2 ) · · · (1 − q 1 )(1 − q 3 ) · · ·\par
\noindent 所以\par
\noindent X X\par
\noindent pd (n)q n = po (n)q n .\par
}
\end{sourceblock}


\begin{explainblock}
定理段不要只看结论，要看它把哪两个计数方式连在一起。组合学证明最常见的技巧是：左边按一种标准数，右边按另一种标准数，然后说明两边数的是同一批对象。

\smallskip
\textbf{初学者检查点：}
\begin{itemize}[leftmargin=2em,itemsep=0.25em]
\item 先复述定理的假设，不要把适用范围扩大。
\item 找出证明中的基本计数原则：乘法、加法、双射、递推、容斥或生成函数。
\item 检查最后的等式化简是否和前面的对象解释一致。
\end{itemize}

\smallskip
\textbf{本章读法提醒：} 第九章分拆题先判断部件是否允许重复、是否要求不同、是否限制奇偶。
\end{explainblock}


\begin{sourceblock}{第 9 章续读}
{\small\sloppy
\noindent n≥0 n≥0\par
}
\end{sourceblock}


\begin{explainblock}
这一段是原书论述链的一部分。读的时候把它当作桥梁：它通常负责把上一个定义、例子或定理，引到下一个计数方法。

\smallskip
\textbf{初学者检查点：}
\begin{itemize}[leftmargin=2em,itemsep=0.25em]
\item 找出这一段承接的上一个概念。
\item 找出它准备引出的下一个概念。
\item 把中间的关键等式或构造单独抄出来。
\end{itemize}

\smallskip
\textbf{本章读法提醒：} 第九章分拆题先判断部件是否允许重复、是否要求不同、是否限制奇偶。
\end{explainblock}


\begin{sourceblock}{定理得证.}
{\small\sloppy
\noindent 定理得证.\par
\noindent 下面我们给出欧拉分拆定理两种组合证明.\par
\noindent 证明: 法 1: (Sylvester) 我们想要建立 Ho (n) 到 Hd (n) 之间的双射 S. 首\par
\noindent 先我们考虑 Ho (n) 中元素 λ = \{λ1 , λ2 , . . . , λk \} 的 Ferrers 图并将每一行居中排\par
\noindent 列, 我们首先将第一行的右 λ12+1 个点与中轴线上的点相连, 再将第一行左 λ12−1 个\par
\noindent 点与中轴线左边第一列中的点相连, 再将第二行 λ22−1 个点与中轴线右边第一列中\par
\noindent 还未被连上的点相连, 以此类推我们可以按点的相连关系唯一地得到 n 的一个新\par
\noindent 划分 S(λ) = µ = \{µ1 , µ2 , . . . , µr \}. 将相连得到的图形成为一个 hook, 其横向部\par
\noindent 分称为 arm, 纵向部分称为 leg. 如下图我们将 λ = \{7, 7, 5, 5, 5, 5, 1\} 变成了\par
\noindent µ = \{9, 7, 6, 4, 2\}.\par
\noindent 注意到在得到的 µ 中, µ2k−1 − µ2k 为其对应的 leg 的差加 1, 而 µ2k 的 arm 不短\par
\noindent 于 µ2k+1 且 µ2k 的 leg 要比 µ2k+1 的 leg 长 1, 故 µ1 > µ2 > · · · µr , 即 µ ∈ Hd (n).\par
\noindent 我们接下来只需说明任给 Hd (n) 中的元素 µ, 我们可以唯一地找到与之对应的\par
\noindent λ = S −1 (µ). 对于给定的 µ = \{µ1 , µ2 , . . . , µr \}, 我们首先计算其对应的 S −1 的中\par
\noindent 轴线长度 l, 实际上我们注意到 l = µ1 − µ2 + µ3 − µ4 + · · · µ2r−1 − µ2r, 其中若 µ\par
\noindent 只含有奇数个部分, 则在其最后加上一个空部分. 当 l 确定后, 我们可以唯一地画出\par
\noindent S −1 (µ) 的 Ferrers 图, 从而可以唯一地求出 λ = S −1 (µ).\par
\noindent 法 2: (Glaisher) 我们现在用另一种方式来构造 Ho (n) 到 Hd (n) 的双射 G. 对任\par
}
\end{sourceblock}


\begin{explainblock}
定理段不要只看结论，要看它把哪两个计数方式连在一起。组合学证明最常见的技巧是：左边按一种标准数，右边按另一种标准数，然后说明两边数的是同一批对象。

\smallskip
\textbf{初学者检查点：}
\begin{itemize}[leftmargin=2em,itemsep=0.25em]
\item 先复述定理的假设，不要把适用范围扩大。
\item 找出证明中的基本计数原则：乘法、加法、双射、递推、容斥或生成函数。
\item 检查最后的等式化简是否和前面的对象解释一致。
\end{itemize}

\smallskip
\textbf{本章读法提醒：} 第九章分拆题先判断部件是否允许重复、是否要求不同、是否限制奇偶。
\end{explainblock}


\begin{sourceblock}{第 9 章续读}
{\small\sloppy
\noindent 意 λ ∈ Ho (n), 可唯一地将 µ 表示为 µ = \{r1k1 , r2k2 , . . . , rsks \}. 其中 ri 为奇数, ki 为 ri\par
\noindent P\par
\noindent 在 µ 中出现的次数. 将 ki 写成二进制表示 ki = j≥0 aj 2j , ai = 0, 1. 令 µij = 2j ri ,\par
\noindent 从而唯一地得到一个 n 的划分 G(λ) = µ = \{µij \}i,j≥0 . 例如 λ = \{35 , 54 , 73 \}, 则由\par
\noindent 5 = 22 +1, 4 = 22 , 3 = 2+1 推出 G(λ) = \{22 ·3, 3, 22 ·5, 2·7, 7\} = \{20, 14, 12, 7, 3\}.\par
\noindent 下面要说明对任意的 µ = \{µ1 , µ2 , · · · , µr \} ∈ D(n), 我们都可以唯一地确定\par
\noindent −1\par
\noindent G (µ) = λ. 因为对任意的正整数 k, 分解 k = 2t b, 2 ∤ b 是唯一的, 对一个奇数 b 设\par
\noindent P µi\par
\noindent 所有以 b 为最大奇因子的 µi 构成集合 U , 则可定义 b 在 λ 中出现的次数为 b\par
\noindent µi ∈U\par
\noindent 即可唯一地得到 G−1 (µ) = λ.\par
}
\end{sourceblock}


\begin{explainblock}
这一段是原书论述链的一部分。读的时候把它当作桥梁：它通常负责把上一个定义、例子或定理，引到下一个计数方法。

\smallskip
\textbf{初学者检查点：}
\begin{itemize}[leftmargin=2em,itemsep=0.25em]
\item 找出这一段承接的上一个概念。
\item 找出它准备引出的下一个概念。
\item 把中间的关键等式或构造单独抄出来。
\end{itemize}

\smallskip
\textbf{本章读法提醒：} 第九章分拆题先判断部件是否允许重复、是否要求不同、是否限制奇偶。
\end{explainblock}


\begin{sourceblock}{§ 9.4 欧拉五角数定理}
{\small\sloppy
\noindent § 9.4. 欧拉五角数定理\par
\noindent 这节我们介绍欧拉五角数定理. 我们称 k(3k−1)\par
\noindent 为其计数了类似以下的图形的点数.\par
\noindent 类似地还可以定义 n(n+1)\par
}
\end{sourceblock}


\begin{explainblock}
这一节先不要急着背公式。先问三个问题：对象是什么，哪些限制条件不能丢，最后到底是在数所有对象、还是在数某个统计量的分布。把这三件事说清楚，后面的公式才会有位置。

\smallskip
\textbf{初学者检查点：}
\begin{itemize}[leftmargin=2em,itemsep=0.25em]
\item 把本节出现的新对象写成一句话定义。
\item 把每个公式左边和右边分别翻译成“在数什么”。
\item 找一个最小非平凡例子，手工列出所有对象。
\end{itemize}

\smallskip
\textbf{本章读法提醒：} 第九章分拆题先判断部件是否允许重复、是否要求不同、是否限制奇偶。
\end{explainblock}


\begin{sourceblock}{定理 9.4.1 (欧拉五角数定理)}
{\small\sloppy
\noindent 定理 9.4.1 (欧拉五角数定理)\par
\noindent ∞\par
\noindent Y ∞\par
\noindent X j(3j−1) j(3j+1)\par
\noindent (1 − q j ) = 1 + (−1)j (q 2 +q 2 )\par
\noindent j=1 j=1\par
\noindent = 1 − q − q 2 + q 5 + q 7 − q 12 − q 15 + · · · .\par
\noindent 其中 j(3j−1)\par
\noindent 记 p(n) 为 n 的分拆个数, 我们已经证明过如下式子成立:\par
\noindent ∞\par
\noindent X 1\par
\noindent p(n)q n = ∞ .\par
\noindent Q\par
\noindent n=0 (1 − q j )\par
\noindent j=1\par
\noindent 因此有\par
\noindent ∞\par
\noindent Y ∞\par
}
\end{sourceblock}


\begin{explainblock}
定理段不要只看结论，要看它把哪两个计数方式连在一起。组合学证明最常见的技巧是：左边按一种标准数，右边按另一种标准数，然后说明两边数的是同一批对象。

\smallskip
\textbf{初学者检查点：}
\begin{itemize}[leftmargin=2em,itemsep=0.25em]
\item 先复述定理的假设，不要把适用范围扩大。
\item 找出证明中的基本计数原则：乘法、加法、双射、递推、容斥或生成函数。
\item 检查最后的等式化简是否和前面的对象解释一致。
\end{itemize}

\smallskip
\textbf{本章读法提醒：} 第九章分拆题先判断部件是否允许重复、是否要求不同、是否限制奇偶。
\end{explainblock}


\begin{sourceblock}{第 9 章续读}
{\small\sloppy
\noindent X\par
\noindent (1 − q j ) p(n)q n = 1.\par
\noindent j=1 n=0\par
\noindent 故若欧拉五角数定理成立, 我们就可以得到 P (n) 的递推关系\par
\noindent p(n) = p(n − 1) + p(n − 2) − p(n − 5) − p(n − 7) + · · · . (9.1)\par
}
\end{sourceblock}


\begin{explainblock}
这一段是原书论述链的一部分。读的时候把它当作桥梁：它通常负责把上一个定义、例子或定理，引到下一个计数方法。

\smallskip
\textbf{初学者检查点：}
\begin{itemize}[leftmargin=2em,itemsep=0.25em]
\item 找出这一段承接的上一个概念。
\item 找出它准备引出的下一个概念。
\item 把中间的关键等式或构造单独抄出来。
\end{itemize}

\smallskip
\textbf{本章读法提醒：} 第九章分拆题先判断部件是否允许重复、是否要求不同、是否限制奇偶。
\end{explainblock}


\begin{sourceblock}{例如我们有}
{\small\sloppy
\noindent 例如我们有\par
\noindent p(0) = 1,\par
\noindent p(1) = 1,\par
\noindent p(2) = p(1) + p(0) = 2,\par
\noindent p(3) = p(2) + p(1) = 3,\par
\noindent p(4) = p(3) + p(2) = 5,\par
\noindent p(5) = p(4) + p(3) − p(0) = 7,\par
\noindent p(6) = p(5) + p(4) − p(1) − p(0) = 10.\par
\noindent 接下来我们给出欧拉五角数定理的一个组合证明.\par
\noindent 证明: 记 D 为所有部分所含元素个数均不相同的分拆构成的集合. 我们已经\par
\noindent 证明过如下等式\par
\noindent ∞\par
\noindent Y X\par
\noindent (1 + xq j ) = xl(λ) q |λ| .\par
\noindent j=1 λ∈D\par
\noindent 于是我们有\par
\noindent ∞\par
\noindent Y X ∞\par
}
\end{sourceblock}


\begin{explainblock}
例题是学习这本书最重要的部分。不要只看答案，要把每个限制条件变成一个计数动作：先安排谁、再插入谁、有没有重复、需不需要除以对称性。

\smallskip
\textbf{初学者检查点：}
\begin{itemize}[leftmargin=2em,itemsep=0.25em]
\item 先写出没有限制时的总数。
\item 逐条加入限制条件，观察总数怎样变化。
\item 最后检查有没有把同一种方案重复算了多次。
\end{itemize}

\smallskip
\textbf{本章读法提醒：} 第九章分拆题先判断部件是否允许重复、是否要求不同、是否限制奇偶。
\end{explainblock}


\begin{sourceblock}{第 9 章续读}
{\small\sloppy
\noindent X X\par
\noindent (1 − q j ) = (−1)l(λ) q |λ| = qn (−1)l(λ) .\par
\noindent j=1 λ∈D n=0 λ∈D(n)\par
\noindent 为证明欧拉五角数定理, 只需证明\par
\noindent (\par
\noindent X (−1)j n = j(3j±1) ,\par
\noindent l(λ) 2\par
\noindent (−1) =\par
\noindent λ∈D(n)\par
\noindent 我们用组合方法来证明上式. 我们接下来给出的映射也被称为 Franklin 对合 F , 所\par
\noindent 谓的对合指的是若 λ, F (λ) 均在 F 的定义域里, 则 F 2 (λ) = λ.\par
\noindent 对于一个 D 中的排列 λ = \{λ1 , λ2 , . . . , λn \}, λi > λi + 1, 我们记 Dλ 为连续下\par
\noindent 降的部分数, 即 Dλ 为最小的满足 λi−1 − λi = 1, ∀i ≤ r 但 λr − λr+1 > i 的 r. 记\par
\noindent Sλ = λn 为最小的部分数. 我们对 λ 的 Ferrers 图做如下变换:\par
\noindent a. 若 Dλ < Sλ , 将前 Dλ 行的最后一个元素挪至最后一行下面构成新的一行.\par
\noindent b. 若 Dλ ≥ Sλ , 将最后一行的元素依次挪至前 Sλ 行的末尾.\par
}
\end{sourceblock}


\begin{explainblock}
这一段是原书论述链的一部分。读的时候把它当作桥梁：它通常负责把上一个定义、例子或定理，引到下一个计数方法。

\smallskip
\textbf{初学者检查点：}
\begin{itemize}[leftmargin=2em,itemsep=0.25em]
\item 找出这一段承接的上一个概念。
\item 找出它准备引出的下一个概念。
\item 把中间的关键等式或构造单独抄出来。
\end{itemize}

\smallskip
\textbf{本章读法提醒：} 第九章分拆题先判断部件是否允许重复、是否要求不同、是否限制奇偶。
\end{explainblock}


\begin{sourceblock}{例如对于 λ = \{8, 7, 5, 3\}, 变换如下}
{\small\sloppy
\noindent 例如对于 λ = \{8, 7, 5, 3\}, 变换如下\par
\noindent 这样的变换可以定义映射 F , F (λ) = µ = \{7, 6, 5, 3, 2\}. 而反之也有 F (µ) = λ.\par
\noindent 事实上可以发现 F 是定义在 D 上除去以下两类分拆的的一个对合.\par
\noindent a: Dλ = Sλ = l(λ) 如下图所示\par
\noindent .\par
\noindent 此时 F 没有定义. 容易验证此时 n = |λ| = j(3j−1)\par
\noindent b: Dλ = Sλ − 1 = l(λ) 如下图所示\par
\noindent .\par
\noindent 此时 F (λ) 不属于 D. 容易验证此时 n = |λ| = j(3j+1)\par
\noindent 又 F 作用在不为以上两种情况的 λ 时使 l(λ) 改变了 1, 从而改变了其奇偶\par
\noindent P\par
\noindent 性, 这意味着当 n 6= j(3j±1) 2 时 (−1)l(λ) = 0. n = j(3j±1)\par
\noindent λ∈D(n)\par
\noindent P\par
\noindent (−1)l(λ) = (−1)j .\par
\noindent λ∈D(n)\par
}
\end{sourceblock}


\begin{explainblock}
例题是学习这本书最重要的部分。不要只看答案，要把每个限制条件变成一个计数动作：先安排谁、再插入谁、有没有重复、需不需要除以对称性。

\smallskip
\textbf{初学者检查点：}
\begin{itemize}[leftmargin=2em,itemsep=0.25em]
\item 先写出没有限制时的总数。
\item 逐条加入限制条件，观察总数怎样变化。
\item 最后检查有没有把同一种方案重复算了多次。
\end{itemize}

\smallskip
\textbf{本章读法提醒：} 第九章分拆题先判断部件是否允许重复、是否要求不同、是否限制奇偶。
\end{explainblock}


\begin{sourceblock}{§ 9.5 两个欧拉恒等式}
{\small\sloppy
\noindent § 9.5. 两个欧拉恒等式\par
\noindent 这节我们再介绍两个欧拉恒等式, 它们常被称为欧拉 q-指数公式.\par
}
\end{sourceblock}


\begin{explainblock}
这一节先不要急着背公式。先问三个问题：对象是什么，哪些限制条件不能丢，最后到底是在数所有对象、还是在数某个统计量的分布。把这三件事说清楚，后面的公式才会有位置。

\smallskip
\textbf{初学者检查点：}
\begin{itemize}[leftmargin=2em,itemsep=0.25em]
\item 把本节出现的新对象写成一句话定义。
\item 把每个公式左边和右边分别翻译成“在数什么”。
\item 找一个最小非平凡例子，手工列出所有对象。
\end{itemize}

\smallskip
\textbf{本章读法提醒：} 第九章分拆题先判断部件是否允许重复、是否要求不同、是否限制奇偶。
\end{explainblock}


\begin{sourceblock}{定理 9.5.1 我们有如下两个公式}
{\small\sloppy
\noindent 定理 9.5.1 我们有如下两个公式\par
\noindent ∞\par
\noindent Y X ∞\par
\noindent = , (1)\par
\noindent 1 − zq j (1 − q)(1 − q 2 ) · · · (1 − q k )\par
\noindent j=1 k=0\par
\noindent ∞\par
\noindent Y ∞\par
\noindent X k+1\par
\noindent j zk q( 2 )\par
\noindent (1 − q)(1 − q 2 ) · · · (1 − q k )\par
\noindent j=1 k=0\par
\noindent 证明: 由定理 9.2.2 知, 取集合 S = \{s1 , s2 , · · · , sk \}< , 即表示 \{si \} 是严格递增\par
\noindent 的正整数序列. 设 P (S) 为每个部分都在 S 中的分拆所构成的集合, 则我们有\par
\noindent X 1\par
\noindent xl(λ) q |λ| = .\par
\noindent (1 − xq s1 ) · · · (1 − xq sk )\par
\noindent λ∈P (S)\par
}
\end{sourceblock}


\begin{explainblock}
定理段不要只看结论，要看它把哪两个计数方式连在一起。组合学证明最常见的技巧是：左边按一种标准数，右边按另一种标准数，然后说明两边数的是同一批对象。

\smallskip
\textbf{初学者检查点：}
\begin{itemize}[leftmargin=2em,itemsep=0.25em]
\item 先复述定理的假设，不要把适用范围扩大。
\item 找出证明中的基本计数原则：乘法、加法、双射、递推、容斥或生成函数。
\item 检查最后的等式化简是否和前面的对象解释一致。
\end{itemize}

\smallskip
\textbf{本章读法提醒：} 第九章分拆题先判断部件是否允许重复、是否要求不同、是否限制奇偶。
\end{explainblock}


\begin{sourceblock}{第 9 章续读}
{\small\sloppy
\noindent 记 n 的分拆个数为 p(n), n 的小于 k 个部分的分拆个数为 p≤k (n), n 的恰有 k 个部\par
\noindent 分的分拆个数为 p=k (n). 则由上式及其证明过程容易得到\par
\noindent ∞\par
\noindent X ∞\par
\noindent Y 1\par
\noindent p(n)xn = , (3)\par
\noindent n=0 i=1\par
\noindent ∞\par
\noindent X 1\par
\noindent p≤k (n)xn = (4),\par
\noindent (1 − q)(1 − q 2 ) · · · (1 − q k )\par
\noindent n=0\par
\noindent ∞\par
\noindent X qk\par
\noindent p=k (n)xn = (5).\par
\noindent (1 − q)(1 − q 2 ) · · · (1 − q k )\par
\noindent n=0\par
\noindent 我 们 先 来 证 明 (1) 式, 设 其 左 端 为 F (x, q), 右 端 为 G(x, q), F (x, q) 满 足 方 程\par
}
\end{sourceblock}


\begin{explainblock}
这一段是原书论述链的一部分。读的时候把它当作桥梁：它通常负责把上一个定义、例子或定理，引到下一个计数方法。

\smallskip
\textbf{初学者检查点：}
\begin{itemize}[leftmargin=2em,itemsep=0.25em]
\item 找出这一段承接的上一个概念。
\item 找出它准备引出的下一个概念。
\item 把中间的关键等式或构造单独抄出来。
\end{itemize}

\smallskip
\textbf{本章读法提醒：} 第九章分拆题先判断部件是否允许重复、是否要求不同、是否限制奇偶。
\end{explainblock}


\begin{sourceblock}{第 9 章续读}
{\small\sloppy
\noindent F (x, q) = F (xq, q)/(1 − xq) 及初值 F (x, 0) = F (0, q) = 1. 我们只需要证明 G(x, q)\par
\noindent 满足同样的方程及初值即可. 关于 G(x, q) 我们有\par
\noindent X xq 2k\par
\noindent G(xq, q) =\par
\noindent (1 − q)(1 − q 2 ) · · · (1 − q k )\par
\noindent k≥0\par
\noindent X\par
\noindent xq 2k 1\par
\noindent = −1\par
\noindent (1 − q)(1 − q 2 ) · · · (1 − q k−1 ) 1 − qk\par
\noindent k≥0\par
\noindent = (1 − xq)G(x, q).\par
\noindent 显然也有 G(x, 0) = G(0, q) = 1, 故 (1) 得证.\par
\noindent 接下来证明 (2). (2) 式左端的 xk q n 项的系数计数了将 n 的含有 k 个互不\par
\noindent 相同的部分的划分 λ = \{λ1 , λ2 , . . . , λk \}, λ1 > λ2 > · · · > λk 的个数. 于是\par
\noindent \{λ1 − k, λ2 − k + 1, . . . , λk − 1\} 是 n − k+1\par
\noindent 2 的不超过 k 个部分的划分. 因此由 (4)\par
\noindent 即可证明 (2).\par
}
\end{sourceblock}


\begin{explainblock}
这一段是原书论述链的一部分。读的时候把它当作桥梁：它通常负责把上一个定义、例子或定理，引到下一个计数方法。

\smallskip
\textbf{初学者检查点：}
\begin{itemize}[leftmargin=2em,itemsep=0.25em]
\item 找出这一段承接的上一个概念。
\item 找出它准备引出的下一个概念。
\item 把中间的关键等式或构造单独抄出来。
\end{itemize}

\smallskip
\textbf{本章读法提醒：} 第九章分拆题先判断部件是否允许重复、是否要求不同、是否限制奇偶。
\end{explainblock}


\begin{sourceblock}{习 题 九}
{\small\sloppy
\noindent 习 题 九\par
}
\end{sourceblock}


\begin{explainblock}
习题段不要先追求技巧。先给题目分类：它是在问排列、组合、递推、生成函数、容斥，还是结构计数。分类正确，工具通常就只剩一两个候选。

\smallskip
\textbf{初学者检查点：}
\begin{itemize}[leftmargin=2em,itemsep=0.25em]
\item 判断题目所属工具。
\item 列出变量和限制条件。
\item 先做小规模 n=2、3、4 的例子，确认直觉。
\end{itemize}

\smallskip
\textbf{本章读法提醒：} 第九章分拆题先判断部件是否允许重复、是否要求不同、是否限制奇偶。
\end{explainblock}


\begin{sourceblock}{习题 9.1 设 n 为正整数, 证明: n 的只有奇数部分可重复的分拆个数, 等于 n}
{\small\sloppy
\noindent 习题 9.1 设 n 为正整数, 证明: n 的只有奇数部分可重复的分拆个数, 等于 n\par
\noindent 的没有部分重复次数严格超过 3 次的分拆个数.\par
}
\end{sourceblock}


\begin{explainblock}
习题段不要先追求技巧。先给题目分类：它是在问排列、组合、递推、生成函数、容斥，还是结构计数。分类正确，工具通常就只剩一两个候选。

\smallskip
\textbf{初学者检查点：}
\begin{itemize}[leftmargin=2em,itemsep=0.25em]
\item 判断题目所属工具。
\item 列出变量和限制条件。
\item 先做小规模 n=2、3、4 的例子，确认直觉。
\end{itemize}

\smallskip
\textbf{本章读法提醒：} 第九章分拆题先判断部件是否允许重复、是否要求不同、是否限制奇偶。
\end{explainblock}


\begin{sourceblock}{习题 9.2 证明: n 的自共轭分拆的个数等于 n 的各部分都是奇数且两两不同}
{\small\sloppy
\noindent 习题 9.2 证明: n 的自共轭分拆的个数等于 n 的各部分都是奇数且两两不同\par
\noindent 的分拆的个数.\par
}
\end{sourceblock}


\begin{explainblock}
习题段不要先追求技巧。先给题目分类：它是在问排列、组合、递推、生成函数、容斥，还是结构计数。分类正确，工具通常就只剩一两个候选。

\smallskip
\textbf{初学者检查点：}
\begin{itemize}[leftmargin=2em,itemsep=0.25em]
\item 判断题目所属工具。
\item 列出变量和限制条件。
\item 先做小规模 n=2、3、4 的例子，确认直觉。
\end{itemize}

\smallskip
\textbf{本章读法提醒：} 第九章分拆题先判断部件是否允许重复、是否要求不同、是否限制奇偶。
\end{explainblock}


\begin{sourceblock}{习题 9.3 记 Am (n) 为 n 的仅有不被 2m 整除的元素可以在部分中重复出现的}
{\small\sloppy
\noindent 习题 9.3 记 Am (n) 为 n 的仅有不被 2m 整除的元素可以在部分中重复出现的\par
\noindent 划分构成的集合. 记 Bm (n) 为 n 的每一部分出现次数均不超过 2m+1 − 1 的划分构\par
\noindent 成的集合. 如对于 n = 4 我们有\par
\noindent P (4) = \{\{4\}, \{3, 1\}, \{2, 2\}, \{2, 1, 1\}, \{1, 1, 1, 1\}\},\par
\noindent A1 (4) = \{\{4\}, \{3, 1\}, \{2, 1, 1\}, \{1, 1, 1, 1\}\},\par
\noindent B1 (4) = \{\{4\}, \{3, 1\}, \{2, 2\}, \{2, 1, 1\}\},\par
\noindent 其中 \{2, 2\} 因为多次出现了被 2 整除的元素 2 而不在 A1 (4) 中, \{1, 1, 1, 1\} 含有\par
\noindent 4 > 3 个 1, 故不在 B1 (4) 中.\par
\noindent 求证: 对任意的 m, n, \#Am (n) = \#Bm (n).\par
}
\end{sourceblock}


\begin{explainblock}
习题段不要先追求技巧。先给题目分类：它是在问排列、组合、递推、生成函数、容斥，还是结构计数。分类正确，工具通常就只剩一两个候选。

\smallskip
\textbf{初学者检查点：}
\begin{itemize}[leftmargin=2em,itemsep=0.25em]
\item 判断题目所属工具。
\item 列出变量和限制条件。
\item 先做小规模 n=2、3、4 的例子，确认直觉。
\end{itemize}

\smallskip
\textbf{本章读法提醒：} 第九章分拆题先判断部件是否允许重复、是否要求不同、是否限制奇偶。
\end{explainblock}


\begin{sourceblock}{习题 9.4 设 n, k 都是正整数, f (n, k) 定义如下: 对将 n 写成恰好 k 个非负整}
{\small\sloppy
\noindent 习题 9.4 设 n, k 都是正整数, f (n, k) 定义如下: 对将 n 写成恰好 k 个非负整\par
\noindent 数的有序和的每种方法, 设 s 是这 k 个整数之积, 则 f (n, k) 是用这种方法得到的所\par
\noindent 有 s 这和. 求 f (n, k) 的普通型生成函数和 f (n, k) 的表达式.\par
}
\end{sourceblock}


\begin{explainblock}
习题段不要先追求技巧。先给题目分类：它是在问排列、组合、递推、生成函数、容斥，还是结构计数。分类正确，工具通常就只剩一两个候选。

\smallskip
\textbf{初学者检查点：}
\begin{itemize}[leftmargin=2em,itemsep=0.25em]
\item 判断题目所属工具。
\item 列出变量和限制条件。
\item 先做小规模 n=2、3、4 的例子，确认直觉。
\end{itemize}

\smallskip
\textbf{本章读法提醒：} 第九章分拆题先判断部件是否允许重复、是否要求不同、是否限制奇偶。
\end{explainblock}


\begin{sourceblock}{习题 9.5}
{\small\sloppy
\noindent 习题 9.5\par
\noindent 设 n, k, h 都是非负整数, f (n, k, h) 是将 n 表示为 k 个有序整数之和\par
\noindent P\par
\noindent 且每个整数都不小于 h 的方法数, 求 n≥0 f (n, k, h)xn .\par
}
\end{sourceblock}


\begin{explainblock}
习题段不要先追求技巧。先给题目分类：它是在问排列、组合、递推、生成函数、容斥，还是结构计数。分类正确，工具通常就只剩一两个候选。

\smallskip
\textbf{初学者检查点：}
\begin{itemize}[leftmargin=2em,itemsep=0.25em]
\item 判断题目所属工具。
\item 列出变量和限制条件。
\item 先做小规模 n=2、3、4 的例子，确认直觉。
\end{itemize}

\smallskip
\textbf{本章读法提醒：} 第九章分拆题先判断部件是否允许重复、是否要求不同、是否限制奇偶。
\end{explainblock}


\begin{sourceblock}{习题 9.6 设 f (n) 记为将 n 表示成偶数个正整数之和的方法数与将 n 表示成}
{\small\sloppy
\noindent 习题 9.6 设 f (n) 记为将 n 表示成偶数个正整数之和的方法数与将 n 表示成\par
\noindent 奇数个正整数之和的方法数之差, 求 f (n) 的普通型生成函数并求 f (n).\par
}
\end{sourceblock}


\begin{explainblock}
习题段不要先追求技巧。先给题目分类：它是在问排列、组合、递推、生成函数、容斥，还是结构计数。分类正确，工具通常就只剩一两个候选。

\smallskip
\textbf{初学者检查点：}
\begin{itemize}[leftmargin=2em,itemsep=0.25em]
\item 判断题目所属工具。
\item 列出变量和限制条件。
\item 先做小规模 n=2、3、4 的例子，确认直觉。
\end{itemize}

\smallskip
\textbf{本章读法提醒：} 第九章分拆题先判断部件是否允许重复、是否要求不同、是否限制奇偶。
\end{explainblock}


\begin{sourceblock}{习题 9.7 设 p(n, k) 为将正整数 n 分为 k 个部分的分拆的个数, 证明:}
{\small\sloppy
\noindent 习题 9.7 设 p(n, k) 为将正整数 n 分为 k 个部分的分拆的个数, 证明:\par
\noindent (1) p(n, k) = p(n − 1, k − 1) + p(n − k, k);\par
\noindent (2) p(n + k, k) = p(n, 1) + p(n, 2) + · · · + p(n, k);\par
\noindent (3) p(2n, n) = p(n).\par
}
\end{sourceblock}


\begin{explainblock}
习题段不要先追求技巧。先给题目分类：它是在问排列、组合、递推、生成函数、容斥，还是结构计数。分类正确，工具通常就只剩一两个候选。

\smallskip
\textbf{初学者检查点：}
\begin{itemize}[leftmargin=2em,itemsep=0.25em]
\item 判断题目所属工具。
\item 列出变量和限制条件。
\item 先做小规模 n=2、3、4 的例子，确认直觉。
\end{itemize}

\smallskip
\textbf{本章读法提醒：} 第九章分拆题先判断部件是否允许重复、是否要求不同、是否限制奇偶。
\end{explainblock}


\begin{sourceblock}{习题 9.8 用 pt (n) 来表示 n 拆分成 2 的不同次幂的方法数, 证明}
{\small\sloppy
\noindent 习题 9.8 用 pt (n) 来表示 n 拆分成 2 的不同次幂的方法数, 证明\par
\noindent pt (n) = 1.\par
\noindent Pk\par
}
\end{sourceblock}


\begin{explainblock}
习题段不要先追求技巧。先给题目分类：它是在问排列、组合、递推、生成函数、容斥，还是结构计数。分类正确，工具通常就只剩一两个候选。

\smallskip
\textbf{初学者检查点：}
\begin{itemize}[leftmargin=2em,itemsep=0.25em]
\item 判断题目所属工具。
\item 列出变量和限制条件。
\item 先做小规模 n=2、3、4 的例子，确认直觉。
\end{itemize}

\smallskip
\textbf{本章读法提醒：} 第九章分拆题先判断部件是否允许重复、是否要求不同、是否限制奇偶。
\end{explainblock}


\begin{sourceblock}{习题 9.9 设 p1 , p2 , . . . , pk 是任意给定的正整数, i=1 pi ≤ n, 求正整数 n 的 k}
{\small\sloppy
\noindent 习题 9.9 设 p1 , p2 , . . . , pk 是任意给定的正整数, i=1 pi ≤ n, 求正整数 n 的 k\par
\noindent 部有序分拆中的第 i 个部分 ≥ pi (i = 1, 2, . . . , k) 的分拆个数是多少?\par
}
\end{sourceblock}


\begin{explainblock}
习题段不要先追求技巧。先给题目分类：它是在问排列、组合、递推、生成函数、容斥，还是结构计数。分类正确，工具通常就只剩一两个候选。

\smallskip
\textbf{初学者检查点：}
\begin{itemize}[leftmargin=2em,itemsep=0.25em]
\item 判断题目所属工具。
\item 列出变量和限制条件。
\item 先做小规模 n=2、3、4 的例子，确认直觉。
\end{itemize}

\smallskip
\textbf{本章读法提醒：} 第九章分拆题先判断部件是否允许重复、是否要求不同、是否限制奇偶。
\end{explainblock}


\begin{sourceblock}{习题 9.10 证明: n 的各部分都大于 1 的分拆的个数等于 p(n) − p(n − 1), 并据}
{\small\sloppy
\noindent 习题 9.10 证明: n 的各部分都大于 1 的分拆的个数等于 p(n) − p(n − 1), 并据\par
\noindent 此证明\par
\noindent p(n + 2) + p(n) ≥ 2p(n + 1).\par
}
\end{sourceblock}


\begin{explainblock}
习题段不要先追求技巧。先给题目分类：它是在问排列、组合、递推、生成函数、容斥，还是结构计数。分类正确，工具通常就只剩一两个候选。

\smallskip
\textbf{初学者检查点：}
\begin{itemize}[leftmargin=2em,itemsep=0.25em]
\item 判断题目所属工具。
\item 列出变量和限制条件。
\item 先做小规模 n=2、3、4 的例子，确认直觉。
\end{itemize}

\smallskip
\textbf{本章读法提醒：} 第九章分拆题先判断部件是否允许重复、是否要求不同、是否限制奇偶。
\end{explainblock}


\begin{sourceblock}{习题 9.11 设 n 和 k 都是正整数, 以 P ≤k (n) 表示无一部分大于 k 的 n 的分拆}
{\small\sloppy
\noindent 习题 9.11 设 n 和 k 都是正整数, 以 P ≤k (n) 表示无一部分大于 k 的 n 的分拆\par
\noindent 的个数, 并令 P ≤k (0) = 1, 证明\par
\noindent P ≤k (n) = Pk (n + k).\par
}
\end{sourceblock}


\begin{explainblock}
习题段不要先追求技巧。先给题目分类：它是在问排列、组合、递推、生成函数、容斥，还是结构计数。分类正确，工具通常就只剩一两个候选。

\smallskip
\textbf{初学者检查点：}
\begin{itemize}[leftmargin=2em,itemsep=0.25em]
\item 判断题目所属工具。
\item 列出变量和限制条件。
\item 先做小规模 n=2、3、4 的例子，确认直觉。
\end{itemize}

\smallskip
\textbf{本章读法提醒：} 第九章分拆题先判断部件是否允许重复、是否要求不同、是否限制奇偶。
\end{explainblock}


\begin{sourceblock}{习题 9.12 设 n 和 k 都是正整数, 以 P≤k (n) 表示至多只有 k 个部分的 n 的分}
{\small\sloppy
\noindent 习题 9.12 设 n 和 k 都是正整数, 以 P≤k (n) 表示至多只有 k 个部分的 n 的分\par
\noindent 拆的个数, 并令 P≤k (0) = 1, 证明\par
\noindent P≤k (n) = Pk (n + k).\par
}
\end{sourceblock}


\begin{explainblock}
习题段不要先追求技巧。先给题目分类：它是在问排列、组合、递推、生成函数、容斥，还是结构计数。分类正确，工具通常就只剩一两个候选。

\smallskip
\textbf{初学者检查点：}
\begin{itemize}[leftmargin=2em,itemsep=0.25em]
\item 判断题目所属工具。
\item 列出变量和限制条件。
\item 先做小规模 n=2、3、4 的例子，确认直觉。
\end{itemize}

\smallskip
\textbf{本章读法提醒：} 第九章分拆题先判断部件是否允许重复、是否要求不同、是否限制奇偶。
\end{explainblock}


\begin{sourceblock}{习题 9.13 一个整数的有序分拆是关心部分顺序的分拆, 例如把 4 分成 1 和}
{\small\sloppy
\noindent 习题 9.13 一个整数的有序分拆是关心部分顺序的分拆, 例如把 4 分成 1 和\par
\noindent 2 的有序分拆有 2+2, 2+1+1, 1+2+1, 1+1+2, 1+1+1+1 这五个. 证明对任意的\par
\noindent n ≥ 1, n 分成 1 和 2 的有序分拆的个数为 Fibonacci 数 Fn+1 .\par
}
\end{sourceblock}


\begin{explainblock}
习题段不要先追求技巧。先给题目分类：它是在问排列、组合、递推、生成函数、容斥，还是结构计数。分类正确，工具通常就只剩一两个候选。

\smallskip
\textbf{初学者检查点：}
\begin{itemize}[leftmargin=2em,itemsep=0.25em]
\item 判断题目所属工具。
\item 列出变量和限制条件。
\item 先做小规模 n=2、3、4 的例子，确认直觉。
\end{itemize}

\smallskip
\textbf{本章读法提醒：} 第九章分拆题先判断部件是否允许重复、是否要求不同、是否限制奇偶。
\end{explainblock}


\begin{sourceblock}{习题 9.14 令 c(k, m, n) 表示恰好分成 m 个部分且每个部分都不超过 k 的关}
{\small\sloppy
\noindent 习题 9.14 令 c(k, m, n) 表示恰好分成 m 个部分且每个部分都不超过 k 的关\par
\noindent 于 n 的有序分拆的个数, 固定 k, m, 求:\par
\noindent (1)c(n, m, n),\par
\noindent P∞ n.\par
\noindent (2)c(k, m, n) 的生成函数 n=0 c(k, m, n)x\par
}
\end{sourceblock}


\begin{explainblock}
习题段不要先追求技巧。先给题目分类：它是在问排列、组合、递推、生成函数、容斥，还是结构计数。分类正确，工具通常就只剩一两个候选。

\smallskip
\textbf{初学者检查点：}
\begin{itemize}[leftmargin=2em,itemsep=0.25em]
\item 判断题目所属工具。
\item 列出变量和限制条件。
\item 先做小规模 n=2、3、4 的例子，确认直觉。
\end{itemize}

\smallskip
\textbf{本章读法提醒：} 第九章分拆题先判断部件是否允许重复、是否要求不同、是否限制奇偶。
\end{explainblock}


\begin{sourceblock}{习题 9.15 证明 n 分成没有部分出现恰好一次的分拆数量等于 n 分成模 6 不}
{\small\sloppy
\noindent 习题 9.15 证明 n 分成没有部分出现恰好一次的分拆数量等于 n 分成模 6 不\par
\noindent 同余于 ±1 的部分的分拆数量.\par
}
\end{sourceblock}


\begin{explainblock}
习题段不要先追求技巧。先给题目分类：它是在问排列、组合、递推、生成函数、容斥，还是结构计数。分类正确，工具通常就只剩一两个候选。

\smallskip
\textbf{初学者检查点：}
\begin{itemize}[leftmargin=2em,itemsep=0.25em]
\item 判断题目所属工具。
\item 列出变量和限制条件。
\item 先做小规模 n=2、3、4 的例子，确认直觉。
\end{itemize}

\smallskip
\textbf{本章读法提醒：} 第九章分拆题先判断部件是否允许重复、是否要求不同、是否限制奇偶。
\end{explainblock}


\begin{sourceblock}{习题 9.16 令 n 的分拆为 n = b1 + b2 + · · · + bs , bi ≥ bi+1 的形式, 令 Ar (n) 为}
{\small\sloppy
\noindent 习题 9.16 令 n 的分拆为 n = b1 + b2 + · · · + bs , bi ≥ bi+1 的形式, 令 Ar (n) 为\par
\noindent 所有奇数部分都 ≥ 2r + 1, 而且若 bi − bi+1 为奇数则 bi − bi+1 ≥ 2r + 1 的分拆的个\par
\noindent 数. 令 Br (n) 为每个部分要么是偶数, 要么模 4r + 2 余 2r + 1 的 n 的分拆的个数.\par
\noindent 则 Ar (n) = Br (n).\par
}
\end{sourceblock}


\begin{explainblock}
习题段不要先追求技巧。先给题目分类：它是在问排列、组合、递推、生成函数、容斥，还是结构计数。分类正确，工具通常就只剩一两个候选。

\smallskip
\textbf{初学者检查点：}
\begin{itemize}[leftmargin=2em,itemsep=0.25em]
\item 判断题目所属工具。
\item 列出变量和限制条件。
\item 先做小规模 n=2、3、4 的例子，确认直觉。
\end{itemize}

\smallskip
\textbf{本章读法提醒：} 第九章分拆题先判断部件是否允许重复、是否要求不同、是否限制奇偶。
\end{explainblock}


\begin{sourceblock}{习题 9.17 证 n 明分成每个部分恰好出现 2, 3, 5 次的分拆数量等于 n 分成模}
{\small\sloppy
\noindent 习题 9.17 证 n 明分成每个部分恰好出现 2, 3, 5 次的分拆数量等于 n 分成模\par
\noindent 12 同余于 ±2, ±3, 6 的部分的分拆数量.\par
}
\end{sourceblock}


\begin{explainblock}
习题段不要先追求技巧。先给题目分类：它是在问排列、组合、递推、生成函数、容斥，还是结构计数。分类正确，工具通常就只剩一两个候选。

\smallskip
\textbf{初学者检查点：}
\begin{itemize}[leftmargin=2em,itemsep=0.25em]
\item 判断题目所属工具。
\item 列出变量和限制条件。
\item 先做小规模 n=2、3、4 的例子，确认直觉。
\end{itemize}

\smallskip
\textbf{本章读法提醒：} 第九章分拆题先判断部件是否允许重复、是否要求不同、是否限制奇偶。
\end{explainblock}


\begin{sourceblock}{习题 9.18 令 H 为正整数集合, 用 ‘‘H” 表示部分在 H 中的所有分拆的集合,}
{\small\sloppy
\noindent 习题 9.18 令 H 为正整数集合, 用 ‘‘H” 表示部分在 H 中的所有分拆的集合,\par
\noindent 因此 p(‘‘H”, n) 是其部分都属于 H 的 n 的分拆数. 例如若 H0 为所有正奇数的集\par
\noindent 合, 则 ‘‘H0 ” = O , p(‘‘H0 ”, n) = p(O, n) . 令\par
\noindent X\par
\noindent f (q) = p(‘‘H”, n)q n , (9.2)\par
\noindent n≥0\par
\noindent X\par
\noindent fd (q) = p(‘‘H”(≤ d), n)q n . (9.3)\par
\noindent n≥0\par
\noindent 证明对于 |q| < 1 , 有\par
\noindent Y\par
\noindent f (q) = (1 − q n )−1 , (9.4)\par
\noindent n∈H\par
\noindent Y\par
\noindent fd (q) = (1 + q n + . . . + q dn )\par
\noindent n∈H\par
\noindent Y\par
\noindent = (1 − q (d+1)n )(1 − q n )−1 . (9.5)\par
}
\end{sourceblock}


\begin{explainblock}
习题段不要先追求技巧。先给题目分类：它是在问排列、组合、递推、生成函数、容斥，还是结构计数。分类正确，工具通常就只剩一两个候选。

\smallskip
\textbf{初学者检查点：}
\begin{itemize}[leftmargin=2em,itemsep=0.25em]
\item 判断题目所属工具。
\item 列出变量和限制条件。
\item 先做小规模 n=2、3、4 的例子，确认直觉。
\end{itemize}

\smallskip
\textbf{本章读法提醒：} 第九章分拆题先判断部件是否允许重复、是否要求不同、是否限制奇偶。
\end{explainblock}


\begin{sourceblock}{第 9 章续读}
{\small\sloppy
\noindent n∈H\par
}
\end{sourceblock}


\begin{explainblock}
这一段是原书论述链的一部分。读的时候把它当作桥梁：它通常负责把上一个定义、例子或定理，引到下一个计数方法。

\smallskip
\textbf{初学者检查点：}
\begin{itemize}[leftmargin=2em,itemsep=0.25em]
\item 找出这一段承接的上一个概念。
\item 找出它准备引出的下一个概念。
\item 把中间的关键等式或构造单独抄出来。
\end{itemize}

\smallskip
\textbf{本章读法提醒：} 第九章分拆题先判断部件是否允许重复、是否要求不同、是否限制奇偶。
\end{explainblock}


\begin{sourceblock}{习题 9.19 用 ‘‘H”(≤ d) 表示部分在 H 中且没有任何一部分出现超过 d 次的}
{\small\sloppy
\noindent 习题 9.19 用 ‘‘H”(≤ d) 表示部分在 H 中且没有任何一部分出现超过 d 次的\par
\noindent 所有分拆的集合. 例如若 N 表示所有正整数的集合, 则 p(‘‘N ”(≤ 1), n) = p(D, n).\par
\noindent Nd 记为不能被 d 整除的正整数的集合, 则对于任意 n,\par
\noindent p(‘‘Nd+1 ”, n) = p(‘‘N ”(≤ d), n).\par
}
\end{sourceblock}


\begin{explainblock}
习题段不要先追求技巧。先给题目分类：它是在问排列、组合、递推、生成函数、容斥，还是结构计数。分类正确，工具通常就只剩一两个候选。

\smallskip
\textbf{初学者检查点：}
\begin{itemize}[leftmargin=2em,itemsep=0.25em]
\item 判断题目所属工具。
\item 列出变量和限制条件。
\item 先做小规模 n=2、3、4 的例子，确认直觉。
\end{itemize}

\smallskip
\textbf{本章读法提醒：} 第九章分拆题先判断部件是否允许重复、是否要求不同、是否限制奇偶。
\end{explainblock}


\chapter{图与树的计数}
\begin{roadmapblock}
本章保留原书正文的主要数学骨架，并在每个定义、定理、例题和论述片段后加入解读。
阅读时不要把目标放在“背下答案”上，而是放在“看懂这一步为什么这样数”上。

\smallskip
\textbf{本章最低目标：} 能说出每个公式左边在数什么、右边在数什么，以及为什么两边相等。
\end{roadmapblock}


\begin{sourceblock}{第 10 章正文}
{\small\sloppy
\noindent 图与树的计数\par
\noindent 在 18 世纪的德国哥尼斯堡, 有一个岛位于两条河流的交汇处, 如下图七座桥把这\par
\noindent 几块岛连在了一起. 哥尼斯堡七桥 (the Seven Bridges of Königsberg) 问题就是要找\par
\noindent 到一条把每座桥都只走过一次的路, 一共有多少种这样的路呢? Euler 证明了不存在\par
\noindent 这样的路, 从而解决了问题. Euler 的这篇文章 “Solutio problematis ad geometriam\par
\noindent situs pertinentis” 是在 1741 年发表的, 但由于注明的日期是 1736 年, 我们一般认为\par
\noindent 图论是在 1736 年作为数学的一个分支诞生的.\par
}
\end{sourceblock}


\begin{explainblock}
这一段是原书论述链的一部分。读的时候把它当作桥梁：它通常负责把上一个定义、例子或定理，引到下一个计数方法。

\smallskip
\textbf{初学者检查点：}
\begin{itemize}[leftmargin=2em,itemsep=0.25em]
\item 找出这一段承接的上一个概念。
\item 找出它准备引出的下一个概念。
\item 把中间的关键等式或构造单独抄出来。
\end{itemize}

\smallskip
\textbf{本章读法提醒：} 第十章树计数常靠编码或递归分解；平面图公式要确认是否满足平面嵌入。
\end{explainblock}


\begin{sourceblock}{§ 10.1 图的基本概念}
{\small\sloppy
\noindent § 10.1. 图的基本概念\par
}
\end{sourceblock}


\begin{explainblock}
这一节先不要急着背公式。先问三个问题：对象是什么，哪些限制条件不能丢，最后到底是在数所有对象、还是在数某个统计量的分布。把这三件事说清楚，后面的公式才会有位置。

\smallskip
\textbf{初学者检查点：}
\begin{itemize}[leftmargin=2em,itemsep=0.25em]
\item 把本节出现的新对象写成一句话定义。
\item 把每个公式左边和右边分别翻译成“在数什么”。
\item 找一个最小非平凡例子，手工列出所有对象。
\end{itemize}

\smallskip
\textbf{本章读法提醒：} 第十章树计数常靠编码或递归分解；平面图公式要确认是否满足平面嵌入。
\end{explainblock}


\begin{sourceblock}{定义 10.1.1 一个图 G (Graph) 定义为一个有序对 (V, E), 记为 G = (V, E). 其中}
{\small\sloppy
\noindent 定义 10.1.1 一个图 G (Graph) 定义为一个有序对 (V, E), 记为 G = (V, E). 其中\par
\noindent V 是顶点集 (Vertex), E 由 V 中的顶点组成的无序点对组成, 其元素称为边, 称为边\par
\noindent 集. 用 |V | 表示顶点数, 用 |E| 表示边数. 边记为 uv, 也可记作 e, 即 e = uv, 此时称\par
\noindent u 和 v 是 e 端点, 并称 u 和 v 相邻, u (或 v) 与 e 相关. 顶点集和边集都有限的图,\par
\noindent 称为有限图, 只有一个顶点的图称为平凡图; 边集为空的图称为空图; 顶点数为 n 的\par
\noindent 图形称为 n 阶图; 顶点数为 n 边数为 m 的图称为 (n, m) 图; 连接两个相同顶点的\par
\noindent 边的条数称为边的重数; 重数大于 1 的边称为重边, 端点重合为一点的边称为环; 无\par
\noindent 环无重边的图称为简单图, 其余的图称为复合图.\par
}
\end{sourceblock}


\begin{explainblock}
定义段的任务是定边界。这里要特别注意参数范围、是否允许重复、是否考虑顺序、对象有没有标签。后面的每一道例题和证明，本质上都在反复使用这些边界。

\smallskip
\textbf{初学者检查点：}
\begin{itemize}[leftmargin=2em,itemsep=0.25em]
\item 圈出被定义的对象名称。
\item 圈出参数范围，例如 n、k 是否要求正整数。
\item 判断它是有序对象、无序对象、带标签对象还是结构对象。
\end{itemize}

\smallskip
\textbf{本章读法提醒：} 第十章树计数常靠编码或递归分解；平面图公式要确认是否满足平面嵌入。
\end{explainblock}


\begin{sourceblock}{注记 10.1.2 对任意一个图 G = (V, E), 若我们称 E 是一个集合, 则意味着 G 不}
{\small\sloppy
\noindent 注记 10.1.2 对任意一个图 G = (V, E), 若我们称 E 是一个集合, 则意味着 G 不\par
\noindent 含有重边; 所以一般来说, E 并不一定是一个集合.\par
\noindent 若用点代表顶点, 两点间的连线表示边, 则可将图用 “ 图形” 表示.\par
}
\end{sourceblock}


\begin{explainblock}
注记通常是在补充术语、记号、历史或一个容易忽略的约定。它未必给新公式，但常常决定你后面会不会把符号读错。

\smallskip
\textbf{初学者检查点：}
\begin{itemize}[leftmargin=2em,itemsep=0.25em]
\item 记下新符号的读法。
\item 确认这个约定是否只在本章使用。
\item 把注记里的限制条件补到前面的定义旁边。
\end{itemize}

\smallskip
\textbf{本章读法提醒：} 第十章树计数常靠编码或递归分解；平面图公式要确认是否满足平面嵌入。
\end{explainblock}


\begin{sourceblock}{例 10.1.1 设图 G = (V, E), 其中 V = \{v1 , v2 , v3 , v4 \}, E = \{e1 , e2 , e3 , e4 , e5 , e6 \},}
{\small\sloppy
\noindent 例 10.1.1 设图 G = (V, E), 其中 V = \{v1 , v2 , v3 , v4 \}, E = \{e1 , e2 , e3 , e4 , e5 , e6 \},\par
\noindent e1 = (v1 , v2 ), e2 = (v1 , v3 ), e3 = (v1 , v4 ), e4 = (v2 , v3 ), e5 = (v2 , v3 ), e6 = (v3 , v3 ), 它\par
\noindent 的图形表示如 10.1.\par
\noindent v1\par
\noindent v4\par
\noindent e3\par
\noindent e1 e2\par
\noindent e4\par
\noindent v2 v3 e6\par
\noindent e5\par
}
\end{sourceblock}


\begin{explainblock}
例题是学习这本书最重要的部分。不要只看答案，要把每个限制条件变成一个计数动作：先安排谁、再插入谁、有没有重复、需不需要除以对称性。

\smallskip
\textbf{初学者检查点：}
\begin{itemize}[leftmargin=2em,itemsep=0.25em]
\item 先写出没有限制时的总数。
\item 逐条加入限制条件，观察总数怎样变化。
\item 最后检查有没有把同一种方案重复算了多次。
\end{itemize}

\smallskip
\textbf{本章读法提醒：} 第十章树计数常靠编码或递归分解；平面图公式要确认是否满足平面嵌入。
\end{explainblock}


\begin{sourceblock}{例 10.1.2 若把哥尼斯堡的每一块土地看作一个顶点, 每一座桥看作一条边, 则}
{\small\sloppy
\noindent 例 10.1.2 若把哥尼斯堡的每一块土地看作一个顶点, 每一座桥看作一条边, 则\par
\noindent 哥尼斯堡七桥可用以下图形表示:\par
}
\end{sourceblock}


\begin{explainblock}
例题是学习这本书最重要的部分。不要只看答案，要把每个限制条件变成一个计数动作：先安排谁、再插入谁、有没有重复、需不需要除以对称性。

\smallskip
\textbf{初学者检查点：}
\begin{itemize}[leftmargin=2em,itemsep=0.25em]
\item 先写出没有限制时的总数。
\item 逐条加入限制条件，观察总数怎样变化。
\item 最后检查有没有把同一种方案重复算了多次。
\end{itemize}

\smallskip
\textbf{本章读法提醒：} 第十章树计数常靠编码或递归分解；平面图公式要确认是否满足平面嵌入。
\end{explainblock}


\begin{sourceblock}{定义 10.1.3 设图 G 与 H, 若 V (H) ⊆ V (G), E(H) ⊆ E(G), 则称 H 是 G 的子}
{\small\sloppy
\noindent 定义 10.1.3 设图 G 与 H, 若 V (H) ⊆ V (G), E(H) ⊆ E(G), 则称 H 是 G 的子\par
\noindent 图.\par
}
\end{sourceblock}


\begin{explainblock}
定义段的任务是定边界。这里要特别注意参数范围、是否允许重复、是否考虑顺序、对象有没有标签。后面的每一道例题和证明，本质上都在反复使用这些边界。

\smallskip
\textbf{初学者检查点：}
\begin{itemize}[leftmargin=2em,itemsep=0.25em]
\item 圈出被定义的对象名称。
\item 圈出参数范围，例如 n、k 是否要求正整数。
\item 判断它是有序对象、无序对象、带标签对象还是结构对象。
\end{itemize}

\smallskip
\textbf{本章读法提醒：} 第十章树计数常靠编码或递归分解；平面图公式要确认是否满足平面嵌入。
\end{explainblock}


\begin{sourceblock}{定义 10.1.4 设图 G, 记 dG(v) 为 G 中与顶点 v 相关联的边数, 每个环算作两条}
{\small\sloppy
\noindent 定义 10.1.4 设图 G, 记 dG(v) 为 G 中与顶点 v 相关联的边数, 每个环算作两条\par
\noindent 边, 称 dG(v) 为顶点 v 的度.\par
}
\end{sourceblock}


\begin{explainblock}
定义段的任务是定边界。这里要特别注意参数范围、是否允许重复、是否考虑顺序、对象有没有标签。后面的每一道例题和证明，本质上都在反复使用这些边界。

\smallskip
\textbf{初学者检查点：}
\begin{itemize}[leftmargin=2em,itemsep=0.25em]
\item 圈出被定义的对象名称。
\item 圈出参数范围，例如 n、k 是否要求正整数。
\item 判断它是有序对象、无序对象、带标签对象还是结构对象。
\end{itemize}

\smallskip
\textbf{本章读法提醒：} 第十章树计数常靠编码或递归分解；平面图公式要确认是否满足平面嵌入。
\end{explainblock}


\begin{sourceblock}{定理 10.1.5 若图 G 有 |E| 条边, 顶点分别为 v1 , v2 , . . ., 则}
{\small\sloppy
\noindent 定理 10.1.5 若图 G 有 |E| 条边, 顶点分别为 v1 , v2 , . . ., 则\par
\noindent d(v1 ) + d(v2 ) + · · · = 2|E|.\par
\noindent 证明: 考虑边的端点, 一个度数为 d(v) 的顶点也占了 d(v) 个边的端点, 则整个\par
\noindent 图的边的端点个数为\par
\noindent d(v1 ) + d(v2 ) + · · ·\par
\noindent 同时每个边有两个端点, 这就意味着图 G 中边的端点的总数为 2|E| .\par
}
\end{sourceblock}


\begin{explainblock}
定理段不要只看结论，要看它把哪两个计数方式连在一起。组合学证明最常见的技巧是：左边按一种标准数，右边按另一种标准数，然后说明两边数的是同一批对象。

\smallskip
\textbf{初学者检查点：}
\begin{itemize}[leftmargin=2em,itemsep=0.25em]
\item 先复述定理的假设，不要把适用范围扩大。
\item 找出证明中的基本计数原则：乘法、加法、双射、递推、容斥或生成函数。
\item 检查最后的等式化简是否和前面的对象解释一致。
\end{itemize}

\smallskip
\textbf{本章读法提醒：} 第十章树计数常靠编码或递归分解；平面图公式要确认是否满足平面嵌入。
\end{explainblock}


\begin{sourceblock}{定理 10.1.6 一个图里度数为奇数的顶点个数为偶数.}
{\small\sloppy
\noindent 定理 10.1.6 一个图里度数为奇数的顶点个数为偶数.\par
\noindent 证明: 由定理 10.1.5 知所有顶点度数的和为偶数, 而所有度为偶数的顶点的度\par
\noindent 的和也是偶数, 则所有度为奇数的顶点的度的和必然是偶数, 从而度数为奇数的顶点\par
\noindent 个数为偶数.\par
}
\end{sourceblock}


\begin{explainblock}
定理段不要只看结论，要看它把哪两个计数方式连在一起。组合学证明最常见的技巧是：左边按一种标准数，右边按另一种标准数，然后说明两边数的是同一批对象。

\smallskip
\textbf{初学者检查点：}
\begin{itemize}[leftmargin=2em,itemsep=0.25em]
\item 先复述定理的假设，不要把适用范围扩大。
\item 找出证明中的基本计数原则：乘法、加法、双射、递推、容斥或生成函数。
\item 检查最后的等式化简是否和前面的对象解释一致。
\end{itemize}

\smallskip
\textbf{本章读法提醒：} 第十章树计数常靠编码或递归分解；平面图公式要确认是否满足平面嵌入。
\end{explainblock}


\begin{sourceblock}{定义 10.1.7 称图 G 是正则的, 若对所有 v ∈ V (G), 其度数相同.}
{\small\sloppy
\noindent 定义 10.1.7 称图 G 是正则的, 若对所有 v ∈ V (G), 其度数相同.\par
}
\end{sourceblock}


\begin{explainblock}
定义段的任务是定边界。这里要特别注意参数范围、是否允许重复、是否考虑顺序、对象有没有标签。后面的每一道例题和证明，本质上都在反复使用这些边界。

\smallskip
\textbf{初学者检查点：}
\begin{itemize}[leftmargin=2em,itemsep=0.25em]
\item 圈出被定义的对象名称。
\item 圈出参数范围，例如 n、k 是否要求正整数。
\item 判断它是有序对象、无序对象、带标签对象还是结构对象。
\end{itemize}

\smallskip
\textbf{本章读法提醒：} 第十章树计数常靠编码或递归分解；平面图公式要确认是否满足平面嵌入。
\end{explainblock}


\begin{sourceblock}{定义 10.1.8 图 G 的一条路径 (或通道或通路) 是指一个有限非空序列 w =}
{\small\sloppy
\noindent 定义 10.1.8 图 G 的一条路径 (或通道或通路) 是指一个有限非空序列 w =\par
\noindent v0 e1 v1 e2 v2 · · · ek vk , 它的项交替地为顶点和边, 使得对 1 ≤ i ≤ k, ei 的端点是 vi−1\par
\noindent 与 vi , 称 w 是从顶点 v0 到 vk 的一条途径, 或一条 (v0 , vk ) 途径. 顶点 v0 和 vk 分\par
\noindent 别称为 w 的起点和终点, 而 v1 , v2 , · · · vk−1 称为它的内部顶点, 整数 k 称为 w 的长.\par
\noindent 若途径 w 的边 e1 , e2 , · · · ek , 互不相同, 则称 w 为迹, 这时 w 的长恰好等于边数. 又\par
\noindent 若途径 w 的顶点 v0 , v1 , v2 , · · · vk 也互不相同, 则称 w 为路.\par
}
\end{sourceblock}


\begin{explainblock}
定义段的任务是定边界。这里要特别注意参数范围、是否允许重复、是否考虑顺序、对象有没有标签。后面的每一道例题和证明，本质上都在反复使用这些边界。

\smallskip
\textbf{初学者检查点：}
\begin{itemize}[leftmargin=2em,itemsep=0.25em]
\item 圈出被定义的对象名称。
\item 圈出参数范围，例如 n、k 是否要求正整数。
\item 判断它是有序对象、无序对象、带标签对象还是结构对象。
\end{itemize}

\smallskip
\textbf{本章读法提醒：} 第十章树计数常靠编码或递归分解；平面图公式要确认是否满足平面嵌入。
\end{explainblock}


\begin{sourceblock}{定义 10.1.9 图 G 的两个顶点 u 和 v 称为是连通的, 如果在 G 中存在 (u, v) 路,}
{\small\sloppy
\noindent 定义 10.1.9 图 G 的两个顶点 u 和 v 称为是连通的, 如果在 G 中存在 (u, v) 路,\par
\noindent 规定 u 和 u 是连通的. 如果对 G 中每一对顶点 u, v 都有一条 (u, v) 路, 则称 G 为\par
\noindent 连通图, 否则称为非连通图.\par
}
\end{sourceblock}


\begin{explainblock}
定义段的任务是定边界。这里要特别注意参数范围、是否允许重复、是否考虑顺序、对象有没有标签。后面的每一道例题和证明，本质上都在反复使用这些边界。

\smallskip
\textbf{初学者检查点：}
\begin{itemize}[leftmargin=2em,itemsep=0.25em]
\item 圈出被定义的对象名称。
\item 圈出参数范围，例如 n、k 是否要求正整数。
\item 判断它是有序对象、无序对象、带标签对象还是结构对象。
\end{itemize}

\smallskip
\textbf{本章读法提醒：} 第十章树计数常靠编码或递归分解；平面图公式要确认是否满足平面嵌入。
\end{explainblock}


\begin{sourceblock}{定义 10.1.10 称一条途径是闭的, 如果它有正的长且起点和终点相同. 闭迹也称}
{\small\sloppy
\noindent 定义 10.1.10 称一条途径是闭的, 如果它有正的长且起点和终点相同. 闭迹也称\par
\noindent 为同路, 若一条闭连的起点和内部顶点互不相同则称它为圈, 长为 k 的圈称为 k 圈.\par
}
\end{sourceblock}


\begin{explainblock}
定义段的任务是定边界。这里要特别注意参数范围、是否允许重复、是否考虑顺序、对象有没有标签。后面的每一道例题和证明，本质上都在反复使用这些边界。

\smallskip
\textbf{初学者检查点：}
\begin{itemize}[leftmargin=2em,itemsep=0.25em]
\item 圈出被定义的对象名称。
\item 圈出参数范围，例如 n、k 是否要求正整数。
\item 判断它是有序对象、无序对象、带标签对象还是结构对象。
\end{itemize}

\smallskip
\textbf{本章读法提醒：} 第十章树计数常靠编码或递归分解；平面图公式要确认是否满足平面嵌入。
\end{explainblock}


\begin{sourceblock}{引理 10.1.11 若一个图 G 里所有顶点的度数都不小于 2 , 则 G 中有一个圈.}
{\small\sloppy
\noindent 引理 10.1.11 若一个图 G 里所有顶点的度数都不小于 2 , 则 G 中有一个圈.\par
\noindent 证明: 从任意一个顶点 u 开始, 走过连到 u 的某条边到达另一个顶点 u1 , 由于\par
\noindent u1 的度数不小于 2, 则可从 u1 走过另一条边到达顶点 u2 . 此过程可一直重复下去,\par
\noindent 除非走到了一个之前到过的点 v, 则从第一次到达 v 到第二次到达 v 的过程形成一\par
\noindent 个圈.\par
}
\end{sourceblock}


\begin{explainblock}
定理段不要只看结论，要看它把哪两个计数方式连在一起。组合学证明最常见的技巧是：左边按一种标准数，右边按另一种标准数，然后说明两边数的是同一批对象。

\smallskip
\textbf{初学者检查点：}
\begin{itemize}[leftmargin=2em,itemsep=0.25em]
\item 先复述定理的假设，不要把适用范围扩大。
\item 找出证明中的基本计数原则：乘法、加法、双射、递推、容斥或生成函数。
\item 检查最后的等式化简是否和前面的对象解释一致。
\end{itemize}

\smallskip
\textbf{本章读法提醒：} 第十章树计数常靠编码或递归分解；平面图公式要确认是否满足平面嵌入。
\end{explainblock}


\begin{sourceblock}{定义 10.1.12 若图 G 中没有圈, 则称 G 为森林 (forest), 若连通图 G 中没有圈,}
{\small\sloppy
\noindent 定义 10.1.12 若图 G 中没有圈, 则称 G 为森林 (forest), 若连通图 G 中没有圈,\par
\noindent 则称 G 为树 (tree).\par
}
\end{sourceblock}


\begin{explainblock}
定义段的任务是定边界。这里要特别注意参数范围、是否允许重复、是否考虑顺序、对象有没有标签。后面的每一道例题和证明，本质上都在反复使用这些边界。

\smallskip
\textbf{初学者检查点：}
\begin{itemize}[leftmargin=2em,itemsep=0.25em]
\item 圈出被定义的对象名称。
\item 圈出参数范围，例如 n、k 是否要求正整数。
\item 判断它是有序对象、无序对象、带标签对象还是结构对象。
\end{itemize}

\smallskip
\textbf{本章读法提醒：} 第十章树计数常靠编码或递归分解；平面图公式要确认是否满足平面嵌入。
\end{explainblock}


\begin{sourceblock}{例 10.1.3 如图 10.3, 有 6 个顶点的树有 6 个.}
{\small\sloppy
\noindent 例 10.1.3 如图 10.3, 有 6 个顶点的树有 6 个.\par
}
\end{sourceblock}


\begin{explainblock}
例题是学习这本书最重要的部分。不要只看答案，要把每个限制条件变成一个计数动作：先安排谁、再插入谁、有没有重复、需不需要除以对称性。

\smallskip
\textbf{初学者检查点：}
\begin{itemize}[leftmargin=2em,itemsep=0.25em]
\item 先写出没有限制时的总数。
\item 逐条加入限制条件，观察总数怎样变化。
\item 最后检查有没有把同一种方案重复算了多次。
\end{itemize}

\smallskip
\textbf{本章读法提醒：} 第十章树计数常靠编码或递归分解；平面图公式要确认是否满足平面嵌入。
\end{explainblock}


\begin{sourceblock}{定理 10.1.13 树中的任意两个点都能唯一地连成一条路.}
{\small\sloppy
\noindent 定理 10.1.13 树中的任意两个点都能唯一地连成一条路.\par
\noindent 证明: 利用反证法, 设树 G 的两个顶点 u, v 之间有两条不同的路 P1 和 P2 , 由\par
\noindent 于 P1 和 P2 不同, 设有一条边 e = u1 u2 属于 P1 但不属于 P2 . 则显然图 (P1 ∪ P2 ) − e\par
\noindent 是连通图, 从而包含一条从 u1 到 u2 的路 P , 则图 P 就是 G 中的一个圈, 矛盾.\par
}
\end{sourceblock}


\begin{explainblock}
定理段不要只看结论，要看它把哪两个计数方式连在一起。组合学证明最常见的技巧是：左边按一种标准数，右边按另一种标准数，然后说明两边数的是同一批对象。

\smallskip
\textbf{初学者检查点：}
\begin{itemize}[leftmargin=2em,itemsep=0.25em]
\item 先复述定理的假设，不要把适用范围扩大。
\item 找出证明中的基本计数原则：乘法、加法、双射、递推、容斥或生成函数。
\item 检查最后的等式化简是否和前面的对象解释一致。
\end{itemize}

\smallskip
\textbf{本章读法提醒：} 第十章树计数常靠编码或递归分解；平面图公式要确认是否满足平面嵌入。
\end{explainblock}


\begin{sourceblock}{定理 10.1.14 图 G 是树当且仅当 G 连通且 |E| = |V | − 1.}
{\small\sloppy
\noindent 定理 10.1.14 图 G 是树当且仅当 G 连通且 |E| = |V | − 1.\par
\noindent 证明: 下面我们对顶点 n 作归纳来证明边的个数为 n − 1. 当 n = 1 时, 显然成\par
\noindent 立. 假定对于顶点小于 n(n ≥ 2) 的树均成立. 设具有 n 个顶点的树 T 的边的集合\par
\noindent 为 E(T ), 取 e ∈ E(T ), 将边 e 去掉, 则得到了两个顶点数分别为 n1 , n2 (n1 , n2 ≥ 1)\par
\noindent 的树, 不妨设为 T1 , T2 , 则\par
\noindent n1 + n2 = n,\par
\noindent 所以\par
\noindent n1 < n, n2 < n,\par
\noindent 故由归纳法知\par
\noindent |E(T1 )| = n1 − 1, |E(T2 )| = n2 − 1,\par
\noindent 因此\par
\noindent |E(T )| = 1 + E(T1 ) + E(T2 ) = n1 + n2 − 1 = n − 1.\par
}
\end{sourceblock}


\begin{explainblock}
定理段不要只看结论，要看它把哪两个计数方式连在一起。组合学证明最常见的技巧是：左边按一种标准数，右边按另一种标准数，然后说明两边数的是同一批对象。

\smallskip
\textbf{初学者检查点：}
\begin{itemize}[leftmargin=2em,itemsep=0.25em]
\item 先复述定理的假设，不要把适用范围扩大。
\item 找出证明中的基本计数原则：乘法、加法、双射、递推、容斥或生成函数。
\item 检查最后的等式化简是否和前面的对象解释一致。
\end{itemize}

\smallskip
\textbf{本章读法提醒：} 第十章树计数常靠编码或递归分解；平面图公式要确认是否满足平面嵌入。
\end{explainblock}


\begin{sourceblock}{推论 10.1.15 任何非平凡的树都至少有两个度数为 1 的顶点.}
{\small\sloppy
\noindent 推论 10.1.15 任何非平凡的树都至少有两个度数为 1 的顶点.\par
\noindent 证明: 由定理 10.1.14 知非平凡的树的度数总和\par
\noindent X\par
\noindent d(v) = 2|E| = 2|V | − 2,\par
\noindent v∈V\par
\noindent 而非平凡的树的顶点的度至少为 1, 从而至少有两个度数为 1 的顶点.\par
\noindent 下面给出一个森林的应用.\par
}
\end{sourceblock}


\begin{explainblock}
定理段不要只看结论，要看它把哪两个计数方式连在一起。组合学证明最常见的技巧是：左边按一种标准数，右边按另一种标准数，然后说明两边数的是同一批对象。

\smallskip
\textbf{初学者检查点：}
\begin{itemize}[leftmargin=2em,itemsep=0.25em]
\item 先复述定理的假设，不要把适用范围扩大。
\item 找出证明中的基本计数原则：乘法、加法、双射、递推、容斥或生成函数。
\item 检查最后的等式化简是否和前面的对象解释一致。
\end{itemize}

\smallskip
\textbf{本章读法提醒：} 第十章树计数常靠编码或递归分解；平面图公式要确认是否满足平面嵌入。
\end{explainblock}


\begin{sourceblock}{例 10.1.4 (J. A. Bondy, 1972) 设 A = \{A1 , A2 , . . . , An \} 是 X = \{1, 2, . . . , n\}}
{\small\sloppy
\noindent 例 10.1.4 (J. A. Bondy, 1972) 设 A = \{A1 , A2 , . . . , An \} 是 X = \{1, 2, . . . , n\}\par
\noindent 的 n 个不同子集族, 则存在 x ∈ X 使得 A1 \textbackslash{}\{x\}, A2 \textbackslash{}\{x\}, . . . , An \textbackslash{}\{x\} 互不相同.\par
\noindent 证明: 首先注意到, 若 A, B ⊂ X, A 6= B, 且 A\textbackslash{}\{i\} = B\textbackslash{}\{i\}, 则 A = B ∪ \{i\}\par
\noindent 或 B = A ∪ \{i\}, 即 A 和 B 的对称差 A∆B = (A\textbackslash{}B) ∪ (B\textbackslash{}A) = \{i\}.\par
\noindent 利用反证法, 设对于任意 i ∈ X , 都存在 k = k(i), l = l(i), 1 ≤ k < l ≤ n ,\par
\noindent 使得 Ak \textbackslash{}\{i\} = Al \textbackslash{}\{i\} , 即 Ak ∆Ai = \{i\}. 构造简单图 G : V (G) = X, k(i)l(i) ∈\par
\noindent E(G) ⇐⇒ Ak Al = \{i\}, i = 1, 2, . . . , n. 由刚才的假设知 |E| ≥ n = |V | , 从而 G 中\par
\noindent 有一个圈, 设为 (i1 , i2 , . . . , is , i1 ). 为了方便表示, 令 ij = j = k(j), ij+1 = l(j), 从而\par
\noindent 有\par
\noindent \{s\} =A1 ∆As = (A1 ∆A2 )∆(A2 ∆A3 )∆ · · · ∆(As−1 ∆As )\par
\noindent [\par
\noindent s−1\par
\noindent ⊂ (Aj ∆Aj+1 ) = \{1, 2, . . . , s − 1\}.\par
\noindent j=1\par
\noindent 很明显矛盾, 得证.\par
}
\end{sourceblock}


\begin{explainblock}
例题是学习这本书最重要的部分。不要只看答案，要把每个限制条件变成一个计数动作：先安排谁、再插入谁、有没有重复、需不需要除以对称性。

\smallskip
\textbf{初学者检查点：}
\begin{itemize}[leftmargin=2em,itemsep=0.25em]
\item 先写出没有限制时的总数。
\item 逐条加入限制条件，观察总数怎样变化。
\item 最后检查有没有把同一种方案重复算了多次。
\end{itemize}

\smallskip
\textbf{本章读法提醒：} 第十章树计数常靠编码或递归分解；平面图公式要确认是否满足平面嵌入。
\end{explainblock}


\begin{sourceblock}{定义 10.1.16 G 中的一条走过所有边的迹叫做 G 的 Euler 迹 (Euler trail) .}
{\small\sloppy
\noindent 定义 10.1.16 G 中的一条走过所有边的迹叫做 G 的 Euler 迹 (Euler trail) .\par
\noindent 可以看出, 哥尼斯堡七桥问题实质上是问图 10.2 中是否有 Euler 迹.\par
}
\end{sourceblock}


\begin{explainblock}
定义段的任务是定边界。这里要特别注意参数范围、是否允许重复、是否考虑顺序、对象有没有标签。后面的每一道例题和证明，本质上都在反复使用这些边界。

\smallskip
\textbf{初学者检查点：}
\begin{itemize}[leftmargin=2em,itemsep=0.25em]
\item 圈出被定义的对象名称。
\item 圈出参数范围，例如 n、k 是否要求正整数。
\item 判断它是有序对象、无序对象、带标签对象还是结构对象。
\end{itemize}

\smallskip
\textbf{本章读法提醒：} 第十章树计数常靠编码或递归分解；平面图公式要确认是否满足平面嵌入。
\end{explainblock}


\begin{sourceblock}{定义 10.1.17 把 G 中每条边都经过了至少一次的闭途径称为环游 (tour), 每条边}
{\small\sloppy
\noindent 定义 10.1.17 把 G 中每条边都经过了至少一次的闭途径称为环游 (tour), 每条边\par
\noindent 都经过正好一次的闭途径称为 Euler 环游 (Euler tour), 同时也是闭 Euler 迹.\par
}
\end{sourceblock}


\begin{explainblock}
定义段的任务是定边界。这里要特别注意参数范围、是否允许重复、是否考虑顺序、对象有没有标签。后面的每一道例题和证明，本质上都在反复使用这些边界。

\smallskip
\textbf{初学者检查点：}
\begin{itemize}[leftmargin=2em,itemsep=0.25em]
\item 圈出被定义的对象名称。
\item 圈出参数范围，例如 n、k 是否要求正整数。
\item 判断它是有序对象、无序对象、带标签对象还是结构对象。
\end{itemize}

\smallskip
\textbf{本章读法提醒：} 第十章树计数常靠编码或递归分解；平面图公式要确认是否满足平面嵌入。
\end{explainblock}


\begin{sourceblock}{定义 10.1.18 有 Euler 环游的图称为欧拉的 (Eulerian).}
{\small\sloppy
\noindent 定义 10.1.18 有 Euler 环游的图称为欧拉的 (Eulerian).\par
\noindent 10.2 Euler 公式及应用 197\par
}
\end{sourceblock}


\begin{explainblock}
定义段的任务是定边界。这里要特别注意参数范围、是否允许重复、是否考虑顺序、对象有没有标签。后面的每一道例题和证明，本质上都在反复使用这些边界。

\smallskip
\textbf{初学者检查点：}
\begin{itemize}[leftmargin=2em,itemsep=0.25em]
\item 圈出被定义的对象名称。
\item 圈出参数范围，例如 n、k 是否要求正整数。
\item 判断它是有序对象、无序对象、带标签对象还是结构对象。
\end{itemize}

\smallskip
\textbf{本章读法提醒：} 第十章树计数常靠编码或递归分解；平面图公式要确认是否满足平面嵌入。
\end{explainblock}


\begin{sourceblock}{定理 10.1.19 非空连通图 G 为欧拉的当且仅当 G 中无度数为奇数的顶点.}
{\small\sloppy
\noindent 定理 10.1.19 非空连通图 G 为欧拉的当且仅当 G 中无度数为奇数的顶点.\par
\noindent 证明: 若 G 是欧拉的, 令 C 为 G 中的一个起点和终点是 u 的 Euler 环游. 每\par
\noindent 当一个顶点 v 所为 C 的内部顶点出现时, 都会数两条与之相连的边, 因 Euler 环游\par
\noindent 包含了 G 中所有顶点, 我们知道对任意 v 6= u, d(v) 为偶数. 同时 u 作为起点和终\par
\noindent 点, 可知 d(u) 也是偶数.\par
\noindent 反之, 连通图 G 的顶点度数都是偶数, 则至少是 2, 从而根据引理 10.1.11 知 G\par
\noindent 中包含一个闭迹 C , 从 G 中删去闭迹 C 中的边, 得到的不一定连通的图 G′ 的顶点\par
\noindent 的度仍然都是偶数, 包含一个闭迹 C ′ , 不断重复这个操作, 直到 G 中所有的边都在\par
\noindent 某个闭迹上. 由于 G 是连通图, 我们得到的任何一个闭迹, 必然与其他一个闭迹存在\par
\noindent 公共顶点, 而任意两个由公共顶点的闭迹都能合成一个大的闭迹, 直到只剩下一个闭\par
\noindent 迹, 这个闭迹就是 G 的 Euler 环游, 从而 G 是欧拉的.\par
}
\end{sourceblock}


\begin{explainblock}
定理段不要只看结论，要看它把哪两个计数方式连在一起。组合学证明最常见的技巧是：左边按一种标准数，右边按另一种标准数，然后说明两边数的是同一批对象。

\smallskip
\textbf{初学者检查点：}
\begin{itemize}[leftmargin=2em,itemsep=0.25em]
\item 先复述定理的假设，不要把适用范围扩大。
\item 找出证明中的基本计数原则：乘法、加法、双射、递推、容斥或生成函数。
\item 检查最后的等式化简是否和前面的对象解释一致。
\end{itemize}

\smallskip
\textbf{本章读法提醒：} 第十章树计数常靠编码或递归分解；平面图公式要确认是否满足平面嵌入。
\end{explainblock}


\begin{sourceblock}{推论 10.1.20 连通图 G 中有一 Euler 迹当且仅当 G 中度数为奇数的顶点最多有}
{\small\sloppy
\noindent 推论 10.1.20 连通图 G 中有一 Euler 迹当且仅当 G 中度数为奇数的顶点最多有\par
\noindent 两个.\par
\noindent 证明: 若 G 中有一个 Euler 迹, 则正如定理 10.1.19 中证明的那样, Euler 迹上\par
\noindent 起点和终点以外的顶点的度数都是偶数.\par
\noindent 反之, 若 G 是有最多两个顶点度数为奇数的非平凡连通图, 度数为奇数的顶点\par
\noindent 个数为零时, 根据定理 10.1.19 知 G 有一个 Euler 迹. 度数为奇数的顶点个数为两\par
\noindent 个时, 设这两个顶点分别是 u, v, 在 G 中增加一条连接 u, v 的边 e, 由定理 10.1.19\par
\noindent 知图 G+e 有一个 Euler 环游. 去掉这个 Euler 环游中的边 e, 剩下的即是 G 中的\par
\noindent Euler 迹.\par
\noindent 由于哥尼斯堡每块土地都连出了奇数个桥, 由推论 10.1.20 直接得到并不存在一\par
\noindent 条把每座桥都只走过一次的路.\par
\noindent 1946 年，荷兰数学家 Nicolaas de Bruijn 开始对“超级字符串问题 (superstring\par
\noindent problem)”感兴趣：找到一个最短的循环超级字符串, 使之包含所有可能出现的给定\par
\noindent 字母表中的长度为 k(k − mers) 的子字符串. 例如字母表中只有 0 和 1 的话, 所有\par
\noindent 3 − mers 就是这八个三位的二进制数:000, 001, 010, 011, 100, 101, 110, 111. 这时循\par
\noindent 环超级字符串 0001110100 就包含了所有的 3 − mers 而且是最短的. 但是对于任意\par
\noindent 的字母表和 k, 如何构造这样的字符串呢? de Bruijn 借助 Euler 解决七桥问题的方\par
\noindent 法, 解决了这个问题.\par
}
\end{sourceblock}


\begin{explainblock}
定理段不要只看结论，要看它把哪两个计数方式连在一起。组合学证明最常见的技巧是：左边按一种标准数，右边按另一种标准数，然后说明两边数的是同一批对象。

\smallskip
\textbf{初学者检查点：}
\begin{itemize}[leftmargin=2em,itemsep=0.25em]
\item 先复述定理的假设，不要把适用范围扩大。
\item 找出证明中的基本计数原则：乘法、加法、双射、递推、容斥或生成函数。
\item 检查最后的等式化简是否和前面的对象解释一致。
\end{itemize}

\smallskip
\textbf{本章读法提醒：} 第十章树计数常靠编码或递归分解；平面图公式要确认是否满足平面嵌入。
\end{explainblock}


\begin{sourceblock}{第 10 章续读}
{\small\sloppy
\noindent 下列定义的有向图称为 de Bruijn 图, 记为 B(d, n), n ≥ 1, d ≥ 2 :\par
\noindent V (B(d, n)) = \{x1 x2 · · · xn : xi ∈ \{0, 1, . . . , (d − 1)\}, i = 1, 2, . . . , n\}. 设 x, y ∈\par
\noindent V (B(d, n)), 若 x = x1 x2 · · · xn , 则有向边 (x, y) ∈ E(B(d, n)) ⇐⇒ y = x2 x3 · · · xn α,\par
\noindent 其中 α ∈ \{0, 1, . . . , (d − 1)\}, 并记有向边 (x, y) 为 x1 x2 · · · xn α.\par
\noindent 如图 10.4 为 B(2, 3) 的 de Bruijn 图, 其中的欧拉环游每条边的第一个数字\par
\noindent 按该边在环游中的顺序排成的序列 0000110010111101 称为 de Bruijn 序列, 即为\par
\noindent 4 − mers 的一个最短的循环超级字符串.\par
}
\end{sourceblock}


\begin{explainblock}
这一段是原书论述链的一部分。读的时候把它当作桥梁：它通常负责把上一个定义、例子或定理，引到下一个计数方法。

\smallskip
\textbf{初学者检查点：}
\begin{itemize}[leftmargin=2em,itemsep=0.25em]
\item 找出这一段承接的上一个概念。
\item 找出它准备引出的下一个概念。
\item 把中间的关键等式或构造单独抄出来。
\end{itemize}

\smallskip
\textbf{本章读法提醒：} 第十章树计数常靠编码或递归分解；平面图公式要确认是否满足平面嵌入。
\end{explainblock}


\begin{sourceblock}{§ 10.2 Euler 公式及应用}
{\small\sloppy
\noindent § 10.2. Euler 公式及应用\par
}
\end{sourceblock}


\begin{explainblock}
这一节先不要急着背公式。先问三个问题：对象是什么，哪些限制条件不能丢，最后到底是在数所有对象、还是在数某个统计量的分布。把这三件事说清楚，后面的公式才会有位置。

\smallskip
\textbf{初学者检查点：}
\begin{itemize}[leftmargin=2em,itemsep=0.25em]
\item 把本节出现的新对象写成一句话定义。
\item 把每个公式左边和右边分别翻译成“在数什么”。
\item 找一个最小非平凡例子，手工列出所有对象。
\end{itemize}

\smallskip
\textbf{本章读法提醒：} 第十章树计数常靠编码或递归分解；平面图公式要确认是否满足平面嵌入。
\end{explainblock}


\begin{sourceblock}{定义 10.2.1 (平面图 (planar graph)) 如果一个图能画在平面上使得它的边仅在端}
{\small\sloppy
\noindent 定义 10.2.1 (平面图 (planar graph)) 如果一个图能画在平面上使得它的边仅在端\par
\noindent 点相交, 则称这个图为可嵌入平面的, 或称为平面图. 而如果一个图无论怎样都无法\par
\noindent 1001 6 000 1 0000 010 101 1111 11 111 4 0110\par
\noindent 画在平面上, 并使得不同的边互不交叠, 那么这样的图不是平面图, 或者称为非平面\par
\noindent 图.\par
}
\end{sourceblock}


\begin{explainblock}
定义段的任务是定边界。这里要特别注意参数范围、是否允许重复、是否考虑顺序、对象有没有标签。后面的每一道例题和证明，本质上都在反复使用这些边界。

\smallskip
\textbf{初学者检查点：}
\begin{itemize}[leftmargin=2em,itemsep=0.25em]
\item 圈出被定义的对象名称。
\item 圈出参数范围，例如 n、k 是否要求正整数。
\item 判断它是有序对象、无序对象、带标签对象还是结构对象。
\end{itemize}

\smallskip
\textbf{本章读法提醒：} 第十章树计数常靠编码或递归分解；平面图公式要确认是否满足平面嵌入。
\end{explainblock}


\begin{sourceblock}{定义 10.2.2 (平面图的面) (1) 一个平面图 G 把平面分成若干连续的部分, 这些连}
{\small\sloppy
\noindent 定义 10.2.2 (平面图的面) (1) 一个平面图 G 把平面分成若干连续的部分, 这些连\par
\noindent 续的部分称为 G 的区域, 或 G 的一个面. G 的面组成的集合用 Φ 表示.\par
\noindent (2) 面积有限的区域称为平面图 G 的内部面, 否则称为 G 的外部面.\par
}
\end{sourceblock}


\begin{explainblock}
定义段的任务是定边界。这里要特别注意参数范围、是否允许重复、是否考虑顺序、对象有没有标签。后面的每一道例题和证明，本质上都在反复使用这些边界。

\smallskip
\textbf{初学者检查点：}
\begin{itemize}[leftmargin=2em,itemsep=0.25em]
\item 圈出被定义的对象名称。
\item 圈出参数范围，例如 n、k 是否要求正整数。
\item 判断它是有序对象、无序对象、带标签对象还是结构对象。
\end{itemize}

\smallskip
\textbf{本章读法提醒：} 第十章树计数常靠编码或递归分解；平面图公式要确认是否满足平面嵌入。
\end{explainblock}


\begin{sourceblock}{定理 10.2.3 (欧拉平面图公式) 设 G = (V, E) 是连通平面图, 令 f 表示 G 的面}
{\small\sloppy
\noindent 定理 10.2.3 (欧拉平面图公式) 设 G = (V, E) 是连通平面图, 令 f 表示 G 的面\par
\noindent 数, n 表示顶点数, e 表示边数, 则:\par
\noindent n+f −e=2\par
\noindent 证明: 首先, 如果 G 是一棵树, 那么我们有 e = n − 1, f = 1 . 在这种情况下容\par
\noindent 易验证定理成立.\par
\noindent 如果 G 是不为树的连通平面图. 假设此时欧拉公式不成立, 则存在一个含有最少\par
\noindent 边数的连通的平面图 G , 使得它不满足欧拉公式. 设这个最少边数的连通平面图为\par
\noindent G = (V, E), 面数为 f , 则\par
\noindent n − e + f 6= 2\par
\noindent 因为 G 不是树, 所以存在非割边 l. 显然 G − l 是平面连通图, 此时边数为 e − 1, 顶\par
\noindent 点数为 n 面数为 f − 1.\par
\noindent 由我们的边数最少假设, G − e 满足欧拉公式：\par
\noindent n − (e − 1) + (f − 1) = 2\par
\noindent 10.2 Euler 公式及应用 199\par
\noindent 此即\par
\noindent n−e+f =2\par
\noindent 产生矛盾.\par
\noindent 这一过程也可以这样理解, 对于一棵树, 其顶点、边和面的关系满足欧拉平面图\par
}
\end{sourceblock}


\begin{explainblock}
定理段不要只看结论，要看它把哪两个计数方式连在一起。组合学证明最常见的技巧是：左边按一种标准数，右边按另一种标准数，然后说明两边数的是同一批对象。

\smallskip
\textbf{初学者检查点：}
\begin{itemize}[leftmargin=2em,itemsep=0.25em]
\item 先复述定理的假设，不要把适用范围扩大。
\item 找出证明中的基本计数原则：乘法、加法、双射、递推、容斥或生成函数。
\item 检查最后的等式化简是否和前面的对象解释一致。
\end{itemize}

\smallskip
\textbf{本章读法提醒：} 第十章树计数常靠编码或递归分解；平面图公式要确认是否满足平面嵌入。
\end{explainblock}


\begin{sourceblock}{第 10 章续读}
{\small\sloppy
\noindent 公式, 若此时连接树的任意两顶点形成一条边, 那么边数 e 增加 1 , 但是新产生的边\par
\noindent 产生一条回路, 所以面数也增加 1, 这一关系在不断向此图中添加新的边的过程中始\par
\noindent 终满足, 而每一个连通的平面图都可以由这样的方式 “ 生成”, 从而定理成立.\par
}
\end{sourceblock}


\begin{explainblock}
这一段是原书论述链的一部分。读的时候把它当作桥梁：它通常负责把上一个定义、例子或定理，引到下一个计数方法。

\smallskip
\textbf{初学者检查点：}
\begin{itemize}[leftmargin=2em,itemsep=0.25em]
\item 找出这一段承接的上一个概念。
\item 找出它准备引出的下一个概念。
\item 把中间的关键等式或构造单独抄出来。
\end{itemize}

\smallskip
\textbf{本章读法提醒：} 第十章树计数常靠编码或递归分解；平面图公式要确认是否满足平面嵌入。
\end{explainblock}


\begin{sourceblock}{定义 10.2.4 (完全图) 每一对不同的顶点都有一条边相连的简单图称为完全图.}
{\small\sloppy
\noindent 定义 10.2.4 (完全图) 每一对不同的顶点都有一条边相连的简单图称为完全图.\par
\noindent 在同构意义下, n 个顶点的完全图只有一个, 记为 Kn\par
}
\end{sourceblock}


\begin{explainblock}
定义段的任务是定边界。这里要特别注意参数范围、是否允许重复、是否考虑顺序、对象有没有标签。后面的每一道例题和证明，本质上都在反复使用这些边界。

\smallskip
\textbf{初学者检查点：}
\begin{itemize}[leftmargin=2em,itemsep=0.25em]
\item 圈出被定义的对象名称。
\item 圈出参数范围，例如 n、k 是否要求正整数。
\item 判断它是有序对象、无序对象、带标签对象还是结构对象。
\end{itemize}

\smallskip
\textbf{本章读法提醒：} 第十章树计数常靠编码或递归分解；平面图公式要确认是否满足平面嵌入。
\end{explainblock}


\begin{sourceblock}{例 10.2.1 K5 图形由下图给出}
{\small\sloppy
\noindent 例 10.2.1 K5 图形由下图给出\par
}
\end{sourceblock}


\begin{explainblock}
例题是学习这本书最重要的部分。不要只看答案，要把每个限制条件变成一个计数动作：先安排谁、再插入谁、有没有重复、需不需要除以对称性。

\smallskip
\textbf{初学者检查点：}
\begin{itemize}[leftmargin=2em,itemsep=0.25em]
\item 先写出没有限制时的总数。
\item 逐条加入限制条件，观察总数怎样变化。
\item 最后检查有没有把同一种方案重复算了多次。
\end{itemize}

\smallskip
\textbf{本章读法提醒：} 第十章树计数常靠编码或递归分解；平面图公式要确认是否满足平面嵌入。
\end{explainblock}


\begin{sourceblock}{定义 10.2.5 偶图 (或二部图) 是指一个图, 它的顶点集可以分解为两个 (非空) 子}
{\small\sloppy
\noindent 定义 10.2.5 偶图 (或二部图) 是指一个图, 它的顶点集可以分解为两个 (非空) 子\par
\noindent 集 X 和 Y , 使得每条边都有一个端点在 X 中, 另一个端点在 Y 中; 这样一种分类\par
\noindent (X, Y ) 称为图的一个二分类. 完全偶图是具有二分类 (X, Y ) 的简单偶图, 其中 X\par
\noindent 的每个顶点都与 Y 的每个顶点相连; 若 |X| 而 |X|, 则这样的图记为 Km,n .\par
}
\end{sourceblock}


\begin{explainblock}
定义段的任务是定边界。这里要特别注意参数范围、是否允许重复、是否考虑顺序、对象有没有标签。后面的每一道例题和证明，本质上都在反复使用这些边界。

\smallskip
\textbf{初学者检查点：}
\begin{itemize}[leftmargin=2em,itemsep=0.25em]
\item 圈出被定义的对象名称。
\item 圈出参数范围，例如 n、k 是否要求正整数。
\item 判断它是有序对象、无序对象、带标签对象还是结构对象。
\end{itemize}

\smallskip
\textbf{本章读法提醒：} 第十章树计数常靠编码或递归分解；平面图公式要确认是否满足平面嵌入。
\end{explainblock}


\begin{sourceblock}{例 10.2.2 K3,3 图形由下图给出}
{\small\sloppy
\noindent 例 10.2.2 K3,3 图形由下图给出\par
}
\end{sourceblock}


\begin{explainblock}
例题是学习这本书最重要的部分。不要只看答案，要把每个限制条件变成一个计数动作：先安排谁、再插入谁、有没有重复、需不需要除以对称性。

\smallskip
\textbf{初学者检查点：}
\begin{itemize}[leftmargin=2em,itemsep=0.25em]
\item 先写出没有限制时的总数。
\item 逐条加入限制条件，观察总数怎样变化。
\item 最后检查有没有把同一种方案重复算了多次。
\end{itemize}

\smallskip
\textbf{本章读法提醒：} 第十章树计数常靠编码或递归分解；平面图公式要确认是否满足平面嵌入。
\end{explainblock}


\begin{sourceblock}{定理 10.2.6 K5 , K3,3 都是非平面图}
{\small\sloppy
\noindent 定理 10.2.6 K5 , K3,3 都是非平面图\par
\noindent 证明:\par
\noindent 我们采用欧拉平面图公式来证明 K5 , K3,3 都是非平面图\par
\noindent 我们用 n1 , n2 , n3 , · · · 来表示度为 1, 2, 3, · · · 的点的数量, 从而我们有\par
\noindent X\par
\noindent deg(V ) = n1 + 2n2 + 3n3 + · · · = 2e\par
\noindent 回忆我们曾令 d¯ = 2e\par
\noindent n , 我们称之为平均点度.\par
\noindent 同样地, 令 f1 , f2 , f3 , · · · 表示边界含有图 G 中 1, 2, 3, · · · 条边的面的数量,\par
\noindent 继而有\par
\noindent f1 + 2f2 + 3f3 + · · · = 2e\par
\noindent 总面数为\par
\noindent f = f3 + f4 + f5 · · ·\par
\noindent (除了只含有两个或者三个顶点的树, 至少三条边才能“围成”一个面)\par
\noindent 令 f¯ = 2e 代表平均每个面相邻的边的数量,\par
\noindent f\par
\noindent 若 K3,3 是平面图, n = 6, e = 9 , 由欧拉平面图公式知 f = 5 , 从而\par
\noindent 2e 18\par
}
\end{sourceblock}


\begin{explainblock}
定理段不要只看结论，要看它把哪两个计数方式连在一起。组合学证明最常见的技巧是：左边按一种标准数，右边按另一种标准数，然后说明两边数的是同一批对象。

\smallskip
\textbf{初学者检查点：}
\begin{itemize}[leftmargin=2em,itemsep=0.25em]
\item 先复述定理的假设，不要把适用范围扩大。
\item 找出证明中的基本计数原则：乘法、加法、双射、递推、容斥或生成函数。
\item 检查最后的等式化简是否和前面的对象解释一致。
\end{itemize}

\smallskip
\textbf{本章读法提醒：} 第十章树计数常靠编码或递归分解；平面图公式要确认是否满足平面嵌入。
\end{explainblock}


\begin{sourceblock}{第 10 章续读}
{\small\sloppy
\noindent f¯ = = <4\par
\noindent f 5\par
\noindent 这意味着 K3,3 中必有由 3 条边围成的面, 但是由 K3,3 的结构, 同一组的三个顶点之\par
\noindent 间互相没有边连接, 即不存在三角形, 产生矛盾, 故 K3,3 是非平面图\par
\noindent 若 K5 是平面图, n = 5, e = 10 , 由欧拉平面图公式知 f = 7 , 从而\par
\noindent 2e 20\par
\noindent f¯ = = <3\par
\noindent f 7\par
\noindent 除了只含有两个或者三个顶点的树, 至少三条边才能“围成”一个面, 这产生了矛盾,\par
\noindent 故 K5 是非平面图\par
}
\end{sourceblock}


\begin{explainblock}
这一段是原书论述链的一部分。读的时候把它当作桥梁：它通常负责把上一个定义、例子或定理，引到下一个计数方法。

\smallskip
\textbf{初学者检查点：}
\begin{itemize}[leftmargin=2em,itemsep=0.25em]
\item 找出这一段承接的上一个概念。
\item 找出它准备引出的下一个概念。
\item 把中间的关键等式或构造单独抄出来。
\end{itemize}

\smallskip
\textbf{本章读法提醒：} 第十章树计数常靠编码或递归分解；平面图公式要确认是否满足平面嵌入。
\end{explainblock}


\begin{sourceblock}{推论 10.2.7 设 G 是顶点数 n > 2 的简单平面图, 那么}
{\small\sloppy
\noindent 推论 10.2.7 设 G 是顶点数 n > 2 的简单平面图, 那么\par
\noindent (1) G 最多有 3n − 6 条边.\par
\noindent (2) G 中存在一个度数最多是 5 的顶点.\par
\noindent 证明: 我们假设 G 是连通的. 令 fk 表示 k- 面的个数, 则面数共计有\par
\noindent f = f1 + f2 + f3 + f4 + · · · , (10.1)\par
\noindent 通过面的边界来计算边数, 我们有\par
\noindent 2e = f1 + 2f2 + 3f3 + 4f4 · · · . (10.2)\par
\noindent 因为 G 是简单图, 所以每一个面至少有 3 条边, 所以我们得到:\par
\noindent f = f3 + f4 + f5 + · · · ,\par
\noindent 和\par
\noindent 2e = 3f3 + 4f4 + 5f5 + · · · .\par
\noindent 故\par
\noindent 2e − 3f ≥ 0.\par
\noindent 所以由欧拉公式知:\par
\noindent 3n − 6 = 3e − 3f = e + 2e − 3f ≥ e.\par
\noindent 10.2 Euler 公式及应用 201\par
}
\end{sourceblock}


\begin{explainblock}
定理段不要只看结论，要看它把哪两个计数方式连在一起。组合学证明最常见的技巧是：左边按一种标准数，右边按另一种标准数，然后说明两边数的是同一批对象。

\smallskip
\textbf{初学者检查点：}
\begin{itemize}[leftmargin=2em,itemsep=0.25em]
\item 先复述定理的假设，不要把适用范围扩大。
\item 找出证明中的基本计数原则：乘法、加法、双射、递推、容斥或生成函数。
\item 检查最后的等式化简是否和前面的对象解释一致。
\end{itemize}

\smallskip
\textbf{本章读法提醒：} 第十章树计数常靠编码或递归分解；平面图公式要确认是否满足平面嵌入。
\end{explainblock}


\begin{sourceblock}{推论 (1) 得证.}
{\small\sloppy
\noindent 推论 (1) 得证.\par
\noindent 由 (1) 知顶点的平均度满足\par
\noindent 2e 6n − 12\par
\noindent d¯ = ≤ < 6.\par
\noindent n n\par
}
\end{sourceblock}


\begin{explainblock}
定理段不要只看结论，要看它把哪两个计数方式连在一起。组合学证明最常见的技巧是：左边按一种标准数，右边按另一种标准数，然后说明两边数的是同一批对象。

\smallskip
\textbf{初学者检查点：}
\begin{itemize}[leftmargin=2em,itemsep=0.25em]
\item 先复述定理的假设，不要把适用范围扩大。
\item 找出证明中的基本计数原则：乘法、加法、双射、递推、容斥或生成函数。
\item 检查最后的等式化简是否和前面的对象解释一致。
\end{itemize}

\smallskip
\textbf{本章读法提醒：} 第十章树计数常靠编码或递归分解；平面图公式要确认是否满足平面嵌入。
\end{explainblock}


\begin{sourceblock}{推论 (2) 得证.}
{\small\sloppy
\noindent 推论 (2) 得证.\par
\noindent 我们下面介绍 Jordan 曲线. 一条 Jordan 曲线是指一条连续的, 自身不相交的,\par
\noindent 起点和终点相重合的曲线. 对于平面图, 一个由各条边形成的圈则构成一条 Jordan\par
\noindent 曲线. 关于 Jordan 曲线有一个著名的定理, 叙述如下:\par
}
\end{sourceblock}


\begin{explainblock}
定理段不要只看结论，要看它把哪两个计数方式连在一起。组合学证明最常见的技巧是：左边按一种标准数，右边按另一种标准数，然后说明两边数的是同一批对象。

\smallskip
\textbf{初学者检查点：}
\begin{itemize}[leftmargin=2em,itemsep=0.25em]
\item 先复述定理的假设，不要把适用范围扩大。
\item 找出证明中的基本计数原则：乘法、加法、双射、递推、容斥或生成函数。
\item 检查最后的等式化简是否和前面的对象解释一致。
\end{itemize}

\smallskip
\textbf{本章读法提醒：} 第十章树计数常靠编码或递归分解；平面图公式要确认是否满足平面嵌入。
\end{explainblock}


\begin{sourceblock}{定理 10.2.8 设 J 是平面上的一条 Jordan 曲线, 平面的剩下的部分被分成两个不}
{\small\sloppy
\noindent 定理 10.2.8 设 J 是平面上的一条 Jordan 曲线, 平面的剩下的部分被分成两个不\par
\noindent 相交的开集, 称为 J 的内部和外部, 分别记为 intJ 和 extJ, 并用 IntJ 和 ExtJ 分别\par
\noindent 表示它们的闭包. 若连接 intJ 中的任意点与 extJ 中的任意点的任何曲线必在某一\par
\noindent 点和 J 相交.\par
\noindent 如图:\par
\noindent 这个定理尽管在直观上很明显, 但它的严格证明却十分的困难.\par
\noindent 下面我们利用定理 10.2.8 也能证明 K5 不是平面图.\par
\noindent 假设 K5 是平面图, 不妨设其顶点为 v1 , v2 , v3 , v4 和 v5 . 由于 K5 是完全图, 即\par
\noindent 它任意两个顶点之间都有一条边存在. 因此, 圈 C = v1 v2 v3 v1 是一条平面 Jordan\par
\noindent 曲线, 而点 v4 是必然在 intC 或在 extC 内. 不妨假设 v4 ∈ intC( v4 若在内部可\par
\noindent 用同样的方法处理). 这样, 边 v4 v1 , v4 v2 及 v4 v3 将 intC 分成三个区域不妨设为\par
\noindent intC1 , intC2 和 intC3 . 则平面现在被分成了四个部分, 所以 v5 必在某一个部分. 不\par
\noindent 妨设在 extC(若 v5 ∈ intCi (i = 1, 2, 3) 的情形可类似处理), 则由于 v4 在 intC 内, 从\par
\noindent Jordan 曲线定理可推知边 v4 v5 必然在某一点与 C 相交. 这和假设矛盾.\par
}
\end{sourceblock}


\begin{explainblock}
定理段不要只看结论，要看它把哪两个计数方式连在一起。组合学证明最常见的技巧是：左边按一种标准数，右边按另一种标准数，然后说明两边数的是同一批对象。

\smallskip
\textbf{初学者检查点：}
\begin{itemize}[leftmargin=2em,itemsep=0.25em]
\item 先复述定理的假设，不要把适用范围扩大。
\item 找出证明中的基本计数原则：乘法、加法、双射、递推、容斥或生成函数。
\item 检查最后的等式化简是否和前面的对象解释一致。
\end{itemize}

\smallskip
\textbf{本章读法提醒：} 第十章树计数常靠编码或递归分解；平面图公式要确认是否满足平面嵌入。
\end{explainblock}


\begin{sourceblock}{定理 10.2.9 (1890, Percy John Heaword(英国), 五色定理) 任何简单平面图都是}
{\small\sloppy
\noindent 定理 10.2.9 (1890, Percy John Heaword(英国), 五色定理) 任何简单平面图都是\par
\noindent 5-染色的.\par
\noindent 证明: 设简单平面图 G(V, E), 分别用 c1 , c2 , c3 , c4 , c5 五中颜色对其染色. 我们\par
\noindent 对其顶点数目 n = |V | 作归纳. 显然, 当 n ≤ 5, 结论成立; 现假设 n = k(k ≥ 5) 时\par
\noindent 成立, 我们考虑 n = k + 1 时的情形, 由推论 10.2.7 (2) 知在简单平面图中存在一个\par
\noindent 顶点 v0 , 其度数 d(v0 ) ≤ 5. 设 G′ = G − v0 , 由归纳假设知 G′ 是 5-可染色的. 给 G′\par
\noindent 固定一种 5− 染色方案, 再将 v0 加入 G′ 中得到 G. 在此情况下讨论 v0 的着色.\par
\noindent (1) 若 d(v0 ) ≤ 4, 则 v0 最多关联四个顶点, 所以 v0 可至少一种颜色染色;\par
\noindent (2) 若 d(v0 ) = 5, 设 v0 的邻接点按逆时针排列为 v1 , v2 , v3 , v4 , v5 如图所示:\par
\noindent v1\par
\noindent v5\par
\noindent v0\par
\noindent v4 v2\par
\noindent v3\par
\noindent (a) 若 v1 , v2 , v3 , v4 , v5 的着色小于等于 4, 则 v0 最多邻接 4 种颜色的顶点, 所以\par
\noindent v0 是可着色的;\par
\noindent (b) 若 v1 , v2 , v3 , v4 , v5 分别被染成 c1 , c2 , c3 , c4 , c5 , 此时 v0 的着色就成了问题. 考\par
\noindent 虑如下:\par
}
\end{sourceblock}


\begin{explainblock}
定理段不要只看结论，要看它把哪两个计数方式连在一起。组合学证明最常见的技巧是：左边按一种标准数，右边按另一种标准数，然后说明两边数的是同一批对象。

\smallskip
\textbf{初学者检查点：}
\begin{itemize}[leftmargin=2em,itemsep=0.25em]
\item 先复述定理的假设，不要把适用范围扩大。
\item 找出证明中的基本计数原则：乘法、加法、双射、递推、容斥或生成函数。
\item 检查最后的等式化简是否和前面的对象解释一致。
\end{itemize}

\smallskip
\textbf{本章读法提醒：} 第十章树计数常靠编码或递归分解；平面图公式要确认是否满足平面嵌入。
\end{explainblock}


\begin{sourceblock}{第 10 章续读}
{\small\sloppy
\noindent 令 V13 是染成 c1 或 c3 的顶点集合; (假设 V ′ 是 V 的一个非空子集, 以 V ′ 为顶\par
\noindent 点集, 以两端点都在 V ′ 中的边的全体为边集所做成的子图, 称为 G 的由 V ′ 导出的\par
\noindent 子图, 简称 G 的导出子图), 设 G13 是 V13 在 G′ 的导出子图.\par
\noindent 情形 1. 若 v1 和 v3 在 G13 中不连通, 将 G13 中 v1 所在的连通分支所有的顶点\par
\noindent 颜色互换, 即: 染成 c1 的改成 c3 , 染成 c3 的改成 c1 ; 则得到了 G′ 的另一种 5− 着\par
\noindent 色问题. 此时 v1 和 v3 都染成了 c3 , 所以只需给 v0 染成 c1 即可;\par
\noindent c3\par
\noindent c1\par
\noindent v1\par
\noindent v5 c1\par
\noindent c5 c3\par
\noindent v1 v0\par
\noindent v5\par
\noindent v4 v\par
\noindent v0 c4 c2 2\par
\noindent c3 v3\par
\noindent v4 v2 =⇒\par
\noindent c1\par
}
\end{sourceblock}


\begin{explainblock}
这一段是原书论述链的一部分。读的时候把它当作桥梁：它通常负责把上一个定义、例子或定理，引到下一个计数方法。

\smallskip
\textbf{初学者检查点：}
\begin{itemize}[leftmargin=2em,itemsep=0.25em]
\item 找出这一段承接的上一个概念。
\item 找出它准备引出的下一个概念。
\item 把中间的关键等式或构造单独抄出来。
\end{itemize}

\smallskip
\textbf{本章读法提醒：} 第十章树计数常靠编码或递归分解；平面图公式要确认是否满足平面嵌入。
\end{explainblock}


\begin{sourceblock}{第 10 章续读}
{\small\sloppy
\noindent v3 c3\par
\noindent 情形 2.\par
\noindent 若 v1 和 v3 在 G13 中有连通的路径 P13 . 令 V24 为染成 c2 或 c4 颜色的所有顶点\par
\noindent 的集合, G24 为 V24 在 G′ 中的导出子图. 在 G′ 中 v2 和 v4 的通路 P24 一定不存在,\par
\noindent 否则与 G 是平面图矛盾. 所以对 v2 和 v4 类似情形 1 处理即可得到 G 的一种 5-着\par
\noindent 色.\par
\noindent 综上, 当 n = k + 1 时结论仍然成立. 由归纳原理, 定理得证.\par
}
\end{sourceblock}


\begin{explainblock}
这一段是原书论述链的一部分。读的时候把它当作桥梁：它通常负责把上一个定义、例子或定理，引到下一个计数方法。

\smallskip
\textbf{初学者检查点：}
\begin{itemize}[leftmargin=2em,itemsep=0.25em]
\item 找出这一段承接的上一个概念。
\item 找出它准备引出的下一个概念。
\item 把中间的关键等式或构造单独抄出来。
\end{itemize}

\smallskip
\textbf{本章读法提醒：} 第十章树计数常靠编码或递归分解；平面图公式要确认是否满足平面嵌入。
\end{explainblock}


\begin{sourceblock}{定理 10.2.10 (Sylvester-Gallai 定理) 任意给定 n ≥ 3 个平面上的不共线的点, 一}
{\small\sloppy
\noindent 定理 10.2.10 (Sylvester-Gallai 定理) 任意给定 n ≥ 3 个平面上的不共线的点, 一\par
\noindent 定存在一条直线经过其中恰好两个点.\par
\noindent 证明: 方法一 (由 Euler 公式推出 Sylvester-Gallai 的证明 [2]): 如果我们\par
\noindent 如图 10.5 所示把 R2 平面镶嵌到 R3 中单位球面 S 2 上, 那么 R2 中的每个点对\par
\noindent 10.2 Euler 公式及应用 203\par
\noindent c3\par
\noindent v1 v1\par
\noindent v5 c1\par
\noindent c5 c1 v5\par
\noindent v0 P13\par
\noindent v0\par
\noindent v4 v\par
\noindent c4 c2 2 c3\par
\noindent v4 v2\par
\noindent c3 v3 v3\par
\noindent c1 c1\par
\noindent c3 P24\par
\noindent 应了球面 S 2 上的一对对径点, R2 中的每条直线对应了 S 2 中的一个大圆. 从而\par
}
\end{sourceblock}


\begin{explainblock}
定理段不要只看结论，要看它把哪两个计数方式连在一起。组合学证明最常见的技巧是：左边按一种标准数，右边按另一种标准数，然后说明两边数的是同一批对象。

\smallskip
\textbf{初学者检查点：}
\begin{itemize}[leftmargin=2em,itemsep=0.25em]
\item 先复述定理的假设，不要把适用范围扩大。
\item 找出证明中的基本计数原则：乘法、加法、双射、递推、容斥或生成函数。
\item 检查最后的等式化简是否和前面的对象解释一致。
\end{itemize}

\smallskip
\textbf{本章读法提醒：} 第十章树计数常靠编码或递归分解；平面图公式要确认是否满足平面嵌入。
\end{explainblock}


\begin{sourceblock}{第 10 章续读}
{\small\sloppy
\noindent Sylvester-Gallai 定理等价于如下命题:\par
\noindent 给定任何 n ≥ 3 对球面 S 2 上的不共圆的对径点, 总存在一个大圆恰好经过两对\par
\noindent 对径点.\par
\noindent 对偶地, 我们把每对对径点用对应的球面上的大圆代替. 也就是说, 我们不考虑\par
\noindent ±v ∈ S 2 而是用 Cv := x ∈ S 2 : (x, v) = 0 给定的正交圆去代替它. (如果 v 是球面\par
\noindent 上的北极点, 则这个 Cv 是赤道.)\par
\noindent 那么 Sylvester-Gallai 问题要求我们去证明:\par
\noindent 在球面 S 2 上给定 n ≥ 3 个不共点大圆, 那么一定存在一个点恰好在两个大圆上.\par
\noindent 但是 S 2 上的这些大圆的安排给出了一个简单平面图. 顶点是其中两个大圆的交\par
\noindent 点, 它们又把大圆分开为边. 根据构造, 所有顶点的度数都是偶数, 并且至少是 4. 从\par
\noindent 而得出了 4 度顶点的存在性. 得证!\par
\noindent 方法二 (Leroy Milton Kelly): 设给定点组成的集合为 P, 考虑集合 L 为所有经\par
\noindent 过 P 中至少两点的直线. 在所有满足点 P 不再直线 ℓ 上的有序对 (P, ℓ) 中, 选择一\par
\noindent 对 P0 , ℓ0 使得 P0 到 ℓ0 的距离最短, 令 Q 为直线 ℓ0 上的距离 P0 最近的点 (垂点).\par
\noindent 我们断言 ℓ0 是满足定理的直线.\par
\noindent 如果不是, 那么 ℓ0 至少含有 P 中三个点, 则至少有两点落在 Q 的同一侧, 不妨\par
\noindent 设为 P1 , P2 , 且 P1 落在 Q 和 P2 之间, 不排除 P1 与 Q 重合的可能性. 如图所示, P1\par
\noindent 到过 P0 和 P2 的直线的距离比 P0 到 ℓ0 的距离小, 这与选取的 (P0 , ℓ0 ) 矛盾.\par
}
\end{sourceblock}


\begin{explainblock}
这一段是原书论述链的一部分。读的时候把它当作桥梁：它通常负责把上一个定义、例子或定理，引到下一个计数方法。

\smallskip
\textbf{初学者检查点：}
\begin{itemize}[leftmargin=2em,itemsep=0.25em]
\item 找出这一段承接的上一个概念。
\item 找出它准备引出的下一个概念。
\item 把中间的关键等式或构造单独抄出来。
\end{itemize}

\smallskip
\textbf{本章读法提醒：} 第十章树计数常靠编码或递归分解；平面图公式要确认是否满足平面嵌入。
\end{explainblock}


\begin{sourceblock}{第 10 章续读}
{\small\sloppy
\noindent P0\par
\noindent ℓ1\par
\noindent Q P1 P2 ℓ0\par
}
\end{sourceblock}


\begin{explainblock}
这一段是原书论述链的一部分。读的时候把它当作桥梁：它通常负责把上一个定义、例子或定理，引到下一个计数方法。

\smallskip
\textbf{初学者检查点：}
\begin{itemize}[leftmargin=2em,itemsep=0.25em]
\item 找出这一段承接的上一个概念。
\item 找出它准备引出的下一个概念。
\item 把中间的关键等式或构造单独抄出来。
\end{itemize}

\smallskip
\textbf{本章读法提醒：} 第十章树计数常靠编码或递归分解；平面图公式要确认是否满足平面嵌入。
\end{explainblock}


\begin{sourceblock}{推论 10.2.11 简单多面体的顶点数 V , 棱数 E 和面数 F 满足:}
{\small\sloppy
\noindent 推论 10.2.11 简单多面体的顶点数 V , 棱数 E 和面数 F 满足:\par
\noindent V − E + F = 2,\par
\noindent 同样被称作欧拉公式.\par
\noindent 证明: 下面的证明由柯西贡献.\par
\noindent 去掉多面体的一个面, 保留顶点和边, 通过将去掉的面的边拉长把剩下的面和边\par
\noindent 都拉伸到一个平面上形成一个平面图. 正常的面不再是正常的多边形, 但是, 点, 边\par
\noindent 和面的个数保持不变, 和之前的多面体一样; 故而由平面图上的欧拉公式立证.\par
}
\end{sourceblock}


\begin{explainblock}
定理段不要只看结论，要看它把哪两个计数方式连在一起。组合学证明最常见的技巧是：左边按一种标准数，右边按另一种标准数，然后说明两边数的是同一批对象。

\smallskip
\textbf{初学者检查点：}
\begin{itemize}[leftmargin=2em,itemsep=0.25em]
\item 先复述定理的假设，不要把适用范围扩大。
\item 找出证明中的基本计数原则：乘法、加法、双射、递推、容斥或生成函数。
\item 检查最后的等式化简是否和前面的对象解释一致。
\end{itemize}

\smallskip
\textbf{本章读法提醒：} 第十章树计数常靠编码或递归分解；平面图公式要确认是否满足平面嵌入。
\end{explainblock}


\begin{sourceblock}{命题 10.2.12 仅有五个正多面体, 分别为正四面体, 正六面体, 正八面体, 正十二}
{\small\sloppy
\noindent 命题 10.2.12 仅有五个正多面体, 分别为正四面体, 正六面体, 正八面体, 正十二\par
\noindent 面体和正二十面体.\par
\noindent 证明: 依旧设一正多面体的顶点数为 V , 棱数为 E 及面数为 F . 并设该正多面\par
\noindent 体的每个面是 p 边形及相交于每个顶点的面数为 q. 注意到 p, q 均大于 2. 根据正多\par
\noindent 面体的特性, 我们得到:\par
\noindent 2E = pF, (10.3)\par
\noindent 及\par
\noindent 2E = qV. (10.4)\par
\noindent 10.2 Euler 公式及应用 205\par
\noindent 故而有\par
\noindent 2E 2E\par
\noindent F = , V = .\par
\noindent p q\par
\noindent 由欧拉公式知:\par
\noindent 2E 2E\par
\noindent 2=V −E+F = −E+\par
\noindent q p\par
\noindent = −1+ E\par
}
\end{sourceblock}


\begin{explainblock}
定理段不要只看结论，要看它把哪两个计数方式连在一起。组合学证明最常见的技巧是：左边按一种标准数，右边按另一种标准数，然后说明两边数的是同一批对象。

\smallskip
\textbf{初学者检查点：}
\begin{itemize}[leftmargin=2em,itemsep=0.25em]
\item 先复述定理的假设，不要把适用范围扩大。
\item 找出证明中的基本计数原则：乘法、加法、双射、递推、容斥或生成函数。
\item 检查最后的等式化简是否和前面的对象解释一致。
\end{itemize}

\smallskip
\textbf{本章读法提醒：} 第十章树计数常靠编码或递归分解；平面图公式要确认是否满足平面嵌入。
\end{explainblock}


\begin{sourceblock}{第 10 章续读}
{\small\sloppy
\noindent q p\par
\noindent 2p + 2q − pq\par
\noindent = E,\par
\noindent pq\par
\noindent 即\par
\noindent 2pq\par
\noindent E= , (10.5)\par
\noindent 2p + 2q − pq\par
\noindent 再由 (10.3) 及 (10.4) 可知:\par
\noindent 4q\par
\noindent F = , (10.6)\par
\noindent 2p + 2q − pq\par
\noindent 4p\par
\noindent V = . (10.7)\par
\noindent 2p + 2q − pq\par
\noindent 由上面的公式显然可推知:\par
\noindent 2p + 2q − pq > 0,\par
\noindent 也可以变形为\par
}
\end{sourceblock}


\begin{explainblock}
这一段是原书论述链的一部分。读的时候把它当作桥梁：它通常负责把上一个定义、例子或定理，引到下一个计数方法。

\smallskip
\textbf{初学者检查点：}
\begin{itemize}[leftmargin=2em,itemsep=0.25em]
\item 找出这一段承接的上一个概念。
\item 找出它准备引出的下一个概念。
\item 把中间的关键等式或构造单独抄出来。
\end{itemize}

\smallskip
\textbf{本章读法提醒：} 第十章树计数常靠编码或递归分解；平面图公式要确认是否满足平面嵌入。
\end{explainblock}


\begin{sourceblock}{第 10 章续读}
{\small\sloppy
\noindent pq − 2p − 2q + 4 < 4,\par
\noindent 上式可以因式分解为\par
\noindent (p − 2)(q − 1) < 4, (10.8)\par
\noindent 因为 p − 2 与 q − 1 二者都是正整数, 所以满足 (10.8) 的正整数 (p, q) 对有:\par
\noindent (3, 3), (4, 3), (3, 4), (3, 5), (5, 3).\par
\noindent 由 (10.5), (10.6) 和 (10.7) 知, 当 (p, q) = (3, 3) 时, E = 6, V = 4, F = 4 对应正四面\par
\noindent 体; 当 (p, q) = (4, 3) 时, E = 12, V = 8, F = 6 对应正六面体; 当 (p, q) = (3, 4) 时,\par
\noindent E = 12, V = 6, F = 8 对应正八面体; 当 (p, q) = (5, 3) 时, E = 30, V = 20, F = 12\par
\noindent 对应正十二面体; 当 (p, q) = (3, 5) 时, E = 30, V = 12, F = 20 对应正二十面体. 如\par
\noindent 下图所示:\par
\noindent 这个结论最早出现在 Euclid 的几何原本中, 其中的证明也很巧妙, 叙述如下.\par
\noindent 证明: 在一个正三角形中每一个角是 60◦ , 三个, 四个或五个这样的角能相交构\par
\noindent 成一个三维尖点. 如图 ***.\par
\noindent 如果尝试用六个这样的角去构造一个三维尖点时, 这些角将凑成一个平面, 因为\par
\noindent 围绕这个顶点角之和此时是 360◦ , 而这恰好是平面中六个 60◦ 角. 显然地, 对七个或\par
\noindent 更多个 60◦ 角要构成一个尖点是不可能的.\par
\noindent 同样, 对于正四边形, 其内角为 90◦ , 只有三个角相交构成一个尖点; 对于五边形,\par
\noindent 其内角为 108◦ , 也只有三个角相交构成一个尖点; 对于正六边形, 其内角为 120◦ , 三\par
}
\end{sourceblock}


\begin{explainblock}
这一段是原书论述链的一部分。读的时候把它当作桥梁：它通常负责把上一个定义、例子或定理，引到下一个计数方法。

\smallskip
\textbf{初学者检查点：}
\begin{itemize}[leftmargin=2em,itemsep=0.25em]
\item 找出这一段承接的上一个概念。
\item 找出它准备引出的下一个概念。
\item 把中间的关键等式或构造单独抄出来。
\end{itemize}

\smallskip
\textbf{本章读法提醒：} 第十章树计数常靠编码或递归分解；平面图公式要确认是否满足平面嵌入。
\end{explainblock}


\begin{sourceblock}{第 10 章续读}
{\small\sloppy
\noindent 个角就已经太多了, 正七边形, 八边形等等肯定都不行.\par
\noindent 这就意味着正多面体的顶点只有五种可能: 三个, 四个或五个正三角形围成; 或\par
\noindent 三个正方形围成, 或三个正五边形围成, 即我们前面所提到的: 正四面体, 正六面体,\par
\noindent 正八面体, 正十二面体及正二十面体.\par
}
\end{sourceblock}


\begin{explainblock}
这一段是原书论述链的一部分。读的时候把它当作桥梁：它通常负责把上一个定义、例子或定理，引到下一个计数方法。

\smallskip
\textbf{初学者检查点：}
\begin{itemize}[leftmargin=2em,itemsep=0.25em]
\item 找出这一段承接的上一个概念。
\item 找出它准备引出的下一个概念。
\item 把中间的关键等式或构造单独抄出来。
\end{itemize}

\smallskip
\textbf{本章读法提醒：} 第十章树计数常靠编码或递归分解；平面图公式要确认是否满足平面嵌入。
\end{explainblock}


\begin{sourceblock}{§ 10.3 无标号树的计数}
{\small\sloppy
\noindent § 10.3. 无标号树的计数\par
\noindent 在上节我们已经知道, 树是一个无圈的连通图, 它是一种非常重要的非线性结构.\par
\noindent 计数含有固定个节点的树的个数是一个很困难的问题. 这节我们只考虑二叉树与平\par
\noindent 面树的计数问题. 首先我们给出树的递归性的定义.\par
}
\end{sourceblock}


\begin{explainblock}
这一节先不要急着背公式。先问三个问题：对象是什么，哪些限制条件不能丢，最后到底是在数所有对象、还是在数某个统计量的分布。把这三件事说清楚，后面的公式才会有位置。

\smallskip
\textbf{初学者检查点：}
\begin{itemize}[leftmargin=2em,itemsep=0.25em]
\item 把本节出现的新对象写成一句话定义。
\item 把每个公式左边和右边分别翻译成“在数什么”。
\item 找一个最小非平凡例子，手工列出所有对象。
\end{itemize}

\smallskip
\textbf{本章读法提醒：} 第十章树计数常靠编码或递归分解；平面图公式要确认是否满足平面嵌入。
\end{explainblock}


\begin{sourceblock}{定义 10.3.1 树 (有根树) 是具有 n(n ≥ 1) 个节点且满足下列两个条件的有限集}
{\small\sloppy
\noindent 定义 10.3.1 树 (有根树) 是具有 n(n ≥ 1) 个节点且满足下列两个条件的有限集\par
\noindent T:\par
\noindent 1) 它有一个特殊的节点称为它的根节点 (root node);\par
\noindent 2) 除了根节点, 其他的节点划分成 m 个不相交的有限集 T1 , T2 , . . . , Tm , 而这些\par
\noindent 有限集也都是树 (有根树), 称为 T 的子树. 若这些子树 T1 , T2 , . . . , Tm 是有序的则称\par
\noindent T 为有序树 (Ordered Tree) 或平面树 (Plane Tree) 或 Catalan 树 (Catalan Tree).\par
}
\end{sourceblock}


\begin{explainblock}
定义段的任务是定边界。这里要特别注意参数范围、是否允许重复、是否考虑顺序、对象有没有标签。后面的每一道例题和证明，本质上都在反复使用这些边界。

\smallskip
\textbf{初学者检查点：}
\begin{itemize}[leftmargin=2em,itemsep=0.25em]
\item 圈出被定义的对象名称。
\item 圈出参数范围，例如 n、k 是否要求正整数。
\item 判断它是有序对象、无序对象、带标签对象还是结构对象。
\end{itemize}

\smallskip
\textbf{本章读法提醒：} 第十章树计数常靠编码或递归分解；平面图公式要确认是否满足平面嵌入。
\end{explainblock}


\begin{sourceblock}{例 10.3.1 节点有 4 个的有根树有 4 个 (如图 10.6), 节点有 4 个的平面图有 5}
{\small\sloppy
\noindent 例 10.3.1 节点有 4 个的有根树有 4 个 (如图 10.6), 节点有 4 个的平面图有 5\par
\noindent 个 (如图 10.7).\par
\noindent 根节点 根节点 根节点 根节点\par
\noindent 根节点 根节点 根节点 根节点 根节点\par
\noindent 一些不交树的集合 (可以为空) 称为森林, 因此树也可由森林的概念进行定义.\par
}
\end{sourceblock}


\begin{explainblock}
例题是学习这本书最重要的部分。不要只看答案，要把每个限制条件变成一个计数动作：先安排谁、再插入谁、有没有重复、需不需要除以对称性。

\smallskip
\textbf{初学者检查点：}
\begin{itemize}[leftmargin=2em,itemsep=0.25em]
\item 先写出没有限制时的总数。
\item 逐条加入限制条件，观察总数怎样变化。
\item 最后检查有没有把同一种方案重复算了多次。
\end{itemize}

\smallskip
\textbf{本章读法提醒：} 第十章树计数常靠编码或递归分解；平面图公式要确认是否满足平面嵌入。
\end{explainblock}


\begin{sourceblock}{定义 10.3.2 树 (有根树) 是具有 n(n ≥ 1) 个节点有限集 T , 满足:}
{\small\sloppy
\noindent 定义 10.3.2 树 (有根树) 是具有 n(n ≥ 1) 个节点有限集 T , 满足:\par
\noindent 1) 它有一个特殊的节点称为它的根节点 (root node);\par
\noindent 2) 除根节点外, 其他的节点形成一个森林.\par
\noindent 因此如果删除树的根节点, 就得到了一个森林; 相反的, 对于一个含有有限个树的\par
\noindent 森林, 我们加一个根节点, 使森林中的树都是新节点的子树, 这样我们就得到了一棵\par
\noindent 树.\par
}
\end{sourceblock}


\begin{explainblock}
定义段的任务是定边界。这里要特别注意参数范围、是否允许重复、是否考虑顺序、对象有没有标签。后面的每一道例题和证明，本质上都在反复使用这些边界。

\smallskip
\textbf{初学者检查点：}
\begin{itemize}[leftmargin=2em,itemsep=0.25em]
\item 圈出被定义的对象名称。
\item 圈出参数范围，例如 n、k 是否要求正整数。
\item 判断它是有序对象、无序对象、带标签对象还是结构对象。
\end{itemize}

\smallskip
\textbf{本章读法提醒：} 第十章树计数常靠编码或递归分解；平面图公式要确认是否满足平面嵌入。
\end{explainblock}


\begin{sourceblock}{定义 10.3.3 设 T 是一棵树, r 是它的任意一个节点. 记以 r 为根节点的子树的个}
{\small\sloppy
\noindent 定义 10.3.3 设 T 是一棵树, r 是它的任意一个节点. 记以 r 为根节点的子树的个\par
\noindent 数 dr (T ), 则称 dr (T ) 为 r 的度. 度为 0 的节点称为叶子节点.\par
\noindent 树中的节点可以根据到根的距离分层. 一颗树的层数叫做这棵树的高度. 节点最\par
\noindent 多的那一层的节点数叫做这棵树的宽度. 每条边都有一个特殊的方向: 指向根节点\par
\noindent 的方向, 或者说上一层的方向 (或者相反的, 指向叶节点的方向, 下一层的方向).\par
}
\end{sourceblock}


\begin{explainblock}
定义段的任务是定边界。这里要特别注意参数范围、是否允许重复、是否考虑顺序、对象有没有标签。后面的每一道例题和证明，本质上都在反复使用这些边界。

\smallskip
\textbf{初学者检查点：}
\begin{itemize}[leftmargin=2em,itemsep=0.25em]
\item 圈出被定义的对象名称。
\item 圈出参数范围，例如 n、k 是否要求正整数。
\item 判断它是有序对象、无序对象、带标签对象还是结构对象。
\end{itemize}

\smallskip
\textbf{本章读法提醒：} 第十章树计数常靠编码或递归分解；平面图公式要确认是否满足平面嵌入。
\end{explainblock}


\begin{sourceblock}{例 10.3.2 如图 10.8, 有 7 个节点层数为 3 的树:}
{\small\sloppy
\noindent 例 10.3.2 如图 10.8, 有 7 个节点层数为 3 的树:\par
\noindent 第0层 A\par
\noindent 第1层 B C\par
\noindent 第2层 D E F\par
\noindent 第3层 G\par
\noindent 下面我们首先给出 n 条边的平面树的计数公式.\par
}
\end{sourceblock}


\begin{explainblock}
例题是学习这本书最重要的部分。不要只看答案，要把每个限制条件变成一个计数动作：先安排谁、再插入谁、有没有重复、需不需要除以对称性。

\smallskip
\textbf{初学者检查点：}
\begin{itemize}[leftmargin=2em,itemsep=0.25em]
\item 先写出没有限制时的总数。
\item 逐条加入限制条件，观察总数怎样变化。
\item 最后检查有没有把同一种方案重复算了多次。
\end{itemize}

\smallskip
\textbf{本章读法提醒：} 第十章树计数常靠编码或递归分解；平面图公式要确认是否满足平面嵌入。
\end{explainblock}


\begin{sourceblock}{定理 10.3.4 设具有 n 条边的平面树的个数为 pn , 则}
{\small\sloppy
\noindent 定理 10.3.4 设具有 n 条边的平面树的个数为 pn , 则\par
\noindent p n = Cn = .\par
\noindent n+1 n\par
\noindent 证明: 设\par
\noindent ∞\par
\noindent X\par
\noindent p(x) = pn xn ,\par
\noindent n=0\par
\noindent 下面我们将用赋权的方式给出平面树的生成函数的关系式. 给每条边赋权 x, 则\par
\noindent p(x) = 1 + xp2 (x),\par
\noindent 这与 Catalan 数的生成函数相同, 故 pn = Cn .\par
}
\end{sourceblock}


\begin{explainblock}
定理段不要只看结论，要看它把哪两个计数方式连在一起。组合学证明最常见的技巧是：左边按一种标准数，右边按另一种标准数，然后说明两边数的是同一批对象。

\smallskip
\textbf{初学者检查点：}
\begin{itemize}[leftmargin=2em,itemsep=0.25em]
\item 先复述定理的假设，不要把适用范围扩大。
\item 找出证明中的基本计数原则：乘法、加法、双射、递推、容斥或生成函数。
\item 检查最后的等式化简是否和前面的对象解释一致。
\end{itemize}

\smallskip
\textbf{本章读法提醒：} 第十章树计数常靠编码或递归分解；平面图公式要确认是否满足平面嵌入。
\end{explainblock}


\begin{sourceblock}{定义 10.3.5 二叉树是 n(n ≥ 0) 个节点的有限集, 它或为空集 (n = 0), 或由一个}
{\small\sloppy
\noindent 定义 10.3.5 二叉树是 n(n ≥ 0) 个节点的有限集, 它或为空集 (n = 0), 或由一个\par
\noindent 根节点和两棵二叉树分别称根的左子树和右子树 (互不相交的二叉树) 构成. n 个顶\par
\noindent 点的二叉树的个数记为 bn . 非叶子节点恰有两个孩子的二叉树称为完全二叉树, 具\par
\noindent 有 n 个非叶子节点的完全二叉树的个数记为 Cb (n).\par
\noindent 由定义可知二叉树有五种基本形态 (如图: 10.9):\par
\noindent ∅\par
\noindent ⃝\par
}
\end{sourceblock}


\begin{explainblock}
定义段的任务是定边界。这里要特别注意参数范围、是否允许重复、是否考虑顺序、对象有没有标签。后面的每一道例题和证明，本质上都在反复使用这些边界。

\smallskip
\textbf{初学者检查点：}
\begin{itemize}[leftmargin=2em,itemsep=0.25em]
\item 圈出被定义的对象名称。
\item 圈出参数范围，例如 n、k 是否要求正整数。
\item 判断它是有序对象、无序对象、带标签对象还是结构对象。
\end{itemize}

\smallskip
\textbf{本章读法提醒：} 第十章树计数常靠编码或递归分解；平面图公式要确认是否满足平面嵌入。
\end{explainblock}


\begin{sourceblock}{例 10.3.3 (1) 有 3 个节点的二叉树有 5 个 (如图 10.10), 即 b3 = 5.}
{\small\sloppy
\noindent 例 10.3.3 (1) 有 3 个节点的二叉树有 5 个 (如图 10.10), 即 b3 = 5.\par
\noindent 根节点 根节点 根节点 根节点 根节点\par
\noindent (2) 具有 3 个非叶子节点的完全二叉树的个数有 5 个 (如图 10.11), 即 Cb () = 5.\par
\noindent 根节点 根节点 根节点 根节点 根节点\par
\noindent 下面我们给出完全二叉树和二叉树的计数公式.\par
}
\end{sourceblock}


\begin{explainblock}
例题是学习这本书最重要的部分。不要只看答案，要把每个限制条件变成一个计数动作：先安排谁、再插入谁、有没有重复、需不需要除以对称性。

\smallskip
\textbf{初学者检查点：}
\begin{itemize}[leftmargin=2em,itemsep=0.25em]
\item 先写出没有限制时的总数。
\item 逐条加入限制条件，观察总数怎样变化。
\item 最后检查有没有把同一种方案重复算了多次。
\end{itemize}

\smallskip
\textbf{本章读法提醒：} 第十章树计数常靠编码或递归分解；平面图公式要确认是否满足平面嵌入。
\end{explainblock}


\begin{sourceblock}{定理 10.3.6 设具有 n 个非叶子节点的完全二叉树的个数为 Cb (n), 具有 n 个顶点}
{\small\sloppy
\noindent 定理 10.3.6 设具有 n 个非叶子节点的完全二叉树的个数为 Cb (n), 具有 n 个顶点\par
\noindent 的二叉树的个数为 bn , 则\par
\noindent bn = Cb (n) = Cn = . (10.9)\par
\noindent n+1 n\par
\noindent 证明: 首先, 我们说明 bn = Cb (n), 这是显然的. 因为对任意一棵 n 个节点的\par
\noindent 二叉树都可通过将它的每一个节点的左右孩子补全得到了一棵具有 n 个非叶子节点\par
\noindent 的完全二叉树; 反之, 任意一棵具有 n 个非叶子节点的完全二叉树也可将它的所有\par
\noindent 叶子删掉得到一个具有 n 个节点的二叉树. 这实际上建立了 n 个非叶子节点的完全\par
\noindent 二叉树与有 n 个顶点的二叉树之间的一一对应. 故而 bn = Cb (n).\par
\noindent 下面我们证明 bn = Cn . 由二叉树的定义知 b0 = 1, 当 n ≥ 1 时, 若删除根节点,\par
\noindent 则剩下左右两个二叉树, 所以当 n ≥ 1 时, bn 满足下面的递推关系:\par
\noindent X\par
\noindent n−1\par
\noindent bn = bk bn−1−k .\par
\noindent k=0\par
\noindent 因此 bn 与 Catalan 数 Cn 的递推关系相同且具有相同的初值, 进而 bn = Cn , 故\par
\noindent 等式 (10.9) 成立.\par
\noindent 由定理 10.3.4 和定理 10.3.6 可以得到下面的推论.\par
}
\end{sourceblock}


\begin{explainblock}
定理段不要只看结论，要看它把哪两个计数方式连在一起。组合学证明最常见的技巧是：左边按一种标准数，右边按另一种标准数，然后说明两边数的是同一批对象。

\smallskip
\textbf{初学者检查点：}
\begin{itemize}[leftmargin=2em,itemsep=0.25em]
\item 先复述定理的假设，不要把适用范围扩大。
\item 找出证明中的基本计数原则：乘法、加法、双射、递推、容斥或生成函数。
\item 检查最后的等式化简是否和前面的对象解释一致。
\end{itemize}

\smallskip
\textbf{本章读法提醒：} 第十章树计数常靠编码或递归分解；平面图公式要确认是否满足平面嵌入。
\end{explainblock}


\begin{sourceblock}{推论 10.3.7 设 Cb (n) 为有 n 个非叶子节点的完全二叉树的个数, bn 为 n 个节点}
{\small\sloppy
\noindent 推论 10.3.7 设 Cb (n) 为有 n 个非叶子节点的完全二叉树的个数, bn 为 n 个节点\par
\noindent 的二叉树的个数, pn 为 n 条边的平面树的个数, 则\par
\noindent Cb (n) = bn = pn . (10.10)\par
\noindent 下面我们通过建立双射来证明 n + 1 个顶点的平面树的个数与 n 个顶点的二叉\par
\noindent 树的个数相等.\par
\noindent 证明: 设 T 为任一个具有 n + 1 个顶点的平面树, 从右向左, 从下往上, 依次\par
\noindent 考察 T 的每个非叶子节点 R. 取 R 的最左的孩子作为 R 的左孩子, 并切断与其他\par
\noindent 孩子的连接的边, 同时将这些被切断了边的节点从左到右依次作为它们前一个兄弟\par
\noindent 的右孩子, 最后去掉 T 的根节点, 这样就得到了一个具有 n 个节点的二叉树 (如图\par
\noindent 10.12).\par
\noindent a\par
\noindent b c\par
\noindent a b c a d b\par
\noindent f h e c\par
\noindent d f d\par
\noindent e g e g\par
\noindent f\par
\noindent i\par
}
\end{sourceblock}


\begin{explainblock}
定理段不要只看结论，要看它把哪两个计数方式连在一起。组合学证明最常见的技巧是：左边按一种标准数，右边按另一种标准数，然后说明两边数的是同一批对象。

\smallskip
\textbf{初学者检查点：}
\begin{itemize}[leftmargin=2em,itemsep=0.25em]
\item 先复述定理的假设，不要把适用范围扩大。
\item 找出证明中的基本计数原则：乘法、加法、双射、递推、容斥或生成函数。
\item 检查最后的等式化简是否和前面的对象解释一致。
\end{itemize}

\smallskip
\textbf{本章读法提醒：} 第十章树计数常靠编码或递归分解；平面图公式要确认是否满足平面嵌入。
\end{explainblock}


\begin{sourceblock}{第 10 章续读}
{\small\sloppy
\noindent h i j h i j j g\par
\noindent 上面的过程显然是可逆的. 设 B 为具有 n 个顶点的二叉树, 首先添加一个新的\par
\noindent 节点, 使原来的根节点是它的左孩子, 然后从上往下, 从左往右, 将每个顶点的右孩子\par
\noindent 取下作为这个节点的右兄弟. 这样就得到了一个具有 n + 1 个顶点的平面树.\par
}
\end{sourceblock}


\begin{explainblock}
这一段是原书论述链的一部分。读的时候把它当作桥梁：它通常负责把上一个定义、例子或定理，引到下一个计数方法。

\smallskip
\textbf{初学者检查点：}
\begin{itemize}[leftmargin=2em,itemsep=0.25em]
\item 找出这一段承接的上一个概念。
\item 找出它准备引出的下一个概念。
\item 把中间的关键等式或构造单独抄出来。
\end{itemize}

\smallskip
\textbf{本章读法提醒：} 第十章树计数常靠编码或递归分解；平面图公式要确认是否满足平面嵌入。
\end{explainblock}


\begin{sourceblock}{例 10.3.4 比如 3 个结点二叉树与 3 条边平面树的一一对应 (如图 10.13).}
{\small\sloppy
\noindent 例 10.3.4 比如 3 个结点二叉树与 3 条边平面树的一一对应 (如图 10.13).\par
\noindent m m m m m\par
}
\end{sourceblock}


\begin{explainblock}
例题是学习这本书最重要的部分。不要只看答案，要把每个限制条件变成一个计数动作：先安排谁、再插入谁、有没有重复、需不需要除以对称性。

\smallskip
\textbf{初学者检查点：}
\begin{itemize}[leftmargin=2em,itemsep=0.25em]
\item 先写出没有限制时的总数。
\item 逐条加入限制条件，观察总数怎样变化。
\item 最后检查有没有把同一种方案重复算了多次。
\end{itemize}

\smallskip
\textbf{本章读法提醒：} 第十章树计数常靠编码或递归分解；平面图公式要确认是否满足平面嵌入。
\end{explainblock}


\begin{sourceblock}{§ 10.4 有标号树的计数}
{\small\sloppy
\noindent § 10.4. 有标号树的计数\par
\noindent 这节我们讨论标号树计数的问题. 对一个含有 n 个节点的有根树, 我们用\par
\noindent 1, 2, . . . , n 对其节点进行标号, 要求每个节点标号不同. 这样方式标记的树称为标号\par
\noindent 的有根树, 含有 n(n ≥ 1) 个节点的标号的有根树的个数记为 t(n). 规定 t(0) = 0.\par
}
\end{sourceblock}


\begin{explainblock}
这一节先不要急着背公式。先问三个问题：对象是什么，哪些限制条件不能丢，最后到底是在数所有对象、还是在数某个统计量的分布。把这三件事说清楚，后面的公式才会有位置。

\smallskip
\textbf{初学者检查点：}
\begin{itemize}[leftmargin=2em,itemsep=0.25em]
\item 把本节出现的新对象写成一句话定义。
\item 把每个公式左边和右边分别翻译成“在数什么”。
\item 找一个最小非平凡例子，手工列出所有对象。
\end{itemize}

\smallskip
\textbf{本章读法提醒：} 第十章树计数常靠编码或递归分解；平面图公式要确认是否满足平面嵌入。
\end{explainblock}


\begin{sourceblock}{例 10.4.1 (1) 如图 10.14 是节点标号为 \{1, 2, 3, 4, 5, 6, 7, 8, 9\}, 根结点为 6 的树.}
{\small\sloppy
\noindent 例 10.4.1 (1) 如图 10.14 是节点标号为 \{1, 2, 3, 4, 5, 6, 7, 8, 9\}, 根结点为 6 的树.\par
\noindent (2) 3 个节点标号为 \{1, 2, 3\} 的有根树共有 9 个 (如图 10.15), 即 t(3) = 9.\par
\noindent 根 根 根 根 根 根 根 根 根\par
\noindent 2 3 1 3 1 2 2 3 1 3 1 2\par
}
\end{sourceblock}


\begin{explainblock}
例题是学习这本书最重要的部分。不要只看答案，要把每个限制条件变成一个计数动作：先安排谁、再插入谁、有没有重复、需不需要除以对称性。

\smallskip
\textbf{初学者检查点：}
\begin{itemize}[leftmargin=2em,itemsep=0.25em]
\item 先写出没有限制时的总数。
\item 逐条加入限制条件，观察总数怎样变化。
\item 最后检查有没有把同一种方案重复算了多次。
\end{itemize}

\smallskip
\textbf{本章读法提醒：} 第十章树计数常靠编码或递归分解；平面图公式要确认是否满足平面嵌入。
\end{explainblock}


\begin{sourceblock}{定理 10.4.1 设 t(n) 有 n 个节点有根的标号树的个数, 则 Et (x) 为其指数型生成}
{\small\sloppy
\noindent 定理 10.4.1 设 t(n) 有 n 个节点有根的标号树的个数, 则 Et (x) 为其指数型生成\par
\noindent 函数, 即:\par
\noindent X xn\par
\noindent Et (x) = t(n) . (10.11)\par
\noindent n!\par
\noindent n≥1\par
\noindent 则 Et (x) 满足 Pólya 方程:\par
\noindent Et (x) = x exp (Et (x)). (10.12)\par
\noindent 证明: 设 Tn 为 n 个节点标号树的个数. 因为通过任选取标号树中一个顶点\par
\noindent 作为树根节点得到了一个有根的标号树, 所以若设 n 个顶点的有根标号树的个数为\par
\noindent t(n), 则我们有\par
\noindent t(n) = nTn . (10.13)\par
\noindent (注意: 以特定标号节点为根节点的标号有根树的个数与无根标号树的个数显然相\par
\noindent 同.)\par
\noindent 由一些有根树组成的森林成为有根森林, 若对含有 n(n ≥ 1) 个顶点的有根树森\par
\noindent 林, 用 [n] 进行标号, 得到的有标号树组成的森林, 称为标号的有根森林, 记它的个数\par
\noindent 为 f (n), 并规定 f (0) = 1. 因为任意一个含有 n + 1 个顶点的标号树, 若去掉标号为\par
\noindent n + 1 的顶点及与它相连的边, 得到了一个带标号的森林, 现将原来标记为 n + 1 的\par
}
\end{sourceblock}


\begin{explainblock}
定理段不要只看结论，要看它把哪两个计数方式连在一起。组合学证明最常见的技巧是：左边按一种标准数，右边按另一种标准数，然后说明两边数的是同一批对象。

\smallskip
\textbf{初学者检查点：}
\begin{itemize}[leftmargin=2em,itemsep=0.25em]
\item 先复述定理的假设，不要把适用范围扩大。
\item 找出证明中的基本计数原则：乘法、加法、双射、递推、容斥或生成函数。
\item 检查最后的等式化简是否和前面的对象解释一致。
\end{itemize}

\smallskip
\textbf{本章读法提醒：} 第十章树计数常靠编码或递归分解；平面图公式要确认是否满足平面嵌入。
\end{explainblock}


\begin{sourceblock}{第 10 章续读}
{\small\sloppy
\noindent 节点的孩子作为得到的森林中树的根, 则就得到了一个含有 n 个顶点的标号的有根\par
\noindent 森林. 这个过程显然是可逆的. 因为在一个含有 n 个顶点的标号的有根森林中加一\par
\noindent 个顶点并标记为 n + 1, 再将标号的有根森林中的每一棵树中的根节点作为新加的节\par
\noindent 点孩子, 这样就得到了一个含有 n + 1 个顶点且以 n + 1 为根的有根标号树. 因此,\par
\noindent f (n) = Tn+1 . 所以由 (10.13) 得\par
\noindent t(n + 1)\par
\noindent f (n) = .\par
\noindent n+1\par
\noindent 而由含有 n 个节点的标号的有根森林的定义, 我们显然可得到下面的关系:\par
\noindent X\par
\noindent n X\par
\noindent f (n) = t(\#B1 )t(\#B2 ) · · · t(\#Bk ) (10.14)\par
\noindent k=1 \{B1 ,B2 ,...,Bk \}\par
\noindent 其中 \{B1 , B2 , . . . , Bk \} 是集合 [n] 的 k 部划分, 因此由指数型生成函数复合公式知\par
\noindent ∞ ∞\par
\noindent !\par
\noindent X xn X xn\par
\noindent Ef (x) = f (n) = exp t(n) = exp (Et (x)). (10.15)\par
}
\end{sourceblock}


\begin{explainblock}
这一段是原书论述链的一部分。读的时候把它当作桥梁：它通常负责把上一个定义、例子或定理，引到下一个计数方法。

\smallskip
\textbf{初学者检查点：}
\begin{itemize}[leftmargin=2em,itemsep=0.25em]
\item 找出这一段承接的上一个概念。
\item 找出它准备引出的下一个概念。
\item 把中间的关键等式或构造单独抄出来。
\end{itemize}

\smallskip
\textbf{本章读法提醒：} 第十章树计数常靠编码或递归分解；平面图公式要确认是否满足平面嵌入。
\end{explainblock}


\begin{sourceblock}{第 10 章续读}
{\small\sloppy
\noindent n! n!\par
\noindent n=0 n=1\par
\noindent 又因为 f (n) = t(n+1)\par
\noindent n+1 , 所以\par
\noindent ∞\par
\noindent X ∞\par
\noindent xn 1X xn+1 Et (x)\par
\noindent Ef (x) = f (n) = t(n + 1) = , (10.16)\par
\noindent n! x (n + 1)! x\par
\noindent n=0 n=0\par
\noindent 故由 (10.15) 和 (10.16) 得\par
\noindent Et (x) = x exp (Et (x)).\par
\noindent 下面我们用拉格朗日反演公式求解有根标号树的计数函数.\par
}
\end{sourceblock}


\begin{explainblock}
这一段是原书论述链的一部分。读的时候把它当作桥梁：它通常负责把上一个定义、例子或定理，引到下一个计数方法。

\smallskip
\textbf{初学者检查点：}
\begin{itemize}[leftmargin=2em,itemsep=0.25em]
\item 找出这一段承接的上一个概念。
\item 找出它准备引出的下一个概念。
\item 把中间的关键等式或构造单独抄出来。
\end{itemize}

\smallskip
\textbf{本章读法提醒：} 第十章树计数常靠编码或递归分解；平面图公式要确认是否满足平面嵌入。
\end{explainblock}


\begin{sourceblock}{定理 10.4.2 (Cayley 定理) 设 n 个节点的有根标号树的个数为 t(n), 则}
{\small\sloppy
\noindent 定理 10.4.2 (Cayley 定理) 设 n 个节点的有根标号树的个数为 t(n), 则\par
\noindent t(n) = nn−1 .\par
\noindent 证明: 由定理 10.4.1 知有根标号树的计数函数满足下面方程\par
\noindent Et (x) = x exp (Et (x)),\par
\noindent 因此, 由 (8.16) 取 f (x) = Et (x), G(x) = E(x), k = 1 可得:\par
\noindent t(n) nn−1\par
\noindent n[xn ]Et (x) = n = [xn−1 ]e(nx) = ,\par
\noindent n! (n − 1)!\par
\noindent 所以\par
\noindent t(n) = nn−1 .\par
\noindent 由 Cayley 定理及 10.13 我们可得到下面的推论.\par
}
\end{sourceblock}


\begin{explainblock}
定理段不要只看结论，要看它把哪两个计数方式连在一起。组合学证明最常见的技巧是：左边按一种标准数，右边按另一种标准数，然后说明两边数的是同一批对象。

\smallskip
\textbf{初学者检查点：}
\begin{itemize}[leftmargin=2em,itemsep=0.25em]
\item 先复述定理的假设，不要把适用范围扩大。
\item 找出证明中的基本计数原则：乘法、加法、双射、递推、容斥或生成函数。
\item 检查最后的等式化简是否和前面的对象解释一致。
\end{itemize}

\smallskip
\textbf{本章读法提醒：} 第十章树计数常靠编码或递归分解；平面图公式要确认是否满足平面嵌入。
\end{explainblock}


\begin{sourceblock}{推论 10.4.3 具有 n 个节点无根标号树的个数为:}
{\small\sloppy
\noindent 推论 10.4.3 具有 n 个节点无根标号树的个数为:\par
\noindent Tn = nn−2 . (10.17)\par
\noindent 这里再给这个推论的一个组合证明. 关于其它证明, 读者可以参考 Aigner 和\par
\noindent Ziegler 的 Proofs From The Book [2].\par
\noindent 证明: 显然当 n = 1, 2 时, 定理成立. 当 n > 2 时, 我们通过建立 n(n > 2) 个\par
\noindent 节点的带有标号的无根树与 n − 2 元数组之间的双射来证明 (10.17) 成立.\par
\noindent 对任意一棵带标号的无根树, 由于节点数 n > 2 的树总存在叶子节点 (度为 1 的\par
\noindent 节点). 因此, 可对给定的带标号的无根树作下面三步操作:\par
\noindent (1) 给定一个辅助集合 S, 初值定为空集;\par
\noindent (2) 找出树中标号最小的叶子节点, 并将与它相邻的节点的标号有序放入集合 S\par
\noindent 中, 然后删掉这个叶子节点.\par
\noindent (3) 反复执行操作 (2) 直到只剩两个节点为止.\par
\noindent 这样我们最终得到了一个有序 n − 2 元数组 S. 因此任意给定的带标号的无根树\par
\noindent 都唯一地对应了一个长度为 n − 2 的数组且数组中的每个值都在集合 [n] 中.\par
\noindent 比如下面具有 9 个节点的标号树:\par
\noindent 注意: 在上面的操作中, 我们显然能够发现没有在 n − 2 元数组中出现过的数字\par
\noindent 对应\par
\noindent 恰好就是这棵树 (初始时) 的叶子节点的标号.\par
}
\end{sourceblock}


\begin{explainblock}
定理段不要只看结论，要看它把哪两个计数方式连在一起。组合学证明最常见的技巧是：左边按一种标准数，右边按另一种标准数，然后说明两边数的是同一批对象。

\smallskip
\textbf{初学者检查点：}
\begin{itemize}[leftmargin=2em,itemsep=0.25em]
\item 先复述定理的假设，不要把适用范围扩大。
\item 找出证明中的基本计数原则：乘法、加法、双射、递推、容斥或生成函数。
\item 检查最后的等式化简是否和前面的对象解释一致。
\end{itemize}

\smallskip
\textbf{本章读法提醒：} 第十章树计数常靠编码或递归分解；平面图公式要确认是否满足平面嵌入。
\end{explainblock}


\begin{sourceblock}{第 10 章续读}
{\small\sloppy
\noindent 下面说明上面的过程是可逆的. 我们将对任意给定一个元素值在集合 [n](n ≥ 2)\par
\noindent 的 n − 2 元数组 S 作下面的四步操作:\par
\noindent (1) 引入一个辅助集合 B, T , 它起始值为空集. 并将原始的 n − 2 元数组记为 S;\par
\noindent (2) 在集合 [n] 中划掉集合 B 中的元素和数组 S 的后 n − 2 − |B| 项出现的元素,\par
\noindent 然后在剩余的数字中选择最小的那个数字, 并将它与数组 S 中的第 |B| + 1 项构成\par
\noindent 有序对放入到集合 T 中得到新的集合 T , 同时将这个最小数放入到 B 中得到新的集\par
\noindent 合 B;\par
\noindent (3) 重复步骤 (2) 直到集合 |B| = n − 2;\par
\noindent (4) 将集合 [n] \textbackslash{} B(只有两个元素) 中的元素构成有序对放入到 T 中.\par
\noindent 这样操作后, 我们得到了有 n − 1 个有序对的集合 T . 我们将集合 T 若看作是顶\par
\noindent 点标号为 [n] 的图 G 的边集合, 则由上述操作过程可知 G 中不含圈, 且边的个数恰\par
\noindent 为 n − 1, 因此, 图 G 是 n 个节点的无根标号树.\par
\noindent 比如上例中, 给定数组 (3, 3, 5, 2, 5, 6, 1), 我们可得到的\par
\noindent T = \{(4, 3), (7, 3), (3, 5), (8, 2), (2, 5), (5, 6), (6, 1), (1, 9)\}.\par
\noindent 它对应的树也正如图 10.16 所示.\par
\noindent 因此任何一个 n − 2 元数组都唯一地对应了一棵带有标号的无根树, 从而 n 个节\par
\noindent 点的带有标号的无根树的个数为 nn−2 , 即 Tn = nn−2 .\par
}
\end{sourceblock}


\begin{explainblock}
这一段是原书论述链的一部分。读的时候把它当作桥梁：它通常负责把上一个定义、例子或定理，引到下一个计数方法。

\smallskip
\textbf{初学者检查点：}
\begin{itemize}[leftmargin=2em,itemsep=0.25em]
\item 找出这一段承接的上一个概念。
\item 找出它准备引出的下一个概念。
\item 把中间的关键等式或构造单独抄出来。
\end{itemize}

\smallskip
\textbf{本章读法提醒：} 第十章树计数常靠编码或递归分解；平面图公式要确认是否满足平面嵌入。
\end{explainblock}


\begin{sourceblock}{习 题 十}
{\small\sloppy
\noindent 习 题 十\par
}
\end{sourceblock}


\begin{explainblock}
习题段不要先追求技巧。先给题目分类：它是在问排列、组合、递推、生成函数、容斥，还是结构计数。分类正确，工具通常就只剩一两个候选。

\smallskip
\textbf{初学者检查点：}
\begin{itemize}[leftmargin=2em,itemsep=0.25em]
\item 判断题目所属工具。
\item 列出变量和限制条件。
\item 先做小规模 n=2、3、4 的例子，确认直觉。
\end{itemize}

\smallskip
\textbf{本章读法提醒：} 第十章树计数常靠编码或递归分解；平面图公式要确认是否满足平面嵌入。
\end{explainblock}


\begin{sourceblock}{习题 10.1 有根树的高度 (height) 是从该树从根到别的点的最长的距离. 记}
{\small\sloppy
\noindent 习题 10.1 有根树的高度 (height) 是从该树从根到别的点的最长的距离. 记\par
\noindent rph 为 p 个点, 高度为 h 的有根树的个数, 证明:\par
\noindent (1)rh+1,h = 1,\par
\noindent (2)rh+2,h = h,\par
\noindent h+1\par
\noindent (3)rh+3,h = + h − 1,\par
\noindent h+2 h\par
\noindent (4)rh+4,h = +2 + 2(h − 2),\par
\noindent h+3 h+1 h−1\par
\noindent (5)rh+5,h = +3 +5 + 7h − 15.\par
}
\end{sourceblock}


\begin{explainblock}
习题段不要先追求技巧。先给题目分类：它是在问排列、组合、递推、生成函数、容斥，还是结构计数。分类正确，工具通常就只剩一两个候选。

\smallskip
\textbf{初学者检查点：}
\begin{itemize}[leftmargin=2em,itemsep=0.25em]
\item 判断题目所属工具。
\item 列出变量和限制条件。
\item 先做小规模 n=2、3、4 的例子，确认直觉。
\end{itemize}

\smallskip
\textbf{本章读法提醒：} 第十章树计数常靠编码或递归分解；平面图公式要确认是否满足平面嵌入。
\end{explainblock}


\begin{sourceblock}{习题 10.2 记 sph 为 p 个点, 高度为不超过 h 的有根树的个数, 证明 sn2 的值}
{\small\sloppy
\noindent 习题 10.2 记 sph 为 p 个点, 高度为不超过 h 的有根树的个数, 证明 sn2 的值\par
\noindent 为分拆数 p(n − 1).\par
}
\end{sourceblock}


\begin{explainblock}
习题段不要先追求技巧。先给题目分类：它是在问排列、组合、递推、生成函数、容斥，还是结构计数。分类正确，工具通常就只剩一两个候选。

\smallskip
\textbf{初学者检查点：}
\begin{itemize}[leftmargin=2em,itemsep=0.25em]
\item 判断题目所属工具。
\item 列出变量和限制条件。
\item 先做小规模 n=2、3、4 的例子，确认直觉。
\end{itemize}

\smallskip
\textbf{本章读法提醒：} 第十章树计数常靠编码或递归分解；平面图公式要确认是否满足平面嵌入。
\end{explainblock}


\begin{sourceblock}{习题 10.3 计算有 n + 2 个点, 且根的每个子树的最右边的的路径都有偶数长}
{\small\sloppy
\noindent 习题 10.3 计算有 n + 2 个点, 且根的每个子树的最右边的的路径都有偶数长\par
\noindent (即有偶数条边) 的平面树的个数. 例如, n = 3 时,\par
}
\end{sourceblock}


\begin{explainblock}
习题段不要先追求技巧。先给题目分类：它是在问排列、组合、递推、生成函数、容斥，还是结构计数。分类正确，工具通常就只剩一两个候选。

\smallskip
\textbf{初学者检查点：}
\begin{itemize}[leftmargin=2em,itemsep=0.25em]
\item 判断题目所属工具。
\item 列出变量和限制条件。
\item 先做小规模 n=2、3、4 的例子，确认直觉。
\end{itemize}

\smallskip
\textbf{本章读法提醒：} 第十章树计数常靠编码或递归分解；平面图公式要确认是否满足平面嵌入。
\end{explainblock}


\begin{sourceblock}{习题 10.4 称一棵带标记的二叉树 T 是递增的, 若 T 的标号沿边的方向是递}
{\small\sloppy
\noindent 习题 10.4 称一棵带标记的二叉树 T 是递增的, 若 T 的标号沿边的方向是递\par
\noindent 增的, 即对于每一条边上, 其对应的父亲节点的标号始终大于孩子节点的标号. 证明:\par
\noindent 有 n 个节点的递增二叉树的个数为 n!.\par
}
\end{sourceblock}


\begin{explainblock}
习题段不要先追求技巧。先给题目分类：它是在问排列、组合、递推、生成函数、容斥，还是结构计数。分类正确，工具通常就只剩一两个候选。

\smallskip
\textbf{初学者检查点：}
\begin{itemize}[leftmargin=2em,itemsep=0.25em]
\item 判断题目所属工具。
\item 列出变量和限制条件。
\item 先做小规模 n=2、3、4 的例子，确认直觉。
\end{itemize}

\smallskip
\textbf{本章读法提醒：} 第十章树计数常靠编码或递归分解；平面图公式要确认是否满足平面嵌入。
\end{explainblock}


\begin{sourceblock}{习题 10.5 树可定义为无圈的连通图, 令 T 为一棵树, 顶点集记为 V, |V | ≥ 2,}
{\small\sloppy
\noindent 习题 10.5 树可定义为无圈的连通图, 令 T 为一棵树, 顶点集记为 V, |V | ≥ 2,\par
\noindent 证明 T 至少有两个叶子.\par
}
\end{sourceblock}


\begin{explainblock}
习题段不要先追求技巧。先给题目分类：它是在问排列、组合、递推、生成函数、容斥，还是结构计数。分类正确，工具通常就只剩一两个候选。

\smallskip
\textbf{初学者检查点：}
\begin{itemize}[leftmargin=2em,itemsep=0.25em]
\item 判断题目所属工具。
\item 列出变量和限制条件。
\item 先做小规模 n=2、3、4 的例子，确认直觉。
\end{itemize}

\smallskip
\textbf{本章读法提醒：} 第十章树计数常靠编码或递归分解；平面图公式要确认是否满足平面嵌入。
\end{explainblock}


\begin{sourceblock}{习题 10.6 令 T 是一棵树, v 是 T 的一个叶子节点, 去掉 v 及其相连的边, 得}
{\small\sloppy
\noindent 习题 10.6 令 T 是一棵树, v 是 T 的一个叶子节点, 去掉 v 及其相连的边, 得\par
\noindent 到 T ′ , 证明 T ′ 仍然是一棵树.\par
}
\end{sourceblock}


\begin{explainblock}
习题段不要先追求技巧。先给题目分类：它是在问排列、组合、递推、生成函数、容斥，还是结构计数。分类正确，工具通常就只剩一两个候选。

\smallskip
\textbf{初学者检查点：}
\begin{itemize}[leftmargin=2em,itemsep=0.25em]
\item 判断题目所属工具。
\item 列出变量和限制条件。
\item 先做小规模 n=2、3、4 的例子，确认直觉。
\end{itemize}

\smallskip
\textbf{本章读法提醒：} 第十章树计数常靠编码或递归分解；平面图公式要确认是否满足平面嵌入。
\end{explainblock}


\begin{sourceblock}{习题 10.7 证明: 集合 [n] 上的有根森林的数量为}
{\small\sloppy
\noindent 习题 10.7 证明: 集合 [n] 上的有根森林的数量为\par
\noindent (n + 1)n−1 .\par
}
\end{sourceblock}


\begin{explainblock}
习题段不要先追求技巧。先给题目分类：它是在问排列、组合、递推、生成函数、容斥，还是结构计数。分类正确，工具通常就只剩一两个候选。

\smallskip
\textbf{初学者检查点：}
\begin{itemize}[leftmargin=2em,itemsep=0.25em]
\item 判断题目所属工具。
\item 列出变量和限制条件。
\item 先做小规模 n=2、3、4 的例子，确认直觉。
\end{itemize}

\smallskip
\textbf{本章读法提醒：} 第十章树计数常靠编码或递归分解；平面图公式要确认是否满足平面嵌入。
\end{explainblock}


\begin{sourceblock}{习题 10.8 递增标号树是指一棵标号树, 其标号沿着每条边都是递增的. 设}
{\small\sloppy
\noindent 习题 10.8 递增标号树是指一棵标号树, 其标号沿着每条边都是递增的. 设\par
\noindent π = π1 π2 · · · πn ∈ Sn , 构造 T (π) 为 Sn 到 [n + 1] 上的递增标号树的双射, 并证明: i\par
\noindent 是 π 的下降位或者最后一位, 当且仅当 πi 是 T (π) 中的叶子.\par
}
\end{sourceblock}


\begin{explainblock}
习题段不要先追求技巧。先给题目分类：它是在问排列、组合、递推、生成函数、容斥，还是结构计数。分类正确，工具通常就只剩一两个候选。

\smallskip
\textbf{初学者检查点：}
\begin{itemize}[leftmargin=2em,itemsep=0.25em]
\item 判断题目所属工具。
\item 列出变量和限制条件。
\item 先做小规模 n=2、3、4 的例子，确认直觉。
\end{itemize}

\smallskip
\textbf{本章读法提醒：} 第十章树计数常靠编码或递归分解；平面图公式要确认是否满足平面嵌入。
\end{explainblock}


\begin{sourceblock}{习题 10.9 求有 n 个节点的三叉树的个数.}
{\small\sloppy
\noindent 习题 10.9 求有 n 个节点的三叉树的个数.\par
}
\end{sourceblock}


\begin{explainblock}
习题段不要先追求技巧。先给题目分类：它是在问排列、组合、递推、生成函数、容斥，还是结构计数。分类正确，工具通常就只剩一两个候选。

\smallskip
\textbf{初学者检查点：}
\begin{itemize}[leftmargin=2em,itemsep=0.25em]
\item 判断题目所属工具。
\item 列出变量和限制条件。
\item 先做小规模 n=2、3、4 的例子，确认直觉。
\end{itemize}

\smallskip
\textbf{本章读法提醒：} 第十章树计数常靠编码或递归分解；平面图公式要确认是否满足平面嵌入。
\end{explainblock}


\begin{sourceblock}{习题 10.10 把一个标号树的所有顶点排成一列, 若对于任意 i < j < k < l, 都}
{\small\sloppy
\noindent 习题 10.10 把一个标号树的所有顶点排成一列, 若对于任意 i < j < k < l, 都\par
\noindent 有 ik 和 jl 不同时是边,\par
\noindent 则称该树为不交的. 证明 n 个节点的不交的标号树 f (n) 的\par
\noindent 个数为 2n−1 .\par
\noindent n−1\par
\noindent x\par
}
\end{sourceblock}


\begin{explainblock}
习题段不要先追求技巧。先给题目分类：它是在问排列、组合、递推、生成函数、容斥，还是结构计数。分类正确，工具通常就只剩一两个候选。

\smallskip
\textbf{初学者检查点：}
\begin{itemize}[leftmargin=2em,itemsep=0.25em]
\item 判断题目所属工具。
\item 列出变量和限制条件。
\item 先做小规模 n=2、3、4 的例子，确认直觉。
\end{itemize}

\smallskip
\textbf{本章读法提醒：} 第十章树计数常靠编码或递归分解；平面图公式要确认是否满足平面嵌入。
\end{explainblock}


\begin{sourceblock}{习题 10.11 求 F (x) = 1−x 2 的复合逆.}
{\small\sloppy
\noindent 习题 10.11 求 F (x) = 1−x 2 的复合逆.\par
}
\end{sourceblock}


\begin{explainblock}
习题段不要先追求技巧。先给题目分类：它是在问排列、组合、递推、生成函数、容斥，还是结构计数。分类正确，工具通常就只剩一两个候选。

\smallskip
\textbf{初学者检查点：}
\begin{itemize}[leftmargin=2em,itemsep=0.25em]
\item 判断题目所属工具。
\item 列出变量和限制条件。
\item 先做小规模 n=2、3、4 的例子，确认直觉。
\end{itemize}

\smallskip
\textbf{本章读法提醒：} 第十章树计数常靠编码或递归分解；平面图公式要确认是否满足平面嵌入。
\end{explainblock}


\begin{sourceblock}{习题 10.12 利用双射证明:}
{\small\sloppy
\noindent 习题 10.12 利用双射证明:\par
\noindent 2(2n + 1)Cn = (n + 2)Cn+1 .\par
\noindent P n−1 xn , 根据 y = xey 证明\par
}
\end{sourceblock}


\begin{explainblock}
习题段不要先追求技巧。先给题目分类：它是在问排列、组合、递推、生成函数、容斥，还是结构计数。分类正确，工具通常就只剩一两个候选。

\smallskip
\textbf{初学者检查点：}
\begin{itemize}[leftmargin=2em,itemsep=0.25em]
\item 判断题目所属工具。
\item 列出变量和限制条件。
\item 先做小规模 n=2、3、4 的例子，确认直觉。
\end{itemize}

\smallskip
\textbf{本章读法提醒：} 第十章树计数常靠编码或递归分解；平面图公式要确认是否满足平面嵌入。
\end{explainblock}


\begin{sourceblock}{习题 10.13 令 y = R(x) = n≥1 n n!}
{\small\sloppy
\noindent 习题 10.13 令 y = R(x) = n≥1 n n!\par
\noindent X xn\par
\noindent [1 − R(x)]−1 = 1 + nn .\par
\noindent n!\par
\noindent n≥1\par
}
\end{sourceblock}


\begin{explainblock}
习题段不要先追求技巧。先给题目分类：它是在问排列、组合、递推、生成函数、容斥，还是结构计数。分类正确，工具通常就只剩一两个候选。

\smallskip
\textbf{初学者检查点：}
\begin{itemize}[leftmargin=2em,itemsep=0.25em]
\item 判断题目所属工具。
\item 列出变量和限制条件。
\item 先做小规模 n=2、3、4 的例子，确认直觉。
\end{itemize}

\smallskip
\textbf{本章读法提醒：} 第十章树计数常靠编码或递归分解；平面图公式要确认是否满足平面嵌入。
\end{explainblock}


\begin{sourceblock}{习题 10.14 用拉格朗日反演求}
{\small\sloppy
\noindent 习题 10.14 用拉格朗日反演求\par
\noindent 1− =1+ n+ + ··· .\par
\noindent 的前 n 项和, 比如 n = 3 时有 1 + 32 + 64 = 4.\par
}
\end{sourceblock}


\begin{explainblock}
习题段不要先追求技巧。先给题目分类：它是在问排列、组合、递推、生成函数、容斥，还是结构计数。分类正确，工具通常就只剩一两个候选。

\smallskip
\textbf{初学者检查点：}
\begin{itemize}[leftmargin=2em,itemsep=0.25em]
\item 判断题目所属工具。
\item 列出变量和限制条件。
\item 先做小规模 n=2、3、4 的例子，确认直觉。
\end{itemize}

\smallskip
\textbf{本章读法提醒：} 第十章树计数常靠编码或递归分解；平面图公式要确认是否满足平面嵌入。
\end{explainblock}


\begin{sourceblock}{习题 10.15 证明有理函数 (1 + x)2n−1 (2 + x)−n 的幂级数展开的 xn−1 的系数}
{\small\sloppy
\noindent 习题 10.15 证明有理函数 (1 + x)2n−1 (2 + x)−n 的幂级数展开的 xn−1 的系数\par
\noindent 为 21 . 也就是说满足\par
\noindent J(x)n 1\par
\noindent [xn−1 ] =\par
\noindent 1+x 2\par
\noindent 的幂级数 J(x) = (1 + x)2 /(2 + x).\par
}
\end{sourceblock}


\begin{explainblock}
习题段不要先追求技巧。先给题目分类：它是在问排列、组合、递推、生成函数、容斥，还是结构计数。分类正确，工具通常就只剩一两个候选。

\smallskip
\textbf{初学者检查点：}
\begin{itemize}[leftmargin=2em,itemsep=0.25em]
\item 判断题目所属工具。
\item 列出变量和限制条件。
\item 先做小规模 n=2、3、4 的例子，确认直觉。
\end{itemize}

\smallskip
\textbf{本章读法提醒：} 第十章树计数常靠编码或递归分解；平面图公式要确认是否满足平面嵌入。
\end{explainblock}


\chapter{单峰性与对数凹}
\begin{roadmapblock}
本章保留原书正文的主要数学骨架，并在每个定义、定理、例题和论述片段后加入解读。
阅读时不要把目标放在“背下答案”上，而是放在“看懂这一步为什么这样数”上。

\smallskip
\textbf{本章最低目标：} 能说出每个公式左边在数什么、右边在数什么，以及为什么两边相等。
\end{roadmapblock}


\begin{sourceblock}{第 11 章正文}
{\small\sloppy
\noindent 单峰性与对数凹\par
}
\end{sourceblock}


\begin{explainblock}
这一段是原书论述链的一部分。读的时候把它当作桥梁：它通常负责把上一个定义、例子或定理，引到下一个计数方法。

\smallskip
\textbf{初学者检查点：}
\begin{itemize}[leftmargin=2em,itemsep=0.25em]
\item 找出这一段承接的上一个概念。
\item 找出它准备引出的下一个概念。
\item 把中间的关键等式或构造单独抄出来。
\end{itemize}

\smallskip
\textbf{本章读法提醒：} 第十一章研究的是一串数的形状，不是单个计数答案。
\end{explainblock}


\begin{sourceblock}{§ 11.1 单峰性}
{\small\sloppy
\noindent § 11.1. 单峰性\par
\noindent 到目前为止, 我们已经介绍了很多组合数或者称作组合序列. 例如二项式序列\par
\noindent \{ nk \}nk=0 ; Eulerian 数序列 \{A(n, k)\}nk=0 ; 第二类 Stirling 数 S(n, k) 等等. 这里再次\par
\noindent 列出这三类数的递推关系式\par
\noindent n n−1 n−1\par
\noindent = + ,\par
\noindent k k−1 k\par
\noindent A(n, k) = (n − k − 1)A(n − 1, k − 1) + kA(n − 1, k),\par
\noindent S(n, k) = S(n − 1, k − 1) + kS(n − 1, k).\par
\noindent 当 n = 0, 1, 2, 3, 4, 5 时, 这三类组合数各自的取值列表如下:\par
\noindent k\par
\noindent n\par
\noindent n\par
\noindent k\par
\noindent n\par
\noindent k\par
\noindent n\par
\noindent 观察这些组合序列, 不难发现: 对于固定的 n, 这些序列均先增后减. 这样的性质\par
}
\end{sourceblock}


\begin{explainblock}
这一节先不要急着背公式。先问三个问题：对象是什么，哪些限制条件不能丢，最后到底是在数所有对象、还是在数某个统计量的分布。把这三件事说清楚，后面的公式才会有位置。

\smallskip
\textbf{初学者检查点：}
\begin{itemize}[leftmargin=2em,itemsep=0.25em]
\item 把本节出现的新对象写成一句话定义。
\item 把每个公式左边和右边分别翻译成“在数什么”。
\item 找一个最小非平凡例子，手工列出所有对象。
\end{itemize}

\smallskip
\textbf{本章读法提醒：} 第十一章研究的是一串数的形状，不是单个计数答案。
\end{explainblock}


\begin{sourceblock}{第 11 章续读}
{\small\sloppy
\noindent 在组合上有很多相关的研究, 有兴趣的同学可以参考 [19, 20].\par
}
\end{sourceblock}


\begin{explainblock}
这一段是原书论述链的一部分。读的时候把它当作桥梁：它通常负责把上一个定义、例子或定理，引到下一个计数方法。

\smallskip
\textbf{初学者检查点：}
\begin{itemize}[leftmargin=2em,itemsep=0.25em]
\item 找出这一段承接的上一个概念。
\item 找出它准备引出的下一个概念。
\item 把中间的关键等式或构造单独抄出来。
\end{itemize}

\smallskip
\textbf{本章读法提醒：} 第十一章研究的是一串数的形状，不是单个计数答案。
\end{explainblock}


\begin{sourceblock}{定义 11.1.1 (单峰性 (Unimodal)) 给定一组正实数序列 \{ai \}ni=0 , 若存在 k(0 ≤}
{\small\sloppy
\noindent 定义 11.1.1 (单峰性 (Unimodal)) 给定一组正实数序列 \{ai \}ni=0 , 若存在 k(0 ≤\par
\noindent k ≤ n) 使得\par
\noindent a0 ≤ a1 ≤ · · · ≤ ak ≥ ak+1 ≥ · · · ≥ an ,\par
\noindent 则称序列 \{ai \}ni=0 是单峰的.\par
}
\end{sourceblock}


\begin{explainblock}
定义段的任务是定边界。这里要特别注意参数范围、是否允许重复、是否考虑顺序、对象有没有标签。后面的每一道例题和证明，本质上都在反复使用这些边界。

\smallskip
\textbf{初学者检查点：}
\begin{itemize}[leftmargin=2em,itemsep=0.25em]
\item 圈出被定义的对象名称。
\item 圈出参数范围，例如 n、k 是否要求正整数。
\item 判断它是有序对象、无序对象、带标签对象还是结构对象。
\end{itemize}

\smallskip
\textbf{本章读法提醒：} 第十一章研究的是一串数的形状，不是单个计数答案。
\end{explainblock}


\begin{sourceblock}{例 11.1.1 设 n 是非负整数, 证明: 二项式序列 \{ ni \}i≥0 是单峰的.}
{\small\sloppy
\noindent 例 11.1.1 设 n 是非负整数, 证明: 二项式序列 \{ ni \}i≥0 是单峰的.\par
\noindent 证明: 考虑 \{ ni \}i≥0 中连续两个二项式系数的商. 设 k 为一整数且满足\par
\noindent n\par
\noindent n!\par
\noindent k k!(n−k)! k+1\par
\noindent = = .\par
\noindent n\par
\noindent k+1\par
\noindent n!\par
\noindent (k+1)!(n−k−1)!\par
\noindent n −k\par
\noindent 当 n 为奇数时, 若 k + 1 ≤ n − k, 即 k ≤ n−1\par
\noindent n n\par
\noindent ≤ ,\par
\noindent k k+1\par
\noindent 2 时成立. 若 k + 1 > n − k, 即 k > 2 , 此时\par
\noindent 且等号仅在 k = n−1 n−1\par
\noindent n n\par
}
\end{sourceblock}


\begin{explainblock}
例题是学习这本书最重要的部分。不要只看答案，要把每个限制条件变成一个计数动作：先安排谁、再插入谁、有没有重复、需不需要除以对称性。

\smallskip
\textbf{初学者检查点：}
\begin{itemize}[leftmargin=2em,itemsep=0.25em]
\item 先写出没有限制时的总数。
\item 逐条加入限制条件，观察总数怎样变化。
\item 最后检查有没有把同一种方案重复算了多次。
\end{itemize}

\smallskip
\textbf{本章读法提醒：} 第十一章研究的是一串数的形状，不是单个计数答案。
\end{explainblock}


\begin{sourceblock}{第 11 章续读}
{\small\sloppy
\noindent > .\par
\noindent k k+1\par
\noindent 因此\par
\noindent n n n n n\par
\noindent < < ··· < = > ··· > .\par
\noindent 0 1 (n − 1)/2 (n + 1)/2 n\par
\noindent 当 n 为偶数时, 若 k + 1 ≤ n − k, 即 k ≤ n−1 n−2\par
\noindent n n\par
\noindent < .\par
\noindent k k+1\par
\noindent 若 k + 1 > n − k, 即 k > n−1 n\par
\noindent n n\par
\noindent > .\par
\noindent k k+1\par
\noindent 因此\par
\noindent n n n n n n\par
\noindent < < ··· < < > > ··· > .\par
\noindent 0 1 (n − 2)/2 n/2 (n + 2)/2 n\par
}
\end{sourceblock}


\begin{explainblock}
这一段是原书论述链的一部分。读的时候把它当作桥梁：它通常负责把上一个定义、例子或定理，引到下一个计数方法。

\smallskip
\textbf{初学者检查点：}
\begin{itemize}[leftmargin=2em,itemsep=0.25em]
\item 找出这一段承接的上一个概念。
\item 找出它准备引出的下一个概念。
\item 把中间的关键等式或构造单独抄出来。
\end{itemize}

\smallskip
\textbf{本章读法提醒：} 第十一章研究的是一串数的形状，不是单个计数答案。
\end{explainblock}


\begin{sourceblock}{§ 11.2 对数凹性}
{\small\sloppy
\noindent § 11.2. 对数凹性\par
\noindent 下面我们给出一个比单峰性更强的性质—对数凹性.\par
}
\end{sourceblock}


\begin{explainblock}
这一节先不要急着背公式。先问三个问题：对象是什么，哪些限制条件不能丢，最后到底是在数所有对象、还是在数某个统计量的分布。把这三件事说清楚，后面的公式才会有位置。

\smallskip
\textbf{初学者检查点：}
\begin{itemize}[leftmargin=2em,itemsep=0.25em]
\item 把本节出现的新对象写成一句话定义。
\item 把每个公式左边和右边分别翻译成“在数什么”。
\item 找一个最小非平凡例子，手工列出所有对象。
\end{itemize}

\smallskip
\textbf{本章读法提醒：} 第十一章研究的是一串数的形状，不是单个计数答案。
\end{explainblock}


\begin{sourceblock}{定义 11.2.1 (对数凹性 (Log-concave)) 给定一组实序列 \{ai \}i=0}
{\small\sloppy
\noindent 定义 11.2.1 (对数凹性 (Log-concave)) 给定一组实序列 \{ai \}i=0\par
\noindent n , 若对任意的\par
\noindent k(1 ≤ k ≤ n − 1),\par
\noindent a2k ≥ ak−1 ak+1 , (11.1)\par
\noindent 则称序列 \{ai \}ni=0 是对数凹的, 若 (11.1) 中等号不成立, 则称序列 \{ai \}ni=0 是严格对\par
\noindent 数凹的.\par
}
\end{sourceblock}


\begin{explainblock}
定义段的任务是定边界。这里要特别注意参数范围、是否允许重复、是否考虑顺序、对象有没有标签。后面的每一道例题和证明，本质上都在反复使用这些边界。

\smallskip
\textbf{初学者检查点：}
\begin{itemize}[leftmargin=2em,itemsep=0.25em]
\item 圈出被定义的对象名称。
\item 圈出参数范围，例如 n、k 是否要求正整数。
\item 判断它是有序对象、无序对象、带标签对象还是结构对象。
\end{itemize}

\smallskip
\textbf{本章读法提醒：} 第十一章研究的是一串数的形状，不是单个计数答案。
\end{explainblock}


\begin{sourceblock}{注记 11.2.2 之所以称序列 \{ak \}nk=0 是对数凹的, 是因为当序列 \{ak \}nk=0 的各项都}
{\small\sloppy
\noindent 注记 11.2.2 之所以称序列 \{ak \}nk=0 是对数凹的, 是因为当序列 \{ak \}nk=0 的各项都\par
\noindent 大于 0 时, (11.1) 可以写作:\par
\noindent log(ak−1 ) + log(ak+1 )\par
\noindent log(ak ) ≥ ,\par
\noindent 而在数学分析中凹函数是指定义在区间 (a, b) 的连续函数 f 且对任意 x, y ∈ (a, b),\par
\noindent 都满足\par
\noindent x+y f (x) + f (y)\par
\noindent f ≥ .\par
}
\end{sourceblock}


\begin{explainblock}
注记通常是在补充术语、记号、历史或一个容易忽略的约定。它未必给新公式，但常常决定你后面会不会把符号读错。

\smallskip
\textbf{初学者检查点：}
\begin{itemize}[leftmargin=2em,itemsep=0.25em]
\item 记下新符号的读法。
\item 确认这个约定是否只在本章使用。
\item 把注记里的限制条件补到前面的定义旁边。
\end{itemize}

\smallskip
\textbf{本章读法提醒：} 第十一章研究的是一串数的形状，不是单个计数答案。
\end{explainblock}


\begin{sourceblock}{注记 11.2.3 不等式 (11.1) 也称作 Newton 不等式, 或 Turán 不等式.}
{\small\sloppy
\noindent 注记 11.2.3 不等式 (11.1) 也称作 Newton 不等式, 或 Turán 不等式.\par
}
\end{sourceblock}


\begin{explainblock}
注记通常是在补充术语、记号、历史或一个容易忽略的约定。它未必给新公式，但常常决定你后面会不会把符号读错。

\smallskip
\textbf{初学者检查点：}
\begin{itemize}[leftmargin=2em,itemsep=0.25em]
\item 记下新符号的读法。
\item 确认这个约定是否只在本章使用。
\item 把注记里的限制条件补到前面的定义旁边。
\end{itemize}

\smallskip
\textbf{本章读法提醒：} 第十一章研究的是一串数的形状，不是单个计数答案。
\end{explainblock}


\begin{sourceblock}{例 11.2.1 证明: 二项式系数序列 \{ ni \}ni=0 是严格对数凹的.}
{\small\sloppy
\noindent 例 11.2.1 证明: 二项式系数序列 \{ ni \}ni=0 是严格对数凹的.\par
\noindent 证明: 设 k 为非负整数且 1 ≤ k ≤ n − 1, 因为\par
\noindent n n n!2\par
\noindent =\par
\noindent k−1 k+1 (k − 1)(k + 1)(n − k − 1)(n − k + 1)\par
\noindent n!2 k(n − k)\par
\noindent =\par
\noindent (k!(n − k)!)2 (k + 1)(n − k + 1)\par
\noindent n nk − k 2\par
\noindent = ,\par
\noindent k nk − k 2 + n + 1\par
\noindent 由\par
\noindent nk − k 2\par
\noindent <1\par
\noindent n2 − k 2 + n + 1\par
\noindent 可得 2\par
\noindent n n n\par
\noindent < ,\par
}
\end{sourceblock}


\begin{explainblock}
例题是学习这本书最重要的部分。不要只看答案，要把每个限制条件变成一个计数动作：先安排谁、再插入谁、有没有重复、需不需要除以对称性。

\smallskip
\textbf{初学者检查点：}
\begin{itemize}[leftmargin=2em,itemsep=0.25em]
\item 先写出没有限制时的总数。
\item 逐条加入限制条件，观察总数怎样变化。
\item 最后检查有没有把同一种方案重复算了多次。
\end{itemize}

\smallskip
\textbf{本章读法提醒：} 第十一章研究的是一串数的形状，不是单个计数答案。
\end{explainblock}


\begin{sourceblock}{第 11 章续读}
{\small\sloppy
\noindent k−1 k+1 k\par
\noindent 也即二项式系数序列是严格对数凹的.\par
}
\end{sourceblock}


\begin{explainblock}
这一段是原书论述链的一部分。读的时候把它当作桥梁：它通常负责把上一个定义、例子或定理，引到下一个计数方法。

\smallskip
\textbf{初学者检查点：}
\begin{itemize}[leftmargin=2em,itemsep=0.25em]
\item 找出这一段承接的上一个概念。
\item 找出它准备引出的下一个概念。
\item 把中间的关键等式或构造单独抄出来。
\end{itemize}

\smallskip
\textbf{本章读法提醒：} 第十一章研究的是一串数的形状，不是单个计数答案。
\end{explainblock}


\begin{sourceblock}{定理 11.2.4 设正实数序列 \{ai \}ni=0 是对数凹的, 则 \{ai \}ni=0 是单峰的.}
{\small\sloppy
\noindent 定理 11.2.4 设正实数序列 \{ai \}ni=0 是对数凹的, 则 \{ai \}ni=0 是单峰的.\par
\noindent 证明: (反证法) 设 \{ai \}ni=0 是对数凹, 若它不是单峰的, 则存在相邻的三项满足\par
\noindent ak−1 ≥ ak ≤ ak+1 ,\par
\noindent 并且两个不等号中至少有一个的等号是不成立的, 即\par
\noindent ak ak+1\par
\noindent ≤ 1, ≥ 1,\par
\noindent ak−1 ak\par
\noindent 且两不等号不能同时取得, 所以\par
\noindent ak ak+1\par
\noindent < ,\par
\noindent ak−1 ak\par
\noindent 即\par
\noindent a2k < ak+1 ak−1 ,\par
\noindent 这与对数凹性矛盾.\par
\noindent 下面我们将给出判断实序列是否是对数凹的一个重要工具, 在此之前我们先证明\par
\noindent 下面的引理.\par
}
\end{sourceblock}


\begin{explainblock}
定理段不要只看结论，要看它把哪两个计数方式连在一起。组合学证明最常见的技巧是：左边按一种标准数，右边按另一种标准数，然后说明两边数的是同一批对象。

\smallskip
\textbf{初学者检查点：}
\begin{itemize}[leftmargin=2em,itemsep=0.25em]
\item 先复述定理的假设，不要把适用范围扩大。
\item 找出证明中的基本计数原则：乘法、加法、双射、递推、容斥或生成函数。
\item 检查最后的等式化简是否和前面的对象解释一致。
\end{itemize}

\smallskip
\textbf{本章读法提醒：} 第十一章研究的是一串数的形状，不是单个计数答案。
\end{explainblock}


\begin{sourceblock}{引理 11.2.5 设 f (x) 为一元 n(n ≥ 1) 次实系数多项式, 若 f (x) 所有的根都是实}
{\small\sloppy
\noindent 引理 11.2.5 设 f (x) 为一元 n(n ≥ 1) 次实系数多项式, 若 f (x) 所有的根都是实\par
\noindent 根, 则 f ′ (x) 所有的根也都是实根.\par
\noindent 证明: 当 f (x) 没有重根时, 则 f (x) 有 n 个互异的实根, 由罗尔定理知, f ′ (x)\par
\noindent 至少有 n − 1 个互异的实根.\par
\noindent 又因为 n − 1 次多项式 f ′ (x) 至多有 n − 1 个实根, 所以 f ′ (x) 恰有这 n − 1 个\par
\noindent 互异的实根.\par
\noindent 当 f (x) 有 重 根 时, 不 妨 设 多 项 式 f (x) 有 s 个 单 根 (重 数 为 1 的 根), 记\par
\noindent 为 r1 , r2 , . . . , rs ; 及 m 个 重 数 大 于 1 的 根, 记 为 R1 , R2 , . . . , Rm , 重 数 分 别 为\par
\noindent ℓ1 , ℓ2 , . . . , ℓm (ℓi ≥ 2, 1 ≤ i ≤ m). 则我们有下面的结论:\par
\noindent 1. 因为 f (x) 为 n 次多项式, 且只有实根, 知:\par
\noindent s + ℓ1 + ℓ2 + · · · + ℓm = n;\par
\noindent 2. 对于 1 ≤ i ≤ m, Ri 是 f ′ (x) 的 ℓi − 1 重实根;\par
\noindent 3. 由罗尔定理知, 在 f (x) 相邻的根之间存在 f ′ (x) 的一个实根, 共计 m + s − 1\par
\noindent 个.\par
\noindent 所以 n − 1 次多项式 f ′ (x) 实根的个数是\par
\noindent (m + s − 1) + (ℓ1 − 1) + (ℓ2 − 1) + · · · + (ℓm − 1) = n − 1.\par
\noindent 综上, 多项式 f ′ (x) 的根均是实根.\par
}
\end{sourceblock}


\begin{explainblock}
定理段不要只看结论，要看它把哪两个计数方式连在一起。组合学证明最常见的技巧是：左边按一种标准数，右边按另一种标准数，然后说明两边数的是同一批对象。

\smallskip
\textbf{初学者检查点：}
\begin{itemize}[leftmargin=2em,itemsep=0.25em]
\item 先复述定理的假设，不要把适用范围扩大。
\item 找出证明中的基本计数原则：乘法、加法、双射、递推、容斥或生成函数。
\item 检查最后的等式化简是否和前面的对象解释一致。
\end{itemize}

\smallskip
\textbf{本章读法提醒：} 第十一章研究的是一串数的形状，不是单个计数答案。
\end{explainblock}


\begin{sourceblock}{定理 11.2.6 设实系数多项式}
{\small\sloppy
\noindent 定理 11.2.6 设实系数多项式\par
\noindent p(x) = a0 + a1 x + a2 x2 + · · · + an xn ,\par
\noindent 其中 n ≥ 2 且 an 6= 0. 若 p(x) 的所有根都是实根, 则其系数序列 \{ai \}ni=0 是对数凹\par
\noindent 的.\par
\noindent 证明: 设 k 是任意给定正整数且 1 ≤ k ≤ n − 1.\par
\noindent (1) 当 ak−1 = 0 时, 显然有\par
\noindent a2k ≥ 0 = ak−1 ak+1 .\par
\noindent (2) 当 ak−1 6= 0 时, p(x) 的 k − 1 阶导数为\par
\noindent X\par
\noindent n\par
\noindent i\par
\noindent p (k−1)\par
\noindent (x) = (k − 1)! ai xi−k+1 ,\par
\noindent k−1\par
\noindent i=k−1\par
\noindent 所以 p(k−1) (0) = ak−1 6= 0. 故由引理 11.2.5 知 p(k−1) (x) 有 n − k + 1 个不为 0 的\par
\noindent 实根, 因此,\par
\noindent (k − 1)! X\par
}
\end{sourceblock}


\begin{explainblock}
定理段不要只看结论，要看它把哪两个计数方式连在一起。组合学证明最常见的技巧是：左边按一种标准数，右边按另一种标准数，然后说明两边数的是同一批对象。

\smallskip
\textbf{初学者检查点：}
\begin{itemize}[leftmargin=2em,itemsep=0.25em]
\item 先复述定理的假设，不要把适用范围扩大。
\item 找出证明中的基本计数原则：乘法、加法、双射、递推、容斥或生成函数。
\item 检查最后的等式化简是否和前面的对象解释一致。
\end{itemize}

\smallskip
\textbf{本章读法提醒：} 第十一章研究的是一串数的形状，不是单个计数答案。
\end{explainblock}


\begin{sourceblock}{第 11 章续读}
{\small\sloppy
\noindent n\par
\noindent (k−1) 1 i\par
\noindent p = n−k+1 ai xn−i ,\par
\noindent x x k−1\par
\noindent i=k−1\par
\noindent 也只有 n − k + 1 个不为 0 的实根. 设\par
\noindent X\par
\noindent n\par
\noindent i\par
\noindent g(x) = ai xn−i ,\par
\noindent k−1\par
\noindent i=k−1\par
\noindent 则 g(x) 的所有根都是不为 0 的实数. 又 g(x) 的 n − k − 1 阶导数为\par
\noindent k(n − k)! n − k + 1 k+1\par
\noindent g (n−k−1) (x) = ak−1 x2 + 2ak x + ak+1 .\par
\noindent 再应用罗尔定理知 g (n−k−1) (x) 只有实根, 所以对任意的 k(1 ≤ k ≤ n − 1), 若\par
\noindent ak−1 6= 0 都有:\par
\noindent n−k+1 k+1\par
}
\end{sourceblock}


\begin{explainblock}
这一段是原书论述链的一部分。读的时候把它当作桥梁：它通常负责把上一个定义、例子或定理，引到下一个计数方法。

\smallskip
\textbf{初学者检查点：}
\begin{itemize}[leftmargin=2em,itemsep=0.25em]
\item 找出这一段承接的上一个概念。
\item 找出它准备引出的下一个概念。
\item 把中间的关键等式或构造单独抄出来。
\end{itemize}

\smallskip
\textbf{本章读法提醒：} 第十一章研究的是一串数的形状，不是单个计数答案。
\end{explainblock}


\begin{sourceblock}{第 11 章续读}
{\small\sloppy
\noindent a2k ≥ · ak−1 ak+1 ≥ ak−1 ak+1 .\par
\noindent n−k k\par
\noindent 综上, 序列 \{ak \}nk=0 是对数凹的.\par
\noindent 由该定理还可以直接得到下面的推论.\par
}
\end{sourceblock}


\begin{explainblock}
这一段是原书论述链的一部分。读的时候把它当作桥梁：它通常负责把上一个定义、例子或定理，引到下一个计数方法。

\smallskip
\textbf{初学者检查点：}
\begin{itemize}[leftmargin=2em,itemsep=0.25em]
\item 找出这一段承接的上一个概念。
\item 找出它准备引出的下一个概念。
\item 把中间的关键等式或构造单独抄出来。
\end{itemize}

\smallskip
\textbf{本章读法提醒：} 第十一章研究的是一串数的形状，不是单个计数答案。
\end{explainblock}


\begin{sourceblock}{推论 11.2.7 设实系数多项式}
{\small\sloppy
\noindent 推论 11.2.7 设实系数多项式\par
\noindent p(x) = a0 + a1 x + a2 x2 + · · · + an xn ,\par
\noindent 其中 n ≥ 2 且对 ∀ 1 ≤ i ≤ n, ai 6= 0. 若 p(x) 的所有根都是实根, 则其系数序列\par
\noindent \{ai \}ni=0 是严格对数凹的, 即对任意的 1 ≤ k ≤ n − 1 都有\par
\noindent a2k > ak−1 ak+1 .\par
\noindent 且该序列是单峰的, 其最大值最多有两个.\par
}
\end{sourceblock}


\begin{explainblock}
定理段不要只看结论，要看它把哪两个计数方式连在一起。组合学证明最常见的技巧是：左边按一种标准数，右边按另一种标准数，然后说明两边数的是同一批对象。

\smallskip
\textbf{初学者检查点：}
\begin{itemize}[leftmargin=2em,itemsep=0.25em]
\item 先复述定理的假设，不要把适用范围扩大。
\item 找出证明中的基本计数原则：乘法、加法、双射、递推、容斥或生成函数。
\item 检查最后的等式化简是否和前面的对象解释一致。
\end{itemize}

\smallskip
\textbf{本章读法提醒：} 第十一章研究的是一串数的形状，不是单个计数答案。
\end{explainblock}


\begin{sourceblock}{例 11.2.2 证明: 二项式系数序列 \{ nk \}nk=0 是严格对数凹的.}
{\small\sloppy
\noindent 例 11.2.2 证明: 二项式系数序列 \{ nk \}nk=0 是严格对数凹的.\par
\noindent 证明: 因为\par
\noindent n\par
\noindent X n k\par
\noindent x = (1 + x)n ,\par
\noindent k\par
\noindent k=0\par
\noindent 它的根为 −1, 重数是 n, 且任意二项式系数都不为 0, 所以由推论 11.2.7 知二项式系\par
\noindent 数序列 \{ nk \}nk=0 是严格对数凹的.\par
}
\end{sourceblock}


\begin{explainblock}
例题是学习这本书最重要的部分。不要只看答案，要把每个限制条件变成一个计数动作：先安排谁、再插入谁、有没有重复、需不需要除以对称性。

\smallskip
\textbf{初学者检查点：}
\begin{itemize}[leftmargin=2em,itemsep=0.25em]
\item 先写出没有限制时的总数。
\item 逐条加入限制条件，观察总数怎样变化。
\item 最后检查有没有把同一种方案重复算了多次。
\end{itemize}

\smallskip
\textbf{本章读法提醒：} 第十一章研究的是一串数的形状，不是单个计数答案。
\end{explainblock}


\begin{sourceblock}{例 11.2.3 证明: 第一类无符号 Stirling 数序列 \{c(n, k)\}nk=1 为对数凹的.}
{\small\sloppy
\noindent 例 11.2.3 证明: 第一类无符号 Stirling 数序列 \{c(n, k)\}nk=1 为对数凹的.\par
\noindent 证明: 因为\par
\noindent X\par
\noindent n\par
\noindent c(n, k)xk = x(x + 1)(x + 2) · · · (x + n − 1),\par
\noindent k=1\par
\noindent 它的根为 0, −1, −2, · · · , −n + 1, 都是实根. 由定理 11.2.6 知序列 \{c(n, k)\}nk=1 为对\par
\noindent 数凹的.\par
}
\end{sourceblock}


\begin{explainblock}
例题是学习这本书最重要的部分。不要只看答案，要把每个限制条件变成一个计数动作：先安排谁、再插入谁、有没有重复、需不需要除以对称性。

\smallskip
\textbf{初学者检查点：}
\begin{itemize}[leftmargin=2em,itemsep=0.25em]
\item 先写出没有限制时的总数。
\item 逐条加入限制条件，观察总数怎样变化。
\item 最后检查有没有把同一种方案重复算了多次。
\end{itemize}

\smallskip
\textbf{本章读法提醒：} 第十一章研究的是一串数的形状，不是单个计数答案。
\end{explainblock}


\begin{sourceblock}{例 11.2.4 证明: Eulerian 数序列 \{A(n, k)\}nk=1 是对数凹的.}
{\small\sloppy
\noindent 例 11.2.4 证明: Eulerian 数序列 \{A(n, k)\}nk=1 是对数凹的.\par
\noindent 证明: 我们考虑 Eulerian 多项式:\par
\noindent X\par
\noindent n\par
\noindent An (x) = A(n, k)xk .\par
\noindent k=1\par
\noindent 注意到 An (x) 所有系数都是非负的, 所以 An (x) 不可能有正根, 所以不妨设 x ≤ 0.\par
\noindent 现令\par
\noindent An (x)\par
\noindent Gn (x) = .\par
\noindent (1 − x)n+1\par
\noindent 由 Eulrian 多项式 An (x) 的性质:\par
\noindent An (x) = x(1 − x)A′n−1 (x) + nxAn−1 (x),\par
\noindent 可知:\par
\noindent ′\par
\noindent An (x) An−1 (x)\par
\noindent =x ,\par
\noindent (1 − x)n+1 (1 − x)n\par
}
\end{sourceblock}


\begin{explainblock}
例题是学习这本书最重要的部分。不要只看答案，要把每个限制条件变成一个计数动作：先安排谁、再插入谁、有没有重复、需不需要除以对称性。

\smallskip
\textbf{初学者检查点：}
\begin{itemize}[leftmargin=2em,itemsep=0.25em]
\item 先写出没有限制时的总数。
\item 逐条加入限制条件，观察总数怎样变化。
\item 最后检查有没有把同一种方案重复算了多次。
\end{itemize}

\smallskip
\textbf{本章读法提醒：} 第十一章研究的是一串数的形状，不是单个计数答案。
\end{explainblock}


\begin{sourceblock}{第 11 章续读}
{\small\sloppy
\noindent 即:\par
\noindent Gn (x) = xG′n−1 (x).\par
\noindent 下面我们用归纳法证明 An (x) 的根都是非负实根.\par
\noindent 可以验证 n = 1, 2 时结论成立. 假设结论对小于 n(n > 2) 时都成立, 则 An−1 (x)\par
\noindent 有 n − 1 个非负实根.\par
\noindent 由 Gn−1 (x) 在复平面除去点 x = 1 的区域内都是解析函数, 所以 Gn−1 (x) 在非\par
\noindent 负实轴上有 n − 1 个零点 (计重数), 由罗尔定理以及解析函数的性质, G′n−1 (x) 在非\par
\noindent 负实轴上至少有 n − 2 个零点 (计重数), 注意到\par
\noindent An (x) = x(1 − x)−n−1 G′n−1 (x),\par
\noindent 所以 x = 0 和 G′n−1 (x) 在非负实轴上的零点都是 An (x) 的根, 所以 An (x) 至少有\par
\noindent n − 1 个非负实根, 而 An (x) 为 n 次多项式, 只有 n 个根. 剩下的一个根显然不可\par
\noindent 能是虚根 (因为实系数多项式的虚根成对出现), 故 An (x) 有 n 个非负实根. 由定理\par
\noindent 11.2.6 知序列 \{A(n, k)\}nk=1 为对数凹的.\par
}
\end{sourceblock}


\begin{explainblock}
这一段是原书论述链的一部分。读的时候把它当作桥梁：它通常负责把上一个定义、例子或定理，引到下一个计数方法。

\smallskip
\textbf{初学者检查点：}
\begin{itemize}[leftmargin=2em,itemsep=0.25em]
\item 找出这一段承接的上一个概念。
\item 找出它准备引出的下一个概念。
\item 把中间的关键等式或构造单独抄出来。
\end{itemize}

\smallskip
\textbf{本章读法提醒：} 第十一章研究的是一串数的形状，不是单个计数答案。
\end{explainblock}


\begin{sourceblock}{例 11.2.5 证明第二类 Stirling 数序列 \{S(n, k)\}nk=1 是单峰的.}
{\small\sloppy
\noindent 例 11.2.5 证明第二类 Stirling 数序列 \{S(n, k)\}nk=1 是单峰的.\par
\noindent 证明: 设\par
\noindent X\par
\noindent n\par
\noindent pn (x) = S(n, k)xk .\par
\noindent k=1\par
\noindent 因为第二类 Stiling 数 S(n, k) 满足下面的递推式\par
\noindent S(n, k) = S(n − 1, k − 1) + kS(n − 1, k),\par
\noindent 所以\par
\noindent pn (x) = xpn−1 (x) + xp′n−1 (x).\par
\noindent 令\par
\noindent fn (x) = ex pn (x),\par
\noindent 则\par
\noindent ′\par
\noindent fn (x) = x(ex pn−1 (x))′ = xfn−1 (x). (11.2)\par
\noindent 我们下面用归纳法证明 fn (x) 恰好有 n 个不同的实根, 且除 0 外其它均是负数. 当\par
\noindent n = 1 时, fn (x) = xex 只有 0 是它的零点, 所以成立, 现假设结论对 1, 2, . . . , n − 1\par
\noindent ′\par
}
\end{sourceblock}


\begin{explainblock}
例题是学习这本书最重要的部分。不要只看答案，要把每个限制条件变成一个计数动作：先安排谁、再插入谁、有没有重复、需不需要除以对称性。

\smallskip
\textbf{初学者检查点：}
\begin{itemize}[leftmargin=2em,itemsep=0.25em]
\item 先写出没有限制时的总数。
\item 逐条加入限制条件，观察总数怎样变化。
\item 最后检查有没有把同一种方案重复算了多次。
\end{itemize}

\smallskip
\textbf{本章读法提醒：} 第十一章研究的是一串数的形状，不是单个计数答案。
\end{explainblock}


\begin{sourceblock}{第 11 章续读}
{\small\sloppy
\noindent 都成立, 则由罗尔定理知 fn−1 (x) 有 n − 2 个负的零点. 所以由 (11.2) 知 0 是 fn (x)\par
\noindent ′\par
\noindent 零点, 它还有 n − 2 个与 fn−1 (x) 相同的负零点, 它们在 fn−1 (x) 相邻的每对零点之\par
\noindent ′\par
\noindent 间. 不妨设 fn−1 (x) 的 n − 2 个负的零点中最小的是 x0 , 则 fn−1 (x0 ) = 0, 又因为\par
\noindent lim fn−1 (x) = lim ex pn−1 (x) = 0 = fn−1 (x0 ),\par
\noindent x→−∞ x→−∞\par
\noindent ′\par
\noindent 所以由广义的罗尔定理知在区间 (−∞, x0 ) 内存在唯一一点 x1 使得 fn−1 (x1 ) = 0(其\par
\noindent 唯一性由 (11.2)及 fn (x) 只有 n 个根确定). 因此, fn (x) 有 n 个不同的实根, 也即\par
\noindent pn (x) 的根都是实根. 又因为对任意的 1 ≤ k ≤ n, 第二类 Stirling 数都不为为 0, 所\par
\noindent 以由推论 11.2.7 知 \{S(n, k)\}nk=1 是严格对数凹的. 因此是单峰的且最多有两个最大\par
\noindent 值.\par
\noindent 猜想 11.2.8 (Wegner’s Conjecture) 当 n ≥ 2 时, 第二类 Stirling 数的峰值唯一.\par
}
\end{sourceblock}


\begin{explainblock}
这一段是原书论述链的一部分。读的时候把它当作桥梁：它通常负责把上一个定义、例子或定理，引到下一个计数方法。

\smallskip
\textbf{初学者检查点：}
\begin{itemize}[leftmargin=2em,itemsep=0.25em]
\item 找出这一段承接的上一个概念。
\item 找出它准备引出的下一个概念。
\item 把中间的关键等式或构造单独抄出来。
\end{itemize}

\smallskip
\textbf{本章读法提醒：} 第十一章研究的是一串数的形状，不是单个计数答案。
\end{explainblock}


\chapter{Holonomic functions and Sequences}
\begin{roadmapblock}
本章保留原书正文的主要数学骨架，并在每个定义、定理、例题和论述片段后加入解读。
阅读时不要把目标放在“背下答案”上，而是放在“看懂这一步为什么这样数”上。

\smallskip
\textbf{本章最低目标：} 能说出每个公式左边在数什么、右边在数什么，以及为什么两边相等。
\end{roadmapblock}


\begin{sourceblock}{第 12 章正文}
{\small\sloppy
\noindent Holonomic functions and\par
\noindent Sequences\par
}
\end{sourceblock}


\begin{explainblock}
这一段是原书论述链的一部分。读的时候把它当作桥梁：它通常负责把上一个定义、例子或定理，引到下一个计数方法。

\smallskip
\textbf{初学者检查点：}
\begin{itemize}[leftmargin=2em,itemsep=0.25em]
\item 找出这一段承接的上一个概念。
\item 找出它准备引出的下一个概念。
\item 把中间的关键等式或构造单独抄出来。
\end{itemize}

\smallskip
\textbf{本章读法提醒：} 第十二章关键词是有限阶线性递推和有限阶线性微分方程。
\end{explainblock}


\begin{sourceblock}{§ 12.1 Holonomic 函数}
{\small\sloppy
\noindent § 12.1. Holonomic 函数\par
}
\end{sourceblock}


\begin{explainblock}
这一节先不要急着背公式。先问三个问题：对象是什么，哪些限制条件不能丢，最后到底是在数所有对象、还是在数某个统计量的分布。把这三件事说清楚，后面的公式才会有位置。

\smallskip
\textbf{初学者检查点：}
\begin{itemize}[leftmargin=2em,itemsep=0.25em]
\item 把本节出现的新对象写成一句话定义。
\item 把每个公式左边和右边分别翻译成“在数什么”。
\item 找一个最小非平凡例子，手工列出所有对象。
\end{itemize}

\smallskip
\textbf{本章读法提醒：} 第十二章关键词是有限阶线性递推和有限阶线性微分方程。
\end{explainblock}


\begin{sourceblock}{定义 12.1.1 设 \{cn \}n≥0 是一个序列, 若存在有限个多项式 ar (n), ar−1 (n), . . . , a0 (n) ∈}
{\small\sloppy
\noindent 定义 12.1.1 设 \{cn \}n≥0 是一个序列, 若存在有限个多项式 ar (n), ar−1 (n), . . . , a0 (n) ∈\par
\noindent C[n] 使得\par
\noindent ar (n)cn+r + ar−1 (n)cn+r−1 + . . . + a0 (n)cn = 0\par
\noindent 对所有的 n ∈ N 都成立, 则称序列 \{cn \}∞\par
\noindent n=0 是一个 holonomic 序列.\par
\noindent 下面举一些简单的例子.\par
}
\end{sourceblock}


\begin{explainblock}
定义段的任务是定边界。这里要特别注意参数范围、是否允许重复、是否考虑顺序、对象有没有标签。后面的每一道例题和证明，本质上都在反复使用这些边界。

\smallskip
\textbf{初学者检查点：}
\begin{itemize}[leftmargin=2em,itemsep=0.25em]
\item 圈出被定义的对象名称。
\item 圈出参数范围，例如 n、k 是否要求正整数。
\item 判断它是有序对象、无序对象、带标签对象还是结构对象。
\end{itemize}

\smallskip
\textbf{本章读法提醒：} 第十二章关键词是有限阶线性递推和有限阶线性微分方程。
\end{explainblock}


\begin{sourceblock}{例 12.1.1 Fibonacci 序列 \{Fn \}n≥0 满足递归关系}
{\small\sloppy
\noindent 例 12.1.1 Fibonacci 序列 \{Fn \}n≥0 满足递归关系\par
\noindent Fn+2 = Fn+1 + Fn ,\par
\noindent 且 F0 = 0, F1 = 1. 只需取 a2 (n) = 1, a1 (n) = a0 (n) = −1 立得 Fibonacci 序列是一\par
\noindent 个 holonomic 序列.\par
}
\end{sourceblock}


\begin{explainblock}
例题是学习这本书最重要的部分。不要只看答案，要把每个限制条件变成一个计数动作：先安排谁、再插入谁、有没有重复、需不需要除以对称性。

\smallskip
\textbf{初学者检查点：}
\begin{itemize}[leftmargin=2em,itemsep=0.25em]
\item 先写出没有限制时的总数。
\item 逐条加入限制条件，观察总数怎样变化。
\item 最后检查有没有把同一种方案重复算了多次。
\end{itemize}

\smallskip
\textbf{本章读法提醒：} 第十二章关键词是有限阶线性递推和有限阶线性微分方程。
\end{explainblock}


\begin{sourceblock}{例 12.1.2 序列 \{n!\}n≥0 满足}
{\small\sloppy
\noindent 例 12.1.2 序列 \{n!\}n≥0 满足\par
\noindent (n + 1)! = (n + 1) · n!,\par
\noindent 只需取 a1 (n) = 1, a0 (n) = −n − 1, 可得序列 \{n!\}∞\par
\noindent n=0 是一个 holonomic 序列.\par
}
\end{sourceblock}


\begin{explainblock}
例题是学习这本书最重要的部分。不要只看答案，要把每个限制条件变成一个计数动作：先安排谁、再插入谁、有没有重复、需不需要除以对称性。

\smallskip
\textbf{初学者检查点：}
\begin{itemize}[leftmargin=2em,itemsep=0.25em]
\item 先写出没有限制时的总数。
\item 逐条加入限制条件，观察总数怎样变化。
\item 最后检查有没有把同一种方案重复算了多次。
\end{itemize}

\smallskip
\textbf{本章读法提醒：} 第十二章关键词是有限阶线性递推和有限阶线性微分方程。
\end{explainblock}


\begin{sourceblock}{例 12.1.3 错排数 Dn 满足}
{\small\sloppy
\noindent 例 12.1.3 错排数 Dn 满足\par
\noindent Dn+2 = (n + 1)(Dn+1 + Dn ),\par
\noindent 其中 D0 = 1, D1 = 1. 所以取 a2 (n) = 1, a1 (n) = a0 (n) = −n−1, 可知序列 \{Dn \}n≥1\par
\noindent 是 holonomic 序列.\par
}
\end{sourceblock}


\begin{explainblock}
例题是学习这本书最重要的部分。不要只看答案，要把每个限制条件变成一个计数动作：先安排谁、再插入谁、有没有重复、需不需要除以对称性。

\smallskip
\textbf{初学者检查点：}
\begin{itemize}[leftmargin=2em,itemsep=0.25em]
\item 先写出没有限制时的总数。
\item 逐条加入限制条件，观察总数怎样变化。
\item 最后检查有没有把同一种方案重复算了多次。
\end{itemize}

\smallskip
\textbf{本章读法提醒：} 第十二章关键词是有限阶线性递推和有限阶线性微分方程。
\end{explainblock}


\begin{sourceblock}{例 12.1.4 二项式系数}
{\small\sloppy
\noindent 例 12.1.4 二项式系数\par
\noindent n n!\par
\noindent =\par
\noindent k k!(n − k)!\par
\noindent 226 第 12 章 Holonomic functions and Sequences\par
\noindent 满足\par
\noindent n+1 n\par
\noindent (n + 1 − k) = (n + 1) ,\par
\noindent k k\par
\noindent 所以只需取 a1 (n) = n + 1 − k, a0 (n) = −n − 1, 可知对于固定的 k ≥ 0, 序列\par
\noindent n\par
\noindent k n≥0 是关于 n 的一个 holonomic 序列. 另一方面, 二项式系数满足\par
\noindent n n\par
\noindent (k + 1) = (n − k) ,\par
\noindent k+1 k\par
\noindent n\par
\noindent 所以只需取 a1 (n) = k + 1, a0 (n) = k − n, 可知对于固定的 n ≥ 0, 序列 k k≥0\par
\noindent 是\par
}
\end{sourceblock}


\begin{explainblock}
例题是学习这本书最重要的部分。不要只看答案，要把每个限制条件变成一个计数动作：先安排谁、再插入谁、有没有重复、需不需要除以对称性。

\smallskip
\textbf{初学者检查点：}
\begin{itemize}[leftmargin=2em,itemsep=0.25em]
\item 先写出没有限制时的总数。
\item 逐条加入限制条件，观察总数怎样变化。
\item 最后检查有没有把同一种方案重复算了多次。
\end{itemize}

\smallskip
\textbf{本章读法提醒：} 第十二章关键词是有限阶线性递推和有限阶线性微分方程。
\end{explainblock}


\begin{sourceblock}{第 12 章续读}
{\small\sloppy
\noindent 关于 k 的一个 holonomic 序列.\par
}
\end{sourceblock}


\begin{explainblock}
这一段是原书论述链的一部分。读的时候把它当作桥梁：它通常负责把上一个定义、例子或定理，引到下一个计数方法。

\smallskip
\textbf{初学者检查点：}
\begin{itemize}[leftmargin=2em,itemsep=0.25em]
\item 找出这一段承接的上一个概念。
\item 找出它准备引出的下一个概念。
\item 把中间的关键等式或构造单独抄出来。
\end{itemize}

\smallskip
\textbf{本章读法提醒：} 第十二章关键词是有限阶线性递推和有限阶线性微分方程。
\end{explainblock}


\begin{sourceblock}{例 12.1.5 Catalan 数 Cn 为:}
{\small\sloppy
\noindent 例 12.1.5 Catalan 数 Cn 为:\par
\noindent Cn = .\par
\noindent n+1 n\par
\noindent 故\par
\noindent (n + 2)Cn+1 = 2(2n + 1)Cn .\par
\noindent 只需取 a1 (n) = (n + 2), a0 (n) = −2(2n + 1), 可知序列 \{Cn \}n≥0 是一个 holonomic\par
\noindent 序列.\par
}
\end{sourceblock}


\begin{explainblock}
例题是学习这本书最重要的部分。不要只看答案，要把每个限制条件变成一个计数动作：先安排谁、再插入谁、有没有重复、需不需要除以对称性。

\smallskip
\textbf{初学者检查点：}
\begin{itemize}[leftmargin=2em,itemsep=0.25em]
\item 先写出没有限制时的总数。
\item 逐条加入限制条件，观察总数怎样变化。
\item 最后检查有没有把同一种方案重复算了多次。
\end{itemize}

\smallskip
\textbf{本章读法提醒：} 第十二章关键词是有限阶线性递推和有限阶线性微分方程。
\end{explainblock}


\begin{sourceblock}{定义 12.1.2 设 f (x) ∈ K[[x]] , 如果存在有限个多项式 ar (x), ar−1 (x), . . . , a0 (x) ∈}
{\small\sloppy
\noindent 定义 12.1.2 设 f (x) ∈ K[[x]] , 如果存在有限个多项式 ar (x), ar−1 (x), . . . , a0 (x) ∈\par
\noindent K[x], 且 ar (x) 6= 0, 使得\par
\noindent ar (x)f (n+r) (x) + ar−1 (x)f (n+r−1) (x) + · · · + a1 (x)f ′ (x) + a0 (x)f (x) = 0\par
\noindent 成立, 则称 f (x) 是一个 D-finite 幂级数, 也称 f (x) 是一个 holonomic 函数.\par
\noindent P\par
}
\end{sourceblock}


\begin{explainblock}
定义段的任务是定边界。这里要特别注意参数范围、是否允许重复、是否考虑顺序、对象有没有标签。后面的每一道例题和证明，本质上都在反复使用这些边界。

\smallskip
\textbf{初学者检查点：}
\begin{itemize}[leftmargin=2em,itemsep=0.25em]
\item 圈出被定义的对象名称。
\item 圈出参数范围，例如 n、k 是否要求正整数。
\item 判断它是有序对象、无序对象、带标签对象还是结构对象。
\end{itemize}

\smallskip
\textbf{本章读法提醒：} 第十二章关键词是有限阶线性递推和有限阶线性微分方程。
\end{explainblock}


\begin{sourceblock}{定理 12.1.3 [18, Proposition 6.4.3] 设 f (x) = n≥0 cn xn ∈ K[[x]]. 则 f (x) 是}
{\small\sloppy
\noindent 定理 12.1.3 [18, Proposition 6.4.3] 设 f (x) = n≥0 cn xn ∈ K[[x]]. 则 f (x) 是\par
\noindent holonomic 函数当且仅当 \{cn \}n≥0 是 holonomic 序列.\par
\noindent c\par
\noindent 命 题 12.1.4 序 列 \{cn \}n≥0 是 一 个 holonomic 序 列 当 且 仅 当 n!\par
\noindent n\par
\noindent n≥0\par
\noindent 是一个\par
\noindent holonomic 序列.\par
\noindent 证 明: 先 证 明 必 要 性. 假 设 \{cn \}n≥0 是 一 个 holonomic 序 列, 则 存 在\par
\noindent ar (n), ar−1 (n), . . . , a0 (n) ∈ C[n] 使得\par
\noindent ar (n)cn+r + ar−1 (n)cn+r−1 + . . . + a0 (n)cn = 0\par
\noindent 对所有的 n ∈ N 都成立. 故\par
\noindent (n + r)! cn+r (n + r − 1)! cn+r−1 cn\par
\noindent · ar (n) + · ar−1 (n) + · · · + a0 (n) = 0.\par
\noindent n! (n + r)! n! (n + r − 1)! n!\par
\noindent 对任意 0 ≤ i ≤ r, 令 bi (n) = (n+i)!\par
\noindent n! · · · ai (n) 得\par
\noindent cn+r cn+r−1 cn\par
}
\end{sourceblock}


\begin{explainblock}
定理段不要只看结论，要看它把哪两个计数方式连在一起。组合学证明最常见的技巧是：左边按一种标准数，右边按另一种标准数，然后说明两边数的是同一批对象。

\smallskip
\textbf{初学者检查点：}
\begin{itemize}[leftmargin=2em,itemsep=0.25em]
\item 先复述定理的假设，不要把适用范围扩大。
\item 找出证明中的基本计数原则：乘法、加法、双射、递推、容斥或生成函数。
\item 检查最后的等式化简是否和前面的对象解释一致。
\end{itemize}

\smallskip
\textbf{本章读法提醒：} 第十二章关键词是有限阶线性递推和有限阶线性微分方程。
\end{explainblock}


\begin{sourceblock}{第 12 章续读}
{\small\sloppy
\noindent br (n) + br−1 (n) + · · · + b0 = 0,\par
\noindent (n + r)! (n + r − 1)! n!\par
\noindent 12.1 Holonomic 函数 227\par
\noindent cn ∞\par
\noindent 所以 n! n=0\par
\noindent 是一个 holonomic 序列.\par
\noindent cn\par
\noindent 再证明充分性. 假设 n! n≥0\par
\noindent 是一个 holonomic 序列, 则存在 ar (n), ar−1 (n), . . . , a0 (n) ∈\par
\noindent C[n] 使得\par
\noindent cn+r cn+r−1 cn\par
\noindent ar (n) + ar−1 (n) + · · · + a0 (n) = 0\par
\noindent (n + r)! (n + r − 1)! n!\par
\noindent 对所有的 n ∈ N 都成立. 等式两边同时乘以 (n + r)! 得到\par
\noindent ar (n)cn+r + (n + r)ar−1 (n)cn+r−1 + · · · + (n + r)(n + r − 1) · · · (n + 1) · · · a0 (n)cn = 0.\par
\noindent 对任意 0 ≤ i ≤ r, 令 bi (n) = (n + r)(n + r − 1) · · · (n + i + 1) · ai (n), 可知 \{cn \}n≥0\par
\noindent 是一个 holonomic 序列.\par
}
\end{sourceblock}


\begin{explainblock}
这一段是原书论述链的一部分。读的时候把它当作桥梁：它通常负责把上一个定义、例子或定理，引到下一个计数方法。

\smallskip
\textbf{初学者检查点：}
\begin{itemize}[leftmargin=2em,itemsep=0.25em]
\item 找出这一段承接的上一个概念。
\item 找出它准备引出的下一个概念。
\item 把中间的关键等式或构造单独抄出来。
\end{itemize}

\smallskip
\textbf{本章读法提醒：} 第十二章关键词是有限阶线性递推和有限阶线性微分方程。
\end{explainblock}


\begin{sourceblock}{推论 12.1.5 序列 \{cn \}n≥0 的生成函数是一个 holonomic 函数当且仅当它的指数}
{\small\sloppy
\noindent 推论 12.1.5 序列 \{cn \}n≥0 的生成函数是一个 holonomic 函数当且仅当它的指数\par
\noindent 生成函数是一个 holonomic 函数.\par
}
\end{sourceblock}


\begin{explainblock}
定理段不要只看结论，要看它把哪两个计数方式连在一起。组合学证明最常见的技巧是：左边按一种标准数，右边按另一种标准数，然后说明两边数的是同一批对象。

\smallskip
\textbf{初学者检查点：}
\begin{itemize}[leftmargin=2em,itemsep=0.25em]
\item 先复述定理的假设，不要把适用范围扩大。
\item 找出证明中的基本计数原则：乘法、加法、双射、递推、容斥或生成函数。
\item 检查最后的等式化简是否和前面的对象解释一致。
\end{itemize}

\smallskip
\textbf{本章读法提醒：} 第十二章关键词是有限阶线性递推和有限阶线性微分方程。
\end{explainblock}


\begin{sourceblock}{例 12.1.6 Fibonacci 序列的生成函数}
{\small\sloppy
\noindent 例 12.1.6 Fibonacci 序列的生成函数\par
\noindent x\par
\noindent s(x) =\par
\noindent 是 holonomic 函数.\par
}
\end{sourceblock}


\begin{explainblock}
例题是学习这本书最重要的部分。不要只看答案，要把每个限制条件变成一个计数动作：先安排谁、再插入谁、有没有重复、需不需要除以对称性。

\smallskip
\textbf{初学者检查点：}
\begin{itemize}[leftmargin=2em,itemsep=0.25em]
\item 先写出没有限制时的总数。
\item 逐条加入限制条件，观察总数怎样变化。
\item 最后检查有没有把同一种方案重复算了多次。
\end{itemize}

\smallskip
\textbf{本章读法提醒：} 第十二章关键词是有限阶线性递推和有限阶线性微分方程。
\end{explainblock}


\begin{sourceblock}{例 12.1.7 错排数的生成函数}
{\small\sloppy
\noindent 例 12.1.7 错排数的生成函数\par
\noindent X xn e−x\par
\noindent Dn = ,\par
\noindent n! 1−x\par
\noindent n≥0\par
\noindent 是 holonomic 函数.\par
}
\end{sourceblock}


\begin{explainblock}
例题是学习这本书最重要的部分。不要只看答案，要把每个限制条件变成一个计数动作：先安排谁、再插入谁、有没有重复、需不需要除以对称性。

\smallskip
\textbf{初学者检查点：}
\begin{itemize}[leftmargin=2em,itemsep=0.25em]
\item 先写出没有限制时的总数。
\item 逐条加入限制条件，观察总数怎样变化。
\item 最后检查有没有把同一种方案重复算了多次。
\end{itemize}

\smallskip
\textbf{本章读法提醒：} 第十二章关键词是有限阶线性递推和有限阶线性微分方程。
\end{explainblock}


\begin{sourceblock}{例 12.1.8 序列 \{n!\}∞}
{\small\sloppy
\noindent 例 12.1.8 序列 \{n!\}∞\par
\noindent n=0 的指数生成函数\par
\noindent f (x) =\par
\noindent 1−x\par
\noindent 是 holonomic 函数.\par
}
\end{sourceblock}


\begin{explainblock}
例题是学习这本书最重要的部分。不要只看答案，要把每个限制条件变成一个计数动作：先安排谁、再插入谁、有没有重复、需不需要除以对称性。

\smallskip
\textbf{初学者检查点：}
\begin{itemize}[leftmargin=2em,itemsep=0.25em]
\item 先写出没有限制时的总数。
\item 逐条加入限制条件，观察总数怎样变化。
\item 最后检查有没有把同一种方案重复算了多次。
\end{itemize}

\smallskip
\textbf{本章读法提醒：} 第十二章关键词是有限阶线性递推和有限阶线性微分方程。
\end{explainblock}


\begin{sourceblock}{例 12.1.9 任意给定正整数 n, 由二项式定理有}
{\small\sloppy
\noindent 例 12.1.9 任意给定正整数 n, 由二项式定理有\par
\noindent ∞\par
\noindent X n k\par
\noindent x = (1 + x)n\par
\noindent k\par
\noindent k=0\par
\noindent n ∞\par
\noindent 是一个 holonomic 函数. 另一方面, 任意固定正整数 k, 序列 k n=0\par
\noindent 的指数生成\par
\noindent 函数\par
\noindent ∞ n\par
\noindent X ∞\par
\noindent X ∞ ∞\par
\noindent n x n! xn 1 X xn xk X xn xk ex\par
\noindent = = = =\par
\noindent k n! k!(n − k)! n! k! (n − k)! k! n! k!\par
\noindent n=0 n=k n=k n=0\par
\noindent 是 holonomic 函数.\par
}
\end{sourceblock}


\begin{explainblock}
例题是学习这本书最重要的部分。不要只看答案，要把每个限制条件变成一个计数动作：先安排谁、再插入谁、有没有重复、需不需要除以对称性。

\smallskip
\textbf{初学者检查点：}
\begin{itemize}[leftmargin=2em,itemsep=0.25em]
\item 先写出没有限制时的总数。
\item 逐条加入限制条件，观察总数怎样变化。
\item 最后检查有没有把同一种方案重复算了多次。
\end{itemize}

\smallskip
\textbf{本章读法提醒：} 第十二章关键词是有限阶线性递推和有限阶线性微分方程。
\end{explainblock}


\begin{sourceblock}{第 12 章续读}
{\small\sloppy
\noindent 228 第 12 章 Holonomic functions and Sequences\par
}
\end{sourceblock}


\begin{explainblock}
这一段是原书论述链的一部分。读的时候把它当作桥梁：它通常负责把上一个定义、例子或定理，引到下一个计数方法。

\smallskip
\textbf{初学者检查点：}
\begin{itemize}[leftmargin=2em,itemsep=0.25em]
\item 找出这一段承接的上一个概念。
\item 找出它准备引出的下一个概念。
\item 把中间的关键等式或构造单独抄出来。
\end{itemize}

\smallskip
\textbf{本章读法提醒：} 第十二章关键词是有限阶线性递推和有限阶线性微分方程。
\end{explainblock}


\begin{sourceblock}{例 12.1.10 Catalan 数的生成函数为}
{\small\sloppy
\noindent 例 12.1.10 Catalan 数的生成函数为\par
\noindent c(x) = √\par
\noindent 是 holonomic 函数.\par
}
\end{sourceblock}


\begin{explainblock}
例题是学习这本书最重要的部分。不要只看答案，要把每个限制条件变成一个计数动作：先安排谁、再插入谁、有没有重复、需不需要除以对称性。

\smallskip
\textbf{初学者检查点：}
\begin{itemize}[leftmargin=2em,itemsep=0.25em]
\item 先写出没有限制时的总数。
\item 逐条加入限制条件，观察总数怎样变化。
\item 最后检查有没有把同一种方案重复算了多次。
\end{itemize}

\smallskip
\textbf{本章读法提醒：} 第十二章关键词是有限阶线性递推和有限阶线性微分方程。
\end{explainblock}


\begin{sourceblock}{定理 12.1.6 [11, Theorem 1, 奇异点有限性定理] Holonomic 函数仅有有限多个}
{\small\sloppy
\noindent 定理 12.1.6 [11, Theorem 1, 奇异点有限性定理] Holonomic 函数仅有有限多个\par
\noindent 奇异点.\par
\noindent 这一定理可以帮助我们找到一些不是 holomonic 函数的例子.\par
}
\end{sourceblock}


\begin{explainblock}
定理段不要只看结论，要看它把哪两个计数方式连在一起。组合学证明最常见的技巧是：左边按一种标准数，右边按另一种标准数，然后说明两边数的是同一批对象。

\smallskip
\textbf{初学者检查点：}
\begin{itemize}[leftmargin=2em,itemsep=0.25em]
\item 先复述定理的假设，不要把适用范围扩大。
\item 找出证明中的基本计数原则：乘法、加法、双射、递推、容斥或生成函数。
\item 检查最后的等式化简是否和前面的对象解释一致。
\end{itemize}

\smallskip
\textbf{本章读法提醒：} 第十二章关键词是有限阶线性递推和有限阶线性微分方程。
\end{explainblock}


\begin{sourceblock}{例 12.1.11 整数分拆函数 p(n) 的生成函数为}
{\small\sloppy
\noindent 例 12.1.11 整数分拆函数 p(n) 的生成函数为\par
\noindent X ∞\par
\noindent Y 1\par
\noindent p(n)xn = ,\par
\noindent n≥0 n=1\par
\noindent 它 有 无 穷 多 个 奇 异 点, 所 以 它 不 是 holonomic 函 数. 故 序 列 \{p(n)\}∞\par
\noindent n=0 也 不 是\par
\noindent holonomic 序列.\par
}
\end{sourceblock}


\begin{explainblock}
例题是学习这本书最重要的部分。不要只看答案，要把每个限制条件变成一个计数动作：先安排谁、再插入谁、有没有重复、需不需要除以对称性。

\smallskip
\textbf{初学者检查点：}
\begin{itemize}[leftmargin=2em,itemsep=0.25em]
\item 先写出没有限制时的总数。
\item 逐条加入限制条件，观察总数怎样变化。
\item 最后检查有没有把同一种方案重复算了多次。
\end{itemize}

\smallskip
\textbf{本章读法提醒：} 第十二章关键词是有限阶线性递推和有限阶线性微分方程。
\end{explainblock}


\begin{sourceblock}{例 12.1.12 Bernoulli 数 Bn 的生成函数为}
{\small\sloppy
\noindent 例 12.1.12 Bernoulli 数 Bn 的生成函数为\par
\noindent X xn x\par
\noindent Bn = x ,\par
\noindent n! e −1\par
\noindent n≥0\par
\noindent 它有无穷多个奇异点, 所以它不是 holonomic 函数, \{B(n)\}n≥0 也不是 holonomic 序\par
\noindent 列.\par
}
\end{sourceblock}


\begin{explainblock}
例题是学习这本书最重要的部分。不要只看答案，要把每个限制条件变成一个计数动作：先安排谁、再插入谁、有没有重复、需不需要除以对称性。

\smallskip
\textbf{初学者检查点：}
\begin{itemize}[leftmargin=2em,itemsep=0.25em]
\item 先写出没有限制时的总数。
\item 逐条加入限制条件，观察总数怎样变化。
\item 最后检查有没有把同一种方案重复算了多次。
\end{itemize}

\smallskip
\textbf{本章读法提醒：} 第十二章关键词是有限阶线性递推和有限阶线性微分方程。
\end{explainblock}


\begin{sourceblock}{例 12.1.13 Euler 数 En 的生成函数为}
{\small\sloppy
\noindent 例 12.1.13 Euler 数 En 的生成函数为\par
\noindent X xn\par
\noindent En = tan(x) + sec(x),\par
\noindent n!\par
\noindent n≥0\par
\noindent 它有无穷多个奇异点, 所以它不是 holonomic 函数, 序列 \{A(n)\}n≥0 也不是 holo-\par
\noindent nomic 序列.\par
}
\end{sourceblock}


\begin{explainblock}
例题是学习这本书最重要的部分。不要只看答案，要把每个限制条件变成一个计数动作：先安排谁、再插入谁、有没有重复、需不需要除以对称性。

\smallskip
\textbf{初学者检查点：}
\begin{itemize}[leftmargin=2em,itemsep=0.25em]
\item 先写出没有限制时的总数。
\item 逐条加入限制条件，观察总数怎样变化。
\item 最后检查有没有把同一种方案重复算了多次。
\end{itemize}

\smallskip
\textbf{本章读法提醒：} 第十二章关键词是有限阶线性递推和有限阶线性微分方程。
\end{explainblock}


\begin{sourceblock}{定义 12.1.7 设 K 为一个域. 形式幂级数 f ∈ K[[x]] 被称为代数的, 如果存在不}
{\small\sloppy
\noindent 定义 12.1.7 设 K 为一个域. 形式幂级数 f ∈ K[[x]] 被称为代数的, 如果存在不\par
\noindent 全为零的多项式 p0 (x), . . . , pn (x) ∈ K[x], 对任意的 x 使得\par
\noindent pn (x)f (x)n + · · · + p1 (x)f (x) + p0 (x) = 0. (12.1)\par
}
\end{sourceblock}


\begin{explainblock}
定义段的任务是定边界。这里要特别注意参数范围、是否允许重复、是否考虑顺序、对象有没有标签。后面的每一道例题和证明，本质上都在反复使用这些边界。

\smallskip
\textbf{初学者检查点：}
\begin{itemize}[leftmargin=2em,itemsep=0.25em]
\item 圈出被定义的对象名称。
\item 圈出参数范围，例如 n、k 是否要求正整数。
\item 判断它是有序对象、无序对象、带标签对象还是结构对象。
\end{itemize}

\smallskip
\textbf{本章读法提醒：} 第十二章关键词是有限阶线性递推和有限阶线性微分方程。
\end{explainblock}


\begin{sourceblock}{例 12.1.14 指数函数 f (x) = ex 就不是代数的. 我们可以用反证法证明. 假设}
{\small\sloppy
\noindent 例 12.1.14 指数函数 f (x) = ex 就不是代数的. 我们可以用反证法证明. 假设\par
\noindent f (x) = ex 是代数, 则存在有限多个多项式 pi (x) (0 ≤ i ≤ n) 使得\par
\noindent pn (x)enx + pn−1 (x)e(n−1)x + · · · + p1 (x)ex + p0 (x) = 0, pn (x) 6= 0\par
\noindent 成立, 然后等式两边同时除以 pn (x)enx 得到\par
\noindent pn−1 (x) −x p1 (x) −(n−1)x p0 (x) −nx\par
\noindent 1+ e + ··· + e + e = 0.\par
\noindent pn (x) pn (x) pn (x)\par
\noindent 12.1 Holonomic 函数 229\par
\noindent 最后我们取极限 x → ∞ 后得出矛盾 1 = 0.\par
\noindent Rx\par
}
\end{sourceblock}


\begin{explainblock}
例题是学习这本书最重要的部分。不要只看答案，要把每个限制条件变成一个计数动作：先安排谁、再插入谁、有没有重复、需不需要除以对称性。

\smallskip
\textbf{初学者检查点：}
\begin{itemize}[leftmargin=2em,itemsep=0.25em]
\item 先写出没有限制时的总数。
\item 逐条加入限制条件，观察总数怎样变化。
\item 最后检查有没有把同一种方案重复算了多次。
\end{itemize}

\smallskip
\textbf{本章读法提醒：} 第十二章关键词是有限阶线性递推和有限阶线性微分方程。
\end{explainblock}


\begin{sourceblock}{引理 12.1.8 ( [17]) 给定一 holonomic 函数 f (x), 则 exp 0 f (z)dz 是 holonomic}
{\small\sloppy
\noindent 引理 12.1.8 ( [17]) 给定一 holonomic 函数 f (x), 则 exp 0 f (z)dz 是 holonomic\par
\noindent 函数当且仅当 f (x) 是代数的.\par
\noindent 显然 f (x) = ex 是 holonomic 函数. 由于 f (x) = ex 不是代数的, 则根据引理\par
\noindent 12.1.8, 我们得到 Z x\par
\noindent exp f (z)dz = exp (ex − 1)\par
\noindent 不是 holonomic 函数. 函数 g(x) = exp(ex − 1) 是 Bell 数 Bn 的指数生成函数, 其中\par
\noindent Bell 数为集合 \{1, . . . , n\} 的所有划分的个数. 由定理 12.1.3 和推论 12.1.5 我们得出\par
\noindent Bell 序列 \{Bn \}∞\par
\noindent n=0 不是 holonomic 的.\par
\noindent 230 第 12 章 Holonomic functions and Sequences\par
}
\end{sourceblock}


\begin{explainblock}
定理段不要只看结论，要看它把哪两个计数方式连在一起。组合学证明最常见的技巧是：左边按一种标准数，右边按另一种标准数，然后说明两边数的是同一批对象。

\smallskip
\textbf{初学者检查点：}
\begin{itemize}[leftmargin=2em,itemsep=0.25em]
\item 先复述定理的假设，不要把适用范围扩大。
\item 找出证明中的基本计数原则：乘法、加法、双射、递推、容斥或生成函数。
\item 检查最后的等式化简是否和前面的对象解释一致。
\end{itemize}

\smallskip
\textbf{本章读法提醒：} 第十二章关键词是有限阶线性递推和有限阶线性微分方程。
\end{explainblock}

\backmatter
\chapter*{参考阅读}
\begin{itemize}[leftmargin=2em]
\item 陈永川，《欧拉与计数组合学》，天津大学应用数学中心，2020 年秋季课程讲义。
\item M. Aigner, \emph{A Course in Enumeration}, Springer, 2007.
\item R. P. Stanley, \emph{Enumerative Combinatorics}, Vol. 1--2.
\item R. L. Graham, D. E. Knuth, O. Patashnik, \emph{Concrete Mathematics}.
\end{itemize}
\end{document}
